% REVIEW COPY: proof progress, not a completed proof of Theorem 1.
% STAGE 0 RE-AUDIT (2026-09-17): the original initialized noisy many-neuron
% reversal goal is reopened by the user. Scope A is not accepted as necessary.
% See sbsep:theorem: a noisy many-neuron learned-strong-cohort endpoint theorem
% is proved in an extreme-separation region. Two learned families and a fixed
% concept-ratio regime remain open; no headline scope pivot has been accepted.

\documentclass{article} % For LaTeX2e
\usepackage{iclr2027_conference,times}

% Optional math commands from https://github.com/goodfeli/dlbook_notation.
\input{math_commands.tex}

\usepackage{hyperref}
\usepackage{url}

\usepackage{amsmath}
\usepackage{amssymb}
\usepackage{mathtools}
\usepackage{amsthm}
\usepackage{enumitem}

\theoremstyle{plain}
\newtheorem{theorem}{Theorem}[section]
\newtheorem{proposition}[theorem]{Proposition}
\newtheorem{lemma}[theorem]{Lemma}
\newtheorem{corollary}[theorem]{Corollary}
\newtheorem{conjecture}[theorem]{Conjecture}

\theoremstyle{definition}
\newtheorem{definition}[theorem]{Definition}
\newtheorem{assumption}[theorem]{Assumption}

\theoremstyle{remark}
\newtheorem{remark}[theorem]{Remark}


\title{Emergent Concept Specialization in Two-Layer ReLU Networks:\\Gate Creation, Predictive OOD Normal Forms, and Restricted Swing-By Reversal}

% Authors must not appear in the submitted version. They should be hidden
% as long as the \iclrfinalcopy macro remains commented out below.
% Non-anonymous submissions will be rejected without review.

\author{Antiquus S.~Hippocampus, Natalia Cerebro \& Amelie P. Amygdale \thanks{ Use footnote for providing further information
about author (webpage, alternative address)---\emph{not} for acknowledging
funding agencies.  Funding acknowledgements go at the end of the paper.} \\
Department of Computer Science\\
Cranberry-Lemon University\\
Pittsburgh, PA 15213, USA \\
\texttt{\{hippo,brain,jen\}@cs.cranberry-lemon.edu} \\
\And
Ji Q. Ren \& Yevgeny LeNet \\
Department of Computational Neuroscience \\
University of the Witwatersrand \\
Joburg, South Africa \\
\texttt{\{robot,net\}@wits.ac.za} \\
\AND
Coauthor \\
Affiliation \\
Address \\
\texttt{email}
}

% The \author macro works with any number of authors. There are two commands
% used to separate the names and addresses of multiple authors: \And and \AND.
%
% Using \And between authors leaves it to \LaTeX{} to determine where to break
% the lines. Using \AND forces a linebreak at that point. So, if \LaTeX{}
% puts 3 of 4 authors names on the first line, and the last on the second
% line, try using \AND instead of \And before the third author name.

\newcommand{\fix}{\marginpar{FIX}}
\newcommand{\new}{\marginpar{NEW}}

%\iclrfinalcopy % Uncomment for camera-ready version, but NOT for submission.
\begin{document}


\maketitle

\begin{abstract}
We study gradient-flow training of two-layer ReLU networks on a two-concept
structured identity-mapping problem.  The main noisy many-neuron target is a
high-probability concept-specialization theorem: a strong concept locks first,
its learned residual creates a weak-concept attraction basin through a
finite-sample gate cliff, and a positive weak cohort is subsequently captured
and amplified.  Conditional on specialization, we derive a predictive
near-boundary OOD normal form with explicit coefficient and derivative
remainders.  A complete swing-by reversal sign theorem is proved in a
restricted noiseless fixed-gate subcase.  A general noisy many-neuron reversal
sign theorem in the original two-family regime is not yet proved here.
A restricted noisy many-neuron theorem now certifies, with probability at
least $0.97$, an initial OOD transient, a dominant cohort with learned mass
at least $0.55$, and later strict degradation on a deterministic sector of
width $3.65\cdot10^{-4}$ radians. An explicit regime uses $h\ge10^7$,
$h\varepsilon^2=0.04$, concept ratio $10^{-4}$, and noise $10^{-9}$.
On the certified time interval the error minimum precedes the declared
specialization time; specialization is not proved to cause the reversal.
The weak concept remains unlearned. R0 and the earlier coherent theorem
are retained, and the original two-learned-family goal remains open. This file is a proof-progress
review copy and marks the remaining trajectory-level obligations explicitly.
\end{abstract}

\section{Submission of conference papers to ICLR 2027}

ICLR requires electronic submissions, processed by
\url{https://openreview.net/}. See ICLR's website for more instructions.

If your paper is ultimately accepted, the statement {\tt
  {\textbackslash}iclrfinalcopy} should be inserted to adjust the
format to the camera ready requirements.

The format for the submissions is a variant of the NeurIPS format.
Please read carefully the instructions below, and follow them
faithfully.

\subsection{Style}

Papers to be submitted to ICLR 2027 must be prepared according to the
instructions presented here.

%% Please note that we have introduced automatic line number generation
%% into the style file for \LaTeXe. This is to help reviewers
%% refer to specific lines of the paper when they make their comments. Please do
%% NOT refer to these line numbers in your paper as they will be removed from the
%% style file for the final version of accepted papers.

Authors are required to use the ICLR \LaTeX{} style files obtainable at the
ICLR website. Please make sure you use the current files and
not previous versions. Tweaking the style files may be grounds for rejection.

\subsection{Retrieval of style files}

The style files for ICLR and other conference information are available online at:
\begin{center}
   \url{http://www.iclr.cc/}
\end{center}
The file \verb+iclr2027_conference.pdf+ contains these
instructions and illustrates the
various formatting requirements your ICLR paper must satisfy.
Submissions must be made using \LaTeX{} and the style files
\verb+iclr2027_conference.sty+ and \verb+iclr2027_conference.bst+ (to be used with \LaTeX{}2e). The file
\verb+iclr2027_conference.tex+ may be used as a ``shell'' for writing your paper. All you
have to do is replace the author, title, abstract, and text of the paper with
your own.

The formatting instructions contained in these style files are summarized in
sections \ref{gen_inst}, \ref{headings}, and \ref{others} below.

\section{General formatting instructions}
\label{gen_inst}

The text must be confined within a rectangle 5.5~inches (33~picas) wide and
9~inches (54~picas) long. The left margin is 1.5~inch (9~picas).
Use 10~point type with a vertical spacing of 11~points. Times New Roman is the
preferred typeface throughout. Paragraphs are separated by 1/2~line space,
with no indentation.

Paper title is 17~point, in small caps and left-aligned.
All pages should start at 1~inch (6~picas) from the top of the page.

Authors' names are
set in boldface, and each name is placed above its corresponding
address. The lead author's name is to be listed first, and
the co-authors' names are set to follow. Authors sharing the
same address can be on the same line.

Please pay special attention to the instructions in section \ref{others}
regarding figures, tables, acknowledgments, and references.


There will be a strict upper limit of \textbf{9 pages} for the main text of the initial submission, with unlimited additional pages for citations. This limit will be expanded to \textbf{10 pages} for rebuttal/camera ready.

\section{Headings: first level}
\label{headings}

First level headings are in small caps,
flush left and in point size 12. One line space before the first level
heading and 1/2~line space after the first level heading.

\subsection{Headings: second level}

Second level headings are in small caps,
flush left and in point size 10. One line space before the second level
heading and 1/2~line space after the second level heading.

\subsubsection{Headings: third level}

Third level headings are in small caps,
flush left and in point size 10. One line space before the third level
heading and 1/2~line space after the third level heading.

\section{Citations, figures, tables, references}
\label{others}

These instructions apply to everyone, regardless of the formatter being used.

\subsection{Citations within the text}

Citations within the text should be based on the \texttt{natbib} package
and include the authors' last names and year (with the ``et~al.'' construct
for more than two authors). When the authors or the publication are
included in the sentence, the citation should not be in parenthesis using \verb|\citet{}| (as
in ``See \citet{Hinton06} for more information.''). Otherwise, the citation
should be in parenthesis using \verb|\citep{}| (as in ``Deep learning shows promise to make progress
towards AI~\citep{Bengio+chapter2007}.'').

The corresponding references are to be listed in alphabetical order of
authors, in the \textsc{References} section. As to the format of the
references themselves, any style is acceptable as long as it is used
consistently.

\subsection{Footnotes}

Indicate footnotes with a number\footnote{Sample of the first footnote} in the
text. Place the footnotes at the bottom of the page on which they appear.
Precede the footnote with a horizontal rule of 2~inches
(12~picas).\footnote{Sample of the second footnote}

\subsection{Figures}

All artwork must be neat, clean, and legible. Lines should be dark
enough for purposes of reproduction; art work should not be
hand-drawn. The figure number and caption always appear after the
figure. Place one line space before the figure caption, and one line
space after the figure. The figure caption is lower case (except for
first word and proper nouns); figures are numbered consecutively.

Make sure the figure caption does not get separated from the figure.
Leave sufficient space to avoid splitting the figure and figure caption.

You may use color figures.
However, it is best for the
figure captions and the paper body to make sense if the paper is printed
either in black/white or in color.
\begin{figure}[h]
\begin{center}
%\framebox[4.0in]{$\;$}
\fbox{\rule[-.5cm]{0cm}{4cm} \rule[-.5cm]{4cm}{0cm}}
\end{center}
\caption{Sample figure caption.}
\end{figure}

\subsection{Tables}

All tables must be centered, neat, clean and legible. Do not use hand-drawn
tables. The table number and title always appear before the table. See
Table~\ref{sample-table}.

Place one line space before the table title, one line space after the table
title, and one line space after the table. The table title must be lower case
(except for first word and proper nouns); tables are numbered consecutively.

\begin{table}[t]
\caption{Sample table title}
\label{sample-table}
\begin{center}
\begin{tabular}{ll}
\multicolumn{1}{c}{\bf PART}  &\multicolumn{1}{c}{\bf DESCRIPTION}
\\ \hline \\
Dendrite         &Input terminal \\
Axon             &Output terminal \\
Soma             &Cell body (contains cell nucleus) \\
\end{tabular}
\end{center}
\end{table}

\section{Default Notation}

In an attempt to encourage standardized notation, we have included the
notation file from the textbook, \textit{Deep Learning}
\cite{goodfellow2016deep} available at
\url{https://github.com/goodfeli/dlbook_notation/}.  Use of this style
is not required and can be disabled by commenting out
\texttt{math\_commands.tex}.


\centerline{\bf Numbers and Arrays}
\bgroup
\def\arraystretch{1.5}
\begin{tabular}{p{1in}p{3.25in}}
$\displaystyle a$ & A scalar (integer or real)\\
$\displaystyle \va$ & A vector\\
$\displaystyle \mA$ & A matrix\\
$\displaystyle \tA$ & A tensor\\
$\displaystyle \mI_n$ & Identity matrix with $n$ rows and $n$ columns\\
$\displaystyle \mI$ & Identity matrix with dimensionality implied by context\\
$\displaystyle \ve^{(i)}$ & Standard basis vector $[0,\dots,0,1,0,\dots,0]$ with a 1 at position $i$\\
$\displaystyle \text{diag}(\va)$ & A square, diagonal matrix with diagonal entries given by $\va$\\
$\displaystyle \ra$ & A scalar random variable\\
$\displaystyle \rva$ & A vector-valued random variable\\
$\displaystyle \rmA$ & A matrix-valued random variable\\
\end{tabular}
\egroup
\vspace{0.25cm}

\centerline{\bf Sets and Graphs}
\bgroup
\def\arraystretch{1.5}

\begin{tabular}{p{1.25in}p{3.25in}}
$\displaystyle \sA$ & A set\\
$\displaystyle \R$ & The set of real numbers \\
$\displaystyle \{0, 1\}$ & The set containing 0 and 1 \\
$\displaystyle \{0, 1, \dots, n \}$ & The set of all integers between $0$ and $n$\\
$\displaystyle [a, b]$ & The real interval including $a$ and $b$\\
$\displaystyle (a, b]$ & The real interval excluding $a$ but including $b$\\
$\displaystyle \sA \backslash \sB$ & Set subtraction, i.e., the set containing the elements of $\sA$ that are not in $\sB$\\
$\displaystyle \gG$ & A graph\\
$\displaystyle \parents_\gG(\ervx_i)$ & The parents of $\ervx_i$ in $\gG$
\end{tabular}
\vspace{0.25cm}


\centerline{\bf Indexing}
\bgroup
\def\arraystretch{1.5}

\begin{tabular}{p{1.25in}p{3.25in}}
$\displaystyle \eva_i$ & Element $i$ of vector $\va$, with indexing starting at 1 \\
$\displaystyle \eva_{-i}$ & All elements of vector $\va$ except for element $i$ \\
$\displaystyle \emA_{i,j}$ & Element $i, j$ of matrix $\mA$ \\
$\displaystyle \mA_{i, :}$ & Row $i$ of matrix $\mA$ \\
$\displaystyle \mA_{:, i}$ & Column $i$ of matrix $\mA$ \\
$\displaystyle \etA_{i, j, k}$ & Element $(i, j, k)$ of a 3-D tensor $\tA$\\
$\displaystyle \tA_{:, :, i}$ & 2-D slice of a 3-D tensor\\
$\displaystyle \erva_i$ & Element $i$ of the random vector $\rva$ \\
\end{tabular}
\egroup
\vspace{0.25cm}


\centerline{\bf Calculus}
\bgroup
\def\arraystretch{1.5}
\begin{tabular}{p{1.25in}p{3.25in}}
% NOTE: the [2ex] on the next line adds extra height to that row of the table.
% Without that command, the fraction on the first line is too tall and collides
% with the fraction on the second line.
$\displaystyle\frac{d y} {d x}$ & Derivative of $y$ with respect to $x$\\ [2ex]
$\displaystyle \frac{\partial y} {\partial x} $ & Partial derivative of $y$ with respect to $x$ \\
$\displaystyle \nabla_\vx y $ & Gradient of $y$ with respect to $\vx$ \\
$\displaystyle \nabla_\mX y $ & Matrix derivatives of $y$ with respect to $\mX$ \\
$\displaystyle \nabla_\tX y $ & Tensor containing derivatives of $y$ with respect to $\tX$ \\
$\displaystyle \frac{\partial f}{\partial \vx} $ & Jacobian matrix $\mJ \in \R^{m\times n}$ of $f: \R^n \rightarrow \R^m$\\
$\displaystyle \nabla_\vx^2 f(\vx)\text{ or }\mH( f)(\vx)$ & The Hessian matrix of $f$ at input point $\vx$\\
$\displaystyle \int f(\vx) d\vx $ & Definite integral over the entire domain of $\vx$ \\
$\displaystyle \int_\sS f(\vx) d\vx$ & Definite integral with respect to $\vx$ over the set $\sS$ \\
\end{tabular}
\egroup
\vspace{0.25cm}

\centerline{\bf Probability and Information Theory}
\bgroup
\def\arraystretch{1.5}
\begin{tabular}{p{1.25in}p{3.25in}}
$\displaystyle P(\ra)$ & A probability distribution over a discrete variable\\
$\displaystyle p(\ra)$ & A probability distribution over a continuous variable, or over
a variable whose type has not been specified\\
$\displaystyle \ra \sim P$ & Random variable $\ra$ has distribution $P$\\% so thing on left of \sim should always be a random variable, with name beginning with \r
$\displaystyle  \E_{\rx\sim P} [ f(x) ]\text{ or } \E f(x)$ & Expectation of $f(x)$ with respect to $P(\rx)$ \\
$\displaystyle \Var(f(x)) $ &  Variance of $f(x)$ under $P(\rx)$ \\
$\displaystyle \Cov(f(x),g(x)) $ & Covariance of $f(x)$ and $g(x)$ under $P(\rx)$\\
$\displaystyle H(\rx) $ & Shannon entropy of the random variable $\rx$\\
$\displaystyle \KL ( P \Vert Q ) $ & Kullback-Leibler divergence of P and Q \\
$\displaystyle \mathcal{N} ( \vx ; \vmu , \mSigma)$ & Gaussian distribution %
over $\vx$ with mean $\vmu$ and covariance $\mSigma$ \\
\end{tabular}
\egroup
\vspace{0.25cm}

\centerline{\bf Functions}
\bgroup
\def\arraystretch{1.5}
\begin{tabular}{p{1.25in}p{3.25in}}
$\displaystyle f: \sA \rightarrow \sB$ & The function $f$ with domain $\sA$ and range $\sB$\\
$\displaystyle f \circ g $ & Composition of the functions $f$ and $g$ \\
  $\displaystyle f(\vx ; \vtheta) $ & A function of $\vx$ parametrized by $\vtheta$.
  (Sometimes we write $f(\vx)$ and omit the argument $\vtheta$ to lighten notation) \\
$\displaystyle \log x$ & Natural logarithm of $x$ \\
$\displaystyle \sigma(x)$ & Logistic sigmoid, $\displaystyle \frac{1} {1 + \exp(-x)}$ \\
$\displaystyle \zeta(x)$ & Softplus, $\log(1 + \exp(x))$ \\
$\displaystyle || \vx ||_p $ & $\normlp$ norm of $\vx$ \\
$\displaystyle || \vx || $ & $\normltwo$ norm of $\vx$ \\
$\displaystyle x^+$ & Positive part of $x$, i.e., $\max(0,x)$\\
$\displaystyle \1_\mathrm{condition}$ & is 1 if the condition is true, 0 otherwise\\
\end{tabular}
\egroup
\vspace{0.25cm}



\section{Final instructions}
Do not change any aspects of the formatting parameters in the style files.
In particular, do not modify the width or length of the rectangle the text
should fit into, and do not change font sizes (except perhaps in the
\textsc{References} section; see below). Please note that pages should be
numbered.

\section{Preparing PostScript or PDF files}

Please prepare PostScript or PDF files with paper size ``US Letter'', and
not, for example, ``A4''. The -t
letter option on dvips will produce US Letter files.

Consider directly generating PDF files using \verb+pdflatex+
(especially if you are a MiKTeX user).
PDF figures must be substituted for EPS figures, however.

Otherwise, please generate your PostScript and PDF files with the following commands:
\begin{verbatim}
dvips mypaper.dvi -t letter -Ppdf -G0 -o mypaper.ps
ps2pdf mypaper.ps mypaper.pdf
\end{verbatim}

\subsection{Margins in LaTeX}

Most of the margin problems come from figures positioned by hand using
\verb+\special+ or other commands. We suggest using the command
\verb+\includegraphics+
from the graphicx package. Always specify the figure width as a multiple of
the line width as in the example below using .eps graphics
\begin{verbatim}
   \usepackage[dvips]{graphicx} ...
   \includegraphics[width=0.8\linewidth]{myfile.eps}
\end{verbatim}
or % Apr 2009 addition
\begin{verbatim}
   \usepackage[pdftex]{graphicx} ...
   \includegraphics[width=0.8\linewidth]{myfile.pdf}
\end{verbatim}
for .pdf graphics.
See section~4.4 in the graphics bundle documentation (\url{http://www.ctan.org/tex-archive/macros/latex/required/graphics/grfguide.ps})

A number of width problems arise when LaTeX cannot properly hyphenate a
line. Please give LaTeX hyphenation hints using the \verb+\-+ command.

\subsection*{AI use statement}

(This section is \textbf{required} and does not count toward the page limit.)

In this work, we used generative AI tools for [tasks with required disclosure].
We have not used generative AI tools for [other tasks with required disclosure],
and [the rest of the required disclosure tasks] are not applicable to this work.
Additionally, we used generative AI tools for [tasks with recommended
disclosure]. We have reviewed all AI-assisted work. [Elaborate. For example, “we
checked LLM-generated research ideas for potential plagiarism through a manual
literature survey”, “LLM-generated code was verified and tested for correctness
by 2 authors”, etc.]. We take responsibility for the final content of this work,
including text, claims or artifacts produced with the aid of generative AI.

See the ICLR 2027 AI Policy for Authors for more details. This statement should
not be more than 1 page.

\subsection*{Ethics statement}

(This section is \textbf{recommended} and does not count toward the page limit.)

If authors feel that their paper submission raises questions regarding the Code
of Ethics, they are encouraged to include a paragraph of Ethics Statement (at
the end of the main text before references) to address potential concerns where
appropriate. Topics include, but are not limited to, studies that involve human
subjects, practices to data set releases, potentially harmful insights,
methodologies and applications, potential conflicts of interest and sponsorship,
discrimination/bias/fairness concerns, privacy and security issues, legal
compliance, and research integrity issues (e.g., IRB, documentation, research
ethics). This statement should not be more than 1 page.

\subsection*{Reproducibility statement}

(This section is \textbf{recommended} and does not count toward the page limit.)

It is important that the work published in ICLR is reproducible. Authors are
strongly encouraged to include a paragraph-long Reproducibility Statement at the
end of the main text (before references) to discuss the efforts that have been
made to ensure reproducibility. This paragraph should not itself describe
details needed for reproducing the results, but rather reference the parts of
the main paper, appendix, and supplemental materials that will help with
reproducibility. For example, for novel models or algorithms, a link to an
anonymous downloadable source code can be submitted as supplementary materials;
for theoretical results, clear explanations of any assumptions and a complete
proof of the claims can be included in the appendix; for any datasets used in
the experiments, a complete description of the data processing steps can be
provided in the supplementary materials. Each of the above are examples of
things that can be referenced in the reproducibility statement.

\subsubsection*{Author Contributions}
If you'd like to, you may include  a section for author contributions as is done
in many journals. This is optional and at the discretion of the authors.

\subsubsection*{Acknowledgments}
Use unnumbered third level headings for the acknowledgments. All
acknowledgments, including those to funding agencies, go at the end of the paper.


\section{A sector-wide quantitative swing-by theorem}
\label{sbsector:main}

The following result strengthens the noiseless one-neuron reversal
theorem to a common sector, with uniform amplitude and time bounds.
The two-neuron extension below fits both training concepts. A further
corollary gives a fixed positive probability under planar isotropic
initialization. The leakage scales are inherited from initialization;
this is not the high-probability noisy many-neuron theorem.

Consider the planar ReLU autoencoder
\[
 f(x)=\sum_i w_i(u_i^\top x)_+,\qquad
 \mathcal L=\frac14\bigl(\|f(\mu_1e_1)-\mu_1e_1\|^2
                         +\|f(\mu_2e_2)-\mu_2e_2\|^2\bigr),
 \quad \mu_1,\mu_2>0.
\]
Time is gradient-flow time for this normalization. Probe error is
$E(\xi,t)=\|f_t(\xi)-\xi\|^2/2$.

\begin{theorem}[A whole sector with quantitative reversal]
\label{sbsector:theorem}
Take one neuron and
\[
 u(0)=w(0)=(a_0,-\eta),\qquad
 a_0,\eta>0,\qquad a_0^2+\eta^2\le\frac1{16}.
\]
Write $\lambda=\mu_1^2/2$, $\beta=1-\eta^2$, and, for every
$z\in[1/2,1]$, let
\[
 k_z=\sqrt{1-\eta^2z^2},\qquad \xi_z=(\eta z,-k_z).
\]
The following conclusions hold for the actual gradient flow.
\begin{enumerate}[label=(\roman*)]
\item Both training gates and all these probe gates are fixed for every
finite time. There are increasing functions $a,c$ and a decreasing
positive function $D$ such that
\[
 u=(a,-\eta),\quad w=(c,-\eta D),\qquad
 M=wu^\top=
 \begin{pmatrix}ac&-\eta c\\-\eta aD&\eta^2D\end{pmatrix}.
\]
The same matrix $M$ acts on the entire probe sector. Its four scaled
coefficients are bounded by $2$, their time derivatives by $5\lambda$,
and transverse input/output energy is at most $2\eta^2$.

\item Define the finite hitting times
\[
 t_-:=\inf\{t:a(t)=1/4\},\qquad
 t_+:=\inf\{t:a(t)=1/\sqrt2\}.
\]
For every $z\in[1/2,1]$, $\dot E(\xi_z,t)<0$ on $[0,t_-]$ and
$\dot E(\xi_z,t)>0$ on $[t_+,\infty)$. Every global minimizing time
$t_{\rm rev}(z)$ is a strict local minimum and satisfies
\begin{equation}\label{sbsector:time}
 \frac1\lambda\log\frac1{4a_0}
 \le t_-<t_{\rm rev}(z)<t_+
 \le\frac1{2\lambda\beta}
 \log\frac{\beta-a_0^2}{a_0^2(2\beta-1)}.
\end{equation}
No uniqueness of the minimum is asserted.

\item If $E_\infty(\xi_z)=\lim_{t\to\infty}E(\xi_z,t)$, then
\begin{equation}\label{sbsector:amplitude}
 \frac14\eta^2
 < E_\infty(\xi_z)-\min_{t\ge0}E(\xi_z,t)
 <4\eta^2.
\end{equation}
The sector has angular width
\[
 \arcsin\eta-\arcsin(\eta/2),\qquad
 \frac\eta2\le\arcsin\eta-\arcsin(\eta/2)\le\eta.
\]
Thus its width is $\Theta(\eta)$ and its rebound amplitude is
$\Theta(\eta^2)$, with constants uniform in $a_0,\eta,z$ in the stated
range. The upper and lower times in \eqref{sbsector:time} are of order
$\lambda^{-1}\log(1/a_0)$ as $a_0\to0$.
\end{enumerate}
\end{theorem}

\begin{corollary}[Two learned concepts and two reversal sectors]
\label{sbsector:two-concept}
Take two neurons initialized by
\[
 u_1(0)=w_1(0)=(a_0,-\eta),\qquad
 u_2(0)=w_2(0)=(-\zeta,b_0),
\]
where $a_0,b_0,\eta,\zeta>0$ and
$a_0^2+\eta^2\le1/16$, $b_0^2+\zeta^2\le1/16$.
Both training reconstructions converge to their targets. The conclusions
of Theorem~\ref{sbsector:theorem} hold for the strong probes
$(\eta z,-\sqrt{1-\eta^2z^2})$, $z\in[1/2,1]$.
They also hold, with $(a_0,\eta,\lambda)$ replaced by
$(b_0,\zeta,\mu_2^2/2)$, for the weak probes
$(-\sqrt{1-\zeta^2z^2},\zeta z)$.
The sector widths and rebound amplitudes are respectively
$\Theta(\eta),\Theta(\eta^2)$ and
$\Theta(\zeta),\Theta(\zeta^2)$.
\end{corollary}

\begin{corollary}[A positive-probability isotropic initialization result]
\label{sbsector:isotropic}
In the same noiseless planar dataset, take $h=2$ and independently
initialize $u_i(0)=w_i(0)=\varepsilon\widehat u_i$, with
$\widehat u_i$ uniform on $\mathbb S^1$ and $0<\varepsilon\le1/4$.
With probability $1/8$, one neuron starts in quadrant IV and the other
in quadrant II. On that event, both concepts are fitted and
Corollary~\ref{sbsector:two-concept} holds with the positive coordinate
magnitudes supplied by initialization, after relabeling if necessary.

There is a subevent of probability $1/72$ on which both learned
concepts have a probe sector of width between $\varepsilon/4$ and
$\varepsilon$, and every probe in each sector has rebound amplitude
strictly between $\varepsilon^2/16$ and $3\varepsilon^2$.
For each concept $p$ and every such probe, every global minimizing time
satisfies the deterministic bounds
\begin{equation}\label{sbsector:isotropic-time}
 \frac1{\lambda_p}\log\frac1{2\sqrt3\varepsilon}
 <t_{\rm rev}
 <\frac8{15\lambda_p}\log\frac{32}{7\varepsilon^2},
 \qquad \lambda_p=\mu_p^2/2.
\end{equation}
The probability constants are lower guarantees from explicit angular
events; no claim is made about the remaining initializations.
\end{corollary}

\begin{remark}[What is and is not being derived]
The training and probe gates, amplification, matrix grading, bounded
coefficient rates, reversal signs, and amplitude bounds are conclusions
from the displayed deterministic initialization. The sizes $\eta$ and
$\zeta$ are inherited from that initialization. In
Corollary~\ref{sbsector:isotropic} they are random initial coordinate
magnitudes, comparable to $\varepsilon$ on the smaller event. Their
values are not derived from noisy weak-basin creation or a large
random population of neurons. The two-neuron result fits both concepts,
but its two pure training sectors decouple; it does not establish the
noisy creation-and-capture mechanism.
\end{remark}


\bibliography{iclr2027_conference}
\bibliographystyle{iclr2027_conference}

\appendix
\section{Appendix}

\section{Notation}

Throughout the two-concept proof, let
\[
\mathcal S
:=
\operatorname{span}\{e_1,e_2\}.
\]
For $p\in\{1,2\}$, define the within-concept-span transverse projection
\[
\Pi_{\perp p}x
:=
P_{\mathcal S}x
-
\langle x,e_p\rangle e_p.
\]
Thus $\Pi_{\perp p}$ removes only the $e_p$ coordinate inside
$\mathcal S$, whereas $P_{\mathcal S^\perp}$ removes the entire concept
span.

Define the in-span unit circle
\[
\mathbb S(\mathcal S)
:=
\left\{
v\in\mathcal S:
\|v\|=1
\right\}.
\]
For $\phi\in\mathbb R$, write
\[
\widetilde u(\phi)
:=
\cos\phi\,e_1+\sin\phi\,e_2.
\]
For an angular set
\[
J\subset\mathbb R/(2\pi\mathbb Z),
\]
write
\[
\widetilde u(J)
:=
\left\{
\widetilde u(\phi):
\phi\in J
\right\}
\subset
\mathbb S(\mathcal S).
\]

\section{Proof of Theorem 1}

%%%%%%%%%%%%%%%%%%%%%%%%%%%%%%%%%%%%%%%%%%%%%%%%%%%%%%%%%%%%%%%%%%%%%%%%%%%%%%%
% T1.1: Strengthened master event
%%%%%%%%%%%%%%%%%%%%%%%%%%%%%%%%%%%%%%%%%%%%%%%%%%%%%%%%%%%%%%%%%%%%%%%%%%%%%%%

\begin{lemma}[Strengthened two-concept master event]
\label{lem:Hstrong}
Fix $\delta\in(0,1)$ and consider the two-concept SIM model
\[
x_k^{(1)}
=
\mu_1 e_1+\xi_k^{(1)},
\qquad
x_k^{(2)}
=
\mu_2 e_2+\xi_k^{(2)},
\]
where
\[
\xi_k^{(p)}
\stackrel{\mathrm{i.i.d.}}{\sim}
\mathcal N
\!\left(
0,
\operatorname{diag}(\sigma_1^2,\sigma_2^2,0,\ldots,0)
\right),
\]
independently across $p\in\{1,2\}$ and $k\in[n]$, and assume
\[
\sigma_1>0,
\qquad
\sigma_2>0.
\]

Let
\[
\mathcal S=\operatorname{span}\{e_1,e_2\},
\qquad
\mu_{\min}:=\min\{\mu_1,\mu_2\},
\qquad
\mu_{\max}:=\max\{\mu_1,\mu_2\},
\]
and
\[
\sigma_{\min}:=\min\{\sigma_1,\sigma_2\},
\qquad
\sigma_{\max}:=\max\{\sigma_1,\sigma_2\}.
\]

For $\phi\in\mathbb R$, write
\[
\widetilde u(\phi)
=
\cos\phi\,e_1+\sin\phi\,e_2
\]
and define the directional noise scale
\[
\mathsf v(\phi)^2
:=
\sigma_1^2\cos^2\phi
+
\sigma_2^2\sin^2\phi.
\]
Then
\[
0<\sigma_{\min}
\le
\mathsf v(\phi)
\le
\sigma_{\max}
\qquad
\forall\phi\in\mathbb R.
\]

Set
\[
L_\delta
:=
\sqrt{2\log\frac{48n}{\delta}},
\]
\[
R_{\rm off}
:=
\sigma_{\max}L_\delta,
\qquad
X_\delta
:=
\mu_{\max}
+
\sqrt{2}\,\sigma_{\max}L_\delta,
\]
and
\[
C_{\rm abs}
:=
\sqrt{\frac2\pi}
+
C
\sqrt{
\frac{\log(24/\delta)}{n}
}.
\]

The halfspace class used below has VC dimension at most $2$, whereas
the centered-slab class has VC dimension at most $4$.  We use the single
uniform deviation scale
\[
\Delta_{\rm VC}
:=
C
\sqrt{
\frac{4+\log(24/\delta)}{n}
},
\]
which is the larger of the two corresponding VC bounds.
Here and below $C>0$ denotes a sufficiently large universal constant.

Assume the SNR condition
\[
\mu_{\min}
\ge
2\sigma_{\max}L_\delta.
\]

Then there exists an event
\[
\mathcal H_{\rm strong}(\delta)
\]
with
\[
\Pr\!\left[
\mathcal H_{\rm strong}(\delta)
\right]
\ge
1-\delta
\]
on which all of the following hold simultaneously.

\begin{enumerate}[label=(\roman*),leftmargin=2.4em]

\item
\textbf{Uniform coordinate-noise envelope.}
For every $p,q\in\{1,2\}$ and $k\in[n]$,
\[
|(\xi_k^{(p)})_q|
\le
\sigma_qL_\delta.
\]
Consequently,
\[
(x_k^{(p)})_p
\ge
\frac{\mu_p}{2},
\qquad
\|\Pi_{\perp p}x_k^{(p)}\|
\le
R_{\rm off},
\]
and
\[
\|x_k^{(p)}\|
\le
X_\delta.
\]

\item
\textbf{Uniform absolute empirical first moments.}
For every $p,q\in\{1,2\}$,
\[
\frac1n
\sum_{k=1}^n
|(\xi_k^{(p)})_q|
\le
C_{\rm abs}\sigma_q.
\]
Hence, for $q\neq p$,
\[
\frac1n
\sum_{k=1}^n
|(x_k^{(p)})_q|
\le
C_{\rm abs}\sigma_q,
\]
while
\[
\frac1n
\sum_{k=1}^n
|(x_k^{(p)})_p|
\le
\mu_p+C_{\rm abs}\sigma_p.
\]

\item
\textbf{Uniform halfspace-count concentration.}
For
\[
H_\phi
:=
\left\{
x:
\widetilde u(\phi)^\top x>0
\right\},
\]
define
\[
\widehat P_p(H_\phi)
:=
\frac1n
\sum_{k=1}^n
\mathbf 1\{
x_k^{(p)}\in H_\phi
\}.
\]
Then
\[
\sup_{\phi\in\mathbb R}
\left|
\widehat P_p(H_\phi)
-
P_p(H_\phi)
\right|
\le
\Delta_{\rm VC},
\qquad
p=1,2.
\]
Moreover,
\[
P_1(H_\phi)
=
\Phi\!\left(
\frac{\mu_1\cos\phi}{\mathsf v(\phi)}
\right),
\qquad
P_2(H_\phi)
=
\Phi\!\left(
\frac{\mu_2\sin\phi}{\mathsf v(\phi)}
\right).
\]

\item
\textbf{Uniform moving-slab anti-concentration.}
For
\[
S_{\phi,\tau}
:=
\left\{
x:
\left|
\widetilde u(\phi)^\top x
\right|
\le\tau
\right\},
\qquad
\tau\ge0,
\]
define
\[
\widehat P_p(S_{\phi,\tau})
:=
\frac1n
\sum_{k=1}^n
\mathbf 1\{
x_k^{(p)}\in S_{\phi,\tau}
\}.
\]
Then
\[
\sup_{\phi\in\mathbb R,\ \tau\ge0}
\left|
\widehat P_p(S_{\phi,\tau})
-
P_p(S_{\phi,\tau})
\right|
\le
\Delta_{\rm VC},
\qquad
p=1,2.
\]
Furthermore,
\[
P_p(S_{\phi,\tau})
\le
\sqrt{\frac2\pi}
\frac{\tau}{\mathsf v(\phi)},
\]
and hence
\[
\boxed{
\widehat P_p(S_{\phi,\tau})
\le
\sqrt{\frac2\pi}
\frac{\tau}{\mathsf v(\phi)}
+
\Delta_{\rm VC}
}
\]
simultaneously for every $p\in\{1,2\}$,
$\phi\in\mathbb R$, and $\tau\ge0$.
In particular,
\[
\boxed{
\widehat P_p(S_{\phi,\tau})
\le
\sqrt{\frac2\pi}
\frac{\tau}{\sigma_{\min}}
+
\Delta_{\rm VC}.
}
\]

\end{enumerate}
\end{lemma}


\begin{proof}
We construct $\mathcal H_{\rm strong}(\delta)$ as the intersection of
four events and then apply a union bound.

%%%%%%%%%%%%%%%%%%%%%%%%%%%%%%%%%%%%%%%%%%%%%%%%%%%%%%%%%%%%%%%%%%%%%%%%%%%%%%%
\paragraph{Step 1: coordinatewise Gaussian envelope.}
%%%%%%%%%%%%%%%%%%%%%%%%%%%%%%%%%%%%%%%%%%%%%%%%%%%%%%%%%%%%%%%%%%%%%%%%%%%%%%%

For every fixed $(p,q,k)$,
\[
\frac{(\xi_k^{(p)})_q}{\sigma_q}
\sim
\mathcal N(0,1).
\]
Therefore the standard Gaussian tail bound gives
\[
\Pr\!\left(
|(\xi_k^{(p)})_q|
>
\sigma_qL_\delta
\right)
\le
2e^{-L_\delta^2/2}.
\]
By the definition
\[
L_\delta
=
\sqrt{2\log\frac{48n}{\delta}},
\]
we have
\[
2e^{-L_\delta^2/2}
=
\frac{\delta}{24n}.
\]

There are $4n$ triples $(p,q,k)$, since there are two clusters,
two noisy coordinates, and $n$ samples per cluster. Hence
\[
\Pr\!\left[
\max_{p,q,k}
\frac{|(\xi_k^{(p)})_q|}{\sigma_q}
>
L_\delta
\right]
\le
4n\frac{\delta}{24n}
=
\frac{\delta}{6}.
\]

Thus, except on an event of probability at most $\delta/6$,
\[
|(\xi_k^{(p)})_q|
\le
\sigma_qL_\delta
\]
simultaneously for every $p,q,k$.

On this event,
\[
(x_k^{(p)})_p
=
\mu_p+(\xi_k^{(p)})_p
\ge
\mu_p-\sigma_pL_\delta.
\]
Since
\[
\sigma_pL_\delta
\le
\sigma_{\max}L_\delta
\le
\frac{\mu_{\min}}2
\le
\frac{\mu_p}{2},
\]
the SNR assumption implies
\[
(x_k^{(p)})_p
\ge
\frac{\mu_p}{2}.
\]

Because there are only two concept coordinates, if $q\neq p$ then
\[
\Pi_{\perp p}x_k^{(p)}
=
(\xi_k^{(p)})_q e_q,
\]
and therefore
\[
\|\Pi_{\perp p}x_k^{(p)}\|
\le
\sigma_{\max}L_\delta
=
R_{\rm off}.
\]

Finally,
\[
x_k^{(p)}
=
\mu_pe_p+\xi_k^{(p)},
\]
so
\[
\|x_k^{(p)}\|
\le
\mu_p+\|\xi_k^{(p)}\|.
\]
The coordinatewise envelope gives
\[
\|\xi_k^{(p)}\|
\le
\sqrt{
\sigma_1^2+\sigma_2^2
}\,L_\delta
\le
\sqrt2\,\sigma_{\max}L_\delta.
\]
Hence
\[
\|x_k^{(p)}\|
\le
\mu_{\max}
+
\sqrt2\,\sigma_{\max}L_\delta
=
X_\delta.
\]

%%%%%%%%%%%%%%%%%%%%%%%%%%%%%%%%%%%%%%%%%%%%%%%%%%%%%%%%%%%%%%%%%%%%%%%%%%%%%%%
\paragraph{Step 2: absolute empirical first moments.}
%%%%%%%%%%%%%%%%%%%%%%%%%%%%%%%%%%%%%%%%%%%%%%%%%%%%%%%%%%%%%%%%%%%%%%%%%%%%%%%

Fix $p,q$ and define
\[
Y_k
:=
\frac{|(\xi_k^{(p)})_q|}{\sigma_q}.
\]
Then $Y_k$ is a folded standard Gaussian and
\[
\mathbb EY_k
=
\sqrt{\frac2\pi}.
\]

Because the map
\[
z\mapsto |z|
\]
is $1$-Lipschitz, Gaussian concentration implies that each centered
variable
\[
Y_k-\mathbb EY_k
\]
is sub-Gaussian with a universal sub-Gaussian parameter. Since the
$Y_k$ are independent, their average
\[
\frac1n
\sum_{k=1}^n
\left(
Y_k-\mathbb EY_k
\right)
\]
is sub-Gaussian with parameter of order $n^{-1/2}$.
The standard sub-Gaussian tail bound therefore gives, for a sufficiently
large universal $C$,
\[
\Pr\!\left[
\frac1n\sum_{k=1}^nY_k
>
\sqrt{\frac2\pi}
+
C
\sqrt{
\frac{\log(24/\delta)}{n}
}
\right]
\le
\frac{\delta}{24}.
\]

Union bounding over the four pairs $(p,q)$ gives total failure
probability at most
\[
4\frac{\delta}{24}
=
\frac{\delta}{6}.
\]
Thus
\[
\frac1n
\sum_{k=1}^n
|(\xi_k^{(p)})_q|
\le
C_{\rm abs}\sigma_q
\]
simultaneously for every $p,q$.

If $q\neq p$, then
\[
(x_k^{(p)})_q
=
(\xi_k^{(p)})_q,
\]
which gives the off-axis conclusion. On the own coordinate,
\[
|(x_k^{(p)})_p|
=
|\mu_p+(\xi_k^{(p)})_p|
\le
\mu_p+|(\xi_k^{(p)})_p|,
\]
and hence
\[
\frac1n
\sum_{k=1}^n
|(x_k^{(p)})_p|
\le
\mu_p+C_{\rm abs}\sigma_p.
\]

%%%%%%%%%%%%%%%%%%%%%%%%%%%%%%%%%%%%%%%%%%%%%%%%%%%%%%%%%%%%%%%%%%%%%%%%%%%%%%%
\paragraph{Step 3: uniform halfspace counts.}
%%%%%%%%%%%%%%%%%%%%%%%%%%%%%%%%%%%%%%%%%%%%%%%%%%%%%%%%%%%%%%%%%%%%%%%%%%%%%%%

Consider
\[
\mathscr H
=
\left\{
H_\phi:
\phi\in\mathbb R
\right\},
\qquad
H_\phi
=
\left\{
x:
\widetilde u(\phi)^\top x>0
\right\}.
\]
This is a subclass of the homogeneous halfspaces in $\mathbb R^2$.
Homogeneous halfspaces in $\mathbb R^2$ have VC dimension $2$, so
\[
d_{\rm VC}(\mathscr H)\le2.
\]

The standard uniform VC inequality therefore yields, for every fixed
cluster $p$,
\[
\sup_{\phi}
\left|
\widehat P_p(H_\phi)-P_p(H_\phi)
\right|
\le
C
\sqrt{
\frac{
2+\log(24/\delta)
}{n}
}
\le
\Delta_{\rm VC}
\]
except on an event of probability at most $\delta/12$.

Here the use of $\log(24/\delta)$ is deliberately slightly more
conservative than is required for failure probability $\delta/12$.
Union bounding over the two clusters gives failure probability at most
$\delta/6$.

It remains to compute the population probabilities. For concept $1$,
\[
\widetilde u(\phi)^\top x^{(1)}
=
\mu_1\cos\phi
+
\cos\phi\,\xi_1
+
\sin\phi\,\xi_2,
\]
whereas for concept $2$,
\[
\widetilde u(\phi)^\top x^{(2)}
=
\mu_2\sin\phi
+
\cos\phi\,\xi_1
+
\sin\phi\,\xi_2.
\]
In both cases the centered Gaussian component has variance
\[
\mathsf v(\phi)^2
=
\sigma_1^2\cos^2\phi
+
\sigma_2^2\sin^2\phi.
\]
Therefore
\[
P_1(H_\phi)
=
\Phi\!\left(
\frac{\mu_1\cos\phi}{\mathsf v(\phi)}
\right)
\]
and
\[
P_2(H_\phi)
=
\Phi\!\left(
\frac{\mu_2\sin\phi}{\mathsf v(\phi)}
\right).
\]

%%%%%%%%%%%%%%%%%%%%%%%%%%%%%%%%%%%%%%%%%%%%%%%%%%%%%%%%%%%%%%%%%%%%%%%%%%%%%%%
\paragraph{Step 4: uniform moving-slab anti-concentration.}
%%%%%%%%%%%%%%%%%%%%%%%%%%%%%%%%%%%%%%%%%%%%%%%%%%%%%%%%%%%%%%%%%%%%%%%%%%%%%%%

Consider the class
\[
\mathscr S
=
\left\{
S_{\phi,\tau}:
\phi\in\mathbb R,\ \tau\ge0
\right\},
\]
where
\[
S_{\phi,\tau}
=
\left\{
x:
|\widetilde u(\phi)^\top x|
\le\tau
\right\}.
\]

We first bound its VC dimension. The slab condition is equivalent to
\[
|\widetilde u(\phi)^\top x|
\le\tau
\quad\Longleftrightarrow\quad
(\widetilde u(\phi)^\top x)^2-\tau^2\le0.
\]
Writing
\[
\widetilde u(\phi)
=
(\cos\phi,\sin\phi),
\]
we have
\[
(\widetilde u(\phi)^\top x)^2-\tau^2
=
\cos^2\phi\,x_1^2
+
2\sin\phi\cos\phi\,x_1x_2
+
\sin^2\phi\,x_2^2
-
\tau^2.
\]

Introduce the quadratic feature map
\[
\Psi(x)
=
\begin{pmatrix}
x_1^2\\
2x_1x_2\\
x_2^2\\
1
\end{pmatrix}
\in\mathbb R^4.
\]
Then
\[
x\in S_{\phi,\tau}
\]
if and only if
\[
\left\langle
\begin{pmatrix}
\cos^2\phi\\
\sin\phi\cos\phi\\
\sin^2\phi\\
-\tau^2
\end{pmatrix},
\Psi(x)
\right\rangle
\le0.
\]
Thus $\mathscr S$ is a subclass of homogeneous halfspaces in the
four-dimensional feature space. Consequently
\[
d_{\rm VC}(\mathscr S)\le4.
\]
In particular, the same representation handles the threshold parameter
$\tau$ through the coefficient $-\tau^2$; no separate discretization
or parameter-counting argument is required.

The uniform VC inequality therefore gives, for each fixed cluster $p$,
\[
\sup_{\phi,\tau}
\left|
\widehat P_p(S_{\phi,\tau})
-
P_p(S_{\phi,\tau})
\right|
\le
C
\sqrt{
\frac{
4+\log(24/\delta)
}{n}
}
=
\Delta_{\rm VC}
\]
except on an event of probability at most $\delta/12$.

Again, $\log(24/\delta)$ is chosen slightly more conservatively than
necessary for this failure level. A union bound over $p=1,2$ gives
total failure probability at most $\delta/6$.

Now fix $p,\phi,\tau$. The scalar random variable
\[
Y_p(\phi)
:=
\widetilde u(\phi)^\top x^{(p)}
\]
is Gaussian with variance
\[
\mathsf v(\phi)^2>0.
\]
Its mean depends on $p$ and $\phi$, but is irrelevant for the
anti-concentration estimate. Its density satisfies
\[
\sup_y f_{Y_p(\phi)}(y)
=
\frac{1}{
\sqrt{2\pi}\,\mathsf v(\phi)
}.
\]
Hence
\begin{align*}
P_p(S_{\phi,\tau})
&=
\Pr\!\left(
|Y_p(\phi)|\le\tau
\right)
\\
&=
\int_{-\tau}^{\tau}
f_{Y_p(\phi)}(y)\,dy
\\
&\le
2\tau
\sup_y f_{Y_p(\phi)}(y)
\\
&=
\sqrt{\frac2\pi}
\frac{\tau}{\mathsf v(\phi)}.
\end{align*}

Combining this population bound with the uniform VC deviation gives
\[
\widehat P_p(S_{\phi,\tau})
\le
\sqrt{\frac2\pi}
\frac{\tau}{\mathsf v(\phi)}
+
\Delta_{\rm VC}.
\]
Since
\[
\mathsf v(\phi)\ge\sigma_{\min}>0,
\]
we also obtain
\[
\widehat P_p(S_{\phi,\tau})
\le
\sqrt{\frac2\pi}
\frac{\tau}{\sigma_{\min}}
+
\Delta_{\rm VC}.
\]

%%%%%%%%%%%%%%%%%%%%%%%%%%%%%%%%%%%%%%%%%%%%%%%%%%%%%%%%%%%%%%%%%%%%%%%%%%%%%%%
\paragraph{Step 5: conclusion.}
%%%%%%%%%%%%%%%%%%%%%%%%%%%%%%%%%%%%%%%%%%%%%%%%%%%%%%%%%%%%%%%%%%%%%%%%%%%%%%%

Each of the four preceding event families fails with probability at most
$\delta/6$. Therefore
\[
\Pr\!\left[
\mathcal H_{\rm strong}(\delta)^c
\right]
\le
\frac{4\delta}{6}
=
\frac{2\delta}{3}
<
\delta.
\]
Thus
\[
\Pr\!\left[
\mathcal H_{\rm strong}(\delta)
\right]
\ge
1-\delta.
\]
\end{proof}


\begin{remark}[Downstream role of the master event]
\label{rem:Hstrong-downstream}
The halfspace and slab conclusions serve different purposes.

First, the uniform halfspace bound controls the empirical firing fraction
of a candidate direction:
\[
\widehat P_p(H_\phi)
=
\frac1n
\sum_{k=1}^n
\mathbf 1\{
\widetilde u(\phi)^\top x_k^{(p)}>0
\},
\]
uniformly in $\phi$. It will be used in the weak-seed survival and capture
arguments and in Lemma~\ref{lem:gated-moment-concentration}, where
population gated moments are transferred to the empirical dataset.

Second, the uniform slab bound controls the set of samples whose gate
decision can change under a small angular perturbation. Because the bound
is uniform in $\phi$, it can later be evaluated at a random,
data-dependent moving direction such as
\[
\phi=\phi_W(t)
\]
simultaneously for every time $t$.

Finally, on $\mathcal H_{\rm strong}(\delta)$,
\[
\|x_k^{(p)}\|
\le
X_\delta
\qquad
\forall p,k.
\]
Consequently every gated quadratic function of the form
\[
x
\longmapsto
x_r x_c
\mathbf 1\{
\widetilde u(\phi)^\top x>0
\}
\]
has deterministic sample envelope
\[
\left|
x_r x_c
\mathbf 1\{
\widetilde u(\phi)^\top x>0
\}
\right|
\le
X_\delta^2.
\]
Thus, conditional on the envelope supplied by
$\mathcal H_{\rm strong}(\delta)$, the later uniform gated-moment
analysis is a bounded-envelope empirical-process problem whose indicator
component is the VC class $\{H_\phi\}_{\phi\in\mathbb R}$.
\end{remark}


%%%%%%%%%%%%%%%%%%%%%%%%%%%%%%%%%%%%%%%%%%%%%%%%%%%%%%%%%%%%%%%%%%%%%%%%%%%%%%%
% Probability-one general-position event
%%%%%%%%%%%%%%%%%%%%%%%%%%%%%%%%%%%%%%%%%%%%%%%%%%%%%%%%%%%%%%%%%%%%%%%%%%%%%%%

\begin{lemma}[Almost-sure finite-sample general position]
\label{lem:general-position}

Under the two-concept SIM model with
\[
\sigma_1>0,
\qquad
\sigma_2>0,
\]
define $\mathcal G_{\mathrm{pos}}$ to be the event that

\begin{enumerate}[label=(\roman*),leftmargin=2em]
\item every training sample is nonzero;
\item no two distinct training samples are collinear through the origin.
\end{enumerate}

Then
\[
\boxed{
\mathbb P(\mathcal G_{\mathrm{pos}})=1.
}
\]

On $\mathcal G_{\mathrm{pos}}$, every one-dimensional linear subspace of
\[
\mathcal S=\operatorname{span}\{e_1,e_2\}
\]
contains at most one training sample. Consequently, uniformly over all
\[
\widetilde u\in\mathcal S,
\qquad
\|\widetilde u\|=1,
\]
one has
\[
\#\left\{
(p,k):
\widetilde u^\top x_k^{(p)}=0
\right\}
\le1,
\]
and hence
\[
\boxed{
\mathbb E_n
\left[
\mathbf 1
\{
\widetilde u^\top x=0
\}
\right]
\le
\frac1{2n}.
}
\]

The bound holds simultaneously for all directions and therefore also for
random, data-dependent, time-varying directions.
\end{lemma}

\begin{proof}
Each training sample has a density with respect to two-dimensional Lebesgue
measure on $\mathcal S$. Hence
\[
\mathbb P
\left(
x_k^{(p)}=0
\right)
=
0.
\]

Fix two distinct training samples $X$ and $Y$. Conditional on
\[
X=x\neq0,
\]
the set
\[
\operatorname{span}\{x\}
\]
is a one-dimensional linear subspace of $\mathcal S$ and therefore has
two-dimensional Lebesgue measure zero. Since $Y$ has a density on
$\mathcal S$,
\[
\mathbb P
\left(
Y\in\operatorname{span}\{X\}
\mid X
\right)
=
0
\]
almost surely. Thus
\[
\mathbb P
\left(
X
\text{ and }
Y
\text{ are collinear through the origin}
\right)
=
0.
\]

A finite union over all samples and all distinct sample pairs proves
\[
\mathbb P(\mathcal G_{\mathrm{pos}})=1.
\]

On $\mathcal G_{\mathrm{pos}}$, for every unit vector
$\widetilde u\in\mathcal S$, the set
\[
\left\{
x\in\mathcal S:
\widetilde u^\top x=0
\right\}
\]
is a one-dimensional linear subspace and hence contains at most one observed
sample. Since $\mathbb E_n$ assigns mass $1/(2n)$ to each sample,
\[
\mathbb E_n
\left[
\mathbf 1
\{
\widetilde u^\top x=0
\}
\right]
\le
\frac1{2n}.
\]
Because this statement holds simultaneously for every one-dimensional
subspace of $\mathcal S$, it may be evaluated at arbitrary random,
data-dependent directions.
\end{proof}


%%%%%%%%%%%%%%%%%%%%%%%%%%%%%%%%%%%%%%%%%%%%%%%%%%%%%%%%%%%%%%%%%%%%%%%%%%%%%%%
% Empirical second-moment control
%%%%%%%%%%%%%%%%%%%%%%%%%%%%%%%%%%%%%%%%%%%%%%%%%%%%%%%%%%%%%%%%%%%%%%%%%%%%%%%

\begin{lemma}[Empirical second-moment spectral bound]
\label{lem:second-moment}

Define
\[
\Sigma_n
:=
\mathbb E_n[xx^\top]
=
\frac1{2n}
\sum_{p=1}^2
\sum_{k=1}^n
x_k^{(p)}x_k^{(p)\top}.
\]
The corresponding population second-moment matrix is
\[
\Sigma_\star
:=
\frac12
\sum_{p=1}^2
\mathbb E
\left[
x^{(p)}x^{(p)\top}
\right]
=
\begin{pmatrix}
\mu_1^2/2+\sigma_1^2 & 0\\
0 & \mu_2^2/2+\sigma_2^2
\end{pmatrix}
\]
on $\mathcal S$, and vanishes on $\mathcal S^\perp$.

Set
\[
\lambda_\star
:=
\lambda_{\max}(\Sigma_\star)
=
\max_{p\in\{1,2\}}
\left(
\frac{\mu_p^2}{2}+\sigma_p^2
\right).
\]
For $\delta\in(0,1)$ define
\[
\Delta_\Sigma(\delta)
:=
C
\left[
\mu_{\max}\sigma_{\max}
\sqrt{\frac{\log(8/\delta)}{n}}
+
\sigma_{\max}^2
\left(
\sqrt{\frac{\log(8/\delta)}{n}}
+
\frac{\log(8/\delta)}{n}
\right)
\right],
\]
where $C>0$ is a sufficiently large universal constant.

Then there exists an event
\[
\mathcal G_\Sigma(\delta)
\]
with
\[
\mathbb P(\mathcal G_\Sigma(\delta))
\ge
1-\delta
\]
such that, on $\mathcal G_\Sigma(\delta)$,
\[
\boxed{
\|\Sigma_n-\Sigma_\star\|_{_{\mathrm{op}}}
\le
\Delta_\Sigma(\delta).
}
\]
Consequently,
\[
\boxed{
\lambda_{\max}(\Sigma_n)
\le
\lambda_\star+\Delta_\Sigma(\delta).
}
\]
\end{lemma}

\begin{proof}
For $p\in\{1,2\}$ write
\[
\overline\xi^{(p)}
:=
\frac1n
\sum_{k=1}^n
\xi_k^{(p)},
\qquad
\widehat\Sigma_\xi^{(p)}
:=
\frac1n
\sum_{k=1}^n
\xi_k^{(p)}\xi_k^{(p)\top},
\]
and let
\[
\Sigma_\xi
=
\operatorname{diag}
(\sigma_1^2,\sigma_2^2)
\]
on $\mathcal S$.

Expanding
\[
x_k^{(p)}
=
\mu_p e_p+\xi_k^{(p)}
\]
gives
\[
\Sigma_n-\Sigma_\star
=
\frac12
\sum_{p=1}^2
\left[
\mu_p
\left(
e_p\overline\xi^{(p)\top}
+
\overline\xi^{(p)}e_p^\top
\right)
+
\left(
\widehat\Sigma_\xi^{(p)}
-
\Sigma_\xi
\right)
\right].
\]

Since
\[
\overline\xi^{(p)}
\sim
\mathcal N
\left(
0,\frac1n\Sigma_\xi
\right),
\]
standard Gaussian concentration in the fixed two-dimensional space
$\mathcal S$ gives, simultaneously for $p=1,2$,
\[
\|\overline\xi^{(p)}\|
\le
C\sigma_{\max}
\sqrt{
\frac{\log(8/\delta)}{n}
}
\]
except on an event of probability at most $\delta/2$.

Likewise, standard Gaussian sample-covariance concentration in fixed
dimension gives, simultaneously for $p=1,2$,
\[
\left\|
\widehat\Sigma_\xi^{(p)}
-
\Sigma_\xi
\right\|_{_{\mathrm{op}}}
\le
C\sigma_{\max}^2
\left(
\sqrt{
\frac{\log(8/\delta)}{n}
}
+
\frac{\log(8/\delta)}{n}
\right)
\]
except on an event of probability at most $\delta/2$.

Using
\[
\left\|
e_p a^\top+a e_p^\top
\right\|_{_{\mathrm{op}}}
\le
2\|a\|,
\]
the displayed decomposition yields
\[
\|\Sigma_n-\Sigma_\star\|_{_{\mathrm{op}}}
\le
\Delta_\Sigma(\delta)
\]
after enlarging the universal constant $C$. Weyl's inequality then gives
the claimed upper bound on $\lambda_{\max}(\Sigma_n)$.
\end{proof}

\begin{remark}[No ambient-dimension dependence]
\label{rem:second-moment-dimension}

All training samples are supported on the fixed two-dimensional concept span
$\mathcal S$. Consequently, the concentration estimate in
Lemma~\ref{lem:second-moment} is a fixed-dimensional Gaussian
sample-covariance bound, and the universal constant $C$ is independent of
the ambient dimension $d$.
\end{remark}



%%%%%%%%%%%%%%%%%%%%%%%%%%%%%%%%%%%%%%%%%%%%%%%%%%%%%%%%%%%%%%%%%%%%%%%%%%%%%%%
% Theorem 1 confidence allocation
%%%%%%%%%%%%%%%%%%%%%%%%%%%%%%%%%%%%%%%%%%%%%%%%%%%%%%%%%%%%%%%%%%%%%%%%%%%%%%%

\paragraph{Confidence allocation.}

Fix
\[
\boxed{
\delta_{\mathrm{strong}}
:=
\frac{\delta}{2},
\qquad
\delta_{\Sigma}
:=
\frac{\delta}{8},
\qquad
\delta_{\mathrm{init}}
:=
\frac{\delta}{8},
\qquad
\delta_{\mathrm{gate}}
:=
\frac{\delta}{4}.
}
\]
Thus
\[
\boxed{
\delta_{\mathrm{strong}}
+
\delta_{\Sigma}
+
\delta_{\mathrm{init}}
+
\delta_{\mathrm{gate}}
=
\delta.
}
\]

The four budgets are reserved respectively for the strengthened finite-sample
data event, empirical second-moment concentration, initialization coverage,
and the later gated-moment concentration event. The probability-one event
\[
\mathcal G_{\mathrm{pos}}
\]
consumes no confidence budget.

%%%%%%%%%%%%%%%%%%%%%%%%%%%%%%%%%%%%%%%%%%%%%%%%%%%%%%%%%%%%%%%%%%%%%%%%%%%%%%%
% Data-only event
%%%%%%%%%%%%%%%%%%%%%%%%%%%%%%%%%%%%%%%%%%%%%%%%%%%%%%%%%%%%%%%%%%%%%%%%%%%%%%%

\begin{definition}[Data event]
\label{def:Hdata}

Define
\[
\boxed{
\mathcal H_{\mathrm{data}}(\delta)
:=
\mathcal H_{\mathrm{strong}}(\delta_{\mathrm{strong}})
\cap
\mathcal G_{\Sigma}(\delta_{\Sigma})
\cap
\mathcal G_{\mathrm{pos}}.
}
\]
\end{definition}

The invocation of Lemma~\ref{lem:Hstrong} requires the standing SNR
condition
\begin{equation}
\boxed{
\mu_{\min}
\ge
2\sigma_{\max}L_{\delta_{\mathrm{strong}}},
}
\label{eq:Hdata-snr}
\end{equation}
where
\[
L_{\delta_{\mathrm{strong}}}
:=
\sqrt{
2\log
\frac{48n}{\delta_{\mathrm{strong}}}
}.
\]
Under
\[
\delta_{\mathrm{strong}}
=
\frac{\delta}{2},
\]
this becomes
\[
L_{\delta_{\mathrm{strong}}}
=
\sqrt{
2\log
\frac{96n}{\delta}
}.
\]

Equivalently, condition~\eqref{eq:Hdata-snr} requires
\begin{equation}
\boxed{
n
\le
n_{\max}^{\mathrm{SNR}}
:=
\frac{\delta_{\mathrm{strong}}}{48}
\exp
\left(
\frac{\mu_{\min}^2}
{8\sigma_{\max}^2}
\right)
=
\frac{\delta}{96}
\exp
\left(
\frac{\mu_{\min}^2}
{8\sigma_{\max}^2}
\right).
}
\label{eq:nmax-snr}
\end{equation}

By Lemma~\ref{lem:Hstrong},
\[
\mathbb P
\left(
\mathcal H_{\mathrm{strong}}(\delta_{\mathrm{strong}})
\right)
\ge
1-\delta_{\mathrm{strong}},
\]
by Lemma~\ref{lem:second-moment},
\[
\mathbb P
\left(
\mathcal G_{\Sigma}(\delta_{\Sigma})
\right)
\ge
1-\delta_{\Sigma},
\]
and by Lemma~\ref{lem:general-position},
\[
\mathbb P(\mathcal G_{\mathrm{pos}})=1.
\]
Therefore
\[
\boxed{
\mathbb P
\left(
\mathcal H_{\mathrm{data}}(\delta)
\right)
\ge
1-\delta_{\mathrm{strong}}-\delta_{\Sigma}
=
1-\frac{5\delta}{8}.
}
\]

The event $\mathcal H_{\mathrm{data}}(\delta)$ contains only training-data
randomness. It is not the final theorem-level event.

\begin{remark}[Confidence-budget invariant]
\label{rem:confidence-budget-invariant}
The top-level allocation
\[
\delta
=
\delta_{\mathrm{strong}}
+
\delta_{\Sigma}
+
\delta_{\mathrm{init}}
+
\delta_{\mathrm{gate}}
\]
is fixed throughout the proof. If a later argument requires several
stochastic events of the same type, those events subdivide the corresponding
reserved budget internally. In particular, all stochastic events introduced
in the gated-moment analysis must have total failure probability at most
\[
\delta_{\mathrm{gate}}.
\]
No later lemma repartitions the top-level confidence budgets.
\end{remark}

%%%%%%%%%%%%%%%%%%%%%%%%%%%%%%%%%%%%%%%%%%%%%%%%%%%%%%%%%%%%%%%%%%%%%%%%%%%%%%%
% Deterministic theorem-level envelopes
%%%%%%%%%%%%%%%%%%%%%%%%%%%%%%%%%%%%%%%%%%%%%%%%%%%%%%%%%%%%%%%%%%%%%%%%%%%%%%%

\paragraph{Deterministic theorem-level envelopes.}

Define
\[
\boxed{
R_{\mathrm{off},\delta_{\mathrm{strong}}}
:=
\sigma_{\max}L_{\delta_{\mathrm{strong}}},
}
\]
\[
\boxed{
X_{\delta_{\mathrm{strong}}}
:=
\mu_{\max}
+
\sqrt2\,\sigma_{\max}L_{\delta_{\mathrm{strong}}},
}
\]
\[
\boxed{
C_{\mathrm{abs},\delta_{\mathrm{strong}}}
:=
\sqrt{\frac2\pi}
+
C
\sqrt{
\frac{
\log(24/\delta_{\mathrm{strong}})
}{n}
},
}
\]
and
\[
\boxed{
\Delta_{\mathrm{VC},\delta_{\mathrm{strong}}}
:=
C
\sqrt{
\frac{
4+\log(24/\delta_{\mathrm{strong}})
}{n}
}.
}
\]

For statements involving the fixed OOD probe $x^\star$, define
\[
\boxed{
X_\delta^\star
:=
\max
\left\{
X_{\delta_{\mathrm{strong}}},
\,
\|x^\star\|
\right\}.
}
\]

Also define
\[
\lambda_\star
:=
\max_{p\in\{1,2\}}
\left(
\frac{\mu_p^2}{2}
+
\sigma_p^2
\right),
\]
and
\[
\boxed{
\overline\lambda_\delta
:=
\lambda_\star
+
\Delta_\Sigma(\delta_{\Sigma}).
}
\]

On $\mathcal H_{\mathrm{data}}(\delta)$,
\[
\boxed{
\max_{x\in\mathcal D}
\|x\|
\le
X_{\delta_{\mathrm{strong}}},
}
\]
\[
\boxed{
\max_{x\in\mathcal D\cup\{x^\star\}}
\|x\|
\le
X_\delta^\star,
}
\]
and
\[
\boxed{
\lambda_{\max}(\Sigma_n)
\le
\overline\lambda_\delta.
}
\]

Purely training-dynamics estimates use $X_{\delta_{\mathrm{strong}}}$
whenever possible. Statements involving the OOD probe $x^\star$ use
$X_\delta^\star$.


%%%%%%%%%%%%%%%%%%%%%%%%%%%%%%%%%%%%%%%%%%%%%%%%%%%%%%%%%%%%%%%%%%%%%%%%%%%%%%%
% T1.1b: uniform finite-sample gated moments
%%%%%%%%%%%%%%%%%%%%%%%%%%%%%%%%%%%%%%%%%%%%%%%%%%%%%%%%%%%%%%%%%%%%%%%%%%%%%%%

\begin{lemma}[Uniform gated second moments with entrywise concentration]
\label{lem:gated-moment-concentration}

Work with the confidence split fixed above, and set
\[
L_s:=L_{\delta_{\mathrm{strong}}}.
\]
For $p,q\in\{1,2\}$ define the deterministic coordinate envelopes
\[
\boxed{
B_{p,q}
:=
\mu_p\mathbf 1\{q=p\}
+
\sigma_qL_s.
}
\]
Thus
\[
B_{1,1}=\mu_1+\sigma_1L_s,
\qquad
B_{1,2}=\sigma_2L_s,
\]
and
\[
B_{2,1}=\sigma_1L_s,
\qquad
B_{2,2}=\mu_2+\sigma_2L_s.
\]

For the population comparison only, define the coordinatewise clipped
variables
\[
\boxed{
\bar x_q^{(p)}
:=
\operatorname{clip}
\bigl(
 x_q^{(p)};
 -B_{p,q},B_{p,q}
\bigr).
}
\]
On $\mathcal H_{\mathrm{strong}}(\delta_{\mathrm{strong}})$, every empirical
sample obeys
\[
\bar x_{k,q}^{(p)}=x_{k,q}^{(p)}.
\]

For $p\in\{1,2\}$ and $\phi\in\mathbb R/(2\pi\mathbb Z)$ define
\[
\boxed{
\widehat A_p(\phi)
:=
\mathbb E_{p,n}
\left[
xx^\top
\mathbf 1\{
\widetilde u(\phi)^\top x>0
\}
\right],
}
\]
where
\[
\mathbb E_{p,n}[\varphi]
:=
\frac1n\sum_{k=1}^n\varphi(x_k^{(p)}),
\]
and define the clipped population target
\[
\boxed{
\bar A_{p,\star}(\phi)
:=
\mathbb E_p
\left[
\bar x^{(p)}\bar x^{(p)\top}
\mathbf 1\{
\widetilde u(\phi)^\top x^{(p)}>0
\}
\right].
}
\]
Write
\[
\widehat a_{rs}^{(p)}(\phi)
:=
e_r^\top\widehat A_p(\phi)e_s,
\qquad
\bar a_{rs,\star}^{(p)}(\phi)
:=
e_r^\top\bar A_{p,\star}(\phi)e_s.
\]

Define the entrywise envelopes
\[
\boxed{
E_{rs}^{(p)}
:=
B_{p,r}B_{p,s},
}
\]
and the entrywise concentration scales
\[
\boxed{
\varepsilon_{rs}^{(p)}
:=
C E_{rs}^{(p)}
\sqrt{
\frac{
\log(Cn)+\log(C/\delta_{\mathrm{gate}})
}{n}
}.
}
\tag{gate-entry-error}\label{eq:gate-entry-error}
\]

Then there exists an event $\mathcal G_{\mathrm{gate}}$ with
\[
\boxed{
\mathbb P(\mathcal G_{\mathrm{gate}})
\ge
1-\delta_{\mathrm{gate}}
}
\]
such that on
\[
\mathcal H_{\mathrm{strong}}(\delta_{\mathrm{strong}})
\cap
\mathcal G_{\mathrm{gate}}
\]
the following hold simultaneously.

\begin{enumerate}[label=\textnormal{(\roman*)},leftmargin=2.4em]

\item
\textbf{Whole-circle uniform gated-moment concentration.}
For every $p,r,s\in\{1,2\}$,
\[
\boxed{
\sup_{\phi\in\mathbb R/(2\pi\mathbb Z)}
\left|
\widehat a_{rs}^{(p)}(\phi)
-
\bar a_{rs,\star}^{(p)}(\phi)
\right|
\le
\varepsilon_{rs}^{(p)}.
}
\tag{uniform-entry}
\]
No population-tail correction is present: the deterministic comparison
object is the clipped population moment itself.

If
\[
\varepsilon_{\mathrm{op}}^{(p)}
:=
\max\left\{
\varepsilon_{11}^{(p)}+\varepsilon_{12}^{(p)},
\varepsilon_{22}^{(p)}+\varepsilon_{12}^{(p)}
\right\},
\]
then
\[
\boxed{
\sup_\phi
\left\|
\widehat A_p(\phi)
-
\bar A_{p,\star}(\phi)
\right\|_{\mathrm{op}}
\le
\varepsilon_{\mathrm{op}}^{(p)}.
}
\tag{uniform-op}
\]

\item
\textbf{Partially-gated concept-$1$ band.}
Define
\[
\Gamma_{1,n}(\phi)
:=
\widetilde u(\phi)^\top
\widehat A_1(\phi)
\widetilde u(\phi),
\]
and
\[
\bar\Gamma_{1,\star}(\phi)
:=
\widetilde u(\phi)^\top
\bar A_{1,\star}(\phi)
\widetilde u(\phi).
\]
For any deterministic angular band $J$,
\[
\boxed{
\sup_{\phi\in J}\Gamma_{1,n}(\phi)
\le
\sup_{\phi\in J}\bar\Gamma_{1,\star}(\phi)
+
\varepsilon_{\mathrm{op}}^{(1)}.
}
\tag{band-Gamma}
\]
Consequently,
\[
\boxed{
\sup_{\phi\in J}\sqrt{\Gamma_{1,n}(\phi)}
\le
\sqrt{
\sup_{\phi\in J}\bar\Gamma_{1,\star}(\phi)
+
\varepsilon_{\mathrm{op}}^{(1)}
}.
}
\]
In particular, this applies to the concept-$1$ partially-gated band whose
inner boundary is the same $\theta_C$ appearing in the pure-sector
condition below.

\item
\textbf{Moving-angle modulus from the master slab event.}
Let
\[
\phi,\phi'\in\mathbb R/(2\pi\mathbb Z),
\qquad
\eta:=\operatorname{dist}_{\angle}(\phi,\phi'),
\]
and set
\[
\tau_\eta
:=
2X_{\delta_{\mathrm{strong}}}
\sin\frac\eta2.
\]
Then for every $p,r,s\in\{1,2\}$,
\[
\boxed{
\begin{aligned}
&\left|
\widehat a_{rs}^{(p)}(\phi)
-
\widehat a_{rs}^{(p)}(\phi')
\right|
\\
&\qquad
\le
E_{rs}^{(p)}
\left[
\sqrt{\frac2\pi}
\frac{\tau_\eta}{\mathsf v(\phi)}
+
\Delta_{\mathrm{VC},\delta_{\mathrm{strong}}}
\right]
\\
&\qquad
\le
E_{rs}^{(p)}
\left[
\sqrt{\frac2\pi}
\frac{
2X_{\delta_{\mathrm{strong}}}\sin(\eta/2)
}{\sigma_{\min}}
+
\Delta_{\mathrm{VC},\delta_{\mathrm{strong}}}
\right].
\end{aligned}
}
\tag{gate-modulus}\label{eq:gate-modulus}
\]
The same estimate holds with $\phi'$ as the slab center, so one may use the
better of the two variance denominators.  The bound is valid at random,
data-dependent moving angles because Lemma~\ref{lem:Hstrong} already gives
uniform control simultaneously over all slab centers and widths.

\item
\textbf{Single-sample boundary jumps.}
On $\mathcal G_{\mathrm{pos}}$, at every angular gate crossing at most one
training sample changes status. Therefore
\[
\boxed{
\left|
\widehat a_{rs}^{(p)}(\phi_0^+)
-
\widehat a_{rs}^{(p)}(\phi_0^-)
\right|
\le
\frac{E_{rs}^{(p)}}{n}.
}
\tag{gate-jump}
\]
In particular, the off-diagonal concept-$1$ boundary jump is at most
\[
\boxed{
\frac{
(\mu_1+\sigma_1L_s)\sigma_2L_s
}{n},
}
\]
rather than $X_{\delta_{\mathrm{strong}}}^2/n$.

\item
\textbf{Inner-region deterministic constants.}
For a later deterministic interval
\[
J_{\mathrm{in}}
=
[-\theta_G,\theta_H],
\]
define
\[
\underline a_{11}^{(p)}
:=
\inf_{\phi\in J_{\mathrm{in}}}
\bar a_{11,\star}^{(p)}(\phi),
\]
\[
\overline a_{22}^{(p)}
:=
\sup_{\phi\in J_{\mathrm{in}}}
\bar a_{22,\star}^{(p)}(\phi),
\]
and
\[
\overline a_{12}^{(p)}
:=
\sup_{\phi\in J_{\mathrm{in}}}
\left|
\bar a_{12,\star}^{(p)}(\phi)
\right|.
\]
Then uniformly over $\phi\in J_{\mathrm{in}}$,
\[
\boxed{
\widehat a_{11}^{(p)}(\phi)
\ge
\underline a_{11}^{(p)}
-
\varepsilon_{11}^{(p)},
}
\]
\[
\boxed{
\widehat a_{22}^{(p)}(\phi)
\le
\overline a_{22}^{(p)}
+
\varepsilon_{22}^{(p)},
}
\]
and
\[
\boxed{
\left|
\widehat a_{12}^{(p)}(\phi)
\right|
\le
\overline a_{12}^{(p)}
+
\varepsilon_{12}^{(p)}.
}
\]
These are the data-only constants used below to construct the inner-region
anisotropy and branch margins.

\item
\textbf{Explicit gated-moment sample-size floor.}
Let $\epsilon_{rs,\mathrm{adm}}^{(p)}>0$ be the error budget assigned
downstream to entry $(r,s)$ of cluster $p$. A sufficient condition for
\[
\varepsilon_{rs}^{(p)}
\le
\epsilon_{rs,\mathrm{adm}}^{(p)}
\]
is
\[
\boxed{
n
\ge
C
\frac{(E_{rs}^{(p)})^2}{(\epsilon_{rs,\mathrm{adm}}^{(p)})^2}
\left[
\log\left(
C\frac{E_{rs}^{(p)}}{\epsilon_{rs,\mathrm{adm}}^{(p)}}
\right)
+
\log\frac{C}{\delta_{\mathrm{gate}}}
\right].
}
\tag{gate-sample-floor}
\]
Define $n_{\mathrm{gate}}$ as the maximum of the right-hand side over the
entries actually used downstream. The final admissible region contains
\[
\boxed{n\ge n_{\mathrm{gate}}.}
\]
For the concept-$1$ off-diagonal entry, the sharper scale is
\[
\boxed{
n_{12}^{(1)}
\gtrsim
\frac{
(\mu_1+\sigma_1L_s)^2\sigma_2^2L_s^2
}{(\epsilon_{12,\mathrm{adm}}^{(1)})^2}
\log
\frac{C}{
\delta_{\mathrm{gate}}\epsilon_{12,\mathrm{adm}}^{(1)}
},
}
\tag{offdiag-sample-floor}
\]
up to universal constants and the harmless logarithmic dependence hidden in
solving for $n$.  This replaces the much looser common-envelope scaling
$X_{\delta_{\mathrm{strong}}}^4/(\epsilon_{12,\mathrm{adm}}^{(1)})^2$.

\end{enumerate}
\end{lemma}

\begin{proof}
For fixed $p,r,s$, consider the bounded function class
\[
\mathcal F_{p,rs}
=
\left\{
x\mapsto
\bar x_r^{(p)}\bar x_s^{(p)}
\mathbf 1\{
\widetilde u(\phi)^\top x>0
\}:
\phi\in\mathbb R
\right\}.
\]
Its deterministic envelope is $E_{rs}^{(p)}$.  The indicator component is
the class of homogeneous planar halfspaces, of universal VC dimension.
Multiplication by the fixed bounded function
$\bar x_r^{(p)}\bar x_s^{(p)}$ preserves a universal VC-type entropy bound.
The bounded VC maximal inequality therefore yields
\[
\sup_\phi
\left|
(\mathbb P_{p,n}-P_p)
\left[
\bar x_r^{(p)}\bar x_s^{(p)}
\mathbf 1\{
\widetilde u(\phi)^\top x>0
\}
\right]
\right|
\le
\varepsilon_{rs}^{(p)}
\]
simultaneously over $p,r,s$, except on an event of total probability at
most $\delta_{\mathrm{gate}}$.

On $\mathcal H_{\mathrm{strong}}(\delta_{\mathrm{strong}})$, the coordinate
envelope gives $|x_{k,q}^{(p)}|\le B_{p,q}$, so
$\bar x_{k,q}^{(p)}=x_{k,q}^{(p)}$ for every empirical sample. Hence the
empirical clipped moment is exactly $\widehat a_{rs}^{(p)}$, proving
\textnormal{(i)} without any population-tail correction.

Part~\textnormal{(ii)} follows by evaluating the operator-norm bound in the
direction $\widetilde u(\phi)$.

For \textnormal{(iii)}, if the gates at $\phi$ and $\phi'$ disagree on a
training sample $x$, then
\[
\left|
\widetilde u(\phi)^\top x
\right|
\le
\|\widetilde u(\phi)-\widetilde u(\phi')\|\,\|x\|
\le
2X_{\delta_{\mathrm{strong}}}\sin\frac\eta2.
\]
Thus every changed sample lies in $S_{\phi,\tau_\eta}$.  On the same event,
$|x_rx_s|\le E_{rs}^{(p)}$, and therefore
\[
\left|
\widehat a_{rs}^{(p)}(\phi)
-
\widehat a_{rs}^{(p)}(\phi')
\right|
\le
E_{rs}^{(p)}\widehat P_p(S_{\phi,\tau_\eta}).
\]
The uniform moving-slab conclusion of Lemma~\ref{lem:Hstrong} gives
\eqref{eq:gate-modulus}.

For \textnormal{(iv)}, general position implies that at any gate boundary
at most one empirical summand changes. Its absolute $(r,s)$ entry is at
most $E_{rs}^{(p)}/n$. Part~\textnormal{(v)} is immediate from
\textnormal{(i)}. Finally, \textnormal{(vi)} follows by solving
\eqref{eq:gate-entry-error} for $n$ and enlarging the universal constant to
absorb the implicit $\log n$ term.
\end{proof}

\begin{remark}[Why the off-diagonal entry is the binding gated-moment scale]
The strong inner-region fixed point is controlled by an off-diagonal to
diagonal ratio.  Accordingly, a single common operator-norm tolerance would
pay the large $X_{\delta_{\mathrm{strong}}}^2$ envelope for the very small
$(1,2)$ entry and can make the sample-size floor artificially dominant.
Lemma~\ref{lem:gated-moment-concentration} retains the coordinate structure:
on concept $1$ the $(1,2)$ envelope is only
\[
(\mu_1+\sigma_1L_s)\sigma_2L_s.
\]
A variance-sensitive VC--Bernstein refinement could further replace much of
the $L_s$ dependence by the natural scale $\mu_1\sigma_2$, but the bounded
entrywise form above is the theorem-level baseline used below.
\end{remark}


%%%%%%%%%%%%%%%%%%%%%%%%%%%%%%%%%%%%%%%%%%%%%%%%%%%%%%%%%%%%%%%%%%%%%%%%%%%%%%%
% T1.2: Clarke gradient-flow regularity
%%%%%%%%%%%%%%%%%%%%%%%%%%%%%%%%%%%%%%%%%%%%%%%%%%%%%%%%%%%%%%%%%%%%%%%%%%%%%%%

\begin{lemma}[Global Clarke gradient flow and measurable ReLU selectors]
\label{lem:gf-regularity}

Fix the finite two-concept training set
\[
\mathcal D
=
\left\{
x_k^{(p)}:
p\in\{1,2\},\ k\in[n]
\right\}
\subset\mathbb R^d,
\]
and consider the two-layer ReLU network
\[
f(\theta;x)
=
\sum_{i=1}^h
w_i\,\sigma(u_i^\top x),
\qquad
\sigma(z)=\max\{0,z\}.
\]
Write
\[
\rho_\theta(x)
:=
f(\theta;x)-x
\]
and
\[
\mathbb E_n[\varphi(x)]
:=
\frac1{2n}
\sum_{p=1}^2
\sum_{k=1}^n
\varphi(x_k^{(p)}).
\]
Thus
\[
\mathcal L(\theta)
=
\frac12
\mathbb E_n
\left[
\|\rho_\theta(x)\|^2
\right].
\]

Let $X_{\max}$ denote any deterministic envelope satisfying
\[
\|x\|\le X_{\max}
\qquad
\forall x\in\mathcal D.
\]

Then
\[
\mathcal L:\mathbb R^{2dh}\to\mathbb R
\]
is locally Lipschitz and finitely piecewise polynomial, hence piecewise
$C^\infty$.

For every initial condition $\theta(0)$, there exists an absolutely
continuous solution of
\[
\dot\theta(t)
\in
-\partial^\circ\mathcal L(\theta(t))
\qquad
\text{for a.e. }t\ge0.
\]
Every maximal absolutely continuous solution is global and satisfies the
growth bound in part~\textnormal{(iii)} below.

Fix one such global trajectory
\[
\theta:[0,\infty)\to\mathbb R^{2dh}
\]
once and for all. Then the following hold.

\begin{enumerate}[label=(\roman*),leftmargin=2.4em]

\item
\textbf{Measurable selected neuronwise equations.}

There exist Lebesgue-measurable functions
\[
\alpha_{i,k}^{(p)}:[0,\infty)\to[0,1],
\qquad
i\in[h],\quad
p\in\{1,2\},\quad
k\in[n],
\]
such that, writing
\[
\alpha_i(t;x_k^{(p)})
:=
\alpha_{i,k}^{(p)}(t),
\]
one has for every $t\ge0$
\[
\alpha_i(t;x_k^{(p)})
\in
\begin{cases}
\{1\},
&
u_i(t)^\top x_k^{(p)}>0,
\\[1mm]
[0,1],
&
u_i(t)^\top x_k^{(p)}=0,
\\[1mm]
\{0\},
&
u_i(t)^\top x_k^{(p)}<0.
\end{cases}
\]
For a zero training sample we fix
\[
\alpha_i(t;0)=0.
\]

For a.e. $t\ge0$ and every $i\in[h]$,
\[
\dot w_i(t)
=
-
\mathbb E_n
\left[
\rho_{\theta(t)}(x)
\sigma(u_i(t)^\top x)
\right],
\]
and
\[
\dot u_i(t)
=
-
\mathbb E_n
\left[
\langle
\rho_{\theta(t)}(x),
w_i(t)
\rangle
\alpha_i(t;x)x
\right].
\]
Moreover, for every $t\ge0$ and every training sample,
\[
\boxed{
\alpha_i(t;x)u_i(t)^\top x
=
\sigma(u_i(t)^\top x).
}
\]

Consequently,
\[
Q_i^\alpha(t)
:=
\mathbb E_n
\left[
xx^\top\alpha_i(t;x)
\right]
\]
is Lebesgue measurable and obeys the operator sandwich
\[
\boxed{
0
\preceq
Q_i^\alpha(t)
\preceq
\Sigma_n
}
\qquad
\forall t\ge0.
\]

\item
\textbf{Classical flow away from activation boundaries.}

Suppose that at some $t_\star\ge0$,
\[
u_i(t_\star)^\top x_k^{(p)}
\neq0
\]
for every $i,p,k$ with $x_k^{(p)}\neq0$. Then there exists a relative open
interval
\[
I\subset[0,\infty),
\qquad
t_\star\in I,
\]
on which all nontrivial training-sample preactivation signs remain fixed.

On $I$,
\[
\theta\in C^1(I),
\]
and the classical gradient-flow equations hold pointwise:
\[
\dot w_i(t)
=
-
\mathbb E_n
\left[
\rho_{\theta(t)}(x)
\sigma(u_i(t)^\top x)
\right],
\]
\[
\dot u_i(t)
=
-
\mathbb E_n
\left[
\langle
\rho_{\theta(t)}(x),
w_i(t)
\rangle
\mathbf 1\{u_i(t)^\top x>0\}x
\right].
\]

\item
\textbf{Universal growth bound and explicit finite-horizon Lipschitz
constant.}

The assertions in this part hold for every absolutely continuous Clarke
gradient-flow solution, not merely the trajectory fixed above.

For a.e. $t$ on the solution's interval of existence,
\[
\frac{d}{dt}\|\theta(t)\|^2
\le
X_{\max}^2.
\]
Consequently,
\[
\boxed{
\|\theta(t)\|^2
\le
\|\theta(0)\|^2
+
X_{\max}^2t.
}
\]

For a finite horizon $\mathsf T>0$, define
\[
\Theta_{\mathsf T}
:=
\sqrt{
\|\theta(0)\|^2
+
X_{\max}^2\mathsf T
}.
\]
Then
\[
\sup_{0\le t\le\mathsf T}
\|\theta(t)\|
\le
\Theta_{\mathsf T},
\]
and
\[
\|f(\theta(t);x)\|
\le
\frac12
X_{\max}
\Theta_{\mathsf T}^2,
\]
\[
\|\rho_{\theta(t)}(x)\|
\le
X_{\max}
\left(
1+\frac12\Theta_{\mathsf T}^2
\right)
\qquad
\forall x\in\mathcal D.
\]

Therefore, for a.e. $t\in[0,\mathsf T]$,
\[
\boxed{
\|\dot\theta(t)\|
\le
\Lambda_{\mathsf T}
:=
X_{\max}^2
\left(
1+\frac12\Theta_{\mathsf T}^2
\right)
\Theta_{\mathsf T}.
}
\]
In particular, the trajectory is
$\Lambda_{\mathsf T}$-Lipschitz on $[0,\mathsf T]$.

\item
\textbf{Absolute-continuity calculus and angular coordinates.}

Every parameter block
\[
u_i(\cdot),
\qquad
w_i(\cdot)
\]
is locally absolutely continuous.

If
\[
F:\mathbb R^{2dh}\to\mathbb R
\]
is $C^1$, then $F(\theta(\cdot))$ is locally absolutely continuous and
\[
\frac{d}{dt}F(\theta(t))
=
\left\langle
\nabla F(\theta(t)),
\dot\theta(t)
\right\rangle
\]
for a.e. $t$.

More generally, if $F$ is locally Lipschitz, then
\[
F(\theta(\cdot))
\]
is locally absolutely continuous.

Let
\[
\mathcal S
=
\operatorname{span}\{e_1,e_2\},
\qquad
v_i=P_{\mathcal S}u_i,
\qquad
r_i=\|v_i\|.
\]
Then $r_i(\cdot)$ is locally absolutely continuous.

On any interval $J$ on which
\[
r_i(t)\ge c>0,
\]
define
\[
\widetilde u_i(t)
:=
\frac{v_i(t)}{r_i(t)}.
\]
Then $\widetilde u_i$ is locally absolutely continuous on $J$, and one may
choose a continuous locally absolutely continuous angular lift $\phi_i$ such
that
\[
\widetilde u_i(t)
=
\cos\phi_i(t)e_1
+
\sin\phi_i(t)e_2.
\]
Set
\[
\widetilde u_i^\angle(t)
=
-\sin\phi_i(t)e_1
+
\cos\phi_i(t)e_2.
\]

All later radial and angular differential identities are understood as a.e.
identities on intervals where the corresponding projected radius is strictly
positive.

\item
\textbf{Selection-free consequences.}

The assertions in this part hold for every absolutely continuous Clarke
gradient-flow solution and every compatible selector representation, not
merely the trajectory fixed above.

For every neuron $i$,
\[
\boxed{
\frac{d}{dt}
\left(
\|w_i(t)\|^2-\|u_i(t)\|^2
\right)
=
0
}
\]
for a.e. $t$. Thus exact neuronwise balance is independent of the selector
values chosen at zero preactivations.

If all training samples lie in $\mathcal S$, then
\[
\boxed{
P_{\mathcal S^\perp}\dot u_i(t)=0
}
\]
for a.e. $t$, and therefore
\[
P_{\mathcal S^\perp}u_i(t)
\equiv
P_{\mathcal S^\perp}u_i(0).
\]

Assume all training samples lie in $\mathcal S$ and $r_i(t)>0$. Then
\[
u_i(t)^\top x
=
r_i(t)\widetilde u_i(t)^\top x,
\]
and therefore
\[
\boxed{
\alpha_i(t;x)
\widetilde u_i(t)^\top x
=
\bigl(
\widetilde u_i(t)^\top x
\bigr)_+.
}
\]
Consequently, projected radial identities in which $\alpha_i$ occurs only
through the product
\[
\alpha_i(t;x)
\widetilde u_i(t)^\top x
\]
are selection-free.

By contrast, projected angular identities contain
\[
\alpha_i(t;x)
\widetilde u_i^\angle(t)^\top x.
\]
When
\[
\widetilde u_i(t)^\top x=0,
\]
one may still have
\[
\widetilde u_i^\angle(t)^\top x\neq0.
\]
Thus angular dynamics retain genuine selector dependence on exact activation
boundaries.

\end{enumerate}
\end{lemma}

\begin{proof}\leavevmode

\paragraph{Step 1: finite piecewise-polynomial structure.}

Index the $2n$ training samples by $a$ and write
\[
x_a\in\mathcal D.
\]

If $x_a=0$, then, because the network has no biases,
\[
f(\theta;0)=0.
\]
Since the target is the identity, this sample contributes identically zero
loss for every parameter value and may be discarded from the activation-cell
decomposition.

For every remaining $x_a\neq0$ and every neuron $i$, the switching set
\[
\{\theta:u_i^\top x_a=0\}
\]
is a genuine parameter hyperplane. The finitely many such hyperplanes
partition parameter space into finitely many open activation cells.

On a fixed cell there are constants
\[
\tau_{ia}\in\{0,1\}
\]
such that
\[
\sigma(u_i^\top x_a)
=
\tau_{ia}u_i^\top x_a.
\]
Hence
\[
f(\theta;x_a)
=
\sum_{i=1}^h
\tau_{ia}w_i(u_i^\top x_a)
\]
is polynomial in the parameters on that cell, and therefore so is
$\mathcal L$.

The loss is continuous across activation boundaries. Since only finitely
many polynomial branches occur, their gradients are uniformly bounded on
every compact parameter set. Thus $\mathcal L$ is locally Lipschitz and
finitely piecewise polynomial.

\paragraph{Step 2: pointwise Clarke subgradients admit compatible ReLU
selector representations.}

Fix $\theta$. Let
\[
\mathcal N
:=
\bigcup_{i,a:\,x_a\neq0}
\left\{
\vartheta:
u_i^\top x_a=0
\right\}.
\]
This is a finite union of parameter hyperplanes and therefore has Lebesgue
measure zero.

By the standard characterization of the Clarke generalized gradient, the
limiting-gradient representation of
\[
\partial^\circ\mathcal L(\theta)
\]
is unchanged if one excludes an arbitrary Lebesgue-null set from the
approximating differentiability points; see Clarke~(1983),
Theorem~2.5.1. We may therefore restrict the approximating sequence to
points outside $\mathcal N$, hence to open activation cells.

There are finitely many activation patterns. Thus, from every convergent
sequence of such differentiability points, one may pass to a subsequence on
which the activation pattern
\[
\tau=(\tau_{ia})_{i,a}
\]
is constant. If $\mathcal L_\tau$ denotes the polynomial extension from this
cell, continuity of its polynomial gradient gives
\[
\lim_{m\to\infty}
\nabla\mathcal L(\theta_m)
=
\nabla\mathcal L_\tau(\theta).
\]

It follows that every
\[
\zeta\in\partial^\circ\mathcal L(\theta)
\]
admits a finite convex representation
\[
\zeta
=
\sum_{r=1}^m
\lambda_r
\nabla\mathcal L_{\tau^{(r)}}(\theta),
\qquad
\lambda_r\ge0,
\qquad
\sum_{r=1}^m\lambda_r=1.
\]

At $\theta$, all such branches have the same network value and residual.
Indeed, if
\[
u_i^\top x_a=0,
\]
then independently of $\tau_{ia}\in\{0,1\}$,
\[
\tau_{ia}u_i^\top x_a
=
0
=
\sigma(u_i^\top x_a).
\]

Ordinary differentiation of branch $\tau$ gives
\[
\nabla_{w_i}\mathcal L_\tau(\theta)
=
\mathbb E_n
\left[
\rho_\theta(x)
\tau_i(x)u_i^\top x
\right]
=
\mathbb E_n
\left[
\rho_\theta(x)
\sigma(u_i^\top x)
\right],
\]
whereas
\[
\nabla_{u_i}\mathcal L_\tau(\theta)
=
\mathbb E_n
\left[
\langle
\rho_\theta(x),
w_i
\rangle
\tau_i(x)x
\right].
\]

For $x_a\neq0$, define
\[
\alpha_i(x_a)
:=
\sum_{r=1}^m
\lambda_r\tau_{ia}^{(r)},
\]
and set
\[
\alpha_i(0):=0.
\]
Then
\[
\alpha_i(x_a)=1
\quad\text{if }u_i^\top x_a>0,
\]
\[
\alpha_i(x_a)=0
\quad\text{if }u_i^\top x_a<0,
\]
and
\[
\alpha_i(x_a)\in[0,1]
\quad\text{if }u_i^\top x_a=0.
\]

Taking the corresponding blocks of the convex combination gives
\[
\zeta_{w_i}
=
\mathbb E_n
\left[
\rho_\theta(x)
\sigma(u_i^\top x)
\right],
\]
\[
\zeta_{u_i}
=
\mathbb E_n
\left[
\langle
\rho_\theta(x),
w_i
\rangle
\alpha_i(x)x
\right].
\]
Furthermore,
\[
\boxed{
\alpha_i(x)u_i^\top x
=
\sigma(u_i^\top x)
}
\]
for every training sample.

\paragraph{Step 3: local existence, universal growth bound, and global
continuation.}

Since $\mathcal L$ is locally Lipschitz, the multifunction
\[
\theta
\longmapsto
-\partial^\circ\mathcal L(\theta)
\]
has nonempty compact convex values, is locally bounded, and is upper
semicontinuous. The standard existence theorem for upper-semicontinuous
differential inclusions with nonempty compact convex values therefore gives
a local absolutely continuous solution from every initial condition; see,
for example, Aubin--Cellina, Theorem~2.1.4.

Let $\theta(\cdot)$ be any such local solution. At a.e. time at which
$\dot\theta(t)$ exists and the differential inclusion holds, set
\[
\zeta(t):=-\dot\theta(t)
\in
\partial^\circ\mathcal L(\theta(t)).
\]
By Step~2, $\zeta(t)$ admits a compatible selector representation. Therefore
\[
\begin{aligned}
\langle\theta,\zeta\rangle
&=
\sum_i
\langle
w_i,\zeta_{w_i}
\rangle
+
\sum_i
\langle
u_i,\zeta_{u_i}
\rangle
\\
&=
\mathbb E_n
\left[
\sum_i
\langle\rho,w_i\rangle
\sigma(u_i^\top x)
\right]
\\
&\quad+
\mathbb E_n
\left[
\sum_i
\langle\rho,w_i\rangle
\alpha_i(x)u_i^\top x
\right].
\end{aligned}
\]
Since
\[
\alpha_i(x)u_i^\top x
=
\sigma(u_i^\top x),
\]
the terms agree, giving
\[
\boxed{
\langle\theta,\zeta\rangle
=
2\mathbb E_n
\left[
\langle
\rho_\theta(x),
f(\theta;x)
\rangle
\right].
}
\]

Because
\[
\rho_\theta(x)
=
f(\theta;x)-x,
\]
\[
\begin{aligned}
2\langle
\rho_\theta(x),
f(\theta;x)
\rangle
&=
2\|f(\theta;x)\|^2
-
2\langle x,f(\theta;x)\rangle
\\
&=
2\left\|
f(\theta;x)-\frac{x}{2}
\right\|^2
-
\frac12\|x\|^2
\\
&\ge
-\frac12X_{\max}^2.
\end{aligned}
\]
Hence
\[
\langle\theta,\zeta\rangle
\ge
-\frac12X_{\max}^2.
\]
Therefore
\[
\frac{d}{dt}\|\theta(t)\|^2
=
-2\langle\theta(t),\zeta(t)\rangle
\le
X_{\max}^2
\]
for a.e. $t$, and
\[
\|\theta(t)\|^2
\le
\|\theta(0)\|^2
+
X_{\max}^2t.
\]

This bound holds for every local Clarke solution. In particular, no maximal
solution can escape to infinity in finite time, so the local solution
continues globally.

\paragraph{Step 4: measurable compatible selectors along the fixed
trajectory.}

Now fix one global trajectory $\theta(\cdot)$.

There exists a Lebesgue-measurable set
\[
G\subset[0,\infty)
\]
of full measure such that, for every $t\in G$, $\dot\theta(t)$ exists and
\[
-\dot\theta(t)
\in
\partial^\circ\mathcal L(\theta(t)).
\]

For $t\in G$, collect the finitely many selector coordinates into
\[
\alpha
=
\bigl(
\alpha_i(x_k^{(p)})
\bigr)_{i,p,k}
\in
[0,1]^{2nh},
\]
and let
\[
\mathcal A(t)
\subset
[0,1]^{2nh}
\]
be the set of all $\alpha$ satisfying the ReLU sign constraints,
\[
\alpha_i(x_k^{(p)})=0
\qquad
\text{if }x_k^{(p)}=0,
\]
and the compatibility equations
\[
-\dot u_i(t)
=
\mathbb E_n
\left[
\langle
\rho_{\theta(t)}(x),
w_i(t)
\rangle
\alpha_i(x)x
\right],
\qquad
i\in[h].
\]

Step~2 gives
\[
\mathcal A(t)\neq\varnothing
\qquad
\forall t\in G.
\]
Every $\mathcal A(t)$ is compact and convex.

The trajectory is continuous and its a.e. derivative is Lebesgue measurable.
Hence the coefficients in the defining linear equalities and inequalities
are Lebesgue measurable in $t$. Consequently,
\[
\operatorname{gph}\mathcal A
=
\left\{
(t,\alpha):
t\in G,\ 
\alpha\in\mathcal A(t)
\right\}
\]
is measurable with respect to the completed Lebesgue sigma-algebra on $G$
and the Borel sigma-algebra on $[0,1]^{2nh}$.

Since the underlying Lebesgue sigma-algebra is complete and
$\mathcal A$ is closed-valued, Rockafellar--Wets,
Theorem~14.8, implies that
\[
t\longmapsto\mathcal A(t)
\]
is a measurable multifunction. Rockafellar--Wets,
Corollary~14.6, therefore yields a Lebesgue-measurable selector
\[
t\longmapsto\alpha(t)\in\mathcal A(t),
\qquad
t\in G.
\]

On the null complement $G^c$, define
\[
\alpha_i(t;x)
=
\mathbf 1
\{
u_i(t)^\top x>0
\}
\]
for $x\neq0$, and set
\[
\alpha_i(t;0)=0.
\]
The resulting selector is Lebesgue measurable on $[0,\infty)$ and obeys the
stated sign constraints at every time.

The neuronwise equations hold on $G$, hence for a.e. $t$, while
\[
\alpha_i(t;x)u_i(t)^\top x
=
\sigma(u_i(t)^\top x)
\]
holds for every $t$ and every training sample.

Finally,
\[
Q_i^\alpha(t)
=
\frac1{2n}
\sum_{p,k}
x_k^{(p)}
x_k^{(p)\top}
\alpha_i(t;x_k^{(p)})
\]
is Lebesgue measurable, and for every $v\in\mathbb R^d$,
\[
v^\top Q_i^\alpha(t)v
=
\mathbb E_n
\left[
(v^\top x)^2
\alpha_i(t;x)
\right]
\ge0.
\]
Thus
\[
Q_i^\alpha(t)\succeq0.
\]
Moreover, since
\[
0\le\alpha_i(t;x)\le1,
\]
for every $v\in\mathbb R^d$,
\[
\begin{aligned}
v^\top
\left(
\Sigma_n-Q_i^\alpha(t)
\right)v
&=
\mathbb E_n
\left[
(v^\top x)^2
\left(
1-\alpha_i(t;x)
\right)
\right]
\\
&\ge0.
\end{aligned}
\]
Hence
\[
\boxed{
0
\preceq
Q_i^\alpha(t)
\preceq
\Sigma_n.
}
\]

\paragraph{Step 5: classical flow away from activation boundaries.}

Suppose all nontrivial training-sample preactivations are nonzero at
$t_\star$. Since the trajectory is continuous and only finitely many
preactivations occur, all their signs remain fixed on some relative open
interval
\[
I\ni t_\star.
\]

On the corresponding activation region, $\mathcal L$ is smooth. Therefore
\[
\dot\theta(t)
=
-\nabla\mathcal L(\theta(t))
\]
for a.e. $t\in I$.

By absolute continuity,
\[
\theta(t)
=
\theta(t_0)
-
\int_{t_0}^t
\nabla\mathcal L(\theta(s))\,ds,
\qquad
t,t_0\in I.
\]
The integrand is continuous on $I$. Hence
\[
\theta\in C^1(I),
\]
and
\[
\dot\theta(t)
=
-\nabla\mathcal L(\theta(t))
\]
holds pointwise throughout $I$.

\paragraph{Step 6: explicit finite-horizon Lipschitz bound.}

Fix $\mathsf T>0$ and set
\[
\Theta_{\mathsf T}^2
=
\|\theta(0)\|^2
+
X_{\max}^2\mathsf T.
\]
By Step~3,
\[
\|\theta(t)\|
\le
\Theta_{\mathsf T}
\qquad
0\le t\le\mathsf T.
\]

For every training sample,
\[
\begin{aligned}
\|f(\theta;x)\|
&\le
\sum_i
\|w_i\|
|u_i^\top x|
\\
&\le
X_{\max}
\sum_i
\|w_i\|\|u_i\|
\\
&\le
\frac{X_{\max}}2
\sum_i
\left(
\|w_i\|^2+\|u_i\|^2
\right)
\\
&=
\frac12
X_{\max}
\|\theta\|^2.
\end{aligned}
\]
Thus
\[
\|\rho_\theta(x)\|
\le
X_{\max}
\left(
1+\frac12\|\theta\|^2
\right)
\le
R_{\mathsf T},
\]
where
\[
R_{\mathsf T}
:=
X_{\max}
\left(
1+\frac12\Theta_{\mathsf T}^2
\right).
\]

Using the neuronwise equations and $0\le\alpha_i\le1$,
\[
\|\dot w_i(t)\|
\le
R_{\mathsf T}
X_{\max}
\|u_i(t)\|,
\]
and
\[
\|\dot u_i(t)\|
\le
R_{\mathsf T}
X_{\max}
\|w_i(t)\|.
\]
Hence
\[
\begin{aligned}
\|\dot\theta(t)\|^2
&=
\sum_i
\left(
\|\dot w_i(t)\|^2
+
\|\dot u_i(t)\|^2
\right)
\\
&\le
R_{\mathsf T}^2
X_{\max}^2
\|\theta(t)\|^2
\\
&\le
R_{\mathsf T}^2
X_{\max}^2
\Theta_{\mathsf T}^2.
\end{aligned}
\]
Therefore
\[
\|\dot\theta(t)\|
\le
X_{\max}^2
\left(
1+\frac12\Theta_{\mathsf T}^2
\right)
\Theta_{\mathsf T}
=
\Lambda_{\mathsf T}
\]
for a.e. $t\in[0,\mathsf T]$.

\paragraph{Step 7: absolute-continuity calculus.}

Absolute continuity of the parameter blocks follows from absolute continuity
of $\theta$.

For $F\in C^1$, the ordinary chain rule for absolutely continuous curves
gives
\[
\frac{d}{dt}F(\theta(t))
=
\left\langle
\nabla F(\theta(t)),
\dot\theta(t)
\right\rangle
\]
for a.e. $t$.

If $F$ is locally Lipschitz, then its restriction to a compact neighborhood
of any finite-horizon trajectory is Lipschitz. Hence
\[
F\circ\theta
\]
is locally absolutely continuous.

In particular,
\[
r_i(t)
=
\|P_{\mathcal S}u_i(t)\|
\]
is locally absolutely continuous. On an interval on which
\[
r_i\ge c>0,
\]
normalization is locally Lipschitz, so
\[
\widetilde u_i
=
\frac{P_{\mathcal S}u_i}{r_i}
\]
is locally absolutely continuous. Since $\widetilde u_i$ takes values in
$\mathbb S^1$, it admits a continuous locally absolutely continuous angular
lift $\phi_i$.

\paragraph{Step 8: selection-free identities.}

Let $\theta(\cdot)$ be any absolutely continuous Clarke solution and choose
any compatible selector representation at a.e. differentiability time.

Then
\[
\begin{aligned}
\frac12
\frac{d}{dt}
\left(
\|w_i\|^2-\|u_i\|^2
\right)
&=
-
\mathbb E_n
\left[
\langle\rho,w_i\rangle
\sigma(u_i^\top x)
\right]
\\
&\quad+
\mathbb E_n
\left[
\langle\rho,w_i\rangle
\alpha_i(t;x)
u_i^\top x
\right]
\\
&=0.
\end{aligned}
\]
Thus balance is universal and selection-free.

If all data lie in $\mathcal S$, then
\[
P_{\mathcal S^\perp}\dot u_i
=
-
\mathbb E_n
\left[
\langle\rho,w_i\rangle
\alpha_i(t;x)
P_{\mathcal S^\perp}x
\right]
=
0.
\]

Finally, if $r_i>0$ and $x\in\mathcal S$, then
\[
u_i^\top x
=
r_i
\widetilde u_i^\top x,
\]
so
\[
\alpha_i(t;x)
\widetilde u_i^\top x
=
\frac1{r_i}
\sigma(u_i^\top x)
=
(\widetilde u_i^\top x)_+.
\]

No corresponding cancellation eliminates $\alpha_i$ from
\[
\alpha_i(t;x)
(\widetilde u_i^\angle)^\top x
\]
on an exact activation boundary.

This completes the proof.
\end{proof}


%%%%%%%%%%%%%%%%%%%%%%%%%%%%%%%%%%%%%%%%%%%%%%%%%%%%%%%%%%%%%%%%%%%%%%%%%%%%%%%
% Flow-selection convention
%%%%%%%%%%%%%%%%%%%%%%%%%%%%%%%%%%%%%%%%%%%%%%%%%%%%%%%%%%%%%%%%%%%%%%%%%%%%%%%

\begin{remark}[Flow and selector convention]
\label{rem:gf-selector-convention}

The primary nonsmooth object fixed throughout the paper is one global
absolutely continuous Clarke trajectory
\[
\theta(\cdot).
\]
On the full-measure set where its derivative exists and the differential
inclusion holds,
\[
-\dot\theta(t)
\in
\partial^\circ\mathcal L(\theta(t))
\]
is determined by the trajectory.

Lemma~\ref{lem:gf-regularity} fixes one measurable compatible ReLU-selector
representation
\[
\alpha_i(t;x)
\]
of this a.e. velocity. At simultaneous zero preactivations the representation
need not be unique. All later quantities involving $\alpha_i$, in particular
\[
Q_i^\alpha(t),
\]
refer to this same fixed measurable representation.
\end{remark}


%%%%%%%%%%%%%%%%%%%%%%%%%%%%%%%%%%%%%%%%%%%%%%%%%%%%%%%%%%%%%%%%%%%%%%%%%%%%%%%
% Boundary correction for later angular dynamics
%%%%%%%%%%%%%%%%%%%%%%%%%%%%%%%%%%%%%%%%%%%%%%%%%%%%%%%%%%%%%%%%%%%%%%%%%%%%%%%

\begin{remark}[Boundary contribution in projected angular dynamics]
\label{rem:angular-boundary-selection}

Work on the data event $\mathcal H_{\mathrm{data}}(\delta)$ of
Definition~\ref{def:Hdata}. Since
\[
\mathcal H_{\mathrm{data}}(\delta)
\subset
\mathcal G_{\mathrm{pos}},
\]
the finite-sample general-position conclusion of
Lemma~\ref{lem:general-position} is available throughout.

Assume all training samples lie in $\mathcal S$, and consider an interval on
which
\[
r_i(t)\ge c>0.
\]
Then Lemma~\ref{lem:gf-regularity}\textnormal{(iv)} gives, for a.e. $t$,
\[
\boxed{
\dot\phi_i(t)
=
-\frac1{r_i(t)}
\mathbb E_n
\left[
\langle
\rho_{\theta(t)}(x),
w_i(t)
\rangle
\mathbf 1
\{
\widetilde u_i(t)^\top x>0
\}
\widetilde u_i^\angle(t)^\top x
\right]
+
E_i^{\mathrm{bd}}(t).
}
\]

Here
\[
\Delta_i(t;x)
:=
\alpha_i(t;x)
-
\mathbf 1
\{
\widetilde u_i(t)^\top x>0
\},
\]
and
\[
E_i^{\mathrm{bd}}(t)
:=
-\frac1{r_i(t)}
\mathbb E_n
\left[
\langle
\rho_{\theta(t)}(x),
w_i(t)
\rangle
\Delta_i(t;x)
\widetilde u_i^\angle(t)^\top x
\right].
\]
The quantity $\Delta_i(t;x)$ vanishes unless
\[
\widetilde u_i(t)^\top x=0.
\]

Suppose additionally that
\[
s_i(t)
:=
\|P_{\mathcal S}w_i(t)\|
>0,
\]
and write
\[
z_i(t)
:=
P_{\mathcal S}w_i(t)
=
s_i(t)\widetilde w_i(t),
\qquad
w_{i,\perp}(t)
:=
P_{\mathcal S^\perp}w_i(t).
\]
Since
\[
\langle\rho,w_i\rangle
=
s_i
\langle\rho,\widetilde w_i\rangle
+
\left\langle
P_{\mathcal S^\perp}\rho,
w_{i,\perp}
\right\rangle,
\]
decompose
\[
E_i^{\mathrm{bd}}
=
E_{i,\parallel}^{\mathrm{bd}}
+
E_{i,\perp}^{\mathrm{bd}},
\]
where
\[
E_{i,\parallel}^{\mathrm{bd}}
=
-\frac{s_i}{r_i}
\mathbb E_n
\left[
\langle
\rho,
\widetilde w_i
\rangle
\Delta_i
(\widetilde u_i^\angle)^\top x
\right],
\]
and
\[
E_{i,\perp}^{\mathrm{bd}}
=
-\frac1{r_i}
\mathbb E_n
\left[
\left\langle
P_{\mathcal S^\perp}\rho,
w_{i,\perp}
\right\rangle
\Delta_i
(\widetilde u_i^\angle)^\top x
\right].
\]

By Lemma~\ref{lem:general-position}, at most one of the $2n$ training
samples lies on the relevant activation boundary. Moreover, on
$\mathcal H_{\mathrm{data}}(\delta)$,
\[
\|x\|
\le
X_{\delta_{\mathrm{strong}}}
\qquad
\forall x\in\mathcal D.
\]
Hence
\[
|\Delta_i(t;x)|\le1,
\qquad
|\widetilde u_i^\angle(t)^\top x|
\le
X_{\delta_{\mathrm{strong}}}.
\]
Therefore
\[
\boxed{
|E_{i,\parallel}^{\mathrm{bd}}(t)|
\le
\frac{s_i(t)}{r_i(t)}
\frac{X_{\delta_{\mathrm{strong}}}}{2n}
\max_{x\in\mathcal D}
\left|
\left\langle
\rho_{\theta(t)}(x),
\widetilde w_i(t)
\right\rangle
\right|.
}
\]

Thus the in-span boundary correction carries exactly the same
\[
\frac{s_i}{r_i}
\]
prefactor as the principal projected angular field. In particular, no
$\|w_i\|/r_i$ factor and hence no spurious ambient-dimension factor is
introduced.

The perpendicular boundary contribution satisfies
\[
\boxed{
|E_{i,\perp}^{\mathrm{bd}}(t)|
\le
\frac{X_{\delta_{\mathrm{strong}}}}{2n\,r_i(t)}
\max_{x\in\mathcal D}
\left|
\left\langle
P_{\mathcal S^\perp}\rho_{\theta(t)}(x),
P_{\mathcal S^\perp}w_i(t)
\right\rangle
\right|.
}
\]

This is a $1/(2n)$ boundary contribution to the same perpendicular remainder
that must already be controlled in the projected angular dynamics.
\end{remark}


%%%%%%%%%%%%%%%%%%%%%%%%%%%%%%%%%%%%%%%%%%%%%%%%%%%%%%%%%%%%%%%%%%%%%%%%%%%%%%%
% T1.3: exact balance -- corollary of T1.2
%%%%%%%%%%%%%%%%%%%%%%%%%%%%%%%%%%%%%%%%%%%%%%%%%%%%%%%%%%%%%%%%%%%%%%%%%%%%%%%

\begin{corollary}[Exact neuronwise balance]
\label{cor:exact-balance}

Suppose the initialization is neuronwise balanced:
\[
\|w_i(0)\|
=
\|u_i(0)\|
\qquad
\forall i\in[h].
\]
Then every absolutely continuous Clarke gradient-flow solution satisfies
\[
\boxed{
\|w_i(t)\|
=
\|u_i(t)\|
}
\qquad
\forall i\in[h],
\quad
\forall t\ge0.
\]

In particular, under the aligned initialization
\[
w_i(0)
=
u_i(0)
=
\varepsilon\widehat u_i,
\]
exact balance holds for every neuron for all time.
\end{corollary}

\begin{proof}
By Lemma~\ref{lem:gf-regularity}\textnormal{(v)},
\[
\|w_i(t)\|^2-\|u_i(t)\|^2
=
\|w_i(0)\|^2-\|u_i(0)\|^2
=
0.
\]
Since both norms are nonnegative, the conclusion follows.
\end{proof}


%%%%%%%%%%%%%%%%%%%%%%%%%%%%%%%%%%%%%%%%%%%%%%%%%%%%%%%%%%%%%%%%%%%%%%%%%%%%%%%
% T1.4: perpendicular dynamics -- corollary of T1.2
%%%%%%%%%%%%%%%%%%%%%%%%%%%%%%%%%%%%%%%%%%%%%%%%%%%%%%%%%%%%%%%%%%%%%%%%%%%%%%%

\begin{corollary}[Perpendicular invariance and stacked output contraction]
\label{cor:perp-dynamics}

Assume that every training sample lies in
\[
\mathcal S
=
\operatorname{span}\{e_1,e_2\}.
\]
Then, along every absolutely continuous Clarke gradient-flow solution, the
following statements hold independently of the particular compatible ReLU
selector representation.

\begin{enumerate}[label=(\roman*),leftmargin=2.4em]

\item
For every neuron $i$,
\[
\boxed{
P_{\mathcal S^\perp}u_i(t)
=
P_{\mathcal S^\perp}u_i(0)
}
\qquad
\forall t\ge0.
\]

\item
Write
\[
w_{i,\perp}(t)
:=
P_{\mathcal S^\perp}w_i(t)
\]
and stack the perpendicular output components as
\[
W_\perp(t)
:=
\begin{bmatrix}
w_{1,\perp}(t)
&
\cdots
&
w_{h,\perp}(t)
\end{bmatrix}
\in\mathbb R^{d\times h}.
\]
Define
\[
G(t)
=
\bigl(
G_{ij}(t)
\bigr)_{i,j=1}^h,
\]
where
\[
G_{ij}(t)
:=
\mathbb E_n
\left[
\sigma(u_i(t)^\top x)
\sigma(u_j(t)^\top x)
\right].
\]
Then
\[
\boxed{
G(t)\succeq0
}
\qquad
\forall t\ge0,
\]
and, for a.e. $t\ge0$,
\[
\boxed{
\dot W_\perp(t)
=
-W_\perp(t)G(t).
}
\]

Consequently,
\[
\boxed{
\frac{d}{dt}
\|W_\perp(t)\|_F^2
=
-2
\operatorname{tr}
\left(
G(t)
W_\perp(t)^\top W_\perp(t)
\right)
\le0
}
\]
for a.e. $t$, and hence
\[
\boxed{
\|W_\perp(t)\|_F
\le
\|W_\perp(0)\|_F
}
\qquad
\forall t\ge0.
\]

Under the aligned initialization
\[
w_i(0)
=
\varepsilon\widehat u_i,
\qquad
\|\widehat u_i\|=1,
\]
this yields
\[
\boxed{
\|W_\perp(t)\|_F
\le
\|W_\perp(0)\|_F
\le
\varepsilon\sqrt h
}
\qquad
\forall t\ge0.
\]

\end{enumerate}

The output-side contraction is a \emph{stacked Frobenius-norm statement}.
It does not imply that any individual quantity
\[
\|w_{i,\perp}(t)\|
\]
is nonincreasing in time.
\end{corollary}

\begin{proof}
Part~\textnormal{(i)} is exactly
Lemma~\ref{lem:gf-regularity}\textnormal{(v)}.

For part~\textnormal{(ii)}, the $w_i$ equation from
Lemma~\ref{lem:gf-regularity}\textnormal{(i)} contains no Clarke selector:
\[
\dot w_i
=
-
\mathbb E_n
\left[
\rho(x)
\sigma(u_i^\top x)
\right].
\]
Thus the following calculation is selection-free.

Since every training sample lies in $\mathcal S$,
\[
P_{\mathcal S^\perp}\rho(x)
=
P_{\mathcal S^\perp}f(x)
=
\sum_{j=1}^h
\sigma(u_j^\top x)
w_{j,\perp}.
\]
Therefore, for a.e. $t$,
\[
\begin{aligned}
\dot w_{i,\perp}
&=
-
\mathbb E_n
\left[
P_{\mathcal S^\perp}\rho(x)
\sigma(u_i^\top x)
\right]
\\
&=
-
\sum_{j=1}^h
\mathbb E_n
\left[
\sigma(u_i^\top x)
\sigma(u_j^\top x)
\right]
w_{j,\perp}
\\
&=
-
\sum_{j=1}^h
G_{ij}w_{j,\perp}.
\end{aligned}
\]
Since $G=G^\top$, this is equivalently
\[
\dot W_\perp
=
-W_\perp G.
\]

For every $a\in\mathbb R^h$,
\[
a^\top G a
=
\mathbb E_n
\left[
\left(
\sum_{i=1}^h
a_i
\sigma(u_i^\top x)
\right)^2
\right]
\ge0.
\]
Hence
\[
G\succeq0.
\]

Consequently,
\[
\begin{aligned}
\frac12
\frac{d}{dt}
\|W_\perp\|_F^2
&=
\operatorname{tr}
\left(
W_\perp^\top
\dot W_\perp
\right)
\\
&=
-
\operatorname{tr}
\left(
W_\perp^\top
W_\perp
G
\right)
\\
&=
-
\operatorname{tr}
\left(
G
W_\perp^\top
W_\perp
\right)
\\
&\le0,
\end{aligned}
\]
because
\[
G\succeq0,
\qquad
W_\perp^\top W_\perp\succeq0.
\]
Absolute continuity gives
\[
\|W_\perp(t)\|_F
\le
\|W_\perp(0)\|_F.
\]

Finally,
\[
\begin{aligned}
\|W_\perp(0)\|_F^2
&=
\sum_{i=1}^h
\|P_{\mathcal S^\perp}w_i(0)\|^2
\\
&\le
\sum_{i=1}^h
\|w_i(0)\|^2
\\
&=
h\varepsilon^2,
\end{aligned}
\]
so
\[
\|W_\perp(0)\|_F
\le
\varepsilon\sqrt h.
\]

No corresponding per-neuron monotonicity follows from the trace inequality,
because the dynamics couple the columns of $W_\perp$ through the generally
nondiagonal Gram matrix $G(t)$.
\end{proof}

Fix once and for all a universal perturbative-output constant
\[
\boxed{
c_f\in(0,1).
}
\]
Its numerical value will be chosen sufficiently small to satisfy the later
angular-field and locking inequalities.

%%%%%%%%%%%%%%%%%%%%%%%%%%%%%%%%%%%%%%%%%%%%%%%%%%%%%%%%%%%%%%%%%%%%%%%%%%%%%%%
% T1.5: general-horizon spectral norm/output control
%%%%%%%%%%%%%%%%%%%%%%%%%%%%%%%%%%%%%%%%%%%%%%%%%%%%%%%%%%%%%%%%%%%%%%%%%%%%%%%

\begin{lemma}[General-horizon spectral norm and output bounds]
\label{lem:norm-output}

Work on the data event $\mathcal H_{\mathrm{data}}(\delta)$ of
Definition~\ref{def:Hdata}, and assume the aligned-balanced initialization
\[
w_i(0)=u_i(0)=\varepsilon\widehat u_i,
\qquad
\|\widehat u_i\|=1.
\]
For compactness write
\[
X:=X_{\delta_{\mathrm{strong}}},
\qquad
\lambda:=\overline\lambda_\delta.
\]
For $t\ge0$ define
\[
\boxed{
\chi(t)
:=
\frac{X^2h\varepsilon^2}{\lambda}
\left(e^{2\lambda t}-1\right).
}
\]

Fix a deterministic horizon $H>0$ such that
\[
\boxed{\chi(H)<1.}
\]
Then every absolutely continuous Clarke gradient-flow solution satisfies,
for every $0\le t\le H$,
\[
\boxed{
S(t)
:=
\sum_{j=1}^h\|u_j(t)\|^2
\le
\overline S(t)
:=
\frac{h\varepsilon^2e^{2\lambda t}}{1-\chi(t)}
\le
\frac{h\varepsilon^2e^{2\lambda t}}{1-\chi(H)}.
}
\]
Moreover, for every $i\in[h]$,
\[
\boxed{
\|u_i(t)\|
=
\|w_i(t)\|
\le
\frac{\varepsilon e^{\lambda t}}{\sqrt{1-\chi(t)}}
\le
\frac{\varepsilon e^{\lambda t}}{\sqrt{1-\chi(H)}}.
}
\]

The network output obeys
\[
\boxed{
F(t)
:=
\sup_{x\in\mathcal D}\|f_t(x)\|
\le
X\overline S(t)
\le
\frac{Xh\varepsilon^2e^{2\lambda t}}{1-\chi(H)}.
}
\]
Consequently,
\[
\boxed{
\sup_{x\in\mathcal D}\|\rho_t(x)\|
\le
X+F(t).
}
\]
In addition,
\[
\boxed{
2X\int_0^tF(s)\,ds
\le
\log\frac1{1-\chi(t)}
\le
\log\frac1{1-\chi(H)}.
}
\]

Now fix any $\overline\chi\in(0,1)$ and suppose
\[
\boxed{\chi(H)\le\overline\chi.}
\]
Then the entire interval $[0,H]$ admits the horizon-level output ceiling
\[
\boxed{
F(t)
\le
F_H^\star
:=
\frac{\lambda}{X}
\frac{\chi(H)}{1-\chi(H)}
\frac1{1-e^{-2\lambda H}}
\le
\frac{\lambda}{X}
\frac{\overline\chi}{1-\overline\chi}
\frac1{1-e^{-2\lambda H}}.
}
\]
Equivalently, $\chi(H)\le\overline\chi$ is exactly
\[
\boxed{
H
\le
H_{\overline\chi}
:=
\frac1{2\lambda}
\log\left(
1+\frac{\overline\chi\lambda}{X^2h\varepsilon^2}
\right).
}
\]
In particular, the Riccati comparison remains finite precisely for
\[
\boxed{
H
<
T_{\mathrm{blow}}
:=
\frac1{2\lambda}
\log\left(
1+\frac{\lambda}{X^2h\varepsilon^2}
\right).
}
\]
Finally, if
\[
F_H^\star\le c_f\mu_{\min},
\]
then
\[
\boxed{
\sup_{x\in\mathcal D}\|f_t(x)\|
\le
c_f\mu_p,
\qquad p=1,2,
\qquad 0\le t\le H.
}
\]
\end{lemma}

\begin{proof}
By Corollary~\ref{cor:exact-balance},
\[
\|w_i(t)\|=\|u_i(t)\|
\qquad
\forall i,\ t\ge0.
\]
The argument used in the original early-time proof gives, for a.e. $t$,
\[
\frac{d}{dt}\|u_i(t)\|^2
\le
2\left(\lambda+X^2S(t)\right)\|u_i(t)\|^2.
\]
Summing over $i$ yields
\[
\boxed{
\dot S(t)
\le
2\lambda S(t)+2X^2S(t)^2,
\qquad
S(0)=h\varepsilon^2.
}
\]
Let $\overline S$ solve the corresponding equality.  Its explicit solution is
\[
\overline S(t)
=
\frac{h\varepsilon^2e^{2\lambda t}}
{1-\frac{X^2h\varepsilon^2}{\lambda}(e^{2\lambda t}-1)}
=
\frac{h\varepsilon^2e^{2\lambda t}}{1-\chi(t)}.
\]
Because $\chi$ is increasing and $\chi(H)<1$, scalar comparison gives the
stated bound on $S$ throughout $[0,H]$.

For $y_i(t):=\|u_i(t)\|^2$, the same linear comparison against
$m(t):=\overline S(t)/h$ yields
\[
y_i(t)
\le
\frac{\overline S(t)}h
=
\frac{\varepsilon^2e^{2\lambda t}}{1-\chi(t)}.
\]
Taking square roots and using exact balance proves the neuronwise norm bound.

Next,
\[
\|f_t(x)\|
\le
\sum_{j=1}^h\|w_j(t)\|\,|u_j(t)^\top x|
\le
X\sum_{j=1}^h\|w_j(t)\|\|u_j(t)\|
=
XS(t),
\]
which proves the output estimate and hence the residual estimate.

For the integrated estimate set $D(t):=1-\chi(t)$.  Then
\[
D'(t)=-2X^2h\varepsilon^2e^{2\lambda t},
\qquad
\overline S(t)=-\frac{D'(t)}{2X^2D(t)}.
\]
Therefore
\[
2X\int_0^tF(s)\,ds
\le
2X^2\int_0^t\overline S(s)\,ds
=
-\int_0^t\frac{D'(s)}{D(s)}\,ds
=
\log\frac1{1-\chi(t)}.
\]

Finally,
\[
h\varepsilon^2
=
\frac{\lambda\chi(H)}{X^2(e^{2\lambda H}-1)},
\]
so for $t\le H$,
\[
F(t)
\le
\frac{\lambda}{X}
\frac{\chi(H)}{1-\chi(H)}
\frac{e^{2\lambda t}}{e^{2\lambda H}-1}
\le
\frac{\lambda}{X}
\frac{\chi(H)}{1-\chi(H)}
\frac1{1-e^{-2\lambda H}}.
\]
The remaining claims follow algebraically from $\chi(H)\le\overline\chi$.
\end{proof}

\begin{corollary}[Original spectral early-time specialization]
\label{cor:norm-output-early}
Assume
\[
\sqrt h\,\varepsilon<1
\]
and define
\[
T_{\mathrm{early}}
:=
\frac1{2\lambda}\log\frac1{\sqrt h\,\varepsilon},
\qquad
\chi_\varepsilon
:=
\frac{X^2\sqrt h\,\varepsilon}{\lambda}.
\]
Then
\[
\chi(T_{\mathrm{early}})
=
\frac{X^2}{\lambda}
\left(\sqrt h\,\varepsilon-h\varepsilon^2\right)
<
\chi_\varepsilon.
\]
Consequently, if $\chi_\varepsilon\le1/2$, all conclusions of
Lemma~\ref{lem:norm-output} hold on $[0,T_{\mathrm{early}}]$, and in
particular
\[
S(t)
\le
\frac{\sqrt h\,\varepsilon}{1-\chi_\varepsilon},
\]
\[
\|u_i(t)\|
=
\|w_i(t)\|
\le
\frac{\varepsilon^{1/2}h^{-1/4}}{\sqrt{1-\chi_\varepsilon}},
\]
and
\[
F(t)
\le
\frac{X\sqrt h\,\varepsilon}{1-\chi_\varepsilon}.
\]
\end{corollary}

%%%%%%%%%%%%%%%%%%%%%%%%%%%%%%%%%%%%%%%%%%%%%%%%%%%%%%%%%%%%%%%%%%%%%%%%%%%%%%%
% Core two-concept admissible region
%%%%%%%%%%%%%%%%%%%%%%%%%%%%%%%%%%%%%%%%%%%%%%%%%%%%%%%%%%%%%%%%%%%%%%%%%%%%%%%

\begin{definition}[Core two-concept admissible region]
\label{def:R2c-core}

Define
\[
\mathfrak R_{\mathrm{2c}}^{\mathrm{core}}
\]
to be the set of parameters
\[
\boxed{
(\mu_1,\mu_2,\sigma_1,\sigma_2,d,h,n,\varepsilon,\delta)
}
\]
satisfying the data-event SNR condition
\[
\boxed{
\mu_{\min}
\ge
2\sigma_{\max}
\sqrt{2\log\frac{96n}{\delta}}.
}
\]
Equivalently,
\[
\boxed{
n
\le
n_{\max}^{\mathrm{SNR}}
:=
\frac{\delta}{96}
\exp\left(\frac{\mu_{\min}^2}{8\sigma_{\max}^2}\right).
}
\]
\end{definition}

The final admissible region $\mathfrak R_{\mathrm{2c}}$ is obtained by
intersecting $\mathfrak R_{\mathrm{2c}}^{\mathrm{core}}$ with the explicit
dynamical and geometric conditions introduced below.

In particular, the final region additionally contains a lower sample-size
bound
\[
n\ge n_{\min},
\qquad
n_{\min}\ge n_{\mathrm{gate}},
\]
where $n_{\mathrm{gate}}$ is the entrywise gated-moment floor from
Lemma~\ref{lem:gated-moment-concentration}.  The $(1,2)$ gated entry is a
serious candidate for the binding lower sample-size condition because it
sets the inner strong-branch fixed point.  The full sample-size requirement
takes the explicit two-sided form
\[
\boxed{
n_{\min}\le n\le n_{\max}^{\mathrm{SNR}}.
}
\]
The final region must therefore be nonempty after intersecting the lower
concentration floor with the upper coordinate-envelope SNR ceiling.

The projected-trapping stage requires the existence of a deterministic
$O(1)$ trapping time $t_{\mathrm{trap}}=T$, a common spectral-envelope
horizon $H\ge T$, and a fixed $\overline\chi\in(0,1)$ such that
\[
\boxed{
\chi(H)
=
\frac{X_{\delta_{\mathrm{strong}}}^2h\varepsilon^2}
{\overline\lambda_\delta}
\left(e^{2\overline\lambda_\delta H}-1\right)
\le
\overline\chi,
}
\]
together with the radial, angular-gap, strong-cone, and weak-retention
inequalities of Lemma~\ref{lem:early-projected-trapping}.  The subsequent
post-trapping locking stage uses the same common horizon through
$t_{\mathrm{lock}}$.  Initialization smallness therefore enters through the
existence of these finite admissible horizons rather than through the
auxiliary scale $T_{\mathrm{early}}$.

The final region also contains the seed-arc width condition from
Lemma~\ref{lem:init-angular-coverage},
\[
\boxed{
h
\ge
\frac{32\pi}{\ell_{\mathrm{seed}}}
\log\frac{2}{\delta_{\mathrm{init}}},
}
\]
and the later dynamical radial-survival and locking constraints, including
the locking-stage lower bound on $n$ induced by the final angular scale
$\eta_1^u$.  The rebuilt strong-progress package below adds the pure-sector
feasibility condition
\[
b_1b_2<m_1^{\mathrm{data}}m_2^{\mathrm{data}},
\]
together with the forced inner-cone/discriminant and jump-floor conditions of
Proposition~\ref{sbnew:inner}, with strict mismatch/perpendicular forcing
room.  The later basin-creation stage uses a separate post-B1 progress
ceiling $m_{\max}$ and the regulator first-hitting interface in
Proposition~\ref{sbgen:first-hit}, followed by
Proposition~\ref{prop:gate-cliff-creation}; no load-bearing static-fold or
$\Delta_{\mathrm{branch}}$ margin is retained.


%%%%%%%%%%%%%%%%%%%%%%%%%%%%%%%%%%%%%%%%%%%%%%%%%%%%%%%%%%%%%%%%%%%%%%%%%%%%%%%
% T1.6: dimension-free angular initialization coverage
%%%%%%%%%%%%%%%%%%%%%%%%%%%%%%%%%%%%%%%%%%%%%%%%%%%%%%%%%%%%%%%%%%%%%%%%%%%%%%%

\begin{lemma}[Dimension-free angular initialization coverage]
\label{lem:init-angular-coverage}

Assume
\[
d\ge2
\]
and the aligned isotropic initialization
\[
w_i(0)
=
u_i(0)
=
\varepsilon\widehat u_i,
\qquad
\widehat u_i
\stackrel{\mathrm{i.i.d.}}{\sim}
\operatorname{Unif}(\mathbb S^{d-1}).
\]
Let
\[
\mathcal S
=
\operatorname{span}\{e_1,e_2\}.
\]

For each $i\in[h]$, define
\[
R_i
:=
\|P_{\mathcal S}\widehat u_i\|.
\]
Then
\[
R_i>0
\qquad
\text{a.s.},
\]
so the normalized in-span direction
\[
q_i
:=
\frac{
P_{\mathcal S}\widehat u_i
}{
\|P_{\mathcal S}\widehat u_i\|
}
\in
\mathbb S(\mathcal S)
\]
is almost surely well defined.

The pairs
\[
(q_i,R_i),
\qquad
i\in[h],
\]
are i.i.d. Moreover,
\[
q_i
\sim
\operatorname{Unif}(\mathbb S(\mathcal S)),
\]
and $q_i$ is independent of $R_i$. For $d>2$,
\[
\boxed{
R_i^2
\sim
\operatorname{Beta}
\left(
1,\frac{d-2}{2}
\right),
}
\]
whereas for $d=2$,
\[
R_i=1
\qquad
\text{a.s.}
\]

Consequently,
\[
r_i(0)
=
s_i(0)
=
\varepsilon R_i,
\]
and
\[
\boxed{
\widetilde u_i(0)
=
\widetilde w_i(0)
=
q_i
}
\qquad
\text{a.s.}
\]

Define the dimension-adapted radial threshold
\[
\boxed{
\tau_d
:=
\begin{cases}
1,
&
d=2,
\\[2mm]
\sqrt{
1-2^{-2/(d-2)}
},
&
d>2.
\end{cases}
}
\]
Then
\[
\mathbb P(R_i\ge\tau_d)
=
\begin{cases}
1,
&
d=2,
\\[1mm]
\frac12,
&
d>2,
\end{cases}
\]
and in particular
\[
\boxed{
\mathbb P(R_i\ge\tau_d)
\ge
\frac12
}
\qquad
\forall d\ge2.
\]

Let
\[
J_1^{\mathrm{seed}},
\qquad
J_W^{\mathrm{seed}}
\subset
\mathbb R/(2\pi\mathbb Z)
\]
be deterministic disjoint angular arcs of positive lengths
\[
\ell_1
:=
|J_1^{\mathrm{seed}}|>0,
\qquad
\ell_W
:=
|J_W^{\mathrm{seed}}|>0.
\]
Define
\[
p_1
:=
\frac{\ell_1}{2\pi},
\qquad
p_W
:=
\frac{\ell_W}{2\pi}.
\]

Define the fixed initialization cohorts jointly by angle and projected
radius:
\[
\boxed{
A_1^{\mathrm{seed}}
:=
\left\{
i\in[h]:
q_i\in\widetilde u(J_1^{\mathrm{seed}}),
\quad
R_i\ge\tau_d
\right\},
}
\]
and
\[
\boxed{
W^{\mathrm{seed}}
:=
\left\{
i\in[h]:
q_i\in\widetilde u(J_W^{\mathrm{seed}}),
\quad
R_i\ge\tau_d
\right\}.
}
\]
Since the arcs are disjoint,
\[
\boxed{
A_1^{\mathrm{seed}}
\cap
W^{\mathrm{seed}}
=
\varnothing.
}
\]

For
\[
q_d
:=
\mathbb P(R_i\ge\tau_d),
\]
angle--radius independence gives
\[
\pi_{1,d}
:=
\mathbb P(i\in A_1^{\mathrm{seed}})
=
p_1q_d,
\]
and
\[
\pi_{W,d}
:=
\mathbb P(i\in W^{\mathrm{seed}})
=
p_Wq_d.
\]
Hence
\[
\pi_{1,d}
\ge
\frac{p_1}{2},
\qquad
\pi_{W,d}
\ge
\frac{p_W}{2},
\]
uniformly in $d$.

The cohort sizes satisfy
\[
|A_1^{\mathrm{seed}}|
\sim
\operatorname{Bin}(h,\pi_{1,d}),
\qquad
|W^{\mathrm{seed}}|
\sim
\operatorname{Bin}(h,\pi_{W,d}).
\]
Therefore
\[
\boxed{
\mathbb P
\left(
|A_1^{\mathrm{seed}}|
<
\frac12\pi_{1,d}h
\right)
\le
\exp\left(-\frac{\pi_{1,d}h}{8}\right),
}
\]
and
\[
\boxed{
\mathbb P
\left(
|W^{\mathrm{seed}}|
<
\frac12\pi_{W,d}h
\right)
\le
\exp\left(-\frac{\pi_{W,d}h}{8}\right).
}
\]
Using
\[
\pi_{1,d}\ge\frac{p_1}{2},
\qquad
\pi_{W,d}\ge\frac{p_W}{2},
\]
it follows that
\[
\boxed{
\mathbb P
\left(
|A_1^{\mathrm{seed}}|
\ge
\frac{p_1}{4}h,
\quad
|W^{\mathrm{seed}}|
\ge
\frac{p_W}{4}h
\right)
\ge
1
-
\exp\left(-\frac{p_1h}{16}\right)
-
\exp\left(-\frac{p_Wh}{16}\right).
}
\]

Let
\[
\ell_{\mathrm{seed}}
:=
\min\{\ell_1,\ell_W\},
\qquad
c_{\mathrm{seed}}
:=
\frac{\ell_{\mathrm{seed}}}{8\pi}.
\]
Then
\[
\boxed{
\mathbb P
\left(
|A_1^{\mathrm{seed}}|,
|W^{\mathrm{seed}}|
\ge
c_{\mathrm{seed}}h
\right)
\ge
1
-
2\exp
\left(
-\frac{\ell_{\mathrm{seed}}}{32\pi}h
\right).
}
\]

For these deterministic arcs, define the initialization-coverage event
\[
\boxed{
\mathcal H_{\mathrm{init}}
\left(
J_1^{\mathrm{seed}},
J_W^{\mathrm{seed}}
\right)
:=
\left\{
|A_1^{\mathrm{seed}}|
\ge
c_{\mathrm{seed}}h
\right\}
\cap
\left\{
|W^{\mathrm{seed}}|
\ge
c_{\mathrm{seed}}h
\right\}.
}
\]
Consequently,
\[
\boxed{
\mathbb P
\left[
\mathcal H_{\mathrm{init}}
\left(
J_1^{\mathrm{seed}},
J_W^{\mathrm{seed}}
\right)^c
\right]
\le
2\exp
\left(
-\frac{\ell_{\mathrm{seed}}}{32\pi}h
\right).
}
\]
Once the later angular-field analysis has fixed admissible seed arcs and
therefore fixed $\ell_{\mathrm{seed}}$, the width condition
\begin{equation}
\boxed{
h
\ge
\frac{32\pi}{\ell_{\mathrm{seed}}}
\log
\frac{2}{\delta_{\mathrm{init}}}
}
\label{eq:init-width}
\end{equation}
implies
\[
\boxed{
\mathbb P
\left[
\mathcal H_{\mathrm{init}}
\left(
J_1^{\mathrm{seed}},
J_W^{\mathrm{seed}}
\right)
\right]
\ge
1-\delta_{\mathrm{init}}.
}
\]
Condition~\eqref{eq:init-width} is therefore a pending admissible-region
constraint until the basin-margin analysis fixes $J_1^{\mathrm{seed}}$ and
$J_W^{\mathrm{seed}}$.

On
\[
\mathcal H_{\mathrm{init}}
\left(
J_1^{\mathrm{seed}},
J_W^{\mathrm{seed}}
\right),
\]
one has
\[
|A_1^{\mathrm{seed}}|,
|W^{\mathrm{seed}}|
\ge
c_{\mathrm{seed}}h,
\]
and every member of either fixed seed cohort satisfies
\[
\boxed{
r_i(0)
=
s_i(0)
\ge
\varepsilon\tau_d.
}
\]
The initialization event depends only on initialization randomness and is
independent of the training-data randomness. This independence is not needed
for the final union bound.
\end{lemma}

\begin{proof}
Let
\[
g_i
\sim
\mathcal N(0,I_d)
\]
independently. By the Gaussian representation of Haar measure on the sphere,
\[
\widehat u_i
\stackrel{d}{=}
\frac{g_i}{\|g_i\|}.
\]
Decompose
\[
g_i
=
g_{i,\mathcal S}
+
g_{i,\perp},
\qquad
g_{i,\mathcal S}
=
g_{i1}e_1+g_{i2}e_2.
\]
Then
\[
P_{\mathcal S}\widehat u_i
=
\frac{
g_{i,\mathcal S}
}{
\|g_i\|
}.
\]
Since
\[
(g_{i1},g_{i2})
\sim
\mathcal N(0,I_2),
\]
one has $g_{i,\mathcal S}\neq0$ almost surely and hence $R_i>0$ almost
surely.

Moreover,
\[
q_i
=
\frac{
g_{i,\mathcal S}
}{
\|g_{i,\mathcal S}\|
}.
\]
Rotational invariance of the standard Gaussian in $\mathcal S$ gives
\[
q_i
\sim
\operatorname{Unif}(\mathbb S(\mathcal S)).
\]
Writing
\[
g_{i,\mathcal S}
=
\rho_iq_i,
\qquad
\rho_i
=
\|g_{i,\mathcal S}\|,
\]
the polar factorization of the two-dimensional isotropic Gaussian shows that
$q_i$ and $\rho_i$ are independent. Since $g_{i,\perp}$ is independent of
$g_{i,\mathcal S}$, $q_i$ is also independent of
\[
R_i
=
\frac{
\rho_i
}{
\sqrt{
\rho_i^2+\|g_{i,\perp}\|^2
}
}.
\]

For $d>2$,
\[
\rho_i^2
\sim
\chi_2^2,
\qquad
\|g_{i,\perp}\|^2
\sim
\chi_{d-2}^2,
\]
independently, so
\[
R_i^2
=
\frac{
\rho_i^2
}{
\rho_i^2+\|g_{i,\perp}\|^2
}
\sim
\operatorname{Beta}
\left(
1,\frac{d-2}{2}
\right).
\]
If
\[
Y
\sim
\operatorname{Beta}
\left(
1,\frac{d-2}{2}
\right),
\]
then
\[
\mathbb P(Y\ge y)
=
(1-y)^{(d-2)/2}.
\]
Therefore, for $d>2$,
\[
\begin{aligned}
\mathbb P(R_i\ge\tau_d)
&=
\left(
1-\tau_d^2
\right)^{(d-2)/2}
\\
&=
\left(
2^{-2/(d-2)}
\right)^{(d-2)/2}
\\
&=
\frac12.
\end{aligned}
\]
For $d=2$, $P_{\mathcal S}=I$, so $R_i=1=\tau_2$ almost surely.

The aligned initialization gives
\[
P_{\mathcal S}u_i(0)
=
P_{\mathcal S}w_i(0)
=
\varepsilon P_{\mathcal S}\widehat u_i,
\]
hence
\[
r_i(0)
=
s_i(0)
=
\varepsilon R_i
\]
and
\[
\widetilde u_i(0)
=
\widetilde w_i(0)
=
q_i.
\]

Now fix the two deterministic seed arcs. Since angle and radius are
independent,
\[
\begin{aligned}
\mathbb P(i\in A_1^{\mathrm{seed}})
&=
\mathbb P
\left(
q_i\in\widetilde u(J_1^{\mathrm{seed}})
\right)
\mathbb P(R_i\ge\tau_d)
\\
&=
p_1q_d
=
\pi_{1,d},
\end{aligned}
\]
and similarly
\[
\mathbb P(i\in W^{\mathrm{seed}})
=
p_Wq_d
=
\pi_{W,d}.
\]
Independence across neurons gives the two binomial marginal laws.

For
\[
N\sim\operatorname{Bin}(h,\pi),
\]
the multiplicative Chernoff lower-tail bound gives
\[
\mathbb P
\left(
N<\frac12\pi h
\right)
\le
\exp\left(-\frac{\pi h}{8}\right).
\]
Applying this to the two cohorts and taking a union bound proves the displayed
joint count estimate. Since
\[
\pi_{1,d}\ge\frac{p_1}{2},
\qquad
\pi_{W,d}\ge\frac{p_W}{2},
\]
one obtains
\[
\frac12\pi_{1,d}h
\ge
\frac{p_1}{4}h,
\qquad
\frac12\pi_{W,d}h
\ge
\frac{p_W}{4}h,
\]
while
\[
\exp\left(-\frac{\pi_{1,d}h}{8}\right)
\le
\exp\left(-\frac{p_1h}{16}\right),
\qquad
\exp\left(-\frac{\pi_{W,d}h}{8}\right)
\le
\exp\left(-\frac{p_Wh}{16}\right).
\]
Finally,
\[
p_1,p_W
\ge
\frac{\ell_{\mathrm{seed}}}{2\pi},
\]
so
\[
\frac{p_1}{4},
\frac{p_W}{4}
\ge
\frac{\ell_{\mathrm{seed}}}{8\pi}
=
c_{\mathrm{seed}},
\]
and
\[
\exp\left(-\frac{p_1h}{16}\right)
+
\exp\left(-\frac{p_Wh}{16}\right)
\le
2\exp
\left(
-\frac{\ell_{\mathrm{seed}}}{32\pi}h
\right).
\]
This proves the claimed event probability and the width condition. The
initial projected-radius floor follows directly from the cohort definition.
\end{proof}

\begin{remark}[Initial versus dynamical radial floor]
\label{rem:init-vs-dynamic-radius}
Lemma~\ref{lem:init-angular-coverage} builds the initial projected-radius
floor directly into the fixed seed cohorts:
\[
r_i(0)
=
s_i(0)
\ge
\varepsilon\tau_d.
\]
No later argument is permitted to redefine these cohorts by intersecting
them with another initialization event. The later radial-survival analysis
must instead prove that these already-fixed cohort members retain a
sufficiently large projected radius over the relevant dynamical window.
For $d\to\infty$,
\[
\tau_d
=
\sqrt{1-2^{-2/(d-2)}}
\sim
\sqrt{\frac{2\log2}{d}},
\]
so the initial floor has the natural $d^{-1/2}$ scale while the retained seed
fraction remains dimension-free.
\end{remark}

\begin{remark}[Distinct dynamical roles of the strong and weak seed arcs]
\label{rem:seed-basin-margin}
The deterministic initialization arcs
\[
J_1^{\mathrm{seed}},
\qquad
J_W^{\mathrm{seed}}
\]
play different dynamical roles.

The strong seed arc $J_1^{\mathrm{seed}}$ is chosen as a neighborhood of
$\phi=0$ lying inside the early target-dominated attracting region of the
strong concept.  Its endpoints must carry a strictly positive inward angular
margin, and the early projected dynamics must preserve this trapping region
while moving the cohort into a smaller strong cone.

The weak seed arc $J_W^{\mathrm{seed}}$, by contrast, is not assumed to lie
inside an attracting weak-concept basin during the early target-dominated
stage.  Such a basin is created only later, after strong-concept learning has
modified the residual field.  Choose instead a deterministic retention
interval
\[
J_W^{\mathrm{seed}}
\Subset
J_W^{\mathrm{ret}}
=
[\phi_W^-,\phi_W^+]
\]
with one-sided buffers
\[
b_{W,-}
:=
\inf J_W^{\mathrm{seed}}-\phi_W^->0,
\qquad
b_{W,+}
:=
\phi_W^+-\sup J_W^{\mathrm{seed}}>0.
\]
The early dynamics are required only to keep the fixed weak seed cohort
inside $J_W^{\mathrm{ret}}$, while retaining projected radius and projected
input/output matching.

Thus early strong angular trapping and early weak angular retention are
distinct obligations.  No target-only weak attracting margin is assumed.
The later weak-basin creation and capture analysis must connect the retained
weak region to the residual-created weak attraction basin.
\end{remark}

\begin{remark}[Seed cohorts are not final blocks]
\label{rem:init-seeds-not-blocks}
The sets
\[
A_1^{\mathrm{seed}},
\qquad
W^{\mathrm{seed}}
\]
are fixed initialization cohorts and satisfy
\[
A_1^{\mathrm{seed}}\cap W^{\mathrm{seed}}=\varnothing.
\]
They are not identified with the final post-specialization blocks $A_1$ and
$A_2$. Later dynamics may recruit neurons outside the initial seed arcs. No
time-dependent membership is introduced into the seed cohorts themselves.
\end{remark}

\begin{remark}[Joint closure of seed arcs, field margins, retention, and width]
\label{rem:seed-margin-closure}
The deterministic seed arcs
\[
J_1^{\mathrm{seed}},
\qquad
J_W^{\mathrm{seed}}
\]
cannot be chosen independently of the later dynamical remainder bounds.
Their constraints are asymmetric.

For the strong arc, the early target field must provide both a positive
radial margin and a positive angular trapping/coercivity margin.  In the
notation of Lemma~\ref{lem:early-projected-trapping}, the relevant quantities
include
\[
\gamma_{\mathrm{rad}}>0,
\qquad
\lambda_{1,\mathrm{ang}}>0,
\qquad
T_{1,+}<\infty.
\]
Widening $J_1^{\mathrm{seed}}$ increases its initialization mass but generally
weakens angular coercivity and increases the tangential gated moment
$T_{1,+}$; narrowing it has the opposite effect.  Thus the strong seed width
has a genuine two-sided cost.

For the weak arc, the early stage requires a positive radial margin but not a
target-only attracting angular margin.  Instead one must choose
\[
J_W^{\mathrm{seed}}
\Subset
J_W^{\mathrm{ret}}
\]
with positive one-sided buffers $b_{W,-},b_{W,+}$ and verify that the weak
angular drift accumulated up to the early trapping time remains inside these
buffers.  The weak target-only angular field may drive the cohort toward the
strong direction during this stage; the proof only requires that the cohort
remain recruitable when the residual-created weak basin appears later.

Both seed arcs must also retain positive target radial mass.  Schematically,
if
\[
\gamma_p^{\mathrm{seed}}
:=
\inf_{\phi\in J_p^{\mathrm{seed}}}
\Gamma_{\mathrm{tar}}(\phi),
\]
then the early projected dynamics must port these positive radial margins to
the true-residual flow after accounting for learned-output, angular-gap,
boundary, and perpendicular terms.

Accordingly, the proof must close the coupled choices
\[
\boxed{
\begin{array}{c}
\text{strong seed width}
\longleftrightarrow
\text{radial and angular margins}
\longleftrightarrow
\text{allowable early remainders},
\\[2mm]
\text{weak seed width}
\longleftrightarrow
\text{retention buffer and radial margin}
\longleftrightarrow
\text{weak drift accumulated to }t_{\mathrm{trap}}.
\end{array}
}
\]
At the same time, decreasing either seed width decreases
\[
\ell_{\mathrm{seed}}
=
\min\left\{
|J_1^{\mathrm{seed}}|,
|J_W^{\mathrm{seed}}|
\right\},
\]
thereby strengthening the initialization-width requirement
\[
h
\ge
\frac{32\pi}{\ell_{\mathrm{seed}}}
\log\frac{2}{\delta_{\mathrm{init}}}.
\]
The constant-fraction guarantee therefore belongs to the asymptotic theorem
regime; it is not imposed merely to certify a particular experimental width.
A sharper multiplicative-Chernoff allocation can later trade the guaranteed
fraction of a cohort for a smaller width constant without changing its
$\Theta(h)$ scaling.
\end{remark}


%%%%%%%%%%%%%%%%%%%%%%%%%%%%%%%%%%%%%%%%%%%%%%%%%%%%%%%%%%%%%%%%%%%%%%%%%%%%%%%
% T1.7: exact projected polar dynamics and projected mismatch
%%%%%%%%%%%%%%%%%%%%%%%%%%%%%%%%%%%%%%%%%%%%%%%%%%%%%%%%%%%%%%%%%%%%%%%%%%%%%%%

\begin{lemma}[Exact projected polar dynamics]
\label{lem:projected-polar-dynamics}
Work on the data event $\mathcal H_{\mathrm{data}}(\delta)$ and along the
fixed absolutely continuous Clarke trajectory and measurable selector from
Lemma~\ref{lem:gf-regularity}.  Assume the training data lie in
\[
\mathcal S=\operatorname{span}\{e_1,e_2\}.
\]
For each neuron define
\[
v_i:=P_{\mathcal S}u_i,
\qquad
z_i:=P_{\mathcal S}w_i,
\qquad
r_i:=\|v_i\|,
\qquad
s_i:=\|z_i\|.
\]
On any interval on which $r_i,s_i>0$, write
\[
v_i=r_i\widetilde u_i,
\qquad
z_i=s_i\widetilde w_i,
\]
and choose continuous angular lifts
\[
\widetilde u_i=\widetilde u(\phi_i),
\qquad
\widetilde w_i=\widetilde u(\psi_i),
\]
where
\[
\widetilde u(\phi)
:=
\cos\phi\,e_1+\sin\phi\,e_2.
\]
Let $\mathsf J:\mathcal S\to\mathcal S$ be the quarter-turn
\[
\mathsf Je_1=e_2,
\qquad
\mathsf Je_2=-e_1,
\]
and define
\[
\widetilde u^\angle(\phi)
:=
\mathsf J\widetilde u(\phi)
=
-\sin\phi\,e_1+\cos\phi\,e_2,
\]
\[
\widetilde u_i^\angle
:=
\widetilde u^\angle(\phi_i),
\qquad
\widetilde w_i^\angle
:=
\widetilde u^\angle(\psi_i).
\]
Also write
\[
w_{i,\perp}:=P_{\mathcal S^\perp}w_i,
\qquad
u_{i,\perp}:=P_{\mathcal S^\perp}u_i.
\]

For arbitrary angles $\phi,\psi$, define the true-residual local fields
\[
\boxed{
\Gamma_{\mathrm{true}}(\phi,\psi;t)
:=
-\mathbb E_n
\left[
\left\langle
P_{\mathcal S}\rho_t(x),
\widetilde u(\psi)
\right\rangle
(\widetilde u(\phi)^\top x)_+
\right],
}
\]
\[
\boxed{
\mathcal V_{\mathrm{true}}(\phi,\psi;t)
:=
-\mathbb E_n
\left[
\left\langle
P_{\mathcal S}\rho_t(x),
\widetilde u(\psi)
\right\rangle
\mathbf 1\{\widetilde u(\phi)^\top x>0\}
(\widetilde u^\angle(\phi))^\top x
\right],
}
\]
and
\[
\boxed{
\mathcal W_{\mathrm{true}}(\phi,\psi;t)
:=
-\mathbb E_n
\left[
\left\langle
P_{\mathcal S}\rho_t(x),
\widetilde u^\angle(\psi)
\right\rangle
(\widetilde u(\phi)^\top x)_+
\right].
}
\]

Then, for a.e. time on the interval,
\[
\boxed{
\dot z_i
=
-r_i
\mathbb E_n
\left[
P_{\mathcal S}\rho_t(x)
(\widetilde u_i^\top x)_+
\right].
}
\]
Consequently,
\[
\boxed{
\dot s_i
=
r_i\Gamma_{\mathrm{true}}(\phi_i,\psi_i;t),
}
\]
\[
\boxed{
\dot\psi_i
=
\frac{r_i}{s_i}
\mathcal W_{\mathrm{true}}(\phi_i,\psi_i;t).
}
\]

Define the perpendicular radial and angular remainders
\[
\boxed{
R_{i,\perp}^{\mathrm{rad}}(t)
:=
-\mathbb E_n
\left[
\left\langle
P_{\mathcal S^\perp}\rho_t(x),
w_{i,\perp}(t)
\right\rangle
(\widetilde u_i(t)^\top x)_+
\right],
}
\]
and
\[
\boxed{
R_{i,\perp}^{\mathrm{ang}}(t)
:=
-\frac1{r_i(t)}
\mathbb E_n
\left[
\left\langle
P_{\mathcal S^\perp}\rho_t(x),
w_{i,\perp}(t)
\right\rangle
\mathbf 1\{\widetilde u_i(t)^\top x>0\}
(\widetilde u_i^\angle(t))^\top x
\right].
}
\]
With the same boundary correction $E_i^{\mathrm{bd}}$ as in
Remark~\ref{rem:angular-boundary-selection}, one has
\[
\boxed{
\dot r_i
=
s_i\Gamma_{\mathrm{true}}(\phi_i,\psi_i;t)
+
R_{i,\perp}^{\mathrm{rad}}(t),
}
\]
and
\[
\boxed{
\dot\phi_i
=
\frac{s_i}{r_i}
\mathcal V_{\mathrm{true}}(\phi_i,\psi_i;t)
+
R_{i,\perp}^{\mathrm{ang}}(t)
+
E_i^{\mathrm{bd}}(t).
}
\]
Hence the projected angular gap obeys
\[
\boxed{
\frac{d}{dt}(\phi_i-\psi_i)
=
\frac{s_i}{r_i}\mathcal V_{\mathrm{true}}(\phi_i,\psi_i;t)
-
\frac{r_i}{s_i}\mathcal W_{\mathrm{true}}(\phi_i,\psi_i;t)
+
R_{i,\perp}^{\mathrm{ang}}(t)
+
E_i^{\mathrm{bd}}(t).
}
\]

Under the aligned initialization
\[
w_i(0)=u_i(0)=\varepsilon\widehat u_i,
\]
input-perpendicular conservation and exact balance imply the exact conserved
offset
\[
\boxed{
\|u_i(t)\|^2
=
r_i(t)^2
+
\varepsilon^2(1-R_i^2),
}
\]
where
\[
R_i=\|P_{\mathcal S}\widehat u_i\|.
\]
They also imply
\[
\boxed{
S(t)
:=
\sum_{j=1}^h\|u_j(t)\|^2
=
\sum_{j=1}^h r_j(t)^2
+
S_\perp^{(0)},
}
\]
with
\[
S_\perp^{(0)}
:=
\varepsilon^2
\sum_{j=1}^h(1-R_j^2)
\le
h\varepsilon^2.
\]

Define
\[
\boxed{
\Xi_{i,\perp}(t)
:=
\|w_{i,\perp}(t)\|^2
-
\|w_{i,\perp}(0)\|^2.
}
\]
Then
\[
\boxed{
r_i(t)^2-s_i(t)^2
=
\Xi_{i,\perp}(t),
}
\]
and
\[
\boxed{
\dot\Xi_{i,\perp}(t)
=
2r_i(t)R_{i,\perp}^{\mathrm{rad}}(t)
}
\]
for a.e. $t$.

If
\[
F_\perp(t)
:=
\sup_{x\in\mathcal D}
\|P_{\mathcal S^\perp}\rho_t(x)\|,
\]
then, because $x\in\mathcal S$ and
Corollary~\ref{cor:perp-dynamics} gives
$\|W_\perp(t)\|_F\le\varepsilon\sqrt h$,
\[
\boxed{
F_\perp(t)
\le
X_{\delta_{\mathrm{strong}}}
\varepsilon\sqrt h\,
\sqrt{S(t)}.
}
\]
Consequently,
\[
\boxed{
|R_{i,\perp}^{\mathrm{rad}}(t)|
\le
X_{\delta_{\mathrm{strong}}}
F_\perp(t)
\|w_{i,\perp}(t)\|,
}
\]
and
\[
\boxed{
|R_{i,\perp}^{\mathrm{ang}}(t)|
\le
\frac{
X_{\delta_{\mathrm{strong}}}
F_\perp(t)
\|w_{i,\perp}(t)\|
}{r_i(t)}.
}
\]
Thus the single envelope
\[
\boxed{
\mathfrak B_i(t)
:=
\frac{
X_{\delta_{\mathrm{strong}}}
F_\perp(t)
\|w_{i,\perp}(t)\|
}{r_i(t)}
}
\]
controls both normalized perpendicular forcings:
\[
\boxed{
\frac{|R_{i,\perp}^{\mathrm{rad}}(t)|}{r_i(t)}
\le
\mathfrak B_i(t),
\qquad
|R_{i,\perp}^{\mathrm{ang}}(t)|
\le
\mathfrak B_i(t).
}
\]
\end{lemma}

\begin{proof}
The projected output equation follows directly from
Lemma~\ref{lem:gf-regularity}:
\[
\dot z_i
=
-\mathbb E_n
\left[
P_{\mathcal S}\rho_t(x)
\sigma(u_i^\top x)
\right].
\]
Since $x\in\mathcal S$,
\[
u_i^\top x
=
r_i\widetilde u_i^\top x,
\]
so positive homogeneity of ReLU gives
\[
\sigma(u_i^\top x)
=
r_i(\widetilde u_i^\top x)_+.
\]
Taking components of $\dot z_i$ along
$\widetilde w_i$ and $\widetilde w_i^\angle$ proves the $s_i$ and $\psi_i$
equations.

For the input projection,
\[
\dot v_i
=
-\mathbb E_n
\left[
\langle\rho_t,w_i\rangle
\alpha_i(t;x)x
\right].
\]
Decompose
\[
w_i=s_i\widetilde w_i+w_{i,\perp}.
\]
The projected radial component is selection-free because
\[
\alpha_i(t;x)\widetilde u_i^\top x
=
(\widetilde u_i^\top x)_+.
\]
This yields
\[
\dot r_i
=
s_i\Gamma_{\mathrm{true}}
+
R_{i,\perp}^{\mathrm{rad}}.
\]
For the angular component, replace the fixed selector by the hard gate away
from the activation boundary and retain the exact discrepancy
$E_i^{\mathrm{bd}}$ from
Remark~\ref{rem:angular-boundary-selection}.  This gives the stated
$\phi_i$ equation.  Subtracting the $\psi_i$ equation gives the gap identity.

Input-perpendicular conservation gives
\[
\|P_{\mathcal S^\perp}u_i(t)\|^2
=
\varepsilon^2(1-R_i^2).
\]
Combining this with exact balance proves the conserved-offset identity and,
after summation, the identity for $S(t)$.

Exact balance also gives
\[
r_i^2+\|P_{\mathcal S^\perp}u_i(0)\|^2
=
s_i^2+\|w_{i,\perp}(t)\|^2.
\]
Since aligned initialization implies
$w_{i,\perp}(0)=P_{\mathcal S^\perp}u_i(0)$, this is exactly
\[
r_i^2-s_i^2=\Xi_{i,\perp}.
\]
Moreover,
\[
\dot w_{i,\perp}
=
-\mathbb E_n
\left[
P_{\mathcal S^\perp}\rho_t(x)
\sigma(u_i^\top x)
\right],
\]
so
\[
\frac{d}{dt}\|w_{i,\perp}\|^2
=
2r_iR_{i,\perp}^{\mathrm{rad}}.
\]

Finally, because $P_{\mathcal S^\perp}x=0$,
\[
P_{\mathcal S^\perp}\rho_t(x)
=
P_{\mathcal S^\perp}f_t(x)
=
\sum_{j=1}^h
w_{j,\perp}\sigma(u_j^\top x).
\]
Cauchy--Schwarz gives
\[
\|P_{\mathcal S^\perp}f_t(x)\|
\le
\|W_\perp(t)\|_F
\left(
\sum_{j=1}^h\sigma(u_j^\top x)^2
\right)^{1/2}.
\]
Using
$\|W_\perp(t)\|_F\le\varepsilon\sqrt h$ and
$\sigma(u_j^\top x)^2\le X_{\delta_{\mathrm{strong}}}^2\|u_j\|^2$
proves the $F_\perp$ bound.  The two remainder estimates are then immediate
from Cauchy--Schwarz and
$|(\widetilde u_i^\top x)_+|,
|\widetilde u_i^{\angle\top}x|
\le X_{\delta_{\mathrm{strong}}}$.
\end{proof}

\begin{corollary}[Contractive Cartesian projected mismatch]
\label{cor:projected-mismatch}
Under the assumptions of Lemma~\ref{lem:projected-polar-dynamics}, define
\[
\boxed{
\delta_i(t)
:=
P_{\mathcal S}\bigl(u_i(t)-w_i(t)\bigr)
=
v_i(t)-z_i(t).
}
\]
Then for a.e. $t$,
\[
\boxed{
\dot\delta_i(t)
=
-Q_i^\alpha(t)\delta_i(t)
-
\mathcal E_i^f(t),
}
\]
where
\[
\boxed{
\mathcal E_i^f(t)
:=
\mathbb E_n
\left[
\langle f_t(x),w_i(t)\rangle
\alpha_i(t;x)x
\right]
-
\mathbb E_n
\left[
P_{\mathcal S}f_t(x)
\sigma(u_i(t)^\top x)
\right].
}
\]
Consequently,
\[
\boxed{
\frac12\frac{d}{dt}\|\delta_i(t)\|^2
=
-\delta_i(t)^\top Q_i^\alpha(t)\delta_i(t)
-
\langle\delta_i(t),\mathcal E_i^f(t)\rangle.
}
\]
Since $Q_i^\alpha(t)\succeq0$,
\[
\boxed{
\|\delta_i(t)\|
\le
\|\delta_i(t_0)\|
+
\int_{t_0}^t\|\mathcal E_i^f(s)\|\,ds.
}
\]
Under aligned initialization, $\delta_i(0)=0$, and on
$\mathcal H_{\mathrm{data}}(\delta)$,
\[
\boxed{
\|\mathcal E_i^f(t)\|
\le
2X_{\delta_{\mathrm{strong}}}
F(t)\|u_i(t)\|,
}
\]
where
\[
F(t):=
\sup_{x\in\mathcal D}\|f_t(x)\|.
\]
Hence
\[
\boxed{
\|\delta_i(t)\|
\le
2X_{\delta_{\mathrm{strong}}}
\int_0^t
F(s)r_i(s)\,ds
+
2X_{\delta_{\mathrm{strong}}}
\varepsilon\sqrt{1-R_i^2}
\int_0^tF(s)\,ds.
}
\]
\end{corollary}

\begin{proof}
Expanding the projected input and output equations into target and network
parts gives
\[
\dot v_i
=
Q_i^\alpha z_i
-
\mathbb E_n
\left[
\langle f_t,w_i\rangle\alpha_i x
\right],
\]
while
\[
\dot z_i
=
Q_i^\alpha v_i
-
\mathbb E_n
\left[
P_{\mathcal S}f_t\sigma(u_i^\top x)
\right].
\]
Subtracting yields the displayed equation for $\delta_i$.  The energy
identity follows by taking the inner product with $\delta_i$, and the norm
comparison follows from $Q_i^\alpha\succeq0$.

For the forcing term,
\[
\left\|
\mathbb E_n
\left[
\langle f_t,w_i\rangle\alpha_i x
\right]
\right\|
\le
X_{\delta_{\mathrm{strong}}}F(t)\|w_i\|,
\]
and
\[
\left\|
\mathbb E_n
\left[
P_{\mathcal S}f_t\sigma(u_i^\top x)
\right]
\right\|
\le
X_{\delta_{\mathrm{strong}}}F(t)\|u_i\|.
\]
Exact balance gives $\|w_i\|=\|u_i\|$.  Finally,
\[
\|u_i\|
=
\sqrt{r_i^2+\varepsilon^2(1-R_i^2)}
\le
r_i+\varepsilon\sqrt{1-R_i^2},
\]
which proves the final estimate.
\end{proof}

\begin{corollary}[Projected radius and angle transfer]
\label{cor:projected-radius-transfer}
Assume $r_i>0$ and define
\[
\eta_i^{\mathrm{mis}}
:=
\frac{\|\delta_i\|}{r_i}.
\]
If $\eta_i^{\mathrm{mis}}<1$, then
\[
\boxed{
|r_i-s_i|
\le
\|\delta_i\|,
\qquad
(1-\eta_i^{\mathrm{mis}})r_i
\le
s_i
\le
(1+\eta_i^{\mathrm{mis}})r_i.
}
\]
If in addition $s_i>0$, then
\[
\boxed{
\|\delta_i\|^2
=
(r_i-s_i)^2
+
4r_is_i
\sin^2\!\left(
\frac{\angle(\widetilde u_i,\widetilde w_i)}2
\right).
}
\]
Consequently,
\[
\boxed{
\sin\!\left(
\frac{\angle(\widetilde u_i,\widetilde w_i)}2
\right)
\le
\frac{\eta_i^{\mathrm{mis}}}{2\sqrt{1-\eta_i^{\mathrm{mis}}}},
}
\]
and
\[
\boxed{
\|\widetilde u_i-\widetilde w_i\|
\le
\frac{\eta_i^{\mathrm{mis}}}{\sqrt{1-\eta_i^{\mathrm{mis}}}}.
}
\]

Equivalently, if
\[
\vartheta_i
:=
\frac{r_i-s_i}{r_i+s_i},
\]
then the exact balance identity gives
\[
\boxed{
\vartheta_i
=
\frac{\Xi_{i,\perp}}{(r_i+s_i)^2},
}
\]
and, for a.e. $t$,
\[
\boxed{
\dot\vartheta_i
=
-2\Gamma_{\mathrm{true}}(\phi_i,\psi_i;t)\vartheta_i
+
\frac{2s_i}{(r_i+s_i)^2}
R_{i,\perp}^{\mathrm{rad}}.
}
\]
\end{corollary}

\begin{proof}
The first inequality is the reverse triangle inequality applied to
$v_i-z_i$.  Expanding
$\|r_i\widetilde u_i-s_i\widetilde w_i\|^2$ gives the exact angle formula,
and the stated bounds follow from
$s_i\ge(1-\eta_i^{\mathrm{mis}})r_i$.

For $\vartheta_i$, use
\[
(r_i-s_i)(r_i+s_i)
=
r_i^2-s_i^2
=
\Xi_{i,\perp}.
\]
Differentiating $\vartheta_i=(r_i-s_i)/(r_i+s_i)$ and substituting the exact
$r_i,s_i$ equations from Lemma~\ref{lem:projected-polar-dynamics} gives the
last identity.
\end{proof}

\begin{remark}[Asymmetry of projected matching mechanisms]
\label{rem:projected-dynamics-asymmetry}
The projected input/output mismatch has two structurally different sources.
The squared-radius mismatch is driven only by the perpendicular output
forcing through
\[
\dot\Xi_{i,\perp}
=
2r_iR_{i,\perp}^{\mathrm{rad}},
\]
whereas the angular mismatch is governed by the difference between the two
projected angular fields and has its own restoring target-only dynamics.
Accordingly, later quantitative arguments should not use the Cartesian norm
$\|\delta_i\|$ to charge radial and angular errors at the same rate when the
exact polar identities give sharper separate controls.
\end{remark}

\begin{corollary}[Aligned invariant manifold for the target-only flow]
\label{cor:target-only-aligned-manifold}
Consider the auxiliary target-only objective
\[
\mathcal L_{\mathrm{tar}}(\theta)
:=
-\mathbb E_n\langle x,f_\theta(x)\rangle
\]
with its Clarke gradient flow, using the same selector convention as in
Lemma~\ref{lem:gf-regularity}.  Under aligned initialization
\[
w_i(0)=u_i(0)=\varepsilon\widehat u_i,
\]
one has
\[
\boxed{
w_i(t)=u_i(t)
}
\qquad
\forall i,
\quad
\forall t\ge0.
\]
In particular,
\[
\boxed{
r_i(t)=s_i(t),
\qquad
\phi_i(t)=\psi_i(t)
}
\]
whenever the projected radius is nonzero, and the common projected vector
obeys
\[
\boxed{
\dot v_i(t)
=
\mathbb E_n
\left[
x(v_i(t)^\top x)_+
\right].
}
\]
The restricted vector field is globally Lipschitz with Lipschitz constant at
most $\lambda_{\max}(\Sigma_n)$.

On the aligned manifold the polar target-only fields reduce to
\[
\boxed{
\dot r_i
=
r_i\Gamma_{\mathrm{tar}}(\phi_i),
\qquad
\dot\phi_i
=
V_{\mathrm{tar}}(\phi_i),
}
\]
where
\[
\Gamma_{\mathrm{tar}}(\phi)
=
\mathbb E_n
\left[
(\widetilde u(\phi)^\top x)_+^2
\right]
\ge0,
\]
and
\[
V_{\mathrm{tar}}(\phi)
=
\mathbb E_n
\left[
(\widetilde u(\phi)^\top x)_+
(\widetilde u^\angle(\phi))^\top x
\right].
\]
\end{corollary}

\begin{proof}
For the target-only flow,
\[
\dot w_i
=
\mathbb E_n
\left[
x\sigma(u_i^\top x)\right],
\qquad
\dot u_i
=
\mathbb E_n
\left[
\langle x,w_i\rangle\alpha_i(t;x)x
\right].
\]
Both perpendicular components are constant because $x\in\mathcal S$.
For the projected mismatch $\delta_i=P_{\mathcal S}(u_i-w_i)$, the same
calculation as in Corollary~\ref{cor:projected-mismatch} gives
\[
\dot\delta_i
=
-Q_i^\alpha\delta_i.
\]
Since $\delta_i(0)=0$ and $Q_i^\alpha\succeq0$,
\[
\frac12\frac{d}{dt}\|\delta_i(t)\|^2
=
-\delta_i(t)^\top Q_i^\alpha(t)\delta_i(t)
\le0,
\]
so $\delta_i\equiv0$.  The perpendicular components agree initially and remain
constant, hence $u_i=w_i$ for all time.

Substituting $u_i=w_i$ gives the common projected ODE.  Its Lipschitz bound
follows from
\[
\left\|
\mathbb E_n
\left[
xx^\top\beta(x)
\right]
\right\|_{\mathrm{op}}
\le
\lambda_{\max}(\Sigma_n)
\]
for every measurable weight $\beta(x)\in[0,1]$.  The polar equations follow
by taking radial and angular components.
\end{proof}



%%%%%%%%%%%%%%%%%%%%%%%%%%%%%%%%%%%%%%%%%%%%%%%%%%%%%%%%%%%%%%%%%%%%%%%%%%%%%%%
% T1.8: early projected trapping, radial survival, and weak-seed retention
%%%%%%%%%%%%%%%%%%%%%%%%%%%%%%%%%%%%%%%%%%%%%%%%%%%%%%%%%%%%%%%%%%%%%%%%%%%%%%%

\begin{lemma}[Early projected trapping and weak-seed retention]
\label{lem:early-projected-trapping}
Work on
\[
\mathcal H_{\mathrm{data}}(\delta)
\cap\mathcal G_{\mathrm{gate}}
\cap
\mathcal H_{\mathrm{init}}
\left(
J_1^{\mathrm{seed}},
J_W^{\mathrm{seed}}
\right),
\]
and assume the conclusions of
Lemmas~\ref{lem:norm-output} and
\ref{lem:projected-polar-dynamics}.

For compactness write
\[
X:=X_{\delta_{\mathrm{strong}}},
\qquad
\lambda:=\overline\lambda_\delta,
\qquad
\beta:=2\lambda.
\]
Fix a common spectral-envelope horizon
\[
H>0
\]
and a parameter
\[
\boxed{\overline\chi\in(0,1)}
\]
such that
\[
\boxed{
\chi(H)
=
\frac{X^2h\varepsilon^2}{\lambda}
\left(e^{2\lambda H}-1\right)
\le
\overline\chi.
}
\]
Set
\[
\boxed{D_H:=1-\chi(H)>0.}
\]

Choose angular representatives so that
\[
J_1^{\mathrm{seed}}
=
[-\theta_1,\theta_1],
\qquad
0<\theta_1<\frac\pi2,
\]
and choose a deterministic weak retention interval
\[
J_W^{\mathrm{seed}}
\Subset
J_W^{\mathrm{ret}}
=
[\phi_W^-,\phi_W^+].
\]
Define the one-sided retention buffers
\[
\boxed{
b_{W,-}
:=
\inf J_W^{\mathrm{seed}}-\phi_W^->0,
\qquad
b_{W,+}
:=
\phi_W^+-\sup J_W^{\mathrm{seed}}>0.
}
\]

For
\[
\widetilde u(\phi)
=
\cos\phi\,e_1+\sin\phi\,e_2,
\qquad
\widetilde u^\angle(\phi)
=
-\sin\phi\,e_1+\cos\phi\,e_2,
\]
define the hard-gated empirical moment
\[
\boxed{
A(\phi)
:=
\mathbb E_n
\left[
xx^\top
\mathbf 1\{\widetilde u(\phi)^\top x>0\}
\right].
}
\]
Set
\[
\boxed{
\Gamma_{\mathrm{tar}}(\phi)
:=
\widetilde u(\phi)^\top A(\phi)\widetilde u(\phi),
}
\]
\[
\boxed{
V_{\mathrm{tar}}(\phi)
:=
\widetilde u(\phi)^\top A(\phi)\widetilde u^\angle(\phi),
}
\]
and
\[
\boxed{
T_{\mathrm{tar}}(\phi)
:=
\widetilde u^\angle(\phi)^\top A(\phi)\widetilde u^\angle(\phi).
}
\]
For later two-angle comparisons also define
\[
\boxed{
\Gamma_{\mathrm{tar}}(\phi,\psi)
:=
\mathbb E_n
\left[
\langle x,\widetilde u(\psi)\rangle
(\widetilde u(\phi)^\top x)_+
\right],
}
\]
\[
\boxed{
\mathcal V_{\mathrm{tar}}(\phi,\psi)
:=
\mathbb E_n
\left[
\langle x,\widetilde u(\psi)\rangle
\mathbf 1\{\widetilde u(\phi)^\top x>0\}
(\widetilde u^\angle(\phi))^\top x
\right],
}
\]
and
\[
\boxed{
\mathcal W_{\mathrm{tar}}(\phi,\psi)
:=
\mathbb E_n
\left[
\langle x,\widetilde u^\angle(\psi)\rangle
(\widetilde u(\phi)^\top x)_+
\right].
}
\]
Thus
\[
\Gamma_{\mathrm{tar}}(\phi)
=
\Gamma_{\mathrm{tar}}(\phi,\phi),
\qquad
V_{\mathrm{tar}}(\phi)
=
\mathcal V_{\mathrm{tar}}(\phi,\phi)
=
\mathcal W_{\mathrm{tar}}(\phi,\phi).
\]
Let
\[
\mathcal J_{\mathrm{early}}
:=
J_1^{\mathrm{seed}}
\cup
J_W^{\mathrm{ret}}.
\]
Define the deterministic finite-sample field constants
\[
\boxed{
\gamma_{\mathrm{rad}}
:=
\inf_{\phi\in\mathcal J_{\mathrm{early}}}
\Gamma_{\mathrm{tar}}(\phi),
}
\]
\[
\boxed{
\nu_{\mathrm{early}}
:=
\sup_{\phi\in\mathcal J_{\mathrm{early}}}
|V_{\mathrm{tar}}(\phi)|,
}
\]
\[
\boxed{
\Gamma_+
:=
\sup_{\phi\in\mathcal J_{\mathrm{early}}}
\Gamma_{\mathrm{tar}}(\phi),
\qquad
T_+
:=
\sup_{\phi\in\mathcal J_{\mathrm{early}}}
T_{\mathrm{tar}}(\phi).
}
\]
For the strong arc, use the bias-aware coercivity of
Proposition~\ref{sbnew:bias} below. Define
\[
b_{\rm tar}:=
\left|\frac{\bar a_{12,\star}^{(1)}(0)+\bar a_{12,\star}^{(2)}(0)}2\right|
+\frac{\varepsilon_{12}^{(1)}+\varepsilon_{12}^{(2)}}2.
\]
This is a deterministic bound for $|V_{\rm tar}(0)|$ on the existing
gated-moment event, which is included in the event for this invocation.
The curvature infimum below may be replaced by the deterministic lower
bound in \eqref{sbnew:kappa-cert}. Define
\[
\boxed{
\lambda_{1,\mathrm{ang}}
:=
\inf_{|\phi|\le\theta_1}
\left(\Gamma_{\mathrm{tar}}(\phi)-T_{\mathrm{tar}}(\phi)\right),
}
\]
and
\[
\boxed{
T_{1,+}
:=
\sup_{\phi\in J_1^{\mathrm{seed}}}
T_{\mathrm{tar}}(\phi).
}
\]
For the weak retention interval define
\[
\boxed{
T_{W,+}
:=
\sup_{\phi\in J_W^{\mathrm{ret}}}
T_{\mathrm{tar}}(\phi),
}
\]
\[
\boxed{
M_{W,-}
:=
\sup_{\phi\in J_W^{\mathrm{ret}}}
[-V_{\mathrm{tar}}(\phi)]_+,
\qquad
M_{W,+}
:=
\sup_{\phi\in J_W^{\mathrm{ret}}}
[V_{\mathrm{tar}}(\phi)]_+.
}
\]
Finally define the gated angular-gap damping floor
\[
\boxed{
\kappa_{\mathrm{gap}}
:=
\inf_{\phi\in\mathcal J_{\mathrm{early}}}
\left(
\Gamma_{\mathrm{tar}}(\phi)
+
T_{\mathrm{tar}}(\phi)
\right).
}
\]

On $\mathcal H_{\mathrm{data}}(\delta)$, Lemma~\ref{lem:Hstrong} implies the
explicit lower bound
\[
\boxed{
\kappa_{\mathrm{gap}}
\ge
\frac12
\inf_{\phi\in\mathcal J_{\mathrm{early}}}
\sum_{p=1}^2
\left(
\mu_p-\sigma_pL_{\delta_{\mathrm{strong}}}
\right)^2
\left[
P_p(H_\phi)
-
\Delta_{\mathrm{VC},\delta_{\mathrm{strong}}}
\right]_+,
}
\]
where
\[
P_1(H_\phi)
=
\Phi\!\left(
\frac{\mu_1\cos\phi}{\mathsf v(\phi)}
\right),
\qquad
P_2(H_\phi)
=
\Phi\!\left(
\frac{\mu_2\sin\phi}{\mathsf v(\phi)}
\right).
\]
Assume
\[
\boxed{
\gamma_{\mathrm{rad}}>0,
\qquad
\lambda_{1,\mathrm{ang}}>0,
\qquad
\kappa_{\mathrm{gap}}>0.
}
\]

For $a\in\mathbb R$ define
\[
\mathcal E_a(t)
:=
\begin{cases}
\dfrac{e^{at}-1}{a},&a\neq0,\\[2mm]
t,&a=0,
\end{cases}
\]
and, for $\omega>0$,
\[
\mathcal K_{\omega,a}(t)
:=
\begin{cases}
\dfrac{e^{at}-e^{-\omega t}}{a+\omega},&a+\omega\neq0,\\[2mm]
t e^{-\omega t},&a+\omega=0.
\end{cases}
\]

Lemma~\ref{lem:norm-output} supplies, for every $0\le t\le H$,
\[
\boxed{
F(t)
:=
\sup_{x\in\mathcal D}\|f_t(x)\|
\le
\overline F(t)
:=
F_0e^{\beta t},
}
\]
where
\[
\boxed{
F_0
:=
\frac{Xh\varepsilon^2}{D_H}.
}
\]
Also set
\[
\boxed{
C_B
:=
\frac{X^2h\varepsilon^2}{\sqrt{D_H}}.
}
\]

Fix a candidate trapping horizon
\[
\boxed{0<T\le H.}
\]
The coarse perpendicular envelope used in the first-exit argument is
\[
\boxed{
\overline B^{(0)}(t)
:=
\frac{C_B}{\tau_d}e^{\lambda t},
}
\]
which is valid on any interval on which the radial-floor guard
$r_i(t)\ge r_i(0)$ holds.  Define
\[
\boxed{
P_T^{(0)}
:=
\int_0^T\overline B^{(0)}(s)\,ds
=
\frac{C_B}{\tau_d}\mathcal E_\lambda(T).
}
\]
Assume
\[
\boxed{
4P_T^{(0)}<1.
}
\]
Define
\[
a_{0,-}
:=
\sqrt{1-4P_T^{(0)}},
\qquad
 a_{0,+}
:=
\sqrt{1+4P_T^{(0)}}.
\]
For $0<a_-\le1\le a_+$ set
\[
m_a
:=
\min\{a_-,a_+^{-1}\},
\]
\[
d_a
:=
\max\{a_-^{-1}-a_-,\,a_+-a_+^{-1}\},
\]
and
\[
\omega_\angle(a_-,a_+)
:=
\frac2\pi m_a\kappa_{\mathrm{gap}}.
\]
Define
\[
C_{\angle,0}(a_-,a_+)
:=
d_a\nu_{\mathrm{early}}
+
a_+\frac{X^2}{2n},
\]
\[
C_{\angle,F}(a_-,a_+)
:=
a_-^{-1}\sqrt{\Gamma_+}
+
a_+\sqrt{T_+}
+
a_+\frac{X}{2n},
\]
and
\[
C_{\angle,B}:=1+\frac1{2n}.
\]
The coarse angular-gap envelope is
\[
\boxed{
\begin{aligned}
D_{\angle,0}(T)
&:=
\frac{
C_{\angle,0}(a_{0,-},a_{0,+})
}{
\omega_\angle(a_{0,-},a_{0,+})
}
\left(
1-e^{-\omega_\angle(a_{0,-},a_{0,+})T}
\right)
\\
&\quad+
C_{\angle,F}(a_{0,-},a_{0,+})
F_0
\mathcal K_{\omega_\angle(a_{0,-},a_{0,+}),\beta}(T)
\\
&\quad+
C_{\angle,B}
\frac{C_B}{\tau_d}
\mathcal K_{\omega_\angle(a_{0,-},a_{0,+}),\lambda}(T).
\end{aligned}
}
\]
Assume
\[
\boxed{
D_{\angle,0}(T)<\frac\pi2.
}
\]
Define the coarse true-radial margin
\[
\boxed{
G_{\mathrm{rad},0}(T)
:=
\cos(D_{\angle,0}(T))\gamma_{\mathrm{rad}}
-
\sin(D_{\angle,0}(T))\nu_{\mathrm{early}}
-
\overline F(T)\sqrt{\Gamma_+},
}
\]
and assume the coarse radial-dominance condition
\[
\boxed{
a_{0,-}G_{\mathrm{rad},0}(T)
>
\overline B^{(0)}(T).
}
\]
Define, once and for all from this coarse tier,
\[
\boxed{
m_{\mathrm{rad}}
:=
a_{0,-}G_{\mathrm{rad},0}(T)
-
\overline B^{(0)}(T)
>
0.
}
\]
The sharpened tier below never feeds back into this definition.  Moreover,
using $0\preceq A(\phi)\preceq\Sigma_n$ and hence
$\Gamma_{\mathrm{tar}}(\phi)\le\lambda$,
\[
\boxed{
0<m_{\mathrm{rad}}<\lambda.
}
\]
In particular, $\lambda-m_{\mathrm{rad}}>0$.

Using this fixed coarse radial rate, define
\[
\boxed{
\widehat B(t)
:=
C_B
\left[
e^{\lambda t}
+
\left(
\frac1{\tau_d}-1
\right)
e^{(\lambda-m_{\mathrm{rad}})t}
\right].
}
\]
Distinguish the raw time-integrated perpendicular forcing
\[
\boxed{
\widehat J(t)
:=
\int_0^t\widehat B(s)\,ds
=
C_B
\left[
\mathcal E_\lambda(t)
+
\left(
\frac1{\tau_d}-1
\right)
\mathcal E_{\lambda-m_{\mathrm{rad}}}(t)
\right],
}
\]
from the radius-normalized mismatch budget
\[
\boxed{
\widehat P(t)
:=
C_B
\left[
\mathcal K_{2m_{\mathrm{rad}},\lambda}(t)
+
\left(
\frac1{\tau_d}-1
\right)
\mathcal K_{2m_{\mathrm{rad}},\lambda-m_{\mathrm{rad}}}(t)
\right].
}
\]
Define
\[
\boxed{
\widehat P_T
:=
\sup_{0\le t\le T}\widehat P(t).
}
\]
Since $e^{-2m_{\mathrm{rad}}(t-s)}\le1$,
\[
\boxed{
\widehat P(t)
\le
\widehat J(t)
\le
P_T^{(0)}
\qquad
(0\le t\le T),
}
\]
and hence
\[
\boxed{\widehat P_T\le P_T^{(0)}.}
\]
The $\mathcal K_{2m_{\mathrm{rad}},\cdot}$ kernels are specific to the
radius-normalized mismatch.  The angular Duhamel terms below retain their
own damping kernels and are not to be replaced by these kernels.
Define the sharpened radius-ratio bounds
\[
\boxed{
a_-:=\sqrt{1-2\widehat P_T},
\qquad
a_+:=\sqrt{1+2\widehat P_T}.
}
\]
Let
\[
\widehat\omega_\angle
:=
\omega_\angle(a_-,a_+).
\]
Define the sharpened angular-gap envelope
\[
\boxed{
\begin{aligned}
\widehat D_\angle(T)
&:=
\frac{C_{\angle,0}(a_-,a_+)}{\widehat\omega_\angle}
\left(
1-e^{-\widehat\omega_\angle T}
\right)
\\
&\quad+
C_{\angle,F}(a_-,a_+)
F_0
\mathcal K_{\widehat\omega_\angle,\beta}(T)
\\
&\quad+
C_{\angle,B}C_B
\left[
\mathcal K_{\widehat\omega_\angle,\lambda}(T)
+
\left(
\frac1{\tau_d}-1
\right)
\mathcal K_{\widehat\omega_\angle,\lambda-m_{\mathrm{rad}}}(T)
\right].
\end{aligned}
}
\]
Assume
\[
\boxed{
\widehat D_\angle(T)<\frac\pi2.
}
\]
For brevity below write
\[
D_\angle:=\widehat D_\angle(T).
\]

Define the strong angular contraction rate
\[
\boxed{
\omega_1
:=
a_-\cos(D_\angle)\lambda_{1,\mathrm{ang}}.
}
\]
Set
\[
C_{1,0}
:=
a_+b_{\rm tar}
+a_+\sin(D_\angle)T_{1,+}
+
a_+\frac{X^2}{2n},
\]
\[
C_{1,F}
:=
a_+\sqrt{T_{1,+}}
+
a_+\frac{X}{2n},
\qquad
C_{1,B}:=1+\frac1{2n}.
\]
Define
\[
\boxed{
\begin{aligned}
\Theta_1(t)
&:=
\theta_1e^{-\omega_1t}
+
\frac{C_{1,0}}{\omega_1}
\left(
1-e^{-\omega_1t}
\right)
\\
&\quad+
C_{1,F}F_0\mathcal K_{\omega_1,\beta}(t)
\\
&\quad+
C_{1,B}C_B
\left[
\mathcal K_{\omega_1,\lambda}(t)
+
\left(
\frac1{\tau_d}-1
\right)
\mathcal K_{\omega_1,\lambda-m_{\mathrm{rad}}}(t)
\right].
\end{aligned}
}
\]
Also define
\[
\boxed{
\widehat B_T
:=
\sup_{0\le t\le T}\widehat B(t).
}
\]
Assume the strong-boundary inward condition
\[
\boxed{
\omega_1\theta_1
>
C_{1,0}
+
C_{1,F}\overline F(T)
+
C_{1,B}\widehat B_T.
}
\]

For the weak retention interval set
\[
C_{W,-,0}
:=
a_+
\left(
M_{W,-}
+
\sin(D_\angle)T_{W,+}
\right)
+
a_+\frac{X^2}{2n},
\]
\[
C_{W,+,0}
:=
a_+
\left(
M_{W,+}
+
\sin(D_\angle)T_{W,+}
\right)
+
a_+\frac{X^2}{2n},
\]
\[
C_{W,F}
:=
a_+\sqrt{T_{W,+}}
+
a_+\frac{X}{2n},
\qquad
C_{W,B}:=1+\frac1{2n}.
\]
Define
\[
\boxed{
D_{W,-}(t)
:=
C_{W,-,0}t
+
C_{W,F}F_0\mathcal E_\beta(t)
+
C_{W,B}\widehat J(t),
}
\]
and
\[
\boxed{
D_{W,+}(t)
:=
C_{W,+,0}t
+
C_{W,F}F_0\mathcal E_\beta(t)
+
C_{W,B}\widehat J(t).
}
\]

Call a pair
\[
(\theta_{\mathrm{trap}},T)
\]
\emph{early-admissible relative to the fixed envelope horizon $H$} if
\[
\boxed{
0<\theta_{\mathrm{trap}}<\theta_1,
\qquad
0<T\le H,
}
\]
all preceding smallness and dominance conditions hold, and
\[
\boxed{
\Theta_1(T)<\theta_{\mathrm{trap}},
}
\]
\[
\boxed{
D_{W,-}(T)<b_{W,-},
\qquad
D_{W,+}(T)<b_{W,+}.
}
\]
Assume the early-admissible set is nonempty and fix one such pair.  Define
\[
\boxed{
t_{\mathrm{trap}}:=T.
}
\]

Then for every
\[
i\in A_1^{\mathrm{seed}}\cup W^{\mathrm{seed}}
\]
and every
\[
0\le t\le t_{\mathrm{trap}},
\]
the following hold.

\begin{enumerate}[label=\textnormal{(\roman*)},leftmargin=2.4em]

\item
\textbf{Projected-radius survival and quantitative growth.}
\[
\boxed{
r_i(t)
\ge
r_i(0)e^{m_{\mathrm{rad}}t}
=
\varepsilon R_i e^{m_{\mathrm{rad}}t}
\ge
\varepsilon\tau_d e^{m_{\mathrm{rad}}t}.
}
\]

\item
\textbf{Projected-radius matching.}
\[
\boxed{
\left|
1-\frac{s_i(t)^2}{r_i(t)^2}
\right|
\le
2\widehat P(t)
\le
2\widehat P_T,
}
\]
and hence
\[
\boxed{
a_-
\le
\frac{s_i(t)}{r_i(t)}
\le
a_+.
}
\]

\item
\textbf{Projected input/output angular matching.}
Writing
\[
\zeta_i(t):=\psi_i(t)-\phi_i(t),
\]
one has
\[
\boxed{
|\zeta_i(t)|
\le
D_\angle.
}
\]
Consequently,
\[
\boxed{
\|\widetilde u_i(t)-\widetilde w_i(t)\|
\le
2\sin\frac{D_\angle}{2}.
}
\]
Moreover, for
\[
\delta_i(t)
:=
P_{\mathcal S}(u_i(t)-w_i(t)),
\]
\[
\boxed{
\frac{\|\delta_i(t)\|}{r_i(t)}
\le
\max\{1-a_-,a_+-1\}
+
2\sin\frac{D_\angle}{2}.
}
\]

\item
\textbf{Strong-seed trapping and coarse alignment.}
For every $i\in A_1^{\mathrm{seed}}$,
\[
\boxed{
\phi_i(t)\in J_1^{\mathrm{seed}}
\qquad
\forall t\in[0,t_{\mathrm{trap}}],
}
\]
and
\[
\boxed{
|\phi_i(t_{\mathrm{trap}})|
<
\theta_{\mathrm{trap}}.
}
\]
Therefore
\[
\boxed{
|\psi_i(t_{\mathrm{trap}})|
<
\theta_{\mathrm{trap}}+D_\angle.
}
\]

\item
\textbf{Weak-seed angular retention.}
For every $i\in W^{\mathrm{seed}}$,
\[
\boxed{
\phi_i(t)\in J_W^{\mathrm{ret}}
\qquad
\forall t\in[0,t_{\mathrm{trap}}].
}
\]
More precisely,
\[
\boxed{
\left[
\phi_i(0)-\inf_{0\le s\le t}\phi_i(s)
\right]_+
\le
D_{W,-}(t),
}
\]
and
\[
\boxed{
\left[
\sup_{0\le s\le t}\phi_i(s)-\phi_i(0)
\right]_+
\le
D_{W,+}(t).
}
\]

\end{enumerate}
\end{lemma}

\begin{proof}
Fix an early-admissible pair
$(\theta_{\mathrm{trap}},T)$ and write
$T=t_{\mathrm{trap}}$.

%%%%%%%%%%%%%%%%%%%%%%%%%%%%%%%%%%%%%%%%%%%%%%%%%%%%%%%%%%%%%%%%%%%%%%%%%%%%%%%
\paragraph{Step 1: early-time output and coarse perpendicular envelopes.}
%%%%%%%%%%%%%%%%%%%%%%%%%%%%%%%%%%%%%%%%%%%%%%%%%%%%%%%%%%%%%%%%%%%%%%%%%%%%%%%

Lemma~\ref{lem:norm-output}, evaluated at the common horizon $H$, gives
for $0\le t\le T\le H$,
\[
S(t)
\le
\frac{h\varepsilon^2e^{2\lambda t}}{D_H},
\]
and hence
\[
F(t)
\le
XS(t)
\le
\overline F(t)
=
F_0e^{\beta t}.
\]
Lemma~\ref{lem:projected-polar-dynamics} gives
\[
F_\perp(t)
\le
X\varepsilon\sqrt h\sqrt{S(t)}.
\]
Thus
\[
F_\perp(t)
\le
\frac{Xh\varepsilon^2}{\sqrt{D_H}}
e^{\lambda t}.
\]

On any interval on which
\[
r_i(t)\ge r_i(0)=\varepsilon R_i,
\]
the exact conserved offset gives
\[
\frac{\|w_i(t)\|}{r_i(t)}
=
\frac{\|u_i(t)\|}{r_i(t)}
\le
\sqrt{
1+
\frac{1-R_i^2}{R_i^2}
}
=
\frac1{R_i}
\le
\frac1{\tau_d}.
\]
Consequently, on such an interval,
\[
\mathfrak B_i(t)
\le
\frac{C_B}{\tau_d}e^{\lambda t}
=
\overline B^{(0)}(t).
\]
In particular,
\[
\frac{|R_{i,\perp}^{\mathrm{rad}}|}{r_i}
\le
\overline B^{(0)}(t),
\qquad
|R_{i,\perp}^{\mathrm{ang}}|
\le
\overline B^{(0)}(t).
\]
The dependence of $\overline B^{(0)}$ on the radial-floor guard is precisely
why that guard is included in the continuation argument below.

%%%%%%%%%%%%%%%%%%%%%%%%%%%%%%%%%%%%%%%%%%%%%%%%%%%%%%%%%%%%%%%%%%%%%%%%%%%%%%%
\paragraph{Step 2: exact local-frame identities and gated-mass damping.}
%%%%%%%%%%%%%%%%%%%%%%%%%%%%%%%%%%%%%%%%%%%%%%%%%%%%%%%%%%%%%%%%%%%%%%%%%%%%%%%

Fix angles $\phi,\psi$ and set
\[
\zeta:=\psi-\phi,
\qquad
\nu:=\widetilde u(\phi),
\qquad
\nu^\angle:=\widetilde u^\angle(\phi).
\]
Then
\[
\widetilde u(\psi)
=
\cos\zeta\,\nu
+
\sin\zeta\,\nu^\angle,
\]
and
\[
\widetilde u^\angle(\psi)
=
-\sin\zeta\,\nu
+
\cos\zeta\,\nu^\angle.
\]
Therefore the target-only two-angle fields satisfy
\[
\boxed{
\Gamma_{\mathrm{tar}}(\phi,\psi)
=
\cos\zeta\,\Gamma_{\mathrm{tar}}(\phi)
+
\sin\zeta\,V_{\mathrm{tar}}(\phi),
}
\]
\[
\boxed{
\mathcal V_{\mathrm{tar}}(\phi,\psi)
=
\cos\zeta\,V_{\mathrm{tar}}(\phi)
+
\sin\zeta\,T_{\mathrm{tar}}(\phi),
}
\]
and
\[
\boxed{
\mathcal W_{\mathrm{tar}}(\phi,\psi)
=
-\sin\zeta\,\Gamma_{\mathrm{tar}}(\phi)
+
\cos\zeta\,V_{\mathrm{tar}}(\phi).
}
\]
Moreover,
\[
\boxed{
\Gamma_{\mathrm{tar}}(\phi)
+
T_{\mathrm{tar}}(\phi)
=
\operatorname{tr}_{\mathcal S}A(\phi)
=
\mathbb E_n
\left[
\|x\|^2
\mathbf 1\{\widetilde u(\phi)^\top x>0\}
\right].
}
\]

To prove the explicit lower bound on $\kappa_{\mathrm{gap}}$, set
\[
m_p:=\mu_p-\sigma_pL_{\delta_{\mathrm{strong}}}.
\]
On $\mathcal H_{\mathrm{data}}(\delta)$,
Lemma~\ref{lem:Hstrong} gives
$(x_k^{(p)})_p\ge m_p$, hence
$\|x_k^{(p)}\|^2\ge m_p^2$.  Therefore
\[
\begin{aligned}
\Gamma_{\mathrm{tar}}(\phi)+T_{\mathrm{tar}}(\phi)
&\ge
\frac12\sum_{p=1}^2
m_p^2\widehat P_p(H_\phi)
\\
&\ge
\frac12\sum_{p=1}^2
m_p^2
\left[
P_p(H_\phi)-\Delta_{\mathrm{VC},\delta_{\mathrm{strong}}}
\right]_+.
\end{aligned}
\]
Taking the infimum over $\mathcal J_{\mathrm{early}}$ proves the displayed
bound.

Finally, because the hard-gate indicator lies in $[0,1]$,
\[
0\preceq A(\phi)\preceq\Sigma_n,
\]
so
\[
\Gamma_{\mathrm{tar}}(\phi),
T_{\mathrm{tar}}(\phi)
\le
\overline\lambda_\delta.
\]

%%%%%%%%%%%%%%%%%%%%%%%%%%%%%%%%%%%%%%%%%%%%%%%%%%%%%%%%%%%%%%%%%%%%%%%%%%%%%%%
\paragraph{Step 3: first-exit construction.}
%%%%%%%%%%%%%%%%%%%%%%%%%%%%%%%%%%%%%%%%%%%%%%%%%%%%%%%%%%%%%%%%%%%%%%%%%%%%%%%

For each seed neuron define
\[
\zeta_i(t):=\psi_i(t)-\phi_i(t)
\]
and, whenever $r_i(t)>0$,
\[
Y_i(t)
:=
\frac{|r_i(t)^2-s_i(t)^2|}{r_i(t)^2}.
\]
Let $\tau\le T$ be the first time at which any one of the following occurs:
\begin{enumerate}[label=\textnormal{(\alph*)},leftmargin=2.6em]
\item for some seed neuron,
$r_i(t)<r_i(0)$;
\item for some seed neuron,
$Y_i(t)\ge4P_T^{(0)}$;
\item for some seed neuron,
$|\zeta_i(t)|\ge\pi/2$;
\item for some $i\in A_1^{\mathrm{seed}}$,
$\phi_i(t)\notin J_1^{\mathrm{seed}}$;
\item for some $i\in W^{\mathrm{seed}}$,
$\phi_i(t)\notin J_W^{\mathrm{ret}}$.
\end{enumerate}
At initialization,
\[
r_i(0)=s_i(0)=\varepsilon R_i>0,
\qquad
Y_i(0)=0,
\qquad
\zeta_i(0)=0.
\]
Strong seeds start in $J_1^{\mathrm{seed}}$, and weak seeds start strictly
inside $J_W^{\mathrm{ret}}$.  Hence the continuation interval is nontrivial.
On $[0,\tau)$,
\[
Y_i(t)<4P_T^{(0)},
\]
so
\[
\boxed{
a_{0,-}
<
\frac{s_i(t)}{r_i(t)}
<
a_{0,+}.
}
\]

%%%%%%%%%%%%%%%%%%%%%%%%%%%%%%%%%%%%%%%%%%%%%%%%%%%%%%%%%%%%%%%%%%%%%%%%%%%%%%%
\paragraph{Step 4: coarse angular-gap control.}
%%%%%%%%%%%%%%%%%%%%%%%%%%%%%%%%%%%%%%%%%%%%%%%%%%%%%%%%%%%%%%%%%%%%%%%%%%%%%%%

Write
\[
a_i(t):=\frac{s_i(t)}{r_i(t)}.
\]
From Lemma~\ref{lem:projected-polar-dynamics},
\[
\dot\zeta_i
=
a_i^{-1}\mathcal W_{\mathrm{true}}
-
a_i\mathcal V_{\mathrm{true}}
-
R_{i,\perp}^{\mathrm{ang}}
-
E_i^{\mathrm{bd}}.
\]
Decompose the true fields into target and learned-output parts:
\[
\mathcal V_{\mathrm{true}}
=
\mathcal V_{\mathrm{tar}}-\mathcal O_{V,i},
\qquad
\mathcal W_{\mathrm{true}}
=
\mathcal W_{\mathrm{tar}}-\mathcal O_{W,i},
\]
where
\[
\mathcal O_{V,i}
:=
\mathbb E_n
\left[
\langle P_{\mathcal S}f_t,\widetilde w_i\rangle
\mathbf 1\{\widetilde u_i^\top x>0\}
(\widetilde u_i^\angle)^\top x
\right],
\]
\[
\mathcal O_{W,i}
:=
\mathbb E_n
\left[
\langle P_{\mathcal S}f_t,\widetilde w_i^\angle\rangle
(\widetilde u_i^\top x)_+
\right].
\]
Cauchy--Schwarz yields
\[
|\mathcal O_{V,i}|
\le
F(t)\sqrt{T_{\mathrm{tar}}(\phi_i)},
\qquad
|\mathcal O_{W,i}|
\le
F(t)\sqrt{\Gamma_{\mathrm{tar}}(\phi_i)}.
\]
Using Step~2 gives the exact gap identity
\[
\boxed{
\begin{aligned}
\dot\zeta_i
&=
-\sin\zeta_i
\left(
a_i^{-1}\Gamma_{\mathrm{tar}}(\phi_i)
+
a_iT_{\mathrm{tar}}(\phi_i)
\right)
\\
&\quad+
\cos\zeta_i
\left(
a_i^{-1}-a_i
\right)V_{\mathrm{tar}}(\phi_i)
\\
&\quad-
a_i^{-1}\mathcal O_{W,i}
+
a_i\mathcal O_{V,i}
-
R_{i,\perp}^{\mathrm{ang}}
-
E_i^{\mathrm{bd}}.
\end{aligned}
}
\]
For $|\zeta_i|<\pi/2$,
\[
\operatorname{sgn}(\zeta_i)\sin\zeta_i
\ge
\frac2\pi|\zeta_i|.
\]
Moreover,
\[
a_i^{-1}\Gamma_{\mathrm{tar}}
+
a_iT_{\mathrm{tar}}
\ge
m_a
(\Gamma_{\mathrm{tar}}+T_{\mathrm{tar}}).
\]
The boundary remark gives
\[
|E_{i,\parallel}^{\mathrm{bd}}|
\le
a_i\frac{X}{2n}(X+F(t)),
\qquad
|E_{i,\perp}^{\mathrm{bd}}|
\le
\frac{\overline B^{(0)}(t)}{2n}.
\]
Therefore, on $[0,\tau)$,
\[
\begin{aligned}
D^+|\zeta_i|
&\le
-\omega_\angle(a_{0,-},a_{0,+})|\zeta_i|
+
C_{\angle,0}(a_{0,-},a_{0,+})
\\
&\quad+
C_{\angle,F}(a_{0,-},a_{0,+})\overline F(t)
+
C_{\angle,B}\overline B^{(0)}(t).
\end{aligned}
\]
Since $\zeta_i(0)=0$, scalar comparison gives
\[
|\zeta_i(t)|
\le
D_{\angle,0}(T)
<
\frac\pi2
\qquad
(t<\tau).
\]
Thus event~\textnormal{(c)} cannot be the first exit.

%%%%%%%%%%%%%%%%%%%%%%%%%%%%%%%%%%%%%%%%%%%%%%%%%%%%%%%%%%%%%%%%%%%%%%%%%%%%%%%
\paragraph{Step 5: coarse radial dominance and quantitative growth.}
%%%%%%%%%%%%%%%%%%%%%%%%%%%%%%%%%%%%%%%%%%%%%%%%%%%%%%%%%%%%%%%%%%%%%%%%%%%%%%%

The true radial field has the exact decomposition
\[
\boxed{
\Gamma_{\mathrm{true}}(\phi_i,\psi_i;t)
=
\cos\zeta_i\Gamma_{\mathrm{tar}}(\phi_i)
+
\sin\zeta_iV_{\mathrm{tar}}(\phi_i)
-
\mathcal O_{\Gamma,i},
}
\]
where
\[
\mathcal O_{\Gamma,i}
:=
\mathbb E_n
\left[
\langle P_{\mathcal S}f_t,\widetilde w_i\rangle
(\widetilde u_i^\top x)_+
\right].
\]
Cauchy--Schwarz gives
\[
|\mathcal O_{\Gamma,i}|
\le
F(t)\sqrt{\Gamma_{\mathrm{tar}}(\phi_i)}.
\]
Hence throughout the coarse bootstrap interval,
\[
\Gamma_{\mathrm{true}}(\phi_i,\psi_i;t)
\ge
G_{\mathrm{rad},0}(T).
\]
Using
\[
\dot r_i
=
s_i\Gamma_{\mathrm{true}}
+
R_{i,\perp}^{\mathrm{rad}},
\]
we obtain
\[
\frac{\dot r_i}{r_i}
\ge
m_{\mathrm{rad}}
>
0
\]
for a.e. $t<\tau$.  Therefore
\[
\boxed{
r_i(t)
\ge
r_i(0)e^{m_{\mathrm{rad}}t}
}
\qquad
(t<\tau).
\]
In particular event~\textnormal{(a)} cannot be the first exit.

%%%%%%%%%%%%%%%%%%%%%%%%%%%%%%%%%%%%%%%%%%%%%%%%%%%%%%%%%%%%%%%%%%%%%%%%%%%%%%%
\paragraph{Step 6: sharpened perpendicular envelope and radius-normalized
matching.}
%%%%%%%%%%%%%%%%%%%%%%%%%%%%%%%%%%%%%%%%%%%%%%%%%%%%%%%%%%%%%%%%%%%%%%%%%%%%%%%

Step~5 gives
\[
\frac{\dot r_i(t)}{r_i(t)}
\ge
m_{\mathrm{rad}}>0
\]
for a.e. $t<\tau$.  Since $r_i(t)>0$ and $\log r_i$ is absolutely
continuous on compact subintervals of $[0,\tau)$, integration gives
\[
\boxed{
\frac{r_i(s)}{r_i(t)}
\le
e^{-m_{\mathrm{rad}}(t-s)}
}
\qquad
0\le s\le t<\tau.
\]

The conserved offset identity from
Lemma~\ref{lem:projected-polar-dynamics} gives
\[
\|u_i(t)\|^2
=
r_i(t)^2+\varepsilon^2(1-R_i^2).
\]
Using
\[
r_i(t)
\ge
r_i(0)e^{m_{\mathrm{rad}}t}
=
\varepsilon R_i e^{m_{\mathrm{rad}}t},
\]
we obtain
\[
\frac{\|u_i(t)\|}{r_i(t)}
\le
\sqrt{
1+
\frac{1-R_i^2}{R_i^2}
e^{-2m_{\mathrm{rad}}t}
}.
\]
Set
\[
g_i:=\frac{\sqrt{1-R_i^2}}{R_i}.
\]
Since $x\mapsto\sqrt{1+g_i^2x^2}$ is convex on $[0,1]$,
\[
\sqrt{1+g_i^2x^2}
\le
1+
\left(\sqrt{1+g_i^2}-1\right)x.
\]
Taking $x=e^{-m_{\mathrm{rad}}t}$, using
$\sqrt{1+g_i^2}=1/R_i$, and then exact balance gives
\[
\boxed{
\frac{\|w_i(t)\|}{r_i(t)}
\le
1+
\left(\frac1{R_i}-1\right)e^{-m_{\mathrm{rad}}t}
\le
1+
\left(\frac1{\tau_d}-1\right)e^{-m_{\mathrm{rad}}t}.
}
\]
Because $\|w_{i,\perp}\|\le\|w_i\|$ and
$XF_\perp(t)\le C_Be^{\lambda t}$ by Step~1,
\[
\boxed{
\mathfrak B_i(t)
\le
\widehat B(t)
\le
\overline B^{(0)}(t).
}
\]

Lemma~\ref{lem:projected-polar-dynamics} gives
\[
\Xi_{i,\perp}(0)=0,
\qquad
\dot\Xi_{i,\perp}=2r_iR_{i,\perp}^{\mathrm{rad}},
\]
and
\[
|R_{i,\perp}^{\mathrm{rad}}|
\le
r_i\mathfrak B_i.
\]
Therefore
\[
|\Xi_{i,\perp}(t)|
\le
2\int_0^t r_i(s)^2\widehat B(s)\,ds.
\]
Dividing by $r_i(t)^2$ and applying the radial comparison pointwise in $s$,
\[
Y_i(t)
=
\frac{|\Xi_{i,\perp}(t)|}{r_i(t)^2}
\le
2\int_0^t
e^{-2m_{\mathrm{rad}}(t-s)}\widehat B(s)\,ds.
\]
For $a\in\{\lambda,\lambda-m_{\mathrm{rad}}\}$ one has
$a+2m_{\mathrm{rad}}>0$ and
\[
\int_0^t
e^{-2m_{\mathrm{rad}}(t-s)}e^{as}\,ds
=
\frac{e^{at}-e^{-2m_{\mathrm{rad}}t}}{a+2m_{\mathrm{rad}}}
=
\mathcal K_{2m_{\mathrm{rad}},a}(t).
\]
Hence
\[
\boxed{
Y_i(t)
\le
2\widehat P(t)
\le
2\widehat P_T
\le
2P_T^{(0)}
<
4P_T^{(0)}.
}
\]
Thus event~\textnormal{(b)} cannot be the first exit, and
\[
\boxed{
a_-
\le
\frac{s_i(t)}{r_i(t)}
\le
a_+.
}
\]

The only remaining lossy step in this estimate is
$\|w_{i,\perp}\|\le\|w_i\|$.  The exact identity
\[
\|w_{i,\perp}(t)\|^2
=
\varepsilon^2(1-R_i^2)+r_i(t)^2-s_i(t)^2
\]
instead gives
\[
\frac{\|w_{i,\perp}(t)\|}{r_i(t)}
\le
\frac1{\tau_d}e^{-m_{\mathrm{rad}}t}
+
\sqrt{Y_i(t)}.
\]
We deliberately do not feed this estimate back into the present coarse rate;
it is reserved for the post-trapping continuation, where the inherited
radius mismatch is already quantitatively small.

%%%%%%%%%%%%%%%%%%%%%%%%%%%%%%%%%%%%%%%%%%%%%%%%%%%%%%%%%%%%%%%%%%%%%%%%%%%%%%%
\paragraph{Step 7: sharpened damped angular-gap estimate.}
%%%%%%%%%%%%%%%%%%%%%%%%%%%%%%%%%%%%%%%%%%%%%%%%%%%%%%%%%%%%%%%%%%%%%%%%%%%%%%%

Repeating Step~4 with the sharpened ratio bounds and the sharpened
perpendicular envelope gives
\[
\begin{aligned}
D^+|\zeta_i|
&\le
-\widehat\omega_\angle|\zeta_i|
+
C_{\angle,0}(a_-,a_+)
\\
&\quad+
C_{\angle,F}(a_-,a_+)\overline F(t)
+
C_{\angle,B}\widehat B(t).
\end{aligned}
\]
Since $\zeta_i(0)=0$, Duhamel comparison gives
\[
\boxed{
|\zeta_i(t)|
\le
\widehat D_\angle(T)
=
D_\angle
}
\qquad
(t<\tau).
\]
The damping term is the exact target-only restoring mechanism
\[
-\sin\zeta_i
\left(
a_i^{-1}\Gamma_{\mathrm{tar}}
+
a_iT_{\mathrm{tar}}
\right),
\]
so this estimate does not require a spectral lower bound on
$Q_i^\alpha|_{\mathcal S}$.

%%%%%%%%%%%%%%%%%%%%%%%%%%%%%%%%%%%%%%%%%%%%%%%%%%%%%%%%%%%%%%%%%%%%%%%%%%%%%%%
\paragraph{Step 8: strong angular trapping and coarse convergence.}
%%%%%%%%%%%%%%%%%%%%%%%%%%%%%%%%%%%%%%%%%%%%%%%%%%%%%%%%%%%%%%%%%%%%%%%%%%%%%%%

Fix $i\in A_1^{\mathrm{seed}}$.  The projected angular equation is
\[
\dot\phi_i
=
a_i\mathcal V_{\mathrm{true}}(\phi_i,\psi_i;t)
+
R_{i,\perp}^{\mathrm{ang}}
+
E_i^{\mathrm{bd}}.
\]
Using the exact local-frame identity,
\[
\mathcal V_{\mathrm{true}}
=
\cos\zeta_iV_{\mathrm{tar}}(\phi_i)
+
\sin\zeta_iT_{\mathrm{tar}}(\phi_i)
-
\mathcal O_{V,i}.
\]
The bias-aware target coercivity in Proposition~\ref{sbnew:bias},
with its bias already included in $C_{1,0}$, therefore gives
\[
D^+|\phi_i|
\le
-\omega_1|\phi_i|
+
C_{1,0}
+
C_{1,F}\overline F(t)
+
C_{1,B}\widehat B(t).
\]
At $|\phi_i|=\theta_1$, the right-hand side is at most
\[
-\omega_1\theta_1
+
C_{1,0}
+
C_{1,F}\overline F(T)
+
C_{1,B}\widehat B_T
<0.
\]
Thus the vector field points strictly inward at both endpoints of
$J_1^{\mathrm{seed}}$, so event~\textnormal{(d)} cannot be the first exit.
Duhamel comparison yields
\[
|\phi_i(t)|
\le
\Theta_1(t).
\]
In particular,
\[
|\phi_i(T)|
\le
\Theta_1(T)
<
\theta_{\mathrm{trap}}.
\]
Since $|\psi_i-\phi_i|\le D_\angle$,
\[
|\psi_i(T)|
<
\theta_{\mathrm{trap}}+D_\angle.
\]

%%%%%%%%%%%%%%%%%%%%%%%%%%%%%%%%%%%%%%%%%%%%%%%%%%%%%%%%%%%%%%%%%%%%%%%%%%%%%%%
\paragraph{Step 9: weak-seed retention.}
%%%%%%%%%%%%%%%%%%%%%%%%%%%%%%%%%%%%%%%%%%%%%%%%%%%%%%%%%%%%%%%%%%%%%%%%%%%%%%%

Fix $i\in W^{\mathrm{seed}}$.  On the continuation interval,
$\phi_i(t)\in J_W^{\mathrm{ret}}$, and
\[
\dot\phi_i
=
a_i
\left[
\cos\zeta_iV_{\mathrm{tar}}(\phi_i)
+
\sin\zeta_iT_{\mathrm{tar}}(\phi_i)
-
\mathcal O_{V,i}
\right]
+
R_{i,\perp}^{\mathrm{ang}}
+
E_i^{\mathrm{bd}}.
\]
Taking negative and positive parts gives, respectively,
\[
[-\dot\phi_i]_+
\le
C_{W,-,0}
+
C_{W,F}\overline F(t)
+
C_{W,B}\widehat B(t),
\]
and
\[
[\dot\phi_i]_+
\le
C_{W,+,0}
+
C_{W,F}\overline F(t)
+
C_{W,B}\widehat B(t).
\]
After integration,
\[
\phi_i(0)-\inf_{0\le s\le t}\phi_i(s)
\le
D_{W,-}(t),
\]
\[
\sup_{0\le s\le t}\phi_i(s)-\phi_i(0)
\le
D_{W,+}(t).
\]
Since
\[
D_{W,-}(T)<b_{W,-},
\qquad
D_{W,+}(T)<b_{W,+},
\]
event~\textnormal{(e)} cannot be the first exit.

%%%%%%%%%%%%%%%%%%%%%%%%%%%%%%%%%%%%%%%%%%%%%%%%%%%%%%%%%%%%%%%%%%%%%%%%%%%%%%%
\paragraph{Step 10: closure and Cartesian consequence.}
%%%%%%%%%%%%%%%%%%%%%%%%%%%%%%%%%%%%%%%%%%%%%%%%%%%%%%%%%%%%%%%%%%%%%%%%%%%%%%%

Steps~4--9 rule out every possible first exit before time $T$.  Therefore
\[
\tau=T=t_{\mathrm{trap}}.
\]
All preceding estimates hold on the full interval
$[0,t_{\mathrm{trap}}]$.

The angular matching conclusion follows from
\[
\|\widetilde u_i-\widetilde w_i\|
=
2\left|
\sin\frac{\psi_i-\phi_i}{2}
\right|
\le
2\sin\frac{D_\angle}{2}.
\]
Finally,
\[
\delta_i
=
r_i\widetilde u_i-s_i\widetilde w_i
=
r_i
\left(
\widetilde u_i-a_i\widetilde w_i
\right),
\]
so
\[
\begin{aligned}
\frac{\|\delta_i\|}{r_i}
&\le
\|\widetilde u_i-\widetilde w_i\|
+
|1-a_i|
\\
&\le
2\sin\frac{D_\angle}{2}
+
\max\{1-a_-,a_+-1\}.
\end{aligned}
\]
This proves all conclusions.
\end{proof}

\begin{remark}[Early trapping versus full strong locking]
\label{rem:early-trap-vs-full-lock}
Lemma~\ref{lem:early-projected-trapping} is run at a deterministic
$O(1)$ trapping time
\[
t_{\mathrm{trap}}=T
\]
inside a common general-horizon spectral envelope $H\ge T$ satisfying
\[
\chi(H)\le\overline\chi<1.
\]
Its role is to establish projected-radius survival, input/output matching,
coarse strong trapping, and weak-seed angular retention while keeping the
network output perturbative.  It does not claim that the final strong-locking
scale $\eta_1^u$ has already been reached.

The perpendicular bookkeeping has two distinct integrated quantities.  The
raw time-integrated forcing is
\[
\widehat J(t)
=
C_B\left[
\mathcal E_\lambda(t)
+
(\tau_d^{-1}-1)
\mathcal E_{\lambda-m_{\mathrm{rad}}}(t)
\right],
\]
whereas the radius-normalized mismatch budget is
\[
\widehat P(t)
=
C_B\left[
\mathcal K_{2m_{\mathrm{rad}},\lambda}(t)
+
(\tau_d^{-1}-1)
\mathcal K_{2m_{\mathrm{rad}},\lambda-m_{\mathrm{rad}}}(t)
\right].
\]
The weak-cohort drift budgets use $\widehat J$ because they integrate
$|\dot\phi_i|$ with no radial normalization.  The ratio bounds use
$\widehat P$ because past mismatch is discounted by the quantitative radial
growth.

Since
\[
0\preceq A(\phi)\preceq\Sigma_n,
\]
one has $\Gamma_{\mathrm{tar}}(\phi)\le\lambda$, and the coarse radial rate
satisfies
\[
\boxed{0<m_{\mathrm{rad}}<\lambda.}
\]
Thus $\lambda-m_{\mathrm{rad}}>0$: the transient envelope is genuinely
growing, and neither $\mathcal E_{\lambda-m_{\mathrm{rad}}}$ nor
$\mathcal K_{2m_{\mathrm{rad}},\lambda-m_{\mathrm{rad}}}$ enters a
degenerate branch.

For fixed $O(1)$ trapping time, the transient mismatch contribution scales
schematically as
\[
\widehat P_T^{\mathrm{trans}}
\asymp
X^2h\varepsilon^2\sqrt d.
\]
Thus a target budget $\widehat P_T\lesssim\eta$ is ensured by
\[
\boxed{
\varepsilon
\lesssim
\eta^{1/2}h^{-1/2}d^{-1/4}
}
\]
up to fixed data-geometry constants.  The potential time--dimension trade is
always unfavorable: because $m_{\mathrm{rad}}<\lambda$,
\[
\frac{\lambda}{2(\lambda+m_{\mathrm{rad}})}>\frac14.
\]
Accordingly the proof keeps $T=O(1)$ and claims no optimized dimension
exponent.

The next continuation stage uses the same common spectral envelope through
$t_{\mathrm{lock}}$ and exploits a different effect: mismatch already present
at $t_{\mathrm{trap}}$ is forgotten as
\[
Y_0e^{-2m_{\mathrm{rad}}^{\mathrm{post}}\tau}.
\]
This is the main gain needed for fine locking.  The later weak-basin-creation
stage is deliberately different: it requires the learned residual to become
large enough to change the angular-field root structure.  The perturbative
small-output condition used through $t_{\mathrm{lock}}$ is therefore not a
standing assumption for basin creation.
\end{remark}


%%%%%%%%%%%%%%%%%%%%%%%%%%%%%%%%%%%%%%%%%%%%%%%%%%%%%%%%%%%%%%%%%%%%%%%%%%%%%%%
% T1.9: post-trapping continuation and fine strong locking
%%%%%%%%%%%%%%%%%%%%%%%%%%%%%%%%%%%%%%%%%%%%%%%%%%%%%%%%%%%%%%%%%%%%%%%%%%%%%%%

\begin{lemma}[Post-trapping continuation and fine strong locking]
\label{lem:post-trapping-locking}

Work on the same data and initialization event as
Lemma~\ref{lem:early-projected-trapping}.  Let
\[
t_0:=t_{\mathrm{trap}}
\]
be the coarse trapping time supplied by that lemma.

Fix the primitive post-trapping parameters
\[
\boxed{
\overline\chi\in(0,1),
\qquad
\theta_{\mathrm{post}}\in(\theta_{\mathrm{trap}},\theta_1),
\qquad
\Delta_{\mathrm{post}}>0,
\qquad
Y_{\mathrm g}\in(0,1),
\qquad
\vartheta\in(0,1).
}
\]
Define the common perturbative horizon
\[
\boxed{
H_\star:=t_0+\Delta_{\mathrm{post}}.
}
\]
Assume
\[
\boxed{
\chi(H_\star)\le\overline\chi<1.
}
\]
All constants used in Lemma~\ref{lem:early-projected-trapping} are understood
to have been evaluated using the same common spectral-envelope horizon
$H=H_\star$.

Write
\[
D_\star^{\mathrm{Ric}}:=1-\chi(H_\star),
\]
and define
\[
\boxed{
F_0^\star:=\frac{Xh\varepsilon^2}{D_\star^{\mathrm{Ric}}},
\qquad
C_B^\star:=\frac{X^2h\varepsilon^2}{\sqrt{D_\star^{\mathrm{Ric}}}}.
}
\]
Set
\[
\boxed{
F_\star:=F_0^\star e^{2\lambda H_\star}.
}
\]
Then Lemma~\ref{lem:norm-output} gives
\[
F(t)\le F_\star
\qquad
\forall\,0\le t\le H_\star.
\]

%%%%%%%%%%%%%%%%%%%%%%%%%%%%%%%%%%%%%%%%%%%%%%%%%%%%%%%%%%%%%%%%%%%%%%%%%%%%%%%
% Inherited state
%%%%%%%%%%%%%%%%%%%%%%%%%%%%%%%%%%%%%%%%%%%%%%%%%%%%%%%%%%%%%%%%%%%%%%%%%%%%%%%

Define the inherited deterministic mismatch and angular-gap budgets
\[
\boxed{
Y_0:=2\widehat P(t_0),
\qquad
D_0:=D_\angle.
}
\]
For every
\[
i\in A_1^{\mathrm{seed}}\cup W^{\mathrm{seed}},
\]
Lemma~\ref{lem:early-projected-trapping} gives
\[
\left|1-\frac{s_i(t_0)^2}{r_i(t_0)^2}\right|\le Y_0,
\qquad
|\zeta_i(t_0)|\le D_0,
\]
and
\[
r_i(t_0)\ge\varepsilon R_i e^{m_{\mathrm{rad}}t_0}.
\]
Moreover,
\[
|\phi_i(t_0)|<\theta_{\mathrm{trap}}
\qquad(i\in A_1^{\mathrm{seed}}),
\]
whereas
\[
\phi_i(t_0)\in J_W^{\mathrm{ret}}
\qquad(i\in W^{\mathrm{seed}}).
\]

Set
\[
\boxed{
q_{\perp,0}
:=
\left(\frac1{\tau_d}-1\right)e^{-m_{\mathrm{rad}}t_0}.
}
\]
For the weak cohort define the remaining one-sided buffers
\[
\boxed{
b_{W,-}^{\mathrm{rem}}
:=b_{W,-}-D_{W,-}(t_0),
\qquad
b_{W,+}^{\mathrm{rem}}
:=b_{W,+}-D_{W,+}(t_0).
}
\]
By early admissibility,
\[
b_{W,-}^{\mathrm{rem}}>0,
\qquad
b_{W,+}^{\mathrm{rem}}>0.
\]

%%%%%%%%%%%%%%%%%%%%%%%%%%%%%%%%%%%%%%%%%%%%%%%%%%%%%%%%%%%%%%%%%%%%%%%%%%%%%%%
% Localized target constants
%%%%%%%%%%%%%%%%%%%%%%%%%%%%%%%%%%%%%%%%%%%%%%%%%%%%%%%%%%%%%%%%%%%%%%%%%%%%%%%

Define
\[
\mathcal J_{\mathrm{post}}
:=
[-\theta_{\mathrm{post}},\theta_{\mathrm{post}}]
\cup J_W^{\mathrm{ret}}.
\]
Assume
\[
\mathcal J_{\mathrm{post}}\subseteq\mathcal J_{\mathrm{early}}.
\]
Set
\[
\boxed{
\gamma_{\mathrm{rad}}^{\mathrm{post}}
:=
\inf_{\phi\in\mathcal J_{\mathrm{post}}}\Gamma_{\mathrm{tar}}(\phi),
}
\]
\[
\boxed{
\nu_{\mathrm{post}}
:=
\sup_{\phi\in\mathcal J_{\mathrm{post}}}|V_{\mathrm{tar}}(\phi)|,
}
\]
\[
\boxed{
\Gamma_+^{\mathrm{post}}
:=
\sup_{\phi\in\mathcal J_{\mathrm{post}}}\Gamma_{\mathrm{tar}}(\phi),
\qquad
T_+^{\mathrm{post}}
:=
\sup_{\phi\in\mathcal J_{\mathrm{post}}}T_{\mathrm{tar}}(\phi),
}
\]
and
\[
\boxed{
\kappa_{\mathrm{gap}}^{\mathrm{post}}
:=
\inf_{\phi\in\mathcal J_{\mathrm{post}}}
\left(\Gamma_{\mathrm{tar}}(\phi)+T_{\mathrm{tar}}(\phi)\right).
}
\]
Assume
\[
\gamma_{\mathrm{rad}}^{\mathrm{post}}>0,
\qquad
\kappa_{\mathrm{gap}}^{\mathrm{post}}>0.
\]
For the localized strong cone define
\[
\boxed{
\lambda_{1,\mathrm{ang}}^{\mathrm{post}}
:=
\inf_{|\phi|\le\theta_{\mathrm{post}}}
\left(\Gamma_{\mathrm{tar}}(\phi)-T_{\mathrm{tar}}(\phi)\right)
>0,
}
\]
and
\[
\boxed{
T_{1,+}^{\mathrm{post}}
:=
\sup_{|\phi|\le\theta_{\mathrm{post}}}T_{\mathrm{tar}}(\phi).
}
\]
Localization is monotone in the favorable direction:
\[
\boxed{
\gamma_{\mathrm{rad}}^{\mathrm{post}}\ge\gamma_{\mathrm{rad}},
\quad
\kappa_{\mathrm{gap}}^{\mathrm{post}}\ge\kappa_{\mathrm{gap}},
\quad
\nu_{\mathrm{post}}\le\nu_{\mathrm{early}},
}
\]
\[
\boxed{
\Gamma_+^{\mathrm{post}}\le\Gamma_+,
\quad
T_+^{\mathrm{post}}\le T_+,
\quad
T_{1,+}^{\mathrm{post}}\le T_{1,+}.
}
\]

For $0<a_-\le1\le a_+$ define
\[
m_a(a_-,a_+):=\min\{a_-,a_+^{-1}\},
\]
\[
d_a(a_-,a_+)
:=
\max\{a_-^{-1}-a_-,\,a_+-a_+^{-1}\},
\]
and
\[
\boxed{
\omega_\angle^{\mathrm{post}}(a_-,a_+)
:=
\frac2\pi m_a(a_-,a_+)\kappa_{\mathrm{gap}}^{\mathrm{post}}.
}
\]
Also define
\[
C_{\angle,0}^{\mathrm{post}}(a_-,a_+)
:=
d_a(a_-,a_+)\nu_{\mathrm{post}}+a_+\frac{X^2}{2n},
\]
\[
C_{\angle,F}^{\mathrm{post}}(a_-,a_+)
:=
a_-^{-1}\sqrt{\Gamma_+^{\mathrm{post}}}
+a_+\sqrt{T_+^{\mathrm{post}}}
+a_+\frac{X}{2n},
\]
and
\[
C_{\angle,B}:=1+\frac1{2n}.
\]

%%%%%%%%%%%%%%%%%%%%%%%%%%%%%%%%%%%%%%%%%%%%%%%%%%%%%%%%%%%%%%%%%%%%%%%%%%%%%%%
% Coarse post-trapping tier
%%%%%%%%%%%%%%%%%%%%%%%%%%%%%%%%%%%%%%%%%%%%%%%%%%%%%%%%%%%%%%%%%%%%%%%%%%%%%%%

Define
\[
\boxed{
a_{0,-}^{\mathrm{post}}:=\sqrt{1-Y_{\mathrm g}},
\qquad
a_{0,+}^{\mathrm{post}}:=\sqrt{1+Y_{\mathrm g}}.
}
\]
For $0\le\tau\le\Delta_{\mathrm{post}}$ define
\[
\boxed{
\overline B_{\mathrm{post}}^{(0)}(\tau)
:=
C_B^\star e^{\lambda t_0}(1+q_{\perp,0})e^{\lambda\tau}.
}
\]
Let
\[
\boxed{
B_{\mathrm{post},0}^\star
:=
\overline B_{\mathrm{post}}^{(0)}(\Delta_{\mathrm{post}}).
}
\]
Set
\[
\omega_{\angle,0}^{\mathrm{post}}
:=
\omega_\angle^{\mathrm{post}}(a_{0,-}^{\mathrm{post}},a_{0,+}^{\mathrm{post}}),
\]
and
\[
\begin{aligned}
Q_{\angle,0}^{\mathrm{post}}
:={}&
C_{\angle,0}^{\mathrm{post}}(a_{0,-}^{\mathrm{post}},a_{0,+}^{\mathrm{post}})
\\
&+
C_{\angle,F}^{\mathrm{post}}(a_{0,-}^{\mathrm{post}},a_{0,+}^{\mathrm{post}})
F_\star
+
C_{\angle,B}B_{\mathrm{post},0}^\star.
\end{aligned}
\]
Define
\[
\boxed{
\overline D_{\angle,0}^{\mathrm{post}}
:=
\max\left\{D_0,\frac{Q_{\angle,0}^{\mathrm{post}}}{\omega_{\angle,0}^{\mathrm{post}}}\right\}.
}
\]
Assume
\[
\boxed{
\overline D_{\angle,0}^{\mathrm{post}}<\frac\pi2.
}
\]
This is deliberately a coarse first-exit ceiling, not the refined
post-trapping gap envelope defined below.

Define
\[
\boxed{
G_{\mathrm{rad},0}^{\mathrm{post}}
:=
\cos(\overline D_{\angle,0}^{\mathrm{post}})\gamma_{\mathrm{rad}}^{\mathrm{post}}
-
\sin(\overline D_{\angle,0}^{\mathrm{post}})\nu_{\mathrm{post}}
-
F_\star\sqrt{\Gamma_+^{\mathrm{post}}}.
}
\]
Assume
\[
\boxed{
a_{0,-}^{\mathrm{post}}G_{\mathrm{rad},0}^{\mathrm{post}}
>
B_{\mathrm{post},0}^\star.
}
\]
Freeze
\[
\boxed{
m_{\mathrm{rad}}^{\mathrm{post}}
:=
a_{0,-}^{\mathrm{post}}G_{\mathrm{rad},0}^{\mathrm{post}}
-
B_{\mathrm{post},0}^\star
>0.
}
\]
This rate is defined once from the coarse tier and is never recomputed from
the sharpened quantities below.  Since
$0\preceq A(\phi)\preceq\Sigma_n$ implies
$\Gamma_{\mathrm{tar}}(\phi)\le\lambda$,
\[
\boxed{
0<m_{\mathrm{rad}}^{\mathrm{post}}<\lambda.
}
\]
In particular,
\[
\boxed{\lambda-m_{\mathrm{rad}}^{\mathrm{post}}>0.}
\]

%%%%%%%%%%%%%%%%%%%%%%%%%%%%%%%%%%%%%%%%%%%%%%%%%%%%%%%%%%%%%%%%%%%%%%%%%%%%%%%
% Sharpened radius matching with forgetting
%%%%%%%%%%%%%%%%%%%%%%%%%%%%%%%%%%%%%%%%%%%%%%%%%%%%%%%%%%%%%%%%%%%%%%%%%%%%%%%

Define
\[
\boxed{
\widehat B_{\mathrm{post}}(\tau)
:=
C_B^\star e^{\lambda t_0}
\left[
e^{\lambda\tau}
+
q_{\perp,0}e^{(\lambda-m_{\mathrm{rad}}^{\mathrm{post}})\tau}
\right].
}
\]
Because $\lambda-m_{\mathrm{rad}}^{\mathrm{post}}>0$, this envelope is
increasing.  Set
\[
\boxed{
B_{\mathrm{post}}^\star
:=
\widehat B_{\mathrm{post}}(\Delta_{\mathrm{post}}).
}
\]
Define
\[
\boxed{
\begin{aligned}
\widehat P_{\mathrm{post}}(\tau)
:=
C_B^\star e^{\lambda t_0}
\Big[&
\mathcal K_{2m_{\mathrm{rad}}^{\mathrm{post}},\lambda}(\tau)
\\
&+
q_{\perp,0}
\mathcal K_{2m_{\mathrm{rad}}^{\mathrm{post}},\lambda-m_{\mathrm{rad}}^{\mathrm{post}}}(\tau)
\Big].
\end{aligned}
}
\]
Both exponents are strictly positive, so
$\widehat P_{\mathrm{post}}$ is increasing.  Define
\[
\boxed{
P_{\mathrm{post}}^\star
:=
\widehat P_{\mathrm{post}}(\Delta_{\mathrm{post}}).
}
\]
Assume the primitive mismatch-guard closure condition
\[
\boxed{
Y_0+2P_{\mathrm{post}}^\star<Y_{\mathrm g}.
}
\]
Define
\[
\boxed{
\mathcal Y_{\mathrm{post}}(\tau)
:=
Y_0e^{-2m_{\mathrm{rad}}^{\mathrm{post}}\tau}
+
2\widehat P_{\mathrm{post}}(\tau),
}
\]
and
\[
\boxed{
Y_{\mathrm{post}}^\star
:=
\sup_{0\le\tau\le\Delta_{\mathrm{post}}}\mathcal Y_{\mathrm{post}}(\tau).
}
\]
Then automatically
\[
\boxed{
Y_{\mathrm{post}}^\star
\le
Y_0+2P_{\mathrm{post}}^\star
<Y_{\mathrm g}<1.
}
\]
Define
\[
\boxed{
a_-^{\mathrm{post}}:=\sqrt{1-Y_{\mathrm{post}}^\star},
\qquad
a_+^{\mathrm{post}}:=\sqrt{1+Y_{\mathrm{post}}^\star},
}
\]
and
\[
\boxed{
c_Y^{\mathrm{post}}
:=
\frac1{\sqrt{1-Y_{\mathrm{post}}^\star}}.
}
\]

%%%%%%%%%%%%%%%%%%%%%%%%%%%%%%%%%%%%%%%%%%%%%%%%%%%%%%%%%%%%%%%%%%%%%%%%%%%%%%%
% Refined angular-gap continuation
%%%%%%%%%%%%%%%%%%%%%%%%%%%%%%%%%%%%%%%%%%%%%%%%%%%%%%%%%%%%%%%%%%%%%%%%%%%%%%%

Define
\[
\boxed{
\omega_{\mathrm{gap}}^{\mathrm{post}}
:=
\omega_\angle^{\mathrm{post}}(a_-^{\mathrm{post}},a_+^{\mathrm{post}}).
}
\]
Define the persistent forcing
\[
\boxed{
\begin{aligned}
Q_\infty^{\mathrm{post}}
:={}&
2c_Y^{\mathrm{post}}\nu_{\mathrm{post}}P_{\mathrm{post}}^\star
+
a_+^{\mathrm{post}}\frac{X^2}{2n}
\\
&+
C_{\angle,F}^{\mathrm{post}}(a_-^{\mathrm{post}},a_+^{\mathrm{post}})F_\star
+
C_{\angle,B}B_{\mathrm{post}}^\star.
\end{aligned}
}
\]
Set
\[
\boxed{
D_\infty^{\mathrm{post}}
:=
\frac{Q_\infty^{\mathrm{post}}}{\omega_{\mathrm{gap}}^{\mathrm{post}}}.
}
\]
Define
\[
\boxed{
\begin{aligned}
\mathcal D_{\mathrm{post}}(\tau)
:={}&
D_0e^{-\omega_{\mathrm{gap}}^{\mathrm{post}}\tau}
\\
&+
c_Y^{\mathrm{post}}\nu_{\mathrm{post}}Y_0
\mathcal K_{\omega_{\mathrm{gap}}^{\mathrm{post}},-2m_{\mathrm{rad}}^{\mathrm{post}}}(\tau)
\\
&+
D_\infty^{\mathrm{post}}
\left(1-e^{-\omega_{\mathrm{gap}}^{\mathrm{post}}\tau}\right).
\end{aligned}
}
\]

%%%%%%%%%%%%%%%%%%%%%%%%%%%%%%%%%%%%%%%%%%%%%%%%%%%%%%%%%%%%%%%%%%%%%%%%%%%%%%%
% Exact fine-locking gap and sample-size floor
%%%%%%%%%%%%%%%%%%%%%%%%%%%%%%%%%%%%%%%%%%%%%%%%%%%%%%%%%%%%%%%%%%%%%%%%%%%%%%%

Define the zero-gap persistent remainder
\[
\boxed{
R_{\mathrm{lock}}^{(0)}
:=
a_+^{\mathrm{post}}b_{\rm tar}
+a_+^{\mathrm{post}}\sqrt{T_{1,+}^{\mathrm{post}}}\,F_\star
+
B_{\mathrm{post}}^\star.
}
\]
Define
\[
\boxed{
M_{\mathrm{lock}}
:=
(1-\vartheta)
a_-^{\mathrm{post}}
\lambda_{1,\mathrm{ang}}^{\mathrm{post}}\eta_1^u
-
R_{\mathrm{lock}}^{(0)}.
}
\]
Assume
\[
\boxed{M_{\mathrm{lock}}>0.}
\]
Define the explicit locking-stage sample-size threshold
\[
\boxed{
n_{\mathrm{lock}}
:=
\frac{
a_+^{\mathrm{post}}X^2
+
a_+^{\mathrm{post}}XF_\star
+
B_{\mathrm{post}}^\star
}{2M_{\mathrm{lock}}}.
}
\]
Assume
\[
\boxed{n>n_{\mathrm{lock}}.}
\]
For $D\in[0,\pi/2)$ define
\[
\boxed{
\begin{aligned}
\Psi(D)
:={}&
a_+^{\mathrm{post}}b_{\rm tar}
+a_+^{\mathrm{post}}\sin(D)T_{1,+}^{\mathrm{post}}
+
a_+^{\mathrm{post}}\frac{X^2}{2n}
\\
&+
a_+^{\mathrm{post}}
\left(\sqrt{T_{1,+}^{\mathrm{post}}}+\frac{X}{2n}\right)F_\star
+
\left(1+\frac1{2n}\right)B_{\mathrm{post}}^\star,
\end{aligned}
}
\]
and
\[
\boxed{
\Phi(D)
:=
(1-\vartheta)a_-^{\mathrm{post}}\cos(D)
\lambda_{1,\mathrm{ang}}^{\mathrm{post}}\eta_1^u.
}
\]
The preceding remainder and sample-size conditions are exactly
\[
\boxed{\Psi(0)<\Phi(0).}
\]
Define the exact maximal locking gap
\[
\boxed{
D_{\mathrm{lock}}
:=
\sup\left\{
D\in(0,\pi/2):\Psi(D)\le\Phi(D)
\right\}.
}
\]
Then
\[
\boxed{
\Psi(D_{\mathrm{lock}})
\le
(1-\vartheta)a_-^{\mathrm{post}}\cos(D_{\mathrm{lock}})
\lambda_{1,\mathrm{ang}}^{\mathrm{post}}\eta_1^u.
}
\]
Assume
\[
\boxed{D_\infty^{\mathrm{post}}<D_{\mathrm{lock}}.}
\]
Define
\[
\boxed{
\tau_{\mathrm{gap}}
:=
\inf\left\{
\tau\in[0,\Delta_{\mathrm{post}}]:
\sup_{\tau\le s\le\Delta_{\mathrm{post}}}
\mathcal D_{\mathrm{post}}(s)
\le
D_{\mathrm{lock}}
\right\},
}
\]
with the convention $\inf\varnothing:=+\infty$.  No assumption
$D_{\mathrm{lock}}<D_0$ is imposed.

%%%%%%%%%%%%%%%%%%%%%%%%%%%%%%%%%%%%%%%%%%%%%%%%%%%%%%%%%%%%%%%%%%%%%%%%%%%%%%%
% Strong-cone invariance before fine locking
%%%%%%%%%%%%%%%%%%%%%%%%%%%%%%%%%%%%%%%%%%%%%%%%%%%%%%%%%%%%%%%%%%%%%%%%%%%%%%%

Define
\[
\boxed{
\omega_{1,0}^{\mathrm{post}}
:=
a_{0,-}^{\mathrm{post}}
\cos(\overline D_{\angle,0}^{\mathrm{post}})
\lambda_{1,\mathrm{ang}}^{\mathrm{post}}.
}
\]
Set
\[
\boxed{
\begin{aligned}
Q_{1,0}^{\mathrm{post}}
:={}&
a_{0,+}^{\mathrm{post}}b_{\rm tar}
+a_{0,+}^{\mathrm{post}}
\sin(\overline D_{\angle,0}^{\mathrm{post}})T_{1,+}^{\mathrm{post}}
+
a_{0,+}^{\mathrm{post}}\frac{X^2}{2n}
\\
&+
a_{0,+}^{\mathrm{post}}
\left(\sqrt{T_{1,+}^{\mathrm{post}}}+\frac{X}{2n}\right)F_\star
+
\left(1+\frac1{2n}\right)B_{\mathrm{post},0}^\star.
\end{aligned}
}
\]
Assume
\[
\boxed{
\omega_{1,0}^{\mathrm{post}}\theta_{\mathrm{post}}
>
Q_{1,0}^{\mathrm{post}}.
}
\]
Define
\[
\boxed{
\Theta_1^{\mathrm{post}}(\tau)
:=
\theta_{\mathrm{trap}}e^{-\omega_{1,0}^{\mathrm{post}}\tau}
+
\frac{Q_{1,0}^{\mathrm{post}}}{\omega_{1,0}^{\mathrm{post}}}
\left(1-e^{-\omega_{1,0}^{\mathrm{post}}\tau}\right).
}
\]
Define the time-independent coarse strong-angle ceiling
\[
\boxed{
\theta_{\mathrm{gap}}
:=
\max\left\{
\theta_{\mathrm{trap}},
\frac{Q_{1,0}^{\mathrm{post}}}{\omega_{1,0}^{\mathrm{post}}}
\right\}.
}
\]
Then
\[
\boxed{\theta_{\mathrm{gap}}<\theta_{\mathrm{post}}.}
\]

%%%%%%%%%%%%%%%%%%%%%%%%%%%%%%%%%%%%%%%%%%%%%%%%%%%%%%%%%%%%%%%%%%%%%%%%%%%%%%%
% Weak-cohort retention
%%%%%%%%%%%%%%%%%%%%%%%%%%%%%%%%%%%%%%%%%%%%%%%%%%%%%%%%%%%%%%%%%%%%%%%%%%%%%%%

Define
\[
\boxed{
\begin{aligned}
Q_{W,-}^{\mathrm{post}}
:={}&
a_{0,+}^{\mathrm{post}}
\left[M_{W,-}+\sin(\overline D_{\angle,0}^{\mathrm{post}})T_{W,+}\right]
+
a_{0,+}^{\mathrm{post}}\frac{X^2}{2n}
\\
&+
a_{0,+}^{\mathrm{post}}
\left(\sqrt{T_{W,+}}+\frac{X}{2n}\right)F_\star
+
\left(1+\frac1{2n}\right)B_{\mathrm{post},0}^\star,
\end{aligned}
}
\]
and
\[
\boxed{
\begin{aligned}
Q_{W,+}^{\mathrm{post}}
:={}&
a_{0,+}^{\mathrm{post}}
\left[M_{W,+}+\sin(\overline D_{\angle,0}^{\mathrm{post}})T_{W,+}\right]
+
a_{0,+}^{\mathrm{post}}\frac{X^2}{2n}
\\
&+
a_{0,+}^{\mathrm{post}}
\left(\sqrt{T_{W,+}}+\frac{X}{2n}\right)F_\star
+
\left(1+\frac1{2n}\right)B_{\mathrm{post},0}^\star.
\end{aligned}
}
\]
Assume
\[
\boxed{
\Delta_{\mathrm{post}}Q_{W,-}^{\mathrm{post}}<b_{W,-}^{\mathrm{rem}},
\qquad
\Delta_{\mathrm{post}}Q_{W,+}^{\mathrm{post}}<b_{W,+}^{\mathrm{rem}}.
}
\]

%%%%%%%%%%%%%%%%%%%%%%%%%%%%%%%%%%%%%%%%%%%%%%%%%%%%%%%%%%%%%%%%%%%%%%%%%%%%%%%
% Fine locking and closure
%%%%%%%%%%%%%%%%%%%%%%%%%%%%%%%%%%%%%%%%%%%%%%%%%%%%%%%%%%%%%%%%%%%%%%%%%%%%%%%

Define
\[
\boxed{
\omega_{1,\star}^{\mathrm{post}}
:=
a_-^{\mathrm{post}}\cos(D_{\mathrm{lock}})
\lambda_{1,\mathrm{ang}}^{\mathrm{post}}.
}
\]
By definition of $D_{\mathrm{lock}}$,
\[
\boxed{
\Psi(D_{\mathrm{lock}})
\le
(1-\vartheta)\omega_{1,\star}^{\mathrm{post}}\eta_1^u.
}
\]
Define
\[
\boxed{
\tau_{\mathrm{lock}}
:=
\frac1{\omega_{1,\star}^{\mathrm{post}}}
\left[
\log\left(\frac{\theta_{\mathrm{gap}}}{\vartheta\eta_1^u}\right)
\right]_+.
}
\]
Finally assume the headline closure condition
\[
\boxed{
\tau_{\mathrm{gap}}+\tau_{\mathrm{lock}}
\le
\Delta_{\mathrm{post}}.
}
\]
Because $\inf\varnothing=+\infty$, this condition automatically implies
that the set defining $\tau_{\mathrm{gap}}$ is nonempty.  Define
\[
\boxed{
t_{\mathrm{lock}}
:=
t_0+\tau_{\mathrm{gap}}+\tau_{\mathrm{lock}}.
}
\]

Then the following hold.
\begin{enumerate}[label=\textnormal{(\roman*)},leftmargin=2.4em]
\item \textbf{Post-trapping radial survival.}
For every $i\in A_1^{\mathrm{seed}}\cup W^{\mathrm{seed}}$ and
$0\le\tau\le\Delta_{\mathrm{post}}$,
\[
\boxed{
r_i(t_0+\tau)
\ge
r_i(t_0)e^{m_{\mathrm{rad}}^{\mathrm{post}}\tau}.
}
\]

\item \textbf{Radius-mismatch forgetting.}
For every seed neuron,
\[
\boxed{
\left|1-\frac{s_i(t_0+\tau)^2}{r_i(t_0+\tau)^2}\right|
\le
\mathcal Y_{\mathrm{post}}(\tau)
=
Y_0e^{-2m_{\mathrm{rad}}^{\mathrm{post}}\tau}
+2\widehat P_{\mathrm{post}}(\tau).
}
\]
In particular,
\[
\boxed{
a_-^{\mathrm{post}}
\le
\frac{s_i(t)}{r_i(t)}
\le
a_+^{\mathrm{post}}}
\qquad(t_0\le t\le H_\star).
\]

\item \textbf{Refined projected input/output matching.}
For every seed neuron,
\[
\boxed{
|\zeta_i(t_0+\tau)|\le\mathcal D_{\mathrm{post}}(\tau).
}
\]
Consequently,
\[
\boxed{
|\zeta_i(t)|\le D_{\mathrm{lock}}}
\qquad
\forall t\in[t_0+\tau_{\mathrm{gap}},H_\star].
\]

\item \textbf{Strong-cone persistence and persistent fine locking.}
For every $i\in A_1^{\mathrm{seed}}$,
\[
\boxed{|\phi_i(t)|<\theta_{\mathrm{post}}}
\qquad
\forall t\in[t_0,H_\star],
\]
and
\[
\boxed{|\phi_i(t)|\le\eta_1^u}
\qquad
\forall t\in[t_{\mathrm{lock}},H_\star].
\]
Moreover,
\[
\boxed{|\psi_i(t)|\le\eta_1^u+D_{\mathrm{lock}}}
\qquad
\forall t\in[t_{\mathrm{lock}},H_\star].
\]

\item \textbf{Weak-seed retention through the post-trapping window.}
For every $i\in W^{\mathrm{seed}}$,
\[
\boxed{\phi_i(t)\in J_W^{\mathrm{ret}}}
\qquad
\forall t\in[t_0,H_\star].
\]
\end{enumerate}
\end{lemma}

\begin{proof}
Fix $H_\star=t_0+\Delta_{\mathrm{post}}$ and work throughout with the
general-horizon bounds from Lemma~\ref{lem:norm-output} evaluated at
$H_\star$.  All differential identities below hold for a.e. time along the
fixed absolutely continuous Clarke trajectory.

\paragraph{Step 1: post-trapping first-exit interval.}
Let $\tau_{\mathrm{ex}}\le\Delta_{\mathrm{post}}$ be the first $\tau\ge0$
at which any of the following occurs:
\begin{enumerate}[label=\textnormal{(\alph*)},leftmargin=2.6em]
\item for some seed neuron, $r_i(t_0+\tau)<r_i(t_0)$;
\item for some seed neuron,
\[
Y_i(t_0+\tau)
:=
\left|1-\frac{s_i(t_0+\tau)^2}{r_i(t_0+\tau)^2}\right|
\ge Y_{\mathrm g};
\]
\item for some seed neuron, $|\zeta_i(t_0+\tau)|\ge\pi/2$;
\item for some $i\in A_1^{\mathrm{seed}}$,
$|\phi_i(t_0+\tau)|\ge\theta_{\mathrm{post}}$;
\item for some $i\in W^{\mathrm{seed}}$,
$\phi_i(t_0+\tau)\notin J_W^{\mathrm{ret}}$.
\end{enumerate}
At $\tau=0$, the mismatch guard is strict because
$Y_0+2P_{\mathrm{post}}^\star<Y_{\mathrm g}$, while the remaining guards are
strict by T1.8 and $\theta_{\mathrm{trap}}<\theta_{\mathrm{post}}$.  Hence
the continuation interval is nontrivial.  On
$0\le\tau<\tau_{\mathrm{ex}}$,
\[
a_{0,-}^{\mathrm{post}}
<
\frac{s_i}{r_i}
<
a_{0,+}^{\mathrm{post}}.
\]

\paragraph{Step 2: coarse post-trapping gap envelope.}
On the radial-floor guard,
\[
r_i(t_0+\tau)\ge r_i(t_0)
\ge\varepsilon R_i e^{m_{\mathrm{rad}}t_0}.
\]
The conserved input offset and exact balance give
\[
\frac{\|w_i(t_0+\tau)\|}{r_i(t_0+\tau)}
\le
1+\left(\frac1{R_i}-1\right)e^{-m_{\mathrm{rad}}t_0}
\le
1+q_{\perp,0}.
\]
Together with $XF_\perp(t)\le C_B^\star e^{\lambda t}$,
\[
\mathfrak B_i(t_0+\tau)
\le
\overline B_{\mathrm{post}}^{(0)}(\tau)
\le
B_{\mathrm{post},0}^\star.
\]
The exact gap equation and the same local-frame estimates as in T1.8 give
\[
D^+|\zeta_i|
\le
-\omega_{\angle,0}^{\mathrm{post}}|\zeta_i|
+Q_{\angle,0}^{\mathrm{post}}.
\]
Hence
\[
|\zeta_i(t_0+\tau)|
\le
D_0e^{-\omega_{\angle,0}^{\mathrm{post}}\tau}
+
\frac{Q_{\angle,0}^{\mathrm{post}}}{\omega_{\angle,0}^{\mathrm{post}}}
\left(1-e^{-\omega_{\angle,0}^{\mathrm{post}}\tau}\right)
\le
\overline D_{\angle,0}^{\mathrm{post}}
<\frac\pi2.
\]
Thus the angular-gap guard $|\zeta_i|<\pi/2$ is self-closing on the
continuation interval.  It is retained in the first-exit definition only so
that the gated-trace damping inequality is legitimate up to the putative
first exit; the comparison above then shows that this guard cannot itself
trigger the exit.

\paragraph{Step 3: coarse post-trapping radial growth.}
On the continuation interval,
\[
\Gamma_{\mathrm{true}}
=
\cos\zeta_i\,\Gamma_{\mathrm{tar}}(\phi_i)
+
\sin\zeta_i\,V_{\mathrm{tar}}(\phi_i)
-
\mathcal O_{\Gamma,i},
\]
with
\[
|\mathcal O_{\Gamma,i}|
\le
F_\star\sqrt{\Gamma_+^{\mathrm{post}}}.
\]
Therefore
\[
\Gamma_{\mathrm{true}}
\ge
G_{\mathrm{rad},0}^{\mathrm{post}}.
\]
Using the projected radial equation,
\[
\frac{\dot r_i}{r_i}
\ge
a_{0,-}^{\mathrm{post}}G_{\mathrm{rad},0}^{\mathrm{post}}
-
B_{\mathrm{post},0}^\star
=
m_{\mathrm{rad}}^{\mathrm{post}}.
\]
Thus
\[
\boxed{
r_i(t_0+\tau)
\ge
r_i(t_0)e^{m_{\mathrm{rad}}^{\mathrm{post}}\tau}.
}
\]
The radial-floor guard cannot trigger.

\paragraph{Step 4: radius-normalized forgetting and mismatch closure.}
For $0\le\sigma\le\tau<\tau_{\mathrm{ex}}$,
\[
\frac{r_i(t_0+\sigma)}{r_i(t_0+\tau)}
\le
e^{-m_{\mathrm{rad}}^{\mathrm{post}}(\tau-\sigma)}.
\]
The conserved offset sharpens to
\[
\frac{\|w_i(t_0+\tau)\|}{r_i(t_0+\tau)}
\le
1+q_{\perp,0}e^{-m_{\mathrm{rad}}^{\mathrm{post}}\tau},
\]
so
\[
\mathfrak B_i(t_0+\tau)
\le
\widehat B_{\mathrm{post}}(\tau).
\]
Using $\Xi_{i,\perp}=r_i^2-s_i^2$ and
$\dot\Xi_{i,\perp}=2r_iR_{i,\perp}^{\mathrm{rad}}$,
\[
\begin{aligned}
Y_i(t_0+\tau)
&\le
Y_i(t_0)e^{-2m_{\mathrm{rad}}^{\mathrm{post}}\tau}
\\
&\quad+
2\int_0^\tau
e^{-2m_{\mathrm{rad}}^{\mathrm{post}}(\tau-\sigma)}
\widehat B_{\mathrm{post}}(\sigma)\,d\sigma
\\
&\le
Y_0e^{-2m_{\mathrm{rad}}^{\mathrm{post}}\tau}
+2\widehat P_{\mathrm{post}}(\tau)
=
\mathcal Y_{\mathrm{post}}(\tau).
\end{aligned}
\]
Hence
\[
Y_i(t_0+\tau)
\le
Y_0+2P_{\mathrm{post}}^\star
<
Y_{\mathrm g},
\]
so the mismatch guard cannot trigger.  Moreover,
\[
a_-^{\mathrm{post}}
\le
\frac{s_i}{r_i}
\le
a_+^{\mathrm{post}}.
\]

\paragraph{Step 5: refined gap estimate with inherited-mismatch forgetting.}
Let $a_i:=s_i/r_i$.  Since
\[
|1-a_i^2|=Y_i\le Y_{\mathrm{post}}^\star,
\]
\[
|a_i^{-1}-a_i|
=
\frac{|1-a_i^2|}{a_i}
\le
c_Y^{\mathrm{post}}Y_i.
\]
Step~4 therefore gives
\[
|a_i^{-1}-a_i|
\le
c_Y^{\mathrm{post}}Y_0e^{-2m_{\mathrm{rad}}^{\mathrm{post}}\tau}
+
2c_Y^{\mathrm{post}}\widehat P_{\mathrm{post}}(\tau).
\]
Substituting into the exact gap equation yields
\[
D^+|\zeta_i|
\le
-\omega_{\mathrm{gap}}^{\mathrm{post}}|\zeta_i|
+
c_Y^{\mathrm{post}}\nu_{\mathrm{post}}Y_0
e^{-2m_{\mathrm{rad}}^{\mathrm{post}}\tau}
+
Q_\infty^{\mathrm{post}}.
\]
Duhamel comparison gives
\[
\boxed{
|\zeta_i(t_0+\tau)|
\le
\mathcal D_{\mathrm{post}}(\tau).
}
\]
The inherited mismatch appears only in the decaying convolution and does not
enter the persistent floor $D_\infty^{\mathrm{post}}$.  By definition of
$\tau_{\mathrm{gap}}$,
\[
|\zeta_i(t)|\le D_{\mathrm{lock}}
\qquad
\forall t\in[t_0+\tau_{\mathrm{gap}},H_\star].
\]

\paragraph{Step 6: strong-cone invariance and the coarse strong-angle ceiling.}
Fix $i\in A_1^{\mathrm{seed}}$. The bias-aware inequality of
Proposition~\ref{sbnew:bias} is used, with the bias included in
$Q_{1,0}^{\mathrm{post}}$. On the bootstrap interval,
\[
D^+|\phi_i|
\le
-\omega_{1,0}^{\mathrm{post}}|\phi_i|
+Q_{1,0}^{\mathrm{post}}.
\]
The strong-boundary condition points strictly inward at
$|\phi_i|=\theta_{\mathrm{post}}$.  Since
$|\phi_i(t_0)|<\theta_{\mathrm{trap}}$, scalar comparison gives
\[
\boxed{
|\phi_i(t_0+\tau)|
\le
\Theta_1^{\mathrm{post}}(\tau)
\le
\theta_{\mathrm{gap}}
<
\theta_{\mathrm{post}}.
}
\]
Thus the strong-cone guard cannot trigger.  In particular,
\[
\boxed{
|\phi_i(t_0+\tau_{\mathrm{gap}})|\le\theta_{\mathrm{gap}}.
}
\]

\paragraph{Step 7: weak-seed retention.}
Fix $i\in W^{\mathrm{seed}}$.  The one-sided decomposition from T1.8 gives
\[
[-\dot\phi_i]_+\le Q_{W,-}^{\mathrm{post}},
\qquad
[\dot\phi_i]_+\le Q_{W,+}^{\mathrm{post}}.
\]
Hence the additional left and right excursions over the whole post window
are at most
\[
\Delta_{\mathrm{post}}Q_{W,-}^{\mathrm{post}}
<
b_{W,-}^{\mathrm{rem}},
\qquad
\Delta_{\mathrm{post}}Q_{W,+}^{\mathrm{post}}
<
b_{W,+}^{\mathrm{rem}}.
\]
Thus the weak-retention guard cannot trigger.  All first-exit events have
therefore been excluded through $\Delta_{\mathrm{post}}$.

\paragraph{Step 8: persistent fine strong locking.}
For $t\ge t_0+\tau_{\mathrm{gap}}$, Step~5 gives
$|\zeta_i(t)|\le D_{\mathrm{lock}}$.  Hence
\[
D^+|\phi_i|
\le
-\omega_{1,\star}^{\mathrm{post}}|\phi_i|
+
\Psi(D_{\mathrm{lock}})
\le
-\omega_{1,\star}^{\mathrm{post}}|\phi_i|
+
(1-\vartheta)\omega_{1,\star}^{\mathrm{post}}\eta_1^u.
\]
For every $\tau\ge0$ with
$t_0+\tau_{\mathrm{gap}}+\tau\le H_\star$,
\[
|\phi_i(t_0+\tau_{\mathrm{gap}}+\tau)|
\le
\theta_{\mathrm{gap}}e^{-\omega_{1,\star}^{\mathrm{post}}\tau}
+
(1-\vartheta)\eta_1^u.
\]
For $\tau\ge\tau_{\mathrm{lock}}$ the first term is at most
$\vartheta\eta_1^u$, and therefore
\[
\boxed{
|\phi_i(t)|\le\eta_1^u
\qquad
\forall t\in[t_{\mathrm{lock}},H_\star].
}
\]
Finally,
\[
|\psi_i(t)|
\le
|\phi_i(t)|+|\zeta_i(t)|
\le
\eta_1^u+D_{\mathrm{lock}}
\]
on the same interval.  This proves all conclusions.
\end{proof}

\begin{remark}[Actual versus heuristic locking-gap scale]
\label{rem:actual-locking-gap}
The quantity $D_{\mathrm{lock}}$ is the exact maximal gap compatible with the
fine-locking comparison.  In the idealized zero-remainder limit and as
$\vartheta\downarrow0$, its small-angle scale reduces to
\[
D_{\mathrm{lock}}
\sim
\frac{a_-^{\mathrm{post}}}{a_+^{\mathrm{post}}}
\frac{\lambda_{1,\mathrm{ang}}^{\mathrm{post}}\eta_1^u}
{T_{1,+}^{\mathrm{post}}}.
\]
Thus the earlier heuristic
$\lambda_{1,\mathrm{ang}}(\theta)\eta_1^u/T_{1,+}(\theta)$ is the
zero-remainder, vanishing-$\vartheta$, perfectly matched limit.  The
parameter $\vartheta$ is a genuine tuning parameter: smaller $\vartheta$
buys gap tolerance at the cost of a larger $\log(1/\vartheta)$ contribution
to $\tau_{\mathrm{lock}}$.
\end{remark}

\begin{remark}[Locking-stage sample-size floor]
\label{rem:locking-sample-floor}
The crossing condition $\Psi(0)<\Phi(0)$ contains a finite-sample obstruction
that does not disappear as $\varepsilon\downarrow0$.  When the perturbative
output and perpendicular remainders are lower order, it requires
schematically
\[
\boxed{
n
\gtrsim
\frac{X^2}
{(1-\vartheta)a_-^{\mathrm{post}}
\lambda_{1,\mathrm{ang}}^{\mathrm{post}}\eta_1^u}.
}
\]
The condition $D_\infty^{\mathrm{post}}<D_{\mathrm{lock}}$ imposes a
comparable requirement.  Hence the eventual choice of $\eta_1^u$ in
Theorem~\ref{thm:block-reduction} also determines a locking-stage lower bound
on $n$.
\end{remark}

\begin{remark}[What T1.9 does not yet prove for the weak cohort]
\label{rem:T19-weak-status}
Lemma~\ref{lem:post-trapping-locking} proves angular retention and radial
survival for the fixed weak seed cohort through the post-trapping window, but
it does not yet prove a lower bound on the aggregate projected mass carried
by that cohort.  Such a mass estimate is a separate dependency of the later
weak-basin capture and final-block construction.

The closure condition allows equality, so the present invocation may have
$H_\star=t_{\mathrm{lock}}$.  No later lemma may silently assume a nontrivial
post-lock interval.  If a later stage needs persistence beyond
$t_{\mathrm{lock}}$, the same continuation argument must be invoked with a
larger admissible $\Delta_{\mathrm{post}}$.
\end{remark}

\begin{remark}[Conservatism of the uniform output ceiling]
\label{rem:T19-Fstar-conservative}
The quantity
\[
F_\star=F_0^\star e^{2\lambda H_\star}
\]
is the endpoint value of the Riccati envelope and is used uniformly over
$[0,H_\star]$.  Lemma~\ref{lem:norm-output} also gives
\[
F(t)
\le
\frac{\lambda}{X}
\frac{\overline\chi}{1-\overline\chi}
\frac1{1-e^{-2\lambda H_\star}}
\]
whenever $\chi(H_\star)\le\overline\chi$.  If the gap floor, weak-retention
inequalities, or fine-locking crossing fail only because the uniform output
remainder is too large, the first tightening is to use this
$\overline\chi$-parameterized ceiling and choose $\overline\chi$ smaller.
\end{remark}

%%%%%%%%%%%%%%%%%%%%%%%%%%%%%%%%%%%%%%%%%%%%%%%%%%%%%%%%%%%%%%%%%%%%%%%%%%%%%%%
% T1.10: residual-gated strong-progress package
%%%%%%%%%%%%%%%%%%%%%%%%%%%%%%%%%%%%%%%%%%%%%%%%%%%%%%%%%%%%%%%%%%%%%%%%%%%%%%%

\subsection{Residual-gated strong progress and the inner strong branch}
\label{subsec:T110-progress}

The seed-only lower bound formerly used at this stage is too coarse: its
positive term scales with the fixed strong-seed mass whereas its adverse
term charges the entire non-seed norm mass.  The correct progress variable is
network-level and the correct post-trapping dynamics are residual-gated.  We
therefore rebuild the argument around the exact scalar
$m_1^{\mathrm{eff}}$, the cluster-$1$ cross-output moment $\mathcal X_1$,
and the anisotropy of the residual-gated planar flow.

\paragraph{Network-level strong progress.}
Define
\[
g_1(x):=\mathbf 1\{x_1>0\},
\qquad
\boxed{
\mathsf D_1
:=
\mathbb E_n[g_1(x)x_1^2]
=
\mathbb E_n[(x_1)_+^2].
}
\]
On the master data event,
\[
\boxed{
\mathsf D_1
\ge
\frac12
(\mu_1-\sigma_1L_{\delta_{\mathrm{strong}}})^2
>0.
}
\]
Set
\[
\boxed{
m_1^{\mathrm{eff}}(t)
:=
\frac{
\mathbb E_n[g_1(x)x_1f_{t,1}(x)]
}{\mathsf D_1}.
}
\]
For an angle $\phi$, define
\[
q_\phi(x):=(\widetilde u(\phi)^\top x)_+,
\]
\[
\boxed{
\mathsf H_1(\phi)
:=
\mathbb E_n[(x_1)_+q_\phi(x)],
\qquad
\Gamma_{\mathrm{tar}}(\phi)
:=
\mathbb E_n[q_\phi(x)^2].
}
\]
Then the exact neuronwise decomposition is
\[
\boxed{
m_1^{\mathrm{eff}}(t)
=
\frac1{\mathsf D_1}
\sum_{i=1}^h
r_i(t)z_{i1}(t)\mathsf H_1(\phi_i(t)).
}
\tag{meff-exact}\label{eq:meff-exact}
\]
This identity is used throughout the rebuilt T1.10 package.

%%%%%%%%%%%%%%%%%%%%%%%%%%%%%%%%%%%%%%%%%%%%%%%%%%%%%%%%%%%%%%%%%%%%%%%%%%%%%%%
% T1.10a: scale-free readout barrier
%%%%%%%%%%%%%%%%%%%%%%%%%%%%%%%%%%%%%%%%%%%%%%%%%%%%%%%%%%%%%%%%%%%%%%%%%%%%%%%

\begin{lemma}[Scale-free strong readout barrier]
\label{lem:strong-readout-barrier}
Fix an angle $\phi$ and write
\[
p(x):=(x_1)_+,
\qquad
q_\phi(x):=(\widetilde u(\phi)^\top x)_+.
\]
Define
\[
\alpha_1(\phi)
:=
\frac{\mathsf H_1(\phi)}{\mathsf D_1},
\qquad
q_{\phi,\perp}
:=
q_\phi-\alpha_1(\phi)p,
\]
so that
\[
\langle p,q_{\phi,\perp}\rangle_n=0,
\]
and set
\[
\boxed{
d_{1,\perp}(\phi)^2
:=
\Gamma_{\mathrm{tar}}(\phi)
-
\frac{\mathsf H_1(\phi)^2}{\mathsf D_1}.
}
\]
Let
\[
\psi_t(x):=x_1-f_{t,1}(x).
\]
For every neuron $i$, at a.e. time,
\[
\boxed{
\dot z_{i1}
\ge
r_i
\left[
\mathsf H_1(\phi_i)(1-m_1^{\mathrm{eff}})
-
\|\psi_t\|_n d_{1,\perp}(\phi_i)
\right].
}
\tag{readout-barrier}\label{eq:readout-barrier}
\]
Define the coherence
\[
\boxed{
\kappa_1(\phi)
:=
\frac{
\mathsf H_1(\phi)
}{
\sqrt{\mathsf D_1\Gamma_{\mathrm{tar}}(\phi)}
}.
}
\]
If on an interval
\[
m_1^{\mathrm{eff}}\le m_{\mathrm{prog}}<1,
\qquad
\|\psi_t\|_n\le R_{1,\star},
\]
then $\dot z_{i1}>0$ whenever
\[
\boxed{
\kappa_1(\phi_i)
>
\kappa_{\mathrm{prog}}
:=
\frac{
R_{1,\star}
}{
\sqrt{
R_{1,\star}^2
+
\mathsf D_1(1-m_{\mathrm{prog}})^2
}
}.
}
\]
\end{lemma}

\begin{proof}
The projected output equation gives
\[
\dot z_{i1}
=
r_i\langle\psi_t,q_{\phi_i}\rangle_n.
\]
Using
$q_{\phi_i}=\alpha_1p+q_{\phi_i,\perp}$ and
\[
\langle\psi_t,p\rangle_n
=
\mathsf D_1(1-m_1^{\mathrm{eff}}),
\]
we obtain
\[
\dot z_{i1}
=
r_i
\left[
\mathsf H_1(\phi_i)(1-m_1^{\mathrm{eff}})
+
\langle\psi_t,q_{\phi_i,\perp}\rangle_n
\right].
\]
Cauchy--Schwarz proves \eqref{eq:readout-barrier}.  Solving the resulting
positivity condition in terms of $\kappa_1$ gives the displayed threshold.
\end{proof}

%%%%%%%%%%%%%%%%%%%%%%%%%%%%%%%%%%%%%%%%%%%%%%%%%%%%%%%%%%%%%%%%%%%%%%%%%%%%%%%
% T1.10b: one-sided harmful readout decomposition
%%%%%%%%%%%%%%%%%%%%%%%%%%%%%%%%%%%%%%%%%%%%%%%%%%%%%%%%%%%%%%%%%%%%%%%%%%%%%%%

\begin{lemma}[One-sided harmful readout decomposition]
\label{lem:one-sided-harmful-readout}
Define
\[
\mathsf C_{12}
:=
\mathbb E_n[(x_1)_+|x_2|],
\qquad
\kappa_{12}
:=
\frac{\mathsf C_{12}}{\mathsf D_1}.
\]
For a fixed deterministic neuron set $B\subseteq[h]$, let
\[
\mathcal N_{1,B}
:=
\frac1{\mathsf D_1}
\sum_{j\in B}
r_j[-z_{j1}]_+\mathsf H_1(\phi_j),
\]
\[
\Delta_B
:=
\sum_{j\in B}\|\delta_j\|^2,
\qquad
L_{\parallel,B}
:=
\sum_{j\in B}(r_j^2+s_j^2),
\]
where
\[
\delta_j:=v_j-z_j.
\]
Then, for every $\rho\in(0,1)$,
\[
\boxed{
\mathcal N_{1,B}
\le
\frac1\rho
\sqrt{
\frac{\overline\lambda_\delta}{\mathsf D_1}
}
\Delta_B
+
\frac{\rho+\kappa_{12}}2
L_{\parallel,B}.
}
\tag{harmful-rho}\label{eq:harmful-rho}
\]
Consequently,
\[
\boxed{
\mathcal N_{1,B}
\le
\sqrt{2C_\Delta\Delta_BL_{\parallel,B}}
+
\frac{\kappa_{12}}2L_{\parallel,B},
\qquad
C_\Delta
:=
\sqrt{
\frac{\overline\lambda_\delta}{\mathsf D_1}
}.
}
\tag{harmful-opt}\label{eq:harmful-opt}
\]
If
\[
\varrho_B^2
:=
\frac{\Delta_B}{L_{\parallel,B}},
\]
then
\[
\boxed{
\mathcal N_{1,B}
\le
\left[
\sqrt{2C_\Delta}\,\varrho_B
+
\frac{\kappa_{12}}2
\right]
L_{\parallel,B}.
}
\]
Thus only negative $e_1$ readout is charged as harmful; positive non-seed
readout is retained with its sign.
\end{lemma}

\begin{proof}
The exact representation \eqref{eq:meff-exact} shows that the adverse part of
any subset is obtained by replacing $z_{j1}$ with $-[z_{j1}]_-$.  Since
$z_{j1}=v_{j1}-\delta_{j1}$,
\[
[-z_{j1}]_+
\le
|\delta_{j1}|+[-v_{j1}]_+.
\]
The mismatch term is controlled by Cauchy--Schwarz and
\[
\mathsf H_1(\phi)^2
\le
\mathsf D_1\Gamma_{\mathrm{tar}}(\phi)
\le
\mathsf D_1\overline\lambda_\delta.
\]
For the geometric term, split according to the sign of $\cos\phi$ and use
\[
\mathsf H_1(\phi)
\le
\mathsf D_1[\cos\phi]_+
+
\mathsf C_{12}|\sin\phi|.
\]
Young's inequality
$ab\le a^2/(2\rho)+\rho b^2/2$ and
$2r_js_j\le r_j^2+s_j^2$ yield \eqref{eq:harmful-rho}.  Optimizing in $\rho$
gives \eqref{eq:harmful-opt}.
\end{proof}

%%%%%%%%%%%%%%%%%%%%%%%%%%%%%%%%%%%%%%%%%%%%%%%%%%%%%%%%%%%%%%%%%%%%%%%%%%%%%%%
% T1.10c: residual-gated projected dynamics and mismatch
%%%%%%%%%%%%%%%%%%%%%%%%%%%%%%%%%%%%%%%%%%%%%%%%%%%%%%%%%%%%%%%%%%%%%%%%%%%%%%%

\begin{lemma}[Residual-gated projected dynamics]
\label{lem:residual-gated-mismatch}
For each neuron choose the measurable Clarke selector
$\alpha_i(t;x)\in[0,1]$ from Lemma~\ref{lem:gf-regularity}, and write
\[
f_{\parallel,t}:=P_{\mathcal S}f_t,
\qquad
f_{\perp,t}:=P_{\mathcal S^\perp}f_t.
\]
Define
\[
\boxed{
A_i^\rho(t)
:=
\mathbb E_n
\left[
(x-f_{\parallel,t}(x))x^\top\alpha_i(t;x)
\right],
}
\]
and, for planar vectors $a,b\in\mathcal S$, use
\[
a\times b:=a_1b_2-a_2b_1.
\]
Set
\[
\boxed{
X_i^\rho(t)
:=
\mathbb E_n
\left[
(f_{\parallel,t}(x)\times x)\alpha_i(t;x)
\right],
}
\]
and
\[
\boxed{
P_{i,\perp}^\rho(t)
:=
\mathbb E_n
\left[
\langle f_{\perp,t}(x),w_{i,\perp}(t)\rangle
\alpha_i(t;x)x
\right].
}
\]
Then, for a.e. time,
\[
\boxed{
\dot v_i
=
(A_i^\rho)^\top z_i
-
P_{i,\perp}^\rho,
\qquad
\dot z_i
=
A_i^\rho v_i.
}
\tag{residual-gated-cartesian}\label{eq:residual-gated-cartesian}
\]
Hence, with $\delta_i=v_i-z_i$,
\[
\boxed{
\dot\delta_i
=
-(A_i^\rho)^\top\delta_i
-
r_iX_i^\rho\widetilde u_i^\angle
-
P_{i,\perp}^\rho.
}
\tag{residual-gated-mismatch}\label{eq:residual-gated-mismatch}
\]
Moreover,
\[
\boxed{
\frac12\frac d{dt}\|\delta_i\|^2
=
-\delta_i^\top\operatorname{sym}(A_i^\rho)\delta_i
-r_iX_i^\rho
\langle\delta_i,\widetilde u_i^\angle\rangle
-\langle\delta_i,P_{i,\perp}^\rho\rangle.
}
\tag{mismatch-energy}\label{eq:mismatch-energy}
\]
The exact input-angle equation is
\[
\boxed{
\dot\phi_i
=
V_i^\rho
-
\frac1{r_i}
\delta_i^\top A_i^\rho\widetilde u_i^\angle
-
\frac1{r_i}
\langle P_{i,\perp}^\rho,\widetilde u_i^\angle\rangle,
}
\tag{residual-gated-angle}\label{eq:residual-gated-angle}
\]
where
\[
\boxed{
V_i^\rho
:=
\widetilde u_i^\top A_i^\rho\widetilde u_i^\angle.
}
\]
Finally,
\[
\boxed{
\lambda_{\min}
(\operatorname{sym}A_i^\rho(t))
\ge
-\mu_\rho(t),
\qquad
\mu_\rho(t)
:=
\sqrt{2\overline\lambda_\delta\mathcal L(t)}.
}
\tag{sym-lower}\label{eq:sym-lower}
\]
For a fixed set $B$, writing
\[
D_B:=\sqrt{\Delta_B},
\qquad
\mathsf R_B:=\sum_{i\in B}r_i^2,
\]
one has
\[
\boxed{
D^+D_B
\le
\mu_\rho D_B
+
\left(
\sum_{i\in B}r_i^2|X_i^\rho|^2
\right)^{1/2}
+
\left(
\sum_{i\in B}\|P_{i,\perp}^\rho\|^2
\right)^{1/2}.
}
\tag{mismatch-aggregate}\label{eq:mismatch-aggregate}
\]
\end{lemma}

\begin{proof}
Project the Clarke gradient-flow equations onto $\mathcal S$ and use that
all training inputs lie in $\mathcal S$.  The $z_i$ equation is immediate.
For the $v_i$ equation, the planar residual contributes
$(A_i^\rho)^\top z_i$, while the perpendicular output produces
$P_{i,\perp}^\rho$.  Subtracting the two projected equations gives
\eqref{eq:residual-gated-mismatch}; in two dimensions the antisymmetric part
acts as
$X_i^\rho\mathsf J$, which gives the displayed tangent term.  Taking the
inner product with $\delta_i$ proves \eqref{eq:mismatch-energy}; projecting
$\dot v_i$ onto $\widetilde u_i^\angle/r_i$ proves
\eqref{eq:residual-gated-angle}.

For a unit $q\in\mathcal S$,
\[
q^\top\operatorname{sym}(A_i^\rho)q
=
-\mathbb E_n
\left[
\alpha_i(q^\top x)(q^\top P_{\mathcal S}\rho_t(x))
\right].
\]
Cauchy--Schwarz, the spectral bound on $\Sigma_n$, and
$\mathbb E_n\|\rho_t\|^2=2\mathcal L(t)$ give
\eqref{eq:sym-lower}.  Summing \eqref{eq:mismatch-energy} over $B$ and applying
Cauchy--Schwarz gives \eqref{eq:mismatch-aggregate}.
\end{proof}

%%%%%%%%%%%%%%%%%%%%%%%%%%%%%%%%%%%%%%%%%%%%%%%%%%%%%%%%%%%%%%%%%%%%%%%%%%%%%%%
% T1.10d: pure negative-sector residual reduction
%%%%%%%%%%%%%%%%%%%%%%%%%%%%%%%%%%%%%%%%%%%%%%%%%%%%%%%%%%%%%%%%%%%%%%%%%%%%%%%

\begin{lemma}[Pure negative-sector cluster-$1$ residual reduction]
\label{lem:negative-sector-residual-reduction}
Let
\[
L_s:=L_{\delta_{\mathrm{strong}}},
\qquad
b_1:=\sigma_1L_s,
\qquad
b_2:=\sigma_2L_s,
\]
and
\[
m_1^{\mathrm{data}}:=\mu_1-b_1,
\qquad
m_2^{\mathrm{data}}:=\mu_2-b_2.
\]
Choose
\[
0<\theta_G<\theta_C<\frac\pi2
\]
so that
\[
\boxed{
\cos\theta_C\,m_1^{\mathrm{data}}
-
\sin\theta_C\,b_2
>0,
}
\tag{C1}\label{eq:C1}
\]
and
\[
\boxed{
\sin\theta_G\,m_2^{\mathrm{data}}
-
\cos\theta_G\,b_1
>0.
}
\tag{C2}
\]
Define
\[
\boxed{
J_C^{\mathrm{pure}}
:=
[-\theta_C,-\theta_G].
}
\]
A nonempty such interval exists whenever
\[
\boxed{
b_1b_2
<
m_1^{\mathrm{data}}m_2^{\mathrm{data}}.
}
\tag{pure-sector-SNR}\label{eq:pure-sector-SNR}
\]
In the equal-noise case $\sigma_1=\sigma_2=\sigma$, this condition reduces
exactly to
\[
\boxed{
L_{\delta_{\mathrm{strong}}}
<
\frac{\mu_1\mu_2}{(\mu_1+\mu_2)\sigma}.
}
\tag{pure-sector-SNR-equal}
\]

On $\mathcal H_{\mathrm{strong}}(\delta_{\mathrm{strong}})$, for every
$\phi\in J_C^{\mathrm{pure}}$, every concept-$1$ training sample fires and
every concept-$2$ training sample is off:
\[
\boxed{
\widetilde u(\phi)^\top x_k^{(1)}>0,
\qquad
\widetilde u(\phi)^\top x_k^{(2)}<0
\qquad
\forall k.
}
\tag{pure-gates}\label{eq:pure-gates}
\]
Thus there is no Clarke boundary ambiguity in this sector.

Define the cluster averages
\[
\mathbb E_{p,n}[\varphi]
:=
\frac1n\sum_{k=1}^n\varphi(x_k^{(p)}),
\]
the cluster-$1$ residual matrix
\[
\boxed{
\mathcal R_1(t)
:=
\mathbb E_{1,n}
\left[
(x-P_{\mathcal S}f_t(x))x^\top
\right]
=
\begin{pmatrix}
R_{11}&R_{12}\\
R_{21}&R_{22}
\end{pmatrix},
}
\]
and the cluster-$1$ cross-output moment
\[
\boxed{
\mathcal X_1(t)
:=
\mathbb E_{1,n}
\left[
P_{\mathcal S}f_t(x)\times x
\right]
=
\mathbb E_{1,n}[f_{t,1}x_2-f_{t,2}x_1].
}
\]
For every neuron with $\phi_i\in J_C^{\mathrm{pure}}$,
\[
\boxed{
A_i^\rho
=
\frac12\mathcal R_1,
\qquad
X_i^\rho
=
\frac12\mathcal X_1,
}
\tag{pure-residual-reduction}\label{eq:pure-residual-reduction}
\]
and
\[
\boxed{
P_{i,\perp}^\rho
=
\frac12
\mathbb E_{1,n}
\left[
\langle f_{\perp,t}(x),w_{i,\perp}\rangle x
\right].
}
\]
Moreover,
\[
\boxed{
R_{21}=R_{12}+\mathcal X_1.
}
\tag{cross-identity}\label{eq:cross-identity}
\]

Writing
$c_\phi:=\cos\phi$ and $s_\phi:=\sin\phi$, define
\[
V_C^\rho(\phi;t)
:=
\widetilde u(\phi)^\top
\left(\frac12\mathcal R_1(t)\right)
\widetilde u^\angle(\phi),
\]
\[
G_C^\rho(\phi;t)
:=
\widetilde u(\phi)^\top
\left(\frac12\mathcal R_1(t)\right)
\widetilde u(\phi).
\]
Then for $\phi\in J_C^{\mathrm{pure}}$,
\[
\boxed{
V_C^\rho(\phi;t)
=
\frac12
\left[
|s_\phi c_\phi|(R_{11}-R_{22})
+
R_{12}\cos2\phi
-
\mathcal X_1s_\phi^2
\right],
}
\tag{pure-angular-field}\label{eq:pure-angular-field}
\]
and
\[
\boxed{
G_C^\rho(\phi;t)
=
\frac12
\left[
R_{11}c_\phi^2
+
R_{22}s_\phi^2
+
(2R_{12}+\mathcal X_1)s_\phi c_\phi
\right].
}
\tag{pure-radial-field}\label{eq:pure-radial-field}
\]

Suppose on an interval beginning at $t_a$,
\[
m_1^{\mathrm{eff}}(t)\le m_{\mathrm{prog}},
\qquad
\mathcal L(t)\le\mathcal L_\star:=\mathcal L(t_a).
\]
Define
\[
S_{11}^{(1)}:=\mathbb E_{1,n}[x_1^2],
\quad
S_{22}^{(1)}:=\mathbb E_{1,n}[x_2^2],
\quad
S_{11}^{(2)}:=\mathbb E_{2,n}[x_1^2],
\]
\[
Q_{2,+}:=\mathbb E_{2,n}[(x_1)_+^2].
\]
Then
\[
\boxed{
R_{11}
\ge
\underline R_{11}
:=
(1-m_{\mathrm{prog}})S_{11}^{(1)}
-
m_{\mathrm{prog}}Q_{2,+}
-
\sqrt{Q_{2,+}}
\left(
2\sqrt{\mathcal L_\star}
+
\sqrt{S_{11}^{(2)}}
\right).
}
\tag{R11-lower}\label{eq:R11-lower}
\]
The loss-only off-diagonal bounds
\[
\boxed{
|R_{12}|,|R_{22}|
\le
\varepsilon_{\mathrm{off}}
:=
2\sqrt{\mathcal L_\star S_{22}^{(1)}}
}
\tag{coarse-offdiag}\label{eq:coarse-offdiag}
\]
are valid but deliberately coarse; the structured off-diagonal bound below
is used for the inner-region cone.

If
\[
|\mathcal X_1|\le X_C,
\]
define
\[
\boxed{
\gamma_{\mathrm{ang}}^{\mathrm{pure}}
:=
\frac12
\inf_{\phi\in J_C^{\mathrm{pure}}}
\left\{
|s_\phi c_\phi|
(\underline R_{11}-\varepsilon_{\mathrm{off}})
-
\varepsilon_{\mathrm{off}}|\cos2\phi|
-
X_Cs_\phi^2
\right\},
}
\]
and
\[
\boxed{
\gamma_{\mathrm{rad}}^{\mathrm{pure}}
:=
\frac12
\inf_{\phi\in J_C^{\mathrm{pure}}}
\left\{
\underline R_{11}c_\phi^2
-
\varepsilon_{\mathrm{off}}s_\phi^2
-
(2\varepsilon_{\mathrm{off}}+X_C)|s_\phi c_\phi|
\right\}.
}
\]
Positive margins give angular and radial coercivity on the pure sector.
The outer angle $\theta_C$ is chosen against the radial margin, which is the
binding one in the canonical sizing.
\end{lemma}

\begin{proof}
For concept $1$, the coordinate envelope gives
\[
x_1\ge m_1^{\mathrm{data}},
\qquad
|x_2|\le b_2.
\]
For $\phi\in[-\theta_C,-\theta_G]$,
\[
\widetilde u(\phi)^\top x^{(1)}
\ge
\cos\theta_C\,m_1^{\mathrm{data}}
-
\sin\theta_C\,b_2>0.
\]
For concept $2$,
\[
|x_1|\le b_1,
\qquad
x_2\ge m_2^{\mathrm{data}},
\]
so
\[
\widetilde u(\phi)^\top x^{(2)}
\le
\cos\theta_G\,b_1
-
\sin\theta_G\,m_2^{\mathrm{data}}<0.
\]
This proves \eqref{eq:pure-gates}.  Solving the two strict inequalities for a
nonempty interval gives \eqref{eq:pure-sector-SNR}; the equal-noise
simplification follows by cancellation of the $\sigma^2L_s^2$ terms.

Since the selector is exactly one on concept $1$ and zero on concept $2$,
\eqref{eq:pure-residual-reduction} follows from the definition of
$A_i^\rho$.  The antisymmetric part of $\mathcal R_1$ satisfies
\[
R_{21}-R_{12}
=
\mathbb E_{1,n}[f_{t,1}x_2-f_{t,2}x_1]
=
\mathcal X_1,
\]
which is \eqref{eq:cross-identity}.  Direct expansion in the frame
$(\widetilde u(\phi),\widetilde u^\angle(\phi))$ gives
\eqref{eq:pure-angular-field} and \eqref{eq:pure-radial-field}.

The lower bound \eqref{eq:R11-lower} follows by writing the definition of
$m_1^{\mathrm{eff}}$ separately over the two clusters and controlling the
concept-$2$ correction by Cauchy--Schwarz and loss monotonicity.  The bounds
\eqref{eq:coarse-offdiag} are the corresponding Cauchy--Schwarz estimates on
cluster $1$.  Substitution into the exact frame identities gives the final
two margins.
\end{proof}

\begin{remark}[Pure-sector width and confidence]
The condition \eqref{eq:pure-sector-SNR} is unusual in this proof: making the
uniform coordinate event more stringent increases $L_s$ and can eventually
destroy the pure sector.  Thus the final admissible region contains both a
lower sample-size requirement from concentration and an upper sample-size
requirement from the high-probability coordinate SNR condition.
\end{remark}

%%%%%%%%%%%%%%%%%%%%%%%%%%%%%%%%%%%%%%%%%%%%%%%%%%%%%%%%%%%%%%%%%%%%%%%%%%%%%%%
% T1.10e: angle-resolved whole-network certificate
%%%%%%%%%%%%%%%%%%%%%%%%%%%%%%%%%%%%%%%%%%%%%%%%%%%%%%%%%%%%%%%%%%%%%%%%%%%%%%%

\begin{lemma}[Angle-resolved strong progress and cluster-$1$ self-field]
\label{lem:angle-resolved-progress-self-field}
For the concept-$1$ empirical distribution define
\[
\widehat A_1(\phi)
:=
\mathbb E_{1,n}
\left[
xx^\top
\mathbf 1\{
\widetilde u(\phi)^\top x>0
\}
\right],
\]
\[
\Gamma_1(\phi)
:=
\widetilde u(\phi)^\top\widehat A_1(\phi)\widetilde u(\phi),
\]
and
\[
\boxed{
V_1(\phi)
:=
\widetilde u(\phi)^\top
\widehat A_1(\phi)
\widetilde u^\angle(\phi).
}
\]
Let
\[
\lambda_{1,+}
:=
\left\|
\mathbb E_{1,n}[xx^\top]
\right\|_{\mathrm{op}}.
\]
Then
\[
0\le\Gamma_1(\phi)\le\lambda_{1,+},
\qquad
|V_1(\phi)|\le\lambda_{1,+}.
\]
Define the progress kernel
\[
\boxed{
p_1(\phi)
:=
\cos\phi\,
\frac{\mathsf H_1(\phi)}{\mathsf D_1}.
}
\]
For a fixed set $B$, define
\[
\mathsf R_B:=\sum_{i\in B}r_i^2,
\qquad
\Delta_B:=\sum_{i\in B}\|\delta_i\|^2,
\]
and the radius-weighted angular measure
\[
\boxed{
\nu_B
:=
\sum_{i\in B}r_i^2\delta_{\phi_i}.
}
\]
Then
\[
\boxed{
\mathcal X_{1,B}
:=
\mathbb E_{1,n}
\left[
f_{t,1}^{B}(x)x_2-f_{t,2}^{B}(x)x_1
\right]
=
\int V_1(\phi)\,d\nu_B(\phi)
+
\mathcal E_{X,B},
}
\tag{X-measure}\label{eq:X-measure}
\]
where
\[
\boxed{
|\mathcal E_{X,B}|
\le
2\lambda_{1,+}
\sqrt{\mathsf R_B\Delta_B}.
}
\tag{X-measure-error}
\]
In particular,
\[
\boxed{
\mathcal X_{1,B}
\le
\int[V_1(\phi)]_+\,d\nu_B(\phi)
+
2\lambda_{1,+}\sqrt{\mathsf R_B\Delta_B}.
}
\tag{X-one-sided-measure}\label{eq:X-one-sided-measure}
\]
Thus directions with $V_1(\phi)<0$ help the self-field budget and are not
charged by absolute value.

Likewise the contribution of $B$ to network-level progress satisfies
\[
\boxed{
m_{1,B}^{\mathrm{eff}}
=
\int p_1(\phi)\,d\nu_B(\phi)
+
\mathcal E_{m,B},
}
\tag{m-measure}
\]
with
\[
\boxed{
|\mathcal E_{m,B}|
\le
C_{H,1}\sqrt{\mathsf R_B\Delta_B},
\qquad
C_{H,1}
:=
\sqrt{
\frac{\overline\lambda_\delta}{\mathsf D_1}
}.
}
\tag{m-measure-error}\label{eq:m-measure-error}
\]
For the whole network, with
\[
\nu_t:=\sum_{i=1}^h r_i(t)^2\delta_{\phi_i(t)},
\qquad
\mathsf R(t):=\sum_i r_i(t)^2,
\qquad
\Delta(t):=\sum_i\|\delta_i(t)\|^2,
\]
we therefore have the terminal sufficient conditions
\[
\boxed{
\int p_1\,d\nu_t
-
C_{H,1}\sqrt{\mathsf R(t)\Delta(t)}
\ge
m_{\mathrm{prog}},
}
\tag{terminal-progress-certificate}\label{eq:terminal-progress-certificate}
\]
and
\[
\boxed{
\int[V_1]_+\,d\nu_t
+
2\lambda_{1,+}\sqrt{\mathsf R(t)\Delta(t)}
\le
X_C.
}
\tag{terminal-self-field-certificate}\label{eq:terminal-self-field-certificate}
\]
No seed/contributor/remainder decomposition appears in these terminal
certificates.

Finally, on concept-$1$ data
$x_1=\mu_1+\xi_1$ and $x_2=\xi_2$, so
\[
\boxed{
\mathcal X_1
=
-\mu_1\mathbb E_{1,n}[f_{t,2}]
+
\mathcal R_X,
}
\]
where
\[
\mathcal R_X
:=
\mathbb E_{1,n}[f_{t,1}\xi_2]
-
\mathbb E_{1,n}[f_{t,2}\xi_1].
\]
Thus $\mathcal X_1$ is, up to a noise-suppressed correction, the network's
spurious $e_2$ output on concept-$1$ data.
\end{lemma}

\begin{proof}
For one neuron, write $\zeta_i:=\psi_i-\phi_i$.  The exact cross-output
contribution is
\[
\mathcal X_{1,i}
=
r_is_i
\left[
-\sin\zeta_i\,\Gamma_1(\phi_i)
+\cos\zeta_i\,V_1(\phi_i)
\right].
\]
Since
\[
r_is_i\sin\zeta_i
=v_i\times z_i
=-v_i\times\delta_i,
\]
and
\[
r_is_i\cos\zeta_i
=v_i^\top z_i
=r_i^2-v_i^\top\delta_i,
\]
we obtain
\[
\mathcal X_{1,i}
=
r_i^2V_1(\phi_i)
+
\Gamma_1(\phi_i)(v_i\times\delta_i)
-
V_1(\phi_i)v_i^\top\delta_i.
\]
The error is bounded by
$2\lambda_{1,+}r_i\|\delta_i\|$; summing and applying
Cauchy--Schwarz gives \eqref{eq:X-measure}--\eqref{eq:X-one-sided-measure}.

For progress,
\[
z_{i1}
=r_i\cos\phi_i-\delta_{i1},
\]
so \eqref{eq:meff-exact} becomes
\[
m_{1,B}^{\mathrm{eff}}
=
\sum_{i\in B}r_i^2p_1(\phi_i)
-
\sum_{i\in B}
r_i\delta_{i1}
\frac{\mathsf H_1(\phi_i)}{\mathsf D_1}.
\]
Since
\[
\mathsf H_1(\phi)^2
\le
\mathsf D_1\Gamma_{\mathrm{tar}}(\phi)
\le
\mathsf D_1\overline\lambda_\delta,
\]
Cauchy--Schwarz gives \eqref{eq:m-measure-error}.  The spurious-output
interpretation follows by substituting $x_1=\mu_1+\xi_1$ and $x_2=\xi_2$
into the definition of $\mathcal X_1$.
\end{proof}

\begin{remark}[High-SNR sign geometry of the self-field kernel]
For equal coordinate noise, the population leading term of $V_1$ in the
interior of the concept-$1$ active half-plane is
\[
V_1(\phi)
\approx
-\mu_1^2\sin\phi\cos\phi.
\]
Thus the dominant harmful region is the fourth quadrant
$(-\pi/2,0)$, while positive-angle mass contributes with the opposite sign.
At the exact outer boundary $\phi=-\pi/2$, however,
\[
V_1(-\pi/2)
=
\mathbb E_1[(-x_2)_+x_1]
=
\frac{\mu_1\sigma_2}{\sqrt{2\pi}}
\]
in the equal-noise population model, so the correct theorem-level object is
always $[V_1]_+$ rather than a hard quadrant indicator.
\end{remark}

%%%%%%%%%%%%%%%%%%%%%%%%%%%%%%%%%%%%%%%%%%%%%%%%%%%%%%%%%%%%%%%%%%%%%%%%%%%%%%%
% T1.10f: structured cluster-1 off-diagonal residual
%%%%%%%%%%%%%%%%%%%%%%%%%%%%%%%%%%%%%%%%%%%%%%%%%%%%%%%%%%%%%%%%%%%%%%%%%%%%%%%

\begin{lemma}[Structured cluster-$1$ off-diagonal residual decomposition]
\label{lem:structured-cluster1-offdiag}
At a fixed time, let
\[
\mathcal A_1(t)
:=
\left\{
i:\
\widetilde u_i(t)^\top x_k^{(1)}>0
\ \forall k
\right\}
\]
be the analytical set of neurons fully active on concept $1$.  For
$q\in\{1,2\}$ define
\[
\boxed{
a_t^{(q)}
:=
\sum_{i\in\mathcal A_1(t)}
r_i z_{iq}\widetilde u_i,
}
\]
and on concept-$1$ inputs write
\[
\boxed{
f_{t,q}(x)
=
a_t^{(q)\top}x
+
g_{t,q}(x),
}
\]
where
\[
g_{t,q}(x)
:=
\sum_{i\notin\mathcal A_1(t)}
r_iz_{iq}(\widetilde u_i^\top x)_+.
\]
Then, with
\[
\mathsf s_{rs}
:=
\mathbb E_{1,n}[x_rx_s],
\]
one has the exact identities
\[
\boxed{
R_{12}
=
(1-a_{t,1}^{(1)})\mathsf s_{12}
-
a_{t,2}^{(1)}\mathsf s_{22}
-
\mathbb E_{1,n}[g_{t,1}(x)x_2],
}
\tag{R12-structured}\label{eq:R12-structured}
\]
and
\[
\boxed{
R_{22}
=
\mathsf s_{22}
-
a_{t,1}^{(2)}\mathsf s_{12}
-
a_{t,2}^{(2)}\mathsf s_{22}
-
\mathbb E_{1,n}[g_{t,2}(x)x_2].
}
\tag{R22-structured}\label{eq:R22-structured}
\]

Let $\mathcal P_1(t)$ denote the neurons partially gated on concept $1$;
fully-off neurons contribute zero to $g_{t,q}$.  Define
\[
\mathsf M_{\mathrm{band}}^\times(t)
:=
\sum_{i\in\mathcal P_1(t)}r_i(t)s_i(t),
\]
and
\[
\Gamma_{1,\mathrm{band}}(t)
:=
\sup_{i\in\mathcal P_1(t)}\Gamma_1(\phi_i(t)).
\]
Then for $q\in\{1,2\}$,
\[
\boxed{
\left|
\mathbb E_{1,n}[g_{t,q}(x)x_2]
\right|
\le
\sqrt{\mathsf s_{22}}
\sqrt{\Gamma_{1,\mathrm{band}}(t)}
\mathsf M_{\mathrm{band}}^\times(t).
}
\tag{partial-tail}\label{eq:partial-tail}
\]
In particular, the crude bound
\[
\mathsf M_{\mathrm{band}}^\times
\le
\sum_i r_is_i
\]
is sufficient; no separate band-mass lemma is required.

On the coordinate-envelope event, the concept-$1$ partially-gated band is
contained in the deterministic angular region where
\[
\mu_1\cos\phi-|\cos\phi|b_1-|\sin\phi|b_2
\le0\le
\mu_1\cos\phi+|\cos\phi|b_1+|\sin\phi|b_2.
\]
Any strictly feasible $\theta_C$ in \eqref{eq:C1} is a buffered
purity edge; it need not equal the exact partial-band boundary.  Lemma~\ref{lem:gated-moment-concentration} controls
$\Gamma_{1,\mathrm{band}}$ uniformly on this moving partially-gated region.
\end{lemma}

\begin{proof}
For fully active neurons, ReLU is exactly linear on every concept-$1$ sample,
which gives the decomposition $f_{t,q}=a_t^{(q)\top}x+g_{t,q}$.  Since
\[
\mathbb E_{1,n}[xx_2]
=
(\mathsf s_{12},\mathsf s_{22})^\top,
\]
substitution into
$R_{q2}=\mathbb E_{1,n}[(x_q-f_{t,q})x_2]$ proves
\eqref{eq:R12-structured} and \eqref{eq:R22-structured}.

A fully-off neuron has $(\widetilde u_i^\top x)_+=0$ on concept $1$.  For a
partially gated neuron, Cauchy--Schwarz gives
\[
\left|
\mathbb E_{1,n}
[(\widetilde u_i^\top x)_+x_2]
\right|
\le
\sqrt{\Gamma_1(\phi_i)\mathsf s_{22}}.
\]
Using $|z_{iq}|\le s_i$ and summing gives \eqref{eq:partial-tail}.  The angular
band description follows directly from the coordinate envelope.
\end{proof}

\begin{remark}[Why the structured off-diagonal estimate matters]
The coarse loss bound
$|R_{12}|\le2\sqrt{\mathcal L_\star\mathsf s_{22}}$ discards the fact that
$x_2$ is noise on concept $1$.  In \eqref{eq:R12-structured}, the fully-active
bulk pays only the $\mathsf s_{22}$ scale, while the partially-gated tail is
suppressed both by its narrow angular support and by
$\sqrt{\Gamma_{1,\mathrm{band}}}$.  This is the high-leverage estimate used
to define the inner-region constant $\beta_{12}^{\mathrm{in}}$.
\end{remark}

%%%%%%%%%%%%%%%%%%%%%%%%%%%%%%%%%%%%%%%%%%%%%%%%%%%%%%%%%%%%%%%%%%%%%%%%%%%%%%%
% T1.10g: pure-sector anisotropic amplification
%%%%%%%%%%%%%%%%%%%%%%%%%%%%%%%%%%%%%%%%%%%%%%%%%%%%%%%%%%%%%%%%%%%%%%%%%%%%%%%

\begin{lemma}[Pure-sector transverse control by anisotropic amplification]
\label{lem:pure-sector-anisotropic-amplification}
Let $B\subseteq[h]$ be a fixed set such that
\[
\phi_i(t)\in J_C^{\mathrm{pure}}
\qquad
\forall i\in B,
\quad t\in I=[t_a,t_b].
\]
Write
\[
v_i=a_ie_1+b_ie_2,
\]
and define
\[
\boxed{
M_B:=\sum_{i\in B}a_i^2,
\qquad
Q_B:=\sum_{i\in B}b_i^2,
\qquad
C_B:=\sum_{i\in B}a_ib_i.
}
\]
Thus
\[
M_B=\int\cos^2\phi\,d\nu_B,
\qquad
Q_B=\int\sin^2\phi\,d\nu_B,
\qquad
|C_B|\le\sqrt{M_BQ_B}.
\]
Let
\[
\mathcal P_B
:=
\left(
\sum_{i\in B}\|P_{i,\perp}^\rho\|^2
\right)^{1/2}.
\]
Then, for a.e. $t\in I$,
\[
\boxed{
\dot a_i
=
\frac12
\left[
R_{11}(a_i-\delta_{i1})
+
R_{21}(b_i-\delta_{i2})
\right]
-
P_{i1}^\rho,
}
\]
and
\[
\boxed{
\dot b_i
=
\frac12
\left[
R_{12}(a_i-\delta_{i1})
+
R_{22}(b_i-\delta_{i2})
\right]
-
P_{i2}^\rho.
}
\]
Consequently,
\[
\boxed{
\begin{aligned}
D_-\sqrt{M_B}
&\ge
\frac12R_{11}\sqrt{M_B}
-
\frac12|R_{21}|\sqrt{Q_B}
\\
&\quad
-
\frac12\sqrt{R_{11}^2+R_{21}^2}\sqrt{\Delta_B}
-
\mathcal P_B,
\end{aligned}
}
\tag{M-sqrt-lower}\label{eq:M-sqrt-lower}
\]
and
\[
\boxed{
\begin{aligned}
D^+\sqrt{Q_B}
&\le
\frac12|R_{12}|\sqrt{M_B}
+
\frac12|R_{22}|\sqrt{Q_B}
\\
&\quad
+
\frac12\sqrt{R_{12}^2+R_{22}^2}\sqrt{\Delta_B}
+
\mathcal P_B.
\end{aligned}
}
\tag{Q-sqrt-upper}\label{eq:Q-sqrt-upper}
\]

Whenever $M_B>0$, define
\[
\boxed{
y_B:=\sqrt{Q_B/M_B}.}
\]
Then
\[
\boxed{
\begin{aligned}
D^+y_B
&\le
\frac12|R_{12}|
-
\frac12(R_{11}-|R_{22}|)y_B
+
\frac12|R_{21}|y_B^2
\\
&\quad
+
\frac12
\left[
\sqrt{R_{12}^2+R_{22}^2}
+
y_B\sqrt{R_{11}^2+R_{21}^2}
\right]
\frac{\sqrt{\Delta_B}}{\sqrt{M_B}}
\\
&\quad
+(1+y_B)
\frac{\mathcal P_B}{\sqrt{M_B}}.
\end{aligned}
}
\tag{pure-ratio}\label{eq:pure-ratio}
\]
Suppose on $I$,
\[
R_{11}\ge\underline R_{11}>0,
\quad
|R_{12}|\le\beta_{12},
\quad
|R_{22}|\le\beta_{22},
\quad
|R_{21}|\le\beta_{21},
\]
\[
R_{11}\le\overline R_{11},
\qquad
\frac{\sqrt{\Delta_B}}{\sqrt{M_B}}\le\rho_\delta,
\qquad
\frac{\mathcal P_B}{\sqrt{M_B}}\le\rho_\perp.
\]
Set
\[
\beta_1:=\sqrt{\overline R_{11}^2+\beta_{21}^2},
\qquad
\beta_2:=\sqrt{\beta_{12}^2+\beta_{22}^2},
\]
and
\[
\boxed{
\mathcal F_B(y)
:=
\frac12
\left[
\beta_{12}
-(\underline R_{11}-\beta_{22})y
+\beta_{21}y^2
+(\beta_2+\beta_1y)\rho_\delta
+2(1+y)\rho_\perp
\right].
}
\]
If
\[
\sup_{y\in[\tan\theta_G,\tan\theta_C]}
\mathcal F_B(y)
\le
-\kappa_{\mathrm{exit}}<0,
\]
then while every neuron of $B$ remains in $J_C^{\mathrm{pure}}$,
\[
D^+y_B\le-\kappa_{\mathrm{exit}}.
\]
Because purity itself implies $y_B\ge\tan\theta_G$, the set cannot remain
entirely in the pure sector for longer than
\[
\boxed{
\frac{
(y_B(t_a)-\tan\theta_G)_+
}{\kappa_{\mathrm{exit}}}
}
\]
without the pure-sector hypothesis breaking.  If, simultaneously, the
positive angular coercivity margin from
Lemma~\ref{lem:negative-sector-residual-reduction} dominates the
mismatch/perpendicular angular errors, the outer face $-\theta_C$ is
impossible and the first angular exit is through the inner edge
$-\theta_G$.  Thus the pure-sector ratio estimate is an exit/recruitment
statement, not the terminal invariant cone.

In the ideal rank-one limit
\[
R_{12}=R_{22}=0,
\qquad
\delta_i=0,
\qquad
P_{i,\perp}^\rho=0,
\]
one has
\[
\boxed{
\frac d{dt}Q_B=0.
}
\]
Equivalently,
\[
Q_B
=
\sum_{i\in B}v_{i2}^2
=
\int\sin^2\phi\,d\nu_B
\]
is conserved while $M_B$ grows.  The leading recruitment mechanism is
therefore longitudinal amplification at nearly fixed transverse energy,
not literal angular transport.
\end{lemma}

\begin{proof}
On the pure sector,
$A_i^\rho=\mathcal R_1/2$, so
\[
\dot v_i
=
\frac12\mathcal R_1^\top(v_i-\delta_i)
-
P_{i,\perp}^\rho.
\]
Expanding the two coordinates gives the displayed $a_i,b_i$ equations.
Multiplication by $a_i$ or $b_i$, summation, and Cauchy--Schwarz give
\eqref{eq:M-sqrt-lower} and \eqref{eq:Q-sqrt-upper}.  Combining the upper
Dini derivative of $\sqrt{Q_B}$ with the lower derivative of
$\sqrt{M_B}$ gives \eqref{eq:pure-ratio}.  The exit conclusion follows from
the geometric lower bound $y_B\ge\tan\theta_G$ while all angles lie in the
pure sector.  In the rank-one limit the $b_i$ equations vanish identically,
which proves conservation of $Q_B$.
\end{proof}

\begin{remark}[Extremal test function]
For the ideal diagonal pure-sector flow,
$G^\rho/V^\rho=\cot|\phi|$.  In the general weighted-transport identity
for $\sum_i r_i^2\Phi(\phi_i)$, the extremal choice
\[
\Phi(\phi)=\sin^2\phi
\]
makes the leading transport term vanish exactly.  This is the same
conservation law $\sum_i v_{i2}^2=\mathrm{const}$ derived above and explains
why a generic Gronwall decay argument is unnecessarily lossy.
\end{remark}


\begin{remark}[Headline anisotropy constant and second-order self-field]
Ignoring the explicit mismatch/perpendicular errors, the small stable ratio
of \eqref{eq:pure-ratio} is controlled to leading order by
\[
\boxed{
y_-
\sim
\frac{|R_{12}|}{R_{11}}.
}
\]
This is the headline constant of the recruitment calculation: transverse
cluster-$1$ residual divided by longitudinal cluster-$1$ residual.  Moreover,
\[
R_{21}=R_{12}+\mathcal X_1,
\]
so the adverse self-field $\mathcal X_1$ enters the ratio dynamics through
the quadratic term $R_{21}y^2$.  Once the inner cone is small, sensitivity
to the self-field is therefore second order.
\end{remark}

%%%%%%%%%%%%%%%%%%%%%%%%%%%%%%%%%%%%%%%%%%%%%%%%%%%%%%%%%%%%%%%%%%%%%%%%%%%%%%%
% T1.10h: inner region, pointwise cone, and strong branch interface
%%%%%%%%%%%%%%%%%%%%%%%%%%%%%%%%%%%%%%%%%%%%%%%%%%%%%%%%%%%%%%%%%%%%%%%%%%%%%%%

\begin{remark}[Corrected inner-branch interface]
\label{prop:inner-strong-branch-interface}
The former proposition at this label incorrectly converted an additive
moving-slab error into a Lipschitz coefficient, omitted the negative-side
interval length from its contraction rate, and inferred a classical
branch across empirical jumps. Proposition~\ref{sbnew:inner} below supplies
the valid pointwise cone and pairwise spread result, with explicit jump
floor and initial transient. A smooth reference branch requires the
additional hypotheses of Proposition~\ref{sbnew:weak-comparison}; no
classical empirical branch is claimed here.
\end{remark}

%%%%%%%%%%%%%%%%%%%%%%%%%%%%%%%%%%%%%%%%%%%%%%%%%%%%%%%%%%%%%%%%%%%%%%%%%%%%%%%
% T1.10i: inner residual decomposition and final assembly interface
%%%%%%%%%%%%%%%%%%%%%%%%%%%%%%%%%%%%%%%%%%%%%%%%%%%%%%%%%%%%%%%%%%%%%%%%%%%%%%%

\begin{lemma}[Inner residual-matrix decomposition]
\label{lem:inner-residual-decomposition}
Assume the explicit cluster-1 purity condition
\eqref{sbnew:inner-purity}. For the selected cluster-2 test gate define
\[
\mathcal R_2^{\rm gate}(\phi,t)
=\mathbb E_{2,n}[(x-f_{\parallel,t}(x))x^\top\alpha_\phi(x)].
\]
Then, including selected boundary values,
\[
\mathfrak A^\rho(\phi,t)
=\tfrac12\mathcal R_1(t)+\tfrac12\mathcal R_2^{\rm gate}(\phi,t).
\]
At a boundary the mixture residual jump is
$\pm(x-f_\parallel(x))x^\top/(2n)$, as in
Lemma~\ref{sbnew:jumps}. The entrywise constants, and the learned-output
budgets they require, are stated explicitly in
Lemma~\ref{sbnew:residual-constants}.
\end{lemma}
\begin{proof}
Split the empirical average into its two cluster averages. Purity makes
the first selector identically one. General position leaves a single
changed summand at a test-angle boundary, proving the jump formula.
\end{proof}

\begin{remark}[T1.10--T1.11 stopping-time assembly: still open]
\label{rem:T110-assembly}
The identities and local comparisons above do not yet exclude every
auxiliary face before the first hitting of $m_1^{\rm eff}=m_{\rm prog}$.
The proof additions below discharge global norm and perpendicular bounds,
correct the early centering and finite-sample branch interfaces, and expose
the exact remaining signed-measure and mismatch obligations.  The new
creation package below further separates the expensive B1 inner-cone/spread
stage from a cheaper B2 late-holding stage and reduces true-field weak-basin
creation to the explicit conditions of
Proposition~\ref{prop:gate-cliff-creation}.  This does not by itself prove
that the random-initialization trajectory reaches those conditions.
In particular, the one-sided terminal self-field certificate alone does not
imply the absolute guard $|\mathcal X_1|\le X_C$.  The exact network
regulator now supplies the algebraic generation bridge; what remains is the
signed strong/complement forcing estimate and enough accumulated regulator
clock for Proposition~\ref{sbgen:first-hit}.  Theorem 1 therefore remains an
unproved target, explicitly marked as such below.
\end{remark}


% Self-contained proof additions. Labels use the sbnew prefix.
\subsection{Scope and status}
This supplement was originally developed against the September 13 source and
is now integrated into the authoritative \texttt{main-10.tex}.  The corrected
gated-moment, finite-sample inner-cone, strong-progress, late-holding, and
gate-cliff creation interfaces below are the versions used by this file.
The experimental record is background evidence; none of its numerical
observations is a hypothesis or a proof step here.

\textbf{Theorem 1 is not closed.} The results below repair two invalid
interfaces, give explicit residual and forcing bounds, and finish the
independent tracking, mass, and selector implications. They do not assume
the missing strong-progress or weak-creation conclusions. In particular,
the conditional comparison lemmas below are not new trajectory assumptions
for the intended unconditional theorem.

Throughout, $u(\phi)=(\cos\phi,\sin\phi)$ denotes the unit vector in
$\mathcal S$, and $u^\angle(\phi)=(-\sin\phi,\cos\phi)$; this temporary
short notation denotes the source's $\widetilde u(\phi)$, not a neuron
weight. Put $N=2n$, $X=X_{\delta_{\rm strong}}$, and
$\lambda=\overline\lambda_\delta$. Fix the source's absolutely continuous
Clarke trajectory and compatible selectors. All differential statements
about that trajectory hold almost everywhere and are integrated using
absolute continuity. No claim below assumes $d=2$.

\subsection{Repair of the early strong-angle estimate}
\begin{proposition}[The centering obstruction and its bias-aware repair]
\label{sbnew:bias}
Let $A(\phi)=\mathbb E_n[xx^\top\mathbf1_{u(\phi)^\top x>0}]$,
$V(\phi)=u(\phi)^\top A(\phi)u^\angle(\phi)$, and define
$\Gamma=u^\top Au$, $T=(u^\angle)^\top Au^\angle$.
Then $V$ is locally absolutely continuous and
\begin{equation}\label{sbnew:target-derivative}
 V'(\phi)=T(\phi)-\Gamma(\phi)\quad\hbox{a.e.}
\end{equation}
On the master data event, $V(0)\ne0$ almost surely. Consequently the
uncentered constant used in the supplied T1.8 and T1.9 satisfies
\begin{equation}\label{sbnew:bad-centering}
 \inf_{0<|\phi|\le\theta}
 \frac{-\operatorname{sgn}(\phi)V(\phi)}{|\phi|}=-\infty
 \qquad(\theta>0).
\end{equation}
At population level its bias is already strictly positive:
\begin{equation}\label{sbnew:population-bias}
 V_{\rm pop}(0)=\frac{\mu_2\sigma_1}{2\sqrt{2\pi}}>0.
\end{equation}

The valid replacement is any deterministic $\kappa_\theta>0$ such that
$\Gamma(\phi)-T(\phi)\ge\kappa_\theta$ almost everywhere on
$[-\theta,\theta]$, together with $b_{\rm tar}\ge|V(0)|$. It gives
\begin{equation}\label{sbnew:bias-coercivity}
 \operatorname{sgn}(\phi)V(\phi)
 \le-\kappa_\theta|\phi|+b_{\rm tar}.
\end{equation}
If $a_\mathrm{-}\le s_i/r_i\le a_\mathrm{+}$ and
$|\psi_i-\phi_i|\le D<\pi/2$, the target term in the input-angle equation
therefore contributes at most
\[
 -a_\mathrm{-}\cos D\,\kappa_\theta|\phi_i|
 +a_\mathrm{+}b_{\rm tar}
 +a_\mathrm{+}\sin D\,T_{1,+}
\]
to the upper derivative of $|\phi_i|$.
\end{proposition}
\begin{proof}
For one sample, $V$ is the function
$(u(\phi)^\top x)_+(u^\angle(\phi)^\top x)$. It is continuous,
piecewise smooth, and locally Lipschitz. Differentiating in each open gate
cell gives
$\mathbf1_{u^\top x>0}[(u^\angle\cdot x)^2-(u\cdot x)^2]$.
The gate-boundary value is zero, so there is no jump. Summation proves
\eqref{sbnew:target-derivative}.

Condition on all first coordinates of the samples. The random variable
\[
 V(0)=\frac1{2n}\sum_{p,k}(x_{k1}^{(p)})_+x_{k2}^{(p)}
\]
is Gaussian with conditional variance
$\sigma_2^2(4n^2)^{-1}\sum_{p,k}(x_{k1}^{(p)})_+^2$.
Whenever at least one first coordinate is positive, this variance is
positive. On the master event every cluster-1 first coordinate is positive.
Thus the event $\{V(0)=0\}$ intersected with the master event has probability
zero; no conditioning on the envelope is used in the Gaussian calculation.
Continuity and approach to zero from the side with the sign of $V(0)$
prove \eqref{sbnew:bad-centering}. Independence of the two Gaussian
coordinates gives zero cluster-1 contribution at population level, and
$\mu_2\mathbb E[(\sigma_1 Z)_+]/2$ from cluster 2, proving
\eqref{sbnew:population-bias}.

Integration of $V'\le-\kappa_\theta$ from zero proves
\eqref{sbnew:bias-coercivity} on both sides. Finally use the source's exact
identity $\mathcal V_{\rm tar}(\phi,\psi)
=\cos(\psi-\phi)V(\phi)+\sin(\psi-\phi)T(\phi)$.
\end{proof}

\paragraph{Deterministic certification from T1.1b.}
Let $\bar A=(\bar A_{1,\star}+\bar A_{2,\star})/2$ and
$e_{\rm op}=(\varepsilon_{\rm op}^{(1)}+
\varepsilon_{\rm op}^{(2)})/2$. The following are valid choices:
\begin{align}
 \kappa_\theta&=
 \inf_{|\phi|\le\theta}
 \{u^\top\bar A(\phi)u-(u^\angle)^\top\bar A(\phi)u^\angle\}
 -2e_{\rm op}>0,\label{sbnew:kappa-cert}\\
 b_{\rm tar}&=
 \left|\frac{\bar a_{12,\star}^{(1)}(0)+
 \bar a_{12,\star}^{(2)}(0)}2\right|
 +\frac{\varepsilon_{12}^{(1)}+\varepsilon_{12}^{(2)}}2.
 \label{sbnew:bias-cert}
\end{align}
These use exactly the clipped population targets in T1.1b; no unaccounted
tail error is discarded. No extra probability event is needed.

\paragraph{Exact changes to T1.8 and T1.9.}
Replace each uncentered angular constant by
$\inf_{|\phi|\le\theta}(\Gamma-T)$, or by the certified lower bound
\eqref{sbnew:kappa-cert}. Add $a_+b_{\rm tar}$ to $C_{1,0}$ in T1.8.
In T1.9 add $a_{0,+}^{\rm post}b_{\rm tar}$ to
$Q_{1,0}^{\rm post}$, and $a_+^{\rm post}b_{\rm tar}$ to both
$R_{\rm lock}^{(0)}$ and $\Psi(D)$. Recompute the dependent displayed
budgets with these definitions. The radial and input/output gap
comparisons require no extra bias term, since they already use
$\sup|V_{\rm tar}|$. The one-way coarse-to-sharp radial argument is
preserved. Proposition~\ref{sbnew:bias} is the missing inequality in the
two strong-angle comparison steps, so those steps now follow with the
stated extra forcing. The resulting fine-lock condition necessarily pays
the nonzero bias floor. Nonemptiness of all resulting time and accuracy
conditions is a separate obligation, not asserted here.

\subsection{Global bounds and the remaining mismatch forcing}
\begin{lemma}[A global norm bound without the Riccati cutoff]
\label{sbnew:global}
Let $S=\sum_i\|u_i\|^2=\sum_i\|w_i\|^2$ and
$C_x=\mathbb E_n\|x\|^2$. Then
\begin{equation}\label{sbnew:global-S}
 \dot S=2\mathbb E_n[(x-f)\cdot f]
 =\frac{C_x}{2}-2\mathbb E_n\|f-x/2\|^2,
 \qquad S(t)\le S(0)+\frac{C_x}{2}t.
\end{equation}
In particular $S(t)\le h\varepsilon^2+\lambda t$, since the data have
rank at most two. On any finite interval the residual matrices obey
\begin{equation}\label{sbnew:A-norm}
 \|A_i^\rho(t)\|_{\rm op}\le
 K:=\sqrt{2\lambda\mathcal L(0)}.
\end{equation}
These conclusions consume balance, the selector identity, and loss
dissipation from T1.2--T1.4, and do not extend the T1.5 Riccati estimate
beyond its horizon.
\end{lemma}
\begin{proof}
Differentiate $S$, sum the selected input equations, and use
$\alpha_i u_i^\top x=\sigma(u_i^\top x)$ to obtain the first identity.
Complete the square for the second. For unit planar vectors $a,b$,
\[
 |a^\top A_i^\rho b|
 \le(\mathbb E_n\|x-f_\parallel\|^2)^{1/2}
       (\mathbb E_n(b^\top x)^2)^{1/2}
 \le\sqrt{2\mathcal L(0)\lambda}.
\]
This proves \eqref{sbnew:A-norm}.
\end{proof}

\begin{lemma}[Aggregate perpendicular forcing and normalized mismatch]
\label{sbnew:forcing}
Write $\mathsf R=\sum_i r_i^2$, $D_B=(\sum_{i\in B}\|\delta_i\|^2)^{1/2}$,
$P_B=(\sum_{i\in B}\|P_{i,\perp}^\rho\|^2)^{1/2}$, and
\[
 Z_B=\left(\sum_{i\in B}r_i^2|X_i^\rho|^2\right)^{1/2}.
\]
For any fixed $B\subseteq[h]$,
\begin{align}
 \|f_\perp\|_n&\le\varepsilon\sqrt{h\lambda\mathsf R},\label{sbnew:fperp}\\
 P_B&\le\lambda\varepsilon\sqrt{h\mathsf R}
           \|W_{\perp,B}\|_F
       \le\lambda h\varepsilon^2\sqrt{\mathsf R},\label{sbnew:P-bound}\\
 D_B(t)&\le e^{K(t-a)}D_B(a)
       +\int_a^t e^{K(t-s)}[Z_B(s)+P_B(s)]\,ds.\label{sbnew:D-bound}
\end{align}
If, on a specified interval, a positive longitudinal amplitude
$U_B=\sqrt{M_B}$ satisfies $\dot U_B/U_B\ge g_B$, then
$d_B=D_B/U_B$ satisfies
\begin{equation}\label{sbnew:normalized-D}
 \dot d_B\le(K-g_B)d_B+Z_B/U_B+P_B/U_B\quad\hbox{a.e.}
\end{equation}
If $g_B>K$ and the last two ratios sum to at most $z_B$, this implies
\[
 d_B(t)\le e^{-(g_B-K)(t-a)}d_B(a)
 +\frac{z_B}{g_B-K}(1-e^{-(g_B-K)(t-a)}).
\]
When $g_B\le K$, the corresponding growing or linear Duhamel kernel must
be used instead; a damping rate is not available from this estimate.
\end{lemma}
\begin{proof}
Write $f_\perp=W_\perp a(x)$, with $a_i(x)=\sigma(v_i^\top x)$.
Then $\sum_i\|a_i\|_n^2\le\lambda\mathsf R$ and
$\|W_\perp\|_F\le\varepsilon\sqrt h$. For unit $q\in\mathcal S$,
\[
 |q\cdot P_{i,\perp}^\rho|
 \le\sqrt\lambda\,\|w_{i,\perp}\|\,\|f_\perp\|_n.
\]
Take the supremum in $q$, square, and sum. The source's residual-mismatch
energy inequality gives $\dot D_B\le KD_B+Z_B+P_B$ in the regularized
norm sense and hence \eqref{sbnew:D-bound}. The quotient rule proves
\eqref{sbnew:normalized-D} wherever $U_B>0$.
\end{proof}

\paragraph{What this closes and what it does not.}
The whole-network perpendicular forcing per $\sqrt{\mathsf R}$ is now
bounded explicitly by $\lambda h\varepsilon^2$, including in ambient
dimension larger than two. For a cohort normalized by $\sqrt{M_B}$, the
additional ratio $\sqrt{\mathsf R/M_B}$ must still be paid. On the pure
sector, $Z_B=|\mathcal X_1|\sqrt{\mathsf R_B}/2$. For the whole network
there is no such reduction. The elementary bound
$|X_i^\rho|\le2K$ gives only $Z_{[h]}\le2K\sqrt{\mathsf R}$, an order-one
forcing. Thus initialization smallness and perpendicular contraction alone
do not close the small global mismatch budget in the terminal certificate.

\subsection{Explicit residual constants from structured output bounds}
\begin{lemma}[Residual bounds with all learned-output inputs exposed]
\label{sbnew:residual-constants}
On an interval $I$, let $F_p$ denote the neurons strictly active on every
cluster-$p$ sample at the time being considered, and write
\[
 B_p=\sum_{i\in F_p}z_i v_i^\top,\qquad
 f_\parallel(x)=B_px+g_p(x)\quad(x\hbox{ in cluster }p).
\]
Neurons active on no sample, including all-zero preactivations, contribute
zero to $g_p$. Let $Z_p\ge\sum_{i\ {
m partial\ on}\ p}r_is_i$, and let
$G_p$ bound $\Gamma_p(\phi_i)$ for all such partial neurons. For either
cluster and every test gate $\alpha\in[0,1]$,
\begin{equation}\label{sbnew:tail}
 |\mathbb E_{p,n}[g_{p,r}(x)x_s\alpha(x)]|
 \le Z_p\sqrt{G_p\,s_{ss}^{(p)}},
 \qquad s_{rs}^{(p)}=\mathbb E_{p,n}[x_rx_s].
\end{equation}

Suppose $|(B_p)_{rs}|\le K_{rs}^{(p)}$ on $I$.
Define
\begin{align*}
 b_{12}^{(1)}&=(1+K_{11}^{(1)})|s_{12}^{(1)}|
 +K_{12}^{(1)}s_{22}^{(1)}+Z_1\sqrt{G_1s_{22}^{(1)}},\\
 b_{22}^{(1)}&=K_{21}^{(1)}|s_{12}^{(1)}|
 +(1+K_{22}^{(1)})s_{22}^{(1)}+Z_1\sqrt{G_1s_{22}^{(1)}}.
\end{align*}
Then $|R_{12}|\le b_{12}^{(1)}$ and $|R_{22}|\le b_{22}^{(1)}$.

For an inner interval $[-\theta_G,\theta_H]$, require the explicit
cluster-1 purity condition
\begin{equation}\label{sbnew:inner-purity}
 \cos\theta_*\,m_1^{\rm data}-\sin\theta_*\,b_2>0,
 \qquad\theta_*:=\max\{\theta_G,\theta_H\}<\pi/2.
\end{equation}
Let $l_{11},u_{11},u_{22},u_{12}$ be, respectively, a lower bound for
the selected cluster-2 $a_{11}$, upper bounds for $a_{11},a_{22}$, and
an upper bound for $|a_{12}|$, uniformly over this interval.
Let $R_{11}\ge L_{11}$ and $(B_2)_{22}\ge-\ell_{22}$ on $I$.
Put $T_s=Z_2\sqrt{G_2s_{ss}^{(2)}}$. The following constants bound
every inner residual matrix, including selected boundary values:
\begin{align}
 \alpha_{11}&=\tfrac12\{L_{11}+l_{11}
       -K_{11}^{(2)}u_{11}-K_{12}^{(2)}u_{12}-T_1\},\label{sbnew:alpha11}\\
 \alpha_{22}&=\tfrac12\{b_{22}^{(1)}+(1+\ell_{22})u_{22}
       +K_{21}^{(2)}u_{12}+T_2\},\label{sbnew:alpha22}\\
 \beta_{12}^{\rm in}&=\tfrac12\{b_{12}^{(1)}
       +(1+K_{11}^{(2)})u_{12}+K_{12}^{(2)}u_{22}+T_2\},\label{sbnew:beta12}\\
 \beta_{21}^{\rm in}&=\tfrac12\{b_{12}^{(1)}+X_C
       +K_{21}^{(2)}u_{11}+(1+K_{22}^{(2)})u_{12}+T_1\},\label{sbnew:beta21}
\end{align}
where the last line requires $|\mathcal X_1|\le X_C$.
Thus $A_{11}^\rho-A_{22}^\rho\ge d_H:=\alpha_{11}-\alpha_{22}$.
\end{lemma}
\begin{proof}
For a partial neuron, Cauchy--Schwarz gives
$|\mathbb E_{p,n}[q_{\phi_i}x_s\alpha]|\le
\sqrt{\Gamma_p(\phi_i)s_{ss}^{(p)}}$.
Summation proves \eqref{sbnew:tail}. Expanding
$(I-B_1)\mathbb E_{1,n}[xx^\top]-\mathbb E_{1,n}[g_1x^\top]$
gives the first two residual bounds, retaining the noise-coordinate
factor $s_{22}^{(1)}$.

Condition \eqref{sbnew:inner-purity} makes cluster 1 fully active for
every angle under consideration. Consequently
\[
 A^\rho=\tfrac12\mathcal R_1+
 \tfrac12\{(I-B_2)\widehat A_2^\alpha-
              \mathbb E_{2,n}[g_2x^\top\alpha]\}.
\]
For its $(1,1)$ entry, bound
$a_{11}-(B_2)_{11}a_{11}-(B_2)_{12}a_{12}$ from below.
For $(2,2)$, use $a_{22}\ge0$ and
$(B_2)_{22}\ge-\ell_{22}$ to bound it from above.
Triangle inequalities give the other two entries, and
$|R_{21}|\le|R_{12}|+|\mathcal X_1|$ gives the final line.
\end{proof}

\paragraph{How these constants can be fixed before use.}
T1.1b supplies $l_{11},u_{11},u_{22},u_{12}$ from clipped population
infima/suprema and entrywise errors. Adding $E_{rs}^{(2)}/n$ to the
corresponding errors covers selected boundary values. This addition is a
deterministic consequence of general position, not a new event.
The source's displayed $(R11\text{-lower})$ supplies $L_{11}$; its
initial loss can be bounded by
$\mathcal L(0)\le\tfrac12(1+h\varepsilon^2)^2C_x$.
Every empirical moment in the preceding formulas can likewise be replaced
by a bound on the existing data event.

For a time ceiling $H$, Proposition~\ref{sbnew:global} gives the safe choices
$Z_p\le S_H:=h\varepsilon^2+\lambda H$ and $K_{rs}^{(p)}\le S_H$.
The useful one-sided diagonal estimate is
\begin{equation}\label{sbnew:diagonal-bulk}
 (B_p)_{rr}=\sum_{i\in F_p}(v_{ir}^2-\delta_{ir}v_{ir})
 \ge-\sqrt{\mathsf R\Delta}.
\end{equation}
Hence $\ell_{22}=\sqrt{R_*\Delta_*}$ under guards
$\mathsf R\le R_*$ and $\Delta\le\Delta_*$.
Sharper off-diagonal $K_{12}^{(2)}$ requires a proved angular-moment
bound; it cannot be inferred from $m_1^{\rm eff}$ alone.
Substituting the safe choices is valid but does not establish positive
discriminant or the nonemptiness of the intended theorem region.

\paragraph{A correct deterministic partial band.}
For cluster $p$, let $q\ne p$. The partial band is contained in
\begin{equation}\label{sbnew:partial-band}
 |u_p(\phi)|m_p^{\rm data}\le |u_q(\phi)|b_q.
\end{equation}
Indeed, outside this band the mean coordinate dominates every noise
realization in the coordinate box, so all signs agree. In the band,
$q_\phi(x)\le2\mu_p b_q/m_p^{\rm data}$ on cluster $p$, giving
$G_p\le(2\mu_pb_q/m_p^{\rm data})^2$, or the sharper T1.1b bound on that
same deterministic band. A strictly feasible $\theta_C$ is a buffered
purity edge; it need not be the exact inner boundary of the partial band.

\subsection{A valid finite-sample inner-cone and spread result}
\begin{lemma}[Residual jumps are not target-moment jumps]
\label{sbnew:jumps}
At a test-angle boundary for one sample $x$, the two residual matrices
differ by
\begin{equation}\label{sbnew:residual-jump}
 \Delta A^\rho=\pm\frac1{2n}(x-f_\parallel(x))x^\top.
\end{equation}
Thus an entry jump is at most
$(|x_r|+|f_r(x)|)|x_s|/(2n)$; T1.1b's target bound alone does not apply.
For the aligned test angular field $V^\rho=u^\top A^\rho u^\angle$,
\begin{equation}\label{sbnew:V-jump}
 \Delta V^\rho=\mp\frac1{2n}
 (u(\phi_0)\cdot f_\parallel(x))(u^\angle(\phi_0)\cdot x).
\end{equation}
In general this is nonzero.
\end{lemma}
\begin{proof}
Only one empirical summand changes. At the crossing $u(\phi_0)\cdot x=0$,
which removes the target part from the scalar jump but not the learned
output part. This proves both identities.
\end{proof}

\begin{proposition}[Inner cone and pairwise contraction with a jump floor]
\label{sbnew:inner}
Use the ratio coordinate $\xi=\tan\phi$ on
$[-g,h]=[-\tan\theta_G,\tan\theta_H]$, and put $T_*:=\max(g,h)$.
Assume the preceding residual bounds hold throughout a stopped interval
on which the relevant projected radii remain positive.
Let $|\mathcal E_i^\xi|\le e_\xi$ in the exact ratio equation
\begin{equation}\label{sbnew:ratio}
 \dot\xi_i=b(\xi_i,t)+\mathcal E_i^\xi,\qquad
 b(\xi,t)=A_{12}^\rho-(A_{11}^\rho-A_{22}^\rho)\xi-A_{21}^\rho\xi^2.
\end{equation}
The error bound follows, without a gate approximation, from
\begin{equation}\label{sbnew:ratio-error}
 |\mathcal E_i^\xi|\le
 (1+\xi_i^2)
 \left(K\frac{\|\delta_i\|}{r_i}+
                  \frac{\|P_{i,\perp}^\rho\|}{r_i}\right).
\end{equation}
Define
\[
 F(y)=\beta_{12}^{\rm in}+e_\xi-d_Hy+\beta_{21}^{\rm in}y^2.
\]
Choose $0<r_c<\min(g,h)$ and $\kappa_c>0$ with
\begin{equation}\label{sbnew:cone-conditions}
 \max_{y\in[r_c,T_*]}F(y)\le-\kappa_c.
\end{equation}
Then $[-g,h]$ is invariant for each such neuron. Provided the stopped
interval lasts that long, every entering neuron reaches $[-r_c,r_c]$ within at most
$(|\xi_i(t_a)|-r_c)_+/\kappa_c$, and the latter cone is invariant.
For a convex $F$, condition \eqref{sbnew:cone-conditions} is equivalent to
checking its two endpoints. A necessary part of this condition is the
strict forced discriminant
$d_H^2>4\beta_{21}^{\rm in}(\beta_{12}^{\rm in}+e_\xi)$ when
$\beta_{21}^{\rm in}>0$.

Inside each open gate cell,
\begin{equation}\label{sbnew:cell-contract}
 \partial_\xi b\le-\lambda_0,
 \qquad \lambda_0:=d_H-2\beta_{21}^{\rm in}T_*.
\end{equation}
Suppose the sum of positive scalar jumps between any $a<b$ in the interval,
including endpoint selector allowances, is bounded by
\begin{equation}\label{sbnew:jump-modulus}
 J_+(a,b;t)\le L_J(b-a)+j_J,
 \qquad\lambda_c:=\lambda_0-L_J>0.
\end{equation}
Then each pair of neurons remaining in this interval obeys
\begin{equation}\label{sbnew:pairwise}
 |\xi_i(t)-\xi_j(t)|\le
 e^{-\lambda_c(t-a)}|\xi_i(a)-\xi_j(a)|
 +\frac{j_J+2e_\xi}{\lambda_c}(1-e^{-\lambda_c(t-a)}).
\end{equation}
The same bound holds for their angular distance, since $\arctan$ is
1-Lipschitz. This proves a family spread with an initial transient and an
additive gate floor; it does not assert a classical empirical root.
\end{proposition}
\begin{proof}
The quotient rule applied to
$\dot v_i=(A_i^\rho)^\top(v_i-\delta_i)-P_{i,\perp}^\rho$
proves \eqref{sbnew:ratio}. Since $v_{i1}=r_i/\sqrt{1+\xi_i^2}$,
the source's ratio-error formula is exactly \eqref{sbnew:ratio-error}.
For $y=|\xi_i|$, the upper derivative is bounded by $F(y)$ outside the
inner cone. At the outer interval endpoints the signs are inward, and on
$[r_c,T_*]$ the magnitude decreases at speed at least $\kappa_c$.
Scalar comparison and a first-exit argument prove both invariances and
the entry time. This requires control on the whole interval; a negative
value of $F$ at one small radius alone is insufficient.

At fixed training time the test-gate matrix is constant in every gate cell,
so differentiating the polynomial $b$ gives
$-(A_{11}-A_{22})-2A_{21}\xi$ and proves
\eqref{sbnew:cell-contract}. Integrating cell by cell and adding jumps yields
\[
 b(b,t)-b(a,t)\le-\lambda_c(b-a)+j_J.
\]
This also covers all selected endpoint values by the stated endpoint
allowance. Apply this inequality in the order of $\xi_i,\xi_j$, add their
two error bounds, and integrate the scalar comparison inequality. At
coincident angles selected velocities are covered by the same $j_J$.
\end{proof}

\paragraph{An explicit but potentially loose jump allowance.}
If $\sup_{x\in\mathcal D}\|f_\parallel(x,t)\|\le F_H$, then each scalar
ratio jump is at most
$c_J/(2n)$, where $c_J=(1+T_*^2)F_H X$; the target part vanishes at the
crossing. The master moving-slab bound therefore permits
\begin{equation}\label{sbnew:jump-explicit}
 L_J=c_J\sqrt{2/\pi}\,X/\sigma_{\min},\qquad
 j_J=c_J\bigl(\Delta_{\rm VC}+1/n\bigr).
\end{equation}
To see this, all crossed samples lie in the slab of width
$X|b-a|$ around $u(\arctan a)$, and averaging the two cluster counts
gives the bound; $1/n$ safely covers endpoint selections.
One may set $F_H=XS_H$ from \eqref{sbnew:global-S}. These constants need
not make $\lambda_c>0$. A useful sharper bound must exploit the signed
or directional learned-output jump, rather than divide an additive
$\Delta_{\rm VC}$ error by a vanishing angular distance.

\paragraph{The frozen quadratic is not the moving branch.}
For fixed coefficients $a=A_{12}$, $d=A_{11}-A_{22}>0$, $c=A_{21}$,
the small root of $a-d\xi-c\xi^2$ is
\[
 \xi_* =\frac{2a}{d+\sqrt{d^2+4ca}}
 =\frac ad-\frac{ca^2}{d^3}
   +O\!\left(\frac{c^2|a|^3}{d^5}\right)
\]
when $|4ca/d^2|$ is bounded strictly below one. The continuous formula
covers $c=0$. In the empirical field, the coefficients depend on the test
gate. This root must also lie in the cell furnishing those coefficients;
otherwise it is not a root of the actual piecewise field.
Strict negativity of the cell derivatives does not repair an upward jump,
nor do three signs imply a classical root across a downward jump.

%%%%%%%%%%%%%%%%%%%%%%%%%%%%%%%%%%%%%%%%%%%%%%%%%%%%%%%%%%%%%%%%%%%%%%%%%%%%%%%
% B2: late holding after the corrected B1 inner-cone/spread stage
%%%%%%%%%%%%%%%%%%%%%%%%%%%%%%%%%%%%%%%%%%%%%%%%%%%%%%%%%%%%%%%%%%%%%%%%%%%%%%%

\begin{proposition}[Late holding of an already-specialized strong family]
\label{prop:late-strong-holding}
Fix a deterministic interval
\[
I_{\mathrm{hold}}=[t_{\sharp},t_{\mathrm{cert}}]
\]
and a fixed neuron set $B_1$.  Write
\[
\xi_i:=\frac{v_{i2}}{v_{i1}}
\]
whenever $v_{i1}>0$.  Suppose that at $t_{\sharp}$,
\[
v_{i1}(t_{\sharp})>0,
\qquad
-\bar\xi_-<\xi_i(t_{\sharp})<\bar\xi_+
\qquad(i\in B_1).
\]
Assume that, whenever $t\in I_{\mathrm{hold}}$ and
$-\bar\xi_-\le\xi_i\le\bar\xi_+$, the exact ratio equation
\eqref{sbnew:ratio} obeys
\[
\dot\xi_i
=
A_{12,i}^\rho-\Delta_i^\rho\xi_i-A_{21,i}^\rho\xi_i^2+\mathcal E_i^\xi,
\qquad
\Delta_i^\rho:=A_{11,i}^\rho-A_{22,i}^\rho,
\]
with
\[
\Delta_i^\rho\ge d_{\mathrm{hold}}>0,
\qquad
\underline A_{12}^{\mathrm{hold}}
\le A_{12,i}^\rho
\le\overline A_{12}^{\mathrm{hold}},
\]
\[
|A_{21,i}^\rho|\le\beta_{21}^{\mathrm{hold}},
\qquad
|\mathcal E_i^\xi|\le\varepsilon_{\mathrm{hold}}.
\]
If
\[
\boxed{
\overline A_{12}^{\mathrm{hold}}
-d_{\mathrm{hold}}\bar\xi_+
+\beta_{21}^{\mathrm{hold}}\bar\xi_+^2
+\varepsilon_{\mathrm{hold}}<0,
}
\tag{late-hold-plus}\label{eq:late-hold-plus}
\]
and
\[
\boxed{
\underline A_{12}^{\mathrm{hold}}
+d_{\mathrm{hold}}\bar\xi_-
-\beta_{21}^{\mathrm{hold}}\bar\xi_-^2
-\varepsilon_{\mathrm{hold}}>0,
}
\tag{late-hold-minus}\label{eq:late-hold-minus}
\]
then
\[
\boxed{
-\bar\xi_-<\xi_i(t)<\bar\xi_+
\qquad
\forall i\in B_1,\quad \forall t\in I_{\mathrm{hold}}.
}
\tag{late-hold-cone}\label{eq:late-hold-cone}
\]
No forced discriminant, pairwise contraction rate, or empirical branch
continuation is required on $I_{\mathrm{hold}}$.
\end{proposition}

\begin{proof}
At the positive face $\xi_i=\bar\xi_+$,
\[
\dot\xi_i
\le
\overline A_{12}^{\mathrm{hold}}
-d_{\mathrm{hold}}\bar\xi_+
+\beta_{21}^{\mathrm{hold}}\bar\xi_+^2
+\varepsilon_{\mathrm{hold}}<0.
\]
At the negative face $\xi_i=-\bar\xi_-$,
\[
\dot\xi_i
\ge
\underline A_{12}^{\mathrm{hold}}
+d_{\mathrm{hold}}\bar\xi_-
-\beta_{21}^{\mathrm{hold}}\bar\xi_-^2
-\varepsilon_{\mathrm{hold}}>0.
\]
The first-exit argument excludes both faces. At empirical test-gate
boundaries the parameter trajectory remains continuous and the same face
argument uses the compatible one-sided selected vector fields.
\end{proof}

\begin{lemma}[Directional physical ratio field and reciprocal upper-face transfer]
\label{lem:late-directional-ratio}
For a strong receiver with physical input ratio $\xi=v_2/v_1$ and selected
residual matrix $A=A^\rho$, put
\[
 b_{\rm phys}(\xi)
 :=A_{12}-(A_{11}-A_{22})\xi-A_{21}\xi^2.
\]
Then exactly
\[
\boxed{
 b_{\rm phys}(\xi)
 =(1,\xi)A(-\xi,1)^\top
 =\mathbb E_n\!\left[
 \bigl((x_1+\xi x_2)-(f_1+\xi f_2)\bigr)
 (x_2-\xi x_1)\alpha_\xi
 \right].
}
\tag{physical-directional-field}\label{eq:physical-directional-field}
\]
If $\xi>0$ and $\tau=1/\xi$, then the receiver gate equals the T1.11
cotangent gate,
\[
 \mathbf 1\{x_1+\xi x_2>0\}
 =\mathbf 1\{\tau x_1+x_2>0\},
\]
and hence
\[
\boxed{
 b_{\rm phys}(\xi)=\xi^2Q^{\rm true}(1/\xi).
}
\tag{reciprocal-upper-face}\label{eq:reciprocal-upper-face}
\]
This reciprocal transfer is structurally an upper-face identity only.  For
$\xi<0$, multiplication by $1/\xi$ reverses the inequality, so the lower face
continues to use the one-sided condition \eqref{eq:late-hold-minus} (or the
corresponding full-minus-selected decomposition), not
\eqref{eq:reciprocal-upper-face}.
\end{lemma}
\begin{proof}
The matrix identity is a direct expansion.  For $\xi>0$, multiply the gate
inequality by $\tau=1/\xi$ and expand the definition of the T1.11 cotangent
quadratic.
\end{proof}

\begin{remark}[Asymmetric late holding is the intended B2 interface]
\label{rem:asymmetric-B2-holding}
The load-bearing late region need not be a symmetric noise-scale box.  The
intended high-SNR interface is asymmetric: retain a noise-scale lower input
face and output floor, allow a wider positive physical input face
$\xi\le u$, and use the longitudinal/balance barrier on the output upper
side.  The positive input face is sharpened by
\eqref{eq:reciprocal-upper-face}; the negative input face remains
\eqref{eq:late-hold-minus}.  This avoids proving late synchronization of all
output slopes and does not require the optional $c_{12}^O$ coherence route.
\end{remark}

\begin{remark}[B1 contraction versus B2 holding]
\label{rem:B1-B2-split}
The full conditions of Proposition~\ref{sbnew:inner} are used only in B1,
through a branch-locking checkpoint at which the strong family is already
inside the smaller holding cone.  The later B2 stage pays only the face
conditions above.  The positive face condition is, to leading order,
\[
\bar\xi_+>\frac{\overline A_{12}^{\mathrm{hold}}}{d_{\mathrm{hold}}},
\]
so the cone need only contain the moving small strong ratio.  When
$A_{12}^\rho>0$, the principal field on the negative face already points
inward; only the explicit quadratic and forcing allowances in
\eqref{eq:late-hold-minus} remain to be paid.

The exact error estimate inherited from \eqref{sbnew:ratio-error} is
\[
\boxed{
|\mathcal E_i^\xi|
\le
(1+\xi_i^2)
\left(
K\frac{\|\delta_i\|}{r_i}
+\frac{\|P_{i,\perp}^\rho\|}{r_i}
\right).
}
\tag{late-hold-error}\label{eq:late-hold-error}
\]
Thus widening the late holding cone relaxes the mismatch/perpendicular
forcing requirement without requiring a late lower bound on the B1
pairwise-contraction rate.  The creation argument below treats strong-family
angular displacement through a signed kernel rather than through a crude
feature-defect norm, so a several-degree holding cone is not automatically
charged as an $O(\bar\xi_+)$ creation error.

This proposition supplies only outer holding.  It does not by itself give the
one-signed strong coherence needed to use the $c_{12}^O$ feedback channel.
Subsection~\ref{sbcoh:scope} refines B2 at the natural noise-scale tilt: the
collective branch is controlled by a total derivative, while pairwise spread
is controlled by the strictly weaker own-tilt derivative.  That refinement is
optional for a $c_{21}^O$-only post-B1 certificate.
\end{remark}

\begin{remark}[Separate branch-locking and creation progress scales]
\label{rem:separate-progress-levels}
The quantity $m_{\mathrm{prog}}$ in the still-open strong-progress first-exit
package is a B1 \emph{branch-locking} target.  It is not the weak-basin
creation threshold.  The later B2 interval instead carries a deterministic
progress ceiling
\[
\boxed{m_{\max}>m_{\mathrm{prog}},}
\]
while Proposition~\ref{sbgen:first-hit} certifies creation at the first time
$b(t)$ reaches the moving reserved middle threshold $B_M(m(t))$.  Thus the
certification time need not coincide with a prescribed level set of $m$.
The B2 progress ceiling is not fixed a priori.  In the corrected high-SNR
ledger it must be chosen jointly with the normalized creation charges.  The
zero-charge creation/progress lower threshold is approximately $m=0.5296$,
while late holding and mixed-sign rate separation require
$m<1-\ell_2/\ell_1=5/9$.  The feasible ceiling is therefore selected only
after the charged creation inequality is evaluated; $m_{\max}=0.50$ is not a
load-bearing recommendation.
\end{remark}

\subsection{What a strong-progress first-exit proof still requires}
The source's whole-network identities are retained, with
$\nu_t=\sum_i r_i^2\delta_{\phi_i}$ and
$p_1(\phi)=\cos\phi\,\mathsf H_1(\phi)/\mathsf D_1$.
Define the continuous state functionals
\begin{align*}
 P_-(t)&=\int p_1\,d\nu_t-C_{H,1}\sqrt{\mathsf R\Delta},\\
 C_+(t)&=\int[V_1]_+\,d\nu_t+2\lambda_{1,+}\sqrt{\mathsf R\Delta},\\
 C_-(t)&=\int[-V_1]_+\,d\nu_t+2\lambda_{1,+}\sqrt{\mathsf R\Delta}.
\end{align*}
The target kernel $V_1$ is continuous by the proof of
Proposition~\ref{sbnew:bias}; the products extend continuously through
zero projected radii. These identities give
\begin{equation}\label{sbnew:signed-certificate}
 m_1^{\rm eff}\ge P_-,\qquad -C_-\le\mathcal X_1\le C_+.
\end{equation}
In particular $C_+\le X_C$ proves only $\mathcal X_1\le X_C$.
The absolute guard used in the source also needs $C_-\le X_C$, or a
separate lower bound on $\mathcal X_1$. Negative self-field helps the
pure-sector angular lower bound but still affects $|R_{21}|$ in mismatch
and ratio estimates. These are different uses of its sign.

\begin{corollary}[An asymmetric replacement for the self-field guard]
\label{sbnew:asymmetric}
Let $C_{x,1}=\mathbb E_{1,n}\|x\|^2$ and
$X_{\rm loss}=2\sqrt{\mathcal L(0)C_{x,1}}$. Then, at every time,
\[
 |\mathcal X_1(t)|\le X_{\rm loss}.
\]
Thus one may replace the symmetric guard by
$-X_-<\mathcal X_1<X_C$, with any deterministic
$X_->X_{\rm loss}$. Its negative face is automatically excluded by loss
dissipation. In the pure-sector angular and radial lower bounds, only
the upper ceiling $X_C$ is needed. In bounds on $|R_{21}|$ or the pure
mismatch forcing, use $\max(X_C,X_{\rm loss})$ instead. Alternatively,
$|R_{21}|\le2\sqrt{\mathcal L(0)s_{11}^{(1)}}$ follows directly from
loss dissipation.
\end{corollary}
\begin{proof}
Since $f_\parallel=x+\rho_\parallel$, one has
$\mathcal X_1=\mathbb E_{1,n}[\rho_\parallel\times x]$.
Cauchy--Schwarz and
$\mathbb E_{1,n}\|\rho\|^2\le4\mathcal L(0)$ prove the absolute bound.
On the negative pure sector, the coefficients of $\mathcal X_1$ in
both displayed source lower fields are negative, so replacing it by an
upper bound is valid. The $R_{21}$ estimate is another direct
Cauchy--Schwarz inequality.
\end{proof}
This is a rigorous alternative to propagating $C_-$, at the cost of a
potentially larger absolute forcing constant. It does not show that the
smaller symmetric ceiling in the supplied assembly persists, nor that
the resulting discriminant and mismatch budgets close.

\begin{lemma}[Exact whole-network weighted evolution]
\label{sbnew:measure-evolution}
For a continuously differentiable periodic test function $\Phi$, while
the relevant $r_i$ are positive,
\begin{equation}\label{sbnew:measure-derivative}
 \frac d{dt}\int\Phi\,d\nu_t
 =\sum_i r_i^2(2G_i^\rho\Phi(\phi_i)+V_i^\rho\Phi'(\phi_i))
 +\sum_i(2r_i e_i^r\Phi(\phi_i)+r_i^2e_i^\phi\Phi'(\phi_i)),
\end{equation}
where
\[
 G_i^\rho=u_i^{\sim\top}A_i^\rho u_i^\sim,
 \quad e_i^r=-\delta_i^\top A_i^\rho u_i^\sim-u_i^\sim\cdot P_{i,\perp}^\rho,
\]
\[
 e_i^\phi=-r_i^{-1}\delta_i^\top A_i^\rho (u_i^\sim)^\angle
             -r_i^{-1}(u_i^\sim)^\angle\cdot P_{i,\perp}^\rho,
 \qquad u_i^\sim=v_i/r_i.
\]
The magnitude of the second sum is bounded by
\begin{equation}\label{sbnew:measure-error}
 (2\|\Phi\|_\infty+\|\Phi'\|_\infty)
 \{K\sqrt{\mathsf R\Delta}+\lambda h\varepsilon^2\mathsf R\}.
\end{equation}
\end{lemma}
\begin{proof}
Take radial and tangential projections of the exact selected Cartesian
equation, differentiate $r_i^2\Phi(\phi_i)$, and sum. Bound the two
errors by $K\|\delta_i\|+\|P_{i,\perp}^\rho\|$ and apply
Cauchy--Schwarz and \eqref{sbnew:P-bound}. Smooth approximations or the
almost-everywhere chain rule give the corresponding integrated statement
for Lipschitz kernels when needed; no favorable sign follows merely from
this extension.
\end{proof}

\paragraph{The precise unproved dynamical obligation.}
Starting from the corrected early handoff, one must still derive from
primitive parameters a finite $H$ and strict budgets such that, before
$m_1^{\rm eff}=m_{\rm prog}$:
\begin{enumerate}
\item the integral in \eqref{sbnew:D-bound} stays below the chosen global
and pointwise mismatch budgets, including the weighted antisymmetric
forcing $Z_{[h]}$;
\item the learned-output bulk constants in
\eqref{sbnew:alpha11}--\eqref{sbnew:beta21} preserve the forced cone and
spread margins, with an admissible entry schedule for a fixed eventual
strong family;
\item the harmful cross-output face stays strictly inward (or below a
strict integrated comparison envelope) on the entire stopping interval;
if the symmetric ceiling is retained, its negative face also needs a
separate argument instead of the asymmetric replacement above;
\item the whole-network progress functional reaches $m_{\rm prog}$ by
$H$. One sufficient route is a proved lower evolution bound for $P_-$
whose comparison reaches that value by $H$.
\end{enumerate}
Even after B1 reaches $m_{\rm prog}$, the later creation interface additionally
requires B2 holding up to the available progress ceiling, a signed regulator
forcing bound, enough accumulated mass-weighted regulator clock to hit
$B_M(m)$, and a timing bound on the orthogonalized weak-feature coefficient.
These are kept separate from the B1 discriminant/spread budget.
The supplied source proves none of the four B1 uniform conclusions above,
and the new B2/creation guards are explicit predecessor obligations rather
than assumed trajectory facts.
Equation \eqref{sbnew:measure-derivative} exposes the missing evolution
estimate: its principal signed integral is not controlled by the local
pure-sector result for a fixed set. A terminal identity cannot rule out an
earlier guard failure and does not establish that the terminal time exists.
Also, at the first exact progress hitting time, $P_-$ may be strictly less
than $m_{\rm prog}$; the sufficient certificate must not be treated as an
equivalent definition of that hitting time.

These are missing estimates for the intended random initialization, not
counterexamples to the intended dynamical phenomenon. They cannot be
declared admissibility conditions on the trajectory while retaining an
unconditional theorem. The perpendicular part and the finite-time norm
ceiling have been discharged above; the signed angular/mismatch part has
not.

\subsection{Independent coarse-field sign result}
\begin{lemma}[Equal-noise coarse-field asymmetry and wrong-quadrant exclusion]
\label{sbnew:coarse}
For the population coarse model $f_m(x)=m e_1(x_1)_+$, define
\[
 V_0(\phi;m)=\mathbb E[((x-f_m(x))\cdot u(\phi))
               \mathbf1_{u(\phi)\cdot x>0}(u^\angle(\phi)\cdot x)].
\]
Then $V_0(\phi;m)=V_{\rm tar,pop}(\phi)-m\cos\phi\,H'(\phi)$,
where $H(\phi)=\mathbb E[(x_1)_+(u(\phi)\cdot x)_+]$.
On $\pi/2<\phi<\pi$, $H'(\phi)<0$, and hence
\begin{equation}\label{sbnew:coarse-q2-m}
 \partial_m V_0(\phi;m)<0.
\end{equation}
If $\sigma_1=\sigma_2=\sigma$, then
$V_{\rm tar,pop}(\phi)<0$ throughout that quadrant. Thus
$V_0(\phi;m)<0$ there for every $m\ge0$, and the coarse model has no
root in that open quadrant.
\end{lemma}
\begin{proof}
Differentiate the ReLU feature in $H$ under the Gaussian integral; its
boundary contribution vanishes. The formula for $V_0$ follows. In the
second quadrant, write $\cos\phi=-a$, $\sin\phi=b$, with $a,b>0$.
On $x_1>0$ and $-ax_1+bx_2>0$, one has $x_2>ax_1/b>0$ and
$u^\angle\cdot x=-bx_1-ax_2<0$. This event has positive probability,
so $H'<0$, proving \eqref{sbnew:coarse-q2-m}.

For isotropic noise within $\mathcal S$, projections onto $u$ and
$u^\angle$ are independent Gaussians. Put
$h(z)=z\Phi(z)+\varphi(z)$, where $\Phi,\varphi$ are the standard
normal distribution function and density. Then
\[
 V_{\rm tar,pop}(\phi)
 =\tfrac12\{-\mu_1 b\,\sigma h(-\mu_1 a/\sigma)
             -\mu_2 a\,\sigma h(\mu_2 b/\sigma)\}.
\]
Both terms are strictly negative because
$h(z)=\mathbb E[(z+Z)_+]>0$.
Finally $-m\cos\phi H'(\phi)\le0$.
\end{proof}

\paragraph{Scope of the asymmetry claim.}
In the first quadrant, the exact criterion is
$\partial_mV_0>0$ if and only if $H'(\phi)<0$; it is not an automatic
whole-quadrant identity. At zero,
$H'(0)=\mu_2\sigma_1/(2\sqrt{2\pi})>0$, so
$\partial_mV_0(0;m)<0$. Any positive first-quadrant sign claim must be
localized to its intended high-angle window. This coarse calculation
does not supply the true-field weak creation signs near a root.

% T1.11: finite-sample gate-cliff creation
%%%%%%%%%%%%%%%%%%%%%%%%%%%%%%%%%%%%%%%%%%%%%%%%%%%%%%%%%%%%%%%%%%%%%%%%%%%%%%%

\subsection{Finite-sample weak-basin creation by self-regulation and a gate cliff}
\label{subsec:gate-cliff-creation}

\begin{remark}[Why the creation proof has this form]
\label{rem:creation-provenance}
Several simpler routes are deliberately not used.
\begin{enumerate}[label=\textnormal{(\roman*)},leftmargin=2.4em]
\item
The diagonal strong-fit surrogate has a nearly degenerate zero-noise angular
threshold at $1-\lambda_2/\lambda_1$; finite noise unfolds this into a narrow
$O(\sigma^2)$ static fold window.  The true network creates the weak basin
well before that surrogate fold through learned residual deformation and the
finite-sample gate cliff.  The surrogate fold is therefore diagnostic, not
load-bearing.
\item
The displaced seed core alone does not supply enough signed $e_2$ output to
explain the middle lift.  Recruited strong mass is required.
\item
A direct differential inequality for a block coefficient $M_{21}$ has an
adverse $A_{21}^\rho$ feedback term.  The correct interpretation is
self-regulation: the same residual entry that drives cross-output generation
changes sign when the network-level signed cross-output reaches a gated-moment
ratio $b_\star$.
\item
An auxiliary signed input moment $J_+=\sum v_{i1}v_{i2}$ can generate the
correct sign, but transferring it to output through fine Cartesian matching
requires unnecessarily sharp mismatch control.  The final proof works
instead with the network-level signed functional $b_1^{\mathrm{eff}}$.
\item
A one-feature $p=(x_1)_+$ projection makes the displaced strong family look
like an adverse remainder.  Projecting additionally onto
$q=(x_2)_+$ reveals that the genuinely rotated first-quadrant transition
component has a favorable signed kernel.  Structured components are therefore
kept signed rather than charged by absolute norms.
\item
The attracting weak crossing is not a smooth saddle-node.  In the empirical
hard-gated field it is created by a staircase of downward cluster-$1$ gate
switches.  The branch is interpreted by one-sided/Clarke continuation.
\end{enumerate}
The mechanism is therefore fixed.  What remains for an unconditional
Theorem~1 proof is finite-sample transfer plus verification that the actual
random-initialization trajectory reaches the stated stopping-time guards;
those reachability estimates are not assumed away here.
\end{remark}

%%%%%%%%%%%%%%%%%%%%%%%%%%%%%%%%%%%%%%%%%%%%%%%%%%%%%%%%%%%%%%%%%%%%%%%%%%%%%%%
% T1.11a: self-regulated progress coordinates
%%%%%%%%%%%%%%%%%%%%%%%%%%%%%%%%%%%%%%%%%%%%%%%%%%%%%%%%%%%%%%%%%%%%%%%%%%%%%%%

\begin{lemma}[Strong progress and signed cross-output as matched network coordinates]
\label{lem:creation-progress-coordinates}
Let
\[
p(x):=(x_1)_+,
\qquad
\mathsf D_1:=\langle p,p\rangle_n,
\]
and define
\[
\boxed{
m(t):=m_1^{\mathrm{eff}}(t)
=\frac{\langle p,f_{t,1}\rangle_n}{\mathsf D_1},
\qquad
b(t):=b_1^{\mathrm{eff}}(t)
=\frac{\langle p,f_{t,2}\rangle_n}{\mathsf D_1}.
}
\tag{creation-progress-coords}\label{eq:creation-progress-coords}
\]
Set
\[
\boxed{
b_\star^{(n)}
:=
\frac{\langle p,x_2\rangle_n}{\mathsf D_1}.
}
\tag{bstar-emp}\label{eq:bstar-emp}
\]
Then at the strong axis $\phi=0$,
\[
\boxed{
A_{11}^\rho(0,t)
=
\mathsf D_1(1-m(t)),
\qquad
A_{21}^\rho(0,t)
=
\mathsf D_1\bigl(b_\star^{(n)}-b(t)\bigr).
}
\tag{self-regulation-identities}\label{eq:self-regulation-identities}
\]
Moreover the signed cross-output has the exact angle-resolved representation
\[
\boxed{
b(t)
=
\frac1{\mathsf D_1}
\sum_{i=1}^h r_i(t)z_{i2}(t)\mathsf H_1(\phi_i(t)).
}
\tag{beff-exact}\label{eq:beff-exact}
\]
\end{lemma}

\begin{proof}
At $\phi=0$ the probe gate is $\mathbf 1\{x_1>0\}$, hence
\[
A_{11}^\rho(0,t)
=
\langle p,x_1-f_{t,1}\rangle_n
=
\mathsf D_1-\langle p,f_{t,1}\rangle_n,
\]
and
\[
A_{21}^\rho(0,t)
=
\langle p,x_2-f_{t,2}\rangle_n.
\]
The first two identities follow from the definitions.  Finally
$f_{t,2}=\sum_i r_i z_{i2}q_{\phi_i}$ and
$\langle p,q_{\phi_i}\rangle_n=\mathsf H_1(\phi_i)$.
\end{proof}

\begin{remark}[The self-regulated cross-output scale]
\label{rem:bstar-scaling}
The quantity $b_\star^{(n)}$ is a signed network-level target, not the older
nonnegative block-leakage statistic used in the experimental bookkeeping.
There is no identity between those observables.

For the equal-noise population model and high signal-to-noise ratio,
\[
\mathbb E[(x_1)_+x_2]
=
\frac{\mu_2\sigma}{2\sqrt{2\pi}},
\qquad
\mathsf D_1
=
\frac{\mu_1^2}{2}+O(\sigma^2),
\]
and therefore
\[
\boxed{
b_\star^{\mathrm{pop}}
=
\frac{\mu_2\sigma}{\sqrt{2\pi}\,\mu_1^2}
\left[1+O\!\left(\frac{\sigma^2}{\mu_1^2}\right)\right].
}
\tag{bstar-scaling}\label{eq:bstar-scaling}
\]
Thus the $10^{-2}$ cross-output scale in the canonical experiment is
predicted by a gated-moment ratio and scales linearly in $\sigma$ at fixed
concept means.  The associated creation threshold in strong progress is
only approximately noise-independent: after the rescaling
$z=\mu_1\tau/\sigma$, the leading ratio depends on
$m$ and $\mu_2/\mu_1$, with finite-noise corrections of order
$\sigma^2/\mu_1^2$.
\end{remark}

\begin{remark}[The regulator reference level is not an upper barrier]
\label{rem:self-regulation-not-monotonicity}
Equation~\eqref{eq:self-regulation-identities} fixes the sign of the
strong-axis $21$ residual while $b<b_\star^{(n)}$, but it is not a
monotonicity theorem for the network quantity $b(t)$.  The exact whole-network
regulator in Proposition~\ref{sbgen:network-regulator} instead gives
\[
(b_\star^{(n)}-b)'=-\kappa(b_\star^{(n)}-b)-\chi.
\]
The signed forcing $\chi$ can be positive.  Therefore $b_\star^{(n)}$ is a
\emph{regulator reference level}, not an upper barrier: favorable forcing may
drive $b$ past $b_\star^{(n)}$ after the creation certificate has fired.
This is compatible with the later downward sweep of the repelling separator.

The preferred generation interface below no longer asks that $b$ converge to
$b_\star^{(n)}$.  Proposition~\ref{sbgen:first-hit} tracks the moving middle
lift deficit
\[
G_M(t)=B_M(m(t))-b(t)
\]
and proves a first hitting of the creation corridor from an accumulated
regulator clock.  The legacy one-sided quantity
\[
\boxed{
\varepsilon_{\mathrm{reg}}(t)
:=
[b_\star^{(n)}-b(t)]_+
}
\tag{regulation-deficit}\label{eq:regulation-deficit}
\]
remains useful for comparison with earlier budgets, but it is not the
load-bearing upper endpoint of the creation window.
\end{remark}

\begin{remark}[Do not conflate the two weak second moments]
\label{rem:strong-gated-vs-positivepart}
Two repeatedly used population quantities are structurally different:
\[
\mathbb E_n[x_2^2\mathbf 1\{x_1>0\}]
\qquad\text{and}\qquad
\mathbb E_n[(x_2)_+^2].
\]
The first carries the \emph{strong-axis gate}; the second carries the weak
positive-part gate.  On cluster~2 the strong-axis gate contributes a factor
$1/2$ at leading high-SNR order, so these quantities must not be substituted
for one another.  For the canonical equal-noise sizing their high-SNR values
are approximately
\[
1.016875
\qquad\text{and}\qquad
2.016875,
\]
respectively.  The factor-of-two difference in the dominant cluster--2 term
has already appeared in several threshold calculations, so all later
regulator and creation constants must retain the gate explicitly.
\end{remark}

\begin{lemma}[Positive-cohort cross-output generation interface]
\label{lem:positive-cohort-generation-interface}
Fix a neuron set $B_+$ on an interval and define its planar block
\[
M_{B_+}:=\sum_{i\in B_+}z_iv_i^\top,
\qquad
b_{B_+}:=(M_{B_+})_{21}
=\sum_{i\in B_+}z_{i2}v_{i1}.
\]
For each $i\in B_+$, the exact coordinate equations imply
\[
\boxed{
\begin{aligned}
\dot b_{B_+}
={}&
\sum_{i\in B_+}
A_{i,21}^\rho(v_{i1}^2+z_{i2}^2)
+
\sum_{i\in B_+}A_{i,22}^\rho v_{i1}v_{i2}
\\
&+
\sum_{i\in B_+}A_{i,11}^\rho z_{i1}z_{i2}
-
\sum_{i\in B_+}z_{i2}P_{i1}^\rho.
\end{aligned}
}
\tag{block-b-generation}\label{eq:block-b-generation}
\]
Hence, if
\[
v_{i1}>0,
\quad
v_{i2}\ge0,
\quad
z_{i1}>0,
\quad
z_{i2}\ge0,
\]
and
\[
A_{i,11}^\rho,A_{i,22}^\rho,A_{i,21}^\rho\ge0,
\qquad
A_{i,21}^\rho\ge a_{21,\star}>0,
\]
then
\[
\boxed{
\dot b_{B_+}
\ge
 a_{21,\star}
 \sum_{i\in B_+}v_{i1}^2
-
\left(\sum_{i\in B_+}z_{i2}^2\right)^{1/2}
\left(\sum_{i\in B_+}|P_{i1}^\rho|^2\right)^{1/2}.
}
\tag{block-b-lower}\label{eq:block-b-lower}
\]

The whole-network signed progress coordinate is nevertheless
\[
b(t)
=
\frac1{\mathsf D_1}
\sum_i r_i z_{i2}\mathsf H_1(\phi_i),
\]
so a block lower bound does not by itself give a lower bound on $b(t)$.
Define the complementary harmful signed mass
\[
\boxed{
\mathcal N_b(t)
:=
\frac1{\mathsf D_1}
\sum_{i\notin B_+}
r_i[-z_{i2}]_+\mathsf H_1(\phi_i).
}
\tag{b-negative-complement}\label{eq:b-negative-complement}
\]
For this cohort-based alternative, obtaining a network-level creation
certificate requires combining \eqref{eq:block-b-lower} with a one-sided
bridge to \eqref{eq:beff-exact} and a bound on $\mathcal N_b$ or a sharper
signed substitute.  The preferred whole-network route is now
Propositions~\ref{sbgen:strong-cone-regulator} and~\ref{sbgen:first-hit}:
it avoids a separate cohort bridge and reduces generation to signed forcing,
complement control, and an accumulated regulator clock.
\end{lemma}

\begin{remark}[Relation to the creation wrong-sign charge]
\label{rem:generation-creation-same-bad-mass}
The generation complement \eqref{eq:b-negative-complement} and the creation
remainder charge are not identical, but they are the same geometric type.
The creation field charges
\[
[-(z_{i2}+\tau_M z_{i1})]_+,
\]
whereas the strong-axis progress coordinate charges $[-z_{i2}]_+$.
Inside the first-quadrant input cone, the coarse mismatch guard
\eqref{eq:coarse-sign-guard} eliminates the former.  A useful final proof
should reuse the same sign/mismatch bookkeeping when controlling the latter
rather than introduce an unrelated absolute-mass estimate.
\end{remark}


%%%%%%%%%%%%%%%%%%%%%%%%%%%%%%%%%%%%%%%%%%%%%%%%%%%%%%%%%%%%%%%%%%%%%%%%%%%%%%%
% T1.11b: exact middle quadratic and joint creation gap
%%%%%%%%%%%%%%%%%%%%%%%%%%%%%%%%%%%%%%%%%%%%%%%%%%%%%%%%%%%%%%%%%%%%%%%%%%%%%%%

\begin{lemma}[Exact middle quadratic in the two progress coordinates]
\label{lem:middle-two-progress}
For $\tau>0$, let
\[
\alpha_\tau(x):=\mathbf 1\{\tau x_1+x_2>0\},
\qquad
h_\tau(x):=(x_1-\tau x_2)\alpha_\tau(x).
\]
Define
\[
G_{rs}^{(n)}(\tau):=\langle x_rx_s,\alpha_\tau\rangle_n,
\]
\[
H_{11}^{(n)}(\tau):=\langle p,x_1\alpha_\tau\rangle_n,
\qquad
H_{12}^{(n)}(\tau):=\langle p,x_2\alpha_\tau\rangle_n,
\]
and
\[
\boxed{
H_\tau^{(n)}
:=
\langle p,h_\tau\rangle_n
=
H_{11}^{(n)}-\tau H_{12}^{(n)}.
}
\]
Suppose $H_\tau^{(n)}>0$.  Define
\[
\boxed{
\begin{aligned}
b_{\mathrm{lift}}^{(n)}(\tau,m)
:={}&
\frac{1}{H_\tau^{(n)}}
\Big[
G_{12}^{(n)}
+
\bigl(G_{11}^{(n)}-G_{22}^{(n)}-mH_{11}^{(n)}\bigr)\tau
\\
&\hspace{18mm}
-
\bigl(G_{12}^{(n)}-mH_{12}^{(n)}\bigr)\tau^2
\Big].
\end{aligned}
}
\tag{blift-emp}\label{eq:blift-emp}
\]
Write
\[
f_{t,1}=m(t)p+r_{1,t},
\qquad
f_{t,2}=b(t)p+r_{2,t},
\qquad
\langle p,r_{j,t}\rangle_n=0.
\]
At
\[
\phi_M=\frac\pi2-\arctan\tau,
\]
the true cotangent quadratic satisfies exactly
\[
\boxed{
Q_M^{\mathrm{true}}(t)
=
H_\tau^{(n)}
\bigl[b(t)-b_{\mathrm{lift}}^{(n)}(\tau,m(t))\bigr]
+
\mathcal E_Q(\tau,t),
}
\tag{middle-Q-decomp}\label{eq:middle-Q-decomp}
\]
where
\[
\boxed{
\mathcal E_Q(\tau,t)
=
\langle r_{2,t}+\tau r_{1,t},h_\tau\rangle_n.
}
\tag{middle-EQ}\label{eq:middle-EQ}
\]
Finally define the empirical creation gap at a deterministic certification
level $m_{\mathrm{cert}}$ by
\[
\boxed{
\Delta_{\mathrm{cr}}^{(n)}(\tau,m_{\mathrm{cert}})
:=
b_\star^{(n)}
-
b_{\mathrm{lift}}^{(n)}(\tau,m_{\mathrm{cert}}).
}
\tag{creation-gap}\label{eq:creation-gap}
\]
\end{lemma}

\begin{proof}
Substitute $f_1=mp+r_1$ and $f_2=bp+r_2$ into the exact high-angle
quadratic
\[
Q_M
=
A_{12}^\rho\tau^2
-(A_{11}^\rho-A_{22}^\rho)\tau
-A_{21}^\rho.
\]
Collecting the coefficients of $b$ gives $H_\tau^{(n)}$ and the displayed
threshold; the remaining terms combine as
\[
\langle r_1,\tau x_1\alpha_\tau-\tau^2x_2\alpha_\tau\rangle_n
+
\langle r_2,x_1\alpha_\tau-\tau x_2\alpha_\tau\rangle_n,
\]
which is \eqref{eq:middle-EQ}.
\end{proof}

\begin{remark}[Concentrate the creation gap jointly]
\label{rem:joint-gap-concentration}
The finite-sample error entering creation is the fluctuation of
$\Delta_{\mathrm{cr}}^{(n)}$, not the sum of separate error bars for
$b_\star^{(n)}$ and $b_{\mathrm{lift}}^{(n)}$.  These quantities are smooth
functions of a common vector of gated empirical moments and share the same
dominant cluster-$1$ cross-moment fluctuation.  The parameter region should
therefore include a direct concentration bound
\[
\boxed{
\Delta_{\mathrm{cr}}^{(n)}
\ge
\Delta_{\mathrm{cr}}^{\mathrm{pop}}
-
\varepsilon_{\mathrm{gap}}.
}
\tag{joint-gap-bound}\label{eq:joint-gap-bound}
\]
This preserves the covariance cancellation that is lost by charging the two
ratios independently. Proposition~\ref{sbgen:joint-gap} now supplies an
explicit joint influence bound with the ratio remainder; use
Corollary~\ref{sbgen:gap-clipping} if the population gap is unclipped.
\end{remark}

%%%%%%%%%%%%%%%%%%%%%%%%%%%%%%%%%%%%%%%%%%%%%%%%%%%%%%%%%%%%%%%%%%%%%%%%%%%%%%%
% T1.11c: fixed auxiliary gate concentration
%%%%%%%%%%%%%%%%%%%%%%%%%%%%%%%%%%%%%%%%%%%%%%%%%%%%%%%%%%%%%%%%%%%%%%%%%%%%%%%

\begin{corollary}[Uniform gated moments with one fixed auxiliary halfspace]
\label{cor:fixed-intersection-gated-moments}
Fix a deterministic unit vector $\nu\in\mathcal S$ and define
\[
\widehat a_{rs}^{(p;\nu)}(\phi)
:=
\mathbb E_{p,n}
\left[
x_rx_s
\mathbf 1\{\widetilde u(\phi)^\top x>0\}
\mathbf 1\{\nu^\top x>0\}
\right].
\]
Let $\bar a_{rs,\star}^{(p;\nu)}(\phi)$ be the corresponding clipped
population moment.  Split the already-reserved gate confidence budget
$\delta_{\mathrm{gate}}$ into finitely many internal sub-budgets for the
single-gate event, the fixed-intersection events, the joint creation-gap
moments, and the switch-ratio DKW event. After enlarging only universal
constants in the corresponding sample-size conditions, the theorem-level
confidence allocation is unchanged and the strengthened data event satisfies
\[
\boxed{
\sup_{\phi}
\left|
\widehat a_{rs}^{(p;\nu)}(\phi)
-
\bar a_{rs,\star}^{(p;\nu)}(\phi)
\right|
\le
\varepsilon_{rs}^{(p;\nu)}
}
\tag{fixed-gate-concentration}\label{eq:fixed-gate-concentration}
\]
simultaneously for $p,r,s\in\{1,2\}$ and for the finite collection of fixed
auxiliary gates used below.

Moreover, if the moving gate changes between $\phi$ and $\phi'$, the changed
samples lie in
\[
S_{\phi,\tau_{\phi,\phi'}}\cap\{\nu^\top x>0\}
\subseteq
S_{\phi,\tau_{\phi,\phi'}},
\]
so the moving-slab modulus from
Lemma~\ref{lem:gated-moment-concentration} carries over with no larger slab
probability.
\end{corollary}

\begin{proof}
On the fixed net, multiplication by the deterministic indicator
$\mathbf 1\{\nu^\top x>0\}$ only shrinks the coordinate envelope.  For the
continuation between net points, the fixed gate never moves, so any changed
summand must already belong to the moving slab of the original proof.  The
same net plus slab argument therefore applies.  A finite union over
the fixed auxiliary halfspaces is paid inside the internal split of
$\delta_{\mathrm{gate}}$ and does not change the theorem-level probability
total.
\end{proof}

%%%%%%%%%%%%%%%%%%%%%%%%%%%%%%%%%%%%%%%%%%%%%%%%%%%%%%%%%%%%%%%%%%%%%%%%%%%%%%%
% T1.11d: signed two-feature decomposition
%%%%%%%%%%%%%%%%%%%%%%%%%%%%%%%%%%%%%%%%%%%%%%%%%%%%%%%%%%%%%%%%%%%%%%%%%%%%%%%

\begin{lemma}[Coarse first-quadrant output-sign transfer]
\label{lem:first-quadrant-output-sign}
Fix $0<\tau<1$.  If
\[
v_i=r_i(\cos\phi_i,\sin\phi_i),
\qquad
0\le\phi_i\le\frac\pi2,
\]
then
\[
\boxed{
z_{i2}+\tau z_{i1}
\ge
r_i(\sin\phi_i+\tau\cos\phi_i)
-
\sqrt{1+\tau^2}\,\|\delta_i\|.
}
\tag{output-sign-transfer}\label{eq:output-sign-transfer}
\]
Consequently
\[
\boxed{
\frac{\|\delta_i\|}{r_i}
\le
\frac{\tau}{\sqrt{1+\tau^2}}
\quad\Longrightarrow\quad
z_{i2}+\tau z_{i1}\ge0
}
\tag{coarse-sign-guard}\label{eq:coarse-sign-guard}
\]
uniformly over the entire first quadrant.
\end{lemma}

\begin{proof}
Since $z_i=v_i-\delta_i$, Cauchy--Schwarz gives the first display.  On
$[0,\pi/2]$ the minimum of
$\sin\phi+\tau\cos\phi$ is $\tau$; its interior critical point is a maximum.
\end{proof}

\begin{lemma}[Signed strong-family shape correction]
\label{lem:signed-strong-creation-kernel}
For fixed $\tau>0$, define
\[
\mathcal J_S^{(n)}(\phi;\tau)
:=
\left\langle
q_\phi-
\frac{\mathsf H_1(\phi)}{\mathsf D_1}p,
\ h_\tau
\right\rangle_n.
\]
For a fixed strong family $B_S$, its non-$p$ contribution to
\eqref{eq:middle-EQ} is exactly
\[
\boxed{
\mathcal E_S(\tau,t)
=
\sum_{i\in B_S}
r_i(z_{i2}+\tau z_{i1})
\mathcal J_S^{(n)}(\phi_i;\tau).
}
\tag{strong-shape-exact}\label{eq:strong-shape-exact}
\]
If $z_{i2}+\tau z_{i1}\ge0$ and
$0\le\phi_i\le\bar\phi_S$, define
\[
\rho_S^{(n)}(\tau,\bar\phi_S)
:=
\sup_{0\le\phi\le\bar\phi_S}
\frac{[-\mathcal J_S^{(n)}(\phi;\tau)]_+}{\mathsf H_1(\phi)}.
\]
Then
\[
\boxed{
\mathcal E_S(\tau,t)
\ge
-\rho_S^{(n)}(\tau,\bar\phi_S)
\mathsf D_1
\bigl(b_S(t)+\tau m_S(t)\bigr),
}
\tag{strong-shape-bound}\label{eq:strong-shape-bound}
\]
where $m_S,b_S$ are the strong-family contributions to the two progress
coordinates.  Moreover
\[
\mathcal J_S^{(n)}(0;\tau)=0,
\]
and
\[
\boxed{
\partial_\phi\mathcal J_S^{(n)}(0;\tau)
=
\left\langle
x_2\mathbf 1\{x_1>0\}
-b_\star^{(n)}p,
\ h_{\tau,\perp p}
\right\rangle_n,
}
\tag{strong-shape-derivative}\label{eq:strong-shape-derivative}
\]
with
$h_{\tau,\perp p}:=h_\tau-\langle p,h_\tau\rangle_np/\mathsf D_1$.
\end{lemma}

\begin{proof}
Subtracting the $p$-projection of each strong feature gives the exact first
display.  The bound follows after replacing the kernel by
$-\rho_S^{(n)}\mathsf H_1$.  Finally
$\partial_\phi q_\phi|_{\phi=0}=x_2\mathbf 1\{x_1>0\}$ almost everywhere,
and differentiating the $p$-projection produces $b_\star^{(n)}p$.
\end{proof}

\begin{lemma}[Signed two-feature transition kernel]
\label{lem:signed-two-feature-transition-kernel}
Let
\[
q(x):=(x_2)_+,
\qquad
P_{12}:=P_{\operatorname{span}\{p,q\}}
\]
be empirical $L^2$ projection.  Define
\[
\widetilde q_\phi:=(I-P_{12})q_\phi,
\qquad
h_{\tau,\perp12}:=(I-P_{12})h_\tau,
\]
\[
\boxed{
J_{12}^{(n)}(\phi;\tau)
:=
\langle\widetilde q_\phi,h_{\tau,\perp12}\rangle_n.
}
\tag{J12-emp}\label{eq:J12-emp}
\]
Equivalently, if
\[
G_{pq}^{(n)}
:=
\begin{pmatrix}
\langle p,p\rangle_n&\langle p,q\rangle_n\\
\langle p,q\rangle_n&\langle q,q\rangle_n
\end{pmatrix},
\]
\[
v_n(\phi)
:=
\binom{\langle p,q_\phi\rangle_n}{\langle q,q_\phi\rangle_n},
\qquad
w_n(\tau)
:=
\binom{\langle p,h_\tau\rangle_n}{\langle q,h_\tau\rangle_n},
\]
and assume $G_{pq}^{(n)}$ is invertible. Then
\[
\boxed{
J_{12}^{(n)}(\phi;\tau)
=
\langle q_\phi,h_\tau\rangle_n
-v_n(\phi)^\top(G_{pq}^{(n)})^{-1}w_n(\tau).
}
\tag{J12-moment-form}\label{eq:J12-moment-form}
\]

Choose a compact interior interval
\[
I_R=[\phi_S,\phi_W]\Subset(0,\pi/2)
\]
and assume the deterministic population admissibility condition
\[
\boxed{
\inf_{\phi\in I_R}
J_{12}^{\mathrm{pop}}(\phi;\tau)
\ge j_\star>0.
}
\tag{J-adm}\label{eq:J-adm}
\]
If Corollary~\ref{cor:fixed-intersection-gated-moments} and the empirical
Gram-matrix event imply
\[
\sup_{\phi\in I_R}
|J_{12}^{(n)}(\phi;\tau)-J_{12}^{\mathrm{pop}}(\phi;\tau)|
\le\varepsilon_J<j_\star,
\]
then
\[
\boxed{
J_{12}^{(n)}(\phi;\tau)\ge j_\star-\varepsilon_J>0
\qquad
\forall\phi\in I_R.
}
\tag{J12-positive}\label{eq:J12-positive}
\]
Consequently any neuron with
$\phi_i\in I_R$ and $z_{i2}+\tau z_{i1}\ge0$ contributes nonnegatively to
the two-feature remainder in the middle quadratic.
\end{lemma}

\begin{proof}
The projection formula is the standard Gram-matrix expression for
$P_{12}$.  All scalar entries are ordinary or fixed-intersection gated
moments, so the uniform empirical bound follows from
Corollary~\ref{cor:fixed-intersection-gated-moments}, the second-moment event,
and inverse perturbation for $G_{pq}^{(n)}$.  The sign conclusion is then
immediate.
\end{proof}

\begin{remark}[Canonical sign profile is sizing evidence, not a theorem]
For the canonical population parameters and the middle test angle selected
near the creation window, numerical Gaussian quadrature gives
$J_{12}^{\mathrm{pop}}(\phi;\tau)>0$ throughout the open first quadrant and a
small normalized kernel.  This motivates \eqref{eq:J-adm}, but the theorem
uses only the explicit Gaussian-integral admissibility constant $j_\star$ on
a compact interval.  No global analytic sign theorem over all SIM parameters
is assumed. Theorem~\ref{sbgen:J-positive} supplies an explicit analytic
high-SNR subregion, and Corollary~\ref{sbgen:J-finite} gives its
finite-sample transfer. Their conservative bounds do not certify the
canonical noise level by themselves.
\end{remark}

\begin{lemma}[Weak-feature and residual decomposition of the middle error]
\label{lem:middle-signed-decomposition}
Let $B_S$ be the late strong family and split
\[
r_j=s_j+u_j,
\]
where
\[
s_j
:=
\sum_{i\in B_S}
r_i z_{ij}
\left(
q_{\phi_i}-\frac{\mathsf H_1(\phi_i)}{\mathsf D_1}p
\right).
\]
Let
\[
q_\perp
:=
q-\frac{\langle p,q\rangle_n}{\mathsf D_1}p,
\qquad
D_{2|1}:=\|q_\perp\|_n^2,
\]
and decompose the non-strong remainder as
\[
u_j=\gamma_{j,R}q_\perp+e_{j,R},
\qquad
e_{j,R}\perp\operatorname{span}\{p,q\}.
\]
Define
\[
K_\tau^{(n)}:=\langle q_\perp,h_\tau\rangle_n,
\qquad
\lambda_q^{(n)}:=-\frac{K_\tau^{(n)}}{H_\tau^{(n)}}.
\]
Assume $K_\tau^{(n)}<0$, so $\lambda_q^{(n)}>0$. Then
\[
\boxed{
\begin{aligned}
Q_M^{\mathrm{true}}
={}&
H_\tau^{(n)}
\bigl[b-b_{\mathrm{lift}}^{(n)}(\tau,m)\bigr]
+
\mathcal E_S
\\
&+
K_\tau^{(n)}(\gamma_{2,R}+\tau\gamma_{1,R})
+
\langle e_{2,R}+\tau e_{1,R},h_{\tau,\perp12}\rangle_n.
\end{aligned}
}
\tag{middle-signed-full}\label{eq:middle-signed-full}
\]
The last term is the sum over non-strong neurons of
$r_i(z_{i2}+\tau z_{i1})J_{12}^{(n)}(\phi_i;\tau)$.
Thus interior first-quadrant neurons satisfying the coarse sign guard
\eqref{eq:coarse-sign-guard} are favorable under
\eqref{eq:J12-positive}; only endpoint strips, off-quadrant mass, or
wrong-sign output mass need an absolute-value charge.
\end{lemma}

\begin{proof}
This is orthogonal projection in empirical $L^2$ applied to
\eqref{eq:middle-EQ}, followed by the neuronwise identity
$f_{2,i}+\tau f_{1,i}=r_i(z_{i2}+\tau z_{i1})q_{\phi_i}$.
\end{proof}

%%%%%%%%%%%%%%%%%%%%%%%%%%%%%%%%%%%%%%%%%%%%%%%%%%%%%%%%%%%%%%%%%%%%%%%%%%%%%%%
% T1.11e: finite-sample staircase
%%%%%%%%%%%%%%%%%%%%%%%%%%%%%%%%%%%%%%%%%%%%%%%%%%%%%%%%%%%%%%%%%%%%%%%%%%%%%%%

\begin{lemma}[Exact downward jump at a cluster-$1$ gate switch]
\label{lem:gate-cliff-jump}
Fix $t$ and let $\phi_k$ be a simple cluster-$1$ gate switch for sample
$x_k$, with the sample turning on-to-off as $\phi$ increases.  Let
$u_k=\widetilde u(\phi_k)$ and $u_k^\angle=\widetilde u^\angle(\phi_k)$.
Then
\[
\boxed{
V_{\mathrm{true}}(\phi_k^+;t)
-
V_{\mathrm{true}}(\phi_k^-;t)
=
\frac1{2n}
\langle f_t(x_k),u_k\rangle
(u_k^{\angle\top}x_k).
}
\tag{gate-cliff-jump}\label{eq:gate-cliff-jump}
\]
For a cluster-$1$ switch near $\pi/2$,
$u_k^{\angle\top}x_k<0$; hence the jump is strictly downward whenever
$\langle f_t(x_k),u_k\rangle>0$.
\end{lemma}

\begin{proof}
Immediately before the switch the sample contributes
$(2n)^{-1}\langle x_k-f_t(x_k),u_k\rangle
(u_k^{\angle\top}x_k)$.  At the switch $u_k^\top x_k=0$, so the target part
vanishes and removing the sample gives the displayed jump.
\end{proof}

\begin{lemma}[One-dimensional switch-density certificate]
\label{lem:switch-density-DKW}
For cluster-$1$ samples define
\[
R_k:=-\frac{x_{k2}^{(1)}}{x_{k1}^{(1)}}.
\]
Let $F_R$ be the population CDF and $F_{R,n}$ the empirical CDF.  On the DKW
event
\[
\sup_r|F_{R,n}(r)-F_R(r)|
\le
\varepsilon_{\mathrm{DKW}}
:=
\sqrt{\frac{\log(2/\delta_{\mathrm{sw}})}{2n}},
\]
every interval $I=[a,b]$ satisfies
\[
\boxed{
\frac1n\#\{k:R_k\in I\}
\ge
F_R(b)-F_R(a)-2\varepsilon_{\mathrm{DKW}}.
}
\tag{switch-count-DKW}\label{eq:switch-count-DKW}
\]
If the ratio density obeys $f_R\ge f_\star$ on $I$, then
\[
\frac1n\#\{k:R_k\in I\}
\ge
f_\star|I|-2\varepsilon_{\mathrm{DKW}}.
\]
No higher-dimensional VC argument is needed because the switch set is
indexed by the scalar statistic $R_k$.
\end{lemma}

%%%%%%%%%%%%%%%%%%%%%%%%%%%%%%%%%%%%%%%%%%%%%%%%%%%%%%%%%%%%%%%%%%%%%%%%%%%%%%%
% T1.11f: deterministic gate-cliff creation assembly
%%%%%%%%%%%%%%%%%%%%%%%%%%%%%%%%%%%%%%%%%%%%%%%%%%%%%%%%%%%%%%%%%%%%%%%%%%%%%%%

\begin{proposition}[Deterministic gate-cliff creation at a certification time]
\label{prop:gate-cliff-creation}
Fix a deterministic progress ceiling
$m_{\max}>m_{\mathrm{prog}}$ and a certification time
$t_{\mathrm{cert}}\ge t_{\sharp}$ satisfying
\[
m(t_{\mathrm{cert}})\le m_{\max}.
\]
Choose
\[
\phi_L<\phi_M<\frac\pi2<\phi_+,
\qquad
\phi_M=\frac\pi2-\arctan\tau_M,
\]
with $\tau_M>0$, and let
$H_M:=H_{\tau_M}^{(n)}>0$.
Assume the following conditions at $t=t_{\mathrm{cert}}$.

\paragraph{(A) Data-side transition geometry.}
The interior two-feature population kernel satisfies
\eqref{eq:J-adm}, and its finite-sample transfer satisfies
\eqref{eq:J12-positive}.  All empirical moments entering
$b_{\mathrm{lift}}^{(n)}(\tau_M,m(t_{\mathrm{cert}}))$ are evaluated on
the same fixed data event.

\paragraph{(B) Dynamical guards.}
The B1 branch-contraction stage has placed the strong family inside the late
holding cone, Proposition~\ref{prop:late-strong-holding} carries it through
$t_{\mathrm{cert}}$, and all first-quadrant neurons charged by the signed
transition argument satisfy
\[
\boxed{
\frac{\|\delta_i(t_{\mathrm{cert}})\|}{r_i(t_{\mathrm{cert}})}
\le
\eta_{\mathrm{sign}}
<
\frac{\tau_M}{\sqrt{1+\tau_M^2}}.
}
\tag{GC-sign-guard}\label{eq:GC-sign-guard}
\]
No convergence of $b(t)$ to $b_\star^{(n)}$ is assumed here.  The preferred
predecessor is the first-hitting regulator certificate in
Proposition~\ref{sbgen:first-hit}.

\paragraph{(C) Direct signed middle-lift budget.}
Define the realized middle lift margin
\[
\boxed{
\Delta_M^{\mathrm{hit}}
:=
b(t_{\mathrm{cert}})
-
b_{\mathrm{lift}}^{(n)}
\bigl(\tau_M,m(t_{\mathrm{cert}})\bigr).
}
\tag{GC-lift-margin}\label{eq:GC-lift-margin}
\]
The strong-family signed kernel obeys, in $b$-units,
\[
\frac{[-\mathcal E_S]_+}{H_M}
\le
\varepsilon_S.
\]
The weak-feature coefficient obeys
\[
\boxed{
\Gamma_W(t_{\mathrm{cert}})
:=
(\gamma_{2,R}+\tau_M\gamma_{1,R})_+
\le
\Gamma_W^{\max}.
}
\tag{GC-weak-timing}\label{eq:GC-weak-timing}
\]
The endpoint-strip, off-quadrant, wrong-sign output, and residual finite-sample
charges sum to
$\varepsilon_{\mathrm{strip}}+\varepsilon_{\mathrm{bad}}$.
Assume
\[
\boxed{
\Delta_M^{\mathrm{hit}}
-
\varepsilon_S
-
\lambda_q^{(n)}\Gamma_W^{\max}
-
\varepsilon_{\mathrm{strip}}
-
\varepsilon_{\mathrm{bad}}
>
\frac{(1+\tau_M^2)c_M}{H_M}
}
\tag{GC-middle-budget}\label{eq:GC-middle-budget}
\]
for some $c_M>0$.  Then
\[
\boxed{
V_{\mathrm{true}}(\phi_M;t_{\mathrm{cert}})\ge c_M.
}
\tag{GC-middle-sign}\label{eq:GC-middle-sign}
\]

\paragraph{(D) Lower and post-cliff negative endpoints.}
At the fully-on lower point write
\[
Q_-(\tau_L)
=
a_{12}\tau_L^2-d_-\tau_L-a_{21},
\qquad
d_-:=a_{11}-a_{22}.
\]
Assume
\[
\boxed{
d_-\tau_L
>
|a_{12}|\tau_L^2+|a_{21}|+(1+\tau_L^2)c_L
}
\tag{GC-left}\label{eq:GC-left}
\]
for some $c_L>0$.  Then
$V_{\mathrm{true}}(\phi_L;t_{\mathrm{cert}})\le-c_L$.

At the post-cliff point, assume the coordinate-envelope geometry turns every
cluster-$1$ gate off and every cluster-$2$ gate on.  Hence
$A^\rho=\mathcal R_2/2$ there.  Writing
\[
Q_+(\tau_R)
=
a_{12}^+\tau_R^2-d_+\tau_R-a_{21}^+,
\qquad
d_+:=a_{22}^+-a_{11}^+,
\]
assume
\[
\boxed{
d_+\tau_R
>
|a_{12}^+|\tau_R^2+|a_{21}^+|+(1+\tau_R^2)c_R
}
\tag{GC-right}\label{eq:GC-right}
\]
for some $c_R>0$.  Then
$V_{\mathrm{true}}(\phi_+;t_{\mathrm{cert}})\le-c_R$.

\paragraph{(E) Gate-cliff staircase.}
Let $I_{\mathrm{cliff}}\subset[\phi_M,\phi_+]$ be the cluster-$1$
on-to-off switch band.  Suppose every smooth activation cell satisfies
\[
\partial_\phi V_{\mathrm{true}}\le L_{\mathrm{cell}},
\]
and every cluster-$1$ switch satisfies
\[
\frac12
\langle f_{t_{\mathrm{cert}}}(x_k),u_k\rangle
|u_k^{\angle\top}x_k|
\ge j_{\mathrm{sw}}>0.
\]
Suppose also that the DKW event and the ratio-density lower bound imply that
every subinterval $J\subset I_{\mathrm{cliff}}$ of width at least $w_0$
has
\[
\frac{N_{\mathrm{sw}}(J)}{n}
\ge
\rho_{\mathrm{sw}}^{\mathrm{eff}}|J|,
\qquad
\rho_{\mathrm{sw}}^{\mathrm{eff}}
:=
\rho_{\mathrm{sw}}
-
\frac{2\varepsilon_{\mathrm{DKW}}}{w_0}>0.
\]
If
\[
\boxed{
j_{\mathrm{sw}}\rho_{\mathrm{sw}}^{\mathrm{eff}}
>
L_{\mathrm{cell}}+\lambda_{\mathrm{cliff}}
}
\tag{GC-staircase}\label{eq:GC-staircase}
\]
for some $\lambda_{\mathrm{cliff}}>0$, then the selected empirical field has
a coarse strictly downward staircase across the cliff band.

\paragraph{Conclusion.}
Under the preceding conditions,
\[
\boxed{
V_{\mathrm{true}}(\phi_L;t_{\mathrm{cert}})<0,
\qquad
V_{\mathrm{true}}(\phi_M;t_{\mathrm{cert}})>0,
\qquad
V_{\mathrm{true}}(\phi_+;t_{\mathrm{cert}})<0.
}
\tag{GC-three-signs}\label{eq:GC-three-signs}
\]
There is a lower repelling selected crossing between $\phi_L$ and $\phi_M$.
Across the gate cliff there is an upper attracting selected crossing, unique
up to the resolution $w_0$ under \eqref{eq:GC-staircase}.  At a smooth
crossing the usual derivative sign gives transversality; at a gate jump the
crossing is interpreted by the filled one-sided/Clarke vertical segment.
Thus the learned residual has created a weak-concept attraction basin.
\end{proposition}

\begin{proof}
The middle sign follows directly from
Lemma~\ref{lem:middle-signed-decomposition}, the realized lift margin
\eqref{eq:GC-lift-margin}, the signed strong and transition kernels, and the
budget \eqref{eq:GC-middle-budget}.  The endpoint signs are the exact
cotangent quadratics.  For the upper crossing,
Lemma~\ref{lem:gate-cliff-jump} gives a downward jump $j_{\mathrm{sw}}/n$ per
switch and Lemma~\ref{lem:switch-density-DKW} supplies the number of switches.
The total downward jump mass therefore beats the positive smooth-cell drift
by \eqref{eq:GC-staircase}.  The lower negative-to-positive and upper
positive-to-negative sign changes give the stated selected crossings.
\end{proof}

\begin{remark}[Canonical diagnostics versus the rigorous high-SNR ledger]
\label{rem:canonical-creation-sizing}
The familiar numerical table at $(\mu_1,\mu_2,\sigma)=(3,2,0.15)$ has
$e=\sigma/\mu_2=0.075$ and is retained only as an experimental sizing
diagnostic.  It is not a point in the rigorous signed-feedback packet, whose
current sufficient assumptions include $e\le1/40=0.025$.  In particular the
values
\[
 \Delta_{\rm cr}^{\rm pop}(0.47)\simeq3.10\times10^{-3},\qquad
 \Delta_{\rm cr}^{\rm pop}(0.50)\simeq6.15\times10^{-3}
\]
must not be combined with the $e\le1/40$ feedback theorem as though they
belonged to one primitive parameter point.

For the common high-SNR analysis set
\[
 \tau_M=\frac{ze}{\varrho},\qquad \varrho=\frac{\mu_1}{\mu_2}.
\]
At fixed $(z,m)$,
\[
\boxed{
 \frac{\Delta_{\rm cr}^{\rm pop}(z,m)}e
 \longrightarrow
 \frac{\varphi(0)}{\varrho^2}
 -\frac{g(z)}{\varrho\Phi(z)}
 +\frac{z}{\varrho^3\Phi(z)}
 +\frac{mz}{\varrho},
 \qquad g(z)=\varphi(z)+z\Phi(z).
}
\tag{CR0}\label{eq:CR0}
\]
For $\varrho=3/2$, optimizing the right side gives the diagnostics
\[
\begin{array}{c|c|c}
 m&z_*&\Delta_{\rm cr}^{\rm pop}/e\\ \hline
 0.50&2.16&0.0811\\
 0.52&2.38&0.1113\\
 0.53&2.53&0.1276\\
 0.54&2.74&0.1451\\
 0.55&3.13&0.1645
\end{array}
\]
and therefore positive $O(e)$ population creation margins at the same
high-SNR scale used by the feedback proof.

The creation/progress clock gives a stricter zero-charge comparison.  In the
ideal high-SNR regulator scaling, the worst deficit at progress $m$ is
$\varphi(0)\varrho^{-2}\sqrt{1-m}$ in $e$-normalized units.  Define
\[
 \Delta_{\rm PC}^{(0)}(z,m)
 :=\mathrm{CR}_0(z,m)-\frac{\varphi(0)}{\varrho^2}\sqrt{1-m}.
\]
For $m_*=0.535$ and $z_*=2.62736$,
\[
 \mathrm{CR}_0(z_*,m_*)=0.136200,\qquad
 \Delta_{\rm PC}^{(0)}(z_*,m_*)=0.015292.
\]
The zero-charge lower feasibility threshold is approximately $m=0.5296$.
The upper threshold is the single rate ratio
\[
 m<1-\frac{\ell_2}{\ell_1}
   =1-\varrho^{-2}
   =\frac59.
\]
Thus the uncharged late interval is approximately
$(0.5296,5/9)$; it is not appropriate to select $m_{\max}$ before the
normalized creation charges are priced.

Equivalently, the largest possible total normalized creation charge is the
zero-charge progress margin evaluated at the $5/9$ ceiling, approximately
$0.0594$.  The previously quoted weak-feature cost
$2.8\times10^{-3}$ at $e=0.075$ corresponds to $0.0373$ after division by
$e$, but its actual high-SNR order must be established before inserting that
number into the rigorous ledger.  Determining whether this charge is truly
$O(e)$ or is $o(e)$ is a load-bearing remaining calculation.  Sampling,
clipping, strong-shape, middle-residual, endpoint-strip and staircase charges
must be entered in the same normalized ledger.
\end{remark}


\subsection{Independent weak-branch comparison and projected mass}
\begin{proposition}[Smooth reference branch, robust signs, and tracking]
\label{sbnew:weak-comparison}
Let $B(\phi,t)$ be a reference scalar field on a compact angular window
and a time interval $[a,H]$, continuously differentiable in $\phi$ and
Lipschitz in time. Suppose fixed $\phi_L<\phi_M<\phi_R$ satisfy
\[
 B(\phi_L,t)\le-c,\qquad B(\phi_M,t)\ge c,
 \qquad B(\phi_R,t)\le-c
\]
uniformly, and that root neighborhoods have respectively
$\partial_\phi B\ge\lambda_R>0$ and
$\partial_\phi B\le-\lambda_W<0$, with opposite endpoint signs
inside each neighborhood. There are unique roots in these neighborhoods,
a repelling reference branch and an attracting reference branch $b_W(t)$.
If the time-Lipschitz constant is $M_t$, then
\begin{equation}\label{sbnew:branch-speed}
 |\dot b_W|\le v_W:=M_t/\lambda_W\quad\hbox{a.e.}
\end{equation}

Let a fixed cohort satisfy the exact perturbed angular equations
$\dot\phi_i=B(\phi_i,t)+e_i(t)$, $|e_i|\le e_W$.
On the tube $|\phi-b_W(t)|\le R_W$, suppose
$\partial_\phi B\le-\lambda_W$. If
\begin{equation}\label{sbnew:tracking-cond}
 \phi_i(a)\in[b_W(a)-R_W,b_W(a)+R_W],\quad
 E_W:=\frac{e_W+v_W}{\lambda_W}<R_W,
\end{equation}
then the tube is invariant and
\begin{equation}\label{sbnew:tracking}
 |\phi_i(t)-b_W(t)|\le
 E_W+(|\phi_i(a)-b_W(a)|-E_W)e^{-\lambda_W(t-a)}.
\end{equation}
For $E_W<\eta_W<R_W$, all these neurons enter the smaller tube by
\begin{equation}\label{sbnew:capture-time}
 t_{\rm cap}\le a+\lambda_W^{-1}
 \left[\log\frac{R_W-E_W}{\eta_W-E_W}\right]_+,
\end{equation}
provided the right side is at most $H$.
\end{proposition}
\begin{proof}
The intermediate value theorem and strict monotonicity give each reference
root and uniqueness in its stated neighborhood. For two times, apply
the mean value theorem in angle and the Lipschitz bound in time to their
root equations. This bounds root displacement by
$M_t|t-s|/\lambda_W$, proving \eqref{sbnew:branch-speed}.
For $y_i=|\phi_i-b_W|$, the mean value theorem gives
$\dot y_i\le-\lambda_Wy_i+e_W+v_W$ almost everywhere.
The strict inequality at $y_i=R_W$ excludes first exit. Integrating gives
\eqref{sbnew:tracking}; solving for $y_i\le\eta_W$ proves
\eqref{sbnew:capture-time}.
\end{proof}

\paragraph{Finite-sample meaning and the interface from gate-cliff creation.}
If Proposition~\ref{prop:gate-cliff-creation} is invoked after its explicit
predecessor guards have been proved, it supplies a finite-sample lower
repelling selected crossing and an upper attracting gate-cliff crossing at
the creation time. This is deliberately a one-sided/Clarke statement rather
than a claim that the empirical hard-gated field is globally smooth.

To invoke the smooth reference comparison above on a later interval, one
must still construct a reference field $B$, prove its approximation error and
time modulus, and connect its attracting branch to the selected gate-cliff
crossing. Alternatively the theorem may continue with generalized selected
equilibria, but their one-sided continuation and transversality require the
corresponding comparison. Thus the creation \emph{mechanism} and deterministic
sign package are explicit, while the later continuation interface remains a
separate obligation.

The inclusion in \eqref{sbnew:tracking-cond} is also substantive: early weak
retention inside a larger interval does not prove that a positive fixed cohort
lies on the capturable side of the new separatrix. The comparison above
finishes capture inside an established attracting tube; transport of a fixed
retained cohort onto that side remains part of the capture proof.

\begin{lemma}[Weak self-field, radius compatibility, and mass time]
\label{sbnew:mass}
For a fixed cohort $W$, write
$L_W=\sum_{i\in W}(r_i^2+s_i^2)$ and
$\Xi_W=\sum_{i\in W}r_is_i\le L_W/2$.
For any unit test directions and selected gate, its output contribution
to the angular field has magnitude at most
\begin{equation}\label{sbnew:weak-self}
 \lambda\Xi_W\le\lambda L_W/2.
\end{equation}
Thus $\lambda L_W/2<c-e_{\rm field}$ is sufficient to keep removal of
this cohort from reversing a sign margin of $c-e_{\rm field}$.

Suppose, on $[t_{\rm cap},H]$, a fixed $W$ has $|W|\ge c_Wh$,
$r_i(t_{\rm cap})\ge r_0>0$, $\dot r_i/r_i\ge\gamma_W>0$,
and $\|\delta_i\|/r_i\le\eta_\delta<1$ for every member.
Then
\begin{align}
 \mathsf R_W(t)&\ge c_Wh r_0^2e^{2\gamma_W(t-t_{\rm cap})},
 \label{sbnew:weak-mass}\\
 \operatorname{tr}\sum_{i\in W}z_iv_i^\top
 &\ge(1-\eta_\delta)\mathsf R_W(t),\label{sbnew:trace-mass}\\
 \angle(v_i,z_i)&\le
 2\arcsin\!\left(\min\{1,\eta_\delta/(2\sqrt{1-\eta_\delta})\}\right).
 \label{sbnew:output-angle}
\end{align}
Given an initialization-independent target mass $M_0>0$, define
\begin{equation}\label{sbnew:tW}
 t_W=t_{\rm cap}+
 \frac1{2\gamma_W}\left[\log\frac{M_0}{c_Wh r_0^2}\right]_+.
\end{equation}
If $t_W\le H$, then $\mathsf R_W(t_W)\ge M_0$.
The same conclusions apply to any fixed enlarged block whose every
additional member has the same radius, mismatch, and tracking guarantees.
An angular-only membership test does not establish these extra guarantees.
\end{lemma}
\begin{proof}
Cauchy--Schwarz gives $\|f_W\|_n\le\sqrt\lambda\,\Xi_W$;
the test tangential feature has norm at most $\sqrt\lambda$, proving
\eqref{sbnew:weak-self}. Integrate the radial lower bound and sum to prove
\eqref{sbnew:weak-mass}. Since
$v_i\cdot z_i=r_i^2-v_i\cdot\delta_i\ge(1-\eta_\delta)r_i^2$,
\eqref{sbnew:trace-mass} follows. Also $s_i\ge(1-\eta_\delta)r_i$ and
\[
 \|\delta_i\|^2=(r_i-s_i)^2+
 4r_is_i\sin^2\!\left(\tfrac12\angle(v_i,z_i)\right),
\]
which proves \eqref{sbnew:output-angle}. Formula \eqref{sbnew:tW} is the
direct solution of the mass inequality.
\end{proof}

\begin{remark}[Intended pure-Q2 to gate-cliff capture route]
\label{rem:weak-capture-route}
The initialization lemma permits $J_W^{\rm seed}$ to be any deterministic
positive-width arc, while T1.8--T1.9 require only retention inside a larger
$J_W^{\rm ret}$ before the residual-created basin exists.  The intended
unconditional capture proof chooses this fixed weak seed arc inside the
cluster--2-pure quadrant-II sector.  In that sector the $e_1$ transverse
energy is conserved at leading order while the $e_2$ coordinate amplifies;
moreover the coarse learned-strong field satisfies
$V_0(\phi;m)<0$ and $\partial_mV_0(\phi;m)<0$ throughout
$\pi/2<\phi<\pi$.  Thus increasing strong progress drives the retained weak
cohort toward $e_2$, not away from the future gate-cliff basin.

The missing quantitative step is an ordering theorem: the fixed retained
cohort must survive in the pure sector until creation, cross onto the
capturable side of the newly created repelling separator, and enter the
attracting tube before its continuation guards expire.  Once this is proved,
Proposition~\ref{sbnew:weak-comparison} supplies tracking and
Lemma~\ref{sbnew:mass} supplies any fixed initialization-independent target
weak mass $M_0>0$ after the logarithmic mass-growth time.  No separate
positive-mass mechanism is missing after capture.
\end{remark}

\subsection{Completed selector and moving-boundary implications}
\begin{lemma}[Fixed-family selectors with ambient correction]
\label{sbnew:selectors}
Let $A$ be a fixed family on an interval $I$, with
$r_i(t)\ge r_*(t)>0$ and
$\operatorname{dist}_\angle(\phi_i(t),c(t))\le\eta(t)\le\pi$.
Set $d_\eta(t)=2\sin(\eta(t)/2)$ and write $x=x_{\mathcal S}+x_\perp$.
Then the gates of every member of $A$ agree with
$\mathbf1_{u(c(t))\cdot x_{\mathcal S}>0}$ whenever
\begin{equation}\label{sbnew:ambient-boundary}
 |u(c(t))\cdot x_{\mathcal S}|>
 d_\eta(t)\|x_{\mathcal S}\|+
 \frac{\varepsilon}{r_*(t)}\|x_\perp\|.
\end{equation}
For the training data, define
$\tau(t)=Xd_\eta(t)$ and
$B(t)=\{x:|u(c(t))\cdot x|\le\tau(t)\}$.
The existing master slab event gives, simultaneously for all $t$,
\begin{equation}\label{sbnew:boundary-mass}
 \mathbb P_n(B(t))\le
 \min\left\{1,\sqrt{2/\pi}\frac{\tau(t)}{\sigma_{\min}}
                   +\Delta_{\rm VC}\right\}.
\end{equation}
For two fixed families, the mass of the union of their boundary layers
is bounded by the sum of the two right sides. No time discretization and
no additional probability budget are required.
\end{lemma}
\begin{proof}
Input-perpendicular conservation gives
$\|u_{i,\perp}(t)\|=\|u_{i,\perp}(0)\|\le\varepsilon$.
Consequently
\[
 \left|\frac{u_i(t)\cdot x}{r_i(t)}-u(c(t))\cdot x_{\mathcal S}\right|
 \le d_\eta(t)\|x_{\mathcal S}\|
           +\varepsilon\|x_\perp\|/r_*(t).
\]
Under \eqref{sbnew:ambient-boundary} the two scalar signs coincide.
Training samples have $x_\perp=0$ and norm at most $X$.
Apply the uniform slab bound separately to both clusters at the actual
data-dependent angle and width, and average. This proves
\eqref{sbnew:boundary-mass}.
\end{proof}

\paragraph{Fixed blocks and the probability total.}
Choose $A_1,A_2$ once at a specified selection time from neurons with all
the needed simultaneous radius, spread, and mismatch properties; retain
these indices thereafter. If the strong and weak tubes are disjoint at
that time, so are the blocks. Count and mass guarantees follow only from
the proved recruited/captured cohorts and Lemma~\ref{sbnew:mass}; selecting
all neurons in an angular tube is insufficient to certify a uniform
normalized mismatch bound. Under established family spread, the selector
and empirical boundary conclusions are now fully proved by
Lemma~\ref{sbnew:selectors}. Finally the existing allocation gives
\[
 \mathbb P(\mathcal H_{\rm data}\cap\mathcal H_{\rm init}
                    \cap\mathcal G_{\rm gate})
 \ge1-(5\delta/8+\delta/8+\delta/4)=1-\delta.
\]
This probability identity is complete; the dynamical conclusions on that
event are not.

\subsection{Dependency chain, conditions, and required statement changes}
\paragraph{Dependency chain.}
T1.1 and T1.1b feed the bias-aware T1.8--T1.9 handoff and the residual
constants. T1.2--T1.4 feed the global and perpendicular bounds. These feed
the pure-sector exit and the corrected B1 inner cone/spread comparison.
The next open implication remains uniform whole-network strong progress with
signed self-field and mismatch control through the B1 branch-locking level.
If that closes, Proposition~\ref{prop:late-strong-holding} supplies the B2
continuation interface.  The regulator package must then prove the signed
forcing conditions in Lemmas~\ref{sbgen:strong-cone-regulator} and
\ref{sbgen:regulator-transition-kernel} and accumulate the clock required by
Proposition~\ref{sbgen:first-hit} before the weak-feature timing allowance or
B2 guards expire.  If the $c_{12}^O$ feedback channel is retained, the same
interval must also remain below the weak off-diagonal stopping face
$t_\chi$ from Corollary~\ref{sbcoh:chi-deadline}, with the own-tilt spread
margin of Lemma~\ref{sbcoh:spread}.  Proposition~\ref{prop:gate-cliff-creation} then supplies
the finite-sample creation sign pattern and attracting gate-cliff crossing. Separatrix-side
capture, later selected/reference-branch continuation, and the retained weak
cohort's transport remain open. Once those quantitative outputs are supplied,
the tracking and mass comparisons above feed the proved selector and
moving-boundary conclusions, followed by fixed-block bookkeeping.

\paragraph{New conditions, with their logical status.}
The deterministic early conditions are \eqref{sbnew:kappa-cert} and the
source's early/post inequalities with the bias additions stated above. The
local conditional B1 requirements are \eqref{sbnew:inner-purity}, the
bulk/tail budgets in Lemma~\ref{sbnew:residual-constants},
\eqref{sbnew:cone-conditions}, the pointwise bound
\eqref{sbnew:ratio-error}, and \eqref{sbnew:jump-modulus}; no claim that
these dynamical budgets already follow from primitive parameters is made.

The B2/creation interface additionally requires the face inequalities
\eqref{eq:late-hold-plus}--\eqref{eq:late-hold-minus}, the jointly
concentrated creation gap \eqref{eq:joint-gap-bound}, the deterministic
population kernel margin \eqref{eq:J-adm} and its fixed-intersection
finite-sample transfer, the coarse sign guard \eqref{eq:GC-sign-guard}, the
signed regulator conditions \eqref{sbgen:regulator-adm},
\eqref{sbgen:u-positive-condition}, and
\eqref{sbgen:complement-forcing}, the first-hitting corridor/clock conditions
\eqref{sbgen:corridor-positive}--\eqref{sbgen:first-hit-clock}, the
weak-feature timing guard \eqref{eq:GC-weak-timing}, the realized middle
budget \eqref{eq:GC-middle-budget}, the two endpoint signs, and the DKW
staircase condition \eqref{eq:GC-staircase}.  For the stronger $c_{12}^O$
route one additionally requires the field transfer
\eqref{sbcoh:field-transfer}, a positive own-tilt margin
\eqref{sbcoh:own-margin}, and the weak off-diagonal ordering condition
\eqref{sbcoh:chi-before-hit-clock}.  The genuinely trajectory-timed
quantities are the weak-feature coefficient, the weak off-diagonal
$\chi_W$, adverse complement forcing, and the accumulated regulator clock;
the other middle terms are data constants, sampling allowances, or guards
inherited from earlier stages.

The weak comparison then requires the reference-field sign/transversality
and time bounds, initial tube inclusion, \eqref{sbnew:tracking-cond}, a
deadline in \eqref{sbnew:capture-time}, a self-field margin in
\eqref{sbnew:weak-self}, and positive radius/mismatch constants with
\eqref{sbnew:tW} at most the available horizon. These are hypotheses of
proved auxiliary implications, not a proposed replacement for the
unconditional theorem's admissible parameter region. There is no new
condition for the selector implication beyond its stated geometric
hypotheses and the existing master event.

\paragraph{Required theorem-statement corrections.}
Replace ``deterministic stopping times'' by sample-dependent hitting or
selection times, with deterministic upper bounds where proved. The creation
statement must use the finite-sample selected-crossing language of
Proposition~\ref{prop:gate-cliff-creation}: at a gate jump the selected zero
is represented by the filled one-sided/Clarke vertical segment, not by a
fictitious globally smooth empirical root. A later continuously varying
reference branch is a separate comparison object.

Count the strong family's B1 entry and contraction time before assigning its
locking start; late recruitment cannot retroactively imply locking for every
recruited member from the original seed-lock time. Distinguish the B1 branch-locking level $m_{\mathrm{prog}}$ from the
later B2 progress ceiling $m_{\max}$.  Creation is certified at the first
middle-threshold hitting time of Proposition~\ref{sbgen:first-hit}, not
necessarily when $m$ hits a prescribed value.  Pay the strong bias floor and
the initial transient in spread bounds. Assert output-direction tracking only on intervals where
normalized mismatch is proved, using \eqref{sbnew:output-angle}. State
selector agreement either on bounded in-span inputs or with
\eqref{sbnew:ambient-boundary}; a fixed-width slab cannot cover all unbounded
ambient inputs. Define final blocks once using all their required properties.
No positive-fraction weak-cohort or nonvanishing-mass conclusion is weakened
here.

\paragraph{Manual audit priorities.}
Audit the learned-output bounds needed to make
\eqref{sbnew:alpha11}--\eqref{sbnew:beta21} useful, the weighted forcing
$Z_{[h]}$, and the signed self-field faces over the B1 stopping interval.
For the default B2/creation stage, audit the asymmetric physical holding
faces, the late wrong-sign coordinate and mismatch bounds, the sharp
mass-weighted regulator clock \eqref{sbgen:mass-clock}, the signed strong
forcing conditions \eqref{sbgen:regulator-adm} and
\eqref{sbgen:u-positive-condition}, the adverse complement bound
\eqref{sbgen:complement-forcing}, the weak-feature timing coefficient
$\Gamma_W$, the joint creation-gap concentration, fixed-intersection kernel
margin, endpoint strips, output-sign complement, and gate-switch amplitude.
The $c_{12}^O$ weak-rotation/own-spread package is an optional strengthened
route, not a default predecessor.
The retained weak cohort's position relative to the newly created separatrix
and the later selected/reference-branch continuation remain separate capture
issues. Nonemptiness of the final primitive parameter region and the strict
ordering of all final times remain unproved. The experiment values and
successful sizing do not discharge these points.

\begin{remark}[The theorem scope used in this review copy]
\label{rem:scope-A-specialization}
\textbf{Historical scope proposal, reopened by the Stage 0 audit in Section~\ref{sbaudit:scope}.} The following paragraph records the incoming review scope, not a proof that the original noisy reversal goal needs a separate major theory.

The noisy many-neuron headline target is the specialization theorem below,
followed by the proved concept-block reduction and the predictive OOD normal
form.  A general noisy many-neuron reversal-sign theorem is intentionally not
a load-bearing claim in this scope.  The rigorous sign-reversal theorem in
this manuscript is the restricted noiseless fixed-gate result; extending its
sign reversal to the full noisy isotropic model is stated separately as open.
Accordingly no lemma needed only for that general reversal is treated as a
predecessor of the specialization theorem.
\end{remark}

\begin{conjecture}[Two-concept gate emergence and specialization: unproved target]
\label{thm:two-concept-specialization}

Consider the two-layer ReLU network
\[
f_t(x)
=
\sum_{i=1}^h
w_i(t)\sigma(u_i(t)^\top x)
\]
trained by gradient flow on the empirical two-concept SIM dataset
\[
x_k^{(1)}=\mu_1 e_1+\xi_k^{(1)},
\qquad
x_k^{(2)}=\mu_2 e_2+\xi_k^{(2)},
\qquad
\mu_1>\mu_2>0,
\]
with balanced isotropic initialization
\[
w_i(0)=u_i(0)=\varepsilon \widehat u_i,
\qquad
\widehat u_i\stackrel{\mathrm{i.i.d.}}{\sim}\mathrm{Unif}(\mathbb S^{d-1}).
\]

Let
\[
\mathcal S=\operatorname{span}\{e_1,e_2\},
\qquad
v_i=P_{\mathcal S}u_i,
\qquad
z_i=P_{\mathcal S}w_i,
\]
and, whenever these vectors are nonzero, write
\[
r_i=\|v_i\|,
\qquad
s_i=\|z_i\|,
\qquad
\widetilde u_i=\frac{v_i}{r_i},
\qquad
\widetilde w_i=\frac{z_i}{s_i}.
\]

There exist universal constants and an explicit two-concept
parameter region
\[
\boxed{
\mathfrak R_{\mathrm{2c}}
=
\mathfrak R_{\mathrm{2c}}
\left(
\mu_1,\mu_2,\sigma_1,\sigma_2,
d,h,n,\varepsilon,\delta
\right).
}
\]

% Probability bookkeeping:
% the final 1-\delta claim is discharged after the gated-moment event is
% constructed and the total theorem-level event is assembled.
If the model parameters lie in
\(\mathfrak R_{\mathrm{2c}}\), then with probability at least
\(1-\delta\), there exist sample-dependent times
\[
0<t_{\mathrm{lock}}
<
t_{\mathrm{create}}
<
t_{\mathrm{cap}}
\le
t_W
<
T
\]
and fixed neuron sets
\[
A_1,A_2\subset[h],
\qquad
A_1\cap A_2=\varnothing,
\]
such that:

\begin{enumerate}[label=\textnormal{(\roman*)}]

\item
\textbf{Strong-concept locking.}
There exists a scale \(\eta_1^u\ll1\) such that
\[
\angle(\widetilde u_i(t),e_1)
\le
C\eta_1^u,
\qquad
\angle(\widetilde w_i(t),e_1)
\le
C\eta_1^u
\]
for every \(i\in A_1\) and every
\(t\in[t_{\mathrm{lock}},T]\).
Moreover \(A_1\) carries nonvanishing projected mass throughout this
interval.

\item
\textbf{Creation of a weak-concept angular basin.}
At time $t_{\mathrm{create}}$ the true empirical selected angular field has a
lower repelling selected crossing and an upper attracting gate-cliff crossing
in a fixed high-angle window near $e_2$. At a smooth crossing the usual
derivative sign supplies transversality; at a gate boundary the selected zero
is understood by the filled one-sided/Clarke vertical segment, with the upper
crossing certified by a downward gate-cliff staircase.

On $[t_{\mathrm{create}},T]$ these selected crossings admit one-sided
continuations
\[
\phi_R(t)<\phi_W(t)
\]
provided the later continuation/capture guards remain active. When the
continuation lies inside smooth gate cells, the corresponding derivative
margins satisfy
\[
\partial_\phi V_{\mathrm{true}}(\phi_R(t);t)\ge\lambda_R>0,
\qquad
\partial_\phi V_{\mathrm{true}}(\phi_W(t);t)\le-\lambda_W<0.
\]
A smooth reference branch used for tracking is a comparison object and need
not coincide pointwise with a globally smooth empirical root.

\item
\textbf{Positive weak-cohort capture.}
There exists a fixed initialization cohort
\[
W_{\mathrm{cap}}\subset[h],
\qquad
|W_{\mathrm{cap}}|\ge c_{\mathrm{cap}}h,
\]
whose members survive until the basin-creation time and subsequently
enter the attraction basin of \(\phi_W(t)\).
In particular, at time \(t_W\) the fixed weak block
\[
A_2
:=
\left\{
i:
\angle(\widetilde u_i(t_W),
       \widetilde u(\phi_W(t_W)))
\le \rho_W
\right\}
\]
contains \(W_{\mathrm{cap}}\) and carries nonvanishing projected mass.

\item
\textbf{Weak-branch tracking.}
There exists \(\eta_2^u\ll1\) such that
\[
\angle(
\widetilde u_i(t),
\widetilde u(\phi_W(t))
)
\le
C\eta_2^u
\]
and
\[
\angle(
\widetilde w_i(t),
\widetilde u(\phi_W(t))
)
\le
C\eta_2^u
\]
for every \(i\in A_2\) and every \(t\in[t_W,T]\).

\item
\textbf{Concept-level gate geometry away from a boundary layer.}
Define
\[
g_1(x)
=
\mathbf 1\{x_1>0\},
\qquad
g_2(x,t)
=
\mathbf 1\{
\widetilde u(\phi_W(t))^\top x>0
\}.
\]
There exist boundary widths
\[
\tau_1
\asymp
X_\delta^\star\eta_1^u,
\qquad
\tau_2
\asymp
X_\delta^\star\eta_2^u.
\]
such that, for
\[
B_1
=
\{x:|x_1|\le\tau_1\}
\]
and
\[
B_2(t)
=
\{
x:
|\widetilde u(\phi_W(t))^\top x|
\le\tau_2
\},
\]
one has
\[
\mathbf 1\{u_i(t)^\top x>0\}
=
g_1(x)
\qquad
(i\in A_1,\ x\notin B_1)
\]
and
\[
\mathbf 1\{u_i(t)^\top x>0\}
=
g_2(x,t)
\qquad
(i\in A_2,\ x\notin B_2(t)).
\]
Furthermore the empirical mass of the boundary sets is quantitatively
small, uniformly over \(t\in[t_W,T]\).

\end{enumerate}

All sets \(A_1,A_2\) are fixed after their definition; in particular,
no time-dependent neuron membership is used in the subsequent
concept-block dynamics.
\end{conjecture}

\section{Proof of Theorem 2}

\begin{theorem}[Post-specialization concept-block reduction]
\label{thm:block-reduction}
Under the conclusions of
Theorem~\ref{thm:two-concept-specialization}, for every
\(t\in[t_W,T]\) and every \(x\in\mathcal S\),
\[
f_t(x)
=
g_1(x)M^{(1)}(t)x
+
g_2(x,t)M^{(2)}(t)x
+
R(x,t),
\]
where
\[
g_1(x)=\mathbf 1\{x_1>0\},
\qquad
g_2(x,t)=\mathbf 1\{u_W(t)^\top x>0\}.
\]

Moreover,
\[
R
=
R_{\mathrm{gate}}
+
R_{\mathrm{rem}}
+
R_{\perp},
\]
and outside the boundary layer
\[
\mathcal B_t(\rho)
=
\left\{
x:
\operatorname{dist}_{\angle}
\bigl(x,\partial g_1\cup\partial g_2(t)\bigr)
\le \rho
\right\},
\]
one has a uniform estimate
\[
\|R(x,t)\|
\le
C\varepsilon_{\mathrm{red}}\|x\|,
\]
with
\[
\varepsilon_{\mathrm{red}}
=
\varepsilon_{\mathrm{red}}
(\eta_1,\eta_2,\eta_{\mathrm{spread}},
R_{\mathrm{off}},h,n,\delta).
\]

The projected blocks satisfy
\[
\dot M^{(p)}
=
B_\parallel^{(p)}N^{(p)}
+
N^{(p)}C_\parallel^{(p)}
+
R_{\mathrm{gate}}^{(p)}
+
R_\perp^{(p)},
\]
where
\[
N^{(p)}
=
\mathbb E_n
\bigl[
(x-f_t(x))x^\top g_p(x,t)
\bigr].
\]
\end{theorem}

\section{Post-specialization predictive OOD normal form}

\begin{theorem}[Post-specialization predictive two-concept OOD normal form]
\label{thm:swingby}
Assume the conclusions of
Theorems~\ref{thm:two-concept-specialization}
and~\ref{thm:block-reduction}.
Fix \(p\in\{1,2\}\), \(q\neq p\), and a bounded near-boundary coordinate
\(|z|\le Z\).

For every direction
\[
v_{p,s}(z;\eta_p)
=
\eta_p z e_p
+
s\sqrt{1-\eta_p^2z^2}\,e_q
\]
lying outside the gate-smearing layer, the normalized OOD reconstruction
error
\[
E_{p,s}(z,t)
=
\frac12\|f_t(v_{p,s})-v_{p,s}\|^2
\]
admits the expansion
\[
E_{p,s}(z,t)
=
\frac12
+
\eta_p^2\mathcal H_{p,s}(z,t)
+
R_E(z,t),
\]
where
\[
\mathcal H_{p,s}(z,t)
=
\frac12
\bigl[(a_p-1)z+s b_p\bigr]^2
-
s c_p z
-
d_p
-
\frac12 z^2,
\]
and
\[
|R_E(z,t)|
\le
C\eta_p^3
+
C\varepsilon_{\mathrm{red}}.
\]

Moreover,
\[
\partial_t E_{p,s}(z,t)
=
\eta_p^2\mathcal G_{p,s}(z,t)
+
R_{\dot E}(z,t),
\]
where
\[
\mathcal G_{p,s}
=
\bigl[(a_p-1)z+s b_p\bigr]
\bigl[\dot a_p z+s\dot b_p\bigr]
-
s z\dot c_p
-
\dot d_p,
\]
and
\[
|R_{\dot E}|
\le
C\eta_p^3
+
C\varepsilon_{\mathrm{dyn}}.
\]

\end{theorem}

\begin{remark}[Scope of the predictive theorem]
\label{rem:T3-scope-A}
\textbf{Stage 0 qualification.} Section~\ref{sbaudit:scope} reopens the original reversal goal; the incoming scope choice below is provisional. No impossibility or necessity-of-a-second-theory conclusion is asserted.

Theorem~\ref{thm:swingby} is a post-specialization normal-form theorem.  It
provides the leading OOD error profile and its time derivative with explicit
remainders, but it does \emph{not} assert a general noisy many-neuron sign
reversal of $\mathcal G_{p,s}$.  The complete sign-reversal result proved in
this manuscript is the restricted noiseless fixed-gate theorem in
Section~\ref{sbsign:scope}.  A general noisy reversal-sign theorem remains a
separate open dynamical problem and is not load-bearing for the specialization
claim or for Theorem~\ref{thm:block-reduction}.
\end{remark}

\section{Aggregate propagation beyond the predictive normal form}

\label{sbagg:scope}
This addition proves deterministic aggregate estimates for the true empirical
gradient flow. It does not establish that their accuracy budget holds through
the learning and reversal interval. All notation in this section is local.
There are $n$ samples in each cluster, all in $\mathcal S$, and
$\mathbb E_n=(\mathbb E_{1,n}+\mathbb E_{2,n})/2$.
Set $\Sigma_n=\mathbb E_n[xx^\top]$, $\lambda_n=\|\Sigma_n\|_{\rm op}$,
and $\|g\|_n^2=\mathbb E_n\|g(x)\|^2$. Use the unscaled network
$f(x)=\sum_iw_i(u_i^\top x)_+$ and loss
$\mathcal L=\frac12\|x-f\|_n^2$.
Write $v_i=P_{\mathcal S}u_i$, $z_i=P_{\mathcal S}w_i$,
$e=x-f_\parallel$, and $\delta_i=v_i-z_i$.
Compatible selected gates $\alpha_i\in[0,1]$ give
\[
 A_i^\rho=\mathbb E_n[e x^\top\alpha_i],\qquad
 P_{i,\perp}^\rho=\mathbb E_n[x\alpha_i
                  \langle w_{i,\perp},f_\perp(x)\rangle],
 \qquad \dot z_i=A_i^\rho v_i,\quad
 \dot v_i=(A_i^\rho)^\top z_i-P_{i,\perp}^\rho.
\]
For a set $J$, let $M_J=\sum_J z_iv_i^\top$, $B_J=\sum_Jz_iz_i^\top$,
and $C_J=\sum_Jv_iv_i^\top$.
All differential inequalities below hold almost everywhere along the
absolutely continuous trajectory and imply their integrated comparisons.
At zero norms these follow by regularization. The notation $D^+$ is used
only in this almost-everywhere comparison sense, without claiming a
continuous empirical angular field or a derivative at every gate switch.


\begin{lemma}[Cross-residual energy controls weighted residual forcing]
\label{sbagg:lem:cross-residual-controls-forcing}

Fix an interval $I$ and write
\[
e_t(x)
:=
x-f_{\parallel,t}(x)
\in\mathcal S.
\]
For $p\in\{1,2\}$ write
\[
\mathbb E_{p,n}[\varphi]
:=
\frac1n\sum_{k=1}^n\varphi(x_k^{(p)}),
\qquad
s_{rs}^{(p)}
:=
\mathbb E_{p,n}[x_rx_s].
\]
Define the cross-coordinate residual energy
\[
\boxed{
\mathcal L_\times(t)
:=
\frac12
\left(
\mathbb E_{1,n}[e_{t,2}^2]
+
\mathbb E_{2,n}[e_{t,1}^2]
\right).
}
\]
Also define the empirical noise-coordinate scale
\[
\boxed{
\nu_\times^2
:=
\frac12
\left(
s_{22}^{(1)}
+
s_{11}^{(2)}
\right),
}
\]
and
\[
U_1:=\max_{k\in[n]}|x_{k,1}^{(1)}|,
\qquad
U_2:=\max_{k\in[n]}|x_{k,2}^{(2)}|,
\qquad
U_*:=\max\{U_1,U_2\}.
\]

Assume that on $I$
\[
\sup_{x\in\mathcal D}\|e_t(x)\|
\le R_{\infty,*}.
\]

For a fixed neuron set $J$, define the strong-axis weighted forcing
\[
\boxed{
\bigl(\mathcal F_J^{(1)}(t)\bigr)^2
:=
\sum_{i\in J}
\left(
v_{i1}^2|A_{i,21}^\rho|^2
+
z_{i1}^2|A_{i,12}^\rho|^2
\right).
}
\]
If
\[
(C_J)_{11}
=
\sum_{i\in J}v_{i1}^2
\le S_J,
\qquad
(B_J)_{11}
=
\sum_{i\in J}z_{i1}^2
\le S_J,
\]
then
\[
\boxed{
\bigl(\mathcal F_J^{(1)}(t)\bigr)^2
\le
S_J
\left[
U_*^2\mathcal L_\times(t)
+
R_{\infty,*}^2\nu_\times^2
\right].
}
\tag{cross-forcing-1}
\label{sbagg:eq:cross-forcing-1}
\]

Likewise define the weak-axis weighted forcing, equivalently the same
quantity in the ordered basis $(e_2,e_1)$, by
\[
\boxed{
\bigl(\mathcal F_J^{(2)}(t)\bigr)^2
:=
\sum_{i\in J}
\left(
v_{i2}^2|A_{i,12}^\rho|^2
+
z_{i2}^2|A_{i,21}^\rho|^2
\right).
}
\]
If
\[
(C_J)_{22}\le S_J,
\qquad
(B_J)_{22}\le S_J,
\]
then
\[
\boxed{
\bigl(\mathcal F_J^{(2)}(t)\bigr)^2
\le
S_J
\left[
U_*^2\mathcal L_\times(t)
+
R_{\infty,*}^2\nu_\times^2
\right].
}
\tag{cross-forcing-2}
\label{sbagg:eq:cross-forcing-2}
\]

In particular, in either orientation,
\[
\boxed{
\mathcal F_J^{(p)}(t)
\le
\sqrt{S_J}
\left(
U_*\sqrt{\mathcal L_\times(t)}
+
R_{\infty,*}\nu_\times
\right).
}
\tag{cross-forcing-root}
\label{sbagg:eq:cross-forcing-root}
\]
\end{lemma}

\begin{proof}
By definition of the residual-gated matrix,
\[
A_{i,21}^\rho
=
\mathbb E_n[e_{t,2}(x)x_1\alpha_i(t;x)],
\qquad
A_{i,12}^\rho
=
\mathbb E_n[e_{t,1}(x)x_2\alpha_i(t;x)].
\]
Set
\[
Y_{21}(x)
:=
\left(
v_{i1}e_{t,2}(x)x_1\alpha_i(t;x)
\right)_{i\in J}.
\]
Then Jensen's inequality in $\ell_2(J)$ gives
\[
\begin{aligned}
\sum_{i\in J}
v_{i1}^2|A_{i,21}^\rho|^2
&=
\left\|
\mathbb E_nY_{21}
\right\|_{\ell_2(J)}^2
\\
&\le
\mathbb E_n
\left[
e_{t,2}^2x_1^2
\sum_{i\in J}
v_{i1}^2\alpha_i(t;x)^2
\right]
\\
&\le
(C_J)_{11}
\mathbb E_n[e_{t,2}^2x_1^2],
\end{aligned}
\]
because $0\le\alpha_i\le1$.  Similarly,
\[
\sum_{i\in J}
z_{i1}^2|A_{i,12}^\rho|^2
\le
(B_J)_{11}
\mathbb E_n[e_{t,1}^2x_2^2].
\]
Hence
\[
\bigl(\mathcal F_J^{(1)}\bigr)^2
\le
S_J
\left\{
\mathbb E_n[e_{t,2}^2x_1^2]
+
\mathbb E_n[e_{t,1}^2x_2^2]
\right\}.
\]

Split the first expectation over the two clusters:
\[
\begin{aligned}
\mathbb E_n[e_{t,2}^2x_1^2]
&=
\frac12
\mathbb E_{1,n}[e_{t,2}^2x_1^2]
+
\frac12
\mathbb E_{2,n}[e_{t,2}^2x_1^2]
\\
&\le
\frac12
U_1^2\mathbb E_{1,n}[e_{t,2}^2]
+
\frac12
R_{\infty,*}^2s_{11}^{(2)}.
\end{aligned}
\]
Likewise,
\[
\mathbb E_n[e_{t,1}^2x_2^2]
\le
\frac12
R_{\infty,*}^2s_{22}^{(1)}
+
\frac12
U_2^2\mathbb E_{2,n}[e_{t,1}^2].
\]
Adding the two inequalities gives
\[
\mathbb E_n[e_{t,2}^2x_1^2]
+
\mathbb E_n[e_{t,1}^2x_2^2]
\le
U_*^2\mathcal L_\times
+
R_{\infty,*}^2\nu_\times^2,
\]
which proves \eqref{sbagg:eq:cross-forcing-1}.

The second orientation is identical after interchanging
$e_1\leftrightarrow e_2$.  Finally,
$\sqrt{a+b}\le\sqrt a+\sqrt b$ gives
\eqref{sbagg:eq:cross-forcing-root}.
\end{proof}

\begin{lemma}[Coupled propagation of cross residual and aggregate axis grading]
\label{sbagg:lem:coupled-cross-transverse}

Fix a finite interval
\[
I=[t_0,t_0+T]
\]
and a fixed partition
\[
[h]=J_1\sqcup J_2.
\]
The membership of $J_1,J_2$ is fixed after $t_0$.

Write
\[
e_t(x):=x-f_{\parallel,t}(x),
\]
and define
\[
X(t)
:=
\sqrt{\mathcal L_\times(t)}
=
\left\{
\frac12
\left(
\mathbb E_{1,n}[e_{t,2}^2]
+
\mathbb E_{2,n}[e_{t,1}^2]
\right)
\right\}^{1/2}.
\]

Define the two aggregate transverse energies
\[
\boxed{
Y_1(t)^2
:=
\sum_{i\in J_1}
\left(
v_{i2}(t)^2+z_{i2}(t)^2
\right),
}
\]
and
\[
\boxed{
Y_2(t)^2
:=
\sum_{i\in J_2}
\left(
v_{i1}(t)^2+z_{i1}(t)^2
\right).
}
\]
Thus $J_1$ is measured relative to the $e_1$ axis and $J_2$ relative to
the $e_2$ axis; no pointwise angular condition is assumed.

Let
\[
\Sigma_p
:=
\mathbb E_{p,n}[xx^\top],
\qquad
\lambda_p:=\|\Sigma_p\|_{\rm op},
\]
and write
\[
s_{rr}^{(p)}:=e_r^\top\Sigma_pe_r.
\]
Define the empirical cross-cluster data constants
\[
\boxed{
\tau_w
:=
\max\left\{
\sqrt{s_{11}^{(1)}s_{11}^{(2)}},
\sqrt{s_{22}^{(1)}s_{22}^{(2)}}
\right\},
}
\]
\[
\boxed{
\Lambda_w
:=
\sqrt{s_{11}^{(1)}s_{22}^{(2)}}
+
\sqrt{s_{22}^{(1)}s_{11}^{(2)}},
}
\]
and
\[
\boxed{
\tau_u
:=
\sqrt{\operatorname{tr}(\Sigma_1\Sigma_2)},
\qquad
\Lambda_u:=\lambda_1+\lambda_2.
}
\]

Assume that on $I$ the already available global bounds give
\[
\sum_{i=1}^h\|v_i(t)\|^2\le S_*,
\qquad
\sum_{i=1}^h\|z_i(t)\|^2\le S_*,
\tag{mass-ceiling}
\label{sbagg:eq:coupled-mass-ceiling}
\]
and
\[
\|e_t\|_n\le R_{n,*},
\qquad
\sup_{x\in\mathcal D}\|e_t(x)\|\le R_{\infty,*}.
\tag{residual-ceiling}
\label{sbagg:eq:coupled-residual-ceiling}
\]
Let
\[
C_x:=\mathbb E_n\|x\|^2,
\qquad
K_*:=\sqrt{\|\Sigma_n\|_{\rm op}}\,R_{n,*}.
\]

Define
\[
\Pi_p(t)
:=
\left(
\sum_{i\in J_p}
\|P_{i,\perp}^\rho(t)\|^2
\right)^{1/2},
\qquad
\Pi(t)
:=
\left(
\sum_{i=1}^h
\|P_{i,\perp}^\rho(t)\|^2
\right)^{1/2}.
\]

Finally put
\[
\nu_\times^2
:=
\frac12
\left(
s_{22}^{(1)}+s_{11}^{(2)}
\right),
\]
and
\[
U_*:=
\max\left\{
\max_k|x_{k,1}^{(1)}|,
\max_k|x_{k,2}^{(2)}|
\right\}.
\]

Then, for a.e. $t\in I$, the following coupled inequalities hold.

First,
\[
\boxed{
D^+Y_p
\le
K_*Y_p
+
\sqrt{S_*}\,U_*X
+
\sqrt{S_*}\,R_{\infty,*}\nu_\times
+
\Pi_p,
\qquad
p=1,2.
}
\tag{coupled-Y}
\label{sbagg:eq:coupled-Y}
\]

Second,
\[
\boxed{
D^+X
\le
a_Y(Y_1+Y_2)
+
a_{\rm data}
+
\sqrt{S_*C_x}\,\Pi,
}
\tag{coupled-X}
\label{sbagg:eq:coupled-X}
\]
where
\[
\boxed{
a_Y
:=
\frac12
R_{n,*}\sqrt{S_*}
(\Lambda_w+\Lambda_u),
}
\]
and
\[
\boxed{
a_{\rm data}
:=
\frac12
R_{n,*}S_*
(\tau_w+\tau_u).
}
\]

Consequently, if
\[
Z(t):=X(t)+Y_1(t)+Y_2(t),
\]
then
\[
\boxed{
D^+Z
\le
L_{\rm cpl}Z
+
F_{\rm cpl}(t),
}
\tag{coupled-Z}
\label{sbagg:eq:coupled-Z}
\]
with
\[
\boxed{
L_{\rm cpl}
:=
\max\left\{
2\sqrt{S_*}U_*,
K_*+a_Y
\right\},
}
\]
and
\[
\boxed{
F_{\rm cpl}(t)
:=
a_{\rm data}
+
2\sqrt{S_*}R_{\infty,*}\nu_\times
+
\sqrt{S_*C_x}\Pi(t)
+
\Pi_1(t)+\Pi_2(t).
}
\]
Therefore
\[
\boxed{
Z(t)
\le
e^{L_{\rm cpl}(t-t_0)}Z(t_0)
+
\int_{t_0}^t
e^{L_{\rm cpl}(t-s)}
F_{\rm cpl}(s)\,ds.
}
\tag{coupled-duhamel}
\label{sbagg:eq:coupled-duhamel}
\]

In particular, if
\[
\Pi(t)\le\Pi_*
\qquad(t\in I),
\]
then
\[
\Pi_1+\Pi_2\le\sqrt2\,\Pi_*
\]
and hence
\[
F_{\rm cpl}(t)\le F_*,
\]
where
\[
\boxed{
F_*
:=
a_{\rm data}
+
2\sqrt{S_*}R_{\infty,*}\nu_\times
+
\left(
\sqrt{S_*C_x}+\sqrt2
\right)\Pi_*.
}
\]
Thus
\[
\boxed{
Z(t)
\le
e^{L_{\rm cpl}T}Z(t_0)
+
F_*
\frac{e^{L_{\rm cpl}T}-1}{L_{\rm cpl}},
}
\tag{coupled-uniform}
\label{sbagg:eq:coupled-uniform}
\]
with the quotient interpreted as $T$ when $L_{\rm cpl}=0$.

Hence a checkpoint bound
\[
\boxed{
Z(t_0)\le z_0\eta
}
\tag{coupled-checkpoint}
\label{sbagg:eq:coupled-checkpoint}
\]
together with a primitive-data/perpendicular budget
\[
\boxed{
F_*\le f_0\eta
}
\tag{coupled-forcing-budget}
\label{sbagg:eq:coupled-forcing-budget}
\]
implies
\[
\boxed{
X(t)+Y_1(t)+Y_2(t)
\le
C_{\rm cpl}\eta
\qquad
(t\in I),
}
\]
where
\[
C_{\rm cpl}
=
e^{L_{\rm cpl}T}z_0
+
f_0
\frac{e^{L_{\rm cpl}T}-1}{L_{\rm cpl}}.
\]

In particular,
\[
\mathcal L_\times(t)
\le
C_{\rm cpl}^2\eta^2,
\]
and Lemma~\ref{sbagg:lem:cross-residual-controls-forcing} gives, for $p=1,2$,
\[
\boxed{
\mathcal F_{J_p}^{(p)}(t)
\le
\eta\sqrt{S_*}
\left(
U_*C_{\rm cpl}
+
R_{\infty,*}\frac{\nu_\times}{\eta}
\right).
}
\]
Thus whenever the budget
\eqref{sbagg:eq:coupled-forcing-budget} is imposed at the intended $\eta$ scale,
the weighted residual forcing required by the aggregate grading lemma is
propagated on the entire interval from the single checkpoint
\eqref{sbagg:eq:coupled-checkpoint}.
\end{lemma}

\begin{proof}\leavevmode

\paragraph{Step 1: transverse-energy inequalities.}

For $J_1$, apply the exact transverse-energy identity in the ordered basis
$(e_1,e_2)$.  The residual-gated projected equations give
\[
D^+Y_1
\le
K_*Y_1+\mathcal F_{J_1}^{(1)}+\Pi_1.
\]
For $J_2$, use the ordered basis $(e_2,e_1)$:
\[
D^+Y_2
\le
K_*Y_2+\mathcal F_{J_2}^{(2)}+\Pi_2.
\]
The operator bound
\[
\|A_i^\rho\|_{\rm op}
\le
\sqrt{\|\Sigma_n\|_{\rm op}}\,
\|e_t\|_n
\le K_*
\]
follows directly from Cauchy--Schwarz.

Lemma~\ref{sbagg:lem:cross-residual-controls-forcing}, together with the global
mass ceiling, gives
\[
\mathcal F_{J_p}^{(p)}
\le
\sqrt{S_*}
\left(
U_*X+R_{\infty,*}\nu_\times
\right),
\]
which proves \eqref{sbagg:eq:coupled-Y}.

\paragraph{Step 2: exact planar residual evolution.}

For
\[
q_i(x):=(v_i^\top x)_+,
\]
the residual-gated Cartesian equations imply
\[
\dot z_i
=
\mathbb E_n[e_t(y)q_i(y)]
\]
and
\[
\dot v_i
=
\mathbb E_n[
y\,\alpha_i(y)\,e_t(y)^\top z_i]
-
P_{i,\perp}^\rho.
\]
Therefore
\[
\dot e_t(x)
=
-\mathcal K_t e_t(x)
+
G_{\perp,t}(x),
\]
where
\[
(\mathcal K_tg)(x)
=
\mathbb E_n[K_t(x,y)g(y)]
\]
with
\[
\boxed{
K_t(x,y)
=
\sum_{i=1}^h
\left[
q_i(x)q_i(y)I_{\mathcal S}
+
\alpha_i(x)\alpha_i(y)
(x^\top y)z_i z_i^\top
\right],
}
\tag{selected-ntk}
\label{sbagg:eq:selected-ntk}
\]
and
\[
\boxed{
G_{\perp,t}(x)
=
\sum_{i=1}^h
z_i\alpha_i(x)
x^\top P_{i,\perp}^\rho.
}
\tag{perp-residual-forcing}
\label{sbagg:eq:perp-residual-forcing}
\]
The operator $\mathcal K_t$ is self-adjoint and positive semidefinite for
every selected gate configuration, including selected boundary values.

Let $P_\times$ denote the sample-dependent orthogonal projection
\[
(P_\times g)(x_k^{(1)})
=
g_2(x_k^{(1)})e_2,
\qquad
(P_\times g)(x_k^{(2)})
=
g_1(x_k^{(2)})e_1,
\]
and set $P_\parallel=I-P_\times$.  Then
\[
X=\|P_\times e_t\|_n.
\]
Hence
\[
\begin{aligned}
\frac12\frac d{dt}X^2
={}&
-\langle P_\times e_t,
\mathcal K_tP_\times e_t\rangle_n
\\
&-
\langle P_\times e_t,
\mathcal K_tP_\parallel e_t\rangle_n
+
\langle P_\times e_t,G_{\perp,t}\rangle_n .
\end{aligned}
\]
The first term is nonpositive.  Therefore, with
\[
\kappa_\times(t)
:=
\|P_\times\mathcal K_tP_\parallel\|_{\rm op},
\]
regularization at $X=0$ gives
\[
D^+X
\le
\kappa_\times\|P_\parallel e_t\|_n
+
\|G_{\perp,t}\|_n
\le
R_{n,*}\kappa_\times
+
\|G_{\perp,t}\|_n.
\]

By Cauchy--Schwarz over neurons,
\[
\|G_{\perp,t}(x)\|
\le
\left(\sum_i\|z_i\|^2\right)^{1/2}
\|x\|
\left(\sum_i\|P_{i,\perp}^\rho\|^2\right)^{1/2},
\]
and consequently
\[
\boxed{
\|G_{\perp,t}\|_n
\le
\sqrt{S_*C_x}\,\Pi(t).
}
\tag{Gperp-bound}
\label{sbagg:eq:Gperp-bound}
\]

\paragraph{Step 3: output-weight part of the cross-NTK.}

Write
\[
\mathcal K_t=\mathcal K_t^w+\mathcal K_t^u
\]
for the two terms in \eqref{sbagg:eq:selected-ntk}.  Since
$\mathcal K_t^w$ is scalar in output coordinates, the map
$P_\times\mathcal K_t^wP_\parallel$ only couples the two clusters.
Therefore
\[
\|P_\times\mathcal K_t^wP_\parallel\|_{\rm op}
\le
\frac12
\sum_i
\|q_i\|_{1,n}\|q_i\|_{2,n}.
\]
Now
\[
\|q_i\|_{1,n}
\le
|v_{i1}|\sqrt{s_{11}^{(1)}}
+
|v_{i2}|\sqrt{s_{22}^{(1)}},
\]
and
\[
\|q_i\|_{2,n}
\le
|v_{i1}|\sqrt{s_{11}^{(2)}}
+
|v_{i2}|\sqrt{s_{22}^{(2)}}.
\]
Thus
\[
\boxed{
\|P_\times\mathcal K_t^wP_\parallel\|_{\rm op}
\le
\frac12
\left[
\tau_w\sum_i\|v_i\|^2
+
\Lambda_w
\sum_i|v_{i1}v_{i2}|
\right].
}
\tag{Kw-cross}
\label{sbagg:eq:Kw-cross}
\]

\paragraph{Step 4: input-weight part of the cross-NTK.}

For each neuron and cluster define
\[
T_{i,p}g
:=
\mathbb E_{p,n}[
\alpha_i(x)xg(x)
],
\qquad
g\in L^2(\widehat P_{p,n}).
\]
Then
\[
T_{i,p}T_{i,p}^*
=
\mathbb E_{p,n}[
\alpha_i(x)^2xx^\top
]
\preceq
\Sigma_p,
\]
so
\[
\|T_{i,p}\|_{\rm op}^2\le\lambda_p.
\]
Moreover,
\[
\begin{aligned}
\|T_{i,1}^*T_{i,2}\|_{\rm op}^2
&\le
\|T_{i,1}^*T_{i,2}\|_{\rm HS}^2
\\
&=
\operatorname{tr}
\left[
(T_{i,1}T_{i,1}^*)
(T_{i,2}T_{i,2}^*)
\right]
\\
&\le
\operatorname{tr}(\Sigma_1\Sigma_2)
=
\tau_u^2.
\end{aligned}
\]
The four correct-to-wrong blocks of the $i$th input-weight kernel carry,
respectively, coefficients
\[
z_{i1}z_{i2},
\qquad
z_{i2}^2,
\qquad
z_{i1}^2,
\qquad
z_{i1}z_{i2}.
\]
Hence
\[
\boxed{
\|P_\times\mathcal K_t^uP_\parallel\|_{\rm op}
\le
\frac12
\left[
\tau_u\sum_i\|z_i\|^2
+
(\lambda_1+\lambda_2)
\sum_i|z_{i1}z_{i2}|
\right].
}
\tag{Ku-cross}
\label{sbagg:eq:Ku-cross}
\]

\paragraph{Step 5: replace absolute mixed masses by aggregate transverse
energies.}

Since $[h]=J_1\sqcup J_2$,
\[
\begin{aligned}
\sum_i|v_{i1}v_{i2}|
&\le
\sqrt{
\sum_{i\in J_1}v_{i1}^2
}
\sqrt{
\sum_{i\in J_1}v_{i2}^2
}
\\
&\quad+
\sqrt{
\sum_{i\in J_2}v_{i2}^2
}
\sqrt{
\sum_{i\in J_2}v_{i1}^2
}
\\
&\le
\sqrt{S_*}(Y_1+Y_2).
\end{aligned}
\]
The identical argument gives
\[
\sum_i|z_{i1}z_{i2}|
\le
\sqrt{S_*}(Y_1+Y_2).
\]
Combining \eqref{sbagg:eq:Kw-cross},
\eqref{sbagg:eq:Ku-cross}, and the mass ceiling yields
\[
\boxed{
\kappa_\times
\le
\frac12S_*(\tau_w+\tau_u)
+
\frac12\sqrt{S_*}
(\Lambda_w+\Lambda_u)
(Y_1+Y_2).
}
\tag{cross-ntk-bound}
\label{sbagg:eq:cross-ntk-bound}
\]
Substitution into the preceding estimate for $D^+X$ proves
\eqref{sbagg:eq:coupled-X}.

\paragraph{Step 6: scalar closure.}

Add \eqref{sbagg:eq:coupled-X} and the two inequalities
\eqref{sbagg:eq:coupled-Y}.  The coefficient of $X$ is
$2\sqrt{S_*}U_*$, while the coefficient of $Y_1+Y_2$ is
$K_*+a_Y$.  Thus
\[
D^+Z
\le
L_{\rm cpl}Z+F_{\rm cpl},
\]
which is \eqref{sbagg:eq:coupled-Z}.  Scalar Duhamel comparison gives
\eqref{sbagg:eq:coupled-duhamel}.  Since $J_1,J_2$ are disjoint,
\[
\Pi_1+\Pi_2
\le
\sqrt2
\left(
\Pi_1^2+\Pi_2^2
\right)^{1/2}
=
\sqrt2\,\Pi,
\]
giving \eqref{sbagg:eq:coupled-uniform}.  The final claims follow by
\eqref{sbagg:eq:coupled-checkpoint}--\eqref{sbagg:eq:coupled-forcing-budget}.
\end{proof}


\subsection{A sharper comparison and the required accuracy condition}
\begin{remark}[Available global ceilings do not require extra trajectory hypotheses]
\label{sbagg:global-ceilings}
Under the source's balanced initialization and perpendicular contraction,
the ceilings in the coupled lemma have explicit choices. Put
$X_{\max}=\max_{\mathcal D}\|x\|$ and $s_0=h\varepsilon^2$.
On $[t_0,t_0+T]$ one may use
\[
 S_*=s_0+\tfrac12C_x(t_0+T),\qquad
 R_{n,*}=\sqrt{C_x}(1+s_0),\qquad
 R_{\infty,*}=(1+S_*)X_{\max},\qquad
 \Pi_*=\lambda_n h\varepsilon^2\sqrt{S_*}.
\]
Indeed balance and the exact mass identity give
$\dot S=C_x/2-2\|f-x/2\|_n^2\le C_x/2$.
The pointwise estimate $\|f(x)\|\le S\|x\|$ and full loss dissipation
give the residual ceilings. Finally
$\|W_\perp\|_F\le\varepsilon\sqrt h$ implies
$\|f_\perp\|_n\le\varepsilon\sqrt{h\lambda_n S_*}$, and
Cauchy--Schwarz gives
$\Pi\le\sqrt{\lambda_n}\|f_\perp\|_n\|W_\perp\|_F$.
These choices establish finite ceilings; their quantitative usefulness
in the accuracy condition below is a separate question.
\end{remark}

\begin{remark}[The smaller input-kernel coefficient is also valid]
\label{sbagg:kernel-sharpening}
The preceding proof safely uses $\Lambda_u=\lambda_1+\lambda_2$.
In fact all its conclusions remain true with
$\Lambda_u=\max\{\lambda_1,\lambda_2\}$.
To see this, the correct-to-wrong input-kernel block of one neuron,
with the common cluster normalization understood, is
\[
 \frac12\begin{pmatrix}
 z_{i1}z_{i2}T_{i,1}^*T_{i,1}&z_{i2}^2T_{i,1}^*T_{i,2}\\
 z_{i1}^2T_{i,2}^*T_{i,1}&z_{i1}z_{i2}T_{i,2}^*T_{i,2}
 \end{pmatrix}.
\]
The block-diagonal part has norm at most
$|z_{i1}z_{i2}|\max_p\lambda_p/2$.
The off-diagonal part has norm at most
$\tau_u\max\{z_{i1}^2,z_{i2}^2\}/2
\le\tau_u\|z_i\|^2/2$.
Add these bounds and sum over neurons. Thus the proposed replacement
of the maximum by the sum was sufficient but not necessary.
\end{remark}

\begin{corollary}[Two-dimensional comparison and finite accuracy budget]
\label{sbagg:budget}
Use the constants of Lemma~\ref{sbagg:lem:coupled-cross-transverse}, and set
\[
 Y=Y_1+Y_2,\quad c=\sqrt{S_*}U_*,\quad
 b=\sqrt{S_*}R_{\infty,*}\nu_\times,\quad
 H=\begin{pmatrix}0&a_Y\\2c&K_*\end{pmatrix}.
\]
Then, componentwise,
\[
 \binom{X(t)}{Y(t)}\le
 e^{H(t-t_0)}\binom{X(t_0)}{Y(t_0)}+
 \int_{t_0}^t e^{H(t-s)}
 \binom{a_{\rm data}+\sqrt{S_*C_x}\Pi(s)}
 {2b+\Pi_1(s)+\Pi_2(s)}\,ds.
\]
In particular a sufficient, fully quantitative condition for
$X+Y\le C_{\rm tar}\eta$ on $[t_0,t_0+T]$ is
\begin{equation}\label{sbagg:accuracy-budget}
 e^{L_{\rm cpl}T}Z(t_0)+F_*\mathcal E_{L_{\rm cpl}}(T)
 \le C_{\rm tar}\eta,
 \qquad
 \mathcal E_a(T)=\begin{cases}(e^{aT}-1)/a,&a>0,\\T,&a=0.\end{cases}
\end{equation}
The matrix comparison can replace the scalar left side by a smaller bound.
If $\eta$ varies, an $O(\eta)$ claim requires a constant $C_{\rm tar}$
uniform in that variation. Finite constants for each fixed interval do
not establish such uniformity when $T$ depends on $\eta$ or $\varepsilon$.
\end{corollary}
\begin{proof}
Add the two $Y_p$ inequalities and retain the separate $X$ inequality.
The off-diagonal entries of $H$ are nonnegative, so its exponential
preserves componentwise order; variation of constants proves the first
claim. The scalar comparison proves \eqref{sbagg:accuracy-budget}.
The largest eigenvalue of $H$ is
$(K_*+\sqrt{K_*^2+8a_Yc})/2$; this improvement still permits growth.
\end{proof}

\begin{corollary}[Consequent grading and coefficient rates]
\label{sbagg:rates}
Fix $p$ and use the ordered basis $(e_p,e_{3-p})$. Suppose on the same
interval $Y_p\le\eta C$ and $\mathcal F_{J_p}^{(p)}+\Pi_p\le\eta f$,
with $\eta>0$ fixed in time. Then write
\[
 M_{J_p}=\begin{pmatrix}a&\eta b_0\\\eta c_0&\eta^2d\end{pmatrix}.
\]
The bounds below hold uniformly, with derivatives almost everywhere:
\[
 |a|\le S_*,\quad |b_0|,|c_0|\le\sqrt{S_*}C,\quad |d|\le C^2/2,
\]
\[
 |\dot a|\le2K_*S_*+\sqrt{S_*}\eta f,\qquad
 |\dot d|\le K_*C^2+Cf,
\]
\[
 |\dot b_0|,|\dot c_0|
 \le K_*\sqrt{2S_*}C+K_*\eta C^2+\sqrt{S_*}f.
\]
Also $(B_{J_p})_{22},(C_{J_p})_{22}\le\eta^2C^2$ in this ordered basis.
The forcing constant can be chosen as
\[
 f=\sqrt{S_*}\left(U_*C_{\rm tar}
              +R_{\infty,*}\frac{\nu_\times}{\eta}\right)
       +\frac{\Pi_*}{\eta}
\]
whenever \eqref{sbagg:accuracy-budget} holds.
\end{corollary}
\begin{proof}
Cauchy--Schwarz gives the static bounds. The exact identity
\[
 \dot M_J=\sum_{i\in J}
 \left(A_i^\rho v_iv_i^\top+z_iz_i^\top A_i^\rho
                         -z_i(P_{i,\perp}^\rho)^\top\right)
\]
gives
$|\dot M_{22}|\le K_*Y_p^2+Y_p(\mathcal F_{J_p}^{(p)}+\Pi_p)$.
For each off-diagonal entry the terms are bounded by
$K_*\sqrt{2S_*}Y_p+K_*Y_p^2+sqrt{S_*}(\mathcal F_{J_p}^{(p)}+\Pi_p)$.
Finally $|\dot M_{11}|\le2K_*S_*+\sqrt{S_*}\Pi_p$.
Divide by the indicated fixed powers of $\eta$ and apply the Jensen lemma.
These are bounds for the partition matrices, not automatically for the
matrix of an arbitrary probe activation cell.
\end{proof}

\subsection{Checkpoint construction and its precise limitation}


\begin{lemma}[Optimal two-axis partition and aggregate checkpoint]
\label{sbagg:lem:optimal-axis-checkpoint}

Fix a time $t_\sharp$ and write
\[
v_i=P_{\mathcal S}u_i,
\qquad
z_i=P_{\mathcal S}w_i.
\]
For each neuron define
\[
d_{i\to1}
:=
v_{i2}^2+z_{i2}^2,
\qquad
d_{i\to2}
:=
v_{i1}^2+z_{i1}^2,
\]
and define the aggregate axis defect
\[
\boxed{
\mathfrak D_{\rm ax}(t_\sharp)
:=
\sum_{i=1}^h
\min\{d_{i\to1},d_{i\to2}\}.
}
\]

For a partition
\[
[h]=J_1\sqcup J_2
\]
set
\[
\mathcal T_{J_1}^{(1)}
:=
\sum_{i\in J_1}
(v_{i2}^2+z_{i2}^2),
\]
and
\[
\mathcal T_{J_2}^{(2)}
:=
\sum_{i\in J_2}
(v_{i1}^2+z_{i1}^2).
\]
Then
\[
\boxed{
\mathfrak D_{\rm ax}(t_\sharp)
=
\min_{J_1\sqcup J_2=[h]}
\left[
\mathcal T_{J_1}^{(1)}
+
\mathcal T_{J_2}^{(2)}
\right].
}
\tag{axis-defect-min}
\label{sbagg:eq:axis-defect-min}
\]

One minimizing partition is
\[
\boxed{
J_1^\sharp
=
\{i:d_{i\to1}\le d_{i\to2}\},
\qquad
J_2^\sharp=[h]\setminus J_1^\sharp,
}
\tag{optimal-axis-partition}
\label{sbagg:eq:optimal-axis-partition}
\]
with arbitrary tie breaking.

Consequently,
\[
\boxed{
\sqrt{\mathcal T_{J_1^\sharp}^{(1)}}
+
\sqrt{\mathcal T_{J_2^\sharp}^{(2)}}
\le
\sqrt{2\mathfrak D_{\rm ax}(t_\sharp)}.
}
\tag{axis-checkpoint-root}
\label{sbagg:eq:axis-checkpoint-root}
\]

Let also
\[
X_\sharp
:=
\sqrt{\mathcal L_\times(t_\sharp)}.
\]
If, for some fixed scale $\eta>0$,
\[
\boxed{
\mathcal L_\times(t_\sharp)
\le c_\times^2\eta^2,
\qquad
\mathfrak D_{\rm ax}(t_\sharp)
\le c_{\rm ax}^2\eta^2,
}
\tag{aggregate-checkpoint}
\label{sbagg:eq:aggregate-checkpoint}
\]
then the fixed partition
\[
J_1:=J_1^\sharp,
\qquad
J_2:=J_2^\sharp
\]
satisfies
\[
\boxed{
X_\sharp
+
\sqrt{\mathcal T_{J_1}^{(1)}(t_\sharp)}
+
\sqrt{\mathcal T_{J_2}^{(2)}(t_\sharp)}
\le
\left(c_\times+\sqrt2\,c_{\rm ax}\right)\eta.
}
\tag{coupled-checkpoint-ready}
\label{sbagg:eq:coupled-checkpoint-ready}
\]

Thus the checkpoint hypothesis required by
Lemma~\ref{sbagg:lem:coupled-cross-transverse}
is equivalent, up to the universal factor $\sqrt2$, to the two scalar
aggregate bounds in \eqref{sbagg:eq:aggregate-checkpoint}.
\end{lemma}

\begin{proof}
For any fixed partition $J_1\sqcup J_2=[h]$,
\[
\mathcal T_{J_1}^{(1)}
+
\mathcal T_{J_2}^{(2)}
=
\sum_{i\in J_1}d_{i\to1}
+
\sum_{i\in J_2}d_{i\to2}.
\]
The choice of group for one neuron does not affect the cost of any other
neuron.  Hence the minimum is obtained independently for every $i$ by
choosing the smaller of $d_{i\to1}$ and $d_{i\to2}$.  This proves
\eqref{sbagg:eq:axis-defect-min} and
\eqref{sbagg:eq:optimal-axis-partition}.

For the minimizing partition,
\[
\mathcal T_{J_1^\sharp}^{(1)}
+
\mathcal T_{J_2^\sharp}^{(2)}
=
\mathfrak D_{\rm ax}.
\]
Therefore
\[
\begin{aligned}
\sqrt{\mathcal T_{J_1^\sharp}^{(1)}}
+
\sqrt{\mathcal T_{J_2^\sharp}^{(2)}}
&\le
\sqrt{
2\left(
\mathcal T_{J_1^\sharp}^{(1)}
+
\mathcal T_{J_2^\sharp}^{(2)}
\right)
}
\\
&=
\sqrt{2\mathfrak D_{\rm ax}},
\end{aligned}
\]
proving \eqref{sbagg:eq:axis-checkpoint-root}.  The final statement follows
immediately from \eqref{sbagg:eq:aggregate-checkpoint}.
\end{proof}


\begin{corollary}[Cohort and primitive checkpoint certificates]
\label{sbagg:primitive}
At a fixed checkpoint let $[h]=B_1\sqcup B_2\sqcup R$ and define
\[
 Q_1=\sum_{B_1}v_{i2}^2,\quad Q_2=\sum_{B_2}v_{i1}^2,\quad
 \Delta_p=\sum_{B_p}\|v_i-z_i\|^2,\quad
 S_R=\sum_R(\|v_i\|^2+\|z_i\|^2).
\]
Then
\[
 \mathfrak D_{\rm ax}\le
 2Q_1+\Delta_1+2\sqrt{Q_1\Delta_1}
 +2Q_2+\Delta_2+2\sqrt{Q_2\Delta_2}+S_R/2.
\]
Independently, balance gives
$\mathfrak D_{\rm ax}(t)\le S(t)=\sum_i\|u_i(t)\|^2
=\sum_i\|w_i(t)\|^2$.
Suppose valid early bounds on $[0,H_\star]$ are
\[
 S(t)\le \frac{h\varepsilon^2e^{2\lambda H_\star}}
                         {D_\star^{\rm Ric}}=:S_\sharp^*,
 \qquad \sup_{x\in\mathcal D}\|f_t(x)\|\le F_\star,
 \qquad D_\star^{\rm Ric}>0.
\]
For any $t_\sharp\le H_\star$, the additional conditions
\[
 S_\sharp^*\le c_{\rm ax}^2\eta^2,\qquad
 \nu_\times+F_\star\le c_\times\eta
\]
imply, for the minimizing partition frozen at $t_\sharp$,
\[
 Z(t_\sharp)\le(c_\times+\sqrt2c_{\rm ax})\eta.
\]
This certifies a checkpoint only. It neither proves
\eqref{sbagg:accuracy-budget} on a later learning interval nor supplies
probe gate orientation, longitudinal mass, or a reversal sign.
\end{corollary}
\begin{proof}
On $B_p$ the output transverse norm is at most
$\sqrt{Q_p}+\sqrt{\Delta_p}$; square and add the input transverse energy.
On $R$, the smaller cost is at most half the total squared norm.
Applying this last observation to every neuron gives
$\mathfrak D_{\rm ax}\le S$.
The triangle inequality gives
$X=\|P_\times(x-f_\parallel)\|_n\le\nu_\times+F_\star$.
Apply Lemma~\ref{sbagg:lem:optimal-axis-checkpoint}.
\end{proof}

\begin{remark}[Why small initial mass does not close late propagation]
\label{sbagg:small-mass-example}
Consider the deterministic planar data $x^{(1)}=e_1,x^{(2)}=e_2$ with
one balanced neuron $u=w=\sqrt{S/2}(1,1)^\top$. This is an algebraic
example within the balanced equations, not a counterexample to a
high-probability assertion under isotropic random initialization.
Both training gates are on, and the exact gradient flow gives
\[
 \dot S=S(1-S),\qquad
 S(t)=\frac{S(0)e^t}{1-S(0)+S(0)e^t}.
\]
Here $\nu_\times=\tau_w=\tau_u=\Pi=0$, so $F_*=0$, while
\[
 \mathfrak D_{\rm ax}=S,\qquad Y_1+Y_2=\sqrt S,\qquad X=S/2.
\]
Taking $S(0)=\eta^2<1/2$ gives the required small checkpoint
$Z(0)=\eta+\eta^2/2$. At
$t_\eta=\log((1-\eta^2)/\eta^2)$, however,
\[
 S(t_\eta)=1/2,\qquad Z(t_\eta)=1/\sqrt2+1/4.
\]
Thus a checkpoint and an $O(\eta)$ forcing budget alone cannot give
uniform $O(\eta)$ grading on a growing learning horizon.
The proved comparison correctly allows its constant to grow.
To derive useful late-time grading one must certify the finite accuracy
budget, or retain more favorable signed or gate-dependent growth terms.
This calculation does not show that such a sharper route is impossible.
\end{remark}

\subsection{The weak pure-sector calculation}
Put $\widetilde u(\phi)=(\cos\phi,\sin\phi)$. The following is conditional
on the displayed coordinate envelopes; it asserts no retention on its own.
Here $L_{\delta_{\mathrm{strong}}}>0$ denotes the supplied envelope factor.
The coordinate-envelope event means that every cluster-1 sample satisfies
$x_1\ge m_1^{\rm data}$ and $|x_2|\le b_2$, and every cluster-2 sample
satisfies $|x_1|\le b_1$ and $x_2\ge m_2^{\rm data}$, with the constants
defined in the next lemma.


\begin{lemma}[Quadrant-II cluster-$2$ pure-sector residual reduction]
\label{sbagg:lem:cluster2-pure-residual-reduction}

Use the notation
\[
b_1:=\sigma_1L_{\delta_{\mathrm{strong}}},
\qquad
b_2:=\sigma_2L_{\delta_{\mathrm{strong}}},
\]
\[
m_1^{\mathrm{data}}:=\mu_1-b_1,
\qquad
m_2^{\mathrm{data}}:=\mu_2-b_2,
\]
and let $0<\theta_G<\theta_C<\pi/2$ satisfy
\[
\cos\theta_C\,m_1^{\mathrm{data}}
-\sin\theta_C\,b_2>0,
\]
\[
\sin\theta_G\,m_2^{\mathrm{data}}
-\cos\theta_G\,b_1>0.
\]

Define
\[
\boxed{
J_{C,2}^{\mathrm{pure}}
:=
[\pi-\theta_C,\pi-\theta_G].
}
\]

Then, on the coordinate-envelope event, for every
$\phi\in J_{C,2}^{\mathrm{pure}}$,
\[
\boxed{
\widetilde u(\phi)^\top x_k^{(1)}<0,
\qquad
\widetilde u(\phi)^\top x_k^{(2)}>0
\qquad
\forall k.
}
\]

Consequently, if
\[
\mathcal R_2(t)
:=
\mathbb E_{2,n}
\left[
(x-f_{\parallel,t}(x))x^\top
\right],
\]
then every neuron satisfying
$\phi_i(t)\in J_{C,2}^{\mathrm{pure}}$ obeys
\[
\boxed{
A_i^\rho(t)=\frac12\mathcal R_2(t).
}
\]

Moreover, if
\[
\mathcal X_2(t)
:=
\mathbb E_{2,n}
[f_{t,1}(x)x_2-f_{t,2}(x)x_1],
\]
then
\[
\boxed{
R_{21}^{(2)}
=
R_{12}^{(2)}+\mathcal X_2.
}
\]

Define the aligned angular proxy by $V_{C,2}^\rho=\widetilde u^\top(\mathcal R_2/2)\widetilde u^\angle$, where $\widetilde u^\angle=(-\sin\phi,\cos\phi)$. It is not the actual mismatched neuron angle velocity. For $\phi\in J_{C,2}^{\mathrm{pure}}$,
\[
\boxed{
V_{C,2}^\rho(\phi;t)
=
\frac12
\left[
-|s_\phi c_\phi|
(R_{22}^{(2)}-R_{11}^{(2)})
+
R_{12}^{(2)}\cos2\phi
-
\mathcal X_2s_\phi^2
\right].
}
\]
\end{lemma}

\begin{proof}
Write $\phi=\pi-\theta$ with
$\theta\in[\theta_G,\theta_C]$.
Then
\[
\cos\phi=-\cos\theta,
\qquad
\sin\phi=\sin\theta.
\]
For concept $1$,
\[
\widetilde u(\phi)^\top x^{(1)}
\le
-\cos\theta\,m_1^{\mathrm{data}}
+\sin\theta\,b_2
\le
-\cos\theta_C\,m_1^{\mathrm{data}}
+\sin\theta_C\,b_2
<0.
\]
For concept $2$,
\[
\widetilde u(\phi)^\top x^{(2)}
\ge
-\cos\theta\,b_1
+\sin\theta\,m_2^{\mathrm{data}}
\ge
-\cos\theta_G\,b_1
+\sin\theta_G\,m_2^{\mathrm{data}}
>0.
\]
Hence the selector equals zero on cluster $1$ and one on cluster $2$,
which gives $A_i^\rho=\mathcal R_2/2$.

The cross identity follows by subtraction:
\[
R_{21}^{(2)}-R_{12}^{(2)}
=
\mathbb E_{2,n}[f_1x_2-f_2x_1]
=
\mathcal X_2.
\]
Finally,
\[
\widetilde u(\phi)^\top\mathcal R_2\widetilde u^\angle(\phi)
=
c_\phi s_\phi(R_{22}^{(2)}-R_{11}^{(2)})
+
R_{12}^{(2)}\cos2\phi
-
\mathcal X_2s_\phi^2,
\]
and $c_\phi s_\phi<0$ on $J_{C,2}^{\mathrm{pure}}$.
\end{proof}


\begin{corollary}[Weak pure-sector absolute transverse energy]
\label{sbagg:weak-Q}
Let a fixed nonempty set $B$ remain in $J_{C,2}^{\rm pure}$ on $[a,b]$.
Write $R_{rs}=R_{rs}^{(2)}$ and define
\[
 M=\sum_{B}v_{i2}^2,\quad Q=\sum_Bv_{i1}^2,\quad
 \Delta=\sum_B\|\delta_i\|^2,\quad
 \mathcal P=\left(\sum_B\|P_{i,\perp}^\rho\|^2\right)^{1/2}.
\]
The exact coordinate equations are
\[
 \dot v_{i1}=\tfrac12[R_{11}(v_{i1}-\delta_{i1})
                       +R_{21}(v_{i2}-\delta_{i2})]-P_{i1},
\]
\[
 \dot v_{i2}=\tfrac12[R_{12}(v_{i1}-\delta_{i1})
                       +R_{22}(v_{i2}-\delta_{i2})]-P_{i2}.
\]
Consequently, in the almost-everywhere comparison sense,
\[
 \frac{d}{dt}\sqrt M\ge
 \tfrac12R_{22}\sqrt M-\tfrac12|R_{12}|\sqrt Q
 -\tfrac12\sqrt{R_{12}^2+R_{22}^2}\sqrt\Delta-\mathcal P,
\]
\[
 \frac{d}{dt}\sqrt Q\le
 \tfrac12|R_{21}|\sqrt M+\tfrac12|R_{11}|\sqrt Q
 +\tfrac12\sqrt{R_{11}^2+R_{21}^2}\sqrt\Delta+\mathcal P.
\]
If $|R_{11}|\le\beta_{11}$, $|R_{21}|\le\beta_{21}$,
$M\le S_*$, $\sqrt\Delta\le D_B$, and $\mathcal P\le P_B$, then
\[
 \sqrt{Q(t)}\le e^{\beta_{11}(t-a)/2}\sqrt{Q(a)}
       +H_B\mathcal E_{\beta_{11}/2}(t-a),
\]
\[
 H_B=\tfrac12\beta_{21}\sqrt{S_*}
      +\tfrac12\sqrt{\beta_{11}^2+\beta_{21}^2}D_B+P_B.
\]
If $\sup_{\mathcal D}\|f_t\|\le F_\star$, valid choices are
\[
 \beta_{11}=s_{11}^{(2)}+F_\star\sqrt{s_{11}^{(2)}},\qquad
 \beta_{21}=|s_{12}^{(2)}|+F_\star\sqrt{s_{11}^{(2)}}.
\]
In the limit $R_{11}=R_{12}=R_{21}=0$, $\delta_i=P_i=0$,
one has $\dot Q=0$ and $\dot M=R_{22}M$.
\end{corollary}
\begin{proof}
Substitute $A_i^\rho=\mathcal R_2/2$ in the input equation. Take the
Euclidean norm over $B$ in each coordinate and use Cauchy--Schwarz for
the mismatch and perpendicular terms. Integrate the resulting scalar
inequality for $\sqrt Q$. Finally
$R_{r1}=s_{r1}^{(2)}-\mathbb E_{2,n}[f_r x_1]$ and
$|\mathbb E_{2,n}[f_r x_1]|\le F_\star\sqrt{s_{11}^{(2)}}$.
\end{proof}

\begin{remark}[Retention and the interface to Theorem 3 remain separate]
The pure arc is antipodal to the cluster-1-only arc under the same envelope
conditions. If a valid invocation of the bias-corrected T1.8--T1.9 proves
retention and its arcs satisfy
\[
 J_W^{\rm seed}\Subset J_W^{\rm ret}\Subset J_{C,2}^{\rm pure},
\]
then the pure-sector identities apply on that invocation's horizon.
But changing the arcs changes their buffers, angular suprema, and retention
budgets. Their compatibility, seed coverage, and a nonempty admissible
parameter region must be checked; no unconditional handoff is proved here.

The optimal partition is not a partition by actual probe gates. For a fixed
probe cell, $M_\omega=\sum_i\omega_i z_iv_i^\top$ can contain longitudinal
mass from both axis groups. Axis defect alone neither controls its required
transverse diagonal nor the signs of the coefficient evolution. One must
establish the cell's active-group contributions and then prove the
$\mathcal G$ sign margins and derivative remainder needed for reversal.
No identification of $\eta$ with the noise scale is used above.
\end{remark}

%%%%%%%%%%%%%%%%%%%%%%%%%%%%%%%%%%%%%%%%%%%%%%%%%%%%%%%%%%%%%%%%%%%%%%%%%%%%%%%
% Post-B1 strong-tilt coherence and weak-rotation clock
%%%%%%%%%%%%%%%%%%%%%%%%%%%%%%%%%%%%%%%%%%%%%%%%%%%%%%%%%%%%%%%%%%%%%%%%%%%%%%%

\subsection{Post-B1 strong-tilt coherence: population branch, spread, and weak rotation}
\label{sbcoh:scope}

This subsection records a refinement of the B2 interface that is useful when
one wants to retain the $c_{12}^O$ feedback channel in addition to $c_{21}^O$.
It separates two different notions that were previously conflated:
\emph{collective branch stability}, in which the strong block and the test
neuron move together, and \emph{pairwise spread contraction}, in which the
collective network state is held fixed while one strong neuron's own tilt is
varied.  The second condition is strictly stronger.  The calculations below
also identify the retained weak block's signed off-diagonal coefficient as a
separate pre-creation stopping quantity.  None of these statements closes the
random-initialization trajectory or Theorem~1 by itself.

To avoid notation collisions, write
\[
 \epsilon_\sigma:=\frac{\sigma}{\mu_2},
 \qquad
 \varrho:=\frac{\mu_1}{\mu_2}>1,
 \qquad
 \lambda_2:=\frac{\mu_2^2}{2},
\]
and, only in the initialization-scaling lemma below, write
$\varepsilon_0:=\varepsilon$ for the manuscript's initialization amplitude.
Let $\varphi$ and $\Phi$ denote the standard-normal density and distribution
function and set $c_0=\varphi(0)=1/\sqrt{2\pi}$.

\begin{lemma}[Exact positive-side tilt kernel and centered $A_{12}$ identity]
\label{sbcoh:tilt-kernel}
Set
\[
 p=(x_1)_+,
 \qquad
 D=\|p\|_n^2,
 \qquad
 b_*:=\frac{\langle p,x_2\rangle_n}{D},
 \qquad
 \alpha_0=\mathbf 1_{\{x_1>0\}},
\]
\[
 h_{\rm tilt}:=x_2\alpha_0-b_*p,
 \qquad
 D_{\rm tilt}:=\|h_{\rm tilt}\|_n^2.
\]
For
\[
 q_\phi(x):=(\widetilde u(\phi)^\top x)_+,
 \qquad
 J_{\rm tilt}(\phi):=\langle h_{\rm tilt},q_\phi\rangle_n,
\]
one has $\langle p,h_{\rm tilt}\rangle_n=0$.  If
$0\le\phi\le\bar\phi<\pi/2$ lies in a cone on which every training sample
with $x_1>0$ is active, then exactly
\begin{equation}\label{sbcoh:tilt-kernel-exact}
 \boxed{
 J_{\rm tilt}(\phi)=D_{\rm tilt}\sin\phi.
 }
\end{equation}
Moreover, if
\[
 m:=\frac{\langle p,f_1\rangle_n}{D},
\]
then at the strong axis
\begin{equation}\label{sbcoh:A12-axis}
 \boxed{
 A_{12}^\rho(0)
 =Db_*(1-m)
 -\sum_i r_i z_{i1}J_{\rm tilt}(\phi_i).
 }
\end{equation}
Consequently, for a positive cone-pure strong block $B$,
\begin{equation}\label{sbcoh:A12-strong-exact}
 \sum_{i\in B}r_i z_{i1}J_{\rm tilt}(\phi_i)
 =D_{\rm tilt}M_{12,B}^{\rm raw},
 \qquad
 M_{12,B}^{\rm raw}:=\sum_{i\in B}z_{i1}v_{i2},
\end{equation}
and, with $J_B:=\sum_{i\in B}v_{i1}v_{i2}$ and
$Z_B:=\sum_{i\in B}v_{i2}^2$,
\begin{equation}\label{sbcoh:raw-vs-input-cross}
 \boxed{
 |M_{12,B}^{\rm raw}-J_B|
 \le\sqrt{Z_B\Delta_{B,1}},
 \qquad
 \Delta_{B,1}:=\sum_{i\in B}\delta_{i1}^2.
 }
\end{equation}
\end{lemma}

\begin{proof}
On $\{x_1\le0\}$ both $p$ and $h_{\rm tilt}$ vanish.  On the support of
$h_{\rm tilt}$, the assumed purity gives
$q_\phi=x_1\cos\phi+x_2\sin\phi$.  Hence
\[
 J_{\rm tilt}(\phi)
 =\cos\phi\langle h_{\rm tilt},p\rangle_n
 +\sin\phi\langle h_{\rm tilt},x_2\alpha_0\rangle_n
 =D_{\rm tilt}\sin\phi,
\]
because
$\langle h_{\rm tilt},x_2\alpha_0\rangle_n=\|h_{\rm tilt}\|_n^2$.
Write $f_1=mp+r_1$ with $\langle p,r_1\rangle_n=0$.  Then
\[
 A_{12}^\rho(0)
 =\langle x_1-f_1,x_2\alpha_0\rangle_n
 =Db_*(1-m)-\langle r_1,h_{\rm tilt}\rangle_n.
\]
Indeed, if
$H_1(\phi):=\langle p,q_\phi\rangle_n$, then
\[
 r_1=\sum_i r_i z_{i1}
 \left(q_{\phi_i}-\frac{H_1(\phi_i)}{D}p\right).
\]
Using $h_{\rm tilt}\perp p$ therefore yields
\eqref{sbcoh:A12-axis}.  On $B$,
$r_i\sin\phi_i=v_{i2}$, so \eqref{sbcoh:tilt-kernel-exact} gives
\eqref{sbcoh:A12-strong-exact}.  Finally
$M_{12,B}^{\rm raw}-J_B=-\sum_B\delta_{i1}v_{i2}$ and Cauchy--Schwarz gives
\eqref{sbcoh:raw-vs-input-cross}.
\end{proof}

\begin{remark}[The moving slab is part of the field, not a perturbation]
\label{sbcoh:slab-not-error}
At the natural strong tilt
$\xi=\epsilon_\sigma s$, a cluster-$2$ gate switches an $O(1)$ fraction of
that cluster when $s=O(1)$.  Therefore a moving-slab estimate is not a small
relative error at the scale of the $A_{12}$ source.  On the positive side the
kernel identity \eqref{sbcoh:tilt-kernel-exact} removes the need for a kernel
modulus entirely.  In the population reduction below the cluster-$2$ slab is
kept with its sign inside the scalar field.
\end{remark}

\begin{proposition}[Coherent high-SNR population strong-tilt field]
\label{sbcoh:population-field}
Consider the equal-noise two-cluster population model and a compact rescaled
band $0\le s\le s_+$ with $\epsilon_\sigma s_+\ll1$.  Suppose the strong
block $B$ is cone-pure on cluster $1$ and coherent at tilt
\[
 \xi=\epsilon_\sigma s.
\]
Let
\[
 M_B:=M_{11,B}^{\rm raw}.
\]
The load-bearing matched reduction is the case in which the progress outside
$B$ is placed in a separate remainder and
\begin{equation}\label{sbcoh:matched-progress}
 M_B=m+o(1).
\end{equation}
This is the natural coherent regime: by Lemma~\ref{sbgap:progress},
$m_B^{\rm eff}=M_B+b_*M_{12,B}^{\rm raw}$ and
$b_*M_{12,B}^{\rm raw}=O(\epsilon_\sigma^2)$ on a noise-scale coherent tube.

Model a cluster-$2$-pure coherent weak block by the planar coefficient matrix
\[
 M_W=
 \begin{pmatrix}
 S_W-w&-c_W\\
 -c_W&w
 \end{pmatrix},
 \qquad
 c_W\ge0,
 \qquad
 \chi_W:=\frac{c_W}{\epsilon_\sigma}.
\]
Then, after division of the strong tilt equation by $\mu_2^2$, the principal
matched population field is
\begin{equation}\label{sbcoh:C1}
 \boxed{
 2\mathfrak F_{\rm coll}(s)
 =(1-m)K_\varrho(s)+W(s),
 }
\end{equation}
where
\begin{equation}\label{sbcoh:C2}
 K_\varrho(s):=\varphi(s)+s\Phi(s)-\varrho^2s,
 \qquad
 W(s):=-(S_W-w)\varphi(s)+(\chi_W-ws)\Phi(s).
\end{equation}
Uniformly on a fixed compact band, the omitted finite-noise, non-$B$, mismatch,
and sampling terms are collected in a remainder
$\mathcal R_{\rm pop}(s,t)$ that must be bounded separately in the finite-sample
transfer.

Differentiating along the collective diagonal, where the block tilt and the
test tilt move together, gives
\begin{equation}\label{sbcoh:C3}
 \boxed{
 2\partial_s^{\rm tot}\mathfrak F_{\rm coll}
 =-\varrho^2(1-m)
 +(1-w-m)\Phi(s)
 +[\chi_W+(S_W-2w)s]\varphi(s).
 }
\end{equation}
Holding the common network state fixed and varying only one strong neuron's
own gate gives instead
\begin{equation}\label{sbcoh:own-partial}
 \boxed{
 2\partial_s^{\rm own}\mathfrak F
 =-\varrho^2(1-m)
 +(1-w)\Phi(s)
 +[\chi_W+(S_W-2w)s]\varphi(s).
 }
\end{equation}
Thus
\begin{equation}\label{sbcoh:own-total-split}
 \boxed{
 2\partial_s^{\rm tot}\mathfrak F
 =2\partial_s^{\rm own}\mathfrak F-m\Phi(s).
 }
\end{equation}
\end{proposition}

\begin{remark}[Matched branch and canonical rescaled zero]
\label{sbcoh:matched-branch}
With no weak block and no non-$B$ progress, \eqref{sbcoh:C1} becomes
\[
 2\mathfrak F_{\rm coll}(s)=(1-m)K_\varrho(s).
\]
Because
$K_\varrho'(s)=\Phi(s)-\varrho^2<0$, the collective branch has a unique
positive zero $s_\varrho$ determined by
\begin{equation}\label{sbcoh:s-rho}
 \varphi(s_\varrho)+s_\varrho\Phi(s_\varrho)
 =\varrho^2s_\varrho.
\end{equation}
For the canonical ratio $\varrho=3/2$,
\[
 \boxed{s_\varrho\approx0.23419,\qquad
 \Phi(s_\varrho)\approx0.5926.}
\]
Hence the principal population prediction is
$\xi_{\rm eq}=(\sigma/\mu_2)s_\varrho$ rather than a universal constant times
$\mu_2\sigma/\mu_1^2$.  The latter is only the large-$\varrho$ expansion
$s_\varrho=\varphi(0)/(\varrho^2-1/2)+O(\varrho^{-6})$ followed by one more
large-$\varrho$ simplification.  This is a derived population prediction,
not a post-hoc theorem claim about the frozen experiment.
\end{remark}

\begin{remark}[Collective stability versus spread contraction]
\label{sbcoh:two-derivatives}
The favorable $-m\Phi(s)$ term in \eqref{sbcoh:own-total-split} is collective
self-regulation and does not contract differences between two neurons when the
common network state is frozen.  Pairwise coherence must therefore use
\eqref{sbcoh:own-partial}, not \eqref{sbcoh:C3}.  In particular, with no weak
block a sufficient full-positive-band population condition is
\[
 \varrho^2(1-m)>1,
 \qquad\text{i.e.}\qquad
 m<1-\varrho^{-2}.
\]
For $\varrho=3/2$ this is $m<5/9$.  On a local band
$0\le s\le s_+$ it sharpens to
$m<1-\Phi(s_+)/\varrho^2$ before weak/error charges; at the canonical branch
center this threshold is approximately $0.7366$.
\end{remark}

\begin{lemma}[Finite-sample spread contraction from the own partial]
\label{sbcoh:spread}
Suppose on a stopped B2 interval and a compact rescaled band $I_s$ every
strong neuron satisfies
\begin{equation}\label{sbcoh:field-transfer}
 \frac{\dot s_i}{\mu_2^2}
 =\mathfrak F(s_i;\mathcal C_t)+e_i(t),
 \qquad
 |e_i(t)|\le\varepsilon_{\rm add},
\end{equation}
where $\mathcal C_t$ is common to the strong block.  Assume
\begin{equation}\label{sbcoh:own-margin}
 \sup_{t,s\in I_s}
 2\partial_s^{\rm own}\mathfrak F(s;\mathcal C_t)
 \le-\mathfrak b
 \qquad(\mathfrak b>0).
\end{equation}
Then for every pair $i,j$,
\begin{equation}\label{sbcoh:spread-dini}
 D^+|s_i-s_j|
 \le-\frac{\mu_2^2\mathfrak b}{2}|s_i-s_j|
      +2\mu_2^2\varepsilon_{\rm add},
\end{equation}
and therefore, for $t\ge a$,
\begin{equation}\label{sbcoh:spread-floor}
 \boxed{
 |s_i(t)-s_j(t)|
 \le e^{-\mu_2^2\mathfrak b(t-a)/2}|s_i(a)-s_j(a)|
 +\frac{4\varepsilon_{\rm add}}{\mathfrak b}.
 }
\end{equation}
\end{lemma}

\begin{proof}
Subtract the two equations in \eqref{sbcoh:field-transfer}, apply the mean
value theorem in the own-tilt coordinate with the collective state held fixed,
and use \eqref{sbcoh:own-margin}.  Scalar comparison gives
\eqref{sbcoh:spread-floor}.
\end{proof}

\begin{remark}[Canonical own-partial and additive-floor budgets]
\label{sbcoh:canonical-spread-budget}
At $\varrho=3/2$ and $m\le1/2$, dropping favorable $w$-terms gives the
conservative doubled own-partial budget
\begin{equation}\label{sbcoh:canonical-own-budget}
 \boxed{
 \mathfrak b
 \ge\frac18
 -c_0\chi_W^{\max}
 -\varphi(1)S_W^{\max}
 -\mathcal B_{\rm nb}
 -\mathcal B_{\rm coeff}.
 }
\end{equation}
Here
$\sup_{s\ge0}s\varphi(s)=\varphi(1)$ and
$(1/8)/c_0\approx0.3133$, so if all other charges vanished one could take
$\chi_W^{\max}<0.3133$.  Preserving merely a positive slope margin is not the
tight condition, however.  If one wants a positive tube with, for example,
upper/lower tilt ratio at most two around $s_\varrho$, a convenient sufficient
condition is
\[
 \frac{4\varepsilon_{\rm add}}{\mathfrak b}\lesssim\frac{s_\varrho}{3}.
\]
At $\mathfrak b=1/8$ and $s_\varrho=0.23419$ this is
\begin{equation}\label{sbcoh:additive-floor-canonical}
 \boxed{
 \varepsilon_{\rm add}\lesssim2.44\times10^{-3}.
 }
\end{equation}
Thus a generic uniform concentration bound of size
$C\epsilon_\sigma^{-1}\sqrt{d/n}$ would require, only at the level of this
coarse comparison,
\[
 n\gg 1.7\times10^5 C^2 d\left(\frac{\mu_2}{\sigma}\right)^2.
\]
This large constant is a proof-budget warning, not a claim that the mechanism
requires such a sample size.
\end{remark}

\begin{remark}[Only the transverse mismatch needs noise-scale accuracy]
\label{sbcoh:directional-mismatch}
The mismatch contribution to the exact ratio equation can be kept directional.
From
$\dot v=A^\top v-A^\top\delta-P$, the part entering $\dot\xi$ is
\begin{equation}\label{sbcoh:directional-mismatch-eq}
 \frac{(A^\top\delta)_2-\xi(A^\top\delta)_1}{v_1}
 =\frac{(A_{22}-\xi A_{21})\delta_2
       +(A_{12}-\xi A_{11})\delta_1}{v_1}.
\end{equation}
On a noise-scale branch, $A_{22}=O(\mu_2^2)$ while
$A_{12}-\xi A_{11}=O(\epsilon_\sigma\mu_2^2)$.  After rescaling
$s=\xi/\epsilon_\sigma$, the $\delta_2$ term is amplified by
$1/\epsilon_\sigma$ but the $\delta_1$ term is not.  Thus the finite-sample
coherence interface should impose a noise-scale bound on
$|\delta_{i2}|/r_i$ and only an ordinary smallness bound on
$|\delta_{i1}|/r_i$, rather than spending the isotropic estimate
\eqref{eq:late-hold-error} in both directions.
\end{remark}

\begin{lemma}[Exact signed weak off-diagonal evolution]
\label{sbcoh:weak-cross}
Let a fixed cohort $W$ remain in $J_{C,2}^{\rm pure}$.  By
Lemma~\ref{sbagg:lem:cluster2-pure-residual-reduction}, every member has
$A_i^\rho=A^W:=\mathcal R_2/2$.  Set
\[
 M_W^{\rm raw}:=\sum_{i\in W}z_iv_i^\top,
 \qquad
 C_W:=-M_{12,W}^{\rm raw},
\]
\[
 C_W^v:=\sum_{i\in W}v_iv_i^\top,
 \qquad
 C_W^z:=\sum_{i\in W}z_iz_i^\top,
 \qquad
 D_W:=(C_W^v)_{22}+(C_W^z)_{11}.
\]
Then exactly, almost everywhere,
\begin{equation}\label{sbcoh:weak-cross-ode}
 \boxed{
 \dot C_W
 =(A_{11}^W+A_{22}^W)C_W-A_{12}^W D_W+\mathcal R_{C,W},
 }
\end{equation}
where
\begin{equation}\label{sbcoh:weak-cross-rem}
 \mathcal R_{C,W}
 =-A_{11}^W\sum_W\delta_{i1}v_{i2}
 +A_{22}^W\sum_W z_{i1}\delta_{i2}
 +\sum_W z_{i1}P_{i2}.
\end{equation}
Consequently, with
$\Delta_W=\sum_W\|\delta_i\|^2$ and
$\mathcal P_W=(\sum_W\|P_{i,\perp}^\rho\|^2)^{1/2}$,
\begin{equation}\label{sbcoh:weak-cross-rem-bound}
 |\mathcal R_{C,W}|
 \le(|A_{11}^W|+|A_{22}^W|)\sqrt{S_W\Delta_W}
      +\sqrt{S_W}\,\mathcal P_W
\end{equation}
whenever $S_W$ dominates the displayed weak coordinate energies.
\end{lemma}

\begin{proof}
Differentiate
$M_W^{\rm raw}=\sum_W z_iv_i^\top$ using
$\dot z_i=A^Wv_i$ and
$\dot v_i=(A^W)^\top z_i-P_i$.  This gives
\[
 \dot M_W^{\rm raw}=A^W C_W^v+C_W^zA^W-\sum_Wz_iP_i^\top.
\]
Take the $(1,2)$ entry and use
\[
 (C_W^v)_{12}=M_{12,W}^{\rm raw}+\sum_W\delta_{i1}v_{i2},
\]
\[
 (C_W^z)_{12}=M_{12,W}^{\rm raw}-\sum_Wz_{i1}\delta_{i2}.
\]
Negating the resulting $M_{12,W}^{\rm raw}$ equation proves
\eqref{sbcoh:weak-cross-ode}--\eqref{sbcoh:weak-cross-rem}; Cauchy--Schwarz
gives the stated bound.
\end{proof}

\begin{proposition}[Signed weak-rotation clock]
\label{sbcoh:weak-rotation}
In the population coherent reduction, define
\[
 K_1(s):=\varphi(s)+s\Phi(s),
 \qquad
 \chi_W:=\frac{C_W}{\epsilon_\sigma}.
\]
Uniformly on a compact B2 band, the pure weak residual coefficients have the
leading form
\begin{equation}\label{sbcoh:weak-A12-pop}
 A_{12}^W
 =\epsilon_\sigma\lambda_2[\chi_W-mK_1(s)]
 +\mathcal E_{12}^W,
\end{equation}
\begin{equation}\label{sbcoh:weak-trace-pop}
 A_{11}^W+A_{22}^W
 =\lambda_2(1-w)+\mathcal E_{\rm tr}^W.
\end{equation}
If the cluster-$2$ wrong output satisfies
$\|f_1\|_{L^2(\text{cluster }2)}\le C_{1,2}\sigma$, then
$\mathcal E_{\rm tr}^W=O(\sigma^2)$; the manuscript's generic
$F_\star$ bound alone is not sufficient for this sharper order.
Dividing \eqref{sbcoh:weak-cross-ode} by $\epsilon_\sigma$ gives
\begin{equation}\label{sbcoh:chi-exact-interface}
 \boxed{
 \dot\chi_W
 =[\lambda_2(1-w)+\mathcal E_{\rm tr}^W]\chi_W
 +\lambda_2D_W[mK_1(s)-\chi_W]
 +\mathcal E_{\chi,W},
 }
\end{equation}
where
\[
 \mathcal E_{\chi,W}
 =-\frac{D_W}{\epsilon_\sigma}\mathcal E_{12}^W
  +\frac{\mathcal R_{C,W}}{\epsilon_\sigma}.
\]
In the ideal rank-one weak block $D_W=S_W$, so the leading signed equation is
\begin{equation}\label{sbcoh:chi-leading}
 \boxed{
 \dot\chi_W
 =\lambda_2[(1-w-S_W)\chi_W+S_WmK_1(s)].
 }
\end{equation}
In particular the strong output forces $\chi_W$ in the adverse positive
direction even from $\chi_W=0$.
\end{proposition}

\begin{remark}[Weak-feature timing does not replace the $\chi_W$ face]
\label{sbcoh:gamma-vs-chi}
For the sign convention
$M_W=\left(\begin{smallmatrix}S_W-w&-c_W\\-c_W&w\end{smallmatrix}\right)$,
the creation coefficient has, to leading order,
\[
 \Gamma_W=(w-\tau_Mc_W)_+.
\]
For the opposite off-diagonal sign it is
$(w+\tau_Mc_W)_+$ and then
$\Gamma_W\le\Gamma_W^{\max}$ implies
$\chi_W\lesssim\Gamma_W^{\max}/(\tau_M\epsilon_\sigma)$.
Thus the existing weak-feature timing guard is not a two-sided control of
$\chi_W$; even in the sign where a relation exists, it is useful at the
noise-scale coherence level only if $\Gamma_W^{\max}=O(\epsilon_\sigma)$.
The retained quadrant-II sign is precisely the case in which increasing
$c_W$ can reduce the positive-part creation coefficient while worsening the
strong spread field.
\end{remark}

\begin{corollary}[Fixed-angle weak mass is noise-scale before the coherence face]
\label{sbcoh:weak-mass-compat}
Suppose the retained weak cohort has input angles
$\phi_i=\pi-\theta_i$ with
$\theta_i\in[\theta_G,\theta_C]\Subset(0,\pi/2)$.  Put
\[
 S_W:=\sum_{i\in W}r_i^2,
 \qquad
 q_W:=\sin\theta_G\cos\theta_C.
\]
Then
\begin{equation}\label{sbcoh:CW-angle-lower}
 C_W
 \ge q_WS_W-\sqrt{S_W\Delta_{W,1}},
 \qquad
 \Delta_{W,1}:=\sum_W\delta_{i1}^2.
\end{equation}
Hence, if
$q_W^\sharp:=q_W-\sqrt{\Delta_{W,1}/S_W}>0$ and
$\chi_W\le\chi_{\max}$, then
\begin{equation}\label{sbcoh:weak-mass-noise-scale}
 \boxed{
 S_W\le\frac{\epsilon_\sigma\chi_{\max}}{q_W^\sharp}.
 }
\end{equation}
Thus a fixed-angle retained weak cohort can remain both nontrivially rotated
and harmless to strong coherence only while its pre-creation mass is itself
of noise scale.  This is consistent with the manuscript's chronology, in
which weak capture and later mass growth occur after basin creation.
\end{corollary}

\begin{proof}
For $\phi_i=\pi-\theta_i$,
$-v_{i1}v_{i2}=r_i^2\sin\theta_i\cos\theta_i\ge q_Wr_i^2$.
Since $z_{i1}=v_{i1}-\delta_{i1}$,
\[
 C_W=-\sum_Wz_{i1}v_{i2}
 \ge q_WS_W-\left|\sum_W\delta_{i1}v_{i2}\right|,
\]
and Cauchy--Schwarz proves the claim.
\end{proof}

\begin{corollary}[A deterministic weak-rotation deadline]
\label{sbcoh:chi-deadline}
Suppose on $[t_a,t_a+H]$ that
$m\le m_{\max}$, $s\le s_+$, $\chi_W\le\chi_{\max}$,
\[
 \lambda_2(1-w)+\mathcal E_{\rm tr}^W\le L_\chi^+,
 \qquad
 D_W\le D_W^+,
 \qquad
 |\mathcal E_{\chi,W}|\le E_\chi^+.
\]
Dropping the favorable term $-\lambda_2D_W\chi_W$ in
\eqref{sbcoh:chi-exact-interface} gives
\begin{equation}\label{sbcoh:chi-duhamel}
 \boxed{
 \chi_W(t_a+H)
 \le e^{L_\chi^+H}\chi_W(t_a)
 +\frac{F_\chi^+}{L_\chi^+}(e^{L_\chi^+H}-1),
 }
\end{equation}
where
\[
 F_\chi^+:=\lambda_2D_W^+m_{\max}K_1(s_+)+E_\chi^+.
\]
Consequently the first time
$t_\chi:=\inf\{t\ge t_a:\chi_W(t)=\chi_{\max}\}$
obeys $t_\chi>t_a+H$ whenever the right side of
\eqref{sbcoh:chi-duhamel} is strictly below $\chi_{\max}$.
If the regulator clock satisfies
\begin{equation}\label{sbcoh:chi-before-hit-clock}
 \boxed{
 \Lambda(t_a,t_\chi)
 >\log\frac{G_M(t_a)+\delta_{\rm gen}}{\delta_{\rm gen}},
 }
\end{equation}
then Proposition~\ref{sbgen:first-hit} certifies
$t_{\rm cert}<t_\chi$.
\end{corollary}

\begin{lemma}[Initialization/noise ordering interface]
\label{sbcoh:init-noise}
Assume an early comparison supplies constants
$\kappa_W,\kappa_B,m_a,g_W,g_B>0$ such that
\[
 C_W(0)\le\kappa_W\varepsilon_0^2,
 \qquad
 M_B(0)\ge\kappa_B\varepsilon_0^2,
\]
\[
 D^+\log C_W\le g_W,
 \qquad
 D^-\log M_B\ge g_B
\]
until B2 entry $t_a$, where $M_B(t_a)\ge m_a$.  Put
\[
 \nu:=\frac{g_W}{g_B}.
\]
Then
\begin{equation}\label{sbcoh:init-chi}
 \boxed{
 \chi_W(t_a)
 \le
 \frac{\kappa_W\kappa_B^{-\nu}m_a^\nu}{\epsilon_\sigma}
 \varepsilon_0^{2(1-\nu)}.
 }
\end{equation}
If the principal early rates are
$g_W=\lambda_2+o(1)$ and $g_B=2\lambda_1+o(1)$ with
$\lambda_1=\mu_1^2/2$, then
\[
 \nu=\frac1{2\varrho^2}+o(1).
\]
For $\varrho=3/2$, $\nu=2/9+o(1)$ and the initialization exponent is
$2(1-\nu)=14/9+o(1)$.
\end{lemma}

\begin{proof}
The lower strong-growth comparison implies
\[
 t_a\le g_B^{-1}\log\frac{m_a}{\kappa_B\varepsilon_0^2}.
\]
Insert this time in the upper weak-cross comparison
$C_W(t_a)\le\kappa_W\varepsilon_0^2e^{g_Wt_a}$ and divide by
$\epsilon_\sigma$.
\end{proof}

Combining Lemma~\ref{sbcoh:init-noise} with
Corollary~\ref{sbcoh:chi-deadline} yields the explicit sufficient ordering
condition
\begin{equation}\label{sbcoh:init-noise-hit}
 \boxed{
 \frac{\kappa_W\kappa_B^{-\nu}m_a^\nu}{\epsilon_\sigma}
 \varepsilon_0^{2(1-\nu)}e^{L_\chi^+H_{\rm hit}}
 +\frac{F_\chi^+}{L_\chi^+}(e^{L_\chi^+H_{\rm hit}}-1)
 <\chi_{\max}.
 }
\end{equation}
At canonical signal ratio this has the headline scaling
\[
 \varepsilon_0^{14/9}e^{\lambda_2H_{\rm hit}}
 \lesssim \epsilon_\sigma\chi_{\max},
\]
up to the explicit cohort constants and finite-sample forcing allowances.
This is a genuine initialization-scale versus noise-scale condition for
retaining the $c_{12}^O$ coherence channel through creation.

\begin{remark}[Frozen canonical audit: weak rotation is not the observed bottleneck]
\label{sbcoh:canonical-audit}
This paragraph is sizing evidence only and is not used by the theorem.  The
frozen canonical diagnostic uses
\[
 \mu_1=3,\quad \mu_2=2,\quad \sigma=0.15,
 \quad h=200,\quad n=2000\ \text{per cluster},
 \quad \varepsilon_0=10^{-3}.
\]
Thus $\epsilon_\sigma=0.075$ and $\lambda_2=2$.  The frozen sweep records
creation at approximately
\[
 t_{\rm create}=1.06411,
 \qquad
 m_{\rm create}=0.43750,
 \qquad
 L_W(t_{\rm create})=1.7466\times10^{-3}.
\]
Even pessimistically using the full interval from initialization to creation
as $H_{\rm hit}$,
\[
 \varepsilon_0^{14/9}e^{\lambda_2t_{\rm create}}
 \approx1.81\times10^{-4},
 \qquad
 \epsilon_\sigma(0.3)=2.25\times10^{-2},
\]
so the first term in \eqref{sbcoh:init-noise-hit} leaves a factor about
$124$ before cohort constants.  The generic bound
$|M_{12,W}^{\rm raw}|\le L_W/2$ is always valid.  In the matched planar
regime it sharpens, with an explicit mismatch remainder, to
\[
 \boxed{
 C_W=-M_{12,W}^{\rm raw}
 \le \frac{L_W}{4}
 +\sqrt{\Delta_{W,1}Q_W^v},
 \qquad Q_W^v:=\sum_{i\in W}v_{i2}^2.
 }
\]
Indeed
$C_W=-\sum_Wv_{i1}v_{i2}+\sum_W\delta_{i1}v_{i2}$,
$|\sum_Wv_{i1}v_{i2}|\le\frac12\sum_W\|v_i\|^2=L_W/4$
in the matched balanced planar limit, and the second term is controlled by
Cauchy--Schwarz.  Ignoring only this displayed mismatch remainder, the frozen
sizing therefore gives the leading estimate
\[
 \widehat\chi_W(t_{\rm create})
 \le\frac{L_W(t_{\rm create})}{4\epsilon_\sigma}
 \approx5.82\times10^{-3}.
\]
This sharpening is not a replacement for the exact remainder bound.  In
either form weak rotation is not the apparent canonical pre-creation
obstruction.  The quantitatively severe issue in this route is
the additive spread floor \eqref{sbcoh:additive-floor-canonical} and the
corresponding directional mismatch/sampling transfer.
\end{remark}

\begin{remark}[Status and fallback to the $c_{21}^O$ channel]
\label{sbcoh:status}
The population branch and the signed weak-cross equation identify a coherent
$c_{12}^O$ route, but a full finite-sample theorem still needs
\eqref{sbcoh:field-transfer} with an additive error small enough for
\eqref{sbcoh:additive-floor-canonical}, together with the directional
$\delta_2$ control in Remark~\ref{sbcoh:directional-mismatch}.  If those
requirements are too expensive, the direct signed-gap certificate should use
the $c_{21}^O$ channel as the load-bearing post-B1 feedback and treat
$c_{12}^O$ as an optional strengthening under the stronger coherence regime.
This fallback does not invalidate the B2 holding proposition or the regulator
first-hitting creation interface. The coupled-ratio continuation in
Section~\ref{sbpos:scope} supplies an alternative feedback certificate,
and Theorem~\ref{sbpos:aggregate-mean} proves a population holding
region using the weighted output-tilt mean. Its initialized empirical
entry and transfer remain open.
\end{remark}


\section{Signed transverse propagation and direct probe dynamics}
\label{sbsign:scope}
This addition proves an exact signed energy identity for the empirical
network, a relative-growth comparison, an aggregate transfer to actual
probe gates, and a complete fixed-gate reversal theorem in a noiseless
one-neuron subcase. The general noisy, isotropically initialized theorem
remains open. All notation is local to this section.

Samples lie in $\mathcal S=\operatorname{span}\{e_1,e_2\}$, with equal
cluster weights and $\mathbb E_n=(\mathbb E_{1,n}+\mathbb E_{2,n})/2$.
Set $\lambda_n=\|\mathbb E_nxx^\top\|_{\rm op}$ and
$\|g\|_n^2=\mathbb E_n\|g\|^2$. The loss is
$\mathcal L=\frac12\|x-f\|_n^2$, where
$f(x)=\sum_iw_i(u_i^\top x)_+$.
Write $v_i=P_{\mathcal S}u_i$, $z_i=P_{\mathcal S}w_i$, $e=x-f_\parallel$.
Compatible selectors satisfy $0\le\alpha_i\le1$ and
$\alpha_i v_i^\top x=(v_i^\top x)_+$, and give
\[
 A_i=\mathbb E_n[e x^\top\alpha_i],\qquad
 \dot z_i=A_iv_i,\qquad \dot v_i=A_i^\top z_i-P_i,
 \qquad
 P_i=\mathbb E_n[x\alpha_i\langle w_{i,\perp},f_\perp\rangle].
\]
Fix $[h]=J_1\sqcup J_2$ and let $p(i)$ be the assigned axis and
$q(i)=3-p(i)$ its transverse axis. Define
\[
 T=\sum_i(v_{i,q(i)}^2+z_{i,q(i)}^2),\quad Y=\sqrt T,\quad
 \Pi=\Big(\sum_i\|P_i\|^2\Big)^{1/2},\quad
 S=\sum_i\|u_i\|^2=\sum_i\|w_i\|^2.
\]
Differential statements hold almost everywhere on the absolutely
continuous selected trajectory; norm comparisons are regularized at zero.

\subsection{The off-diagonal feedback has an exact negative term}
\begin{lemma}[Signed transverse-energy identity]
\label{sbsign:energy}
Define the diagonal and off-diagonal parts of the planar forward map by
\[
 f_D(x)=\sum_i\alpha_i(x)
   (z_{i1}v_{i1}x_1e_1+z_{i2}v_{i2}x_2e_2),
\]
\[
 f_O(x)=\sum_i\alpha_i(x)
   (z_{i1}v_{i2}x_2e_1+z_{i2}v_{i1}x_1e_2),
 \qquad f_\parallel=f_D+f_O,
\]
and let $f_{TT}(x)=\sum_i\alpha_i z_{i,q(i)}v_{i,q(i)}x_{q(i)}e_{q(i)}$.
Then exactly
\begin{equation}\label{sbsign:identity}
 \frac12\dot T=-\|f_O\|_n^2+langle x-f_D,f_O\rangle_n
       +2\sum_i A_{i,q(i)q(i)}z_{i,q(i)}v_{i,q(i)}
       -\sum_i v_{i,q(i)}P_{i,q(i)}.
\end{equation}
Put $m_\times=\mathbb E_n|x_1x_2|$. For $T>0$ define the signed rate
\[
 \kappa(t)=\frac{2\sum_i A_{i,q(i)q(i)}z_{i,q(i)}v_{i,q(i)}
                         -\|f_O\|_n^2}{T},
\]
and set $\kappa=0$ when $T=0$. Then
\begin{equation}\label{sbsign:Y}
 \dot Y\le\kappa Y+\sqrt{2S}(1+S)m_\times+\Pi.
\end{equation}
In particular no separate cross-residual differential comparison is needed.
If $\|A_i\|\le K$, then $|\kappa|\le K+2\lambda_nS$.

For $\nu_\times^2=(s_{22}^{(1)}+s_{11}^{(2)})/2$,
the wrong-coordinate residual satisfies the static estimate
\begin{equation}\label{sbsign:X}
 X:=\|P_\times e\|_n
 \le(1+S)\nu_\times+\sqrt{2\lambda_n S}\,Y.
\end{equation}
\end{lemma}
\begin{proof}
Differentiate $T/2$ using the projected equations. The output-coordinate
and input-coordinate terms together give
$\langle e,f_O+2f_{TT}\rangle_n$; the input equation also contributes
$-\sum_i v_{i,q(i)}P_{i,q(i)}$.
Substitute $e=x-f_D-f_O$ in the first inner product to get
\eqref{sbsign:identity}.

Write $f_D=(d_1(x)x_1,d_2(x)x_2)$ and
$f_O=(o_{12}(x)x_2,o_{21}(x)x_1)$. Balance and Cauchy--Schwarz give
$|d_r(x)|\le S$. Moreover, writing
$Q_{12}=\sum_i|z_{i1}v_{i2}|$ and $Q_{21}=\sum_i|z_{i2}v_{i1}|$,
\[
 Q_{12}+Q_{21}\le
 \Big(\sum_i(v_{i,p(i)}^2+z_{i,p(i)}^2)\Big)^{1/2}\sqrt T
 \le\sqrt{2ST}.
\]
Thus $|\langle x-f_D,f_O\rangle_n|
\le(1+S)m_\times\sqrt{2ST}$.
The perpendicular term is bounded by $\sqrt T\Pi$. Dividing by $Y$
proves the comparison, with regularization at zero.
Also $\|f_O\|_n^2\le\lambda_n(Q_{12}^2+Q_{21}^2)\le2\lambda_nST$,
which proves the bound on $\kappa$.
Finally $\|P_\times(x-f_D)\|_n\le(1+S)\nu_\times$; use the last bound
on $f_O$ to obtain \eqref{sbsign:X}.
\end{proof}

\begin{remark}[What remains inside the signed rate]
For $t_i=v_{i,q(i)}^2+z_{i,q(i)}^2$ and
$\delta_{i,q(i)}=v_{i,q(i)}-z_{i,q(i)}$,
\[
 \kappa T=\sum_i A_{i,q(i)q(i)}t_i
       -\sum_i A_{i,q(i)q(i)}\delta_{i,q(i)}^2-\|f_O\|_n^2.
\]
One cannot replace the diagonal residuals by their positive parts without
also paying their negative parts against the mismatch squares.
This identity isolates the remaining signed polarization estimate;
it does not assert that its rate is favorable for every partition.
\end{remark}

\begin{corollary}[Comparison relative to mass growth]
\label{sbsign:clock}
Suppose $S>0$ and write $g(t)=\dot S/(2S)$ and $r=Y/\sqrt S$.
Then
\[
 \dot r\le(\kappa-g)r+\sqrt2(1+S)m_\times+\Pi/\sqrt S.
\]
If on $[a,b]$, $S\le S_*$, $g-\kappa\ge\gamma>0$, and
$\Pi/\sqrt S\le p_*$, then
\begin{equation}\label{sbsign:relative}
 Y(t)\le\sqrt{S_*}\left[
 e^{-\gamma(t-a)}\frac{Y(a)}{\sqrt{S(a)}}
 +\frac{\sqrt2(1+S_*)m_\times+p_*}{\gamma}
           (1-e^{-\gamma(t-a)})\right].
\end{equation}
The existing perpendicular bound allows $p_*=\lambda_nh\varepsilon^2$.
Unlike the earlier bound, its forcing does not acquire a factor
exponential in the learning-window length. The signed gap
$g-\kappa\ge\gamma$ is an additional dynamical obligation, not a
consequence of the global norm ceiling.
\end{corollary}
\begin{proof}
Differentiate $Y/\sqrt S$ and apply \eqref{sbsign:Y}. Integrate the scalar
inequality. The perpendicular estimate is
$\Pi\le\lambda_nh\varepsilon^2\sqrt S$.
\end{proof}

\begin{corollary}[A spectral-gap certificate in the noiseless planar model]
\label{sbsign:spectral}
Take the two planar training samples $\mu_1e_1,\mu_2e_2$, set
$\lambda_j=\mu_j^2/2$, and choose
$J_2=\{i:v_{i1}(0)<0\}$, $J_1=[h]\setminus J_2$.
Then every member of $J_2$ remains inactive on cluster 1.
Write $m_{rj}=f_r(\mu_je_j)/\mu_j$.
At times where $0<T\le S$, $0\le m_{22}\le1$,
$c_*S\le m_{11}\le\bar m<1$, one has
\[
 g-\kappa\ge\lambda_1c_*(1-\bar m)-\lambda_2.
\]
Thus $\lambda_1c_*(1-\bar m)>\lambda_2$ is a sufficient signed-gap
condition. The fit-to-mass bound $m_{11}\ge c_*S$ and the other displayed
conditions must still be proved along the chosen trajectory.
\end{corollary}
\begin{proof}
If $v_{i1}<0$, then its cluster-1 gate is off and cluster 2 has $x_1=0$,
so $\dot v_{i1}=0$. Hence its initially negative coordinate stays fixed.
Put $E_O=\lambda_1m_{21}^2+\lambda_2m_{12}^2$ and
$L=\sum_{i\in J_1}\alpha_i(\mu_2e_2)z_{i2}v_{i2}$.
The mass identity and Lemma~\ref{sbsign:energy} give
\[
 g=\frac{\lambda_1m_{11}(1-m_{11})+
               \lambda_2m_{22}(1-m_{22})-E_O}{S},
 \qquad
 \kappa T=2\lambda_2(1-m_{22})L-E_O.
\]
Since $2|L|\le T$, subtract these formulas and bound from below by
\[
 \frac{\lambda_1m_{11}(1-m_{11})+
                  \lambda_2m_{22}(1-m_{22})}{S}
       -\lambda_2(1-m_{22})+E_O(1/T-1/S).
\]
The last term is nonnegative. Discard the nonnegative weak progress
term and use the bounds on $m_{11}$.
\end{proof}

\subsection{Actual probe gates can be controlled by aggregate input signs}
\begin{lemma}[Negative input mass and its integral certificate]
\label{sbsign:negative}
Fix $q\in\{1,2\}$, let $p=3-q$, and assume every cluster-$q$ sample has
$x_q>0$. Put
\[
 N_q=\sum_{i\in J_q}(v_{iq})_-^2,\quad
 Q_q=\sum_{i\in J_q}v_{ip}^2,\quad
 R_j=\|e\|_{j,n},\quad
 b_p=R_p\sqrt{s_{qq}^{(p)}S},\quad b_q=R_q\sqrt{s_{pp}^{(q)}S}.
\]
Then
\[
 \dot N_q\le(b_p+2\Pi)\sqrt{N_q}+b_q\sqrt{Q_q}.
\]
Consequently, setting
$A(t)=\int_a^t b_q(s)\sqrt{Q_q(s)}\,ds$ and
$B(t)=\int_a^t(b_p(s)+2\Pi(s))\,ds$,
\begin{equation}\label{sbsign:negative-integral}
 \sqrt{N_q(t)}\le\sqrt{N_q(a)+A(t)}+B(t)/2.
\end{equation}
At the stipulated initialization $N_q(0)\le h\varepsilon^2$.
If $\sqrt{Q_q}\le C\eta$, $\int b_p,\int b_q,\int\Pi=O(\eta)$,
and $N_q(a)=O(\eta^2)$, this proves $N_q=O(\eta^2)$.
For example, $g=\dot S/(2S)\ge g_0>0$ implies
$\int_a^t\sqrt{S(s)}\,ds\le\sqrt{S(t)}/g_0$, making these integral
conditions explicit in the noise moments and perpendicular budget.
\end{lemma}
\begin{proof}
Differentiate the negative-part square. On cluster $p$ use
Cauchy--Schwarz over neurons and samples. On cluster $q$, whenever
$v_{iq}<0$ and $\alpha_i>0$, the gate condition implies
$|v_{iq}|x_q\le|v_{ip}||x_p|$. This also holds for a selected zero gate
boundary. The factor $2$ from differentiating cancels the mixture factor
$1/2$. The two clusters therefore give respectively $b_p\sqrt{N_q}$
and $b_q\sqrt{Q_q}$. The perpendicular contribution is at most
$2\sqrt{N_q}\Pi$.
The square of $\sqrt{N_q(a)+A(t)}+B(t)/2$ is a supersolution;
regularize its first square root if it vanishes.
The mass integral follows by integrating
$\sqrt{S(s)}\le\sqrt{S(t)}e^{-g_0(t-s)}$.
\end{proof}

\begin{lemma}[Grading of the actual active set near a negative concept axis]
\label{sbsign:probe}
Fix $p$, put $q=3-p$, and consider a unit probe
$\xi=\eta z e_p-k e_q$, where $|z|\le Z$,
$k=\sqrt{1-\eta^2z^2}\ge k_0>0$.
At each time let $\omega_i=\mathbf1_{v_i^\top\xi>0}$ and
$\beta=\eta Z/k_0$. Write
$T_p=\sum_{J_p}(v_{iq}^2+z_{iq}^2)$ and
$T_q=\sum_{J_q}(v_{ip}^2+z_{ip}^2)$.
Neuronwise balance and perpendicular input invariance imply
\begin{equation}\label{sbsign:probe-energy}
 T_\omega^{(p)}:=\sum_i\omega_i(v_{iq}^2+z_{iq}^2)
 \le T_p+(1+2\beta^2)T_q+2N_q+h\varepsilon^2.
\end{equation}
Hence aggregate transverse energy and negative input mass of order
$\eta^2$ imply the actual probe matrix's normal-form grading.
No agreement of all neurons with a common probe gate is assumed.
\end{lemma}
\begin{proof}
For an active member of $J_q$ with $v_{iq}\ge0$, activation gives
$v_{iq}\le\beta|v_{ip}|$. Members with $v_{iq}<0$ contribute to $N_q$.
Thus $\sum_{J_q}\omega_i v_{iq}^2\le\beta^2Q_q+N_q$.
Balance gives $z_{iq}^2\le\|w_i\|^2=\|u_i\|^2
=v_{ip}^2+v_{iq}^2+\|u_{i,\perp}(0)\|^2$.
Sum, add the $J_p$ contribution, and use
$\sum_i\|u_{i,\perp}(0)\|^2\le h\varepsilon^2$.
For $M_\omega=\sum_i\omega_i z_iv_i^\top$, the resulting bounds are
$|M_{pp}|\le S$, $|M_{pq}|,|M_{qp}|\le\sqrt{S T_\omega^{(p)}}$,
and $|M_{qq}|\le T_\omega^{(p)}/2$.
\end{proof}

\begin{proposition}[Pointwise error expansion through probe gate changes]
\label{sbsign:switching}
Fix one probe $\xi$ as in Lemma~\ref{sbsign:probe}. At each time form
the actual active matrix $M_\omega$ and its frozen-selector rate
\[
 M_\omega^\dagger=\sum_i\omega_i
 (A_iv_iv_i^\top+z_iz_i^\top A_i-z_iP_i^\top).
\]
Write these, in the ordered basis $(e_p,e_q)$, as
\[
 M_\omega=\begin{pmatrix}a&\eta b\\\eta c&\eta^2d\end{pmatrix},\qquad
 M_\omega^\dagger=\begin{pmatrix}a^\dagger&\eta b^\dagger\\
 \eta c^\dagger&\eta^2d^\dagger\end{pmatrix}.
\]
Assume these eight coefficients are uniformly bounded. Then the direct
normal form holds pointwise and its derivative holds almost everywhere:
\[
 E(\xi,t)=\tfrac12+\eta^2H+O(\eta^4)+\tfrac12\|P_\omega\xi\|^2,
\]
\[
 \dot E(\xi,t)=\eta^2G+O(\eta^4)
                +\langle P_\omega\xi,P_\omega^\dagger\xi\rangle,
\]
where $A=(a-1)z-b$ and
\[
 H=\tfrac12 A^2+cz-d-\tfrac12z^2,\qquad
 G=A(a^\dagger z-b^\dagger)+zc^\dagger-d^\dagger.
\]
The remainder constants depend only on the coefficient bounds, $Z,k_0$.
The rates have the required bounds if \eqref{sbsign:probe-energy} is
$O(\eta^2)$, $S,K$ are bounded, and
$\mathcal F_\omega^{(p)}+\Pi=O(\eta)$.
The Jensen bound and \eqref{sbsign:X} supply the latter from
$Y=O(\eta)$, $\nu_\times=O(\eta)$ and the perpendicular budget.

This is a pointwise-in-probe statement: its coefficients may depend on
$z$ through the active set. It does not assert a common polynomial profile
across different activation cells, or that these coefficients are globally
absolutely continuous. Strict derivative sign margins imply monotonicity
of $E$ even across probe gate changes.
\end{proposition}
\begin{proof}
Let $P_\omega=\sum_i\omega_iw_{i,\perp}v_i^\top$ and define its dagger
rate by differentiating the factors with the same frozen selectors.
For an absolutely continuous scalar $q(t)=v_i(t)^\top\xi$,
$\frac{d}{dt}q_+=\mathbf1_{q>0}\dot q$ almost everywhere, including
the fact that $\dot q=0$ almost everywhere on $\{q=0\}$.
Therefore $f_\parallel(\xi)=M_\omega\xi$ and
$\dot f_\parallel(\xi)=M_\omega^\dagger\xi$ almost everywhere, with the
analogous ambient identity. Squaring the residual gives the expansion.
Explicitly, with $B=cz-dk$, its exact fourth-order term is
\[
 K_E=dz^2+\frac{z^2(bA-cz)}{1+k}
             +\tfrac12 B^2+\frac{\eta^2b^2z^4}{2(1+k)^2}.
\]
Substitution of dagger rates in this polynomial gives the exact
$K_E^\dagger$ in the derivative identity. Bounded coefficients and rates
bound both expressions, without differentiating $\omega_i$.
For the rate bound, write $Y_\omega=\sqrt{T_\omega^{(p)}}$ and
$F=\mathcal F_\omega^{(p)}$. Directly from the displayed matrix rate,
\[
 |(M_\omega^\dagger)_{qq}|
 \le K Y_\omega^2+Y_\omega(F+\Pi),
\]
\[
 |(M_\omega^\dagger)_{pq}|,|(M_\omega^\dagger)_{qp}|
 \le K\sqrt{2S}Y_\omega+K Y_\omega^2+\sqrt S(F+\Pi),
\]
and $|(M_\omega^\dagger)_{pp}|\le2KS+\sqrt S\Pi$.
Finally $E$ itself is absolutely continuous, so almost-everywhere strict
sign bounds integrate to strict monotonicity.
\end{proof}

\subsection{A complete noiseless fixed-gate reversal theorem}
\begin{theorem}[Fixed-gate swing-by in a balanced pure-sector subcase]
\label{sbsign:exact-reversal}
Take $d=2$, $h=1$, noiseless training samples $\mu_1 e_1,\mu_2 e_2$
with equal weights, and initialization
\[
 u(0)=w(0)=(a_0,-\eta),\qquad
 a_0,\eta>0,\quad s_0=a_0^2+\eta^2\le1/16.
\]
Put $\lambda=\mu_1^2/2$, $k=\sqrt{1-\eta^2}$, and choose the fixed
unit probe $\xi=(\eta,-k)$. Then both training gates and this probe gate
are fixed for all finite times. The strong training reconstruction
converges to $\mu_1e_1$. The OOD error $E(\xi,t)$ has negative derivative
initially and positive derivative for all sufficiently large times.
In particular it has an interior strict local minimum.

On the entire trajectory the actual probe matrix has uniformly bounded
normal-form coefficients and rates at scale $\eta$, independently of
the time needed to learn. Its transverse input/output energy is at most
$2\eta^2$. This is a deterministic subcase, not a high-probability
statement under isotropic initialization, and it does not fit concept 2.
\end{theorem}
\begin{proof}
Write $u=(a,-\eta)$ and $w=(c,-\eta D)$. The second training gate is
off and the first is on as long as $a>0$. The exact equations are
\begin{equation}\label{sbsign:one-neuron-ode}
 \dot a=\lambda[c-a(a^2+\eta^2)],\qquad
 \dot c=\lambda a(1-ac),\qquad
 \dot D=-\lambda a^2D,
 \qquad a(0)=c(0)=a_0,\quad D(0)=1.
\end{equation}
Balance gives $c^2+\eta^2D^2=a^2+\eta^2$.
Let $A>0$ be the unique root
\[
 A^2(A^2+\eta^2)=1,\qquad
 A^2=\frac{\sqrt{\eta^4+4}-\eta^2}{2},\qquad C=1/A.
\]
One has $a_0<A<1$. Put $h=c-a(a^2+\eta^2)$.
Initially $h=a_0(1-s_0)>0$. On a boundary $h=0$, $0<a<A$,
\[
 \dot h=\lambda a[1-a^2(a^2+\eta^2)]>0.
\]
Thus $\dot a>0$ while $a<A$. At $a=A$, balance and $D>0$ imply
$c<\sqrt{A^2+\eta^2}=C$, hence $h<0$, so this boundary cannot be
reached from the preceding region at finite time. Also
$ac\le a\sqrt{a^2+\eta^2}<1$, so $c$ increases.
Since $a\ge a_0$, $D$ decreases to zero. Bounded monotone convergence
and \eqref{sbsign:one-neuron-ode} imply
\[
 a\longrightarrow A,\qquad c\longrightarrow C,\qquad D\longrightarrow0.
\]
This proves the gate claims, because $a>0$, the second input coordinate
is the negative constant $-\eta$, and
$u^\top\xi=\eta(a+k)>0$.

The actual probe matrix is
\begin{equation}\label{sbsign:exact-matrix}
 M=\begin{pmatrix}ac&-\eta c\\-\eta aD&\eta^2D\end{pmatrix}.
\end{equation}
Here $a\le1$, $c\le\sqrt{1+\eta^2}\le2$, $D\le1$,
$0<\dot a\le2\lambda$, $0<\dot c\le\lambda$,
and $|\dot D|\le\lambda$. The four scaled coefficients are bounded
by $2$, and their rates by $5\lambda$. Also
$T=\eta^2(1+D^2)\le2\eta^2$ and
$\dot T=-2\lambda\eta^2a^2D^2$.
This is the negative off-diagonal term of \eqref{sbsign:identity},
with zero diagonal transverse residual and zero data noise.

For the chosen probe, exactly
\begin{equation}\label{sbsign:exact-error}
 E=\frac12\left\{
 \eta^2[c(a+k)-1]^2+[k-\eta^2D(a+k)]^2\right\}.
\end{equation}
At initialization put $q_0=\eta(a_0+k)>0$. Direct differentiation gives
\[
 \dot E(0)=\lambda\eta a_0q_0
 \left[a_0(a_0+k)(1-s_0)-(1-a_0^2)
                      +a_0k-(1-s_0)^2\right].
\]
The bracket is at most
$2a_0^2+2a_0-2+2s_0\le-5/4$, so $\dot E(0)<0$.

For completeness the late-time dominance used next follows directly
from a two-variable stable scalar comparison. Eliminate $c$ by balance:
\[
 \dot a=F(a,D):=\lambda[sqrt{a^2+\eta^2(1-D^2)}-a(a^2+\eta^2)].
\]
At $(A,0)$, $F_a=-\lambda(2A^2+\eta^2)$, $F_D=0$, while
$\dot D/D\to-\lambda A^2$. Choose rates strictly between
$\lambda A^2$ and $\min\{\lambda(2A^2+\eta^2),2\lambda A^2\}$.
Taylor's formula and scalar comparison for $A-a$ show
$A-a=o(D)$: explicitly, for arbitrarily small $\epsilon>0$ at late times,
\[
 \frac{d}{dt}(A-a)\le-[\lambda(2A^2+\eta^2)-\epsilon](A-a)+C_\eta D^2,
\]
and $D$ is bounded above and below by exponentials with rates
$\lambda A^2\mp\epsilon$. Integrating this inequality proves the claim.
It also gives $\dot a=o(D)$, $\dot c=o(D)$.
Differentiating \eqref{sbsign:exact-error} now yields
\[
 \lim_{t\to\infty}\frac{\dot E}{\eta^2D}
       =\lambda A^2 k(A+k)>0.
\]
Thus the derivative is positive eventually. Before that time $E$ is
analytic and initially strictly decreasing. A minimum on a sufficiently
large compact time interval is interior; analyticity and nonconstancy
give a strict local minimum. No remainder estimate is used to assert
the sign reversal.
\end{proof}

\begin{corollary}[A second probe has a peak before its gate turns off]
\label{sbsign:switch-peak}
Under the same model and initialization, choose
$Z\in(A,1/\sqrt{s_0})$ and put
$k_Z=\sqrt{1-\eta^2Z^2}$, $\xi_Z=(-\eta Z,-k_Z)$.
The probe is initially active and turns off at the unique finite time
when $a=k_Z/Z<A$. Its error has a strict interior local maximum before
that switch, exceeds $1/2$ near the switch from the left, and equals
$1/2$ thereafter.
\end{corollary}
\begin{proof}
The preactivation is $q=\eta(k_Z-Za)$; the stated range of $Z$ is
equivalent to $a_0<k_Z/Z<A$. Since $a$ increases to $A$, the switch
exists and is unique. To determine the error near it, let $H=c-aD$.
Initially $H=0$, and on $H=0$ the equations give
\[
 \dot H=\lambda a[1-D^2+(a^2+\eta^2)D]>0.
\]
Thus $H>0$ at every positive finite time.
At the switch $w^\top\xi_Z=-\eta Z(c-aD)<0$ and
$\dot q=-\eta Z\dot a<0$. Therefore
$\dot E(\tau^-)=-\langle w,\xi_Z\rangle\dot q<0$.
Just before the switch,
$E-1/2=\|w\|^2q^2/2-q\langle w,\xi_Z\rangle>0$.
Initially $w=u$, so
$E(0)=1/2-(1-s_0/2)q(0)^2<1/2$.
The maximum on the pre-switch compact interval is therefore interior
and, by analyticity, a strict local maximum. After the switch $q_+=0$.
\end{proof}

\subsection{The remaining general-model step}
The signed identity and actual-gate transfer hold for the true noisy
empirical equations. The complete reversal theorem above is restricted
to one balanced neuron in a noiseless pure sector; it is not the intended
isotropic many-neuron theorem and does not explain all empirical peaks.
For the general theorem, the key remaining propagation task is now to
prove a useful bound on $g-\kappa$ (or directly integrate the signed
identity) on a learning interval reached from initialization. The negative
input-mass integral certificate then provides a route to probe grading
without per-neuron gate freezing. The actual noisy coefficient evolution
still requires its own reversal sign margins. The fixed-gate subcase
demonstrates that this mechanism exists; it does not supply those margins
for the general dataset by analogy.

\section{Proof of the sector theorem}
\label{sbsector:proof}

\subsection{Exact gradient flow and invariant region}
As long as the first training gate is on and the second is off,
the exact equations, starting with $a(0)=c(0)=a_0$, $D(0)=1$, are
\begin{equation}\label{sbsector:ode}
 \dot a=\lambda[c-a(a^2+\eta^2)],\qquad
 \dot c=\lambda a(1-ac),\qquad
 \dot D=-\lambda a^2D.
\end{equation}
For example, the strong-sample output is $\mu_1a(c,-\eta D)$;
the input-gradient equation and the balance identity below give the
first equation. The second input coordinate is the constant $-\eta$.
In the reduced equations the derivative of
$c^2+\eta^2D^2-a^2-\eta^2$ is $-2\lambda a^2$ times that quantity.
It vanishes initially, giving balance:
\begin{equation}\label{sbsector:balance}
 c^2+\eta^2D^2=a^2+\eta^2.
\end{equation}
Let $a_\infty>0$ be defined by
\[
 a_\infty^2(a_\infty^2+\eta^2)=1.
\]
Then $a_0<a_\infty<1$ and $a_\infty>1/\sqrt2$.
Set $h=c-a(a^2+\eta^2)$. Initially $h=a_0(1-a_0^2-\eta^2)>0$.
On $h=0$ with $0<a<a_\infty$,
\[
 \dot h=\lambda a[1-a^2(a^2+\eta^2)]>0.
\]
Thus $h>0$ while $a<a_\infty$. At any hypothetical first finite hit
of $a_\infty$, $D>0$ and balance would give
$c<\sqrt{a_\infty^2+\eta^2}=1/a_\infty$, hence $h<0$, a contradiction.
Consequently
\[
 0<a_0\le a<a_\infty<1,\quad
 a\le c\le\sqrt{a^2+\eta^2},\quad ac<1,\quad
 \dot a>0,\quad\dot c>0,\quad 0<D\le1.
\]
These bounded quantities extend globally; $a>0$ keeps the first gate on,
and the constant negative second input coordinate keeps the second off.
Since $a\ge a_0$, $D\to0$. Monotonicity, balance, and
\eqref{sbsector:ode} imply $a\to a_\infty$ and $c\to1/a_\infty$.

For every probe under consideration,
$u^\top\xi_z=\eta(az+k_z)>0$. Thus its exact matrix is the matrix in
the theorem. The bounds $a\le1$, $c\le2$, $D\le1$,
$\dot a\le2\lambda$, $\dot c\le\lambda$, and
$|\dot D|\le\lambda$ prove the coefficient and rate bounds by the
product rule. Transverse energy equals $\eta^2(1+D^2)$.

\subsection{Exact error identity and sign estimates}
For this proof abbreviate $k=k_z$ and set
\[
 \mathcal A=c(az+k)-z,
 \qquad \mathcal T=-D(az+k),
 \qquad B=k+\eta^2\mathcal T.
\]
There is no ambient output component. Exactly,
\begin{equation}\label{sbsector:error}
 E=\frac{\eta^2}{2}\mathcal A^2+\frac12B^2,
 \qquad
 \frac{\dot E}{\eta^2}
 =\mathcal A\dot{\mathcal A}+B\dot{\mathcal T}.
\end{equation}
Throughout the trajectory,
\begin{equation}\label{sbsector:basic}
 k\ge\frac{\sqrt{15}}4>\frac34,\qquad
 B\ge k-\eta^2(1+k)\ge\frac58,
 \qquad\dot{\mathcal A}>0.
\end{equation}
Here $az+k\le2$, $\eta^2\le1/16$, and
$\dot{\mathcal A}=\dot c(az+k)+cz\dot a$.
The leakage derivative is exactly
\begin{equation}\label{sbsector:leakage}
 \dot{\mathcal T}
 =\lambda D\{ka^2+z(2a^3+a\eta^2-c)\}.
\end{equation}

If $a\le1/4$, then
$c\le\sqrt{1/8}<3/8$ and $ac\le3/32$. Hence
\[
 \mathcal A=ck-(1-ac)z\le\frac38-\frac{29}{64}
 =-\frac5{64}.
\]
Using $c\ge a$, $z\ge1/2$, and $\eta^2\le1/16$ in
\eqref{sbsector:leakage} also gives
\[
 \dot{\mathcal T}
 \le\lambda Da\{a-z(1-2a^2-\eta^2)\}
 \le-\frac5{32}\lambda Da.
\]
Thus, on $[0,t_-]$,
\begin{equation}\label{sbsector:early-sign}
 \dot E\le-\frac{25}{256}\lambda\eta^2Da<0.
\end{equation}

For $a\ge1/\sqrt2$, the inequality
$c\le a+\eta^2/(2a)$ yields
\[
 2a^3+a\eta^2-c
 \ge(2a^2-1)\left(a+\frac{\eta^2}{2a}\right)\ge0.
\]
Therefore
\[
 \dot{\mathcal T}\ge\lambda Dka^2,
 \qquad
 \mathcal A\ge ak-(1-a^2)
 \ge\sqrt{15/32}-1/2>0.
\]
Equations \eqref{sbsector:error}--\eqref{sbsector:basic} imply, for
every $t\ge t_+$,
\begin{equation}\label{sbsector:late-sign}
 \dot E\ge\frac{15}{32}\lambda\eta^2a^2D>0.
\end{equation}
Unlike an asymptotic comparison, this estimate applies from an explicit
progress level and retains the natural decaying rate $a^2D$.

\subsection{Reached times and a uniform rebound amplitude}
The invariant-region proof guarantees that $t_-$ and $t_+$ are finite.
Since $(a+c)'\le\lambda(a+c)$ and $c\ge a$,
\[
 a(t)\le a_0e^{\lambda t},
 \qquad t_-\ge\lambda^{-1}\log(1/(4a_0)).
\]
In the other direction, $c\ge a$ gives
\[
 \dot a\ge\lambda a(\beta-a^2),\qquad \beta=1-\eta^2.
\]
Comparison with the scalar logistic equation for $a^2$ through $t_+$
gives the upper bound in \eqref{sbsector:time}.

The error is analytic at every finite time because the relevant gates
are fixed and the equations are analytic. It decreases through $t_-$
and increases from $t_+$ onwards. Its global minimum is therefore attained
in $(t_-,t_+)$. A nonconstant real analytic function has isolated local
minima, so every global minimizing time is a strict local minimum.

To bound $D(t_+)$ from below, integrate with respect to the increasing
variable $a$ up to $1/\sqrt2$:
\[
 \frac{d\log D}{da}=-\frac{\lambda a^2}{\dot a}
 \ge-\frac{a}{\beta-a^2},
 \qquad
 D(t_+)\ge
 \sqrt{\frac{\beta-1/2}{\beta-a_0^2}}
 \ge\frac{\sqrt7}{4}.
\]
Integrating \eqref{sbsector:late-sign}, using
$-\dot D=\lambda a^2D$, proves
\[
 E_\infty-E(t_+)
 \ge\frac{15}{32}\eta^2D(t_+)
 \ge\frac{15\sqrt7}{128}\eta^2>\frac14\eta^2.
\]
This also bounds the rebound from the smaller global minimum below.
For the upper bound, expand \eqref{sbsector:error}:
\[
 E-\frac12
 =\eta^2\left(\frac{\mathcal A^2}{2}-\frac{z^2}{2}
                         +k\mathcal T\right)
   +\frac{\eta^4}{2}\mathcal T^2\ge-\frac52\eta^2.
\]
At infinity, $ac\to1$ and $D\to0$, so
\[
 E_\infty=\frac{k^2}{2}(1+\eta^2/a_\infty^2),
 \qquad
 E_\infty-\frac12\le\frac{17}{32}\eta^2,
\]
where $a_\infty^{-2}=a_\infty^2+\eta^2\le17/16$.
Thus the rebound is at most $(97/32)\eta^2<4\eta^2$.
Finally, integrating the derivative of $\arcsin$ on
$[\eta/2,\eta]$, with $\eta\le1/4$, proves the sector-width bounds.
This completes the proof of Theorem~\ref{sbsector:theorem}.

\subsection{Proof of the two-concept extension}
As long as the two neurons keep their pure training gates, the first
sees only $\mu_1e_1$ and the second only $\mu_2e_2$. Their equations
are two independent copies of \eqref{sbsector:ode}, with the coordinates
swapped for the second neuron. The first has input $(a,-\eta)$ and
the second $(-\zeta,b)$ with $a,b>0$ and constant $\eta,\zeta>0$;
the preceding invariant-region proof therefore preserves those gates
for all finite time. Each training output converges to its own target.
At a strong probe the second preactivation is
$-\zeta\eta z-bk_z<0$; at a weak probe the first is likewise negative.
Each probe thus has exactly the one-neuron error already analyzed.
This proves Corollary~\ref{sbsector:two-concept} without a neglected
complement contribution.

For Corollary~\ref{sbsector:isotropic}, two prescribed quadrants have
probability $(1/4)^2$ in a prescribed neuron order. The two possible
orders are disjoint, giving $1/8$. For the smaller event restrict the
quadrant-IV angle to $[-\pi/3,-\pi/6]$ and the quadrant-II angle to
$[2\pi/3,5\pi/6]$. Each arc has probability $1/12$, so the two orders
have total probability $1/72$. On these arcs all four relevant coordinate
magnitudes lie in $[\varepsilon/2,\sqrt3\varepsilon/2]$.
Substitution into the width and amplitude bounds gives the claims.
In the upper time bound use
$\beta\ge15/16$, $2\beta-1\ge7/8$, and
$a_0^2\ge\varepsilon^2/4$ (or $b_0^2$ for the other concept).
The lower bound uses $a_0,b_0\le\sqrt3\varepsilon/2$.
This proves \eqref{sbsector:isotropic-time}.

\section{Closing the direct-probe obligations in this regime}
\label{sbsector:obligations}

\begin{proposition}[A reached crossing and a late leakage margin]
\label{sbsector:crossing}
Under Theorem~\ref{sbsector:theorem}, each $\mathcal A_z$ has a unique
zero in $(t_-,t_+)$. On $[0,t_+]$,
\[
 \dot{\mathcal A}_z\ge\lambda a_0/3.
\]
On any compact interval $[t_+,H]$, $H>t_+$, one has
\[
 \dot{\mathcal T}_z\ge\frac38\lambda e^{-\lambda H}>0,
 \qquad
 \dot E\ge\frac{15}{64}\lambda\eta^2e^{-\lambda H}>0.
\]
These estimates are uniform over $z\in[1/2,1]$.
\end{proposition}
\begin{proof}
Before $t_+$, $ac\le\sqrt{(1/2)(1/2+1/16)}=3/(4\sqrt2)$.
Thus
\[
 \dot{\mathcal A}\ge k\dot c
 \ge\frac34\lambda a_0\left(1-\frac3{4\sqrt2}\right)
 >\lambda a_0/3.
\]
The signs of $\mathcal A$ at $t_-$ and $t_+$ and its strict
monotonicity prove the crossing claim. Since $a^2\le1$,
$D(t)\ge e^{-\lambda t}$. Insert this, $a^2\ge1/2$, and
$k\ge3/4$ into the late estimates.
\end{proof}

The error minimum need not coincide with the zero of $\mathcal A$.
We have not assumed that the retained leakage term is uniformly small
near that zero. Instead both terms in \eqref{sbsector:error} have
favorable signs on the certified early and late intervals. This closes
the reversal proof without an artificial small-remainder hypothesis.

\begin{proposition}[Uniform normal form and derivative accuracy]
\label{sbsector:profile}
Under the same hypotheses define
\[
 H_z=\frac12[(ac-1)z+c]^2-D(az+1)-\frac12z^2,
 \qquad G_z=\frac{dH_z}{dt}.
\]
For every $t\ge0$ and $z\in[1/2,1]$,
\[
 \left|E(\xi_z,t)-\frac12-\eta^2H_z(t)\right|\le20\eta^4,
 \qquad
 |\dot E(\xi_z,t)-\eta^2G_z(t)|\le80\lambda\eta^4.
\]
The same bounds hold on each of the two sectors in
Corollary~\ref{sbsector:two-concept}, with its own parameters.
\end{proposition}
\begin{proof}
Here is an exact expansion with explicit estimates. For a graded matrix
with entries $(\alpha,\eta b;\eta c_*,\eta^2d)$, put
\[
 A_0=(\alpha-1)z+s b,\qquad T_0=c_*z+s d k.
\]
The exact planar identity is
\[
 E=\frac12+\eta^2\left(\frac12A_0^2-s c_*z-d-\frac12z^2\right)
                  +\eta^4K,
\]
where
\[
 K=dz^2-\frac{s z^2(bA_0-c_*z)}{1+k}
       +\frac12T_0^2+\frac{\eta^2b^2z^4}{2(1+k)^2}.
\]
This follows by squaring the exact residual and using
$k-1=-\eta^2z^2/(1+k)$. In our case
$(\alpha,b,c_*,d)=(ac,-c,-aD,D)$ and $s=-1$, giving $H_z$ above.
Each coefficient is bounded by $2$ and its rate by $5\lambda$.
Consequently $|A_0|\le5$, $|\dot A_0|\le10\lambda$,
$|T_0|\le4$, $|\dot T_0|\le10\lambda$, and $1+k\ge7/4$.
Termwise estimation gives
\[
 |K|\le2+\frac{12}{7/4}+8+\frac{1/4}{2(7/4)^2}<20.
\]
Differentiation with the probe fixed yields
\[
 \dot K=\dot d z^2
 -\frac{s z^2(\dot b A_0+b\dot A_0-\dot c_*z)}{1+k}
 +T_0\dot T_0
 +\frac{\eta^2b\dot b z^4}{(1+k)^2},
\]
so
\[
 |\dot K|\le\lambda\left(5+\frac{50}{7/4}+40
                                  +\frac{10/16}{(7/4)^2}\right)<80\lambda.
\]
All scaled coefficient bounds have been proved along the actual
trajectory; they are not additional hypotheses of this proposition.
\end{proof}

\begin{remark}[Profile versus a universal limiting curve]
The displayed $H_z$ is explicit in the actual, initialization-determined
coefficient dynamics, and the error bound is uniform for all training
times. Its coefficients can still depend on $\eta$ and $a_0$.
This proposition does not additionally claim convergence to a single
parameter-independent curve under a time rescaling. Its reversal proof
uses the exact signs, not the accuracy bound at a vanishing sign margin.
\end{remark}

\section{Restricted reversal theorem and scope of the general noisy reversal}
The sector, amplitude, finite-time reversal, coefficient-rate, and
derivative-remainder obligations are now closed in the stated noiseless
initialization regime. The two-concept extension removes the limitation
that concept 2 is left unfitted, and the isotropic corollary derives the
required pure sectors with fixed positive probability for two neurons.
Neither gives a high-probability result or treats nonzero data noise.
For the noisy isotropic many-neuron model, a reached learning interval with
useful relative grading, actual-probe control, and reversal signs remains a
separate open problem.  Under the present theorem scope this general reversal
is not a predecessor of the specialization theorem: the noisy many-neuron
load-bearing target is concept specialization and the resulting predictive
normal form.  Full weak-basin creation, capture, tracking and positive weak
mass remain necessary for that specialization target.  None of the general
reversal claims follows from the restricted result by analogy.

\section{General-model creation: regulator and data-side advances}
\label{sbgen:scope}

This continuation concerns the noisy many-neuron model.  It proves an exact
whole-network regulation equation, a sharp strong-family damping clock, a
signed two-feature decomposition of the regulator forcing, and a
first-hitting reduction from generation to the gate-cliff creation corridor.
It also proves an analytic high-SNR region for the population transition
kernel and an explicit joint concentration bound for the creation gap.  The
remaining generation work is now trajectory-level: verify the strong
input-side sign condition, control the adverse complement forcing, and show
that enough mass-weighted clock accumulates before the B2 and weak-feature
guards expire.  The complete initialization-to-creation theorem is still
open.

\section{A whole-network equation for the regulation deficit}
\label{sbgen:regulation-section}

Let $\langle\cdot,\cdot\rangle_n$ denote the equally weighted empirical
inner product over both clusters. Write
\[
 f(x)=\sum_i w_i a_i(x),\quad a_i(x)=(u_i^\top x)_+,
 \quad e(x)=x-f(x),\quad p(x)=(x_1)_+,\quad D=\|p\|_n^2>0.
\]
All vector output quantities below include the ambient coordinates.
Along the selected gradient flow let $g_i(x)\in[0,1]$ be compatible
training selectors. Every derivative statement is almost everywhere;
the integrated statements follow by absolute continuity.
Define
\[
 c_i=\frac{\langle p,a_i\rangle_n}{D},\qquad
 \ell_i=\frac{\mathbb E_n[p g_i x]}{D},\qquad
 k_i(x)=g_i(x)\ell_i^\top x,
\]
\[
 y=\frac{\mathbb E_n[p f]}D,\quad
 y_*=\frac{\mathbb E_n[p x]}D=(1,b_*,0,\ldots,0),\quad
 d=y_*-y,\quad R(x)=e(x)-p(x)d.
\]
Thus $\mathbb E_n[pR]=0$, $y_1=m$, $y_2=b$, and $d_2=b_*-b$.

\begin{proposition}[Exact projected regulator and its defect]
\label{sbgen:network-regulator}
Set
\[
 K=D\sum_i\left[c_i^2I+\|\ell_i\|^2w_iw_i^\top\right],
\]
\[
 a_i^\circ=a_i-c_i p,\qquad
 k_i^\circ=k_i-\|\ell_i\|^2p,
\]
\[
 r=\sum_i c_i\mathbb E_n[R a_i^\circ]
       +\sum_i w_i\mathbb E_n[(w_i^\top R)k_i^\circ].
\]
Then, exactly,
\begin{equation}\label{sbgen:vector-regulator}
 \dot y=Kd+r,\qquad K\succeq D\Big(\sum_i c_i^2\Big)I.
\end{equation}
In particular, put
\[
 \kappa=K_{22},\qquad
 \chi=K_{21}(1-m)-\sum_{j\ge3}K_{2j}y_j+r_2.
\]
The signed regulation deficit satisfies
\begin{equation}\label{sbgen:scalar-regulator}
 \frac{d}{dt}(b_*-b)=-\kappa(b_*-b)-\chi.
\end{equation}
If $S=\sum_i\|w_i\|^2>0$, then
\begin{equation}\label{sbgen:clock-lower}
 \kappa\ge D\sum_i c_i^2\ge\frac{Dm^2}{S}.
\end{equation}
No fixed training mask, fixed cohort, or frozen-kernel approximation is
used in these identities.
\end{proposition}
\begin{proof}
The exact gradient equations are
$\dot w_i=\mathbb E_n[e a_i]$ and
$\dot u_i=\mathbb E_n[g_i x(w_i^\top e)]$.
Differentiating the projected output yields
\[
 \dot y=\sum_i c_i\mathbb E_n[e a_i]
             +\sum_i w_i\mathbb E_n[(w_i^\top e)k_i].
\]
Substitute $e=pd+R$. The two required identities are
\[
 \langle p,a_i\rangle_n=Dc_i,\qquad
 \langle p,k_i\rangle_n=D\|\ell_i\|^2.
\]
Orthogonality of $R$ to $p$ gives the displayed $r$ and $K$.
The matrix is positive semidefinite by its explicit sum of squares.
Its second row gives \eqref{sbgen:scalar-regulator}. Finally,
$m=\sum_i w_{i1}c_i$, so Cauchy--Schwarz gives
$m^2\le S\sum_i c_i^2$. The derivation holds with the compatible
selectors at gate changes and integrates along the same trajectory.
\end{proof}

\begin{corollary}[A signed generation-to-network certificate]
\label{sbgen:deficit-certificate}
Let $\Lambda(a,t)=\int_a^t\kappa(s)\,ds$ and
$\epsilon(t)=[b_*-b(t)]_+$. Then
\begin{equation}\label{sbgen:deficit-duhamel}
 \epsilon(t)\le e^{-\Lambda(a,t)}\epsilon(a)
       +\int_a^t e^{-\Lambda(s,t)}[-\chi(s)]_+\,ds.
\end{equation}
In particular, if $\chi\ge-\epsilon_g\kappa$ on this interval, then
\[
 \epsilon(t)\le\epsilon_g+
                  e^{-\Lambda(a,t)}(\epsilon(a)-\epsilon_g).
\]
If $\epsilon(a)>\epsilon_g$ and the permitted deficit is
$\epsilon_{\rm req}>\epsilon_g$, the sufficient clock condition is
\[
 \Lambda(a,t)\ge
 \left[\log\frac{\epsilon(a)-\epsilon_g}
                       {\epsilon_{\rm req}-\epsilon_g}\right]_+.
\]
When $m\ge m_0>0$ and $S\le S_*$ on the interval, the lower bound
$\Lambda(a,t)\ge Dm_0^2(t-a)/S_*$ is available.
\end{corollary}
\begin{proof}
Variation of constants applied to the signed scalar equation gives an
exact formula before taking the positive part. Replacing its initial
value and adverse forcing by their positive parts gives
\eqref{sbgen:deficit-duhamel}. Integrate
$\kappa(s)e^{-\Lambda(s,t)}$ to obtain the floor comparison. The last
claim is \eqref{sbgen:clock-lower} integrated in time.
\end{proof}

\begin{proposition}[An explicit directional forcing budget]
\label{sbgen:forcing-budget}
Define the vector-valued, $p$-orthogonal training function
\[
 Q_2(x)=e_2\sum_i c_i a_i^\circ(x)
                   +\sum_i w_{i2}w_i k_i^\circ(x).
\]
Then $r_2=\mathbb E_n[R^\top Q_2]$. Suppose the training data lie in
$\operatorname{span}\{e_1,e_2\}$, let
$\lambda_n=\|\mathbb E_nxx^\top\|_{\rm op}$, and use exact balance
$\sum_i\|u_i\|^2=\sum_i\|w_i\|^2=S$ together with the established
bound $\|W_\perp\|_F\le\varepsilon\sqrt h$.
On a window with $0\le m\le1$,
\begin{equation}\label{sbgen:forcing-charge}
 [-\chi]_+\le
 (1-m)[-K_{21}]_+
 +\|R\|_n\|Q_2\|_n
 +h\varepsilon^2\lambda_n S\sqrt{\lambda_n/D}.
\end{equation}
The sharper signed quantity $\chi$ may be used directly instead of
this norm bound. In dimension two the last term is zero.
Orthogonal projection also gives
$\|R\|_n^2=2\mathcal L-D\|d\|^2\le2\mathcal L$.
\end{proposition}
\begin{proof}
The formula for $r_2$ is a regrouping of the exact sums.
Cauchy--Schwarz gives its norm bound. For the ambient term,
\[
 \|\ell_i\|^2\le\lambda_n/D,\quad
 \sum_i c_i^2\le\lambda_n S/D,\quad
 \|y_\perp\|\le\varepsilon\sqrt h\sqrt{\lambda_n S/D}.
\]
The first inequality follows by testing $\ell_i$ against unit vectors
and applying Cauchy--Schwarz with $g_i\le1$; the second follows from
$c_i^2\le\|a_i\|_n^2/D$.
Consequently
\[
 \left|\sum_{j\ge3}K_{2j}y_j\right|
 \le\lambda_n\sqrt S\,\varepsilon\sqrt h\,\|y_\perp\|
 \le h\varepsilon^2\lambda_n S\sqrt{\lambda_n/D}.
\]
Taking the adverse part of the remaining first-coordinate term gives
\eqref{sbgen:forcing-charge}.
\end{proof}

This closes the algebraic bridge from whole-network motion to a deficit
bound.  The next lemmas exploit the angular structure of the same exact
regulator rather than charging its defect by a global residual norm.

\begin{lemma}[Strong-cone signed regulator and sharp damping clock]
\label{sbgen:strong-cone-regulator}
Continue with the notation of Proposition~\ref{sbgen:network-regulator} and
assume the training data lie in
$\mathcal S=\operatorname{span}\{e_1,e_2\}$.  Write
\[
a_i=r_iq_{\phi_i},
\qquad
q_\phi(x):=(\widetilde u(\phi)^\top x)_+,
\qquad
\mathsf H_1(\phi):=\langle p,q_\phi\rangle_n.
\]
Then
\[
c_i=\frac{r_i\mathsf H_1(\phi_i)}D
\]
and the scalar damping decomposes exactly as
\[
\boxed{
\kappa
=
\underbrace{\frac1D\sum_i r_i^2\mathsf H_1(\phi_i)^2}_{\kappa_w}
+
\underbrace{D\sum_i\|\ell_i\|^2w_{i2}^2}_{\kappa_u}
}
\tag{sbgen-kappa-split}\label{sbgen:kappa-split}
\]
with both terms nonnegative.

Let $B$ be a strong family on a time interval $I$ and suppose
\[
0\le\phi_i(t)\le\bar\phi<\frac\pi2,
\qquad i\in B,
\qquad t\in I.
\]
Set
\[
\Pi_B(t):=\sum_{i\in B}r_i(t)^2,
\qquad
H_{1,-}:=\inf_{0\le\phi\le\bar\phi}\mathsf H_1(\phi).
\]
Then
\[
\boxed{
\kappa(t)
\ge
\frac{H_{1,-}^2}{D}\Pi_B(t).
}
\tag{sbgen-sharp-clock}\label{sbgen:sharp-clock}
\]
If
\[
\lambda_n:=\lambda_{\max}(\Sigma_n),
\]
then ReLU Lipschitzness gives the deterministic lower bound
\[
\boxed{
H_{1,-}
\ge
D-2\sqrt{D\lambda_n}\sin\frac{\bar\phi}{2}.
}
\tag{sbgen-H1-lower}\label{sbgen:H1-lower}
\]

The forcing also splits exactly as
\[
\boxed{\chi=\chi_w+\chi_u,}
\tag{sbgen-chi-split}\label{sbgen:chi-split}
\]
where
\[
\boxed{
\chi_w
=
\frac1D\sum_i
r_i^2\mathsf H_1(\phi_i)
\langle R_2,q_{\phi_i}\rangle_n
}
\tag{sbgen-chi-w}\label{sbgen:chi-w}
\]
and, writing
\[
y_\perp:=P_{\mathcal S^\perp}y,
\]
\[
\boxed{
\begin{aligned}
\chi_u
={}&
D(1-m)\sum_i\|\ell_i\|^2w_{i1}w_{i2}
\\
&-
D\sum_i\|\ell_i\|^2w_{i2}
\langle w_{i,\perp},y_\perp\rangle
\\
&+
\sum_iw_{i2}\langle w_i^\top R,k_i^\circ\rangle_n.
\end{aligned}
}
\tag{sbgen-chi-u}\label{sbgen:chi-u}
\]
In dimension two the middle line vanishes.  For a neuron set $B$, write
$\chi_{w,B}$ and $\chi_{u,B}$ for the corresponding sums in
\eqref{sbgen:chi-w} and \eqref{sbgen:chi-u} restricted to $i\in B$.

Define
\[
F_t(\phi):=\langle R_2,q_\phi\rangle_n.
\]
Since $\langle p,R_2\rangle_n=0$ and $q_0=p$,
\[
F_t(0)=0.
\]
For each fixed $t$, $F_t$ is absolutely continuous in $\phi$, and for a.e.
$\phi$,
\[
\boxed{
F_t'(\phi)
=
\left\langle
R_2,
(\widetilde u^\angle(\phi)^\top x)
\mathbf1\{\widetilde u(\phi)^\top x>0\}
\right\rangle_n.
}
\tag{sbgen-Fprime}\label{sbgen:Fprime}
\]
Thus, with
\[
\alpha_2(t)
:=
\operatorname*{ess\,inf}_{0\le\phi\le\bar\phi}F_t'(\phi),
\]
if
\[
\alpha_2(t)\ge\underline\alpha_2>0
\qquad(t\in I),
\]
then the strong-family output-side forcing obeys
\[
\boxed{
\chi_{w,B}(t)
\ge
\frac{H_{1,-}\underline\alpha_2}{D}
\sum_{i\in B}r_i(t)^2\phi_i(t)
\ge0.
}
\tag{sbgen-chi-w-positive}\label{sbgen:chi-w-positive}
\]

For the input-side forcing set
\[
\ell_-^2
:=
\inf_{t\in I,i\in B}\|\ell_i(t)\|^2,
\qquad
\ell_+^2
:=
\sup_{t\in I,i\in B}\|\ell_i(t)\|^2,
\]
\[
K_\circ
:=
\sup_{t\in I,i\in B}\|k_i^\circ(t)\|_n,
\qquad
R_\star:=\sup_{t\in I}\|R_t\|_n,
\]
\[
\beta_B
:=
\sup_{t\in I,i\in B}
\frac{\|w_i(t)\|}{w_{i1}(t)},
\qquad
\gamma_{\perp,B}
:=
\sup_{t\in I,i\in B}
\frac{\|w_{i,\perp}(t)\|}{w_{i1}(t)}.
\]
Suppose
\[
w_{i1}>0,
\qquad
w_{i2}\ge0
\qquad(i\in B,t\in I),
\]
$m(t)\le m_{\max}<1$, and
\[
\boxed{
D(1-m_{\max})\ell_-^2
\ge
D\ell_+^2\gamma_{\perp,B}
\sup_{t\in I}\|y_\perp(t)\|
+
\beta_BR_\star K_\circ.
}
\tag{sbgen-u-positive-condition}\label{sbgen:u-positive-condition}
\]
Then
\[
\boxed{\chi_{u,B}(t)\ge0\qquad(t\in I).}
\tag{sbgen-chi-u-positive}\label{sbgen:chi-u-positive}
\]
At exact strong-axis alignment,
\[
\ell_0=(1,b_*)
\]
and
\[
\boxed{
k_0^\circ
=
b_*\mathbf1\{x_1>0\}(x_2-b_*x_1),
}
\tag{sbgen-kcirc-axis}\label{sbgen:kcirc-axis}
\]
so the residual-centered input-side defect carries an explicit factor $b_*$
at the axis.

Finally define
\[
\chi_{\rm comp}
:=
\chi-\chi_{w,B}-\chi_{u,B}.
\]
If
\[
\boxed{
\chi_{\rm comp}(t)\ge-\epsilon_g\kappa(t)
\qquad(t\in I),
}
\tag{sbgen-complement-forcing}\label{sbgen:complement-forcing}
\]
then
\[
\boxed{\chi(t)\ge-\epsilon_g\kappa(t)}
\tag{sbgen-total-forcing}\label{sbgen:total-forcing}
\]
throughout $I$.  Thus the specialized strong family is favorable and only
the complement can consume the regulator budget.
\end{lemma}

\begin{proof}
The first identity follows from
$c_i=r_i\mathsf H_1(\phi_i)/D$ and the exact formula for $K_{22}$ in
Proposition~\ref{sbgen:network-regulator}.  Restricting the nonnegative
$\kappa_w$ sum to $B$ proves \eqref{sbgen:sharp-clock}.

For \eqref{sbgen:H1-lower}, ReLU is one-Lipschitz, so
\[
\|q_\phi-q_0\|_n
\le
\sqrt{\lambda_n}\|\widetilde u(\phi)-e_1\|
=
2\sqrt{\lambda_n}\sin\frac{|\phi|}{2}.
\]
Because $q_0=p$ and $\|p\|_n=\sqrt D$,
\[
\mathsf H_1(\phi)
\ge
D-\sqrt D\|q_\phi-p\|_n,
\]
which yields the displayed bound.

For the forcing, take the second coordinate in the exact residual term $r$
of Proposition~\ref{sbgen:network-regulator}.  Since
$\langle p,R_2\rangle_n=0$,
\[
c_i\langle R_2,a_i^\circ\rangle_n
=
\frac{r_i^2\mathsf H_1(\phi_i)}D
\langle R_2,q_{\phi_i}\rangle_n,
\]
which is \eqref{sbgen:chi-w}.  Moreover
\[
K_{21}(1-m)
=
D(1-m)\sum_i\|\ell_i\|^2w_{i1}w_{i2},
\]
and, because $d_j=-y_j$ on ambient coordinates,
\[
-\sum_{j\ge3}K_{2j}y_j
=
-D\sum_i\|\ell_i\|^2w_{i2}
\langle w_{i,\perp},y_\perp\rangle.
\]
Adding the second-coordinate input-feature part of $r$ gives
\eqref{sbgen:chi-u}.

The function $q_\phi(x)$ is Lipschitz in $\phi$ for every fixed training
sample, hence $F_t$ is absolutely continuous.  Its a.e. derivative is
\eqref{sbgen:Fprime}.  Since $F_t(0)=0$, the lower derivative bound implies
$F_t(\phi)\ge\underline\alpha_2\phi$ on the strong cone.  Substitution into
\eqref{sbgen:chi-w} proves \eqref{sbgen:chi-w-positive}.

For the input-side term, Cauchy--Schwarz gives
\[
|\langle w_i^\top R,k_i^\circ\rangle_n|
\le
\|w_i\|R_\star K_\circ
\]
and
\[
|\langle w_{i,\perp},y_\perp\rangle|
\le
\|w_{i,\perp}\|\|y_\perp\|.
\]
Factoring $w_{i1}w_{i2}\ge0$ from the resulting lower bound shows that
\eqref{sbgen:u-positive-condition} implies
\eqref{sbgen:chi-u-positive}.  At $\phi=0$,
\[
\ell_0
=
D^{-1}\mathbb E_n[p\mathbf1\{x_1>0\}x]
=(1,b_*),
\]
and direct substitution into
$k_0^\circ=k_0-\|\ell_0\|^2p$ gives
\eqref{sbgen:kcirc-axis}.  The final claim follows by adding the strong-family
and complement contributions.
\end{proof}

\begin{lemma}[Two-feature lower bound for the regulator tangent margin]
\label{sbgen:regulator-transition-kernel}
Let
\[
q=(x_2)_+,
\qquad
C:=\langle p,q\rangle_n,
\qquad
q_\perp:=q-\frac CDp,
\qquad
D_{2|1}:=\|q_\perp\|_n^2>0,
\]
and let $P_{12}$ be empirical $L^2$ projection onto
$\operatorname{span}\{p,q\}$.  Define
\[
r_2^{\rm reg}:=f_2-bp,
\qquad
\gamma_2^{\rm reg}
:=
\frac{\langle q_\perp,r_2^{\rm reg}\rangle_n}{D_{2|1}},
\qquad
e_2^{\rm reg}:=(I-P_{12})r_2^{\rm reg}.
\]
Then
\[
r_2^{\rm reg}=\gamma_2^{\rm reg}q_\perp+e_2^{\rm reg},
\qquad
e_2^{\rm reg}\perp\operatorname{span}\{p,q\}.
\]
For
\[
g_\varphi(x)
:=
(\widetilde u^\angle(\varphi)^\top x)
\mathbf1\{\widetilde u(\varphi)^\top x>0\},
\]
define
\[
T_{\rm reg}(\varphi)
:=
\langle x_2-b_*p,g_\varphi\rangle_n,
\qquad
K_q(\varphi):=\langle q_\perp,g_\varphi\rangle_n,
\]
\[
J_{\rm reg}(\theta,\varphi)
:=
\left\langle
(I-P_{12})q_\theta,
(I-P_{12})g_\varphi
\right\rangle_n.
\]
Writing
\[
\beta_i:=r_i z_{i2},
\]
one has the exact identity
\[
\boxed{
F_t'(\varphi)
=
T_{\rm reg}(\varphi)
-
\gamma_2^{\rm reg}K_q(\varphi)
-
\sum_i\beta_iJ_{\rm reg}(\phi_i,\varphi).
}
\tag{sbgen-Fprime-twofeature}\label{sbgen:Fprime-twofeature}
\]

For
\[
L_q(\theta):=\langle q_\perp,q_\theta\rangle_n,
\]
one also has
\[
\boxed{
D_{2|1}\gamma_2^{\rm reg}
=
\sum_i\beta_iL_q(\phi_i).
}
\tag{sbgen-gamma-measure}\label{sbgen:gamma-measure}
\]

Fix $0<\bar\phi<\pi/2$ and let
\[
\mathcal G
:=
\{i:0\le\phi_i\le\pi/2,\ \beta_i\ge0\}.
\]
Suppose, on the relevant deterministic angular region,
\[
L_q(\theta)\ge0,
\]
and there is a finite constant $\rho_{\rm reg}$ such that
\[
\boxed{
[J_{\rm reg}(\theta,\varphi)]_+
\le
\rho_{\rm reg}L_q(\theta)
}
\tag{sbgen-Jreg-profile}\label{sbgen:Jreg-profile}
\]
for every $0\le\theta\le\pi/2$ and $0\le\varphi\le\bar\phi$; at points with
$L_q(\theta)=0$, this condition includes
$[J_{\rm reg}(\theta,\varphi)]_+=0$.
Define
\[
\mathcal B_q
:=
\sum_{i\notin\mathcal G}
|\beta_i|\,|L_q(\phi_i)|,
\]
\[
\mathcal B_J
:=
\sup_{0\le\varphi\le\bar\phi}
\sum_{i\notin\mathcal G}
|\beta_i|\,|J_{\rm reg}(\phi_i,\varphi)|,
\]
\[
T_-:=\inf_{0\le\varphi\le\bar\phi}T_{\rm reg}(\varphi),
\qquad
K_{q,+}:=\sup_{0\le\varphi\le\bar\phi}[K_q(\varphi)]_+.
\]
Then
\[
\boxed{
\alpha_2(t)
\ge
T_-
-
[\gamma_2^{\rm reg}(t)]_+
\bigl(K_{q,+}+\rho_{\rm reg}D_{2|1}\bigr)
-
\rho_{\rm reg}\mathcal B_q
-
\mathcal B_J.
}
\tag{sbgen-alpha2-bound}\label{sbgen:alpha2-bound}
\]
Consequently, if
\[
[\gamma_2^{\rm reg}(t)]_+\le\Gamma_2
\]
and
\[
\boxed{
T_-
>
\Gamma_2
\bigl(K_{q,+}+\rho_{\rm reg}D_{2|1}\bigr)
+
\rho_{\rm reg}\mathcal B_q
+
\mathcal B_J,
}
\tag{sbgen-regulator-adm}\label{sbgen:regulator-adm}
\]
then $\alpha_2(t)>0$, so the strong-family output-side forcing in
\eqref{sbgen:chi-w-positive} is favorable.
\end{lemma}

\begin{proof}
Since
\[
R_2=x_2-f_2-(b_*-b)p
=x_2-b_*p-r_2^{\rm reg},
\]
orthogonal decomposition of $r_2^{\rm reg}$ gives
\[
F_t'(\varphi)
=
T_{\rm reg}(\varphi)
-
\gamma_2^{\rm reg}K_q(\varphi)
-
\langle e_2^{\rm reg},(I-P_{12})g_\varphi\rangle_n.
\]
Because
\[
f_2=\sum_i r_i z_{i2}q_{\phi_i},
\]
\[
e_2^{\rm reg}
=(I-P_{12})f_2
=
\sum_i\beta_i(I-P_{12})q_{\phi_i},
\]
which proves \eqref{sbgen:Fprime-twofeature}.  Taking the $q_\perp$
coefficient of $r_2^{\rm reg}$ gives \eqref{sbgen:gamma-measure}.

For the good set,
\[
\sum_{i\in\mathcal G}\beta_i
[J_{\rm reg}(\phi_i,\varphi)]_+
\le
\rho_{\rm reg}
\sum_{i\in\mathcal G}\beta_iL_q(\phi_i).
\]
Using \eqref{sbgen:gamma-measure},
\[
\sum_{i\in\mathcal G}\beta_iL_q(\phi_i)
\le
D_{2|1}[\gamma_2^{\rm reg}]_+ +\mathcal B_q.
\]
The complement contributes at worst $\mathcal B_J$.  Finally
$-\gamma_2^{\rm reg}K_q\ge-[\gamma_2^{\rm reg}]_+K_{q,+}$, which yields
\eqref{sbgen:alpha2-bound} and the stated positivity criterion.
\end{proof}

\begin{remark}[Relation to the creation weak-feature clock]
The coefficient $\gamma_2^{\rm reg}$ is the whole-network weak-feature
coefficient of $f_2-bp$.  It is not automatically identical to the
non-strong coefficient $\gamma_{2,R}$ in
Lemma~\ref{lem:middle-signed-decomposition}, because that lemma first removes
the late strong-family shape contribution.  The two are linked by an
explicit strong-shape term.  Thus the same weak-fit timing package can be
reused, but no equality between $\gamma_2^{\rm reg}$ and
$\gamma_{2,R}$ is assumed without that bookkeeping.
\end{remark}

\begin{proposition}[Regulator clock and first hitting of the creation corridor]
\label{sbgen:first-hit}
Fix deterministic angles
\[
0<\phi_L<\phi_M<\frac\pi2,
\qquad
\tau_L:=\cot\phi_L>\tau_M:=\cot\phi_M>0,
\]
and deterministic $b$-unit allowances $C_M,C_L\ge0$.  At fixed $\tau$,
Lemma~\ref{lem:middle-two-progress} gives the exact affine identity
\[
b_{\rm lift}^{(n)}(\tau,m)
=
b_{\rm lift}^{(n)}(\tau,0)-m\tau.
\]
Define
\[
\boxed{
B_M(m):=b_{\rm lift}^{(n)}(\tau_M,m)+C_M,
\qquad
B_L(m):=b_{\rm lift}^{(n)}(\tau_L,m)-C_L.
}
\tag{sbgen-corridor}\label{sbgen:corridor}
\]
Then
\[
B_L(m)-B_M(m)
=
B_L(0)-B_M(0)-m(\tau_L-\tau_M),
\]
so the corridor width decreases with $m$.

Let $I=[t_a,t_b]$ be a post-specialization interval on which
\[
\dot m(t)\ge0\quad\text{a.e.},
\qquad
m(t)\le m_{\max},
\]
and suppose
\[
\boxed{B_M(m_{\max})<B_L(m_{\max}).}
\tag{sbgen-corridor-positive}\label{sbgen:corridor-positive}
\]
Assume also
\[
\chi(t)\ge-\epsilon_g\kappa(t)
\qquad(t\in I)
\]
and define
\[
\boxed{
\delta_{\rm gen}
:=
\inf_{m\in[m(t_a),m_{\max}]}
\bigl[b_*-B_M(m)-\epsilon_g\bigr]
>0.
}
\tag{sbgen-delta-gen}\label{sbgen:delta-gen}
\]
Let
\[
G_M(t):=B_M(m(t))-b(t)
\]
and suppose $G_M(t_a)>0$.  Before the first zero of $G_M$,
\[
\boxed{
\dot G_M(t)
\le
-\kappa(t)\bigl(G_M(t)+\delta_{\rm gen}\bigr).
}
\tag{sbgen-G-ineq}\label{sbgen:G-ineq}
\]
Hence, with
\[
\Lambda(t_a,t):=\int_{t_a}^t\kappa(s)\,ds,
\]
\[
\boxed{
G_M(t)+\delta_{\rm gen}
\le
\bigl(G_M(t_a)+\delta_{\rm gen}\bigr)e^{-\Lambda(t_a,t)}.
}
\tag{sbgen-G-duhamel}\label{sbgen:G-duhamel}
\]
If
\[
\boxed{
\Lambda(t_a,t_b)
>
\log\frac{G_M(t_a)+\delta_{\rm gen}}{\delta_{\rm gen}},
}
\tag{sbgen-first-hit-clock}\label{sbgen:first-hit-clock}
\]
then there is a first time
$t_{\rm cert}\in(t_a,t_b)$ such that
\[
\boxed{
b(t_{\rm cert})=B_M(m(t_{\rm cert})).}
\tag{sbgen-first-hit}\label{sbgen:first-hit-eq}
\]
At that time
\[
B_M(m(t_{\rm cert}))<B_L(m(t_{\rm cert})).
\]
Thus the middle lift has reached the reserved margin $C_M$ before the lower
left-endpoint threshold is exhausted.  If $C_M$ dominates the middle adverse
charges in \eqref{eq:GC-middle-budget}, $C_L$ certifies the left negative
endpoint, and the remaining guards of
Proposition~\ref{prop:gate-cliff-creation} hold, then the gate-cliff creation
certificate fires at $t_{\rm cert}$.

If the strong family satisfies Lemma~\ref{sbgen:strong-cone-regulator}, then
\[
\boxed{
\Lambda(t_a,t_b)
\ge
\int_{t_a}^{t_b}
\frac{H_{1,-}(s)^2}{D}\Pi_B(s)\,ds.
}
\tag{sbgen-mass-clock}\label{sbgen:mass-clock}
\]
Thus the load-bearing generation clock is the accumulated strong-family
mass-weighted clock, not the cruder pointwise lower bound $Dm^2/S$.
\end{proposition}

\begin{proof}
Since
$B_M(m)=B_M(0)-\tau_Mm$,
\[
\dot B_M=-\tau_M\dot m.
\]
The exact regulator equation is equivalent to
\[
\dot b=\kappa(b_*-b)+\chi.
\]
Therefore
\[
\begin{aligned}
\dot G_M
&=
-\tau_M\dot m
-\kappa(b_*-b)-\chi
\\
&=
-\tau_M\dot m
-\kappa G_M
-\kappa[b_*-B_M(m)]
-\chi.
\end{aligned}
\]
Using $\dot m\ge0$, $\chi\ge-\epsilon_g\kappa$, and
\eqref{sbgen:delta-gen} gives \eqref{sbgen:G-ineq}.  Multiplication by the
integrating factor $e^{\Lambda(t_a,t)}$ gives
\eqref{sbgen:G-duhamel}.  If no zero occurred before $t_b$, then
$G_M(t_b)>0$, whereas \eqref{sbgen:first-hit-clock} and
\eqref{sbgen:G-duhamel} force $G_M(t_b)<0$, a contradiction.  Continuity
therefore gives the first hitting equality.

Because $m(t_{\rm cert})\le m_{\max}$ and the corridor width decreases in
$m$, \eqref{sbgen:corridor-positive} gives
$B_M(m(t_{\rm cert}))<B_L(m(t_{\rm cert}))$.  The final clock bound is
\eqref{sbgen:sharp-clock} integrated in time.
\end{proof}

\begin{remark}[Creation is a first-hitting event, not an invariant window]
The regulator reference $b_*$ need not be an upper barrier.  After the first
middle-threshold hit, favorable signed forcing may continue to increase $b$
and move the lower repelling crossing to smaller angles.  The creation
certificate is therefore asserted at $t_{\rm cert}$; it does not require
$b$ to remain forever between two fixed lift thresholds.  This is the
proof-level counterpart of the observed early downward separator sweep.
\end{remark}

\begin{remark}[Joint gap concentration feeds the first-hitting margin]
At fixed deterministic $\tau_M$,
\[
b_*-B_M(m)
=
\Delta_{\rm cr}^{(n)}(\tau_M,m)-C_M.
\]
By Proposition~\ref{sbgen:joint-gap}, the sampling error in
$\Delta_{\rm cr}^{(n)}(\tau_M,m)$ is the same for every real $m$ on one event.
Thus the lower bound in \eqref{sbgen:delta-gen} is obtained from one joint
creation-gap concentration inequality, the clipping correction, and the
deterministic allowance $C_M$; no union bound over progress values is needed.
\end{remark}


\section{An analytic region for the transition-kernel sign}
\label{sbgen:J-section}

Consider equal-noise population SIM with means $\mu_1e_1,\mu_2e_2$ and
covariance $\sigma^2I_2$, equal cluster weights, and
$0<\mu_2\le\mu_1$. Put
\[
 r=\mu_2/\mu_1,\quad \delta=\sigma/\mu_1,\quad
 \tau=z\delta,\quad z>0,\quad
 I_\theta=[\theta,\pi/2-\theta],\quad 0<\theta<\pi/4.
\]
Let $P$ be population $L^2$ projection onto
$\operatorname{span}\{(x_1)_+,(x_2)_+\}$ and set
\[
 q_\phi(x)=(\cos\phi\,x_1+\sin\phi\,x_2)_+,\quad
 h_\tau(x)=(x_1-\tau x_2)\mathbf1\{\tau x_1+x_2>0\},
\]
\[
 J_{12}^{\rm pop}(\phi;\tau)=\langle(I-P)q_\phi,h_\tau\rangle.
\]
Write $\varphi,\Phi$ for the standard normal density and distribution
function, and define
\[
 c_N=1/\sqrt{2\pi},\quad v_N=\sqrt{1/2-c_N^2},\quad
 F(z)=\varphi(z)-c_N(1-\Phi(z)).
\]

\begin{theorem}[Uniform high-SNR positivity with explicit constants]
\label{sbgen:J-positive}
Set $\delta_0=r/4$ and
\[
 C_\ell=\frac{\sqrt3}{4}
       \left(1+\frac1r+\frac1{r\sin\theta}\right),\quad
 C_R=C_\ell+\frac3r,
\]
\[
 H_0=\sqrt{\frac12+\delta_0^2
                   +z^2\delta_0^2\left(\frac{r^2}2+\delta_0^2\right)},
\]
\[
 \begin{split}
 C_h=\frac12\bigg[&v_N(1+z\delta_0)
 +c_Nz\left((z+c_N)\sqrt{2/\pi}+z\delta_0\right)\\
 &+v_N\sqrt{1+z^2(r^2+\delta_0^2)}\bigg],
 \qquad C=C_RH_0+C_h.
 \end{split}
\]
For every $0<\delta\le\delta_0$, uniformly in $\phi\in I_\theta$,
\begin{equation}\label{sbgen:J-expansion}
 \left|\frac{J_{12}^{\rm pop}(\phi;z\delta)}{\mu_1^2}
       -\frac\delta2\sin\phi\,F(z)\right|\le C\delta^2.
\end{equation}
Moreover $F(z)\ge\varphi(z)/2>0$. Thus the primitive region
\begin{equation}\label{sbgen:J-region}
 0<\frac\sigma{\mu_1}\le
 \min\left\{\frac r4,\frac{\sin\theta\,F(z)}{4C}\right\}
\end{equation}
is nonempty and certifies
\begin{equation}\label{sbgen:J-margin}
 \inf_{\phi\in I_\theta}J_{12}^{\rm pop}(\phi;\tau)
 \ge\frac{\mu_1\sigma\sin\theta F(z)}4>0.
\end{equation}
This is a compact-interior theorem. It does not assert a uniform positive
margin up to the two axes, where the projected feature vanishes.
\end{theorem}

\begin{proof}
Scale $x=\mu_1Y$ and work in the joint probability space of the cluster
label and independent standard normals $Z_1,Z_2$. The two equally likely
values of $Y$ are
$(1+\delta Z_1,\delta Z_2)$ and
$(\delta Z_1,r+\delta Z_2)$. All inner products below use this mixture.
Let $N(t)=(-t)_+$, $p=(Y_1)_+$, $q=(Y_2)_+$, and
$L_\phi=\cos\phi\,Y_1+\sin\phi\,Y_2$.
Define
\[
 \ell_\phi=\cos\phi\,p+\sin\phi\,q-(L_\phi)_+,
\]
and let $\ell_{0,\phi}$ equal $\sin\phi N(Z_2)$ on cluster 1 and
$\cos\phi N(Z_1)$ on cluster 2. Then
\begin{equation}\label{sbgen:hinge-approx}
 \|\ell_\phi-\delta\ell_{0,\phi}\|_2\le C_\ell\delta^2,
 \qquad\|\ell_{0,\phi}\|_2=1/2.
\end{equation}
Indeed, for $m>0$,
$(-(m+\delta Z))_+\le\delta^2Z^2/(4m)$, so its $L^2$ norm is at
most $\sqrt3\delta^2/(4m)$. Apply this to the positive-mean coordinate
in each cluster and to $L_\phi$, whose mean is at least
$r\sin\theta$ and whose noise variance is $\delta^2$. The triangle
inequality gives \eqref{sbgen:hinge-approx}.

For clarity, compare the actual projection $P$ to the auxiliary
projection $P_0$ onto cluster-label constants. A basis of this latter
space is $p_0=\mathbf1_{\{1\}}$, $q_0=r\mathbf1_{\{2\}}$.
The feature map with columns $(p_0,q_0)$ has smallest singular value
$r/\sqrt2$. ReLU's Lipschitz property gives
\[
 \|p-p_0\|_2\le\delta,\quad\|q-q_0\|_2\le\delta,
 \quad \|(p,q)-(p_0,q_0)\|_{\rm op}\le\sqrt2\delta.
\]
For $\delta\le r/4$ both columns remain independent. Projecting their
difference onto the orthogonal complements gives
\[
 \|(I-P)P_0\|\le2\delta/r,\qquad
 \|(I-P_0)P\|\le4\delta/r,
 \qquad\|P-P_0\|\le6\delta/r.
\]
Since $(I-P)q_\phi=-(I-P)\ell_\phi$, we have
\[
 (I-P)q_\phi=\delta R_{0,\phi}+E_\phi,
 \quad \|E_\phi\|_2\le C_R\delta^2,
\]
where $R_{0,\phi}$ equals
$-\sin\phi(N(Z_2)-c_N)$ on cluster 1 and
$-\cos\phi(N(Z_1)-c_N)$ on cluster 2.
The auxiliary label projection is used only to prove this estimate;
the kernel itself still uses the two observable ReLU features.

In the scaled variables write $h=h_{z\delta}(Y)$, and let $h_0$ equal
$\mathbf1\{Z_2+z>0\}$ on cluster 1 and zero on cluster 2.
Since $\mathbb E[Y_1Y_2]=0$,
$\|h\|_2\le H_0$. Also
\begin{equation}\label{sbgen:h-bound}
 |\langle R_{0,\phi},h-h_0\rangle|\le C_h\delta.
\end{equation}
To see this without losing a square root at the moving gate, on cluster 1
the gate is $\mathbf1\{Z_2+z+z\delta Z_1>0\}$ and the prefactor is
$1+\delta Z_1-z\delta^2Z_2$. The prefactor error contributes at most
$v_N\delta(1+z\delta)$ in absolute expectation. Conditional on $Z_1$,
the changed gate is an interval of length $z\delta|Z_1|$, and throughout
that interval
$|N(Z_2)-c_N|\le z+z\delta|Z_1|+c_N$.
Bounding the normal density by $c_N$ bounds the gate contribution by
\[
 c_Nz\delta\left((z+c_N)\sqrt{2/\pi}+z\delta\right).
\]
On cluster 2, $h=\delta(Z_1-zr-z\delta Z_2)$ times an indicator;
Cauchy--Schwarz bounds the contribution by
$v_N\delta\sqrt{1+z^2(r^2+\delta^2)}$.
Average the two clusters and use $\delta\le\delta_0$ to obtain
\eqref{sbgen:h-bound}.

Finally, direct one-dimensional integration gives
\[
 \langle R_{0,\phi},h_0\rangle
 =\frac{\sin\phi}{2}
     [\varphi(z)-c_N(1-\Phi(z))].
\]
Combining this with the two error estimates proves
\eqref{sbgen:J-expansion}. For $z\ge0$,
\[
 1-\Phi(z)
 =e^{-z^2/2}\int_0^\infty c_Ne^{-t^2/2-zt}\,dt
 \le\tfrac12 e^{-z^2/2},
\]
which proves $F(z)\ge\varphi(z)/2$.
Condition \eqref{sbgen:J-region} leaves at least half the uniform
leading margin, giving \eqref{sbgen:J-margin} after restoring $\mu_1^2$.
\end{proof}

For finite samples the existing moment/inverse-Gram transfer can now
use the explicit margin in \eqref{sbgen:J-margin}; requiring its error
to be smaller than that margin gives positivity. The theorem does not
certify the canonical noise level merely because the canonical
population quadrature has the right sign. Its constants are conservative.

\begin{corollary}[Explicit finite-sample transfer of the sign]
\label{sbgen:J-finite}
Assume the region of Theorem~\ref{sbgen:J-positive}. Let a simultaneous
gated-moment event bound each needed $11,12,22$ entry error by
$e_{11},e_{12},e_{22}$, including the fixed-gate intersections used in
the kernel and any clipping correction. Set
\[
 g=\mu_1^2r^2/8,\quad
 B_2=\mu_1^2[(1+r^2)/2+2\delta^2],\quad s_\tau=\sqrt{1+\tau^2},
\]
\[
 e_G=\max(e_{11}+e_{12},e_{12}+e_{22}),\quad
 e_v=\sqrt{(e_{11}+e_{12})^2+(e_{12}+e_{22})^2},
\]
\[
 e_w=\sqrt{(e_{11}+\tau e_{12})^2+(e_{12}+\tau e_{22})^2},
 \qquad e_c=e_{11}+(1+\tau)e_{12}+\tau e_{22}.
\]
If $e_G<g$, then uniformly on $I_\theta$,
\[
 |J_{12}^{(n)}-J_{12}^{\rm pop}|\le
 \epsilon_J:=e_c+
 \frac{e_v s_\tau B_2+B_2e_w+e_ve_w}{g-e_G}
 +\frac{s_\tau B_2^2e_G}{g(g-e_G)}.
\]
In particular $\epsilon_J<\mu_1\sigma\sin\theta F(z)/4$ certifies
the required empirical positivity. This is a deterministic implication
on the existing moment event, with no extra probability allocation.
\end{corollary}
\begin{proof}
The feature-map estimate in the proof above gives the population Gram
lower eigenvalue at least $g$. The entry errors imply the displayed
operator error $e_G$, vector errors $e_v,e_w$, and scalar error $e_c$
in $J=c-v^\top G^{-1}w$; expand the two linear factors in each gated
moment. Cauchy--Schwarz gives $\|v\|\le B_2$,
$\|w\|\le s_\tau B_2$.
Now $\|G_n^{-1}\|\le1/(g-e_G)$ and
$\|G_n^{-1}-G^{-1}\|\le e_G/[g(g-e_G)]$.
Expand $v_n^\top G_n^{-1}w_n-v^\top G^{-1}w$, retaining the product
of the two vector errors, to get the bound. Combine with
\eqref{sbgen:J-margin}.
\end{proof}

\section{An explicit joint concentration bound for the creation gap}
\label{sbgen:joint-section}

Fix a deterministic $\tau>0$. The creation threshold has a useful exact
simplification. Define the four scalar observables
\[
 d(x)=p(x)^2,\quad c(x)=p(x)x_2,\quad
 h(x)=p(x)(x_1-\tau x_2)\alpha_\tau(x),
\]
\[
 l(x)=(\tau x_1+x_2)(x_1-\tau x_2)\alpha_\tau(x),
 \qquad\alpha_\tau=\mathbf1\{\tau x_1+x_2>0\}.
\]
For either population or empirical means, write these means as
$D,C,H,L$. If $D,H>0$, then
\begin{equation}\label{sbgen:gap-affine}
 b_*=C/D,\qquad b_{\rm lift}(\tau,m)=L/H-m\tau,
 \qquad\Delta(\tau,m)=C/D-L/H+m\tau.
\end{equation}
Indeed, the target part of the source middle numerator factors as $l$,
and its $m$-dependent part is exactly $-m\tau h$.
In particular sampling error in this gap is independent of $m$ at fixed
$\tau$. No union bound over progress values is necessary.

For a rigorous bounded-variable implementation, use independent
radially clipped samples $X^R=X\mathbf1\{\|X\|\le R\}$. All population
constants in the next proposition refer to this clipped distribution.
On the existing event that every sample has norm at most $R$, the
empirical clipped and original moments agree. Do not condition the
independent concentration proof on that envelope event.

\begin{proposition}[Joint influence concentration with ratio remainder]
\label{sbgen:joint-gap}
Take $n$ independent samples from each cluster, $N=2n$, and suppose the
clipped population moments obey $D,H>0$. Set $b_*=C/D$, $\ell_*=L/H$ and
define five scalar random variables on a clipped observation:
\[
 Y_D=d/D,\quad Y_H=h/H,\quad
 U=(c-b_*d)/D,\quad V=(l-\ell_*h)/H,\quad Z=U-V.
\]
Their mixture expectations are $1,1,0,0,0$, respectively. For any of
these variables let
\[
 v_Y=\tfrac12\operatorname{Var}_1(Y)+\tfrac12\operatorname{Var}_2(Y),
 \qquad |Y|\le M_Y.
\]
Valid deterministic bounds are
\[
 M_D=R^2/D,\qquad M_H=\sqrt{1+\tau^2}\,R^2/H,
\]
\[
 M_U=R^2(1/2+|b_*|)/D,
\]
\[
 M_V=\frac{R^2}{H}
       \left(\frac{1+\tau^2}{2}+|\ell_*|\sqrt{1+\tau^2}\right),
 \qquad M_Z=M_U+M_V.
\]
Given failure budget $\delta_g\in(0,1)$, put
\[
 L_g=\log(10/\delta_g),\qquad
 a_Y=\sqrt{\frac{2v_YL_g}{N}}+\frac{4M_YL_g}{3N}.
\]
If $a_D,a_H<1$, then with probability at least $1-\delta_g$,
simultaneously for every real $m$,
\begin{equation}\label{sbgen:gap-bound}
 |\Delta_n(\tau,m)-\Delta(\tau,m)|
 \le a_Z+\frac{a_Da_U}{1-a_D}+\frac{a_Ha_V}{1-a_H}.
\end{equation}
The leading variance is the variance of the \emph{joint} variable
$Z=U-V$. Separate concentration of the two ratios is not used.
\end{proposition}
\begin{proof}
The bounds follow from $p\le\|x\|$, $|px_2|\le\|x\|^2/2$,
$|h|\le\sqrt{1+\tau^2}\|x\|^2$, and
$|l|\le(1+\tau^2)\|x\|^2/2$. The last bound follows from the
orthogonality of $(\tau,1)$ and $(1,-\tau)$.

Center each observation by its own cluster expectation. Each centered
variable is bounded by $2M_Y$, and the variance of the sum is $Nv_Y$.
The bounded-variable Bernstein inequality gives failure at most
$2e^{-L_g}$ for error exceeding $a_Y$. This version follows directly
from the tail bound
\[
 \mathbb P(|\bar Y-\mathbb EY|>t)
 \le2\exp\left[-\frac{Nt^2}{2(v_Y+2M_Yt/3)}\right].
\]
A union bound over five variables gives the stated probability.

Write $d_n=D_n/D-1$, $h_n=H_n/H-1$, and let $U_n,V_n,Z_n$ be the
empirical means of the corresponding variables. Exactly,
\[
 \Delta_n-\Delta
 =\frac{U_n}{1+d_n}-\frac{V_n}{1+h_n}
 =Z_n-\frac{d_nU_n}{1+d_n}+\frac{h_nV_n}{1+h_n}.
\]
The simultaneous bounds and positive denominators give
\eqref{sbgen:gap-bound}. Formula \eqref{sbgen:gap-affine} makes the
same event valid for every $m$.
\end{proof}

For a finite deterministic list of middle angles, allocate $\delta_g$
over that list. An angle chosen from the same data without such a
uniform event needs additional control. The gap budget must use the
clipped population gap or include the population clipping correction;
the displayed bound is not silently centered at the unclipped gap.
The variances are explicit Gaussian integrals of bounded observables,
so these are data-model constants rather than trajectory assumptions.

\begin{corollary}[Explicit clipping correction]
\label{sbgen:gap-clipping}
For unclipped equal-noise planar SIM with $0<\mu_2\le\mu_1$, take
$R>\mu_1$, $t_R=(R-\mu_1)/\sigma$, and set
\[
 A_R=[2\mu_1^2+2\sigma^2(t_R^2+2)]e^{-t_R^2/2}.
\]
Let $D_\infty,C_\infty,H_\infty,L_\infty$ be the unclipped population
moments, $b_\infty=C_\infty/D_\infty$,
$\ell_\infty=L_\infty/H_\infty$, and suppose
$D_\infty>A_R$, $H_\infty>s_\tau A_R$.
Then the population gap changes under clipping by at most
\[
 \epsilon_{\rm clip}=
 \frac{(1/2+|b_\infty|)A_R}{D_\infty-A_R}
 +\frac{[(1+\tau^2)/2+|\ell_\infty|s_\tau]A_R}
          {H_\infty-s_\tau A_R}.
\]
The probability of any clipped observation among the $2n$ samples is
at most $2n e^{-t_R^2/2}$. On the intersection with the event of
Proposition~\ref{sbgen:joint-gap}, the original empirical gap is at
least its unclipped population value minus
$\epsilon_{\rm clip}$ and the right side of \eqref{sbgen:gap-bound}.
Use an already funded data-envelope event, or account for this envelope
failure explicitly; do not add it to the theorem budget without allocation.
\end{corollary}
\begin{proof}
The norm of either population mean is at most $\mu_1$. Thus
$\|X\|>R$ implies $\|Z\|>t_R$ for a standard planar normal $Z$.
Since $\|X\|^2\le2\mu_1^2+2\sigma^2\|Z\|^2$, the exact planar
normal tail and tail second moment give
\[
 \mathbb E[\|X\|^2\mathbf1_{\{\|X\|>R\}}]\le A_R.
\]
Therefore the absolute changes of $D,C,H,L$ are at most
$A_R,A_R/2,s_\tau A_R,(1+\tau^2)A_R/2$.
Apply the elementary ratio identity
$C_R/D_R-C_\infty/D_\infty
=[C_R-b_\infty D_R]/D_R$, and similarly for $L/H$.
This proves the correction. The tail probability and a union bound
give the envelope claim.
\end{proof}

\section{What this closes, and the next dynamical obligation}
The transition-kernel sign now has an analytic nonempty high-SNR region,
and the joint creation-gap sampling allowance has an explicit formula
retaining covariance cancellation and the second-order ratio error.
The whole-network deficit has an exact regulator equation.  The new
strong-cone lemma sharpens its damping clock to the angular-measure bound
\eqref{sbgen:sharp-clock}, separates output- and input-side signed forcing,
and reduces output-side positivity to the two-feature margin
\eqref{sbgen:regulator-adm}.  Proposition~\ref{sbgen:first-hit} then turns
generation into first hitting of the moving middle threshold, rather than
convergence to $b_*$.

The complete noisy creation theorem is still open.  Under the corrected late
architecture, the remaining trajectory work is to prove the gate-aware
wrong-sign and $D_E=O(e)$ mismatch bounds, verify the input-side sign condition
\eqref{sbgen:u-positive-condition} (or a sharper signed substitute), control
the adverse complement in \eqref{sbgen:complement-forcing}, and show that the
mass-weighted clock reaches the perturbed creation/progress corridor before
the asymmetric B2 and weak-feature guards expire.
The left endpoint must remain negative at that first hit, while the existing
right endpoint and staircase conditions must hold simultaneously.  The
progress ceiling is now chosen from the charged high-SNR feasibility window;
no fixed recommendation such as $m_{\max}\simeq0.50$ is load-bearing.
Subsequent weak capture, tracking, and projected mass remain separate
specialization obligations.  General noisy reversal-sign propagation belongs
to the separate open extension described in
Remark~\ref{rem:T3-scope-A}, not to this creation dependency chain.


% September 15 signed-gap continuation; general theorem still open.
\section{A direct signed-gap certificate without a transverse-ratio barrier}
\label{sbgap:scope}

This continuation corrects the proposed groupwise ratio argument and
retains the off-diagonal dissipation in the actual relative-growth rate.
It proves deterministic identities, a counterexample to an unrestricted
ratio-face assertion, and sufficient signed-gap and generation conditions.
It does not discharge these conditions from random initialization.

\subsection{Normalization and the exact feedback identity}

All training samples lie in the plane $\mathcal S$, with balanced neuron
norms. Fix a partition $[h]=J_1\sqcup J_2$ and write $q=3-p$ on $J_p$.
All derivatives below hold almost everywhere along the selected absolutely
continuous gradient-flow trajectory. Use
\[
 S_p=\frac12\sum_{J_p}(\|u_i\|^2+\|w_i\|^2),\qquad
 T_p=\frac12\sum_{J_p}(v_{iq}^2+z_{iq}^2),\qquad
 S=S_1+S_2,\quad T=T_1+T_2.
\]
Thus the transverse energy in the manuscript's signed-energy section is
$2T$, not $T$. Put $\tau_p=T_p/S_p$, $\rho_p=S_p/S$,
$\theta=T_1/T$, and $r=\sqrt{T/S}$ when the denominators are positive.
Let $\delta_i=v_i-z_i$ and
\[
 a_{\perp,p}=\frac{\sum_{J_p}(\|u_{i,\perp}\|^2+
                  \|w_{i,\perp}\|^2)}{2S_p},\qquad
 d_{\parallel,p}=\frac{\sum_{J_p}\delta_{ip}^2}{2S_p}.
\]
For the raw longitudinal matrix entry $M_{pp,p}=\sum_{J_p}z_{ip}v_{ip}$,
\begin{equation}\label{sbgap:efficiency}
 \epsilon_p:=\frac{M_{pp,p}}{S_p}
       =1-\tau_p-a_{\perp,p}-d_{\parallel,p}.
\end{equation}
This follows by expanding $(v_{ip}-z_{ip})^2$. It concerns a raw matrix
entry, not the noisy-data projected progress coordinate.

Let $A_i=\mathbb E_n[(x-f_\parallel)x^\top\alpha_i]$ and
$\dot z_i=A_iv_i$, $\dot v_i=A_i^\top z_i-P_i$.
Define
\[
 L_p=\sum_{J_p}A_{i,pp}z_{ip}v_{ip},\qquad
 Q_p=\sum_{J_p}A_{i,qq}z_{iq}v_{iq},
\]
\[
 C_p=\sum_{J_p}(A_{i,qp}v_{ip}z_{iq}
                       +A_{i,pq}z_{ip}v_{iq}),\qquad
 U_p=\langle f_{J_p,\perp},f_\perp\rangle_n,\qquad
 P_p=\sum_{J_p}v_{iq}P_{iq}.
\]
Then
\begin{equation}\label{sbgap:groups}
 \frac{\dot S_p}{2}=L_p+Q_p+C_p-U_p,\qquad
 \dot T_p=2Q_p+C_p-P_p.
\end{equation}
Indeed, differentiation of the output norm gives the first identity
by neuronwise balance, and direct differentiation of $T_p$ gives the
second. Write $L=L_1+L_2$, $Q=Q_1+Q_2$, $C=C_1+C_2$.

As in the manuscript, let $f_D$ and $f_O$ be the diagonal and off-diagonal
planar forward maps. Set
\[
 F=\|f_O\|_n^2,\qquad K=\langle x-f_D,f_O\rangle_n,\qquad
 U=\|f_\perp\|_n^2,\qquad g=\frac{\dot S}{2S}.
\]
The exact equality $C=K-F$ follows by expanding the off-diagonal residual
terms. With the manuscript's signed rate, in the present normalization,
\[
 \kappa_\perp=\frac{Q}{T}-\frac{F}{2T},
\]
one obtains the following identity.

\begin{lemma}[Retained-feedback signed gap]\label{sbgap:identity}
Whenever $S,T>0$,
\begin{equation}\label{sbgap:gap-exact}
 g-\kappa_\perp
 =\frac{L}{S}
 -(1-T/S)\frac{Q}{T}
 +\left(\frac1{2T}-\frac1S\right)F+\frac KS-\frac US.
\end{equation}
In particular, if $2T\le S$, the displayed $F$ term is nonnegative.
\end{lemma}
\begin{proof}
Sum the mass identities to get $g=(L+Q+K-F-U)/S$ and subtract
$Q/T-F/(2T)$. This proves the equality without discarding any signed term.
\end{proof}

\subsection{Why the proposed ratio barrier is not automatic}

\begin{proposition}[An outward ratio face with a positive grading gap]
\label{sbgap:counterexample}
For noiseless data $3e_1,2e_2$ with equal weights, there is a balanced
three-neuron state with strict training gates, a cluster-2-only weak group,
zero mismatch, positive
longitudinal products, and $\tau_1=\tau_2$, such that
\[
 \frac{d}{dt}\log(\tau_1/\tau_2)>0,\qquad
 g-\kappa_\perp>1.
\]
Thus the structural hypotheses just listed do not imply invariance of
$\tau_1\le\tau_2$, and a positive actual grading gap does not require that
particular face to point inward.
\end{proposition}
\begin{proof}
All quantities below describe an initial state for the smooth gradient
flow in its strict activation cell; it is not asserted to be reached
from the paper's prescribed equal-radius random initialization.
Set
\[
 a=\frac12,\quad e=\frac1{100},\quad s=\frac1{10},\quad
 \ell=e^2,\quad t=\frac{e^2}{1+e^2},\quad S_2=\frac1{1000},
\]
\[
 t_\ell=\frac{\ell^2}{1+\ell^2},\quad
 t_s=\frac{s^2}{1+s^2},\quad
 S_\ell=S_2\frac{t_s-t}{t_s-t_\ell},\quad S_s=S_2-S_\ell.
\]
Both masses are strictly positive. Take $u_i=w_i$ and
\[
 v_1=\sqrt a(1,e),\quad
 v_2=\sqrt{\frac{S_\ell}{1+\ell^2}}(-\ell,1),\quad
 v_3=\sqrt{\frac{S_s}{1+s^2}}(-s,1),
\]
with $J_1=\{1\}$ and $J_2=\{2,3\}$. The first neuron is active on both
samples and the other two only on the second. Here
$S_1=a(1+e^2)$, $T_1=ae^2=tS_1$, and $T_2=tS_2$.
Put
\[
 b=c=ae,\quad d=ae^2,\quad
 f=-\frac{S_\ell\ell}{1+\ell^2}-\frac{S_s s}{1+s^2},
 \quad h=\frac{S_\ell}{1+\ell^2}+\frac{S_s}{1+s^2}.
\]
With $\lambda_1=9/2$, $\lambda_2=2$ and
$R=\lambda_2(1-d-h)$, the exact group quantities are
\[
 L_1=\lambda_1(1-a)a,\quad Q_1=Rd,\quad
 C_1=-\lambda_1c^2-\lambda_2(b+f)b,
\]
\[
 L_2=Rh,\quad Q_2=0,\quad C_2=-\lambda_2(b+f)f.
\]
Let $\mathcal A_p=(L_p+Q_p+C_p)/S_p-(2Q_p+C_p)/(2T_p)$.
Exact rational substitution gives
\[
 \mathcal A_1-\mathcal A_2
   =-\frac{438243077997}{6626125026248}<-\frac3{50}.
\]
Since $(\log(\tau_1/\tau_2))'=-2(\mathcal A_1-\mathcal A_2)$,
the face is outward. For the independent grading check use
\[
 F=\lambda_1c^2+\lambda_2(b+f)^2,\quad
 g=\frac{L_1+L_2+Q_1+C_1+C_2}{S_1+S_2},\quad
 \kappa_\perp=\frac{Q_1}{T_1+T_2}-\frac{F}{2(T_1+T_2)}.
\]
The same rational substitution gives
$1.876<g-\kappa_\perp<1.878$, proving the claim.
\end{proof}

\begin{remark}
This is a counterexample to an implication between proposed hypotheses,
not a counterexample to the intended random-initialization theorem.
Extra trajectory structure may restore a ratio barrier. The direct
certificate below does not require one.
\end{remark}

\subsection{A corrected finite-data certificate}

Suppose $A_{i,pp}\ge\alpha_p$ and $A_{i,qq}\le\beta_p$ on $J_p$,
where $\beta_p\ge0$. Keep the mismatch charges
\[
 E_L=\frac12\sum_p\sum_{J_p}(A_{i,pp}-\alpha_p)\delta_{ip}^2,\qquad
 E_Q=\frac12\sum_p\sum_{J_p}[-A_{i,qq}]_+\delta_{iq}^2.
\]
Both are nonnegative. Direct polarization yields
\begin{equation}\label{sbgap:polarization}
 L\ge\sum_p\alpha_p S_p\epsilon_p-E_L,\qquad
 Q\le\beta_1T_1+\beta_2T_2+E_Q.
\end{equation}
For the first inequality subtract $\alpha_p z_{ip}v_{ip}$ and use
$z_{ip}v_{ip}=(z_{ip}^2+v_{ip}^2-\delta_{ip}^2)/2$.
For the second apply the same expansion and retain the negative part of
$A_{i,qq}$ against $\delta_{iq}^2$. In particular, an upper bound on
$[A_{i,qq}]_+$ alone does not justify dropping $E_Q$.

Let $p(x)=(x_1)_+$, $q(x)=(x_2)_+$,
$D_1=\|p\|_n^2>0$, $D_2=\|q\|_n^2>0$, and define
\[
 c_{21}^O=\frac{\langle p,f_{O,2}\rangle_n}{D_1},\qquad
 c_{12}^O=\frac{\langle q,f_{O,1}\rangle_n}{D_2},\qquad
 F_- =D_1(c_{21}^O)^2+D_2(c_{12}^O)^2.
\]
Separate scalar projections in the two output coordinates give $F\ge F_-$.
There is no orthogonality assumption about the input features $p,q$ here.
Also set $m_\times=\mathbb E_n|x_1x_2|$ and
$\lambda_n=\|\mathbb E_nxx^\top\|_{\rm op}$.

\begin{theorem}[Direct signed-gap certificate]\label{sbgap:certificate}
At times with $0<T/S\le1/2$,
\begin{align}
 g-\kappa_\perp\ \ge\ 
 &\sum_p\rho_p\alpha_p\epsilon_p
 -(1-r^2)\{\theta\beta_1+(1-\theta)\beta_2\}
 +(1-2r^2)\frac{F_-}{2T}-\mathcal E,\label{sbgap:certificate-eq}\\
 \mathcal E:={}&\frac{E_L}{S}+(1-r^2)\frac{E_Q}{T}
        +2(1+S)m_\times r+\lambda_n h\varepsilon^2.\nonumber
\end{align}
If the right side is at least $\gamma>0$ throughout a reached interval,
then the forced comparison of Corollary~\ref{sbsign:clock} applies.  In
particular, for
\[
 r_\perp:=\frac{Y}{\sqrt S}=\sqrt{\frac{2T}{S}},
\]
one has
\[
 D^+r_\perp
 \le-\gamma r_\perp
 +\sqrt2(1+S)m_\times+\Pi/\sqrt S.
\]
Thus the certificate supplies a forced scalar decay inequality rather than
an invariant ordering of $\tau_1,\tau_2$; no such ordering is assumed.
\end{theorem}
\begin{proof}
The signed-energy proof, in the present normalization, gives
$|K|\le2(1+S)m_\times\sqrt{ST}$.
The stacked ambient bound and Cauchy--Schwarz give
$U=\|f_\perp\|_n^2\le\lambda_n S\|W_\perp\|_F^2
\le\lambda_n S h\varepsilon^2$.
Use these estimates, \eqref{sbgap:polarization}, and $F\ge F_-$ in
\eqref{sbgap:gap-exact}; the latter replacement is allowed because
$1-2T/S\ge0$. The relative comparison follows by differentiating
$\sqrt{T/S}$, equivalently rescaling the manuscript's
$\sqrt{2T/S}$ comparison.
\end{proof}

\begin{corollary}[A certificate independent of transverse mass share]
\label{sbgap:theta-free}
Suppose $c_{21}^O\ge\zeta\sqrt{S_1T_1}$ with $\zeta\ge0$. Define
\[
 B_\theta=(1-2r^2)\frac{D_1\zeta^2S_1}{2}
                  -(1-r^2)(\beta_1-\beta_2).
\]
Then
\begin{equation}\label{sbgap:theta-free-eq}
 g-\kappa_\perp\ge
 \sum_p\rho_p\alpha_p\epsilon_p
 -(1-r^2)\beta_2+\min(0,B_\theta)-\mathcal E.
\end{equation}
This eliminates $\theta$ without a ratio barrier.
\end{corollary}
\begin{proof}
Since $T_1/T=\theta$, the retained feedback is at least
$(1-2r^2)D_1\zeta^2S_1\theta/2$. The remaining $\theta$-dependence in
\eqref{sbgap:certificate-eq} is the affine term $\theta B_\theta$.
Its minimum on $[0,1]$ is $\min(0,B_\theta)$.
\end{proof}

\begin{remark}[The direct certificate succeeds at the outward face]
For the state in Proposition~\ref{sbgap:counterexample}, take the exact
diagonal residual bounds, so $E_L=E_Q=0$, and retain only
$D_1(c_{21}^O)^2=\lambda_1c^2$. Corollary~\ref{sbgap:theta-free}
then gives the exact rational lower bound
$g-\kappa_\perp>1.376$, even after minimizing independently over
$0\le\theta\le1$. This demonstrates that the replacement can succeed where
the proposed ratio-face argument fails.
\end{remark}

\begin{remark}[A nonempty numerical certificate box, not a trajectory region]
For example, suppose simultaneously
\[
 \rho_1\ge\frac23,\quad \epsilon_1\ge\frac{49}{50},\quad
 \alpha_1\ge\frac94,\quad \rho_2\alpha_2\epsilon_2\ge0,\quad
 \beta_1\le2,\quad \beta_2\le\frac1{20},
\]
\[
 r^2\le\frac1{100},\quad D_1\ge\frac{22}{5},\quad
 \zeta^2\ge\frac45,\quad S_1\ge\frac25,\quad
 \mathcal E\le\frac1{20}.
\]
The preceding theorem gives
\[
 g-\kappa_\perp\ge
 \frac{147}{100}-2
 +\frac{49}{50}\frac{(22/5)(4/5)(2/5)}2-\frac1{20}
 =\frac{687}{6250}>0.
\]
Here one may bound the transverse penalty by
$1/20+(39/20)\theta$ and retain at least
$(4312/6250)\theta$ of feedback before minimizing in $\theta$.
These are explicit sufficient constants for the deterministic certificate.
No assertion is made that they, the creation guards, or the canonical
experiment's parameters have been jointly realized by the trajectory.
\end{remark}

\subsection{Where the needed feedback can come from}

\begin{lemma}[Strong-family polarization and weak-gate suppression]
\label{sbgap:feedback-source}
Assume every member of $J_1$ is active on all training points with $p>0$.
Suppose $v_{i1}>0$ and
$0<\ell\le v_{i2}/v_{i1}\le u$ on $J_1$. Set
\[
 V=\sum_{J_1}v_{i1}^2,\quad Z=\sum_{J_1}v_{i2}^2,\quad
 \Delta_q=\sum_{J_1}(v_{i2}-z_{i2})^2,\qquad
 k_{\ell,u}=\frac{2\sqrt{\ell u}}{\ell+u}.
\]
Let
\[
 \gamma_i=\frac{\mathbb E_n[p^2\alpha_i]}{D_1},\qquad
 N=\sum_{J_2}\gamma_i[-z_{i2}v_{i1}]_+.
\]
Then
\begin{equation}\label{sbgap:feedback-source-eq}
 c_{21}^O\ge \sqrt V\,(k_{\ell,u}\sqrt Z-\sqrt{\Delta_q})-N.
\end{equation}
If all $J_2$ neurons are inactive on cluster 1, then
\begin{equation}\label{sbgap:weak-suppression}
 N\le 2\gamma_W\sqrt{S_2T_2},\qquad
 \gamma_W=\frac{\mathbb E_{2,n}[p^2]}{2D_1}
       \le\frac{s_{11}^{(2)}}{2D_1}.
\end{equation}
If the inactivity condition holds only on a subset, apply this bound on
that subset and charge the remaining negative products explicitly.
\end{lemma}
\begin{proof}
Since $p x_1=p^2$, expanding $f_{O,2}$ gives
$c_{21}^O=\sum_i\gamma_i z_{i2}v_{i1}$.
The strong-family purity makes $\gamma_i=1$ on $J_1$.
For weights proportional to $v_{i1}^2$, write $t_i=v_{i2}/v_{i1}$.
The elementary inequality
$t_i^2\le(\ell+u)t_i-\ell u$ implies
\[
 (\mathbb E t_i)^2\ge
       \frac{4\ell u}{(\ell+u)^2}\mathbb E t_i^2.
\]
Indeed the quadratic
$m^2-\frac{4\ell u}{(\ell+u)^2}((\ell+u)m-\ell u)$ is a square.
Consequently $\sum v_{i1}v_{i2}\ge k_{\ell,u}\sqrt{VZ}$.
Subtract the mismatch term, bounded by $\sqrt{V\Delta_q}$, and the
adverse complement to get \eqref{sbgap:feedback-source-eq}.
Under cluster-1 inactivity, $0\le\gamma_i\le\gamma_W$.
Finally
$\sum_{J_2}|z_{i2}v_{i1}|
\le\sqrt{\sum z_{i2}^2\sum v_{i1}^2}\le2\sqrt{S_2T_2}$.
\end{proof}

For example, lower bounds $V/S_1\ge v_-$, $Z/T_1\ge z_-$,
an upper bound $\Delta_q/T_1\le d_*^2$, and
$N/\sqrt{S_1T_1}\le n_*$ give the usable coefficient
\[
 \zeta=\sqrt{v_-}(k_{\ell,u}\sqrt{z_-}-d_*)-n_*\ge0,
\]
provided $k_{\ell,u}\sqrt{z_-}-d_*\ge0$.
This is a quantitative polarization requirement. It must be proved on the
chosen strong group; balance alone does not supply it.

\begin{corollary}[Feedback supplied by a positive branch tube]
\label{sbgap:branch-feedback}
In the setting of Lemma~\ref{sbgap:feedback-source}, suppose the held
strong group has
\[
 0<\phi_-\le\phi_i\le\phi_+<\pi/2,\qquad
 0<a_z\le z_{i2}/v_{i2}\le b_z,\qquad
 \sum_{J_1}\|u_{i,\perp}\|^2\le a_{\rm in}S_1,\quad a_{\rm in}<1.
\]
Put
\[
 k_*=\frac{2\sqrt{\tan\phi_-\tan\phi_+}}
                      {\tan\phi_-+\tan\phi_+},\qquad
 \zeta_0=a_z k_*\cos\phi_+
                    \sqrt{\frac{2(1-a_{\rm in})}{1+b_z^2}}.
\]
Then
\[
 c_{21}^O\ge\zeta_0\sqrt{S_1T_1}-N.
\]
If the weak group is inactive on cluster 1, $S_2/S_1\le\mathfrak r_{21}$,
and $\tau_2\le\bar\tau_2$, the sufficient coefficient
\begin{equation}\label{sbgap:branch-zeta}
 \zeta=\zeta_0-
 \frac{2\sqrt2\,\gamma_W\mathfrak r_{21}\sqrt{\bar\tau_2}}
             {\sin\phi_-\sqrt{1-a_{\rm in}}}\ge0
\end{equation}
satisfies the hypothesis of Corollary~\ref{sbgap:theta-free}.
Here $\mathfrak r_{21}$ is a bound on the group-mass ratio.
\end{corollary}
\begin{proof}
The positive transverse comparison gives
$\sum z_{i2}v_{i1}\ge a_z\sum v_{i2}v_{i1}
\ge a_z k_*\sqrt{VZ}$ and
$T_1\le(1+b_z^2)Z/2$.
Balance and the input ambient bound give
$V\ge\cos^2\phi_+(1-a_{\rm in})S_1$, proving the first statement.
Also
$T_1\ge Z/2\ge\frac12\sin^2\phi_-(1-a_{\rm in})S_1$.
Insert this and $T_2=\tau_2S_2$ into
$N\le2\gamma_W\sqrt{S_2T_2}$ to prove \eqref{sbgap:branch-zeta}.
\end{proof}

\begin{remark}[Continuation versus entry]
For the $c_{21}^O$ channel, the positive branch tube and transverse component
comparison can be consumed directly; no separately propagated
transverse-coherence variable is necessary when these stronger pointwise
guards are available.  The $c_{12}^O$ channel is different: retaining it
through B2 requires the own-tilt coherence interface of
Subsection~\ref{sbcoh:scope}.
The weak cancellation is suppressed by the cross-noise moment
$\gamma_W\le s_{11}^{(2)}/(2D_1)$.
This is a post-localization continuation result. If B1 itself requires
relative grading to establish entry, this corollary cannot be used to
justify that earlier step circularly. A certified subfamily also cannot
stand in for all of $J_1$ without accounting for the omitted terms.
\end{remark}

\begin{lemma}[Raw entries versus projected strong progress]
\label{sbgap:progress}
Let $B$ be a cone-pure strong family, so $p\alpha_i=p$ on $B$. Then
\[
 m_B^{\rm eff}:=\frac{\langle p,f_{B,1}\rangle_n}{D_1}
   =M_{11,B}^{\rm raw}+b_*M_{12,B}^{\rm raw}.
\]
Moreover
\[
 |M_{12,B}^{\rm raw}|
 \le\sqrt{\sum_B z_{i1}^2\sum_B v_{i2}^2}.
\]
Thus a progress-share hypothesis involving $m_B^{\rm eff}$ cannot be
substituted into \eqref{sbgap:efficiency} without this correction.
\end{lemma}
\begin{proof}
On the support of $p$, activation equals $v_{i1}x_1+v_{i2}x_2$.
Its $p$-projection divided by $D_1$ is $v_{i1}+b_*v_{i2}$.
Sum after multiplying by $z_{i1}$ and apply Cauchy--Schwarz.
\end{proof}

\subsection{Installing the cone-uniform input-side bound}

For this subsection let $B$ be a fixed strong family and set
$p=(x_1)_+$, $D=\|p\|_n^2$, $b_*=\langle p,x_2\rangle_n/D$.
Write $y=\mathbb E_n[p f]/D$, $y_*=(1,b_*,0,\ldots,0)$,
$R=x-f-p(y_*-y)$ and $m=y_1$. Define
\[
 \ell_i=\mathbb E_n[p\alpha_i x]/D,\qquad
 k_i^\circ=\alpha_i\ell_i^\top x-\|\ell_i\|^2p.
\]
The exact input-side regulator term is
\[
 \chi_{u,B}=\sum_{i\in B}w_{i2}
 \left[D(1-m)\|\ell_i\|^2w_{i1}
 -D\|\ell_i\|^2\langle w_{i,\perp},y_\perp\rangle
 +\langle w_i^\top R,k_i^\circ\rangle_n\right].
\]

\begin{lemma}[Cone-pure centered-feature refinement]\label{sbgap:u-side}
Suppose $0\le\phi_i\le\bar\phi<\pi/2$ on $B$, and that every training
sample with $x_1>0$ is strictly activated throughout this cone.
Put $L_0^2=1+b_*^2$ and
\[
 k_0^\circ=b_*\mathbf1_{\{x_1>0\}}(x_2-b_*x_1),\qquad
 \mathcal W=\{x_1=0\}\ \cup\
 \{x_1<0,\ x_2>0,\ x_1+\tan\bar\phi\,x_2\ge0\},
\]
\[
 K_B=\|k_0^\circ\|_n+
       \max(|b_*|,|b_*-\tan\bar\phi|)
                        \|x_2\mathbf1_{\mathcal W}\|_n.
\]
Assume $w_{i1}\ge w_->0$, $w_{i2}\ge0$,
$\|w_i\|/w_{i1}\le\beta_B$, $m\le m_{\max}<1$,
$\|R\|_n\le R_*$, $\|W_\perp\|_F\le W_*$, and $\|y_\perp\|\le Y_*$.
Then
\[
 \chi_{u,B}\ge\gamma_u\sum_Bw_{i1}w_{i2},\qquad
 \gamma_u=D L_0^2\left(1-m_{\max}-\frac{W_*Y_*}{w_-}\right)
                                      -\beta_BR_*K_B.
\]
In particular $\gamma_u\ge0$ certifies $\chi_{u,B}\ge0$.
If in addition $\|\delta_i\|\le\eta_{\rm mis}r_i$, then the fallback
\[
 \chi_{u,B}\ge-c_u\Pi_B,\quad
 c_u=[-\gamma_u]_+(1+\eta_{\rm mis})
                       (\sin\bar\phi+\eta_{\rm mis}),\qquad
 \Pi_B=\sum_Br_i^2
\]
is valid.
\end{lemma}
\begin{proof}
Purity implies $p\alpha_i=p$, hence $\ell_i=(1,b_*)$.
Therefore $k_i^\circ-k_0^\circ=(\alpha_i-\mathbf1_{\{x_1>0\}})
(x_1+b_*x_2)$. Its support is contained in $\mathcal W$, apart from
points with $x_1=x_2=0$ where the term vanishes. On its $x_1<0$ part,
$-x_1/x_2\in[0,\tan\bar\phi]$; on $x_1=0$ use
$|x_1+b_*x_2|=|b_*x_2|$. Thus the stated $K_B$ bounds
$\|k_i^\circ\|_n$. Including the endpoint permits compatible boundary
selectors, not only strict gates.
Use
$|\langle w_{i,\perp},y_\perp\rangle|\le W_*Y_*$
and
$|\langle w_i^\top R,k_i^\circ\rangle_n|
\le\beta_Bw_{i1}R_*K_B$, and then factor $w_{i1}w_{i2}\ge0$.
Finally
$w_{i1}w_{i2}\le(1+\eta_{\rm mis})
(\sin\bar\phi+\eta_{\rm mis})r_i^2$ proves the fallback.
\end{proof}

The coordinate envelope supplies the purity hypothesis, for example, if
cluster 1 has $x_1\ge m_1^{\rm data}>0$, $|x_2|\le b_2$ and
$m_1^{\rm data}\cos\bar\phi-b_2\sin\bar\phi>0$, while cluster 2 has
$x_2>0$. The existing stacked bound supplies
$W_*\le\varepsilon\sqrt h$ and
$Y_*\le\varepsilon\sqrt h\sqrt{\lambda_nS_*/D}$ on $S\le S_*$.
If $r_i\ge r_-$ and
$a_B=\cos\bar\phi-\eta_{\rm mis}>0$, one may take
\[
 w_-=r_-a_B,\qquad
 \beta_B=\frac{\sqrt{(1+\eta_{\rm mis})^2+(W_*/r_-)^2}}{a_B}.
\]
These substitutions make the input-side sign a stated compatibility
inequality; they do not establish the holding, sign, or mismatch guards.

\begin{lemma}[Two-sided radial mass rate]\label{sbgap:mass-rate}
For a fixed strong family with $r_i>0$, let
$G_i=\widetilde u_i^\top A_i\widetilde u_i$ and suppose
$G_-\le G_i\le G_+$,
$\|A_i\|_{\rm op}\|\delta_i\|/r_i+\|P_i\|/r_i\le e_r$.
Then
\[
 2(G_--e_r)\Pi_B\le\dot\Pi_B\le2(G_++e_r)\Pi_B.
\]
\end{lemma}
\begin{proof}
Since $\dot v_i=A_i^\top(v_i-\delta_i)-P_i$,
$\dot r_i/r_i=G_i-\delta_i^\top A_i\widetilde u_i/r_i
-\widetilde u_i^\top P_i/r_i$.
Multiply the two bounds by $2r_i^2$ and sum.
\end{proof}

\subsection{An integrated generation deadline}

\begin{lemma}[Longitudinal mass share needs no transverse-ratio barrier]
\label{sbgap:mass-share}
On a reached interval $[a,b]$ with $S_1,S_2>0$, suppose
$\dot S_1/(2S_1)\ge g_{1,-}$ and
$\dot S_2/(2S_2)\le g_{2,+}$ for arbitrary real bounds
$g_{1,-},g_{2,+}$.  Then
\begin{equation}\label{sbgap:mass-share-eq}
 \boxed{
 \frac{S_2(t)}{S_1(t)}
 \le\frac{S_2(a)}{S_1(a)}
       e^{-2(g_{1,-}-g_{2,+})(t-a)}.
 }
\end{equation}
Consequently, for every $H>0$ with $[a,a+H]\subset[a,b]$,
\begin{equation}\label{sbgap:mass-share-horizon}
 \boxed{
 \sup_{t\in[a,a+H]}\frac{S_2(t)}{S_1(t)}
 \le\frac{S_2(a)}{S_1(a)}
 e^{2(g_{2,+}-g_{1,-})_+H}.
 }
\end{equation}
Also
\[
 \rho_1(t)\ge
 \frac1{1+(S_2(a)/S_1(a))e^{-2(g_{1,-}-g_{2,+})(t-a)}}.
\]
\end{lemma}
\begin{proof}
Differentiate $\log(S_2/S_1)$ and integrate.  No sign assumption on
$g_{1,-}-g_{2,+}$ is needed.  The finite-horizon estimate follows by
maximizing the exponential factor.  This uses group mass-growth bounds, not
an ordering of $T_1/S_1$ and $T_2/S_2$.
\end{proof}

\begin{proposition}[First hitting from a lower mass-growth rate]
\label{sbgap:deadline}
On a stopped interval beginning at $a$, let $G=B_M(m)-b$ and suppose
$G(a)=G_a>0$, $\Pi_B(a)=\Pi_a>0$, $\tau_M>0$, and
\[
 G'=-\tau_M m'-\kappa\{G+b_*-B_M(m)\}-\chi.
\]
Assume, before first hitting,
\[
 m'\ge0,\quad b_*-B_M(m)\ge\delta_0\ge0,\quad
 \kappa\ge a_\kappa\Pi_B,\quad
 \chi\ge a_\chi\Pi_B,\quad
 \Pi_B'\ge g_-\Pi_B,
\]
where $a_\kappa,g_->0$ and $C=\delta_0+a_\chi/a_\kappa>0$.
If the guards hold through $a+H$, where
\begin{equation}\label{sbgap:deadline-eq}
 H=\frac1{g_-}\log\left[
 1+\frac{g_-}{a_\kappa\Pi_a}
           \log\left(1+\frac{G_a}{C}\right)\right],
\end{equation}
then $G$ has a first zero no later than $a+H$.
If the strict creation corridor and the remaining gate-cliff guards hold
at that hit, the manuscript's creation certificate applies.
\end{proposition}
\begin{proof}
While $G>0$, both $G$ and $\delta_0$ are nonnegative, so the lower damping
bound gives $G'\le-a_\kappa\Pi_B(G+C)$.
Hence
$G(t)+C\le(G_a+C)\exp(-a_\kappa\int_a^t\Pi_B)$.
The lower mass rate implies
$\int_a^{a+H}\Pi_B\ge\Pi_a(e^{g_-H}-1)/g_-$.
At the stated value of $H$, the resulting upper bound for $G$ is zero.
Continuity proves a hit no later than this time.
No division by $\Pi_B'$ and no upper mass-growth rate are needed.
\end{proof}

\begin{remark}[What is still required]
The counterexample rules out treating the proposed ratio barrier as a
consequence of the listed structural conditions alone.
The direct gap certificate offers a different route: certify longitudinal
growth, transverse residual bounds with mismatch charges, and enough
off-diagonal feedback on an actually reached interval.
Its feedback lemma requires strong-group polarization and controlled
exceptional gates; these have not been deduced here from initialization.
Generation also still needs adverse complement forcing, positive radial
growth, progress monotonicity or a replacement allowance, and B2/weak-fit
guards lasting through the deadline.  If $c_{12}^O$ is load-bearing, the
weak off-diagonal face $t_\chi$ and the finite-sample spread floor of
Subsection~\ref{sbcoh:scope} must last through the same deadline.  The frozen
canonical sizing indicates that weak rotation itself is not the observed
pre-creation bottleneck; the sharper obstacle is the additive field-transfer
and directional mismatch budget.  A $c_{21}^O$-only certificate remains the
fallback if that stronger coherence transfer is too expensive.  Weak capture remains a load-bearing specialization step.  General noisy
predictive reversal is a separate open extension under the scope of
Remark~\ref{rem:scope-A-specialization}; it is not a specialization
predecessor.
\end{remark}


% September 16: coupled c21 continuation; general theorem remains open.
\section{The default feedback channel: coupled ratios and fixed gate endpoints}
\label{sbpos:scope}

This continuation targets $c_{21}^O$ directly. It does not use the optional
own-tilt spread theorem. The input and output ratios are evolved together;
their difference is not inserted as an additive error divided by the noise
scale. The resulting positive-cone entry and feedback statement is a
deterministic continuation theorem. Its residual and deadline hypotheses
still require a common initialized-trajectory certificate.
Section~\ref{sbemp:scope} strengthens holding by retaining empirical
moments in the reference field and tracking both weighted tilt means.

\subsection{Exact input--output ratio dynamics}

For a fixed strong family $B$, write
\[
 \xi_i=\frac{v_{i2}}{v_{i1}},\qquad
 \omega_i=\frac{z_{i2}}{z_{i1}},\qquad
 q_i=\frac{z_{i1}}{v_{i1}},
\]
whenever both longitudinal coordinates are positive. The selected empirical
matrix is $A_i=\mathbb E_n[(x-f_\parallel)x^\top\alpha_i]$, so
$\dot z_i=A_iv_i$ and $\dot v_i=A_i^\top z_i-P_i$.
Every differential assertion is almost everywhere along the absolutely
continuous selected flow.

\begin{lemma}[Coupled ratios without a mismatch remainder]
\label{sbpos:ratios}
Exactly,
\begin{align}
 \dot\xi_i
 &=q_i\{A_{12,i}+A_{22,i}\omega_i
                 -\xi_i(A_{11,i}+A_{21,i}\omega_i)\}
       -\frac{P_{i2}-\xi_iP_{i1}}{v_{i1}},\label{sbpos:input-ratio}\\
 \dot\omega_i
 &=q_i^{-1}\{A_{21,i}+A_{22,i}\xi_i
                 -\omega_i(A_{11,i}+A_{12,i}\xi_i)\}.
 \label{sbpos:output-ratio}
\end{align}
On $|\xi_i|,|\omega_i|\le U$, $v_{i1}\ge v_*>0$, neuronwise balance
and the stacked ambient bounds imply, with
$a_*=h\varepsilon^2/v_*^2<1$,
\begin{equation}\label{sbpos:q-bounds}
 \sqrt{\frac{1-a_*}{1+U^2}}\le q_i
           \le\sqrt{1+U^2+a_*}.
\end{equation}
\end{lemma}
\begin{proof}
Apply the quotient rule to the exact projected equations.
For the second assertion, balance gives
\[
 q_i^2(1+\omega_i^2)=1+\xi_i^2+
       \frac{\|u_{i,\perp}\|^2-\|w_{i,\perp}\|^2}{v_{i1}^2}.
\]
Both ambient squared norms lie in $[0,h\varepsilon^2]$:
input ambient components are constant and output ambient energy is bounded
in the stacked norm. Their difference therefore has magnitude at most
$h\varepsilon^2$. This proves the claimed bounds.
\end{proof}

\begin{lemma}[Output cone from longitudinal matching and one-sided balance]
\label{sbpos:longitudinal-barrier}
Suppose a B2 input holding certificate gives $|\xi_i|\le u$ on $[a,a+H]$.
Let $v_{i1}(a)\ge v_*>0$, $z_{i1}(a)>0$, and $q_i(a)>q_0$ for some
$0<q_0<1$. Set
\[
 a_{\rm in}=\frac{\varepsilon^2}{v_*^2},\qquad
 \bar q=\sqrt{1+u^2+a_{\rm in}},\qquad
 W=\sqrt{\frac{1+u^2+a_{\rm in}}{q_0^2}-1}.
\]
On the corresponding stopped interval assume
$A_{11,i}\ge\alpha>0$, $|A_{12,i}|\le\beta_{12}$,
$|A_{21,i}|\le\beta_{21}$, and $\|P_i\|/v_{i1}\le p_0$.
If
\begin{align}
 \alpha(1-q_0^2)-\beta_{12}u-q_0^2\beta_{21}W-q_0p_0&>0,
 \label{sbpos:longitudinal-face}\\
 g_v:=q_0(\alpha-\beta_{21}W)-p_0&>0,\label{sbpos:radial-face}
\end{align}
then throughout the interval
\[
 q_0<q_i\le\bar q,\qquad |\omega_i|\le W,\qquad
 v_{i1}(t)\ge v_{i1}(a)e^{g_v(t-a)}.
\]
In this estimate $a_{\rm in}$ involves $\varepsilon^2$, not
$h\varepsilon^2$.
\end{lemma}
\begin{proof}
The exact longitudinal ratio equation is
\[
 \dot q_i=A_{11,i}(1-q_i^2)+A_{12,i}\xi_i
                  -q_i^2 A_{21,i}\omega_i+q_iP_{i1}/v_{i1}.
\]
Balance gives the one-sided bound
\[
 q_i^2(1+\omega_i^2)
 \le 1+\xi_i^2+\frac{\|u_{i,\perp}(0)\|^2}{v_{i1}^2}
 \le1+u^2+a_{\rm in}.
\]
The output ambient norm has a favorable sign and was discarded.
While $q_i\ge q_0$ and $v_{i1}\ge v_*$ this yields the stated $\bar q,W$.
At the $q_0$ face the displayed longitudinal equation is bounded below
by \eqref{sbpos:longitudinal-face}. The same strict bound holds in a
small neighborhood, using the continuous bound
$|\omega_i|\le\sqrt{(1+u^2+a_{\rm in})/q_i^2-1}$.
Also $\dot v_{i1}/v_{i1}\ge g_v$ there.
The integrated first-exit argument excludes either longitudinal face and
proves all conclusions. The input holding certificate is a separate
hypothesis of this lemma.
\end{proof}

\begin{remark}[Noise-scale feasibility of the longitudinal face]
For $e=\sigma/\mu_2$, suppose
$u=e s_+$, $a_{\rm in}\le e^2 a_0$,
$\beta_{jk}\le\mu_2^2e b_{jk}$, $p_0\le\mu_2^2e^2p$ and
$\alpha\ge\mu_2^2\alpha_0>0$.
Choose $q_0^2=1-k e^2$. Then
\[
 W\le\frac{e\sqrt{s_+^2+a_0+k}}{q_0},
\]
and a sufficient longitudinal-face condition is
\[
 \alpha_0 k>b_{12}s_++b_{21}\sqrt{s_+^2+a_0+k}+p.
\]
The left side grows linearly in $k$, the right side only as $\sqrt{k}$.
Thus some fixed $k$ works, and then all sufficiently small $e$ satisfy
$q_0>0$ and the radial inequality. This establishes feasibility of this
particular face budget, not of the full theorem region.
The entry condition $q_i(a)>q_0$ and the assumed input holding still have
to be delivered; they must not be inferred from this continuation itself.
\end{remark}

\begin{theorem}[Positive-cone entry and persistence from coupled faces]
\label{sbpos:entry}
Fix $U>0$ and a deterministic interval $[a,a+H]$. At $a$, suppose
$v_{i1},z_{i1}>0$, $v_{i1}\ge v_*$, and
$|\xi_i|,|\omega_i|<U$ for every $i\in B$.
Assume $h\varepsilon^2/v_*^2<1$.
On the stopped box suppose all compatible selected matrices obey
\[
 A_{11,i}\ge\alpha>0,\quad A_{22,i}\ge0,\quad
 0<d_-\le A_{11,i}-A_{22,i}\le d_+,
\]
\[
 0<a_-\le A_{12,i},A_{21,i}\le a_+,\qquad
 \max(|A_{12,i}|,|A_{21,i}|)\le\beta,
 \qquad \frac{\|P_i\|}{v_{i1}}\le p_0.
\]
The redundant upper bounds $a_+,\beta$ permit separate face and quadratic
budgets. Let $q_-,q_+$ be the bounds in \eqref{sbpos:q-bounds}, and set
\[
 q_*=\min(q_-,q_+^{-1}),\quad q^*=\max(q_+,q_-^{-1}),\quad
 p_*=\sqrt{1+U^2}\,p_0,
\]
\[
 A_*=q_*(a_--\beta U^2)-p_*,\qquad
 D_-=q_*d_-,\qquad D_+=q^*(d_++\beta U).
\]
Assume
\begin{equation}\label{sbpos:face-budgets}
 A_*>0,\qquad
 q_*(d_-U-a_+-\beta U^2)>p_*,
 \qquad g_v:=q_-(\alpha-\beta U)-p_0>0.
\end{equation}
Choose $0<\ell<\min(U,A_*/D_+)$ and put
$z_a=\min_{i\in B}\{\xi_i(a),\omega_i(a)\}$.
Define
\[
 H_-=\frac1{D_-}\log\left(1+\frac{D_-[-z_a]_+}{A_*}\right),
\]
\[
 H_+=
 \begin{cases}
 \displaystyle\frac1{D_+}
 \log\frac{A_*-D_+\max(z_a,0)}{A_*-D_+\ell},&z_a<\ell,\\
 0,&z_a\ge\ell.
 \end{cases}
\]
If $H_-+H_+\le H$, then
\begin{equation}\label{sbpos:positive-box}
 \ell\le\xi_i(t),\omega_i(t)<U
 \quad\text{for all }i\in B,\quad
 a+H_-+H_+\le t\le a+H.
\end{equation}
The larger box remains invariant before this entry, and
\begin{equation}\label{sbpos:long-growth}
 v_{i1}(t)\ge v_{i1}(a)e^{g_v(t-a)}.
\end{equation}
\end{theorem}
\begin{proof}
While the box and longitudinal floor hold,
$\dot v_{i1}/v_{i1}\ge q_-(\alpha-\beta U)-p_0=g_v$ and
$\dot z_{i1}=v_{i1}(A_{11,i}+A_{12,i}\xi_i)>0$.
Thus a longitudinal coordinate cannot be the first exit, and the balance
bounds for $q_i$ remain available.

Let $z(t)=\min_{i\in B}\{\xi_i(t),\omega_i(t)\}$.
This finite minimum is absolutely continuous; at almost every time its
derivative is an active coordinate derivative. If $z=\xi_i\le0$, then
$\omega_i\ge z$ and $A_{22,i}\ge0$, hence
\[
 A_{12,i}+A_{22,i}\omega_i
    -z(A_{11,i}+A_{21,i}\omega_i)
 \ge a_--d_-z-\beta U^2.
\]
If $0\le z=\xi_i\le\ell$, the corresponding lower bound is
$a_--(d_++\beta U)z$. The same bounds hold with input and output
interchanged. Consequently
\[
 \dot z\ge A_*-D_-z\quad(z\le0),\qquad
 \dot z\ge A_*-D_+z\quad(0\le z\le\ell).
\]
Scalar comparison gives the displayed negative-to-zero and
zero-to-$\ell$ entry times; if the initial minimum is already positive,
use it in the second comparison.

At an upper coordinate face $U$, its companion is at most $U$, and
the corresponding derivative is at most
$q_*(a_+-d_-U+\beta U^2)+p_*<0$.
At the lower face $-U$ the minimum comparison is strictly inward.
These bounds are uniform on the stopped box and all compatible selectors,
so the integrated first-exit argument also covers gate boundaries.
The minimum comparison at $\ell$ prevents subsequent loss of the floor.
\end{proof}

\begin{remark}[What this removes, and what it requires]
The theorem does not require a noise-scale bound on
$|\delta_{i2}|/r_i$, a small own-tilt spread remainder, or a transverse-ratio
barrier. It instead consumes two signed cross-residual lower bounds, an
upper-face budget, and the existing ambient/radial envelopes.
The uniform $a_-$ hypothesis can be restrictive; a positive lower-face
calculation alone does not establish the upper-face inequality.
The latter must be checked, or replaced by valid directional face bounds,
before applying the theorem through a creation deadline.
\end{remark}

\begin{corollary}[Positive entry using an existing input holding certificate]
\label{sbpos:held-input-entry}
Instead of proving a new upper face for both ratios, apply
Lemma~\ref{sbpos:longitudinal-barrier}. Put $U=\max(u,W)$,
$q_-=q_0$, $q_+=\bar q$. Suppose also $A_{22,i}\ge0$,
$d_-\le A_{11,i}-A_{22,i}\le d_+$ with $d_->0$, and
$A_{12,i},A_{21,i}\ge a_->0$ on the held input cone.
Define $A_*,D_\pm$ as in Theorem~\ref{sbpos:entry}, with
$\beta=\max(\beta_{12},\beta_{21})$.
If $A_*>0$ and
$0<\ell<\min(u,W,A_*/D_+)$, the same minimum comparison proves
$\xi_i,\omega_i\ge\ell$ after the displayed entry time, provided it fits
inside the common interval.
The upper-face inequality in \eqref{sbpos:face-budgets} is unnecessary in
this version: the input upper bound is supplied by B2 and the output
upper bound by the longitudinal barrier. The cross-residual bounds need
only be uniform for test input slopes $|\xi|\le u$, not $|\xi|\le W$.
\end{corollary}
\begin{proof}
The longitudinal lemma supplies the radii, ratio bounds and both upper
envelopes. Apply only the minimum-comparison part of
Theorem~\ref{sbpos:entry}; no new upper-face argument is used.
\end{proof}

\subsection{Cross-residual lower bounds from fixed endpoints}

Assume all cluster-1 samples have $x_1>0$ and are strictly active for
$|\xi|\le U$, while all cluster-2 samples have $x_2>0$.
For cluster 2 define the fixed data moments
\[
 H_+=\mathbb E_{2,n}[x_1x_2\mathbf1_{\{x_1+Ux_2\ge0\}}],
 \quad
 H_-=\mathbb E_{2,n}[x_1x_2\mathbf1_{\{x_1-Ux_2>0\}}],
\]
\[
 H_{\min}=\min(H_+,H_-),\qquad
 Y_U=\mathbb E_{2,n}[x_2(x_1+Ux_2)_+].
\]
The endpoint selectors are chosen in the adverse direction. The symbols
$H_\pm$ in this subsection are moments, not the entry times above.
Write $D=\|(x_1)_+\|_n^2$, $b=\langle(x_1)_+,f_2\rangle_n/D$ and
$R_{12}^{(1)}=\mathbb E_{1,n}[(x_1-f_1)x_2]$.

\begin{lemma}[Fixed-endpoint cross-residual certificate]
\label{sbpos:endpoints}
Suppose $v_{i1},z_{i1}>0$ and $|\xi_i|\le U$ on $B$, and
$M_B=\sum_Bz_{i1}v_{i1}\le M_+$.
Set
\[
 E_1^+=\mathbb E_{2,n}[(f_{B^c,1})_+x_2],\qquad
 E_2^-=\frac12\mathbb E_{2,n}
 [(-f_2)_+|x_1|\mathbf1_{\{|x_1|\le Ux_2\}}].
\]
For every test slope $\xi\in[-U,U]$, with any compatible selector,
\begin{align}
 A_{12}(\xi)&\ge
 \frac12\{R_{12}^{(1)}+H_{\min}-M_+Y_U-E_1^+\},
 \label{sbpos:A12-lower}\\
 A_{21}(\xi)&\ge
 \frac12\{s_{12}^{(1)}+H_{\min}\}-Db-E_2^-.
 \label{sbpos:A21-lower}
\end{align}
Thus bounds for two fixed gate moments and one fixed ReLU moment,
together with explicit learned-output charges, suffice for uniform lower
bounds across the whole slope interval.
\end{lemma}
\begin{proof}
On cluster 2, $x_2>0$. As $\xi$ increases from $-U$ to zero,
the gate adds positive $x_1x_2$ terms; from zero to $U$ it adds negative
terms. The minimum is therefore attained at one of the two adverse endpoint
selectors. This proves the target-moment lower bound.

Pointwise on cluster 2, positivity of $z_{i1}v_{i1}$ and monotonicity of
ReLU give
$f_{B,1}(x)\le M_+(x_1+Ux_2)_+$.
The complement costs at most $E_1^+$, since the multiplier $x_2\alpha$
is nonnegative. Cluster 1 is fully active, proving
\eqref{sbpos:A12-lower}.

Let $\alpha_0=\mathbf1_{\{x_1>0\}}$.
The change $x_1(\alpha_\xi-\alpha_0)$ is nonpositive, supported on the
cluster-2 wedge $|x_1|\le Ux_2$, and has magnitude at most $|x_1|$.
Consequently
\[
 \mathbb E_n[f_2x_1\alpha_\xi]
 =Db+\mathbb E_n[f_2x_1(\alpha_\xi-\alpha_0)]
 \le Db+E_2^-.
\]
Subtract this from the target moment to prove \eqref{sbpos:A21-lower}.
All equalities and inequalities also cover selected gate endpoints.
\end{proof}

\subsection{A dimension-free data bound at the natural noise scale}

\begin{proposition}[Joint fixed-endpoint concentration]
\label{sbpos:concentration}
Let the cluster-2 samples be independent copies of
$X_1=\sigma Z_1$, $X_2=\mu+\sigma Z_2$, with independent standard normals.
Fix deterministic $U\ge0$ and $M\ge0$. Put
\[
 t=\frac{U\mu}{\sigma\sqrt{1+U^2}},\qquad
 G_U=\frac{\mu\sigma}{(1+U^2)^{3/2}}\varphi(t),
\]
\[
 Y_U^{\rm pop}=U(\mu^2+\sigma^2)\Phi(t)
                    +\mu\sigma\sqrt{1+U^2}\varphi(t).
\]
The two signed gate moments have the same mean $G_U$, and the ReLU
moment has mean $Y_U^{\rm pop}$. Define
\[
 V_X=\sigma^2(\mu^2+\sigma^2),\qquad
 V_Y=V_X+U^2(\mu^4+6\mu^2\sigma^2+3\sigma^4),\qquad
 V_Z=(\sqrt{V_X}+M\sqrt{V_Y})^2.
\]
Choose $L>0$, let $p_L=\min(1,4e^{-L^2/2})$, and put
\[
 R_X=\sigma L(\mu+\sigma L),\qquad
 R_Y=(\mu+\sigma L)\{\sigma L+U(\mu+\sigma L)\},\quad
 R_Z=R_X+MR_Y.
\]
For $j\in\{X,Z\}$, set
\begin{equation}\label{sbpos:data-errors}
 e_j=\sqrt{\frac{2V_j\log(8/\delta_s)}n}
       +\frac{4R_j\log(8/\delta_s)}{3n}
       +\sqrt{V_jp_L}.
\end{equation}
With probability at least $1-\delta_s-4ne^{-L^2/2}$, simultaneously
\begin{equation}\label{sbpos:moment-lower}
 H_{\min}\ge G_U-e_X,\qquad
 H_{\min}-MY_U\ge G_U-MY_U^{\rm pop}-e_Z.
\end{equation}
Here $H_\pm,Y_U$ denote the empirical fixed-endpoint moments.
If $\mu>\sigma L$, all cluster-2 samples also have positive second
coordinate on this event.
\end{proposition}
\begin{proof}
For $Z=X_1+UX_2$ the truncated Gaussian second-moment formula gives
\[
 \mathbb E[X_1X_2\mathbf1_{\{Z>0\}}]
 =\frac{\mu\sigma}{\sqrt{1+U^2}}\varphi(t)
  -\frac{U^2\mu\sigma}{(1+U^2)^{3/2}}\varphi(t)=G_U.
\]
The same computation for $X_1-UX_2$ gives the identical mean.
Stein's identity, or conditioning on $Z$, gives
$\mathbb E[X_2Z_+]=\mu\mathbb E[Z_+]+U\sigma^2\mathbb P(Z>0)$,
which is the stated $Y_U^{\rm pop}$.

Consider four fixed observables: the two gated products, and each gated
product minus $M$ times the ReLU product
\[ X_2(X_1+UX_2)_+. \]
Their second moments are bounded respectively by $V_X,V_Z$.
Indeed,
\[
 \mathbb E[X_1^2X_2^2]=V_X,\qquad
 \mathbb E[X_2^2(X_1+UX_2)^2]=V_Y;
\]
the latter bounds the squared ReLU product. The $L^2$ triangle inequality
gives $V_Z$.

Clip each observable to zero outside
$\{|Z_1|\le L,\ |Z_2|\le L\}$. These remain independent across samples,
have ranges bounded by $R_X$ or $R_Z$, and second moments bounded as above.
Bernstein's inequality and a union bound for four two-sided events give
the first two terms in \eqref{sbpos:data-errors}. The mean clipping error
is at most $\sqrt{V_jp_L}$ by Cauchy--Schwarz.
All empirical clipped observables equal their original versions except
on an event of probability at most $4ne^{-L^2/2}$.
Their four means are $G_U,G_U,G_U-MY_U^{\rm pop},G_U-MY_U^{\rm pop}$,
which proves \eqref{sbpos:moment-lower}.
\end{proof}

\begin{remark}[The noise-scale sampling improvement is limited but real]
If $U=(\sigma/\mu)s_+$ with fixed $s_+$ and bounded $M,\sigma/\mu$, then
$V_X,V_Z=O(\mu^2\sigma^2)$ and
$R_X,R_Z=O(\mu\sigma)$ up to the displayed envelope factors.
Thus $e_X/(\mu\sigma)$ and $e_Z/(\mu\sigma)$ do not carry an intrinsic
$\mu/\sigma$ multiplier, nor an ambient-dimension factor.
This is a theorem for four fixed directional observables, not a uniform
learned-field or pairwise-spread theorem. The learned-output charges
$R_{12}^{(1)},E_1^+,b,E_2^-$ are not discharged by concentration.
\end{remark}

\subsection{Joint holding: retain the favorable cross-coordinate term}

The positive $A_{21}$ floor above is sufficient but is not necessary for
holding a positive output ratio. Its true lower face contains
$A_{21}+A_{22}\xi$, not $A_{21}$ alone. This distinction matters near
the regulator target. The following replacement checks the actual faces
and permits different input and output bounds.

\begin{lemma}[Coupled rectangular holding with signed residuals]
\label{sbpos:rectangle}
Fix $0<L_u<U_u$ and $0<L_w<U_w$, and suppose initially
$\xi_i\in(L_u,U_u)$, $\omega_i\in(L_w,U_w)$ and $v_{i1},z_{i1}>0$.
On the stopped rectangle assume $q_-\le q_i\le q_+$, with $q_->0$,
and $\|P_i\|/v_{i1}\le p_0$.
Write
\[
 B_u(\xi,\omega)=A_{12}(\xi)+A_{22}(\xi)\omega
                 -\xi\{A_{11}(\xi)+A_{21}(\xi)\omega\},
\]
\[
 B_w(\xi,\omega)=A_{21}(\xi)+A_{22}(\xi)\xi
                 -\omega\{A_{11}(\xi)+A_{12}(\xi)\xi\}.
\]
For all compatible selectors and network states in the stopped region,
suppose the following bounds hold, with strict positive margins:
\begin{align*}
 \inf_{\omega\in[L_w,U_w]} B_u(L_u,\omega)
       &>p_0\sqrt{1+L_u^2}/q_-,\\
 \sup_{\omega\in[L_w,U_w]} B_u(U_u,\omega)
       &<-p_0\sqrt{1+U_u^2}/q_-,\\
 \inf_{\xi\in[L_u,U_u]} B_w(\xi,L_w)&>0,\\
 \sup_{\xi\in[L_u,U_u]} B_w(\xi,U_w)&<0.
\end{align*}
Assume these estimates hold also in fixed small neighborhoods of the
corresponding faces, and the longitudinal coordinates cannot vanish
before another stopping guard fires. Then the rectangle persists up to
that guard. The input-face extrema require only the two endpoint values
of $\omega$, because $B_u$ is affine in $\omega$.
\end{lemma}
\begin{proof}
Insert each face estimate into the exact coupled equations
\eqref{sbpos:input-ratio}--\eqref{sbpos:output-ratio}.
The ambient input error is at most $p_0\sqrt{1+\xi^2}$ and the output
equation has no such term. Each boundary therefore has an inward velocity
with a strict uniform margin. Integrating on a face neighborhood excludes
a first exit, including for measurable compatible gate selections.
\end{proof}

\begin{proposition}[A nonempty common population face budget]
\label{sbpos:population-faces}
Consider the planar equal-noise population model, a positive-longitudinal
strong block containing every neuron, and $e=\sigma/\mu_2\to0$ at fixed
$\rho=\mu_1/\mu_2=3/2$. Write
\[
 \xi_i=e s_i,\quad \omega_i=e t_i,\quad
 a_i=z_{i1}v_{i1}>0,\quad M=\sum_i a_i,\quad
 \beta=\frac{2Db}{e\mu_2^2}.
\]
Here $M$ is raw strong progress; it is not silently identified with the
network projection $m$. Define
\[
 K(s)=\varphi(s)+s\Phi(s),\qquad
 J(s,r)=\varphi(\min(s,r))+r\Phi(\min(s,r)).
\]
Uniformly on compact positive rectangles and bounded $M$, the exact
population equations have the form
\begin{align}
 \frac{2\dot s_i}{\mu_2^2}
 &=q_i\left\{\varphi(s_i)-\sum_j a_jJ(s_i,s_j)
        +t_i\Phi(s_i)-\rho^2(1-M)s_i\right\}+r_{u,i},
 \label{sbpos:population-input}\\
 \frac{2\dot t_i}{\mu_2^2}
 &=q_i^{-1}\{K(s_i)-\beta-\rho^2(1-M)t_i\}+r_{w,i},
 \label{sbpos:population-output}
\end{align}
where $\sup_i(|r_{u,i}|+|r_{w,i}|)=o(1)$ for balanced planar states.
No within-family concentration is imposed.

Take the rational bounds
\[
 L=0.14,\quad U=0.32,\quad l=0.115,\quad W=0.33,
 \quad M\in[0.49,0.50],\quad\beta\in[0.24,0.33].
\]
For $s_i\in[L,U]$, $t_i\in[l,W]$, the four principal inward face
margins, before multiplication by $q_i$ or $q_i^{-1}$, are at least
\begin{align*}
 c_{u,-}&=\varphi(L)-0.50\{\varphi(L)+U\Phi(L)\}
                 +l\Phi(L)-\rho^2(1-0.49)L>0.0118,\\
 c_{u,+}&=\rho^2(1-0.50)U-\varphi(U)
                 +0.49K(L)-W\Phi(U)>0.0062,\\
 c_{w,-}&=K(L)-0.33-\rho^2(1-0.49)l>0.0108,\\
 c_{w,+}&=0.24+\rho^2(1-0.50)W-K(U)>0.0320.
\end{align*}
Consequently the rectangle is invariant while the two progress guards
hold, provided $q_i\in[1/2,2]$ and
$|r_{u,i}|,|r_{w,i}|\le0.003$ throughout that stopped interval.
This last error condition can also be used for a finite-sample transfer,
but is not proved for the general empirical network here.
\end{proposition}
\begin{proof}
On cluster 1, all bounded-tilt gates are active except on a Gaussian tail
whose moments vanish faster than any power of $e$.
Its leading diagonal moment is
$2A_{11}/\mu_2^2=\rho^2(1-M)+o(1)$.
On cluster 2 write $x_1=e\mu_2Z_1$, $x_2=\mu_2(1+eZ_2)$.
The gate converges to $\mathbf1_{\{Z_1+s>0\}}$ and
\[
 \mathbb E[(Z_1+r)_+\mathbf1_{\{Z_1+s>0\}}]=J(s,r).
\]
Thus $2A_{22}/\mu_2^2=\Phi(s)+o(1)$ and
\[
 \frac{2A_{12}}{e\mu_2^2}
   =\varphi(s)-\sum_j a_jJ(s,s_j)+o(1).
\]
The definition $Db=\mathbb E[(x_1)_+f_2]$ includes the entire axis
output projection. Gate changes in this output projection occur only
on cluster 2, up to the cluster-1 tail, and contribute $O(e^3\mu_2^2)$:
there $f_2=O(e^2\mu_2)$ in Gaussian moment norm and $x_1=O(e\mu_2)$.
It follows that
$2A_{21}/(e\mu_2^2)=\varphi(s)-\beta+o(1)$.
The remainders are uniform: Gaussian tails control the unbounded samples;
on a bounded set the only gate discrepancy lies in a slab of vanishing
width, uniformly over a compact slope interval, and the ReLU features
are Lipschitz. The positive weights sum to bounded $M$, so no factor
depending on the number of neurons is introduced.
Balance gives
$q_i^2=(1+e^2s_i^2)/(1+e^2t_i^2)$.
Substitute these estimates into the coupled ratio identities to obtain
\eqref{sbpos:population-input}--\eqref{sbpos:population-output}.

At the lower input face, $J(L,s_j)=\varphi(L)+s_j\Phi(L)$,
so its sum is at most $0.50\{\varphi(L)+U\Phi(L)\}$.
At the upper input face, $J(U,s_j)=K(s_j)\ge K(L)$,
so its sum is at least $0.49K(L)$.
These give $c_{u,-},c_{u,+}$. Since $K'=\Phi>0$, the output lower
and upper faces are bounded using $s=L$ and $s=U$, respectively.
The numerical lower bounds follow by bounding the alternating Taylor
series for $\exp(-s^2/2)$ and
$\Phi(s)=1/2+\varphi(0)\int_0^s\exp(-x^2/2)\,dx$.
Multiplication by $q_i$ or $q_i^{-1}$ loses at most a factor two;
the smallest surviving margin is greater than $0.0031$, which exceeds
the stated additive allowance. Apply Lemma~\ref{sbpos:rectangle}.
Longitudinal positivity persists since $A_{11}$ is uniformly positive,
the off-diagonal contribution is $O(e^2\mu_2^2)$, and the planar
ambient forcing is zero.
\end{proof}

\begin{remark}[What the joint check does and does not close]
The progress guards are compatible with balanced states, rather than
independent formal coefficients. For example, take every $s_i=t_i=0.24$
and $M=0.495$. Then $\beta\to\rho^2M(0.24)=0.2673$, strictly inside
the stated interval, and $q_i=1$. For sufficiently small noise all four
face inequalities have room. This proves nonemptiness of a local
population holding budget, not of the initialized theorem region.

The cross-output ceiling is part of this result: it must not be replaced
by $b=b_*$, for which $\beta\to\varphi(0)\approx0.39894$.
Likewise, the input/output lower bounds at entry, the lower progress
guard, and persistence of the $\beta$ corridor have not been proved
from random initialization. A finite-sample application must bound the
weak/complement, ambient, finite-noise and sampling contributions inside
the displayed $0.003$ error budget and verify that the generation clock
fits before a progress guard exits. Positivity of these four faces alone
does not supply that elapsed time.

For feedback in this asymmetric rectangle, use
$\ell=e\min(L,l)$ and $U_{\rm box}=e\max(U,W)$ in the next subsection.
An output lower face is protected by $A_{21}+A_{22}\xi$, so a uniform
positive $A_{21}$ assumption is not needed for this holding step.
\end{remark}

\begin{theorem}[Aggregate mean closes the population holding corridor]
\label{sbpos:aggregate-mean}
In the all-strong population setting of Proposition~\ref{sbpos:population-faces},
set
\[
 \bar t=\frac{\sum_i a_it_i}{M},\qquad \pi_i=\frac{a_i}{M}.
\]
Uniformly for $0<M\le1/2$ and the same compact positive rectangle,
\begin{equation}\label{sbpos:mean-equation}
 \frac{2\dot{\bar t}}{\mu_2^2}
       =\sum_i\pi_iK(s_i)-\rho^2\bar t+r_{\rm av},
       \qquad r_{\rm av}=o(1).
\end{equation}
Also $\beta=\rho^2M\bar t+o(1)$. Absorb that latter error into
$r_{w,i}$ in \eqref{sbpos:population-output}, so its principal bracket
becomes $K(s_i)-\rho^2\{M\bar t+(1-M)t_i\}$.

The augmented region
\begin{equation}\label{sbpos:augmented-region}
 0.14\le s_i\le0.32,\qquad 0.115\le t_i\le0.33,\qquad
 0.20\le\bar t\le0.29
\end{equation}
is invariant, starting in its interior, while $0<M\le1/2$, provided
$q_i\in[1/2,2]$ and all three normalized error types
$|r_{u,i}|,|r_{w,i}|,|r_{\rm av}|$ are at most $0.003$.
Thus a separate narrow $\beta$ corridor is unnecessary for this
population holding result. For sufficiently small $e$, the all-strong
balanced population equations satisfy these error bounds uniformly on
this region.
\end{theorem}
\begin{proof}
The longitudinal weight equation is exactly
\[
 \frac{\dot a_i}{a_i}
   =A_{11,i}(q_i+q_i^{-1})
       +e s_iA_{12,i}/q_i+e t_iq_iA_{21,i}
   =\mu_2^2\rho^2(1-M)+o(\mu_2^2).
\]
The leading growth rate is common to all weights. Differentiating their
weighted mean therefore gives
\[
 \dot{\bar t}=\sum_i\pi_i\dot t_i+
  \sum_i\pi_i(t_i-\bar t)\frac{\dot a_i}{a_i}.
\]
The second sum is $o(\mu_2^2)$ because the common term cancels.
On cluster 1 the output is, up to its negligible gate tail,
$eM\bar t\,x_1+O(e^2x_2)$; the cluster-2 contribution to $Db$ is
$O(e^3\mu_2^2M)$. Hence
$\beta=\rho^2M\bar t+o(1)$.
Use $q_i=1+o(1)$, sum \eqref{sbpos:population-output} against
$\pi_i$, and combine the $M\bar t$ and $(1-M)\bar t$ terms.
This proves \eqref{sbpos:mean-equation}, uniformly even for small $M>0$:
only normalized positive weights and bounded logarithmic weight errors
were used.

Use $L,U,l,W$ from the preceding proposition. At the input lower face
its principal bracket is bounded below by
\[
 \varphi(L)+l\Phi(L)-\rho^2L
      +M\{\rho^2L-\varphi(L)-U\Phi(L)\}.
\]
The coefficient of $M$ is negative, so the worst case is $M=1/2$.
At the input upper face the corresponding upper bound has coefficient
$\rho^2U-K(L)>0$, so its worst case is again $M=1/2$.
For the output lower and upper faces use $\bar t\le0.29$ and
$\bar t\ge0.20$, respectively. Since $l<0.29$ and $W>0.20$,
both worst cases are $M=1/2$.
Finally $K(L)\le\sum_i\pi_iK(s_i)\le K(U)$.
The resulting six inward margins are bounded below by
\begin{align*}
 \widetilde c_{u,-}
 &=\tfrac12\varphi(L)+(l-U/2)\Phi(L)-\rho^2L/2>0.0150,\\
 \widetilde c_{u,+}
 &=\rho^2U/2-\varphi(U)+K(L)/2-W\Phi(U)>0.0109,\\
 \widetilde c_{w,-}&=K(L)-\rho^2(0.29+l)/2>0.0172,\\
 \widetilde c_{w,+}&=\rho^2(0.20+W)/2-K(U)>0.0170,\\
 c_{{\rm av},-}&=K(L)-\rho^2(0.20)>0.0228,\\
 c_{{\rm av},+}&=\rho^2(0.29)-K(U)>0.0733.
\end{align*}
The same alternating-series bounds certify these numbers. Even after
losing a factor two on the individual faces and charging $0.003$ to
each normalized equation, every face remains strictly inward.
The simultaneous first-exit argument for the rectangle and its weighted
mean proves invariance. The previously derived uniform $o(1)$ errors
establish the final small-noise assertion.
\end{proof}

\begin{remark}[The remaining bridge is entry, transfer, and duration]
This augmented invariant region resolves the population mean-versus-spread
obstruction without an individual matching assumption. Its hypotheses
still include entry with positive input/output tilts and a controlled
weighted mean. It does not prove that isotropic initialization enters
this region, that the empirical weak/complement terms meet the error
budget, or that creation occurs before $M$ reaches $1/2$.
The population feedback and holding statements therefore remain
restricted results. A generation deadline must use the actual trajectory
and must not infer available time from this invariant-region calculation.
\end{remark}

\subsection{Delivering the default feedback after entry}

\begin{corollary}[A positive feedback coefficient without transverse matching]
\label{sbpos:feedback}
On an interval where $\ell\le\xi_i,\omega_i\le U$ and
$q_-\le q_i\le q_+$ for $i\in B$, as delivered by the preceding
entry or holding results, let
\[
 S_B=\frac12\sum_B(\|u_i\|^2+\|w_i\|^2),\quad
 T_B=\frac12\sum_B(v_{i2}^2+z_{i2}^2),\quad
 a_B=\frac{\sum_B(\|u_{i,\perp}\|^2+\|w_{i,\perp}\|^2)}{2S_B}
       \le\bar a_B<1.
\]
Assume the positive cone is active on every training point with $x_1>0$.
For $D=\|(x_1)_+\|_n^2$ set
$\gamma_i=\mathbb E_n[(x_1)_+^2\alpha_i]/D$ and
$N=\sum_{i\notin B}\gamma_i[-z_{i2}v_{i1}]_+$.
Define
\[
 \kappa_q=\min_{q_-\le q\le q_+}\frac{2q}{1+q^2}.
\]
Then
\begin{equation}\label{sbpos:feedback-coefficient}
 c_{21}^O\ge
 \frac{\ell}{U}\kappa_q
       \sqrt{\frac{1-\bar a_B}{1+U^2}}\sqrt{S_BT_B}-N.
\end{equation}
In particular, an upper bound $N\le n_*\sqrt{S_BT_B}$ gives a positive
coefficient whenever the displayed prefactor exceeds $n_*$.
\end{corollary}
\begin{proof}
Put $H_B=\frac12\sum_B(v_{i1}^2+z_{i1}^2)$.
Since $\omega_i\ge\ell$,
$\sum_Bz_{i2}v_{i1}\ge\ell\sum_Bz_{i1}v_{i1}
\ge\ell\kappa_q H_B$.
Also $T_B\le U^2H_B$ and
$H_B/S_B\ge(1-\bar a_B)/(1+U^2)$.
Purity gives $\gamma_i=1$ on $B$; discard favorable complement terms
and subtract $N$ to obtain \eqref{sbpos:feedback-coefficient}.
\end{proof}

For a weak subset $W\subseteq B^c$ inactive on cluster 1, define
$S_W$ by the same full-energy convention and
$T_W=\frac12\sum_W(v_{i1}^2+z_{i1}^2)$. Then
\[
 N_W\le\sqrt{2}\gamma_W\sqrt{S_WT_W},\qquad
 \gamma_W=\frac{\mathbb E_{2,n}[(x_1)_+^2]}{2D}.
\]
Indeed, balance gives $\sum_Wz_{i2}^2\le S_W$, while
$\sum_Wv_{i1}^2\le2T_W$; Cauchy--Schwarz yields $\sqrt{2}$.
Every remaining neuron belongs to an explicit exception sum in $N$.
If $B$ is only a seed subset, $S_B,T_B$ cannot be substituted for the full
strong-assigned group masses in the signed-gap theorem. Instead the
always-valid feedback bound is
$F\ge D[c_B\sqrt{S_BT_B}-N]_+^2$, with $c_B$ the prefactor above,
and it is this bound that enters the full-network gap identity.

\begin{corollary}[Uniform coefficient with a retained weak complement]
\label{sbpos:uniform-feedback}
Suppose $B^c=W\sqcup E$, the weak set $W$ is inactive on cluster 1,
$S_W/S_B\le R_W$, and
$N_E\le n_E\sqrt{S_BT_B}$ on the post-entry interval.
Then a uniform valid coefficient is
\begin{equation}\label{sbpos:zeta-uniform}
 \zeta_*=
 \frac{\ell}{U}\kappa_q\sqrt{\frac{1-\bar a_B}{1+U^2}}
 -\frac{\sqrt{2}\gamma_W R_W}{\ell}
               \sqrt{\frac{1+U^2}{1-\bar a_B}}-n_E.
\end{equation}
If $\zeta_*>0$, then $c_{21}^O\ge\zeta_*\sqrt{S_BT_B}$ throughout
that interval. No small upper bound on $T_W/S_W$ is required.
In particular, when $\ell=e\ell_0$, $U=eU_0$,
$\gamma_W=O(e^2)$ and $R_W$ stays bounded, the weak charge is $O(e)$
whereas the positive coefficient has a nonzero limit, provided
$\kappa_q$ and $1-\bar a_B$ have positive lower bounds and $n_E$ is
smaller than that limit.
\end{corollary}
\begin{proof}
Both ratios are at least $\ell$, so
$T_B\ge\ell^2 H_B$ and
$\sqrt{S_BT_B}\ge\ell S_B\sqrt{(1-\bar a_B)/(1+U^2)}$.
Use $T_W\le S_W$ in $N_W\le\sqrt{2}\gamma_W\sqrt{S_WT_W}$, divide by this
lower bound, and insert the result in \eqref{sbpos:feedback-coefficient}.
\end{proof}

\begin{corollary}[A regulator clock from longitudinal growth]
\label{sbpos:clock}
Suppose all the preceding holding hypotheses last on an interval beginning
at $t_0$ and $\kappa\ge a_\kappa\Pi_B$, with
$\Pi_B=\sum_B\|v_i\|^2$. Set $V_0=\sum_Bv_{i1}(t_0)^2$.
Then
\[
 \int_{t_0}^{t_0+h}\Pi_B(s)\,ds
 \ge \frac{V_0}{2g_v}(e^{2g_vh}-1).
\]
If the signed first-hit comparison is
$G'\le-a_\kappa\Pi_B(G+C)$ before hitting, $G(t_0)=G_0>0$, $C>0$,
then a sufficient additional time is
\[
 H_{\rm hit}=\frac1{2g_v}
 \log\left[1+\frac{2g_v}{a_\kappa V_0}
                   \log\left(1+\frac{G_0}{C}\right)\right].
\]
The creation and holding guards must last through this time.
\end{corollary}
\begin{proof}
Use \eqref{sbpos:long-growth}, $\Pi_B\ge\sum_Bv_{i1}^2$, and integrate.
Insert the resulting integral lower bound into the signed first-hit
integrating-factor comparison.
\end{proof}

\begin{remark}[Remaining initialization-to-creation obligations]
The new module proves positive entry, a feedback coefficient, and a clock
from explicit common residual and holding bounds. It does not prove that
the initialized flow supplies those bounds.
The remaining quantitative check is simultaneous: signed residual faces,
a longitudinal radius floor, complement charges, and available time must
all coexist on the initialized finite-sample trajectory.
Proposition~\ref{sbpos:population-faces} settles local compatibility of
the four tilt faces in a restricted population corridor. It does not
supply finite-sample entry into that corridor or its duration.
If the upper-face or deadline check fails, the theorem cannot be declared
closed merely because the fixed-endpoint lower moments are positive.
An aggregate signed-moment route remains an alternative if this pointwise
continuation cannot meet those inequalities.
\end{remark}

\section{Empirical holding with two tracked means}
\label{sbemp:scope}

The fixed sample enters the reference field in this section exactly.
Population fluctuations of its cross moments are not charged a second
 time to an additive trajectory error. The remaining errors are explicit
finite-noise terms, complement outputs, ambient forcing, and the variation
of longitudinal weights. The result is a deterministic holding certificate;
entry, probability of feasibility, and the generation deadline remain
separate obligations.

\subsection{Convex empirical kernels and the second mean}

Put $e=\sigma/\mu_2>0$. On cluster 2 write
$x_1=e\mu_2 z$, $x_2=\mu_2 y$, and assume $y>0$ for every sample.
For any compatible selector $\alpha_s$ at $z+sy=0$, define
\[
 h(s)=\mathbb E_{2,n}[y^2\alpha_s],\quad
 g(s)=\mathbb E_{2,n}[zy\alpha_s],\quad
 K(s)=\mathbb E_{2,n}[y(z+sy)_+],
\]
\[
 J(s,r)=\mathbb E_{2,n}[y(z+ry)_+\alpha_s].
\]
Thus $K(s)=g(s)+s h(s)$ and $h(s)$ is a selected subgradient of the
continuous, nondecreasing convex function $K$.

\begin{lemma}[Kernel geometry and a mean-input closure]
\label{sbemp:kernel}
For every $s,r\in\mathbb R$, exactly, including compatible endpoint
selections,
\[
 J(s,r)=
 \begin{cases}K(r),&r\le s,\\
 K(s)+(r-s)h(s),&r\ge s.
 \end{cases}
\]
For fixed $s$, the function $r\mapsto J(s,r)$ is convex and nondecreasing.
Let positive weights $\pi_i$ sum to one, with
$s_i\in[L,U]$, $t_i\in[l,W]$, and write
$\bar s=\sum_i\pi_i s_i$, $\bar t=\sum_i\pi_i t_i$.
Set
\[
 h_0=\tfrac12\{h(L)+h(U)\},\quad \Delta h=h(U)-h(L),\quad
 C_h=\tfrac12\max(W-L,U-l)\Delta h,\quad
 D_h=\tfrac14(U-L)\Delta h.
\]
Assume $l\le U$, $L\le W$. Then
\begin{align}
 \left|\sum_i\pi_i(t_i-s_i)h(s_i)-h_0(\bar t-\bar s)\right|
       &\le C_h,\label{sbemp:covariance}\\
 0\le\mathcal D:=\sum_i\pi_iK(s_i)
       -\sum_{i,j}\pi_i\pi_jJ(s_i,s_j)&\le D_h.
 \label{sbemp:defect}
\end{align}
Moreover, for positive contributors of arbitrary slopes $r_j$ and weights
$\nu_j$ summing to one,
\begin{align}
 \sum_j\nu_jJ(L,r_j)&\le g(L)+h(L)\sum_j\nu_j\max(r_j,L),
 \label{sbemp:lagging-lower}\\
 \sum_j\nu_jJ(U,r_j)&\ge J(U,\textstyle\sum_j\nu_j r_j).
 \label{sbemp:lagging-upper}
\end{align}
If their mean lies in $[L,U]$, the right side of the second bound is
$K(\sum_j\nu_jr_j)$, even if some contributors lie above $U$.
\end{lemma}
\begin{proof}
Since $y>0$, the gates are nested in their slope for arbitrary real
$s,r$; no sign or interval assumption on either slope is used. If $r\le s$,
activation of $(z+ry)_+$ implies that of the test gate. If $r\ge s$,
the feature is linear on the active test gate. This proves the two formulas.
They join the convex graph of $K$ to a supporting line with its selected
slope, proving convexity in $r$.
The bound \eqref{sbemp:covariance} follows by subtracting $h_0$ and using
$|t_i-s_i|\le\max(W-L,U-l)$.

For independent $S,R$ with law $\pi$, symmetrization gives
\[
 \mathcal D=\tfrac12\mathbb E\left[
 K(\max(S,R))-K(\min(S,R))
 -|S-R|h(\min(S,R))\right].
\]
The integrand is nonnegative by convexity and at most
$|S-R|\,|h(S)-h(R)|$ by the other supporting-line inequality.
Consequently $\mathcal D\le\operatorname{Cov}(S,h(S))$.
Cauchy--Schwarz and the variance bounds for variables in intervals of
lengths $U-L$ and $\Delta h$ give $D_h$.
Finally, monotonicity gives $J(L,r)\le K(L)$ for $r<L$; use the exact
linear formula for $r\ge L$. Jensen's inequality for the convex function
$J(U,\cdot)$ proves \eqref{sbemp:lagging-upper}.
\end{proof}

\subsection{An exact empirical reference system}

Let $B$ be a fixed strong family, with $v_{i1},z_{i1}>0$,
$\xi_i=e s_i$, $\omega_i=e t_i$, $a_i=z_{i1}v_{i1}$,
$M=\sum_Ba_i>0$, and $\pi_i=a_i/M$.
All cluster-1 samples are assumed strictly active for the held slopes
$[L,U]$. Set $f_E=f_\parallel-f_B$; this includes every omitted neuron.
Define the fixed cluster-1 moments
\[
 a=\frac{\mathbb E_{1,n}x_1^2}{\mu_2^2},\quad
 b=\frac{\mathbb E_{1,n}x_1x_2}{e\mu_2^2},\quad
 c=\frac{\mathbb E_{1,n}x_2^2}{\mu_2^2},\quad
 \overline{st}=\sum_B\pi_i s_it_i,
\]
and the complement moments
$C_{jk}(s)=2\mathbb E_n[f_{E,j}x_k\alpha_s]/\mu_2^2$.
Here $a$ is a data coefficient, distinct from the neuron weights $a_i$.
Also put
\[
 Z(s)=\mathbb E_{2,n}[z^2\alpha_s],\qquad
 J_z(s,r)=\mathbb E_{2,n}[z(z+ry)_+\alpha_s].
\]

\begin{lemma}[Exact residual and dynamical remainders]
\label{sbemp:exact-transfer}
Write $A=A^\rho$ for the true selected residual matrix. Exactly,
\begin{align*}
 2A_{11}/\mu_2^2&=a(1-M)+E_{11},&
 2A_{12}/(e\mu_2^2)&=B_{12}+E_{12},\\
 2A_{21}/(e\mu_2^2)&=B_{21}+E_{21},&
 2A_{22}/\mu_2^2&=h(s)+E_{22},
\end{align*}
where
\[
 B_{12}=(1-M)b+g(s)-\sum_j a_jJ(s,s_j),\qquad
 B_{21}=b+g(s)-aM\bar t,
\]
\begin{align*}
 E_{11}&=-e^2M\bar s b+e^2 Z(s)
             -e^2\sum_j a_jJ_z(s,s_j)-C_{11}(s),\\
 E_{12}&=-cM\bar s-C_{12}(s)/e,\\
 E_{21}&=-e^2M\overline{st}b
             -e^2\sum_j a_jt_jJ_z(s,s_j)-C_{21}(s)/e,\\
 E_{22}&=c-e^2M\bar t b-e^2M\overline{st}c
             -e^2\sum_j a_jt_jJ(s,s_j)-C_{22}(s).
\end{align*}
Use the normalized time $\tau=\mu_2^2t/2$ and a prime for its derivative.
Then
\begin{align}
 s_i'&=q_iF_i+r_{u,i},&
 F_i&=(1-M)b+g(s_i)-\sum_j a_jJ(s_i,s_j)
                  +t_i h(s_i)-a(1-M)s_i,\label{sbemp:input}\\
 t_i'&=q_i^{-1}G_i+r_{w,i},&
 G_i&=b+K(s_i)-a\{M\bar t+(1-M)t_i\},\label{sbemp:output}
\end{align}
with exact errors
\begin{align*}
 r_{u,i}&=q_i\{E_{12}+t_iE_{22}-s_iE_{11}
             -e^2s_it_i(B_{21}+E_{21})\}
             -\frac{2(P_{i2}-e s_iP_{i1})}{e\mu_2^2v_{i1}},\\
 r_{w,i}&=q_i^{-1}\{E_{21}+s_iE_{22}-t_iE_{11}
             -e^2s_it_i(B_{12}+E_{12})\}.
\end{align*}
All entries on the right are evaluated at $s_i$.
\end{lemma}
\begin{proof}
On cluster 1,
$f_{B,1}=M(x_1+e\bar s x_2)$ and
$f_{B,2}=eM\bar t x_1+e^2M\overline{st}x_2$.
On cluster 2 the corresponding outputs are
$e\mu_2\sum_j a_j(z+s_jy)_+$ and
$e^2\mu_2\sum_j a_jt_j(z+s_jy)_+$.
Substitute into each entry of
$A=\mathbb E_n[(x-f_B-f_E)x^\top\alpha_s]$, remembering that
$\mathbb E_n=(\mathbb E_{1,n}+\mathbb E_{2,n})/2$.
The coupled ratio identities then give both normalized equations.
\end{proof}

\begin{lemma}[Exact directional factorization of the complement remainder]
\label{sbemp:directional-complement}
With $C_{rs}(s)=2\mathbb E_n[f_{E,r}x_s\alpha_s]/\mu_2^2$, define
\[
 \mathcal C_u(s,t)
 :=sC_{11}-\frac{C_{12}}e+est\,C_{21}-tC_{22},
\]
\[
 \mathcal C_w(s,t)
 :=tC_{11}+est\,C_{12}-\frac{C_{21}}e-sC_{22}.
\]
Then exactly
\[
\boxed{
 \mathcal C_u(s,t)
 =-\frac{2}{e\mu_2^2}\,
 \mathbb E_n\!\left[
 (f_{E,1}+etf_{E,2})(x_2-esx_1)\alpha_s
 \right],
}
\tag{directional-Cu}\label{sbemp:directional-Cu}
\]
and
\[
\boxed{
 \mathcal C_w(s,t)
 =\frac{2}{e\mu_2^2}\,
 \mathbb E_n\!\left[
 (etf_{E,1}-f_{E,2})(x_1+esx_2)\alpha_s
 \right].
}
\tag{directional-Cw}\label{sbemp:directional-Cw}
\]
Thus the four entries $C_{11},C_{12},C_{21},C_{22}$ enter the exact individual
ratio errors through two directional functionals, not four independent
absolute charges.
\end{lemma}
\begin{proof}
Expand the two displayed expectations and use
$C_{rs}=2\mathbb E_n[f_{E,r}x_s\alpha_s]/\mu_2^2$.  The four terms in the
first identity are respectively
$-C_{12}/e$, $sC_{11}$, $-tC_{22}$ and $est\,C_{21}$; the second identity is
identical term by term with the corresponding signs.
\end{proof}

\begin{remark}[Why the old whole-output late bound loses the relevant scale]
\label{rem:directional-50}
Any late estimate that bounds the four $C_{rs}$ entries separately before
forming \eqref{sbemp:directional-Cu}--\eqref{sbemp:directional-Cw} destroys
the two directional cancellations.  At $e=1/40$ this can lose essentially
two factors of $1/e$: the previously quoted absolute allowance $50.78$ is
about $1493$ times the narrow-face budget $0.0340027$, while $e^{-2}=1600$.
This is a structural artifact of the relaxation, not a large constant that
should be repaired by sharpening the same norm argument.
\end{remark}

\begin{lemma}[Signed reduction of the amplified input complement]
\label{sbemp:signed-input-complement}
At a fixed time let $B=H\sqcup\mathcal L$ be the fixed positive-longitudinal
family and let $E=B^c$.  Let $W\subset E$ be a fixed cluster--$2$-pure weak
cohort and put $E^\circ=E\setminus W$.  Define
\[
 \mathcal X_{12}^{E,+}(s)
 :=\left[
 \frac{1}{e\mu_2^2}
 \mathbb E_{2,n}[f_{E^\circ,1}x_2\alpha_s]
 \right]_+.
\]
Only neurons with $z_{i1}>0$ can contribute adversely to this quantity.
Partition those indices into
\[
 \mathcal A_+
 =\{i\in E^\circ:z_{i1}>0,\ v_{i1}>0,
       \ i\text{ fully active on cluster }1\},
\]
\[
 \mathcal P_+
 =\{i\in E^\circ:z_{i1}>0,\ v_{i1}>0,
       \ i\text{ not fully active on cluster }1\},
\]
\[
 \mathcal S_+
 =\{i\in E^\circ:z_{i1}>0,\ v_{i1}\le0\}.
\]
Put
\[
 \mathcal Z_{\mathcal A}:=\sum_{i\in\mathcal A_+}z_{i1}r_i,
 \qquad
 \mathsf R_{\mathcal P}:=\sum_{\mathcal P_+}r_i^2,
 \qquad
 \mathsf R_{\mathcal S}:=\sum_{\mathcal S_+}r_i^2,
\]
\[
 \Delta_{\mathcal P}:=\sum_{\mathcal P_+}\delta_{i1}^2,
 \qquad
 \Delta_{\mathcal S}:=\sum_{\mathcal S_+}\delta_{i1}^2.
\]
Let
\[
 K_2
 :=\frac{1}{\mu_2^2}
 \sup_{\|u\|=1,\,s}
 \mathbb E_{2,n}[(u^\top x)_+x_2\alpha_s],
\]
and, on the coordinate envelope, put
\[
 \kappa_{\rm band}
 :=\frac{b_2}{\sqrt{(m_1^{\rm data})^2+b_2^2}}.
\]
Then
\begin{equation}\label{sbemp:X12-exception-bound}
 \boxed{
 \mathcal X_{12}^{E,+}(s)
 \le\frac{K_2}{e}\left[
 \mathcal Z_{\mathcal A}
 +\kappa_{\rm band}\mathsf R_{\mathcal P}
 +\sqrt{(\mathsf R_{\mathcal P}+\mathsf R_{\mathcal S})
              (\Delta_{\mathcal P}+\Delta_{\mathcal S})}
 \right].
 }
\end{equation}
Thus the sign-mismatched class $\mathcal S_+$ is a mismatch-energy charge,
not an independent wrong-output moment.  The genuinely new entrant quantity
is $\mathcal Z_{\mathcal A}$.
\end{lemma}
\begin{proof}
For $z_{i1}\le0$ the cluster--$2$ contribution has favorable sign because
$x_2>0$ and $(v_i^\top x)_+\alpha_s\ge0$.  For $z_{i1}>0$, the definition of
$K_2$ bounds the contribution by $(K_2/e)z_{i1}r_i$.
If $i\in\mathcal S_+$, then
$z_{i1}=v_{i1}-\delta_{i1}>0$ with $v_{i1}\le0$, hence
$z_{i1}\le|\delta_{i1}|$.  If $i\in\mathcal P_+$, failure of full
cluster--$1$ activation and the coordinate envelope imply
$\cos\phi_i\le\kappa_{\rm band}$, so
$z_{i1}\le\kappa_{\rm band}r_i+|\delta_{i1}|$.
Summing and applying Cauchy--Schwarz proves
\eqref{sbemp:X12-exception-bound}.
\end{proof}

\begin{lemma}[Netting the empirical cluster--$1$ cross moment]
\label{sbemp:b-netting}
Let $\mathcal F_1\subset E$ be the complement neurons fully active on
cluster $1$, with raw coefficients
\[
 m_{11,E}=\sum_{i\in\mathcal F_1}z_{i1}v_{i1},\qquad
 m_{12,E}=\sum_{i\in\mathcal F_1}z_{i1}v_{i2}.
\]
Let $f_{E,1}^{\rm part}$ denote the output of the remaining partially-gated
cluster--$1$ complement and set
\[
 \mathcal T_{12}^{\rm band}
 :=\frac{1}{e\mu_2^2}
 \mathbb E_{1,n}[f_{E,1}^{\rm part}x_2].
\]
Then, with the manuscript normalization
$\mathbb E_n=(\mathbb E_{1,n}+\mathbb E_{2,n})/2$,
\begin{equation}\label{sbemp:C12-C1-net}
 \boxed{
 \left.\frac{C_{12}}e\right|_{\mathcal C_1}
 =m_{11,E}b+\frac ce\,m_{12,E}+\mathcal T_{12}^{\rm band}.
 }
\end{equation}
Consequently the combination in the input slope equation is exactly
\begin{equation}\label{sbemp:b-netted-input}
 \boxed{
 (1-M)b+E_{12}
 =(1-M-m_{11,E})b-cM\bar s
 -\frac ce\,m_{12,E}-\mathcal T_{12}^{\rm band}
 -\left.\frac{C_{12}}e\right|_{\mathcal C_2}.
 }
\end{equation}
The fixed empirical moment $b$ should therefore be retained with its sign
in the principal field rather than charged again by absolute value inside
$E_{12}$.
\end{lemma}
\begin{proof}
On cluster $1$, the fully-active complement is linear:
$f_{E,1}^{\rm full}=m_{11,E}x_1+m_{12,E}x_2$.
Use the definitions of $b,c$ and the factor $1/2$ in $\mathbb E_n$ to
obtain \eqref{sbemp:C12-C1-net}; substitute it into
$E_{12}=-cM\bar s-C_{12}/e$.
\end{proof}

\begin{remark}[The retained weak cohort has favorable input sign]
\label{sbemp:weak-E12-sign}
For a matched planar cluster--$2$-pure weak cohort with angles
$\phi_i=\pi-\theta_i$,
\[
 E_{12}^W(s)
 =\sum_{i\in W}r_i^2\cos^2\theta_i
 \left\{\frac{\tan\theta_i}{e}h(s)-g(s)\right\}.
\]
Hence
\[
 E_{12}^W(s)
 \ge m_{11}^W
 \left\{\frac{\tan\theta_G}{e}h(s)-g(s)\right\}.
\]
The cluster--$2$ purity condition gives
\[
 \frac{\tan\theta_G}{e}
 >\frac{b_1}{e\,m_2^{\rm data}}.
\]
For the equal-noise canonical envelope $e=0.075$, $L_s=3$ this lower
ratio is $3.871$, so at $s_\rho\approx0.23419$ the bracket is about
$1.906$.  This is sizing information; the finite-sample theorem uses the
exact displayed signed expression.
\end{remark}

The longitudinal logarithmic weight rate $k_i=a_i'/a_i$ is exactly
\[
 k_i=\{a(1-M)+E_{11}\}(q_i+q_i^{-1})
 +e^2s_i(B_{12}+E_{12})/q_i
 +e^2t_iq_i(B_{21}+E_{21})-\frac{2P_{i1}}{\mu_2^2v_{i1}}.
\]
Therefore the two mean equations are
\begin{align}
 \bar s'&=(1-M)\{b+\langle K\rangle-a\bar s\}
       +\langle(t-s)h(s)\rangle+M\mathcal D+r_s,
       \label{sbemp:mean-input}\\
 \bar t'&=b+\langle K\rangle-a\bar t+r_t,
       \label{sbemp:mean-output}
\end{align}
where brackets denote the $\pi$ average and, exactly,
\[
 r_s=\langle(q-1)F+r_u\rangle+\operatorname{Cov}_\pi(s,k),\qquad
 r_t=\langle(q^{-1}-1)G+r_w\rangle+\operatorname{Cov}_\pi(t,k).
\]
In particular, with $\operatorname{osc}k=\max_i k_i-\min_i k_i$,
\begin{align*}
 |r_s|&\le\max_i|q_i-1|\langle|F|\rangle
       +\langle|r_u|\rangle+(U-L)\operatorname{osc}k/4,\\
 |r_t|&\le\max_i|q_i^{-1}-1|\langle|G|\rangle
       +\langle|r_w|\rangle+(W-l)\operatorname{osc}k/4.
\end{align*}
These follow by differentiating the normalized weights and using the
same bounded-variable covariance inequality as above. No frozen training
mask or negligible within-family spread is assumed.

\begin{lemma}[Wide-cone contraction of the two empirical means]
\label{sbemp:wide-mean-contraction}
Fix a stopped rectangle
\[
 s_i\in[L_0,U_0],\qquad t_i\in[l_0,W_0],\qquad 0<M\le M_+<1,
\]
and let $s_\star\in[L_0,U_0]$ solve
\[
 b+K(s_\star)=a s_\star.
\]
Put
\[
 h_-=h(L_0),\qquad h_+=h(U_0),\qquad
 h_0=\frac{h_-+h_+}{2},\qquad \Delta h=h_+-h_-,
\]
\[
 D_K=\frac{U_0-L_0}{4}\Delta h,\qquad
 C_K=\frac12\max\{W_0-L_0,U_0-l_0\}\Delta h,
\]
and
\[
 A_+=(1-M_+)(a-h_+)+h_0.
\]
Assume
\[
 A_+a-h_0h_+>0.
\]
Define
\[
 \lambda_{\rm wc}
 =\frac{A_++a-\sqrt{(A_+-a)^2+4h_0h_+}}{2}>0,
 \qquad
 \vartheta_{\rm wc}=\frac{h_+}{a-\lambda_{\rm wc}},
\]
and the Perron-weighted mean error
\[
 R_{\rm wc}
 :=\max\left\{
 |\bar s-s_\star|,
 \frac{|\bar t-s_\star|}{\vartheta_{\rm wc}}
 \right\}.
\]
Then, almost everywhere on the stopped interval,
\begin{equation}\label{sbemp:wide-mean-Dini}
 \boxed{
 D^+R_{\rm wc}
 \le-\lambda_{\rm wc}R_{\rm wc}+\mathcal E_{\rm wc},
 }
\end{equation}
where one may take
\[
 \boxed{
 \mathcal E_{\rm wc}
 =\max\left\{
 D_K+C_K+|r_s|,
 \frac{D_K+|r_t|}{\vartheta_{\rm wc}}
 \right\}.
 }
\]
Consequently, for $\tau\ge\tau_0$,
\begin{equation}\label{sbemp:wide-mean-Duhamel}
 R_{\rm wc}(\tau)
 \le e^{-\lambda_{\rm wc}(\tau-\tau_0)}R_{\rm wc}(\tau_0)
 +\int_{\tau_0}^{\tau}
 e^{-\lambda_{\rm wc}(\tau-u)}\mathcal E_{\rm wc}(u)\,du.
\end{equation}
No condition-number prefactor appears.
\end{lemma}
\begin{proof}
Write $e_s=\bar s-s_\star$ and $e_t=\bar t-s_\star$.  By convexity,
\[
 \langle K\rangle=K(\bar s)+\Delta_J,
 \qquad 0\le\Delta_J\le D_K,
\]
and convexity gives a secant slope $\kappa\in[h_-,h_+]$ such that
\[
 K(\bar s)-K(s_\star)=\kappa e_s,
 \qquad \kappa\in[h_-,h_+].
\]
Also
\[
 \langle(t-s)h(s)\rangle
 =h_0(e_t-e_s)+\Delta_C,
 \qquad |\Delta_C|\le C_K.
\]
Since $0\le\mathcal D\le D_K$ and
$(1-M)+M=1$, the two exact mean equations imply
\[
 e_s'
 =-\{(1-M)(a-\kappa)+h_0\}e_s+h_0e_t+\eta_s,
 \qquad |\eta_s|\le D_K+C_K+|r_s|,
\]
\[
 e_t'=\kappa e_s-ae_t+\eta_t,
 \qquad |\eta_t|\le D_K+|r_t|.
\]
Taking absolute values gives comparison with the Metzler matrix
\[
 B_+=
 \begin{pmatrix}
 -A_+&h_0\\ h_+&-a
 \end{pmatrix}.
\]
Its determinant is positive by hypothesis and its Perron eigenvalue is
$-\lambda_{\rm wc}$.  The positive right eigenvector may be normalized as
$(1,\vartheta_{\rm wc})$.  The induced weighted max norm therefore gives
\eqref{sbemp:wide-mean-Dini} directly.  Variation of constants proves
\eqref{sbemp:wide-mean-Duhamel}.
\end{proof}

\begin{lemma}[Modulus bound for the longitudinal-weight oscillation]
\label{sbemp:k-modulus}
On a stopped rectangle $|s_i|\le S_0$, $|t_i|\le T_0$ and
$q_-\le q_i\le q_+$, write
\[
 Q(q)=q+q^{-1},\qquad
 E_{11}(s)=E_{11}^{\rm fix}+\widetilde E_{11}(s),
\]
where $E_{11}^{\rm fix}$ is independent of the test slope.  Set
\[
 Q_{\max}=\max_{[q_-,q_+]}Q,\qquad
 \operatorname{osc}Q=\max_{[q_-,q_+]}Q-\min_{[q_-,q_+]}Q,
\]
\[
 H_{12}^\star=\sup_{|s|\le S_0}|B_{12}(s)+E_{12}(s)|,
 \qquad
 H_{21}^\star=\sup_{|s|\le S_0}|B_{21}(s)+E_{21}(s)|,
\]
and
\[
 p_1^\star=\sup_i\frac{|P_{i1}|}{v_{i1}}.
\]
Then the exact logarithmic rate satisfies
\begin{align}\label{sbemp:k-modulus-bound}
 \operatorname{osc}k
 \le{}&
 \Bigl(|a(1-M)+E_{11}^{\rm fix}|
       +\|\widetilde E_{11}\|_\infty\Bigr)\operatorname{osc}Q
 +Q_{\max}\operatorname{osc}_s\widetilde E_{11}\nonumber\\
 &+2e^2S_0q_-^{-1}H_{12}^\star
 +2e^2T_0q_+H_{21}^\star
 +\frac{4p_1^\star}{\mu_2^2}.
\end{align}
In particular, when $q_i=1+O(e^2)$ uniformly,
$\operatorname{osc}Q=O(e^4)$ because $Q'(1)=0$.  Hence a large
slope-independent cluster--$1$ contribution to $E_{11}$ is multiplied only
by a quadratic balance deviation; it must not be charged as
$Q_{\max}\sup|E_{11}|$.
\end{lemma}
\begin{proof}
Apply
$\operatorname{osc}(XY)\le\|X\|_\infty\operatorname{osc}Y
+\|Y\|_\infty\operatorname{osc}X$
to the first term of the exact formula for $k_i$.  The two $e^2$ cross
terms are each bounded by twice their uniform modulus, and the ambient term
has oscillation at most $4p_1^\star/\mu_2^2$.  This proves
\eqref{sbemp:k-modulus-bound}.
\end{proof}

\subsection{Eight empirical faces and a larger error allowance}

For $x\in[L,U]$ let
\[
 K_{\rm sec}(x)=\frac{U-x}{U-L}K(L)+\frac{x-L}{U-L}K(U).
\]
Choose $L<s_-<s_+<U$, $l<t_-<t_+<W$, and $0<M_+<1$.
For an affine function of $m$, minima and maxima below are taken just
at $m=0,M_+$. Define the eight fixed-data margins
\begin{align*}
 d_{u,-}&=\min_m\bigl[(1-m)\{b+g(L)-aL\}+(l-ms_+)h(L)\bigr],\\
 d_{u,+}&=-\max_m\bigl[(1-m)b+g(U)-mK(s_-)+Wh(U)-a(1-m)U\bigr],\\
 d_{w,-}&=b+K(L)-a\{M_+t_++(1-M_+)l\},\\
 d_{w,+}&=a\{M_+t_-+(1-M_+)W\}-b-K(U),\\
 d_{s,-}&=\min_m(1-m)\{b+K(s_-)-as_-\}
                       +h_0(t_--s_-)-C_h,\\
 d_{s,+}&=-\max_m\bigl[(1-m)\{b+K_{\rm sec}(s_+)-as_+\}
                       +h_0(t_+-s_+)+C_h+mD_h\bigr],\\
 d_{t,-}&=b+K(s_-)-at_-,\\
 d_{t,+}&=at_+-b-K_{\rm sec}(s_+).
\end{align*}
Assume $a>0$. Take the endpoints $L,U$ off the finite set of sample
knots, or require these margins uniformly over their compatible endpoint
selections and nearby face values.

\begin{theorem}[Empirical two-mean holding certificate]
\label{sbemp:holding}
Suppose initially the individual slopes and both weighted means lie
strictly inside
\[
 [L,U]\times[l,W],\qquad
 \bar s\in[s_-,s_+],\quad\bar t\in[t_-,t_+].
\]
On the stopped region let $0<M\le M_+$, $q_-\le q_i\le q_+$,
with positive longitudinal coordinates and $q_->0$.
Suppose $|r_{u,i}|\le\epsilon_u$, $|r_{w,i}|\le\epsilon_w$,
$|r_s|\le\epsilon_s$, $|r_t|\le\epsilon_t$.
If
\[
 q_-\min(d_{u,-},d_{u,+})>\epsilon_u,\qquad
 q_+^{-1}\min(d_{w,-},d_{w,+})>\epsilon_w,
\]
\[
 \min(d_{s,-},d_{s,+})>\epsilon_s,\qquad
 \min(d_{t,-},d_{t,+})>\epsilon_t,
\]
all four intervals persist until a progress, longitudinal, purity, or
error guard expires. The interval endpoints may be selected from the
observed data: this is a deterministic certificate and asserts no
probability of feasibility for such a selection.
\end{theorem}
\begin{proof}
At $s_i=L$ use the exact formula
$\sum_j\pi_jJ(L,s_j)=g(L)+\bar s h(L)$ and $t_i\ge l$.
At $s_i=U$, Jensen gives
$\sum_j\pi_jJ(U,s_j)=\langle K\rangle\ge K(\bar s)\ge K(s_-)$.
These prove the input face bounds. Monotonicity of $K$ and the mean
bounds prove the output faces.
For the mean faces, Jensen and the secant bound give
$K(\bar s)\le\langle K\rangle\le K_{\rm sec}(\bar s)$.
Insert \eqref{sbemp:covariance}--\eqref{sbemp:defect} into
\eqref{sbemp:mean-input}, and use $0<M\le M_+<1$.
Equation~\eqref{sbemp:mean-output} gives the other two faces directly.
Subtract the stated error bounds. All eight faces have a strict inward
margin, so the integrated simultaneous first-exit argument applies.
\end{proof}

\begin{theorem}[Lagging-aware empirical two-mean holding]
\label{sbemp:lagging-holding}
Let the fixed positive-longitudinal family be partitioned once as
\[
 B=H\sqcup\mathcal L.
\]
Write
\[
 M_H=\sum_{i\in H}a_i,\qquad
 M_{\mathcal L}=\sum_{j\in\mathcal L}a_j,\qquad M=M_H+M_{\mathcal L},
\]
\[
 \pi_i^H=\frac{a_i}{M_H},\qquad
 \bar s_H=\sum_{i\in H}\pi_i^Hs_i,\qquad
 \bar t_H=\sum_{i\in H}\pi_i^Ht_i.
\]
Assume the held core satisfies
\[
 s_i\in[L,U],\qquad t_i\in[l,W],\qquad
 \bar s_H\in[s_-,s_+],\qquad \bar t_H\in[t_-,t_+],
\]
while the lagging contributors stay in the same cluster--$1$-pure broad
region but need not lie in the narrow rectangle.  Put
\[
 \bar s_{\mathcal L}
 =M_{\mathcal L}^{-1}\sum_{j\in\mathcal L}a_js_j
\]
when $M_{\mathcal L}>0$.  Define the individual-face lagging charges
\[
 \Lambda_{u,-}
 =h(L)\sum_{j\in\mathcal L}a_j(s_j-s_+)_+,
\]
\[
 \Lambda_{u,+}
 =M_{\mathcal L}
 \bigl[K(s_-)-J(U,\bar s_{\mathcal L})\bigr]_+,
\]
\[
 \Lambda_{w,-}
 =a\sum_{j\in\mathcal L}a_j(t_j-t_+)_+,
 \qquad
 \Lambda_{w,+}
 =a\sum_{j\in\mathcal L}a_j(t_--t_j)_+,
\]
with all four equal to zero when $M_{\mathcal L}=0$.  The exact core-mean
lagging terms are
\[
 \Lambda_s
 =\sum_{j\in\mathcal L}a_j
 \left[\langle K(s)\rangle_H
       -\langle J(s,s_j)\rangle_H\right],
 \qquad
 \Lambda_t
 =a\sum_{j\in\mathcal L}a_j(\bar t_H-t_j).
\]
Write
\[
 \Lambda_{s,-}=[-\Lambda_s]_+,
 \quad \Lambda_{s,+}=[\Lambda_s]_+,
 \qquad
 \Lambda_{t,-}=[-\Lambda_t]_+,
 \quad \Lambda_{t,+}=[\Lambda_t]_+.
\]
Then the exact core-mean equations are
\begin{align}
 \bar s_H'={}&(1-M)\{b+\langle K\rangle_H-a\bar s_H\}
 +\langle(t-s)h(s)\rangle_H+M_H\mathcal D_H
 +\Lambda_s+r_{s,H},\label{sbemp:lagging-mean-s}\\
 \bar t_H'={}&b+\langle K\rangle_H-a\bar t_H
 +\Lambda_t+r_{t,H}.\label{sbemp:lagging-mean-t}
\end{align}
Let $d_{u,\pm},d_{w,\pm},d_{s,\pm},d_{t,\pm}$ be the eight principal
margins above, interpreted with the core intervals and the total mass $M$.
Suppose the directional errors obey
\[
 -\epsilon_{u,-}\le r_{u,i}\le\epsilon_{u,+},\qquad
 -\epsilon_{w,-}\le r_{w,i}\le\epsilon_{w,+},
\]
\[
 -\epsilon_{s,-}\le r_{s,H}\le\epsilon_{s,+},\qquad
 -\epsilon_{t,-}\le r_{t,H}\le\epsilon_{t,+}.
\]
If
\[
 q_-\{d_{u,-}-\Lambda_{u,-}\}>\epsilon_{u,-},\qquad
 q_-\{d_{u,+}-\Lambda_{u,+}\}>\epsilon_{u,+},
\]
\[
 q_+^{-1}\{d_{w,-}-\Lambda_{w,-}\}>\epsilon_{w,-},\qquad
 q_+^{-1}\{d_{w,+}-\Lambda_{w,+}\}>\epsilon_{w,+},
\]
\[
 d_{s,-}>\Lambda_{s,-}+\epsilon_{s,-},\qquad
 d_{s,+}>\Lambda_{s,+}+\epsilon_{s,+},
\]
\[
 d_{t,-}>\Lambda_{t,-}+\epsilon_{t,-},\qquad
 d_{t,+}>\Lambda_{t,+}+\epsilon_{t,+},
\]
then the core rectangle and both core-mean intervals persist until another
stopping guard fires.
\end{theorem}
\begin{proof}
At the lower input face,
\[
 \sum_{i\in H}a_iJ(L,s_i)
 =M_H\{g(L)+\bar s_Hh(L)\},
\]
while Lemma~\ref{sbemp:kernel} gives
\[
 J(L,s_j)\le g(L)+h(L)\max(s_j,L)
 \le g(L)+h(L)\{s_++(s_j-s_+)_+\}
\]
for $j\in\mathcal L$.  This yields exactly $\Lambda_{u,-}$.  At the
upper input face, Jensen gives
\[
 \sum_{j\in\mathcal L}a_jJ(U,s_j)
 \ge M_{\mathcal L}J(U,\bar s_{\mathcal L}),
\]
and the core contribution is at least $M_HK(s_-)$.  This yields
$\Lambda_{u,+}$.  The two output charges follow by comparing each lagging
$t_j$ with $t_+$ or $t_-$.

Averaging the exact input equation over $H$ gives
\eqref{sbemp:lagging-mean-s}: the coefficient identity
\[
 (1-M)+M_H=1-M_{\mathcal L}
\]
leaves precisely
$\sum_{\mathcal L}a_j[\langle K\rangle_H-
\langle J(s,s_j)\rangle_H]$.  Averaging the output equation gives
\eqref{sbemp:lagging-mean-t}.  The eight-face proof of
Theorem~\ref{sbemp:holding} now applies after subtracting only the charge
that can hurt each corresponding face.
\end{proof}

\begin{remark}[Fixed maximal handoff enlargement]
\label{sbemp:maximal-handoff}
The partition $B=H\sqcup\mathcal L$ is fixed at a specified handoff time;
it is not a time-dependent membership rule.  It is advantageous to place in
$\mathcal L$ every additional neuron for which cluster--$1$ purity,
positive longitudinal coordinates, and the required radius/mismatch guards
are simultaneously certified through the continuation horizon.  Qualifying
positive mass is then paid through the one-sided $\Lambda$ charges rather
than returned to the $1/e$-amplified complement.  Later entrants remain a
separate stopping quantity.
\end{remark}

\begin{corollary}[A certified population design point for the two-mean region]
\label{sbemp:design}
In the principal population system take
$a=9/4$, $b=0$, $h=\Phi$, $K(s)=\varphi(s)+s\Phi(s)$, and
\[
 (L,U,l,W)=(0.135,0.337,0.122,0.347),\qquad M_+=1/2,
\]
\[
 [s_-,s_+]=[0.193,0.285],\qquad [t_-,t_+]=[0.208,0.264].
\]
The eight principal margins are respectively
\[
\begin{split}
&0.0344359,\quad0.0343462,\quad0.0358221,\quad0.0344910,\\
&0.0347804,\quad0.0431661,\quad0.0348494,\quad0.0349586,
\end{split}
\]
so
\[
 \boxed{d_{\min}>0.03434.}
\]
Consequently a directional error allowance of $0.030$ is compatible with
this population design whenever the relevant longitudinal factors satisfy
\[
 \min(q_-,q_+^{-1})\ge0.99.
\]
This is a population design point only.  It does not certify the empirical
trajectory errors, the entry clock, or simultaneous feasibility of the
remaining guards.
\end{corollary}
\begin{proof}
Substitution into the eight displayed margin formulas gives the stated
values.  Since $0.99(0.03434)>0.030$, the strict face inequalities retain
positive slack.  The two mean intervals contain the matched root
$s_\rho\approx0.23419$.
\end{proof}

\begin{remark}[Balance and the longitudinal factor]
If
$|\|u_{i,\perp}\|^2-\|w_{i,\perp}\|^2|/v_{i1}^2\le a_*$,
balance supplies the sharper explicit bounds
\[
 q_-^2=\frac{1+e^2L^2-a_*}{1+e^2W^2},\qquad
 q_+^2=\frac{1+e^2U^2+a_*}{1+e^2l^2},
\]
provided the first numerator is positive. In the planar case $a_*=0$;
for $e=0.075$ the preceding box has
$\min(q_-,q_+^{-1})>0.999$.
This improves the factor-two allowance without assuming transverse
input/output matching. Ambient forcing remains in the explicit errors.
\end{remark}

\subsection{Feedback, grading, and local stability}

\begin{corollary}[Mean-weighted feedback and the scope of grading]
\label{sbemp:feedback}
Under the purity and energy conventions of the preceding feedback
module, put $U_0=\max(U,W)$ and assume $\bar t\ge t_->0$.
Then exactly $\sum_Bz_{i2}v_{i1}=eM\bar t$, and therefore
\[
 c_{21}^O\ge
 \frac{t_-}{U_0}\kappa_q
 \sqrt{\frac{1-\bar a_B}{1+e^2U_0^2}}\sqrt{S_BT_B}-N.
\]
Also $T_B/S_B\le e^2U_0^2$ directly.
For a complementary group $E$ with the same transverse-energy convention,
\[
 \frac{T_B+T_E}{S_B+S_E}
 \le e^2U_0^2\frac{S_B}{S_B+S_E}
       +\frac{T_E}{S_B+S_E}.
\]
Thus strong-family grading needs no additional signed-gap theorem after
entry into this box; network grading still needs the complement term.
\end{corollary}
\begin{proof}
The exact cross moment is the definition of the weighted mean.
Use $M\ge\kappa_qH_B$, $T_B\le e^2U_0^2H_B$ and
$H_B/S_B\ge(1-\bar a_B)/(1+e^2U_0^2)$ in the earlier feedback proof.
The grading bounds follow by summing the coordinate squares.
\end{proof}

\begin{remark}[What the local stability threshold means]
In the principal population equations the collective fixed point is
$s_i=t_i=s_\rho$ for every $0\le M\le1$, where
$K(s_\rho)=\rho^2s_\rho$.
For $p=\Phi(s_\rho)$ and $a=\rho^2$, the linearized zero-weighted-mean
modes have matrix
\[
 \begin{pmatrix}-a(1-M)&p\\p&-a(1-M)\end{pmatrix},
\]
whereas the collective mean mode has matrix
\[
 \begin{pmatrix}-Mp-a(1-M)&p\\p&-a\end{pmatrix}.
\]
These use normalized time and fixed leading weights. The first matrix
has eigenvalues $-a(1-M)\pm p$, so the local transverse stability
threshold is $M=1-p/a$. The collective matrix has negative trace and
positive determinant for $0\le M\le1$, because $0<p<a$.
A bound using $\Phi(U)$ is a uniform band criterion, not the exact
threshold at $s_\rho$. Neither calculation rules out every nonlinear
invariant box beyond that threshold. Changing one finite-weight neuron
also changes both means; its self contribution must not be silently
omitted when interpreting an individual Jacobian.

For completeness, on the collective diagonal the fixed-point equations
are $0=(1-M)\{K(s)-as\}+(t-s)\Phi(s)$ and $0=K(s)-at$.
They imply $K(s)=as$ and $t=s$. Linearizing and projecting onto the mean
and zero-mean modes yields the displayed matrices. In particular,
input/output matching at the fixed point follows from the coupled
principal dynamics; it was not assumed as a trajectory property.
\end{remark}

\subsection{What the actual T1.9 clock supplies}

Let $a_0=t_0+\tau_{\rm gap}$ be the start of the fine angular comparison
in T1.9, $\omega=\omega_{1,\star}^{\rm post}>0$, and
$\beta_\theta=(1-\vartheta)\eta_1^u$.
On its certified horizon the comparison used there gives
\[
 |\phi_i(a_0+\tau)|\le
 \beta_\theta+(|\phi_i(a_0)|-\beta_\theta)e^{-\omega\tau},
 \qquad |\psi_i-\phi_i|\le D_{\rm lock}.
\]
The first bound is the integrated inequality
$D^+|\phi_i|\le-\omega|\phi_i|+\omega\beta_\theta$.
To certify only the symmetric outer bounds
$|\xi_i|\le eU$, $|\omega_i|\le eW$, put
\[
 \Theta=\min\{\arctan(eU),\arctan(eW)-D_{\rm lock}\}.
\]
If $\Theta>\beta_\theta$, a sufficient additional time for neuron $i$ is
\[
 H_i^{\rm outer}=
 \begin{cases}
 0,&|\phi_i(a_0)|\le\Theta,\\
 \displaystyle\omega^{-1}\log
 \frac{|\phi_i(a_0)|-\beta_\theta}{\Theta-\beta_\theta},
       &|\phi_i(a_0)|>\Theta.
 \end{cases}
\]
It is usable only if $a_0+H_i^{\rm outer}\le H_\star$ and all the T1.9
budgets hold. The angular decay rate is its certified $\omega$, not an
unverified replacement by $\lambda_1-\lambda_2$.

More generally, at time $a_0+\tau$ every neuron failing these two outer
bounds must lie in the fixed-at-$a_0$ set
\[
 \mathcal L(\tau)=\{i:\ |\phi_i(a_0)|>
       \beta_\theta+(\Theta-\beta_\theta)e^{\omega\tau}\}.
\]
If its full balanced neuron energies obey
$S_i(a_0+\tau)\le E(\tau)S_i(a_0)$, then the absolute raw longitudinal
mass of such lagging neurons is at most
\[
 \sum_{\text{outer-bound failures}}|z_{i1}v_{i1}|
       \le E(\tau)\sum_{i\in\mathcal L(\tau)}S_i(a_0),
\]
since $|z_{i1}v_{i1}|\le S_i$ by balance.
This is a deterministic lagging-mass interface, not a numerical estimate
of that mass from initialization. If $\Theta\le\beta_\theta$, the cited
angular comparison does not certify the desired outer radius. Even when
it does, symmetric upper bounds do not imply the positive lower faces or
either weighted-mean interval required by the holding theorem. Those are
additional entry conditions; the T1.9 horizon may also end exactly at its
locking time.


\begin{lemma}[Two-stage sign crossing under the moving mean envelope]
\label{sbemp:sign-crossing}
Work in normalized time $\tau=\mu_2^2t/2$ on the fixed family $B$ and
assume the exact equations of Lemma~\ref{sbemp:exact-transfer}.  Suppose on
a stopped interval
\[
 -S(\tau)\le s_i\le S(\tau),\qquad t_i\ge-D,
\]
\[
 \bar s\le s_\star+\widehat R(\tau),\qquad
 \bar t\le s_\star+\vartheta_{\rm wc}\widehat R(\tau),
 \qquad 0<M\le M_+,
\]
and $q_-\le q_i\le q_+$.  Here $s_\star$ solves
$b+K(s_\star)=as_\star$.  Assume one-sided errors
\[
 r_{u,i}\ge-\epsilon_u^-(\tau),\qquad
 r_{w,i}\ge-\epsilon_w^-(\tau).
\]
Define the exact input crossing margin
\begin{equation}\label{sbemp:input-cross-margin}
 a_{u,\times}(\tau)
 :=\inf_{\substack{0\le m\le M_+\\-S(\tau)\le s\le0}}
 \left[(1-m)\{b+g(s)-as\}
 +\{-D-m(s_\star+\widehat R(\tau))\}h(s)\right],
\end{equation}
and
\[
 \Gamma_u(\tau)=q_-a_{u,\times}(\tau)-\epsilon_u^-(\tau).
\]
If $\Gamma_u>0$ and
\begin{equation}\label{sbemp:input-cross-clock}
 \int_{\tau_0}^{\tau_u}\Gamma_u(v)\,dv\ge S(\tau_0),
\end{equation}
then
\[
 s_i(\tau)\ge0\qquad(i\in B,\ \tau\ge\tau_u)
\]
for as long as the stated guards persist.

After input crossing, define
\begin{equation}\label{sbemp:output-cross-margin}
 a_{w,\times}(\tau)
 :=\inf_{0\le m\le M_+}
 \left[b+K(0)-am\{s_\star+
 \vartheta_{\rm wc}\widehat R(\tau)\}\right],
\end{equation}
\[
 \Gamma_w(\tau)=q_+^{-1}a_{w,\times}(\tau)-\epsilon_w^-(\tau).
\]
If $\Gamma_w>0$ and
\begin{equation}\label{sbemp:output-cross-clock}
 \int_{\tau_w^0}^{\tau_w}\Gamma_w(v)\,dv\ge D
 \qquad(\tau_w^0\ge\tau_u),
\end{equation}
then $t_i(\tau)\ge0$ for all $i\in B$ and $\tau\ge\tau_w$ while the
guards persist.  Both crossing clocks run concurrently with the mean
comparison; they are not added to the mean-entry time.
\end{lemma}
\begin{proof}
Let $z_s=\min_i s_i$.  At almost every time at which $z_s=s_i\le0$,
every $s_j\ge s_i$, so the global kernel identity gives exactly
\[
 \sum_j a_jJ(s_i,s_j)=Mg(s_i)+M\bar s\,h(s_i).
\]
Hence
\[
 F_i=(1-M)\{b+g(s_i)-as_i\}+(t_i-M\bar s)h(s_i)
 \ge a_{u,\times}(\tau).
\]
While $z_s<0$, $D^+[-z_s]_+\le-\Gamma_u(\tau)$, and
\eqref{sbemp:input-cross-clock} forces the minimum to zero.  At the zero
face the same strict lower bound prevents return to the negative half-line.

Afterward let $z_t=\min_i t_i$.  At an active negative minimum,
$s_i\ge0$ and monotonicity of $K$ imply
\[
 G_i\ge b+K(0)-aM\bar t+a(1-M)(-t_i)
 \ge a_{w,\times}(\tau).
\]
While $z_t<0$, $D^+[-z_t]_+\le-\Gamma_w(\tau)$.  Starting from the assumed lower
bound $t_i\ge-D$, \eqref{sbemp:output-cross-clock} proves output crossing
and the strict zero-face margin prevents return.
\end{proof}

\begin{remark}[The pre-crossing output lower bound is a separate guard]
T1.9 alone gives $t_i\ge-T(\tau)$ before input crossing. The smaller
$-\tan(D_{\rm lock})/e$ bound follows from its angle-gap estimate only
after input nonnegativity. Lemma~\ref{sbflux:pre-cross-output} derives a
concurrent output lower envelope and retains favorable forcing through
\eqref{sbflux:reflected-output}; use that envelope in the input margin
unless the stronger pre-crossing bound has been proved separately.
\end{remark}

\begin{remark}[Canonical crossing margins]
\label{sbemp:canonical-crossing}
In the principal Gaussian system with
$a=9/4$, $b=0$, $M_+=1/2$, $s_\star=s_\rho\approx0.23419$ and
$h(0)=1/2$, the input zero-face margin before errors is
\[
 \boxed{0.14092-\frac12D-0.25\widehat R,}
\]
equivalently $0.14092+0.5t_i-0.25\widehat R$ when the actual lower
$t_i$ is retained.  With $\vartheta_{\rm wc}\approx0.54$, the output
zero-face margin is
\[
 \boxed{0.135479-0.6075\widehat R.}
\]
These are roughly $4.5$ times the $0.030$ holding allowance, which is why
sign crossing is used as the robust bridge rather than narrow-box holding.
\end{remark}

\begin{corollary}[Piecewise moving-cone mean contraction across input crossing]
\label{sbemp:piecewise-wide-mean}
Let the coarse T1.9 input-angle envelope in normalized time be
\[
 \Theta(\tau)=\beta_\theta+\theta_{\rm gap}e^{-\nu\tau},
 \qquad \nu=\frac{2\omega_{1,\star}^{\rm post}}{\mu_2^2},
\]
and put
\[
 S(\tau)=\frac{\tan\Theta(\tau)}e,\qquad
 D=\frac{\tan D_{\rm lock}}e,\qquad
 T(\tau)=\frac{\tan(\Theta(\tau)+D_{\rm lock})}{e}.
\]
Then $T(\tau)\ge S(\tau)+D$.  Before the input crossing time $\tau_u$ use
\[
 [-S,S]\times[-T,T].
\]
With
\[
 \Delta h_-=h(S)-h(-S),\qquad
 D_-=\frac{S}{2}\Delta h_-,\qquad
 C_-=\frac{S+T}{2}\Delta h_-,
\]
define
\[
 \mathcal E_-(\tau)
 =\max\left\{D_-+C_-+|r_s|,
 \frac{D_-+|r_t|}{\vartheta_{\rm wc}}\right\}.
\]
After input crossing use
\[
 [0,S]\times[-D,T].
\]
With
\[
 \Delta h_+=h(S)-h(0),\qquad
 D_+=\frac{S}{4}\Delta h_+,\qquad
 C_+=\frac{T}{2}\Delta h_+,
\]
define
\[
 \mathcal E_+(\tau)
 =\max\left\{D_++C_++|r_s|,
 \frac{D_++|r_t|}{\vartheta_{\rm wc}}\right\}.
\]
Assume a common certified wide-cone contraction rate $\lambda_{\rm wc}$
for the moving rectangles.  Then
\begin{equation}\label{sbemp:piecewise-Duhamel}
 \boxed{
 \widehat R(\tau)
 =R_0e^{-\lambda_{\rm wc}\tau}
 +\int_0^\tau e^{-\lambda_{\rm wc}(\tau-v)}
 \left[\mathbf1_{\{v<\tau_u\}}\mathcal E_-(v)
 +\mathbf1_{\{v\ge\tau_u\}}\mathcal E_+(v)\right]dv
 }
\end{equation}
is a valid deterministic mean envelope.

In the symmetric Gaussian specialization $T=S$ before crossing and
$D=0$ after crossing,
\[
 D_-+C_-=3S\{\Phi(S)-1/2\},
 \qquad
 D_++C_+=\frac34S\{\Phi(S)-1/2\},
\]
so input sign crossing quarters the geometric forcing exactly.  Output
crossing may continue in parallel; it is not needed for this $4\times$
mean-forcing improvement.
\end{corollary}

\begin{remark}[Combined entry clock]
\label{sbemp:combined-clock}
Let $\tau_{\rm mean}$ be the first time the envelope
\eqref{sbemp:piecewise-Duhamel} enters the chosen mean window, let
$\tau_w$ be the output sign-crossing time, and let $\tau_\ell$ be the
positive-floor time supplied by the coupled-ratio entry theorem.  The narrow
holding handoff is
\[
 \boxed{\tau_{\rm hand}=\max\{\tau_{\rm mean},\tau_w,\tau_\ell\}.}
\]
The clocks are composed concurrently, not added.  The remaining timing
audit is the strict ordering
\[
 \boxed{
 \tau_{\rm hand}
 <\min\{\tau_{M_+},\tau_{\rm T1.9}\}.
 }
\]
The formerly undefined additional branch deadline is resolved in
Section~\ref{sbguard:scope}: entry uses these early guards; subsequent
feedback must satisfy $t_g<t_{\rm hold}$, and generation has its own
reached continuation exit.
\end{remark}

\begin{remark}[Current status and remaining initialized-trajectory obligations]
\label{sbemp:remaining-entry-obligations}
The deterministic entry architecture is now explicit:
\[
 \text{T1.9 envelope}
 \longrightarrow
 \text{input sign crossing}
 \longrightarrow
 \text{wide-cone mean contraction}
 \longrightarrow
 \text{mean entry}
 \longrightarrow
 \text{lagging-aware holding},
\]
with output crossing and the coupled-ratio positive-entry theorem running in
parallel.  This is a chain of proved deterministic implications; it is not
yet a theorem from random initialization.

The first genuinely new pre-entry stopping quantity is
$\mathcal Z_{\mathcal A}$ in Lemma~\ref{sbemp:signed-input-complement}, or
equivalently the same-sign positive cluster--$2$ wrong-coordinate output
$\mathcal X_{12}^{E,+}$.  At the handoff one should enlarge the fixed
lagging family maximally subject to simultaneous purity, longitudinal,
radius, and mismatch guarantees.  The remaining same-sign mass that enters
later must be controlled by a fixed-family entrant-mass or angular-flux
estimate.  The sign-mismatched class is already reduced to
$\sqrt{\mathsf R_{\mathcal S}\Delta_{\mathcal S}}$.

The remaining quantitative entry terms are the one-sided trajectory errors,
the T1.9 lagging charges $\Lambda_{u,\pm},\Lambda_{w,\pm},\Lambda_s,
\Lambda_t$, the ambient quantities $p_{u,+}$ and $p_1^\star$ outside the
planar case, and any output-crossing charge not bypassed by the fixed-endpoint
certificate of Lemma~\ref{sbpos:endpoints}.  The empirical cross moment $b$
is retained with its sign by Lemma~\ref{sbemp:b-netting}; it is not to be
charged a second time inside $E_{12}$.  Likewise the leading deterministic
$c=e^2$ part of $E_{22}$ should be netted into directional face margins
rather than automatically spent as an absolute error.

The main quantitative risk is the simultaneous clock condition in
Remark~\ref{sbemp:combined-clock}.  The mean relaxation and progress clocks
must overlap rather than stack.  After entry, the least-developed part of
the unconditional theory remains whole-complement regulator forcing with
explicit cone moduli, nonnegative initial regulator gap, progress
monotonicity or a valid replacement, and sufficient damping through the
first-hit creation corridor.  General finite-window effective-matrix
crossing and leakage for Theorem~3 remain separate.  Finally, no common
primitive parameter region satisfying every early, entry, generation, and
creation guard has yet been proved nonempty.
\end{remark}

\section{Fixed-complement transport and the pre-crossing output envelope}
\label{sbflux:scope}

Fix the handoff family $B=H\sqcup\mathcal L$ and retained weak cohort $W$.
Throughout this section $E=[h]\setminus(B\cup W)$ is fixed. Sign sets used
inside sums do not change this membership. We first replace the coarse
entrant mass by its actual cross-coordinate weight, then propagate that
weight and identify the remaining quantitative production bound.

\subsection{The weighted complement quantity}

Use physical time in this subsection and put
\[
 V_E=\sum_{i\in E}(v_{i1})_+^2,\quad
 Q_E=\sum_{i\in E}(v_{i1})_+(v_{i2})_+,\quad
 R_E=\sum_{i\in E}\|v_i\|^2,\quad
 D_{E,1}^2=\sum_{i\in E}\delta_{i1}^2.
\]
With the empirical cluster-2 coordinates $x_1=e\mu_2z$, $x_2=\mu_2y$
from the empirical reference system, define
\[
 G_+(s)=\mathbb E_{2,n}[z_+y\alpha_s],\qquad
 h(s)=\mathbb E_{2,n}[y^2\alpha_s],
\]
and retain the fixed-data constant
\[
 K_2=\mu_2^{-2}\sup_{\|u\|=1,s}
       \mathbb E_{2,n}[(u^\top x)_+x_2\alpha_s].
\]
Assume $x_2>0$ on cluster 2.

\begin{lemma}[Cross-weighted signed complement bound]
\label{sbflux:output}
For every compatible test gate,
\begin{equation}\label{sbflux:output-bound}
 \left[\frac{\mathbb E_{2,n}[f_{E,1}x_2\alpha_s]}{e\mu_2^2}\right]_+
 \le G_+(s)V_E+\frac{h(s)}e Q_E
       +\frac{K_2}e\sqrt{R_E}\,D_{E,1}.
\end{equation}
This covers partial gates, later entrants, and the sign-mismatched class
without a separate wrong-output norm hypothesis.
\end{lemma}
\begin{proof}
Discard favorable terms with $z_{i1}\le0$ and use
$(z_{i1})_+\le(v_{i1})_++|\delta_{i1}|$.
For $v_{i1}>0$,
\[
 (v_i^\top x)_+\le v_{i1}(x_1)_++(v_{i2})_+x_2.
\]
Multiplying by $(v_{i1})_+x_2\alpha_s$ and averaging gives the first two
terms. The remaining neuronwise term is bounded by
$K_2\mu_2^2r_i|\delta_{i1}|$; sum it by Cauchy--Schwarz and divide by
$e\mu_2^2$. No sign restriction on $v_{i1}$ is needed for this remainder.
\end{proof}

The angular representation is exact:
\[
 Q_E=\int w_+(\phi)\,d\nu_E(\phi),\qquad
 w_+(\phi)=(\cos\phi)_+(\sin\phi)_+.
\]
On the first quadrant this is $\tfrac12\sin2\phi$ and on the other
quadrants it is zero. Unlike the indicator of a purity sector, this
Lipschitz weight has no jump at either coordinate axis. Consequently
entry across an axis produces no unaccounted boundary atom. A nonzero
purity boundary inside the first quadrant is included continuously in
this weight, rather than treated as a zero-cost entry point.

\begin{lemma}[Exact weighted transport without a boundary-count assumption]
\label{sbflux:transport}
Write $A_i=A_i^\rho$ and $b_i=A_i^\top\delta_i+P_i$ so that
$\dot v_i=A_i^\top v_i-b_i$.
Let $I_i=\mathbf1_{\{v_{i1}>0,v_{i2}>0\}}$.
Then $Q_E$ is absolutely continuous and, almost everywhere,
\begin{equation}\label{sbflux:transport-exact}
 \dot Q_E=\sum_{i\in E}I_i\left[
 (A_{11,i}+A_{22,i})v_{i1}v_{i2}
 +A_{12,i}v_{i1}^2+A_{21,i}v_{i2}^2
 -v_{i2}b_{i1}-v_{i1}b_{i2}\right].
\end{equation}
The coordinate formula is valid even at times when some $r_i=0$; no
angular chart or crossing-count bound is required.
\end{lemma}
\begin{proof}
Apply the almost-everywhere chain rule to the locally Lipschitz function
$(x,y)\mapsto x_+y_+$. On a coordinate zero set its corresponding
absolutely continuous coordinate has derivative zero almost everywhere.
Thus there is no additional singular term at crossings. In the open first
quadrant differentiate $v_{i1}v_{i2}$ and insert the Cartesian equations.
This is also the Lipschitz-kernel version of the installed weighted
angular-measure evolution, with its signs retained.
\end{proof}

\begin{corollary}[A signed production comparison]
\label{sbflux:production}
Let $\theta(t)$ be any integrable comparison rate and define
\[
 \mathcal J_E^\theta(t)=\left[\sum_{i\in E}I_i
 \{(\operatorname{tr}A_i-\theta)v_{i1}v_{i2}
                 +A_{12,i}v_{i1}^2+A_{21,i}v_{i2}^2\}\right]_+,
\]
\[
 \mathcal B_E^-(t)=\left[-\sum_{i\in E}I_i
                  (v_{i2}b_{i1}+v_{i1}b_{i2})\right]_+.
\]
For $t\ge a_0$,
\begin{equation}\label{sbflux:Q-envelope}
 Q_E(t)\le e^{\int_{a_0}^t\theta}Q_E(a_0)
 +\int_{a_0}^t e^{\int_v^t\theta}
               (\mathcal J_E^\theta+\mathcal B_E^-)(v)\,dv.
\end{equation}
In particular, if $\|A_i\|\le K$, then
\[
 \mathcal B_E^-\le K\sqrt{R_E\Delta_E}+\sqrt{R_E}\,\mathcal P_E,
 \quad \Delta_E=\sum_E\|\delta_i\|^2,\quad
 \mathcal P_E^2=\sum_E\|P_i\|^2.
\]
If instead one has entrywise bounds
$\operatorname{tr}A_i\le\theta$, $[A_{12,i}]_+\le b_{12}$ and
$[A_{21,i}]_+\le b_{21}$ on the active quadrant, then
\[
 \mathcal J_E^\theta\le b_{12}\sum_E I_iv_{i1}^2
                              +b_{21}\sum_E I_iv_{i2}^2.
\]
The signed sum defining $\mathcal J_E^\theta$ is the sharper quantity;
the entrywise bound is only a fallback.
\end{corollary}
\begin{proof}
Subtract $\theta Q_E$ in the exact identity, bound the two remaining
signed sums by their positive parts, and integrate. The mismatch estimate
uses $\|(v_{i2},v_{i1})\|=r_i$ and Cauchy--Schwarz.
\end{proof}

\begin{remark}[A mass-relative version]
If a positive reference mass $P_B$ obeys
$\dot P_B/P_B\ge g_B(t)$, then the same calculation gives
\[
 \frac{Q_E(t)}{P_B(t)}\le
 e^{\int_{a_0}^t(\theta-g_B)}\frac{Q_E(a_0)}{P_B(a_0)}
 +\int_{a_0}^t e^{\int_v^t(\theta-g_B)}
       \frac{\mathcal J_E^\theta+\mathcal B_E^-}{P_B}(v)\,dv.
\]
This retains any anisotropic improvement of the trace rate over the
strong-mass growth rate. It does not impose an ordering on transverse
mass fractions.
\end{remark}

\subsection{What maximal fixed enlargement supplies}

Let the cluster-1 coordinate envelope satisfy
$x_1\ge m_1^{\rm data}>0$, $|x_2|\le b_2$, and define
\[\kappa_{\rm band}=\frac{b_2}{\sqrt{(m_1^{\rm data})^2+b_2^2}}.\]
At handoff partition the fixed complement into neurons not fully active
on cluster 1 and the remaining set $F$ of fully active complement neurons.
Failure of full activation with $v_{i1}>0$ gives
$v_{i1}\le\kappa_{\rm band}r_i$. Hence
\begin{equation}\label{sbflux:handoff-weight}
 Q_E(a_0)\le\kappa_{\rm band}R_{E\setminus F}(a_0)
                      +Q_F(a_0).
\end{equation}
This retains the boundary discount. Maximal enlargement alone does not
set $Q_F(a_0)$ to zero: a currently pure neuron can have failed a required
future radius, mismatch, longitudinal, or purity certificate. Those
failures require a quantitative coverage bound.

For example, the contribution of currently pure neurons with $z_{i1}\le0$
is at most $\sqrt{R_F\Delta_{F,1}}$, since then
$(v_{i1})_+\le|\delta_{i1}|$.
A subset with $r_i\le r_*$ contributes at most its cardinality times
$r_*^2/2$; a subset with $\|\delta_i\|\ge\eta r_i$ contributes at most
$\Delta_F/(2\eta^2)$. These estimates can be summed over specified
admission failures; they do not cover a failure whose only description
is that its future guard has not been proved.

\subsection{The genuinely concurrent output lower bound}

Return to normalized time $\tau=\mu_2^2t/2$. Suppose the input envelope is
$s_i\ge-S(\tau)$, $0<M\le\overline M(\tau)<1$,
$\bar t\le\overline t(\tau)$, $q_-\le q_i\le q_+$, and
$r_{w,i}\ge-\epsilon_w^-(\tau)$ in the exact empirical output equation.
Here $a,b,K$ retain their empirical meanings. Define
\[
 A_0(\tau)=b+K(-S(\tau))
            -a\overline M(\tau)[\overline t(\tau)]_+,
 \quad c_w(\tau)=a(1-\overline M(\tau))/q_+,
\]
\[
 F_w(\tau)=q_-^{-1}[-A_0(\tau)]_++\epsilon_w^-(\tau).
\]

\begin{lemma}[Output lower envelope before input crossing]
\label{sbflux:pre-cross-output}
For $Y=[-\min_{i\in B}t_i]_+$ one has
\[
 D^+Y\le-c_wY+F_w
\]
in the integrated almost-everywhere sense. Therefore
\begin{equation}\label{sbflux:Y-envelope}
 Y(\tau)\le \overline Y(\tau):=
 e^{-\int_{\tau_0}^\tau c_w}Y(\tau_0)
 +\int_{\tau_0}^\tau e^{-\int_v^\tau c_w}F_w(v)\,dv.
\end{equation}
This comparison starts at the same time as mean relaxation and does not
wait for input sign crossing.
\end{lemma}
\begin{proof}
At an active negative output minimum $t_i=-Y$,
\[
 G_i=b+K(s_i)-aM\bar t+a(1-M)Y
       \ge A_0+a(1-\overline M)Y.
\]
Use $q_i^{-1}\ge q_+^{-1}$ for the nonnegative last term.
For the adverse term, $A_0/q_i\ge-q_-^{-1}[-A_0]_+$.
The scalar comparison follows after charging $\epsilon_w^-$.
At $Y=0$ the bound remains valid because $F_w\ge0$.
\end{proof}

The favorable part of $A_0$ can also be retained. Define
\[
 f_w=\epsilon_w^- -q_+^{-1}[A_0]_+
                     +q_-^{-1}[-A_0]_+,\qquad
 C(\tau)=\int_{\tau_0}^\tau c_w,
\]
\[
 Z(\tau)=Y(\tau_0)+\int_{\tau_0}^\tau e^{C(v)}f_w(v)\,dv,
\]
\begin{equation}\label{sbflux:reflected-output}
 \overline Y_{\rm sign}(\tau)
 =e^{-C(\tau)}\left[Z(\tau)-
      \min\left\{0,\inf_{\tau_0\le v\le\tau}Z(v)\right\}\right].
\end{equation}
Then $Y\le\overline Y_{\rm sign}\le\overline Y$.
Indeed, for $Y>0$ the same proof gives $Y'\le-c_wY+f_w$;
at zero the upper bound is $[f_w]_+$. After multiplication by $e^C$,
reflection at zero gives the displayed scalar supersolution. Scalar
comparison proves the assertions. This version retains beneficial target
forcing and can certify an output zero before input crossing. In the
following bound one may replace $\overline Y$ by
$\overline Y_{\rm sign}$ everywhere.

Before input crossing, T1.9 gives $t_i\ge-T(\tau)$, not in general
$t_i\ge-D$ with $D=\tan(D_{\rm lock})/e$.
The admissible pre-crossing lower bound is therefore
\[
 t_i\ge-\min\{T(\tau),\overline Y(\tau)\}.
\]
The input-crossing proof allows this time-dependent quantity in place of
its constant $D$: it only uses the lower bound pointwise and never
differentiates it. After input crossing, the gap bound and monotonicity
of tangent additionally give $t_i\ge-D$. This distinction must be
preserved when using the canonical zero-face margin. The output
lower-envelope calculation, the input-crossing integral, and the moving
mean Duhamel comparison all run on the same time interval.

\subsection{Ambient crossing costs from the installed radius envelopes}

For a neuron in the fixed strong family let $\xi_i=e s_i$ and define
\[
 H_i^\perp=\mathbb E_n[(x_2-\xi_ix_1)^2\alpha_i],\qquad
 p_{u,+}=\max_{i\in B}\left[
 \frac{2(P_{i2}-\xi_iP_{i1})}{e\mu_2^2v_{i1}}\right]_+.
\]
The installed perpendicular energy bound yields directly
\begin{align*}
 p_{u,+}
 &\le\frac{2\varepsilon\sqrt{h\lambda R}}{e\mu_2^2}
       \max_{i\in B}\frac{\|w_{i,\perp}\|\sqrt{H_i^\perp}}{v_{i1}},\\
 p_1^\star:=\max_{i\in B}\frac{|P_{i1}|}{v_{i1}}
 &\le\varepsilon\lambda\sqrt{hR}
       \max_{i\in B}\frac{\|w_{i,\perp}\|}{v_{i1}}.
\end{align*}
Indeed, apply Cauchy--Schwarz to
$P_{i2}-\xi_iP_{i1}
=\mathbb E_n[\langle f_\perp,w_{i,\perp}\rangle
(x_2-\xi_ix_1)\alpha_i]$ and use
$\|f_\perp\|_n\le\varepsilon\sqrt{h\lambda R}$.
The longitudinal estimate uses $\mathbb E_n[x_1^2\alpha_i]\le\lambda$.

In particular, if $v_{i1}\ge v_*(t)>0$ and $|\xi_i|\le u(t)$, then
\[
 H_i^\perp\le H_*(t):=
 \max_{\xi\in\{-u(t),u(t)\}}\mathbb E_n[(x_2-\xi x_1)^2],
\]
\[
 p_{u,+}\le\frac{2h\varepsilon^2\sqrt{\lambda R H_*}}
                       {e\mu_2^2v_*},\qquad
 p_1^\star\le\frac{\lambda h\varepsilon^2\sqrt R}{v_*}.
\]
One may use $R\le h\varepsilon^2+\lambda t$ and the T1.9 radius/cone
lower bound for $v_*$ where that theorem applies; qualifying laggers must
carry their certified radius lower bounds as well. These are explicit
parameter/radius checks, not extra wrong-output assumptions. The costs
vanish in the planar case. Choosing initialization small to pay them must
still be reconciled with the common retention and entry clocks.

\subsection{Following the transport bound into the timing audit}

Let $\overline Q_E$ be the right side of the transport comparison and let
$\overline V_E,\overline R_E,\overline D_{E,1}$ be certified upper bounds.
For the test slopes used before crossing, set
\[
 \mathfrak X_E(t)=\sup_s\left[
 G_+(s)\overline V_E(t)+\frac{h(s)}e\overline Q_E(t)
       +\frac{K_2}e\sqrt{\overline R_E(t)}\,\overline D_{E,1}(t)
 \right].
\]
If $\epsilon_{u,\rm other}^-$ denotes all other adverse input terms after
signed cluster-1 netting, a valid input-crossing audit must in particular
reserve
\begin{equation}\label{sbflux:crossing-budget}
 \Gamma_u(\tau)=q_-a_{u,\times}
 \bigl(\tau;\min(T,\overline Y)\bigr)
 -\epsilon_{u,\rm other}^-(\tau)
 -q_+\mathfrak X_E(2\tau/\mu_2^2)>0,
\end{equation}
and integrate this margin in the same coupled clock as the mean envelope.
There is no extra stage time associated with the transport estimate.

The installed aggregate mismatch inequality supplies
\[
 \sqrt{\Delta_E(t)}\le
 e^{\int_{a_0}^t\mu_\rho}\sqrt{\Delta_E(a_0)}
 +\int_{a_0}^t e^{\int_v^t\mu_\rho}
     \{Z_E(v)+\mathcal P_E(v)\}\,dv,
\]
where $Z_E=(\sum_E r_i^2|X_i^\rho|^2)^{1/2}$ is the weighted
antisymmetric residual forcing. Thus mismatch can be charged to an
installed equation, but that equation still requires a quantitative
bound on $Z_E$ for the excluded fixed family.

\begin{remark}[Transport frontier before the coupled signed-source continuation]
\label{sbflux:frontier}
The source decomposition and coupled propagation requested here are now supplied in Section~\ref{sbcouple:scope}. The remaining initialized audit is \eqref{sbcouple:unresolved-budget}; the following records why the preceding transport identity alone did not close it.

The needed new estimate is a primitive-parameter bound on the right side
of \eqref{sbflux:Q-envelope}, together with the excluded-family mismatch
integral and handoff coverage, small enough for
\eqref{sbflux:crossing-budget} and the concurrent clock. Neither an angular
boundary count nor a new wrong-output norm assumption is needed.

The available operator bound $\|A_i\|\le K$ does give the fallback
$\theta=2K$, $\mathcal J_E^\theta\le K R_E$, but substituting it introduces
$e^{-1}\int e^{\int2K}KR_E$ in the input charge. It does not certify a
noise-scale production estimate. This is a limitation of that bound, not
a proof that the actual signed production is large or that the theorem
region is empty. A sharper proof must retain the signed source sum and
control the excluded-family antisymmetric forcing, or certify enough
handoff coverage to make the remaining complement negligible.

The missing sign information is not supplied merely by balanced SIM
geometry. For a static balanced noiseless SIM state with
$\lambda_1=9/2$, $\lambda_2=2$, take
\[
 u_1=w_1=(1,-2/9),\quad u_2=w_2=(-1/2,1),\quad
 u_3=w_3=r(\eta,1),\qquad r>0,\quad\eta>0.
\]
On the two samples $3e_1,2e_2$, neurons 1 and 2 are respectively
concept-1-only and concept-2-only, while neuron 3 is active on both.
Its residual matrix is exactly
\[
 A_3=\begin{pmatrix}
 -(9/2)r^2\eta^2&1-2r^2\eta\\
 1-(9/2)r^2\eta&-2r^2
 \end{pmatrix}.
\]
For the fixed singleton $E=\{3\}$, at this state $\Delta_E=0$,
$Q_E=r^2\eta$, and direct substitution gives
\[
 \frac{\dot Q_E}{r^2}
 =1+\eta^2-\frac{13}{2}r^2\eta(1+\eta^2)
 \longrightarrow1\qquad(\eta\downarrow0).
\]
Thus a boundary-small weight and zero mismatch do not imply small signed
production, even with the SIM gate geometry. The positive injection comes
from the other neuron's negative coordinate-2 output on concept 1.
This is a static counterexample to that inference, not a trajectory from
the paper's small isotropic initialization, and not a reopening of the
restricted reversal theorem. Its purpose is to locate the missing global
signed-output or coverage estimate precisely.

The distinction is essential: a boundary-small weight is not conserved
by inward angular motion. For a perfectly matched diagonal linear flow
$\dot v=\operatorname{diag}(\lambda_1,\lambda_2)v$ with positive entries,
\[
 Q_i(t)=Q_i(0)e^{(\lambda_1+\lambda_2)t},\qquad
 \tan\phi_i(t)=\tan\phi_i(0)e^{-(\lambda_1-\lambda_2)t}.
\]
The radius growth contributes along with angular transport. This example
only tests the transport inference; it is not a counterexample to the
initialized ReLU theorem. Likewise maximal certified enlargement does not
by itself prove that every excluded pure neuron has small weight.

Until the signed production, mismatch and coverage bounds are supplied,
the strict initialized ordering of handoff versus progress and the T1.9
horizon is not established. Consequently one cannot yet claim to have
reached the hypotheses of the subsequent signed-gap, generation, creation,
or general probe-reversal theorems.
\end{remark}

\section{Signed cohort sources and coupled complement propagation}
\label{sbcouple:scope}

We retain the fixed partition $[h]=B\sqcup W\sqcup E$ and use physical
time. All identities hold almost everywhere along the selected Clarke
trajectory. In particular, sign sets below are weights inside fixed-cohort
sums, not changes to cohort membership. This section closes the algebraic
coupling left in Remark~\ref{sbflux:frontier}; it does not assert initialized
entry or the common stopping-time ordering.

\subsection{Exact signed decomposition and the self-interaction calculation}

Write $\sigma_i(x)=(v_i^\top x)_+$, $I_i=\mathbf1_{\{v_{i1}>0,v_{i2}>0\}}$,
$\widehat v_i=(v_{i2},v_{i1})$, and, for a planar field $F$, set
\[
 \mathscr P_E(F)=-\mathbb E_n\sum_{i\in E}I_i\alpha_i
           (\widehat v_i\cdot x)(v_i\cdot F(x)).
\]
With $f_H=\sum_{j\in H}z_j\sigma_j$, and
$f_E^v=\sum_Ev_j\sigma_j$, $f_E^\delta=\sum_E\delta_j\sigma_j$,
the exact production identity is
\begin{align}
 \dot Q_E={}&-\mathscr P_E(x)+\mathscr P_E(f_B)
       +\mathscr P_E(f_W)+\mathscr P_E(f_E^v)
       -\mathscr P_E(f_E^\delta)\notag\\
 &-\sum_E I_i\widehat v_i\cdot(A_i^\top\delta_i+P_i).
 \label{sbcouple:exact-production}
\end{align}
The same split, before an absolute value is taken, gives
\begin{equation}\label{sbcouple:exact-skew}
 X_i^\rho=X_i(f_B)+X_i(f_W)+X_i(f_E^v)-X_i(f_E^\delta),
 \qquad X_i(F)=\mathbb E_n[(F\times x)\alpha_i].
\end{equation}
There is no target skew: $x\times x=0$ exactly, including noisy samples.
Ambient forcing is $P_i$, not an additional target or planar skew term.
If $E$ is empty all complement charges vanish; below assume $E$ is nonempty.

For vectors, $x_+,x_-$ denote coordinatewise positive and negative parts.
Define the empirical negative-feature moment
\[
 \epsilon_-:=\mathbb E_n[\|x_-\|\,\|x\|].
\]
\begin{lemma}[Matched first-quadrant self-interaction]
\label{sbcouple:self}
At a time when all donor inputs in $E$ are coordinatewise nonnegative,
\begin{equation}\label{sbcouple:self-bound}
 \mathscr P_E(f_E^v)\le\epsilon_-R_E R_I\le\epsilon_-R_E^2,
 \qquad R_I=\sum_E I_i r_i^2.
\end{equation}
In particular the matched self term is nonpositive on nonnegative data.
On the SIM coordinate envelope, with $x_k>0$ on cluster $k$, write
$x=x_ke_k+n e_{3-k}$. If $|n|\le e b_k$, then
\[
 \epsilon_-\le e C_- ,\qquad
 C_-:=\mathbb E_n[b_{k(x)}\|x\|].
\]
This proves a finite-sample $eC_-R_E^2$ bound, with no density or
activation-count assumption. The donor-quadrant condition is essential.
\end{lemma}
\begin{proof}
Here $f_E^v\ge0$ coordinatewise. For every first-quadrant receiver,
$v_i\cdot f_E^v\ge0$, so only $\widehat v_i\cdot x_-$ can contribute
positively to $\mathscr P_E(f_E^v)$. Bound this by $r_i\|x_-\|$,
$v_i\cdot f_E^v$ by $r_i\|f_E^v\|$, and
$\|f_E^v(x)\|\le R_E\|x\|$. Sum and average. Positivity of the
dominant SIM coordinate gives $\|x_-\|\le |n|$, proving the last claim.
\end{proof}

There is a corresponding signed strong-source estimate. Suppose
$z_{j1}\ge0$ for $j\in B$, and set
\[
 U_H=\sum_{j\in H}r_j\|z_j\|,\qquad
 N_B^-:=\sum_{j\in B}r_j[-z_{j2}]_+,
 C_+:=\mathbb E_n[\|x_+\|\,\|x\|].
\]
Splitting both $x$ and $f_B$ into coordinatewise signs gives
\begin{equation}\label{sbcouple:strong-one-sided}
 \mathscr P_E(f_B)\le R_I(C_+N_B^-+\epsilon_-U_B).
\end{equation}
Indeed $[f_{B,1}]_-=0$, $[f_{B,2}]_-\le N_B^-\|x\|$, and
$\|[f_B]_+\|\le U_B\|x\|$; retain the two favorable products before
bounding the remaining two. In the strong ratio coordinates, $|s_j|\le S$
and $t_j\ge-Y$ imply
\[
 N_B^-\le e M Y\sqrt{1+e^2S^2}.
\]
Thus the installed concurrent output envelope controls precisely this
adverse strong production. It does not control the entire strong skew.

\subsection{An axis comparison that retains every actual receiver gate}

For a sample in cluster $k$, set $\bar x=x_ke_k$; assume $x_k>0$.
Define fixed empirical constants
\[
 \chi=\mathbb E_n[\|x-\bar x\|(\|x\|+\|\bar x\|)],\quad
 C_x=\mathbb E_n\|x\|^2,\quad
 \lambda\ge\|\mathbb E_nxx^\top\|_{\rm op},\quad
 c_\delta=\sqrt{\lambda C_x}.
\]
The coordinate envelope gives
$\chi\le e\mathbb E_n[b_{k(x)}(\|x\|+x_{k(x)})]$.
For each actual receiver selector, put
\[
 \ell_{ki}=\mathbb E_n[\mathbf1_{\mathcal C_k}x_k^2\alpha_i(x)],
 \qquad 0\le\ell_{ki}\le\ell_k:=\mathbb E_n[\mathbf1_{\mathcal C_k}x_k^2].
\]
No selector in $\ell_{ki}$ is replaced by its value at $\bar x$.

For a cohort $H$ with its actual output vectors define
\[
 p_{kH}=\sum_H z_{jk}(v_{jk})_+,\qquad
 c_{1H}=\sum_H z_{j2}(v_{j1})_+,\qquad
 c_{2H}=\sum_H z_{j1}(v_{j2})_+.
\]
Thus $f_H(\bar x)=x_1(p_{1H},c_{1H})$ on cluster 1 and
$x_2(c_{2H},p_{2H})$ on cluster 2.

\begin{lemma}[Finite-data cohort comparison]
\label{sbcouple:axis}
For every receiver, replacing $f_H(x)x^\top$ by
$f_H(\bar x)\bar x^\top$ inside its selected average costs at most
$\chi U_H$ in operator norm. Replacing $f_H(x)\times x$ by
$f_H(\bar x)\times\bar x$ costs at most $\chi U_H$ in absolute value.
For $f_E^v$ the same statements hold with $U_H=R_E$.
Moreover
\[
 \|\mathbb E_n[f_E^\delta x^\top\alpha_i]\|_{\rm op}
       \le\lambda\sqrt{R_E}D_E,\qquad
 |X_i(f_E^\delta)|\le c_\delta\sqrt{R_E}D_E.
\]
\end{lemma}
\begin{proof}
ReLU is 1-Lipschitz, so
$\|f_H(x)-f_H(\bar x)\|\le U_H\|x-\bar x\|$ and
$\|f_H(\bar x)\|\le U_H\|\bar x\|$.
Expand the difference of the two products (or determinants) into these
two errors and use $0\le\alpha_i\le1$. For the mismatch field,
$\|f_E^\delta\|_n\le\sqrt{\lambda R_E}D_E$ by Cauchy--Schwarz over
neurons. Pair with a unit projection of $x$, or with $\|x\|$, respectively.
\end{proof}

The missing signed input moments are
\begin{equation}\label{sbcouple:N}
 N_4=\sum_E(v_{i1})_+(-v_{i2})_+,\quad
 N_2=\sum_E(-v_{i1})_+(v_{i2})_+,\quad N=N_4+N_2.
\end{equation}
For matched excluded donors, their two cross outputs on the axes are
\[
 c_{1E}^v=Q_E-N_4,\qquad c_{2E}^v=Q_E-N_2.
\]
In particular, small $Q_E$ and $D_E$ do not control these signed outputs.
Set $c_k=c_{kB}+c_{kW}$ and
\[
 p_1=p_{1B}+p_{1W}+\sum_E(v_{i1})_+^2,\quad
 p_2=p_{2B}+p_{2W}+\sum_E(v_{i2})_+^2,\quad
 U=U_B+U_W+R_E.
\]
The exact residual matrix has the comparison
\begin{equation}\label{sbcouple:A0}
 A_i=A_i^0+\mathcal E_i,\quad
 A_i^0=\begin{pmatrix}
 \ell_{1i}(1-p_1)&-\ell_{2i}(c_2+Q_E-N_2)\\
 -\ell_{1i}(c_1+Q_E-N_4)&\ell_{2i}(1-p_2)
 \end{pmatrix},\quad
 \|\mathcal E_i\|\le\chi(1+U)+\lambda\sqrt{R_E}D_E.
\end{equation}
Here the target replacement costs $\chi$; it is not included in skew.
The corresponding skew formula is
\begin{equation}\label{sbcouple:X0}
 X_i^\rho=\ell_{2i}c_2-\ell_{1i}c_1
       +(\ell_{2i}-\ell_{1i})Q_E+\ell_{1i}N_4-\ell_{2i}N_2+\xi_i,
 \quad |\xi_i|\le\chi U+c_\delta\sqrt{R_E}D_E.
\end{equation}
These formulas retain all mixed-quadrant donors and all gate changes.

\subsection{A coupled comparison with the missing sign source propagated}

Let $J_i=\mathbf1_{\{v_{i1}v_{i2}<0\}}$ and $R_J=\sum_EJ_i r_i^2$.
For $G=I,J$ define
\[
 b_{G2}=\sum_EG_i\ell_{2i}v_{i1}^2,\quad
 b_{G4}=\sum_EG_i\ell_{1i}v_{i2}^2,\quad a_G=b_{G2}+b_{G4},
\]
\[
 \theta_G=\max_{i:G_i=1}
       \{\ell_{1i}(1-p_1)+\ell_{2i}(1-p_2)\},\qquad
 \beta_G=K\sqrt{R_G}+\lambda R_G\sqrt{R_E},
\]
with $\theta_G=0$ for an empty set, and $K$ from \eqref{sbnew:A-norm}.
Put $\mathcal P_E=(\sum_E\|P_i\|^2)^{1/2}$ and
\begin{align*}
 S_Q&=[-b_{I2}c_2-b_{I4}c_1]_+
             +\chi(1+U)R_I+\sqrt{R_I}\mathcal P_E,\\
 S_N&=[b_{J2}c_2+b_{J4}c_1]_+
             +\chi(1+U)R_J+\sqrt{R_J}\mathcal P_E,\\
 \zeta_Q&=\left(\sum_Er_i^2(\ell_{2i}-\ell_{1i})^2\right)^{1/2},\\
 \zeta_N&=\max_{k=1,2}\left(\sum_Er_i^2\ell_{ki}^2\right)^{1/2},\\
 S_D&=\left(\sum_Er_i^2(\ell_{2i}c_2-\ell_{1i}c_1)^2\right)^{1/2}
                    +\chi U\sqrt{R_E}+\mathcal P_E.
\end{align*}

\begin{proposition}[Closed signed-moment/mismatch differential comparison]
\label{sbcouple:comparison}
In the integrated almost-everywhere sense,
\begin{equation}\label{sbcouple:matrix}
 \begin{pmatrix}Q_E\\N\\D_E\end{pmatrix}'\preceq
 \begin{pmatrix}
 \theta_I-a_I&\max(b_{I2},b_{I4})&\beta_I\\
 a_J&\theta_J&\beta_J\\
 \zeta_Q&\zeta_N&\mu_\rho+c_\delta R_E
 \end{pmatrix}
 \begin{pmatrix}Q_E\\N\\D_E\end{pmatrix}
 +\begin{pmatrix}S_Q\\S_N\\S_D\end{pmatrix}.
\end{equation}
In particular no unevaluated $Z_E$ remains. The $N$ equation propagates
the actual omitted sign source; it is not an extra internal radius moment.
If an independent $N$ envelope is available, rows 1 and 3 give the requested
two-variable $(Q_E,D_E)$ comparison with that envelope in their sources.
Setting $N=0$ requires a proved invariant, not merely $N(a_0)=0$.
\end{proposition}
\begin{proof}
For a first-quadrant receiver, insert $A_i^0$ into the exact Cartesian
product derivative. The matched self cross terms sum to
$-a_IQ_E+b_{I2}N_2+b_{I4}N_4$; its diagonal self terms remain in
$\theta_I$. The external cross contribution is exactly
$-b_{I2}c_2-b_{I4}c_1$. The matrix replacement costs
$R_I\{\chi(1+U)+\lambda\sqrt{R_E}D_E\}$, and the receiver mismatch
costs $K\sqrt{R_I}D_E+\sqrt{R_I}\mathcal P_E$.
This proves row 1, retaining self damping.

In either mixed-sign quadrant the product weight is $-v_{i1}v_{i2}$.
Its derivative has the same trace term, but the signs of both cross
terms are reversed. Summing gives $a_JQ_E-b_{J2}N_2-b_{J4}N_4$ plus
$b_{J2}c_2+b_{J4}c_1$. Drop only the two favorable self terms and use
the same error bounds. The coordinatewise positive-part chain rule
handles axes almost everywhere without boundary atoms.

Finally, \eqref{sbcouple:X0} and the triangle inequality in the weighted
Euclidean norm give
\[
 Z_E\le\zeta_QQ_E+\zeta_NN+
 \left(\sum_Er_i^2(\ell_{2i}c_2-\ell_{1i}c_1)^2\right)^{1/2}
 +\chi U\sqrt{R_E}+c_\delta R_ED_E.
\]
Insert this evaluated source into the exact mismatch energy inequality.
All off-diagonal coefficients in \eqref{sbcouple:matrix} are nonnegative.
\end{proof}

\paragraph{Retaining mismatch damping instead of the loss-norm rate.}
Let $d_{ki}=\ell_{ki}(1-p_k)$ and
$h_i=\ell_{2i}(c_2+Q_E-N_2)+\ell_{1i}(c_1+Q_E-N_4)$.
The least eigenvalue of $\operatorname{sym}A_i^0$ is exactly
\[
 m_i=\frac{d_{1i}+d_{2i}-\sqrt{(d_{1i}-d_{2i})^2+h_i^2}}2.
\]
With $m_* =\min_{i\in E}m_i$, the last diagonal coefficient can be
sharpened to
\begin{equation}\label{sbcouple:damping}
 c_\delta R_E+
 \min\{\mu_\rho,-m_*+\chi(1+U)+\lambda\sqrt{R_E}D_E\}.
\end{equation}
This follows directly from the mismatch energy and the operator error in
\eqref{sbcouple:A0}. Under certified caps it provides a cooperative
linear comparison; alternatively keep the displayed quadratic
$\lambda\sqrt{R_E}D_E^2$ in a bootstrap. Positive residual diagonals
can produce damping, but cannot be presumed when constructing the caps.
This avoids paying the global loss-norm growth rate if the signed
residual spectrum is better.

\paragraph{What the strong and weak sources actually cost.}
For the strong family, exactly,
\[
 c_{1B}=eM\bar t,\quad c_{2B}=eM\overline{s_+},\quad
 p_{1B}=M,\quad p_{2B}=e^2M\overline{t s_+}.
\]
Suppose $s_i\in[-S,S]$, $t_i\in[-Y,T]$. Then
\[
 [-c_{1B}]_+\le eMY,\quad c_{2B}\le eMS,\quad
 |c_{1B}|\le eM\max(Y,T),\quad p_{2B}\ge-e^2MYS,
\]
\[
 U_B\le M\sqrt{1+e^2S^2}\sqrt{1+e^2\max(Y,T)^2}.
\]
Use $Y=\min(T,\overline Y_{\rm sign})$ when the installed symmetric
output envelope is $[-T,T]$. If the retained weak family has
$v_1\le0,v_2\ge0,z_1\le0,z_2\ge0$, then
$c_{1W}=p_{1W}=0$, $c_{2W}\le0$, $p_{2W}\ge0$.
Consequently the external positive parts above obey
\begin{align*}
 [-b_{I2}c_2-b_{I4}c_1]_+&\le b_{I2}|c_{2W}|+b_{I4}eMY,\\
 [b_{J2}c_2+b_{J4}c_1]_+&\le b_{J2}eMS+b_{J4}eMT.
\end{align*}
The signed expressions are sharper when cancellation is present. Also
$|c_{2W}|\le U_W$, or, for the matched-plus-mismatch weak decomposition,
$|c_{2W}|\le R_W/2+\sqrt{R_W}D_W$.
For the skew source one has the explicit fallback
\[
 \left(\sum_Er_i^2(\ell_{2i}c_2-\ell_{1i}c_1)^2\right)^{1/2}
 \le\sqrt{R_E}\{\ell_2(eMS+|c_{2W}|)+\ell_1eM\max(Y,T)\}.
\]
Positive strong tilt is therefore a real noise-scale mismatch source
even when its negative-output production source has vanished.

\subsection{Solving the coupling and consuming it in the entry clock}

Let $C(t)$ denote the cooperative matrix in \eqref{sbcouple:matrix},
using \eqref{sbcouple:damping} if better. Given certified coefficient
and source envelopes $C(t)\preceq\overline C(t)$ and
$S(t)\preceq\overline S(t)$, let $\Phi$ be the fundamental matrix
of $\overline C$. Then the joint supersolution is
\begin{equation}\label{sbcouple:solution}
 \overline z(t)=\Phi(t,a_0)\overline z(a_0)
             +\int_{a_0}^t\Phi(t,s)\overline S(s)\,ds,
 \qquad z=(Q_E,N,D_E)^\top\preceq\overline z.
\end{equation}
For constant envelopes $\Phi(t,s)=\exp((t-s)\overline C)$.
Equivalently, positive weights $w$ with
$w^\top\overline C\preceq\gamma w^\top$ give the scalar Duhamel
bound for $w_QQ_E+w_NN+w_DD_E$. The weights must be chosen for this
matrix; a generic scalar Gronwall is not substituted for its coupling.
These assertions follow from variation of constants and positivity of
the fundamental matrix of a cooperative system. For a bootstrap, certify
the coefficient caps on a closed rectangle and check that
\eqref{sbcouple:solution} remains strictly inside it up to the common
deadline.

\begin{lemma}[Common-mass normalization of the signed complement system]
\label{sbcouple:mass-relative}
Let $M(t)>0$ be a strong longitudinal reference mass and put
$r_M=M'/M$.  Define
\[
 \widehat Q:=\frac{Q_E}{M},\qquad
 \widehat N:=\frac{N}{M},\qquad
 \widehat D:=\frac{D_E}{M}.
\]
Then \eqref{sbcouple:matrix} implies
\[
\boxed{
 \begin{pmatrix}\widehat Q\\\widehat N\\\widehat D\end{pmatrix}'
 \preceq
 \begin{pmatrix}
 \theta_I-a_I-r_M&\max(b_{I2},b_{I4})&\beta_I\\
 a_J&\theta_J-r_M&\beta_J\\
 \zeta_Q&\zeta_N&d_D-r_M
 \end{pmatrix}
 \begin{pmatrix}\widehat Q\\\widehat N\\\widehat D\end{pmatrix}
 +\frac1M\begin{pmatrix}S_Q\\S_N\\S_D\end{pmatrix},
}
\tag{mass-relative-comparison}\label{sbcouple:mass-relative-comparison}
\]
where $d_D$ is either the displayed mismatch diagonal in
\eqref{sbcouple:matrix} or the retained-damping replacement
\eqref{sbcouple:damping}.  Dividing all three variables by the same $M$
leaves every off-diagonal coefficient unchanged and subtracts $r_M$ from
every diagonal; there is no time-dependent weight derivative.

If $r_M\ge g_M$ and a constant positive left weight $\omega$ satisfies
\[
 \omega^\top\overline C_{\rm rel}\preceq
 -\kappa_{\rm rel}\omega^\top
\]
on a stopped rectangle, then
\[
 V_{\rm rel}:=\omega_Q\widehat Q+\omega_N\widehat N+\omega_D\widehat D
\]
obeys the scalar Duhamel bound
\[
\boxed{
 V_{\rm rel}(t)
 \le e^{-\kappa_{\rm rel}(t-a)}V_{\rm rel}(a)
 +\int_a^t e^{-\kappa_{\rm rel}(t-s)}
 \frac{\omega^\top S(s)}{M(s)}\,ds.
}
\tag{mass-relative-Duhamel}\label{sbcouple:mass-relative-Duhamel}
\]
\end{lemma}
\begin{proof}
Differentiate $Q_E/M$, $N/M$ and $D_E/M$ and insert
\eqref{sbcouple:matrix}.  Positivity of the off-diagonal coefficients then
allows the constant-weight scalar comparison.
\end{proof}

\begin{remark}[The $5/9$ threshold is one concept-rate ratio]
\label{sbcouple:spectral-check}
In the ideal principal coefficients,
\[
 g_M\simeq2\ell_1(1-m),\qquad
 \theta_J\lesssim\ell_1(1-m)+\ell_2,
\]
so the mass-relative mixed-sign diagonal satisfies
\[
 \theta_J-g_M\lesssim-\{\ell_1(1-m)-\ell_2\}.
\]
It is therefore favorable precisely when
\[
 m<1-\frac{\ell_2}{\ell_1}.
\]
For $\ell_1=9/2$ and $\ell_2=2$ this is $m<5/9$.  The same number appears in
the collective holding asymptotics because
$\varrho^2=\ell_1/\ell_2$; these are two formulas for the same rate
comparison, not independent confirmations.  Numerically the principal
mixed-sign margins are $0.53125$, $0.385$ and $0.25$ at
$m=0.4375,0.47,0.50$ respectively.  The old observation that the unnormalized
Metzler matrix cannot contract when $\theta_J>0$ remains correct; the common
mass normalization is exactly the missing diagonal subtraction.
\end{remark}

\begin{remark}[An algebraic check that the retained damping can matter]
\label{sbcouple:spectral-check-legacy}
Consider the limiting noiseless matched configuration with a pure strong
family of mass $M$, and first-quadrant excluded inputs approaching the
weak axis with total mass $R$. Evaluate the coefficients at $Q_E=N=D_E=0$
by this one-sided limit; take $\ell_1=9/2$, $\ell_2=2$ and the conservative
target-energy scale $K=c_\delta=\sqrt{\ell_1(\ell_1+\ell_2)}$.
There are no mixed-quadrant receivers in this calculation.
Normalize the two remaining variables as $Q_E/M$ and $D_E/\sqrt M$,
using the ideal strong logarithmic mass rate $g=2\ell_1(1-M)$.
The resulting local comparison matrix is
\[
 C_{\rm rel}=\begin{pmatrix}
 \ell_1(1-M)+\ell_2(1-R)-\ell_1R-g
       &(K+\ell_1R)\sqrt{R/M}\\
 (\ell_1-\ell_2)\sqrt{RM}
       &-\min\{\ell_1(1-M),\ell_2(1-R)\}+KR-g/2
 \end{pmatrix}.
\]
At $M=1/2$, $R=1/10$, its trace and determinant are exactly
\[
 \operatorname{tr}C_{\rm rel}=-\frac{99}{20}+\frac{3\sqrt{13}}{20}<0,
 \qquad \det C_{\rm rel}=\frac{1413}{400}-\frac{51\sqrt{13}}{100}>0.
\]
Its larger eigenvalue is approximately $-0.42511$. This demonstrates the
qualitative benefit of retained damping for this mass-relative
\emph{matched-sector two-variable subsystem}, with $R_J=0$.
Replacing the residual damping by the loss-norm coefficient $K$ instead
gives a larger eigenvalue approximately $3.99817$, destroying even that
matched-sector stabilization.

Neither number implies contraction of the full physical-time
three-variable comparison. Its matrix $C$ is Metzler, so
$s(C)\ge\max_i C_{ii}\ge\theta_J$: shift $C$ to a nonnegative matrix
and use that its spectral radius dominates each diagonal entry.
Consequently a positive $\theta_J$ precludes asymptotic contraction of
that unnormalized comparison.  Lemma~\ref{sbcouple:mass-relative} supplies
the required common-mass normalization and subtracts the strong logarithmic
rate from every diagonal.  The finite-window theorem still needs an
affordable stopped Duhamel budget with the actual source constants. This is a
frozen-coefficient algebraic check, not a trajectory, sign-invariance or
initialized entry certificate. Gate changes, noisy forcing, mixed
quadrants and changing coefficients must still be paid in
\eqref{sbcouple:solution}; Section~\ref{sbwindow:scope} begins that
finite-window calculation from initialization.
\end{remark}

The output charge and crossing margin are now fully specified by
\begin{align}
 \overline{\mathfrak X}_E(t)
 &=\sup_s\left\{G_+(s)\overline V_E(t)
       +\frac{h(s)}e\overline z_1(t)
       +\frac{K_2}e\sqrt{\overline R_E(t)}\,\overline z_3(t)\right\},
       \label{sbcouple:output-clock}\\
 \Gamma_u(\tau)
 &=q_-a_{u,\times}(\tau;\min(T,\overline Y_{\rm sign}))
       -\epsilon_{u,\rm other}^-(\tau)
       -q_+\overline{\mathfrak X}_E(2\tau/\mu_2^2).
       \label{sbcouple:clock}
\end{align}
These are evaluated in the same interval as the piecewise two-mean
comparison. Besides $\Gamma_u>0$, input crossing requires
$\int_{\tau_0}^{\tau_u}\Gamma_u\,d\tau\ge[-\min_Bs_i(\tau_0)]_+$.
The subsequent handoff still uses the maximum of concurrent mean,
output-crossing and floor-entry times, not their sum.

\subsection{Non-implications and the remaining initialized estimate}

\begin{proposition}[Two sign reductions that cannot be used]
\label{sbcouple:obstructions}
Neither of the following is a valid consequence of balanced SIM geometry:
(a) bounding general excluded self-production only by $O(eR_E^2)$ and
terms proportional to $Q_E,D_E$;
(b) bounding the strong skew source only by negative strong output and
negative-feature tails. Nor is matched first-quadrant self-skew confined
to negative-feature tails.
\end{proposition}
\begin{proof}
Use the two noiseless samples $3e_1,2e_2$ with equal weights, so
$\ell_1=9/2$, $\ell_2=2$.
For (a), take $B=W=\varnothing$ and two balanced excluded neurons
$v_1=z_1=(1,-2/9)$, $v_2=z_2=r(\eta,1)$.
Then $Q_E=r^2\eta$, $D_E=0$, $N_4=2/9$, and direct substitution gives
$\dot Q_E/r^2\to1$ as $\eta\downarrow0$.
The data-tail costs vanish and $R_E$ stays bounded. Thus the missing
$N_4$ term cannot be replaced by the proposed terms with uniformly
bounded coefficients. This is a static non-implication, not a trajectory
counterexample to the initialized theorem.

For (b), take a balanced strong neuron
$v_B=a e_1$, $z_B=a(\cos\vartheta,\sin\vartheta)$, $0<\vartheta<\pi/2$,
and a matched first-quadrant receiver. All strong outputs are nonnegative,
but its contribution is
$X_{i,B}=-\ell_1a^2\sin\vartheta\ne0$.
Choosing $\vartheta=O(e)$ makes this a noise-scale forcing, not zero.
Finally, for a single matched first-quadrant excluded neuron $v=z=(a,b)$,
its self contribution is $X_{i,E}=(\ell_2-\ell_1)ab$.
It is nonzero on strictly nonnegative data. Its production self term,
by contrast, is $-(\ell_1+\ell_2)(a^2+b^2)ab\le0$.
\end{proof}

\begin{lemma}[Initialization data uniform over the eventual fixed partition]
\label{sbcouple:init}
Under balanced isotropic initialization in dimension $d\ge2$, set
$L=\log(3/\delta_{\rm init})$. With probability at least
$1-\delta_{\rm init}$, simultaneously for every subset $E\subseteq[h]$,
even one selected later from the trajectory,
\begin{align*}
 Q_E(0)&\le\varepsilon^2\left\{\frac{h}{2\pi d}
       +\sqrt{\frac{hL}{2d(d+2)}}+\frac L3\right\},\\
 N(0)&\le\varepsilon^2\left\{\frac{h}{\pi d}
       +\sqrt{\frac{hL}{d(d+2)}}+\frac L3\right\},\\
 R_E(0)&\le\varepsilon^2\min\left\{h,\frac{2h}d
       +\sqrt{\frac{8h(d-2)L}{d^2(d+2)}}+\frac{2L}3\right\},
 \qquad D_E(0)=0.
\end{align*}
\end{lemma}
\begin{proof}
For a uniform unit vector $u$, its projected polar angle is uniform and
independent of its projected radius. Direct angular integration and the
spherical second/fourth moments give
\[
 \mathbb E(u_1)_+(u_2)_+=\frac1{2\pi d},\quad
 \mathbb E\{(u_1)_+(u_2)_+\}^2=\frac1{4d(d+2)},
\]
\[
 \mathbb E|u_1u_2|\mathbf1_{u_1u_2<0}=\frac1{\pi d},\quad
 \mathbb E\{|u_1u_2|\mathbf1_{u_1u_2<0}\}^2=\frac1{2d(d+2)},
\]
\[
 \mathbb E(u_1^2+u_2^2)=\frac2d,\quad
 \operatorname{Var}(u_1^2+u_2^2)=\frac{4(d-2)}{d^2(d+2)}.
\]
The first two variables lie in $[0,1/2]$, and the third in $[0,1]$.
Apply the bounded-variable Bernstein bound
$\sum_j(X_j-\mathbb EX_j)\le\sqrt{2h\operatorname{Var}(X_j)L}
+2bL/3$ with range bound $b$, and take a union bound over these three
whole-network sums. Use their second moments as variance upper bounds
in the first two. Every subset sum is dominated by its corresponding
whole-network nonnegative sum, so no union bound over partitions is
needed. Aligned initialization gives $\delta_i(0)=0$ exactly.
\end{proof}

At handoff, the installed coverage bounds control $Q_E(a_0)$ for
specified admission failures; $0\le N(a_0)\le R_E(a_0)/2$ is always valid.
Neither controls the mass of fully pure neurons excluded solely because
their future guards have not been proved. The existing pure-sector
amplification lemma is conditional on a fixed pure interval; it is not
a theorem admitting all such neurons through every later gate change.
T1.9 supplies its seed bounds, not those missing guards for the entire
candidate enlargement. Thus these hypotheses cannot be marked delivered
by maximality of the fixed enrollment.

Lemma~\ref{sbcouple:init} removes uncertainty about the initial values at
time zero. It does not assert those same bounds at $a_0$. One may apply
the general comparison from zero (with its general cohort sources before
the strong/weak guards are reached), or prove a sharper handoff coverage
theorem. In either case the growth and source integrals must be evaluated;
future selection of $E$ is not an excuse to assume it independent of
the initialized directions.

The next display records the entry budget at this stage of the proof.
Section~\ref{sbwindow:scope} supplies a from-zero sufficient calculation
and its explicit timing premises; the downstream relative feedback
frontier is \eqref{sbwindow:next-gap}. The entry budget
contains no hidden $\mathcal J_E^\theta$ or $Z_E$: supply a reached fixed
partition and coefficient/initial-data caps for \eqref{sbcouple:solution}
whose actual output charge satisfies, throughout the pre-crossing window,
\begin{equation}\label{sbcouple:unresolved-budget}
 q_+\sup_s\{eG_+(s)\overline V_E+h(s)\overline z_1
                  +K_2\sqrt{\overline R_E}\,\overline z_3\}
 <e\{q_-a_{u,\times}-\epsilon_{u,\rm other}^-\},
\end{equation}
and whose integrated margin reaches the initial negative slope before the
common progress, T1.9 and branch deadlines. Even $e\to0$ does not remove the
strong contribution to $S_D/e$, which contains $M\bar t$; its integrated
cost must actually be paid. The global radius bound
$R_E\le h\varepsilon^2+\lambda t$ is an admissible fallback, but does
not itself certify \eqref{sbcouple:unresolved-budget} or eligibility of
the excluded pure neurons. No numerical deficit, empty parameter region,
or failure of the initialized dynamics follows from these observations.
The signed-gap, generation, creation and general Theorem~3 still require
their reached-interval hypotheses; none is asserted here by continuation
past this unresolved budget.

\section{From-zero finite-window propagation and an entry certificate}
\label{sbwindow:scope}

The objective here is a finite-window output budget, not contraction of
all three complement variables. We use physical time $t$ for propagation
and normalized time $\tau=\mu_2^2t/2$ for entry. Work on the common SIM coordinate envelope, so the dominant coordinate
is positive on each cluster. The eventual partition is fixed even when
selected retrospectively. No condition $N(0)=0$ or
$N(a_0)=0$ is imposed.

\subsection{The mixed-sign channel as a finite-window source}

Reorder \eqref{sbcouple:matrix} as $x=(Q_E,D_E)^\top$ and $N$, writing
$x'\preceq A x+bN+f$, $N'\le\theta_JN+c^\top x+g$, with
\[
 b=\binom{\max(b_{I2},b_{I4})}{\zeta_N},\qquad
 c=\binom{a_J}{\beta_J},\qquad g=S_N.
\]
All off-diagonal coefficients and sources are nonnegative; the diagonal
of $A$ may retain \eqref{sbcouple:damping}. With
$G_J(t,s)=\exp(\int_s^t\theta_J)$, exactly the scalar comparison gives
\begin{equation}\label{sbwindow:N-duhamel}
 N(t)\le G_J(t,0)N(0)+\int_0^tG_J(t,s)
                    \{a_JQ_E+\beta_JD_E+S_N\}(s)\,ds.
\end{equation}
This retains the nonzero initialization term. The coefficients of the
injection vanish with $R_J$: $a_J\le\ell_*R_J$ and
$\beta_J\le K\sqrt{R_J}+\lambda R_J\sqrt{R_E}$, where
$\ell_*=\max(\ell_1,\ell_2)$.

For clarity, let $\Psi$ be the positive fundamental matrix of a certified
$A$ envelope and let $n_f=G_J(t,0)\bar N(0)+\int_0^tG_JS_N$.
Substitution in the $x$ comparison gives a positive Volterra comparison
\[
 x(t)\preceq x_f(t)+\int_0^t\mathcal K(t,v)x(v)\,dv,
\]
\[
 x_f=\Psi(t,0)\bar x(0)+\int_0^t\Psi(t,s)(f+b n_f)(s)\,ds,
 \quad
 \mathcal K(t,v)=\left\{\int_v^t\Psi(t,s)b(s)G_J(s,v)\,ds\right\}c(v)^\top.
\]
This follows by interchanging two nonnegative integrals. If in a chosen
positive weighted norm $\|\Psi(t,s)\|\le e^{a(t-s)}$,
$\theta_J\le d$, $\|b\|\le B$, and $\|c^\top\|\le C$, then on
$[0,T]$ the integral operator norm is at most
\[
 BC\,\mathcal I(a,d;T),\qquad
 \mathcal I(a,d;T)=\int_0^T\int_0^{T-r}e^{ar+du}\,du\,dr
 \le\tfrac12T^2 e^{T\max(a,d,0)}.
\]
If this is less than one, the usual geometric resolvent bound applies.
Otherwise the finite-window Volterra series still converges for bounded
coefficients; failure of this sufficient small-gain check is not a
trajectory incompatibility. Positive $\theta_J$ is priced by $G_J$, not
set to zero or required to be negative. Coefficients may change at a
reached handoff without adding a stage to the entry clock.

\subsection{A general envelope before any late cohort geometry is reached}

Put $s_0=h\varepsilon^2$ and retain the data spectral bound $\lambda$.
\begin{lemma}[Global exponential mass bound and early mismatch rate]
\label{sbwindow:mass}
For every finite time,
\[
 S(t)\le s(t):=s_0e^{2\lambda t},\qquad R_E(t)\le s(t).
\]
Also $\|A_i\|\le\lambda(1+s(t))$ and
$\lambda_{\min}(\operatorname{sym}A_i)\ge-\lambda s(t)$.
The latter can replace $-\mu_\rho$ in the mismatch comparison.
\end{lemma}
\begin{proof}
For $T_i=\mathbb E_n[xx^\top\alpha_i]$, $0\preceq T_i\preceq\Sigma_n$.
Balance and the selector identity give
\[
 S'=2\langle x,f\rangle_n-2\|f\|_n^2,
 \qquad \langle x,f\rangle_n=\sum_i z_i^\top T_i v_i\le\lambda S.
\]
Drop only the nonpositive square to obtain the exponential bound.
Furthermore $\|f_\parallel\|_n\le\sqrt\lambda S$, so
$\|\mathbb E_n[f_\parallel x^\top\alpha_i]\|\le\lambda S$.
Subtract this matrix from $T_i$ to prove both residual bounds. This
aggregate improvement does not extend the individual-neuron or angular
conclusions of T1.9 beyond their admissible horizon.
\end{proof}

Use whole-network nonnegative quantities
\[
 H=Q_{[h]}+N_{[h]},\qquad D=D_{[h]},\qquad Z=\sqrt{s}\,D.
\]
Then $Q_E,N_E\le H$, $D_E\le D$ for every eventual subset $E$.
Let $\chi_{\rm ax}$ denote the fixed axis-product error called $\chi$ in
Section~\ref{sbcouple:scope}; it is distinct from the time-dependent
Riccati quantity in T1.5. Define
\[
 \Theta=C_x,\qquad b_\times=\sup_{\alpha\ {
m compatible}}
       |\mathbb E_n[x_1x_2\alpha]|,
 \qquad c_\delta=\sqrt{\lambda C_x}.
\]
Here the supremum is over homogeneous input gates and their compatible
boundary selections. All these are fixed empirical data constants.

\begin{lemma}[Whole-network comparison retaining the target PSD structure]
\label{sbwindow:early-system}
Almost everywhere,
\begin{align}
 H'&\le(\Theta+\ell_*s)H+\lambda(1+2s)Z
            +b_\times s+\chi_{\rm ax}s^2+\lambda s_0s,
                  \label{sbwindow:H}\\
 Z'&\le\ell_*sH+\{\lambda+(\lambda+c_\delta)s\}Z
                   +\chi_{\rm ax}s^2+\lambda s_0s.
                  \label{sbwindow:Z}
\end{align}
In particular these are valid from zero, without any $B/W$ sign guards.
\end{lemma}
\begin{proof}
Apply the cohort calculation with $E=[h]$, $B=W=\varnothing$, but retain
$T_i$ exactly. Its contribution to either product weight has trace at
most $\Theta$ and off-diagonal magnitude at most $b_\times$. The two
sign regions are disjoint, so their target forcing sums to at most
$b_\times R$, not twice that amount. The matched self diagonal terms
are nonpositive. The remaining adverse self cross terms are at most
$\ell_*R H$; their axis-replacement error is at most
$\chi_{\rm ax}R^2$. In particular there is no spurious
$\chi_{\rm ax}R$ target-replacement term.

Before applying Cauchy--Schwarz, combine the receiver mismatch sums over
both sign regions. Their total is at most
$\lambda(1+s)\sqrt{s}D$; the donor mismatch costs
$\lambda s^{3/2}D$. The ambient contribution is at most
$\sqrt{s}\mathcal P_{[h]}\le\lambda s_0s$ by
\eqref{sbnew:P-bound}. This proves \eqref{sbwindow:H}.
For mismatch, the evaluated skew bound with no external cohort gives
$Z_{[h]}^{\rm skew}\le\ell_*\sqrt{s}H+
\chi_{\rm ax}s^{3/2}+c_\delta sD$.
Use the $\lambda s$ rate from Lemma~\ref{sbwindow:mass}, multiply by
$\sqrt{s}$, and differentiate $\sqrt{s}$ with logarithmic rate
$\lambda$. This proves \eqref{sbwindow:Z}.
\end{proof}

Assume $\lambda<\Theta<2\lambda$ and define
\[
 \omega_0=\frac\lambda{\Theta-\lambda},\quad
 \kappa_0=2\lambda-\Theta>0,\quad
 \gamma_0=\max\left\{\ell_*(1+\omega_0),
                \frac{2\lambda}{\omega_0}+\lambda+c_\delta\right\}.
\]
For $F=H+\omega_0Z$, the two inequalities imply
\[
 F'\le(\Theta+\gamma_0s)F+
       \{b_\times+(1+\omega_0)\lambda s_0\}s
       +(1+\omega_0)\chi_{\rm ax}s^2.
\]
Choose $H_0$ as the smaller of $s_0/2$ and the sum of the two
high-probability initialization bounds in Lemma~\ref{sbcouple:init}.
Aligned initialization gives $Z(0)=0$, whereas $N(0)$ is included in $H_0$.
An explicit supersolution is
\begin{align}
 \overline F(t)=\exp\left(\frac{\gamma_0(s(t)-s_0)}{2\lambda}\right)
 \Bigg[&H_0e^{\Theta t}
 +\{b_\times+(1+\omega_0)\lambda s_0\}
             \frac{s(t)-s_0e^{\Theta t}}{\kappa_0}\notag\\
 &+(1+\omega_0)\chi_{\rm ax}
             \frac{s(t)^2-s_0^2e^{\Theta t}}{4\lambda-\Theta}\Bigg].
 \label{sbwindow:closed-envelope}
\end{align}
Indeed variation of constants and
$\exp(\gamma_0\int_u^t s)\le\exp(\gamma_0\int_0^t s)$ give the
formula after integrating the two elementary exponentials. Hence,
for every fixed retrospectively selected $E$,
\begin{equation}\label{sbwindow:subset-envelope}
 Q_E,N_E\le\min\{\overline F,s/2\},\quad
 D_E\le\frac{\overline F}{\omega_0\sqrt{s}},\quad
 R_E,V_E\le s.
\end{equation}
The corresponding completely evaluated output charge is
\begin{equation}\label{sbwindow:charge}
 \overline{\mathfrak X}_E(t)
 \le G_*s(t)+\left(h_*+\frac{K_2}{\omega_0}\right)
                         \frac{\overline F(t)}e,
 \quad G_*:=\sup_sG_+(s),\quad h_*:=\sup_s h(s).
\end{equation}
This bound can be used throughout a short window or as the initial data
for the sharper reached-cohort comparison. Neither use requires $R_J=0$.

\subsection{An explicit small-output region for the finite-window charge}

The following deliberately conservative constants supply a verifiable
entry calculation, not a claim that the canonical experiment satisfies
these sufficient bounds. Set $\mu_2=2$, $0<e\le1/40$, $\rho=10^{-4}$,
and suppose the fixed data constants satisfy
\begin{equation}\label{sbwindow:data-size}
 \begin{gathered}
 4.4\le\lambda\le5,\quad 2\le\omega_0\le3,\quad
 \kappa_0\ge2,\quad \ell_*\le5,\quad c_\delta\le6,\quad
 \gamma_0\le20,\\
 b_\times\le2.1e,\quad \chi_{\rm ax}\le12e,\quad
 G_*\le1,\quad h_*\le1.2,\quad K_2\le2.
 \end{gathered}
\end{equation}
For a specified physical horizon $T$ require only the explicit
initialization smallness conditions
\begin{equation}\label{sbwindow:initial-smallness}
 s_0e^{2\lambda T}\le\rho,\qquad
 s_0e^{\Theta T}\le e\rho.
\end{equation}
Then for every $t\le T$, \eqref{sbwindow:closed-envelope} gives
\begin{equation}\label{sbwindow:small-charge}
 \frac{\overline F(t)}e<1.56\,10^{-4},\qquad
 \overline{\mathfrak X}_E(t)<4.44\,10^{-4}.
\end{equation}
For example the first bound follows by substituting
\[
 e^{20\rho/8.8}\left\{\frac\rho2+
                \frac{(2.1+20\rho)\rho}{2}
                    +\frac{48\rho^2}{10.8}\right\}<1.56\,10^{-4}.
\]
The second follows from $\rho+2.2\overline F/e$.
Thus a constant-scale radius cap suffices here; no condition
$R_E/M=O(e)$ has been used. These are bounds at every intermediate time,
not just at the handoff endpoint.

\subsection{Consuming the output bound in the exact entry equations}

Let $B$ be a fixed family whose broad tube and projected radial matching
have been delivered by T1.9 on $[a_0,T]$. No enlargement beyond already
certified members is needed for this certificate. Empty lagging family
$\mathcal L$ is permitted. Suppose on that interval
\begin{equation}\label{sbwindow:tube}
 |s_i|\le S_b:=0.30,\quad |t_i|\le T_b:=0.32,\quad
 \left|1-\frac{\|z_i\|^2}{\|v_i\|^2}\right|\le10^{-4},
 \quad v_{i1},z_{i1}>0.
\end{equation}
Then $|q_i-1|<10^{-4}$: use the displayed radial guard and
$q_i=(\|z_i\|/\|v_i\|)\cos\psi_i/\cos\phi_i$,
with both angles at most $\arctan(eT_b)$ in magnitude.
This is a consequence of the reached T1.9 tube, not a sign condition on E.
All cluster-1 points must be active throughout the broad input cone.

In addition to \eqref{sbwindow:data-size}, assume the empirical reference
quantities satisfy the following data-side inequalities:
\begin{equation}\label{sbwindow:entry-data}
 \begin{gathered}
 2.24\le a\le2.26,\quad |b|\le0.01,\quad
 c\le2e^2,\quad 0\le Z(s)\le2,\quad |b+g(s)|\le1,\\
 b+K(s_*)=as_*,\quad s_*\in[0.233,0.235],\quad K(S_b)\le0.58,\\
 h(S_b)\le0.625,\quad h(S_b)-h(-S_b)\le0.25,\quad
 h(S_b)-h(0)\le0.125,\quad h(0)\le0.51,\quad h(0.135)\le0.56,\\
 \inf_{[-S_b,0]}(b+g(s)-as)\ge0.39,\quad
 b+K(-S_b)\ge0.265,\quad b+K(0)\ge0.39,\\
 \inf_{[0,0.135]}(b+g(s)-as)\ge0.085.
 \end{gathered}
\end{equation}
Unspecified ranges of $s$ in this display are $[-S_b,S_b]$.
These conditions concern the fixed sample and its one-dimensional kernels.
They keep the signed empirical $b$ in the principal field. They are
sufficient data guards, not sampling probabilities by themselves.

For the nonplanar terms require
\begin{equation}\label{sbwindow:ambient}
 \max_B\frac{2|P_{i2}-es_iP_{i1}|}{e\mu_2^2v_{i1}}\le10^{-4},\qquad
 \max_B\frac{2|P_{i1}|}{\mu_2^2v_{i1}}\le10^{-4}.
\end{equation}
These vanish in the planar case. If T1.9 gives
$r_i\ge\varepsilon\tau_d$ throughout the interval, a sufficient explicit
condition for both is
\[
 s_0e^{\lambda T}\le
 \frac{10^{-4}e\mu_2^2\tau_d}
                  {2\lambda\sqrt h(1+e^2S_b^2)}.
\]
Indeed $\|P_i\|\le\lambda s_0\sqrt{s}$ by the aggregate bound and
$v_{i1}\ge\varepsilon\tau_d/\sqrt{1+e^2S_b^2}$, so the left directional
quantity is at most
$2\lambda\sqrt h(1+e^2S_b^2)\sqrt{s_0s}/(e\mu_2^2\tau_d)$.

\begin{lemma}[The required directional and mean errors are delivered]
\label{sbwindow:errors}
Under \eqref{sbwindow:data-size}--\eqref{sbwindow:ambient}, the exact
empirical equations on $B$ have
\[
 |r_{u,i}|,|r_{w,i}|<0.0025,\qquad |r_s|,|r_t|<0.004.
\]
In particular total individual one-sided error allowances $0.020$ are
valid, including the complement charge; it must not be added a second time.
\end{lemma}
\begin{proof}
The whole-network axis comparison, including its mismatch, gives the
physical off-diagonal learned-output moment bound
\[
 B_\times^{\rm out}/e\le
 \ell_*\overline F/e+(\chi_{\rm ax}/e)s
                   +(\lambda/\omega_0)\overline F/e<0.00237.
\]
Every diagonal learned moment is at most $\lambda s$. Relative to the
$M=0$ empirical target field, both normalized ratio errors from the
network are therefore at most
\[
 1.000101\frac2{\mu_2^2}
 \{(1+e^2S_bT_b)\,0.00237+\lambda\rho(S_b+T_b)\}<0.001341.
\]
The exact finite-noise target terms cost at most
$1.000101e^2\{2T_b+2S_b+S_bT_b\}<0.000836$.
Restoring the actual $M$ principal field adds at most
\[
 1.000101\rho\max\{|b|+K(S_b)+aS_b,\,2aT_b\}<0.000145,
\]
where $M\le S\le\rho$. Including \eqref{sbwindow:ambient} gives less
than $0.0025$. No weak sign or complement eligibility has been assumed.

For the exact logarithmic weight rate $k_i$, write its full residual
formula with target diagonal $a+e^2Z(s_i)$ and the learned moments.
The slope-independent $a$ pays only
$\operatorname{osc}(q+q^{-1})\le10^{-8}/(1-10^{-4})$.
The remaining diagonal, cross, and ambient oscillations are bounded by
\[
 2.26\frac{10^{-8}}{1-10^{-4}}+0.002501+0.001001+0.000778+0.0002
 <0.0045.
\]
Here the two middle diagonal terms bound $e^2Z(q+q^{-1})$ and the learned
longitudinal moment respectively. The cross terms use
$e^2(S_b/q_-+T_bq_+)(1+0.00237/2)$, doubled for oscillation.
Also $|F_i|<1.879$, $|G_i|<1.724$. In the exact mean remainders, the
$q-1$ terms are therefore less than $0.000201$ and covariance is at most
$\max(S_b,T_b)\operatorname{osc}k/2<0.00072$.
Together with the individual errors these prove the stated mean bounds.
\end{proof}

\begin{proposition}[A concurrent entry window with a priced complement]
\label{sbwindow:entry}
The numerical duration below starts at a certified broad-tube time. It is not a dimension-uniform bound on the time from initialization to signed feedback.

Assume the preceding data and initialization guards, and that T1.9
certifies \eqref{sbwindow:tube} through some
$T\ge a_0+4.2(2/\mu_2^2)$. Then the fixed family $B$ enters the retuned
positive rectangle and its two mean intervals no later than
\[
 t_{\rm hand}=a_0+4.1(2/\mu_2^2).
\]
Throughout $[0,T]$, $M\le\rho<1/2$, so this handoff is strictly before
the raw progress ceiling. It is also strictly before the stipulated
certified T1.9 horizon. No assertion about an independently undefined
branch stopping time is implicit in this statement.
\end{proposition}
\begin{proof}
Write $v=\mu_2^2(t-a_0)/2$, and use conservative $M_+=0.05$,
$q_-=0.99$, $q_+=1.01$, and individual error $0.020$.
The concurrent negative-output comparison gives
\[
 Y(v)\le[(0.32+d/c)e^{-cv}-d/c]_+,\quad
 c=\frac{2.24(0.95)}{1.01},\quad
 d=\frac{0.265-2.26(0.05)(0.32)}{1.01}-0.020>0.
\]
It vanishes at $v_Y=c^{-1}\log(1+0.32c/d)<0.689$.
At the negative input minimum, \eqref{sbwindow:entry-data} gives
\[
 \Gamma_u(v)\ge A-0.5049Y(v),\quad
 A=0.99\{0.95(0.39)-0.05(0.30)(0.51)\}-0.020=0.3392215.
\]
It is positive even at $v=0$. Moreover
\[
 \int_0^{1.02}\Gamma_u(v)\,dv
 \ge1.02A-0.5049\{0.32/c-dv_Y/c\}>0.3033>0.30.
\]
Thus input crossing occurs by $v_u=1.02$; output nonnegativity was already
certified concurrently. For a nonnegative input minimum below $0.135$,
the same calculation gives slope at least
\[
 \gamma_{u,\ell}=0.99\{0.95(0.085)-0.05(0.30)(0.56)\}-0.020
                  =0.0516265.
\]
For an output minimum below $0.122$, input nonnegativity gives slope at
least
\[
 \gamma_{w,\ell}=\{0.39-2.26(0.05\cdot0.32+0.122)\}/1.01-0.020
                  >0.05734.
\]
Consequently the two positive floors are reached by $v<3.635$ and
$v<3.148$, respectively. The same face inequalities preserve them.

For the two means use the unweighted maximum norm, equivalently Perron
weight $\vartheta_{\rm wc}=1$. Monotonicity of $K$, $K'=h$ almost
everywhere, and $h\le0.625$ give the cooperative contraction rate
\[
 \lambda_m=(1-0.05)(2.24-0.625)=1.53425.
\]
The initial mean error is at most $0.555$. Before input crossing the
geometric forcing is at most
$(S_b+T_b/2)\,0.25=0.115$, and afterward it is at most
$(S_b/4+T_b/2)\,0.125=0.029375$.
Including the proved $0.004$ mean errors, a single concurrent comparison
therefore uses forcing $0.119$ until $v=1.02$ and $0.033375$ thereafter.
Its value is less than $0.178$ at $1.02$, and less than $0.0232$ at $4.1$.
Since $s_*\in[0.233,0.235]$, this places both means strictly within
$[0.193,0.285]$ and $[0.208,0.264]$, respectively. The broad upper bounds
are already smaller than $U=0.337$, $W=0.347$, and the two floors supply
$L=0.135$, $l=0.122$. The inward lower faces and the mean envelope preserve this rectangle
and these mean bounds through $T$. All entry clocks have run in parallel.
\end{proof}

\paragraph{Scope of this initialized certificate.}
The complement and mean error budgets in Proposition~\ref{sbwindow:entry}
are consequences of time-zero initialization and data constants. They
are not extra trajectory hypotheses. The remaining premise is the
specified reached T1.9 broad tube and its duration, which must be checked
using that lemma's primitive admissibility conditions. For a fixed
admissible horizon, \eqref{sbwindow:initial-smallness} and the ambient
inequality have an explicit nonempty small-$\varepsilon$ interval.
This observation does not prove a simultaneous region when the horizon,
weak retention, fine-locking tolerance and later learning deadline also
vary with the primitive parameters.

Section~\ref{sbguard:scope} resolves the formerly undefined branch
symbol into the actual holding and generation exits. The former is
certified on the priced feedback window; the latter still requires a
reached generation interval and its distinct continuation estimates.

\subsection{Consuming entry in the signed-feedback inequality}

The absolute output budget above is not a relative polarization estimate.
The next calculation makes that distinction quantitative. After entry,
$B$ has positive coordinates and is active on all points with $p=(x_1)_+>0$.
Let $\gamma_i=\mathbb E_n[p^2\alpha_i]/D_1\in[0,1]$. For the
complement $B^c$,
\[
 \sum_{B^c}\gamma_i[-z_{i2}v_{i1}]_+
 \le N_{B^c}+\sqrt{R_{B^c}}D_{B^c}.
\]
The matched negative products are precisely the two mixed-sign weights;
the difference is bounded by $|\delta_{i2}v_{i1}|$.
Since $N_B=0$ and
$Q_B=e\sum_Ba_is_i/q_i\ge eM\bar s/q_+$, the joint functional gives,
for $\omega_0\ge1$,
\[
 N_{B^c}+\sqrt{R_{B^c}}D_{B^c}
 \le H-Q_B+Z\le F-Q_B.
\]
Therefore the reached mean box supplies the rigorous bound
\begin{equation}\label{sbwindow:feedback}
 c_{21}^O\ge eM\left(\bar t+\frac{\bar s}{q_+}\right)-\overline F.
\end{equation}
This subtraction is essential: bounding $H$ and $Z$ separately would
charge the same joint envelope twice and would count the favorable B
cross moment as an adverse complement contribution.

Nevertheless, the scalar whole-network envelope is not sharp enough to
certify the next relative inequality using the box's guaranteed means.
It would require
\begin{equation}\label{sbwindow:first-unclosed}
 \frac{\overline F(t)}e
 <\left(0.208+\frac{0.193}{q_+}\right)\underline M_B(t).
\end{equation}
Initialization smallness alone cannot make this relative cost vanish.
For example, in the high-SNR equal-noise population reference with
$\mu_1=3$, $\mu_2=2$, take the exact spectral bound
$\lambda=9/2+\sigma^2$. Then $\kappa_0=5/2$, and the test gate $x_2>0$
gives the exact lower bound
\[
 \frac{b_\times}{e}\ge\frac{\mu_1\mu_2}{2\sqrt{2\pi}}
                      =\frac3{\sqrt{2\pi}}.
\]
Indeed the cluster-1 contribution is
$\mu_1\sigma/(2\sqrt{2\pi})$ and the cluster-2 contribution is zero.
Consequently the explicitly chosen upper envelope itself obeys
\[
 \frac{\overline F(t)}{e s(t)}\ge
 \frac6{5\sqrt{2\pi}}(1-e^{-5t/2})>0.4758
                 \qquad(t\ge2.05).
\]
But $\underline M_B\le M_B\le s$, and the coefficient certified by the
retuned mean box with $q_+=1.0001$ is less than $0.401$.
Thus this particular scalar certificate misses the sufficient feedback
inequality by at least $0.0748\,s(t)$ in units of $\overline F/e$, even
before demanding a quantitative positive $\zeta_*$.
This is a failure of the unsigned whole-network envelope as a
\emph{relative feedback certificate}, not a failure of the initialized
flow, the finite-window entry comparison, or the general signed system.
Its source is the global worst gated target cross moment. Taking
$\varepsilon$ smaller does not repair that coefficient ratio.

The minimal next propagation refinement is available without assuming
mixed-quadrant exclusion. Split $J$ into
$J_4=\{v_1>0,v_2<0\}$ and $J_2=\{v_1<0,v_2>0\}$, always inside fixed E,
and use the definitions of $b_{G2},b_{G4},a_G,\theta_G,\beta_G$ for each
sign set separately. Keeping the favorable self terms gives
\begin{align}
 N_4'&\le(\theta_{J_4}-b_{J_4,4})N_4+a_{J_4}Q_E+
             \beta_{J_4}D_E+S_{J_4},\label{sbwindow:N4}\\
 N_2'&\le(\theta_{J_2}-b_{J_2,2})N_2+a_{J_2}Q_E+
             \beta_{J_2}D_E+S_{J_2},\label{sbwindow:N2}\\
 S_G&=[b_{G2}c_2+b_{G4}c_1]_+
             +\chi_{\rm ax}(1+U)R_G+\sqrt{R_G}\mathcal P_E.\nonumber
\end{align}
To derive these, use the exact mixed-quadrant product derivatives from
Section~\ref{sbcouple:scope}; drop only the cross terms
$-b_{J_4,2}N_2$ and $-b_{J_2,4}N_4$. Together with the separate Q and D
rows this remains a cooperative finite-window comparison. In the
noiseless pure quadrant-IV sector,
$a_{J_4}=b_{J_4,4}=\ell_1\sum_{J_4}v_{i2}^2$; quadrant II has the
mirror formula. These are transverse energies rather than the full
mixed-sector radial masses. The installed pure-sector amplification
lemma can control them while its purity hypotheses hold, but it does
not by itself control moving sign sets through all gate changes.

At this stage the missing estimate is a signed, mass-relative bound on
the harmful complement (or the separate mixed-sign sources), sufficient
for positive $c_{21}^O$ with a quantitative margin. The entry certificate
has not supplied that relative bound. Nor does it establish every
hypothesis of the theta-free gap by itself. The next subsection supplies a different signed estimate and consumes it;
the unsigned comparison above is retained to document its limitation.

\subsection{A signed target refinement that removes the unsigned feedback deficit}

The deficit in \eqref{sbwindow:first-unclosed} comes from charging all
positive and negative target cross moments together. For Gaussian SIM,
negative input coordinates admit a sharper estimate. This refinement
uses the whole network and does not require exclusion of mixed quadrants.

Let $\Sigma_\sigma=\operatorname{diag}(\sigma_1^2,\sigma_2^2)$ and put
\[
 G(v)=\mathbb E[x(v^\top x)_+],\qquad
 \nu=\sup_{\|v\|=1}\|\mathbb E_n[x(v^\top x)_+]-G(v)\|.
\]
This is a fixed data discrepancy. On the uniform gated-moment event,
\[
 \nu\le\tfrac12\sum_{p=1}^2\left\{
 \varepsilon_{\rm op}^{(p)}+
 \mathbb E_p[\|x\|^2\mathbf1_{\{x\text{ outside the clipping event}\}}]
 \right\}.
\]
Here the clipping event and $\varepsilon_{\rm op}^{(p)}$ are those in
Lemma~\ref{lem:gated-moment-concentration}; the population tail correction
must be included because $G$ is the unclipped Gaussian target.
Compatible selections at a sample boundary do not change
$\alpha_i(x)v_i^\top x=(v_i^\top x)_+$.

\begin{lemma}[Negative-coordinate target growth and full-flow propagation]
\label{sbwindow:negative-coordinates}
For $v_k\le0$,
\[
 G_k(v)\ge\sigma_k^2v_k.
\]
Let $Y_k^-=(\sum_i(-v_{ik})_+^2)^{1/2}$ and
$d=D/\sqrt{s}$, $y_k=Y_k^-/\sqrt{s}$. On the early mass envelope,
\begin{align}
 d'&\le-\{\lambda-(\lambda+c_\delta)s\}d
                 +\ell_*H+\chi_{\rm ax}s+\lambda s_0,
                      \label{sbwindow:d-sharp}\\
 D^+y_k&\le-(\lambda-\sigma_k^2)y_k+
                \nu+\lambda s+\lambda(1+s)d+\lambda s_0.
                      \label{sbwindow:negative-comparison}
\end{align}
For the data packet above, $s\le\rho\le10^{-4}$ implies
$d<1.78\rho$. If also $\max_k\sigma_k^2\le0.0025$, then
\begin{equation}\label{sbwindow:negative-envelope}
 y_k(t)\le e^{-b_0t}+\frac{\nu+19\rho}{b_0},\qquad b_0=4.3975.
\end{equation}
These bounds hold from zero for the entire initialized network.
\end{lemma}
\begin{proof}
For $x=\mu_p e_p+\xi$, Gaussian integration by parts gives
\[
 \mathbb E_p[x(v^\top x)_+]
 =\Sigma_\sigma v\,\mathbb P_p(v^\top x>0)
            +\mu_p e_p\,\mathbb E_p[(v^\top x)_+].
\]
One may justify the identity by smoothing the Lipschitz positive-part
function and passing to the limit. The mean term in each coordinate is
nonnegative. Since the mixture's activation probability lies in $[0,1]$,
its covariance term in a nonpositive coordinate is at least
$\sigma_k^2v_k$. This proves the target inequality without taking an
absolute value of its off-diagonal moment.

Write $U_i=\mathbb E_n[f_\parallel x^\top\alpha_i]$ and
$T_i=\mathbb E_n[xx^\top\alpha_i]$. The exact input equation is
\[
 \dot v_i=T_iv_i-U_i^\top v_i-A_i^\top\delta_i-P_i.
\]
The negative-coordinate energy inequality, the target bound, and
Cauchy--Schwarz therefore give
\[
 D^+Y_k^-\le\sigma_k^2Y_k^-+
       (\nu+\lambda s)\sqrt{s}+\lambda(1+s)D+
                       \lambda s_0\sqrt{s}.
\]
This includes every neuron and every changing gate. Dividing by
$\sqrt{s}$ proves \eqref{sbwindow:negative-comparison}.
Equation~\eqref{sbwindow:d-sharp} is \eqref{sbwindow:Z} divided by $s$,
using $Z=sd$. Since $H\le s/2$, $s_0\le\rho$, $\chi_{\rm ax}\le0.3$,
and $d(0)=0$, scalar comparison gives
\[
 d\le\frac{(2.5+0.3+5)\rho}{4.4-11\rho}<1.78\rho.
\]
The remaining forcing in \eqref{sbwindow:negative-comparison} is at most
$\nu+\{5+5(1.0001)(1.78)+5\}\rho<\nu+19\rho$.
Finally $y_k(0)\le1$ deterministically and
$\lambda-\sigma_k^2\ge b_0$, which proves the envelope.
\end{proof}

The harmful complementary cross output now has the bound
\[
 \sum_{B^c}\gamma_i[-z_{i2}v_{i1}]_+
 \le\sqrt{s}(Y_1^-+Y_2^-+D).
\]
Indeed the matched negative product is bounded by
$(v_{i1})_+(-v_{i2})_++(-v_{i1})_+(v_{i2})_+$,
and the mismatch by $|\delta_{i2}v_{i1}|$. Summing and using
Cauchy--Schwarz proves the claim. Consequently, after the reached
positive mean bound $\bar t\ge0.208$,
\begin{equation}\label{sbwindow:signed-feedback}
 c_{21}^O\ge0.208eM-s\{2e^{-b_0t}+11\rho+0.456\nu\}.
\end{equation}
Unlike the unsigned envelope, the normalized adverse target term here
decays exponentially. No $N=0$ invariant has been invoked.

\begin{lemma}[Longitudinal mass for the reached seed family]
\label{sbwindow:seed-mass}
On the broad tube, data, ambient and mass guards above,
\[
 M(t)\ge M(a_0)e^{g_B(t-a_0)},\qquad g_B=8.957.
\]
If $|B|\ge c_{\rm seed}h$ and T1.9 supplies
$r_i(a_0)\ge\varepsilon\tau_d$ for its members, define
\[
 m_0=\frac{0.9999\,c_{\rm seed}\tau_d^2}
                   {1+e^2S_b^2}\,e^{-2\lambda a_0}>0.
\]
Then, with $\zeta_0=1.043$,
\begin{equation}\label{sbwindow:mass-fraction}
 \frac{M(t)}{s(t)}\ge m_0e^{-\zeta_0(t-a_0)}.
\end{equation}
A sharper already-certified lower bound for $M(a_0)/s(a_0)$ may replace
$m_0$; independence of cohort selection is not required.
\end{lemma}
\begin{proof}
Cluster-1 purity gives $A_{i,11}\ge2a-\lambda\rho$.
Both off-diagonal entries have magnitude at most $b_\times+\lambda\rho$.
Thus $a_i=z_{i1}v_{i1}>0$ satisfies
\[
 \dot a_i\ge
 \{2a-\lambda\rho-eT_b(b_\times+\lambda\rho)\}
                        (v_{i1}^2+z_{i1}^2)-|z_{i1}P_{i1}|.
\]
The brace is positive. Use $v_{i1}^2+z_{i1}^2\ge2a_i$ and
$|P_{i1}|/v_{i1}\le0.0002$ from \eqref{sbwindow:ambient} to get
\[
 \dot a_i/a_i\ge
 4(2.24)-2(5)10^{-4}
 -2(0.025)(0.32)\{2.1(0.025)+5\,10^{-4}\}-0.0002
 =8.957952>g_B.
\]
Summation proves the first assertion. The radial and $q_i$ guards give
$M(a_0)\ge0.9999c_{\rm seed}h\varepsilon^2\tau_d^2/(1+e^2S_b^2)$.
Since $2\lambda-g_B\le1.043$, division by $s(t)$ proves
\eqref{sbwindow:mass-fraction}.
\end{proof}

\begin{proposition}[A finite signed-feedback window priced from initialization]
\label{sbwindow:positive-feedback}
This unprojected version has an explicit dimension-dependent feedback time through $m_0$; see \eqref{sbguard:old-time}--\eqref{sbguard:old-rho}. Corollary~\ref{sbguard:projected-feedback} gives a sharper projected-energy version.

Assume the preceding seed and data guards, with
$\max_k\sigma_k^2\le0.0025$, throughout $[a_0,T]$.
For equal noise this last condition follows from $e\le1/40$ and $\mu_2=2$.
Let $t_{\rm hand}=a_0+4.1(2/\mu_2^2)$ and choose
\begin{equation}\label{sbwindow:feedback-time}
 t_g\ge\max\left\{t_{\rm hand},\,
 \frac{\log(8/(0.208e m_0))-\zeta_0a_0}{b_0-\zeta_0}\right\},
 \qquad T>t_g,\quad T\ge a_0+4.2(2/\mu_2^2).
\end{equation}
Take $\rho\le10^{-4}$ in all earlier mass guards and impose
\begin{equation}\label{sbwindow:feedback-smallness}
 \rho\le\frac{0.208e m_0e^{-\zeta_0(T-a_0)}}{44},\qquad
 \nu\le\frac{0.208e m_0e^{-\zeta_0(T-a_0)}}{4(0.456)}.
\end{equation}
Then throughout $[t_g,T]$,
\begin{equation}\label{sbwindow:c21-positive}
 \boxed{c_{21}^O(t)\ge0.052eM(t)>0.}
\end{equation}
The exclusion charge is proved for the whole complement from time zero;
it is not an added purity, relative-mass or mismatch hypothesis on E.
\end{proposition}
\begin{proof}
The entry proof and its inward lower faces remain valid up to the
stipulated $T$. The post-crossing mean envelope decreases to
$0.033375/1.53425<0.0232$, so its positive mean bounds also persist.
Here $b_0-\zeta_0=3.3545>0$.
The time condition ensures
\[
 2e^{-b_0t}\le\tfrac14(0.208e m_0)e^{-\zeta_0(t-a_0)}
 \qquad(t\ge t_g).
\]
Each of $11\rho$ and $0.456\nu$ is at most one quarter of the signal
$0.208e m_0e^{-\zeta_0(t-a_0)}$ throughout $[t_g,T]$, by
\eqref{sbwindow:feedback-smallness}. The total adverse bracket in
\eqref{sbwindow:signed-feedback} is therefore at most
$\tfrac34(0.208)eM/s$ by \eqref{sbwindow:mass-fraction}.
This proves \eqref{sbwindow:c21-positive}.
\end{proof}

This proposition removes the specific coefficient obstruction of
\eqref{sbwindow:first-unclosed}. Its cost is an explicit additional
finite physical-time feedback horizon, smaller early total mass, and
uniform target-moment accuracy. It does not demand that the full
three-variable comparison be asymptotically contracting. The earlier
entry clocks remain concurrent; $t_g$ is a later signed-feedback time.
For fixed admissible $a_0,T,m_0,e$, the scalar upper bounds on $\rho$ and
$\nu$ are strictly positive. This is not yet the joint finite-sample and
T1.9 admissibility audit, or continuation to order-one progress.

\paragraph{The next inequality after consuming the signed source.}
Let $\mathsf F=\|f_O\|_n^2$, to distinguish this energy from the scalar
comparison $F$ above. The new source gives
$\mathsf F\ge D_1(0.052eM)^2$. Therefore the full-network direct gap
certificate, without identifying the seed family with a whole assigned
block, gives the sufficient bound
\begin{equation}\label{sbwindow:next-gap}
 \begin{split}
 g-\kappa_\perp\ge{}&
 \sum_p\rho_p\alpha_p\epsilon_p
 -(1-r^2)\{\theta\beta_1+(1-\theta)\beta_2\}\\
 &+(1-2r^2)\frac{D_1(0.052eM)^2}{2\mathsf T}-\mathcal E,
 \qquad 0<r^2=\mathsf T/\mathsf S\le\tfrac12.
 \end{split}
\end{equation}
Here $\mathsf S,\mathsf T,\rho_p,\theta,\mathcal E$ have exactly the
full-network meanings in Theorem~\ref{sbgap:certificate}; $\mathsf T$ is
transverse energy, not the physical horizon $T$.
It remains to prove that this entire right side has a positive lower
bound on a reached interval, including the longitudinal mass fractions,
polarization and transverse contributions of the complement. The seed
mass lower bound and positive $c_{21}^O$ alone do not certify that claim.
In particular this small-output certificate cannot be substituted for
creation near order-one strong progress. A common primitive region,
late continuation, generation, first-hit creation and general T3 remain
open; the specific unsigned feedback deficit is no longer the frontier.

\section{Continuation deadlines and the dimension audit}
\label{sbguard:scope}

This section audits the interval needed by the preceding feedback result
before using it in a grading theorem. The clock for entering a positive
strong-family box, the clock for positive signed feedback, and the clock
for creating a weak basin have different hypotheses. A creation conclusion
at its eventual hitting time is not an invariant assumed at early handoff.

\subsection{Which branch interval is actually used?}

The relevant dependencies are as follows.
\begin{enumerate}
\item T1.8 and Lemma~\ref{lem:post-trapping-locking} use the common
perturbative horizon $H_\star$, the radial/mismatch/angle guards, and,
for their joint strong--weak conclusion, the two weak angular buffers.
\item Sections~\ref{sbpos:scope}--\ref{sbemp:scope} use positive
longitudinal coordinates, longitudinal ratios, strong purity, a raw
progress ceiling and four exact error bounds. Theorem~\ref{sbemp:holding}
proves the eight faces inward from these hypotheses; it does not supply
the error bounds after the early comparison expires.
\item Theorem~\ref{sbgap:certificate} uses a full-network assignment,
$0<\mathsf T/\mathsf S\le1/2$, residual diagonal bounds and the stated
mismatch/error charges. Its positive gap is a conclusion to prove, not
an already available branch guard. The scalar signed-feedback proof
requires neither this assignment nor this positive gap.
\item Propositions~\ref{sbgen:first-hit} and~\ref{sbgap:deadline} use
progress monotonicity, a positive corridor reserve, a lower regulator
forcing bound, and a mass-weighted clock on an interval on which these
hypotheses have been reached. These are additional dynamical requirements.
\item Proposition~\ref{prop:gate-cliff-creation} requires late strong
holding, the charged-neuron sign/mismatch bound, weak-feature and adverse
shape budgets, two endpoint signs and the staircase margin at the
certification time. Data transition/DKW events do not expire with time;
the residual-dependent inequalities do.
\item The weak-rotation exit $t_\chi$ of Corollary~\ref{sbcoh:chi-deadline}
is an extra requirement for the stronger coherence route using $c_{12}^O$.
It is not used in the default $c_{21}^O$ calculation. Weak retention is
needed for subsequent capture; the new whole-complement feedback proof
does not require a retained weak cohort.
\item Proposition~\ref{sbnew:weak-comparison} requires a reached
three-sign field, transversality, branch-speed and perturbation bounds
on its own interval. It does not create that interval or provide an
independent pre-entry deadline.
\end{enumerate}
Thus the old symbol $\tau_{\rm branch}$ in
Remark~\ref{sbemp:combined-clock} had no defining downstream predicate.
It must be replaced by the concrete holding and generation exits below,
not by an arbitrary deterministic time.

Fix the initialized seed family $B$ and a time $a$ at which its broad
T1.9 tube has actually been reached. The notation $a_0$ in
Section~\ref{sbwindow:scope} means this time $a$. It need not equal
$t_0+\tau_{\rm gap}$: if necessary add the outer-cone entrance time in
Section~\ref{sbemp:scope}, and verify the radial guard at that later time.
No extra entrance time is silently free.

Let $\mathcal B$ be the retuned positive rectangle and two mean intervals
of Corollary~\ref{sbemp:design}. Write $\mathcal H_B(t)$ for their
membership, positive longitudinal coordinates, cluster-1 purity, the
specified $q$ guard, and $M\le M_+$. Define
\begin{equation}\label{sbguard:hold-exit}
 t_{\rm hold}:=\inf\{t\ge t_{\rm hand}:\mathcal H_B(t)
                                      \text{ fails}\}.
\end{equation}
Take $\inf\varnothing=+\infty$. This is the geometric holding exit
actually consumed by the feedback channel. One may certify it either
by the early tube or by the eight empirical face inequalities; it is
not defined to be $H_\star$.
Let $t_{\rm err}$ be the first essential failure of the four certified
error allowances in Theorem~\ref{sbemp:holding}. The word essential
accounts for equations that hold almost everywhere under compatible
Clarke selections. The first-exit proof there excludes geometric exit
strictly before the progress, longitudinal, purity, ratio or error exit.

\begin{proposition}[What the feedback window does and does not order]
\label{sbguard:early-order}
Under all premises of Proposition~\ref{sbwindow:positive-feedback},
including a certified broad tube through its displayed $T>t_g$, one has
\[
 t_g<T\le t_{\rm hold}.
\]
No separate pre-entry smooth-root or branch-speed deadline is required
for this assertion. It does not imply a generation or creation interval
lasting until order-one progress.
\end{proposition}
\begin{proof}
The early comparison supplies the individual errors and a persistent
positive lower face. The broad upper tube is strictly inside the retuned
upper faces. The mean envelope is below $0.0232$ after entry and stays
there. Longitudinal positivity, purity and the $q$ bound are among the
certified early premises, and $M\le\rho<1/2$ throughout. Hence none
of the failures defining $t_{\rm hold}$ occurs before $T$.
No root, weak rotation or weak-cohort gate hypothesis enters this argument.
\end{proof}

For generation, first choose a reached starting time $t_a$ at which the
premises of Proposition~\ref{sbgen:first-hit} are true. Its actual
continuation exit is
\begin{equation}\label{sbguard:generation-exit}
 t_{\rm gen}:=\inf\{t\ge t_a:\ \mathcal G_{\rm gen}(t)
                         \text{ fails before the first corridor hit}\}.
\end{equation}
Here $\mathcal G_{\rm gen}$ consists of precisely: the strong holding
and longitudinal conditions used for $\kappa$; a proved progress-control
predicate keeping $m<m_{\max}$ before the first hit; $\chi\ge-\epsilon_g\kappa$;
and the reserved creation charges/endpoint/staircase inequalities needed at
the eventual hit.  Progress monotonicity is not a primitive guard: it may be
proved as a consequence of the coupled progress comparison, but the preferred
interface below allows creation and progress to move together without
assuming $m'\ge0$.
For a proof that imposes these latter charges throughout the interval,
include those imposed inequalities in $\mathcal G_{\rm gen}$; if they
are instead proved only at the hit, they are separate hitting-time
obligations, not premature exit conditions. The fixed data events and
the deterministic corridor reserve are checked once. Include $t_\chi$
only when the stronger coherence route is actually invoked.

This definition does not assume the creation middle sign before its
first hit. If the generation entry predicates are false at a proposed
$t_a$, the generation theorem cannot start there. In particular,
$a=t_{\rm hand}$ is not automatically a generation entry time.
The required later ordering is the regulator clock before
$t_{\rm gen}$, as in \eqref{sbgen:first-hit-clock}. The available result
$t_g<t_{\rm hold}$ does not prove this later ordering, reachability of
$t_a$, or even persistence of holding from $T$ to that later interval.

\subsection{Late continuation after the directional reduction}

The old eight-face narrow-box continuation is no longer the load-bearing B2
architecture.  It remains useful as an early diagnostic, but the exact
factorizations in Lemma~\ref{sbemp:directional-complement} and
Lemma~\ref{lem:late-directional-ratio} show that its absolute four-entry error
budget is the wrong late object.

For reference, the former sufficient upper-face condition was
\begin{equation}\label{sbguard:missing-face}
 \sup_{i\in B:\,s_i=U}r_{u,i}<q_-d_{u,+}.
\end{equation}
In the original narrow box the right side was approximately $0.0340027$.
Blindly extending the whole-output norm estimate to order-one progress gave an
allowance exceeding $50.78$.  The ratio $50.78/0.0340027\simeq1493$ is close
to $e^{-2}=1600$ at $e=1/40$: the relaxation had destroyed two directional
$e$-cancellations.  Equation~\eqref{sbguard:missing-face} is therefore retained
only as a historical diagnostic and must not be imposed as a primitive
trajectory assumption.

\paragraph{Physical positive-face reserve.}
Let $u>0$ denote the physical positive input-ratio face and put $S=u/e$.
Using the reciprocal T1.11 field and the signed first-quadrant strong-shape
contribution, the high-SNR positive-face reserve has leading form
\[
\boxed{
 \mathfrak R_{\rm hold}(u,m)
 =\frac{\mu_2^2u}{2}\{\varrho^2(1-m)-1\}+o(u).
}
\tag{late-physical-reserve}\label{sbguard:late-physical-reserve}
\]
Since $\ell_2=\mu_2^2/2$ and
$\varrho^2=\ell_1/\ell_2$,
\[
\boxed{
 \mathfrak R_{\rm hold}(u,m)
 =u\{\ell_1(1-m)-\ell_2\}+o(u).
}
\tag{late-shared-gap}\label{sbguard:late-shared-gap}
\]
Thus the physical holding reserve and the mixed-sign relative-growth margin
are the \emph{same} concept-rate gap.  Their common population ceiling is
$m<1-\ell_2/\ell_1=1-\varrho^{-2}=5/9$ for $\varrho=3/2$; the repeated
appearance of $5/9$ is one rate ratio in several formulas, not independent
corroboration.

The useful asymptotic hierarchy is
\[
\boxed{e\ll u\ll1,}
\tag{late-scale-hierarchy}\label{sbguard:late-scale-hierarchy}
\]
with $u$ constrained by $e$-free geometric conditions: cluster--1 purity,
the regulator overlap lower bound, the regulator tangent-sign region,
longitudinal balance, and the signed strong-shape creation budget.  At fixed
physical $u$, the holding reserve is $O(u)$ whereas the high-SNR creation
margin is $O(e)$.  Small $e$ therefore helps holding through the scale
separation $e/u\to0$; it does not magically beat deterministic creation
charges that are themselves $O(e)$.

\paragraph{Upper and lower faces are asymmetric.}
For the positive input face use
\eqref{eq:reciprocal-upper-face}.  Positive first-quadrant contributors are
kept with their favorable sign in the reciprocal T1.11 residual rather than
charged by an absolute complement norm.  For the negative input face retain
the one-sided B2 inequality
\eqref{eq:late-hold-minus}; there is no positive-$\tau$ reciprocal gate
identity there.  The intended B2 region is therefore a hybrid rectangle with
a noise-scale lower input/output floor and a wider positive physical input
face.

\paragraph{Quadrant bookkeeping and the radial blind spot.}
Adding
\[
 P_3:=\sum_E(-v_{i1})_+(-v_{i2})_+
\]
completes the four sign-product quadrants algebraically, but it is not a
load-bearing reciprocal-holding state.  For the matched neuron
$v=z=-e_1$, one has $Q_E=N_2=N_4=P_3=D_E=0$ while its reciprocal positive-face
charge is
\[
 \frac{e\mu_2^2}{2}\{\varphi(0)-\varphi(S)\}>0.
\]
The correct radial completion is therefore the already-defined continuous
negative-coordinate energy
\[
 Y_k^-:=\left(\sum_i(-v_{ik})_+^2\right)^{1/2}.
\]
The late proof must propagate the harmful wrong-sign coordinate mass and the
normalized mismatch, rather than pretend that quadrant products alone cover
the anti-axis boundary.

\paragraph{Progress bridge.}
Let
\[
 c_i:=\frac{\langle p,(v_i^\top x)_+\rangle_n}{D},\qquad
 \kappa_p:=\frac{\langle p,|x_2|\rangle_n}{D}.
\]
ReLU Lipschitzness gives
\[
\boxed{|c_i-(v_{i1})_+|\le\kappa_p|v_{i2}|.}
\tag{PM-neuron}\label{sbguard:PM-neuron}
\]
Consequently, for
$p_1:=\sum_i z_{i1}(v_{i1})_+$ and
$\mathcal A:=\sum_i|v_{i1}v_{i2}|$,
\[
\boxed{
 |m-p_1|
 \le\kappa_p\{\mathcal A+\sqrt R\,D_E\}.
}
\tag{PM}\label{sbguard:PM}
\]
On SIM, $\kappa_p=O(e)$ and $\mathcal A\le R/2$, so bounded projected mass
and a late $D_E=O(e)$ invariant imply $p_1=m+O(e)$ without requiring
$\mathcal A$ itself to be small.

\begin{lemma}[Coupled regulator-clock progress envelope]
\label{sbguard:progress-envelope}
Continue with the exact regulator coordinates
$y=(m,b,y_\perp)$ and $d=y_*-y$.  Define
\[
 \chi_m:=K_{12}(b_*-b)-\sum_{j\ge3}K_{1j}y_j+r_1.
\]
Then the first coordinate of the exact vector regulator gives
\[
\boxed{\dot m=K_{11}(1-m)+\chi_m.}
\tag{progress-exact}\label{sbguard:progress-exact}
\]
On an interval where $\kappa=K_{22}>0$, introduce regulator time
$d\Lambda=\kappa\,dt$.  If
\[
 R_-\le\frac{K_{11}}\kappa\le R_+,
 \qquad
 -\eta_-\kappa\le\chi_m\le\eta_+\kappa,
\]
then
\[
 R_-(1-m)-\eta_-
 \le\frac{dm}{d\Lambda}
 \le R_+(1-m)+\eta_+.
\]
Consequently
\[
\boxed{
 \underline m(\Lambda)
 =1-\frac{\eta_-}{R_-}
 -\left(1-\frac{\eta_-}{R_-}-m_a\right)e^{-R_-\Lambda}
 \le m(\Lambda),
}
\tag{progress-lower}\label{sbguard:progress-lower}
\]
and
\[
\boxed{
 m(\Lambda)\le
 \overline m(\Lambda)
 =1+\frac{\eta_+}{R_+}
 -\left(1+\frac{\eta_+}{R_+}-m_a\right)e^{-R_+\Lambda}.
}
\tag{progress-upper}\label{sbguard:progress-upper}
\]
Let $B_M(m)=B_0-\tau_Mm$, let
$E_a=[d_2(t_a)-\epsilon_g]_+$, and put
$A_{\rm cr}=B_0-b_*+\epsilon_g$.  The scalar regulator comparison gives
$d_2(\Lambda)\le\epsilon_g+E_ae^{-\Lambda}$.  Hence creation must have
occurred by the first $\Lambda_*$ satisfying
\[
\boxed{
 \tau_M\underline m(\Lambda_*)
 =A_{\rm cr}+E_ae^{-\Lambda_*}.
}
\tag{PC-clock}\label{sbguard:PC-clock}
\]
If simultaneously
\[
\boxed{\overline m(\Lambda_*)<m_{\max},}
\tag{PC-ceiling}\label{sbguard:PC-ceiling}
\]
then the progress ceiling cannot be the first pre-creation exit.  No a priori
assumption $m'\ge0$ is used.  The cheaper face contradiction
$B_M(m_{\max})<b_{\rm floor}$ may replace
\eqref{sbguard:PC-clock}--\eqref{sbguard:PC-ceiling} whenever a proved late
lower bound $b\ge b_{\rm floor}$ is available.
\end{lemma}
\begin{proof}
The first display is the first coordinate of
\eqref{sbgen:vector-regulator}.  Divide by $\kappa$ and compare with the two
scalar linear ODEs.  For creation, pre-hit means
$d_2>b_*-B_M(m)=b_*-B_0+\tau_Mm$; combine this with the regulator upper
bound for $d_2$ and the lower progress envelope.
\end{proof}

\paragraph{Order audit for the corrected asymptotic program.}
The late hierarchy must distinguish terms that are genuinely beaten by
$e\downarrow0$ from terms that compete at the same normalized order.  Under
$e\ll u\ll1$, the intended orders are
\[
\begin{array}{c|c}
\text{quantity}&\text{target order}\\ \hline
\text{positive physical holding reserve}&O(u)\\
\text{population creation/progress margin}&O(e)\\
\text{signed strong-shape creation correction}&O(eu^2)\\
\text{joint gated-moment / DKW errors}&o(e)\\
\text{ambient/perpendicular forcing}&o(e)\\
\text{harmful wrong-sign complement}&o(e)\text{ or a controlled }O(e)\text{ constant}\\
\text{late normalized mismatch}&O(e)\\
\text{weak-feature timing charge}&\text{order still to be proved.}
\end{array}
\]
Thus shrinking $e$ does not resolve an $O(e)$ creation charge by itself; such
terms must satisfy the $e$-normalized constant ledger.  By contrast the
physical holding side has $O(u)$ room and becomes asymptotically nonbinding
when $e/u\to0$ once the wrong-sign and mismatch states have their stated
orders.  A constructive primitive region should therefore be nested: first
choose an interior $(m,z)$ with strict normalized creation/progress and rate
gaps, then choose a small physical cone $u$, then take $e\ll u$, and only
then choose initialization and sample size so the remaining terms are lower
order.

\paragraph{Current load-bearing late obligations.}
The first genuinely new deterministic lemmas are now:
\begin{enumerate}[label=\textnormal{(\roman*)}]
\item a gate-aware wrong-sign coordinate barrier that keeps the harmful parts
of $Y_1^-,Y_2^-$ at the noise/initialization scale through B2;
\item a stopped late mismatch estimate of the form $D_E/\sqrt R\le C_De$ (or
a sharper mass-relative variant), using the retained symmetric-residual
damping and the noise-scale skew source;
\item the asymmetric B2 face comparison using the reciprocal positive face and
\eqref{eq:late-hold-minus};
\item a perturbed creation/progress first-hit estimate in the common high-SNR
ledger, with the weak-feature charge, regulator complement and empirical
charges entered at their actual $e$-orders;
\item post-creation transport of a fixed retained weak cohort from the
cluster--2-pure quadrant-II region onto the attracting side of the created
gate-cliff branch, followed by the already-proved tracking and mass-growth
comparison.
\end{enumerate}
No full-network grading theorem and no optional $c_{12}^O$ coherence/spread
route are load-bearing for this default specialization program.  The complete
common primitive region remains open until these stopped estimates and their
time ordering are proved.

\subsection{Dimension dependence of the installed unprojected certificate}

Put $b_0=4.3975$, $\zeta_0=1.043$, $\beta=b_0-\zeta_0=3.3545$ and
$A=0.208$. With $a=a_0$, the old seed bound is
\[
 m_0=C_a\tau_d^2,\qquad
 C_a=\frac{0.9999c_{\rm seed}}{1+e^2S_b^2}e^{-2\lambda a}.
\]
On the logarithmic branch of the old feedback clock, its elapsed time is
\begin{equation}\label{sbguard:old-time}
 t_g-a=\frac{\log(8/(Ae m_0))-b_0a}{\beta}.
\end{equation}
The subtraction is $b_0a$ after changing from absolute to elapsed time;
it is not correct to omit the $a$ dependence when comparing invocations.
Since $\tau_d^2\sim2\log(2)/d$, fixing $a,e,c_{\rm seed},\lambda$
and the other constants gives
\[
 t_g-a=C+\beta^{-1}\log d+o(1),\qquad \beta^{-1}\simeq0.298107.
\]
The entry duration $4.1(2/\mu_2^2)$ has no explicit $d$; it is the duration
\emph{after a certified broad tube} and does not bound this feedback
clock or prove a dimension-uniform initialized theorem. T1.6's
constant-probability angular/radial selection and width guarantee remain
correct: they do not assert a dimension-uniform lower bound for each
projected radius or for its share of the full ambient energy.

Choose $T=t_g+\Delta$, with fixed $\Delta>0$. Substituting the old clock
in the feedback smallness condition gives exactly, on that branch,
\begin{equation}\label{sbguard:old-rho}
 \rho_{\max}=
 \frac{Ae}{44}\left(\frac{Ae}{8}\right)^{\zeta_0/\beta}
 m_0^{1+\zeta_0/\beta}
 \exp\left(\frac{\zeta_0b_0a}{\beta}-\zeta_0\Delta\right),
\end{equation}
provided this is below $10^{-4}$. Thus the old power is exactly
$p=b_0/\beta\simeq1.310926$ for fixed other constants; it is not merely
an empirical power-law fit. The same power occurs in $\nu_{\max}$.
Here $\rho$ is an output/energy cap enforced by initialization. It is
not an empirical concentration error and is not defined as $n^{-1/2}$.
The sample-size cost arises from $\nu$, separately. The existing
uniform operator-moment bound has, up to its clipping tail,
\[
 \nu\le C E_{\rm op}
     \sqrt{\frac{\log(Cn)+\log(C/\delta_{\rm gate})}{n}}
                 +\mathfrak t_{\rm clip}.
\]
Consequently this old sufficient route demands $n$ of order at least
$d^{2p}$ up to the envelope and logarithmic factors when these other
quantities are fixed. It does not establish an intrinsic sample lower
bound. The old mass cap also enforces
$h\varepsilon^2\le\rho e^{-2\lambda T}$, adding the exponent
$2\lambda/\beta$ to the power in its sufficient $\varepsilon^2$ cap.

\subsection{A projected-energy repair of the seed dilution}

There is a different valid mass envelope that changes this conclusion.
Define
\[
 R=\sum_i\|v_i\|^2,\qquad Z_o=\sum_i\|z_i\|^2,\qquad
 P=\tfrac12(R+Z_o),\qquad s_0=h\varepsilon^2.
\]
\begin{lemma}[Projected energy and its initialization scale]
\label{sbguard:projected-mass}
For every finite time,
\[
 R\le P,\quad \sum_i\|v_i\|\|z_i\|\le P,\quad
 P(t)\le P(0)e^{2\lambda t},
\]
where $P(0)=\varepsilon^2\sum_i R_i^2$ with the T1.6 initialized radii.
For $L=\log(1/\delta_P)$, with probability at least $1-\delta_P$,
\begin{equation}\label{sbguard:projected-init}
 P(0)\le\frac{\varepsilon^2}{d}
                 \{2h+\sqrt{8hL}+2L\}.
\end{equation}
In particular $h\ge8L$ implies $P(0)\le4h\varepsilon^2/d$.
\end{lemma}
\begin{proof}
Balance, conservation of input perpendicular coordinates, and decrease
of the aggregate output perpendicular norm give
$R-Z_o=\|W_\perp(t)\|_F^2-\|U_\perp(0)\|_F^2\le0$.
Cauchy--Schwarz gives $\sum r_i\|z_i\|\le\sqrt{RZ_o}\le P$.
Differentiating the projected energies gives exactly
\[
 P'=2\langle x,f_\parallel\rangle_n-2\|f_\parallel\|_n^2
                              -\|f_\perp\|_n^2\le2\lambda P.
\]
Here $\sum_i v_i\cdot P_i=\|f_\perp\|_n^2$, by the selector identity.
For the initialization estimate, the beta distribution of $R_i^2$
gives $\mathbb E(dR_i^2)^k\le2^k k!$ and mean $2$ for all $d\ge2$.
Thus for $0<t<1/2$,
\[
 \log\mathbb E e^{t(dR_i^2-2)}
 \le-2t-\log(1-2t)\le\frac{2t^2}{1-2t}.
\]
Chernoff's bound for the sum gives
$\sum_i dR_i^2\le2h+\sqrt{8hL}+2L$ with the claimed probability.
For $d=2$ the variable is identically $2$, so the bound still holds.
\end{proof}

Set $p_0=4h\varepsilon^2/d$ and $p(t)=p_0e^{2\lambda t}$ on this
initialization event. In the whole-network comparisons in
Section~\ref{sbwindow:scope}, replace the geometric envelope $s(t)$
by $p(t)$, but keep the ambient source parameter $s_0=h\varepsilon^2$.
This distinction is essential: the output perpendicular energy does not
start at scale $h\varepsilon^2/d$. The proof of
\eqref{sbwindow:H}--\eqref{sbwindow:Z} uses only $R\le s$,
$\sum r_i\|z_i\|\le s$, and
$\mathcal P\le\lambda s_0\sqrt{s}$; all three replacements are valid.
The exponent in the explicit envelope uses $p(t)-p_0$, its homogeneous
initial term uses $H(0)\le p_0/2$, and its ambient coefficient remains
$(1+\omega_0)\lambda s_0$.

Require
\begin{equation}\label{sbguard:projected-smallness}
 p_0e^{2\lambda T}\le\rho,\qquad
 p_0e^{\Theta T}\le e\rho,\qquad s_0\le e\rho.
\end{equation}
Then the entry charge constants and the negative-coordinate proof remain
valid with this projected envelope. In particular $s_0\le\rho$ still
justifies the $1.78\rho$ mismatch and $19\rho$ forcing bounds.
Retaining the already proved seed radial growth gives
\begin{equation}\label{sbguard:projected-seed}
 \frac{M(a)}{p(a)}\ge\widetilde m_0:={0.9999c_{\rm seed}\,d\tau_d^2
                      \over4(1+e^2S_b^2)}e^{-2\lambda a+2G_B(a)}.
\end{equation}
Here one may use the already proved radial amplification
$G_B(a)=m_{\rm rad}t_0+m_{\rm rad}^{\rm post}(a-t_0)$ from T1.8--T1.9;
setting $G_B=0$ gives the earlier, weaker bound. Neither lower bound
requires a new trajectory assumption.
For every $d\ge2$, $d\tau_d^2\ge2\log2$. Indeed for $d>2$ use
$1-e^{-x}\ge x/(1+x)$ with $x=2\log2/(d-2)$; $d=2$ is immediate.
Consequently $\widetilde m_0$ has a strictly positive lower bound
independent of $d$ when the other parameters, including $a$, are fixed.

\begin{corollary}[Feedback without the unprojected seed dilution]
\label{sbguard:projected-feedback}
On the projected initialization event and under the same reached early
tube/data guards, Proposition~\ref{sbwindow:positive-feedback} remains
valid with $m_0$ replaced by $\widetilde m_0$, $s$ by $p$, and the
initialization conditions replaced by \eqref{sbguard:projected-smallness}.
Its relative mass rate and constants $b_0,\zeta_0,0.052$ are unchanged.
Thus the explicit $\log d$ delay and $d^{-p}$ feedback caps are removed
from this comparison. This does not prove that the earlier reached time
$a$, T1.9 admissibility or the later creation clock are dimension-uniform.
\end{corollary}
\begin{proof}
All target/residual operator bounds now use $P\le p$.
The negative-coordinate quantities divided by $\sqrt p$ have initial
values at most one, and their physical decay rate remains
$\lambda-\sigma_k^2$. The seed mass bound divides by $p(a)$ instead of
$s(a)$, and $p'/p=2\lambda$. The proofs of the signed complement bound,
longitudinal growth and the three quarter-budget allocations are
otherwise identical. Ambient sources still use $s_0$, as required above.
\end{proof}

\subsection{Continuing the positive box beyond the locking horizon}

The early ordering can be strengthened: T1.9 need only deliver entry.
Its joint retained-weak horizon is not needed throughout the subsequent
small-output feedback clock. Here is a direct replacement for that part
of the continuation, retaining the actual holding exit
\eqref{sbguard:hold-exit} with $M_+=1/2$.

\begin{proposition}[A direct small-output holding continuation]
\label{sbguard:direct-holding}
Suppose the preceding entry proof delivers $\mathcal B$ at
$t_{\rm hand}$, with $q_i\in(0.99,1.01)$ and
$v_{i1}\ge\varepsilon\tau_d/\sqrt{1+e^2S_b^2}$ there.
Require the projected output/initialization caps through a chosen
$T>t_g\ge t_{\rm hand}$, the safe directional ambient bounds
\eqref{sbwindow:ambient}, and the fixed empirical face margins
\[
 \min(d_{u,-},d_{u,+},d_{w,-},d_{w,+},
                  d_{s,-},d_{s,+},d_{t,-},d_{t,+})\ge0.034.
\]
On $s\in[L,U]=[0.135,0.337]$, also assume the fixed-data bounds
\[
 h(s)\le0.64,\quad 0\le K(s)\le0.61,\quad
 0\le Z(s)\le2,\quad |b+g(s)|\le1,
\]
with the other data bounds unchanged. Then through $T$ the positive
box, purity, longitudinal positivity and $0.99<q_i<1.01$ persist.
The individual errors are less than $0.0026$ and the two mean errors
are less than $0.024$. In particular $t_g<T\le t_{\rm hold}$,
without invoking T1.9 beyond entry. The seed mass rate remains at least
$g_B=8.957$, so the projected signed-feedback certificate applies.
\end{proposition}
\begin{proof}
Stop at the first failure of the displayed box, ratio, positivity or
ambient guards. The projected global bound and the same off-diagonal
estimate as before give $B_\times^{\rm out}/e<0.00237$.
For $U=0.337$, $W=0.347$, both $q$ and $q^{-1}$ are below $1.011$.
The individual errors are bounded by
\[
 \begin{split}
 &\frac{1.011}{2}\{(1+e^2UW)0.00237+5\rho(U+W)\}\\
 &\quad+1.011e^2(2W+2U+UW)
        +1.011\rho\max\{0.01+0.61+2.26U,\,2(2.26)W\}
        +10^{-4}<0.0026.
 \end{split}
\]
This uses $\rho\le10^{-4}$ and $e\le0.025$.
The exact logarithmic-weight formula gives
\[
 \operatorname{osc} k<0.0051.
\]
For clarity the five contributions are bounded respectively by
$2.26(0.01)^2/0.99$, $0.002501$, $0.001001$, $0.000866$, and
$0.0002$: these are the constant target diagonal, noise target diagonal,
learned diagonal, off-diagonal and ambient contributions. Their sum is
less than $0.0051$. Moreover $|F_i|<2.01$ and $|G_i|<1.41$.
The exact mean formulas consequently give
\[
 |r_s|\le0.01(2.01)+0.0026+
              \tfrac14(U-L)(0.0051)<0.024,
\]
\[
 |r_t|\le\tfrac{0.01}{0.99}(1.41)+0.0026+
              \tfrac14(W-l)(0.0051)<0.018.
\]
Since $0.99(0.034)>0.024$ and $0.034/1.01>0.024$,
all eight faces are strictly inward.

The longitudinal ratio has the exact physical-time equation
\[
 \dot q_i=A_{11,i}(1-q_i^2)+e s_iA_{12,i}
                    -e t_i q_i^2 A_{21,i}+q_iP_{i1}/v_{i1}.
\]
Purity gives $A_{11,i}\ge4.48-5\rho\ge4.4795$ and
$|A_{12,i}|,|A_{21,i}|\le2.1e+5\rho\le0.053$.
The absolute charge of the last three terms is at most
\[
 0.025(0.337+0.347\cdot1.01^2)(0.053)+1.01(0.0002)<0.0012.
\]
On the lower face the restoring term is greater than $0.0891$.
On the upper face it is less than $-0.0900$, using
$4.4795(1-0.99^2)>0.0891$ and $4.4795(1.01^2-1)>0.0900$. Thus both ratio faces are strictly inward.
Also
\[
 \dot v_{i1}/v_{i1}
 \ge0.99\{4.4795-0.025(0.347)(0.053)\}-0.0002>4.43.
\]
The input longitudinal coordinate therefore cannot vanish; positivity
of $z_{i1}=q_iv_{i1}$ follows. Cluster-1 purity follows from the fixed
coordinate-envelope cone condition at $U$, which must be included in
the data guards. The lower bound for $v_{i1}$ persists, so the aggregate
ambient estimate closes the ambient guard as well, if
\eqref{sbguard:ambient-primitive} is imposed with $S_b$ replaced by $U$.
Finally the longitudinal product growth lower bound is
\[
 4(2.24)-2(5)10^{-4}
 -2(0.025)(0.347)(0.053)-0.0002
 =8.95788045>8.957.
\]
All guarded quantities thus continue through $T$ by the simultaneous
first-exit argument. This proof does not use a weak retention interval.
\end{proof}

The positive-box entry time still has to be delivered by a legitimate
early invocation. This proposition removes the need to extend that
invocation's weak-retention budget to $t_g$.  The later order-one continuation
is no longer formulated through the narrow-box diagnostic
\eqref{sbguard:missing-face}; it uses the asymmetric physical B2 interface,
the reciprocal positive-face identity, and the late wrong-sign/mismatch
estimates described above.  The present output-mass bound remains only an
early sufficient hypothesis.

\subsection{Primitive ledger and the remaining hard gate}

For precision, a safe directional ambient condition can be obtained
without assuming an individual output perpendicular norm is at most
$\varepsilon$. The aggregate estimate gives
$\|P_i\|\le\lambda s_0\sqrt p$ for every $i$. With the T1.9 radial
lower bound $r_i\ge\varepsilon\tau_d$, a sufficient condition for
\eqref{sbwindow:ambient} is
\begin{equation}\label{sbguard:ambient-primitive}
 s_0e^{\lambda T}\le
 \frac{10^{-4}e\mu_2^2\tau_d\sqrt d}
             {4\lambda\sqrt h(1+e^2S_b^2)}.
\end{equation}
This replaces any unjustified per-neuron perpendicular bound by an
aggregate one. In the original unprojected calculation the analogous
safe right side is
$10^{-4}e\mu_2^2\tau_d/[2\lambda\sqrt h(1+e^2S_b^2)]$.
The feedback proposition's explicit ambient hypotheses remain unchanged.

The complete necessary audit for the current proof route consists of:
\begin{enumerate}
\item The master coordinate/slab and gated-moment events, the prescribed
confidence split, T1.6 seed width, and $h\ge8\log(1/\delta_P)$.
For example split the former initialization budget in half: use
$\delta_{\rm init}/2$ for seed coverage and
$\delta_P=\delta_{\rm init}/2$ for projected concentration. Replace
the seed-width logarithm by $\log(4/\delta_{\rm init})$; the total
initialization failure budget is unchanged. Also keep the
core SNR upper bound $n\le(\delta/96)\exp(\mu_{\min}^2/(8\sigma_{\max}^2))$
from Definition~\ref{def:R2c-core}; increasing $n$ is not unrestricted.
\item T1.8 early admissibility evaluated at the same $H_\star$, then
all T1.9 inequalities: $\chi(H_\star)<1$, radial positivity,
$Y_0+2P_{\rm post}^\star<Y_g$, the coarse angular ceiling and cone
inward margin, $M_{\rm lock}>0$, $n>n_{\rm lock}$,
$D_\infty^{\rm post}<D_{\rm lock}$, and
$\tau_{\rm gap}+\tau_{\rm lock}\le\Delta_{\rm post}$.
For its joint retained-weak conclusion also keep
$\Delta_{\rm post}Q_{W,\pm}^{\rm post}<b_{W,\pm}^{\rm rem}$;
these cannot be dropped merely because feedback uses no W.
\item A reached broad-tube time $a$ with
$\Theta_a\le\arctan(eS_b)$,
$\Theta_a+D_{\rm lock}\le\arctan(eT_b)$, and radial discrepancy
$\mathcal Y_{\rm post}(a-t_0)\le10^{-4}$ through entry $t_{\rm hand}\le H_\star$. For continuation to a later $T$, either extend this invocation or use Proposition~\ref{sbguard:direct-holding}.
The bias floor must lie below these target angles. This is the actual
entry premise, not the definition $a=t_0+\tau_{\rm gap}$ alone.
\item The fixed data packet \eqref{sbwindow:data-size}, the empirical
kernel guards \eqref{sbwindow:entry-data}, the projected initialization
caps \eqref{sbguard:projected-smallness}, and the safe ambient cap
\eqref{sbguard:ambient-primitive}.
\item The physical feedback time and $T$ ordering with
$\widetilde m_0$, its $\rho,\nu$ caps, and the finite-sample inequality
$C E_{\rm op}\sqrt{[\log(Cn)+\log(C/\delta_{\rm gate})]/n}
+\mathfrak t_{\rm clip}\le\nu_{\max}$.
Here $E_{\rm op}$ is the maximum coordinate row-sum envelope supplied
by the existing gated-moment lemma; no ambient $d$ occurs in that
planar data concentration bound.
\item Finally, the late wrong-sign coordinate barrier, the $D_E=O(e)$
mismatch invariant, asymmetric physical B2 holding, perturbed
creation/progress first hit, and weak capture/continuation estimates stated in
the corrected late architecture above.
\end{enumerate}
The first five groups are explicit tests in the manuscript's data and early
comparison constants; their simultaneous primitive feasibility has not been
proved by this list.  The sixth is now a small collection of signed stopped
estimates rather than the obsolete eight-face condition
\eqref{sbguard:missing-face}.  Calling the ledger a nonempty primitive region
before those estimates and their time ordering are proved would still be
circular.

A natural route to nonemptiness is the nested hierarchy
\[
 0<e\ll u\ll1,
\]
followed by initialization and empirical errors that are lower order than the
$O(e)$ creation ledger.  For example one may take an initialization amplitude
as a sufficiently high power of $e$ and a polynomial sample size
$n=e^{-q}$ with $q>2$, provided the same choices also obey the existing
exponential-in-$1/e^2$ coordinate-SNR upper bound and every early locking
constraint.  This is a proof strategy, not yet a certified primitive region;
the weak-feature order, late mismatch constant and capture-time ordering must
still be inserted before nonemptiness is claimed.

For the projected comparison alone, the initialization caps can be
written explicitly as
\[
 \varepsilon^2\le\min\left\{
 \frac{d\rho}{4h}e^{-2\lambda T},\quad
 \frac{de\rho}{4h}e^{-\Theta T},\quad
 \frac{e\rho}{h},\quad
 \frac{10^{-4}e\mu_2^2\tau_d\sqrt d}
 {4\lambda h^{3/2}(1+e^2S_b^2)}e^{-\lambda T}\right\}.
\]
In particular, the necessary two-sided sample audit is
\[
 n_{\min}(\nu_{\max},e,\delta,\text{moment budgets})
 \le n\le\frac{\delta}{96}
       \exp\!\left(\frac{\mu_{\min}^2}{8\sigma_{\max}^2}\right).
\]
For equal noise with $\mu_{\min}=\mu_2$, the exponent on the right is
$1/(8e^2)$. The left side includes both the clipped-moment concentration
floor and the tail allowance, as well as $n_{\rm lock}$ and the later
creation floors. In the old unprojected route its $d^{2p}$ growth would
eventually exceed this fixed ceiling at fixed noise. The projected
repair removes that particular dimension-driven conflict, but does not
prove the remaining two-sided test feasible for the chosen constants.
No numerical value for the universal concentration constant is invented.

Other T1.8--T1.9 constraints must still be intersected with these.
This explicit partial region separates initialization smallness from
sample concentration; it is not advertised as the final theorem region.
The old severe $d$ powers arose from comparing a projected seed with
unprojected mass and are removable in this module. Any remaining
$\log d$ dependence through an actually certified $a$ or another guard
must be reported separately, not presumed absent.

The remaining load-bearing route is now more specific.  First prove the
gate-aware wrong-sign coordinate barrier and the late $D_E=O(e)$ mismatch
invariant.  Consume those estimates in the asymmetric physical B2 holding
argument, using the reciprocal T1.11 positive face and the one-sided negative
face.  Then close the perturbed regulator/progress first-hit comparison in the
normalized $O(e)$ creation ledger, followed by the pure-Q2-to-gate-cliff weak
capture ordering and the already-proved tracking/mass-growth interfaces.
Only after these steps should the common primitive region and final time
ordering be certified.  Full-network grading and the optional $c_{12}^O$
coherence route are not default predecessors, and the general noisy reversal
sign is outside the load-bearing specialization theorem scope.

\section{Re-auditing the initialized noisy reversal goal}
\label{sbaudit:scope}
\textbf{Continuation update.} Section~\ref{sbsep:scope} now proves a restricted noisy learned-strong-cohort endpoint theorem. The following Stage 0 status is historical; the original two-family goal remains open.


This continuation supersedes the architectural assertion that a general noisy
reversal necessarily requires a second large dynamical theory. The intended
target remains initialized noisy many-neuron specialization followed by
reversal on a nonempty OOD sector. Neither that target nor its impossibility
has been proved. In particular, an unproved sign estimate is not a proof that
its eventual proof must be long. We first remove unnecessary uniformity and
grading requirements, and then prove the decreasing side directly from
isotropic initialization.

\subsection{Exact effective matrices and the missing signed quantity}

Use the projected equations in Section~\ref{sbsign:scope}, and put
$b_i=w_{i,\perp}$. For a fixed actual probe pattern $c_i\in\{0,1\}$,
define
\[
 M_c=\sum_i c_i z_iv_i^\top,\qquad
 N_c=\sum_i c_i b_iv_i^\top.
\]
Every neuron is included. On the corresponding cell, for $\xi\in\mathcal S$,
\begin{equation}\label{sbaudit:cell}
 f_\parallel(\xi)=M_c\xi,\qquad f_\perp(\xi)=N_c\xi.
\end{equation}
On an interval with this pattern fixed, the exact equations give
\begin{align}
 \dot M_c&=\sum_i c_i
 (A_iv_iv_i^\top+z_iz_i^\top A_i-z_iP_i^\top),
 \label{sbaudit:matrix-rate}\\
 \dot N_c&=\sum_i c_i
 \left\{-\mathbb E_n[f_\perp(x)(v_i^\top x)_+]v_i^\top
           +b_iz_i^\top A_i-b_iP_i^\top\right\}.
 \label{sbaudit:ambient-rate}
\end{align}
The training selectors in $A_i$ remain the actual selectors. Even if all
active neurons have one common training selector, the two Gram matrices
$\sum c_i v_iv_i^\top$ and $\sum c_i z_iz_i^\top$ enter the rate; fixing
the probe pattern alone does not supply them from $M_c$.

For a unit probe $\xi=\eta z e_p+sk e_q$, where
$k=\sqrt{1-\eta^2z^2}$ and $s\in\{-1,1\}$, write
\[
 \mathcal A_\xi=\eta^{-1}(f_p(\xi)-\eta z),\qquad
 \mathcal T_\xi=\eta^{-2}f_q(\xi).
\]
These definitions do not assume bounded scaled coefficients. Exactly,
\begin{align}
 E(\xi)&=\frac12+\eta^2\left\{
 \frac12\mathcal A_\xi^2-sk\mathcal T_\xi-\frac12z^2
 \right\}+\frac{\eta^4}{2}\mathcal T_\xi^2
                  +\frac12\|f_\perp(\xi)\|^2,
 \label{sbaudit:exact-error}\\
 \frac{\dot E(\xi)}{\eta^2}
 &=\mathcal A_\xi\dot{\mathcal A}_\xi
    +(\eta^2\mathcal T_\xi-sk)\dot{\mathcal T}_\xi
    +\eta^{-2}\langle f_\perp(\xi),\dot f_\perp(\xi)\rangle.
 \label{sbaudit:exact-rate}
\end{align}
For $M_c=\left(\begin{smallmatrix}a&\eta b\\
\eta c&\eta^2d\end{smallmatrix}\right)$, in the ordered basis $(e_p,e_q)$,
$\mathcal A_\xi=(a-1)z+skb$ and $\mathcal T_\xi=cz+skd$.
There is no derivative remainder in \eqref{sbaudit:exact-rate}.
With moving probe patterns, the same identities hold almost everywhere
with frozen-selector rates, as in Proposition~\ref{sbsign:switching}.

There is also a formulation that does not grade any matrix. Write
$r_x=x-f(x)$ and $r_\xi=\xi-f(\xi)$, including ambient output coordinates,
and let $c_i=\mathbf1_{u_i^\top\xi>0}$. Set
\begin{equation}\label{sbaudit:kernel}
 \mathsf K(\xi,x)=\sum_i\left[
 (u_i^\top\xi)_+(u_i^\top x)_+ I_d
 +c_i\alpha_i(x)(\xi^\top x)w_iw_i^\top\right].
\end{equation}
The exact full gradient-flow equations imply, almost everywhere,
\begin{equation}\label{sbaudit:alignment}
 \dot f(\xi)=\mathbb E_n[\mathsf K(\xi,x)r_x],\qquad
 \dot E(\xi)=-\mathfrak J_\xi,\qquad
 \mathfrak J_\xi:=\mathbb E_n\langle r_\xi,\mathsf K(\xi,x)r_x\rangle.
\end{equation}
Indeed, differentiate each $w_i(u_i^\top\xi)_+$ and substitute the two
parameter equations. The positive semidefinite training Gram kernel
controls training-loss dissipation; it does not give a sign for this
mixed probe--training pairing. This is an exact instantaneous Jacobian
kernel, not a frozen-kernel or lazy-training approximation.

To compare with regulator information, set $p=(x_1)_+$,
$D=\mathbb E_n p^2>0$, and define
\[
 d=\mathbb E_n[pr_x]/D,\quad r_x^\circ=r_x-pd,\quad
 q_\xi(x)=\mathsf K(\xi,x)^\top r_\xi,\quad
 Q_\xi=\mathbb E_n[pq_\xi(x)].
\]
Then the exact orthogonal decomposition is
\begin{equation}\label{sbaudit:regulator-split}
 \mathfrak J_\xi=Q_\xi^\top d+
 \mathbb E_n\left\langle q_\xi-pQ_\xi/D,r_x^\circ\right\rangle.
\end{equation}
Regulator progress controls coordinates of $d$. It does not itself give
a signed bound for $Q_\xi$ or the last, centered pairing. This is a precise
additional quantity, rather than a demand for global grading or small
unassigned mass. Reciprocal holding and the signed complement estimates
can be used to estimate this pairing; their existence does not already
establish its sign.

\subsection{An endpoint alternative to derivative propagation}

\begin{proposition}[Three snapshots give a whole sector with a rebound]
\label{sbaudit:secants}
Let a finite network follow an absolutely continuous parameter trajectory
on $[t_0,t_2]$. Suppose $t_0<t_1<t_2$ and a unit $\xi_0\in\mathcal S$
satisfies
\begin{equation}\label{sbaudit:two-gaps}
 E(\xi_0,t_0)-E(\xi_0,t_1)\ge a>0,\qquad
 E(\xi_0,t_2)-E(\xi_0,t_1)\ge b>0.
\end{equation}
Put $L=\max_{j=0,1,2}\sum_i\|w_i(t_j)\|\|u_i(t_j)\|$.
For every unit $\xi\in\mathcal S$ with
\[
 \|\xi-\xi_0\|<\frac{\min(a,b)}{4(1+L)^2},
\]
the two gaps are at least $a/2$ and $b/2$. Consequently the error has an
interior minimum on $[t_0,t_2]$, and its rebound from that minimum to $t_2$
is at least $b/2$. The derivative is negative on a set of positive measure
in $(t_0,t_1)$ and positive on a set of positive measure in $(t_1,t_2)$.
No activation pattern need persist between snapshots.
\end{proposition}
\begin{proof}
At a snapshot the map $f$ is $L$-Lipschitz and positively homogeneous.
On unit directions its residual norm and residual Lipschitz constant are
both at most $1+L$. Thus
$|E(\xi,t_j)-E(\xi_0,t_j)|\le(1+L)^2\|\xi-\xi_0\|$.
Apply this twice to each gap. Continuity gives an attained minimum,
and both endpoints have strictly larger error than $t_1$. Absolute
continuity and integration of the derivative give the two sign sets.
The minimum need not be unique or strict; no such claim follows from
the endpoint inequalities alone.
\end{proof}

At each snapshot form the actual matrices for $\xi_0$ and put
\[
 H_j=(M_j-I)^\top(M_j-I)+N_j^\top N_j.
\]
The two inequalities to prove in this alternative are exactly
\begin{equation}\label{sbaudit:matrix-secants}
 \tfrac12\xi_0^\top(H_0-H_1)\xi_0\ge a,
 \qquad
 \tfrac12\xi_0^\top(H_2-H_1)\xi_0\ge b.
\end{equation}
Each $M_j,N_j$ includes every neuron. A fixed three-snapshot cell permits
using the same quadratic forms for nearby probes; the Lipschitz argument
above does not even need that restriction.

For the stronger claim of strict monotonicity on two open time intervals,
one may instead prove strict signs of \eqref{sbaudit:alignment} at two
regular times, with negative $\dot E$ first. Here regular means every
training preactivation and the chosen probe's preactivations are nonzero.
The parameter vector field and the probe derivative are then continuous
locally. Strict signs persist on two time intervals and one common open
probe sector. This local argument requires no uniform gate persistence
between those intervals. Regularity at those later witnesses must be
checked; the endpoint proposition alone does not assert it.

\subsection{The decreasing side is available from isotropic initialization}

\begin{theorem}[High-probability initial descent for a fixed OOD probe]
\label{sbaudit:initial-descent}
Fix nonzero empirical training samples in $\mathcal S$, a unit probe
$\xi\in\mathcal S$, and $d\ge2$. Initialize independently
$u_i=w_i=\varepsilon\widehat u_i$ with $\widehat u_i$ uniform on
$\mathbb S^{d-1}$. Put $s_0=h\varepsilon^2$,
$C_x=\mathbb E_n\|x\|^2$, and
\[
 F(c)=\sqrt{1-c^2}+(\pi-\arccos c)c,\qquad
 J_\xi=\frac1\pi\mathbb E_n\left[
 (\xi^\top x)\|x\| F\left(\frac{\xi^\top x}{\|x\|}\right)\right].
\]
For $\ell=\log(1/\delta)$ define
\[
 b_h=C_x\sqrt{\frac{6\ell}{d(d+2)h}}+
                 \frac{4C_x\ell}{3h}.
\]
With probability at least $1-\delta$ over initialization,
\begin{equation}\label{sbaudit:initial-rate}
 \dot E(\xi,0)\le-\frac{s_0}{d}
 \{J_\xi-2db_h-2dC_xs_0(2+s_0)\}.
\end{equation}
In particular, if $J_\xi\ge J_*>0$ and
\begin{equation}\label{sbaudit:initial-budget}
 2db_h+2dC_xs_0(2+s_0)\le J_*/2,
\end{equation}
then $\dot E(\xi,0)\le-s_0J_*/(2d)<0$. Almost surely all the finitely
many initial training and probe margins are nonzero, so the strict
decrease persists on some positive time interval and a nonempty open
sector containing $\xi$.
\end{theorem}
\begin{proof}
At initialization, $\|f_0(x)\|\le s_0\|x\|$ and
$\|\mathsf K_0(\xi,x)\|\le2s_0\|x\|$. Replacing both residuals by
their targets in \eqref{sbaudit:alignment} therefore costs at most
\[
 2C_xs_0^2(2+s_0).
\]
Because $w_i=u_i$, the two kernel contributions to the target pairing
are identical. With
\[
 Y_i=\mathbb E_n[(\xi^\top x)(\widehat u_i^\top\xi)_+
                                      (\widehat u_i^\top x)_+],
\]
the target pairing is $2\varepsilon^2\sum_iY_i$.
Rotational integration gives
\[
 \mathbb E[(\widehat u^\top\xi)_+(\widehat u^\top x)_+]
 =\frac{\|x\|}{2\pi d}F(\xi^\top x/\|x\|),
 \qquad \mathbb EY_i=J_\xi/(2d).
\]
For completeness this integral follows by writing a standard Gaussian
as its independent radius and uniform direction, using
$\mathbb E\|g\|^2=d$, and integrating the product of two positive
linear forms over their angular overlap in the two-dimensional plane.
The uniform angular integral is
$\{\sin\theta+(\pi-\theta)\cos\theta\}/(4\pi)$; the independent
two-dimensional Gaussian squared radius has mean $2$.

Also $|Y_i|\le C_x$. Apply Jensen with weights proportional to
$\|x\|^2$, and use
\[
 \mathbb E[(\widehat u^\top a)^2(\widehat u^\top b)^2]
 =\frac{1+2(a^\top b)^2}{d(d+2)}\le\frac3{d(d+2)}
\]
for unit $a,b$. It follows that
$\mathbb EY_i^2\le3C_x^2/[d(d+2)]$.
The centered summands have absolute bound $2C_x$, so the one-sided
Bernstein inequality gives $h^{-1}\sum_iY_i\ge\mathbb EY_i-b_h$
with the stated probability. Substitution proves
\eqref{sbaudit:initial-rate}. On the probability-one nonzero-margin
event the finite empirical vector field is smooth near initialization;
continuity gives the last assertion.
\end{proof}

\begin{corollary}[An explicit noisy SIM region for the decreasing side]
\label{sbaudit:noisy-descent}
Consider equal cluster weights and write
$x^{(k)}=\mu_ke_k+\zeta^{(k)}$, with $\mu_1>\mu_2>0$.
Fix $0<r\le1/4$ and the OOD probe
$\xi_r=re_1-\sqrt{1-r^2}e_2$. Suppose each noise vector has norm at most
$a$, and set
\[
 \overline C=(\mu_1+a)^2,\qquad
 J_*:=\frac{11\mu_1^2r}{32\pi}-2(2\mu_1+a)a>0.
\]
Sufficient primitive initialization conditions are
\begin{equation}\label{sbaudit:primitive-descent}
 h\ge\max\left\{
 \frac{1536\overline C^2\ell}{J_*^2},
 \frac{64d\overline C\ell}{3J_*}\right\},\qquad
 h\varepsilon^2\le\min\left\{1,\frac{J_*}{24d\overline C}\right\}.
\end{equation}
Under these conditions Theorem~\ref{sbaudit:initial-descent} applies.
For planar isotropic Gaussian cluster noise with standard deviations
at most $\sigma_{\max}$ and $N$ samples in total, one may take
$a=\sigma_{\max}\sqrt{2\log(N/\delta_D)}$ with probability at least
$1-\delta_D$. Independently sampling the initialization gives total
success probability at least $1-\delta_D-\delta$.
This is an explicit nonempty region for initial descent, not yet a common
parameter region for specialization and reversal.
\end{corollary}
\begin{proof}
Write $k=\sqrt{1-r^2}$. The noiseless quantity is
\[
 J^{(0)}_{\xi_r}=\frac1{2\pi}
 \{\mu_1^2r[k+(\pi-\arccos r)r]
       -\mu_2^2k[r-k\arcsin r]\}.
\]
Here $k\ge3/4$ and $r-k\arcsin r\le r(1-k)\le r^3$.
Since $\mu_2<\mu_1$ and $r^2\le1/16$, this gives
$J^{(0)}_{\xi_r}\ge11\mu_1^2r/(32\pi)$.

To bound the noise change without differentiating the angular formula,
use the expectation representation in the preceding proof. Cauchy--Schwarz
and $\mathbb E(\widehat u^\top v)^2=\|v\|^2/d$ give, for any $x,y$,
\[
 2d\left|\mathbb E_{\widehat u}
 \bigl[(\xi^\top x)(\widehat u^\top\xi)_+(\widehat u^\top x)_+
       -(\xi^\top y)(\widehat u^\top\xi)_+(\widehat u^\top y)_+\bigr]
 \right|
 \le2(\|x\|+\|y\|)\|x-y\|.
\]
Average this inequality with $y=\mu_ke_k$ to get $J_{\xi_r}\ge J_*$.
The two width bounds make the two terms of $2db_h$ at most $J_*/8$
each. The initialization cap makes the remaining term in
\eqref{sbaudit:initial-budget} at most $J_*/4$. The Gaussian assertion
is the two-dimensional radial tail followed by a union bound. Finally,
choose sufficiently small positive noise, then finite $h$, then positive
$\varepsilon$ sufficiently small to see nonemptiness for this corollary.
\end{proof}

\begin{corollary}[A directional noise estimate reaches noise-scale probe arcs]
\label{sbaudit:sharp-noise}
Let $0<r_0<r_1\le1/4$, suppose all sample perturbations have norm at most
$a\le\mu_2/4$, and suppose the empirical mean perturbation of cluster 1
has norm at most $b$. Uniformly for $r_0\le r\le r_1$,
\begin{equation}\label{sbaudit:sharp-J}
 J_{\xi_r}\ge J_{\rm sh}:=\frac1{2\pi}
 \left\{\mu_1(\mu_1-a)r_0-\mu_1b-a^2
              -\frac{(\mu_2r_1+a)^3}{\mu_2-a}\right\}.
\end{equation}
Whenever $J_{\rm sh}>0$, it may replace $J_*$ in the initialization
conditions, with $\overline C=(\mu_1+a)^2$.
For $n_1$ independent isotropic Gaussian samples in cluster 1, the mean
event can be imposed at failure probability $\delta_m$ with
\[
 b=\sigma_1\sqrt{2\log(1/\delta_m)/n_1}.
\]
\end{corollary}
\begin{proof}
The function $F$ in Theorem~\ref{sbaudit:initial-descent} is increasing,
with $F(0)=1$, so $cF(c)\ge c$ for every $-1\le c\le1$.
On cluster 1 this gives a lower bound by
$\mathbb E_{1,n}[(\xi_r^\top x)\|x\|]$ for the numerator of its
contribution to $J_{\xi_r}$. Since
$|\|x\|-\mu_1|\le a$ and $|\xi_r^\top x|\le\mu_1r+a$,
this is at least
$\mu_1^2r-\mu_1b-\mu_1ra-a^2$.

On cluster 2, $\xi_r^\top x<0$. Write its angle with $\xi_r$ as
$\pi-\beta$, with $0\le\beta<\pi/2$. Then
\[
 F(-\cos\beta)=\sin\beta-\beta\cos\beta
 \le\sin\beta(1-\cos\beta)\le\sin^3\beta.
\]
Moreover $|\det(\xi_r,x)|\le\mu_2r+a$ and $\|x\|\ge\mu_2-a$.
Its possibly negative numerator is therefore bounded below by
$-(\mu_2r+a)^3/(\mu_2-a)$. Combine the two clusters with weights
$1/2$, and use $r\ge r_0$ in the positive term and $r\le r_1$ in the
negative term. The mean-event assertion follows from the planar
Gaussian radial tail with variance $\sigma_1^2/n_1$.
\end{proof}

Unlike the absolute perturbation bound, \eqref{sbaudit:sharp-J} is useful
when $r_0,r_1$ themselves have the noise scale. For fixed means and fixed
$0<\kappa_0<\kappa_1$, take $r_j=\kappa_j e$,
$\sigma_{\max}=O(e)$, polynomial sample sizes growing to infinity,
and fixed positive failure probabilities. Then
$a=O(e\sqrt{\log(1/e)})$, $b=o(e)$, and
\[
 J_{\rm sh}=\frac{\mu_1^2\kappa_0}{2\pi}e+o(e)>0.
\]
Thus the initial-descent argument need not choose probes far outside the
noise-scale learned angles. This asymptotic compatibility concerns this
module only; it does not prove the remaining specialization inequalities.

\begin{corollary}[Uniform descent permits a later, trajectory-chosen probe]
\label{sbaudit:uniform-descent}
Fix $0<r_0<r_1\le1/4$ and the closed probe arc
$\mathcal I=\{re_1-\sqrt{1-r^2}e_2:r_0\le r\le r_1\}$.
Use a positive uniform lower bound $J_*$ from the absolute-noise corollary
(evaluated at $r_0$) or from \eqref{sbaudit:sharp-J}.
Let
\[
 \rho=\frac{J_*}{32d\overline C},\quad
 G=1+\left\lceil\frac{\arcsin r_1-\arcsin r_0}{\rho}\right\rceil,
 \qquad \ell=\log(G/\delta).
\]
Impose \eqref{sbaudit:primitive-descent} with this $\ell$.
With initialization probability at least $1-\delta$, every probe on
$\mathcal I$ has initial upper right error derivative at most
$-s_0J_*/(4d)$. There exists a common $\tau>0$ on which all these errors
are strictly decreasing. Thus a later probe chosen using the trajectory,
anywhere in this deterministic arc, is covered by the same event.
The positive time $\tau$ is not quantitatively lower bounded here.
\end{corollary}
\begin{proof}
Use $G$ equally spaced angular grid points. Every point on the arc is
within Euclidean distance $\rho$ of a grid point. The Bernstein argument
holds at every grid point simultaneously with probability at least
$1-\delta$. Directly from its definition,
\[
 |Y_i(\xi)-Y_i(\zeta)|\le2C_x\|\xi-\zeta\|
 \quad(\|\xi\|=\|\zeta\|=1).
\]
Consequently the transport cost in $2d h^{-1}\sum_iY_i$ is at most
$4dC_x\rho\le J_*/8$. Grid fluctuations cost at most $J_*/4$ and
the residual correction costs at most $J_*/4$. The remaining margin
is at least $3J_*/8$, stronger than asserted.

At an initial zero probe preactivation the target-pairing formula is
unchanged for every compatible selector: its potentially ambiguous
second kernel term has a factor $w_i^\top\xi=u_i^\top\xi=0$.
The residual correction is uniformly bounded for all those selectors.
Training margins are nonzero almost surely, so the parameter vector
field is continuous for a positive initial interval. The graph of
compatible probe selectors is closed. Compactness of the probe arc
and of the selector cube, together with the strict uniform initial
margin, therefore gives a common positive interval on which all
compatible error slopes remain negative. Integrate the almost-everywhere
slopes to obtain strict decrease. This step provides existence of an
interval, not an explicit initialized learning clock.
\end{proof}

\subsection{A local exit mechanism that avoids a global sign theorem}

\begin{proposition}[One gate exit supplies a later positive witness]
\label{sbaudit:exit-witness}
At a time $t_*$ suppose all training preactivations are nonzero. For a
unit probe $\xi_*$ suppose exactly one probe preactivation vanishes,
say $q_j=u_j^\top\xi_*=0$, all other probe preactivations are nonzero,
and $\dot q_j<0$. Remove neuron $j$ and denote the remaining probe
output and its time derivative at $t_*$ by $F$ and $\dot F$.
If
\begin{equation}\label{sbaudit:exit-charge}
 \underbrace{\langle F-\xi_*,\dot F\rangle}_{\text{signed background rate}}
 +\underbrace{\langle F-\xi_*,w_j\rangle\dot q_j}_{\text{exit contribution}}
 >0,
\end{equation}
then the actual full-network probe error has strictly positive derivative
on an interval immediately before $t_*$, and on a common nonempty open
probe sector after shortening that interval.
If $\xi_*$ lies in the interior of $\mathcal I$ and
Corollary~\ref{sbaudit:uniform-descent} holds, there is an initial decreasing
interval and a later increasing interval on a common open OOD sector.
\end{proposition}
\begin{proof}
Before the transverse exit, $q_j>0$ and
$f(\xi_*)=F(t)+w_j(t)q_j(t)$. The active-side derivative is
\[
 \langle F+w_jq_j-\xi_*,
       \dot F+\dot w_jq_j+w_j\dot q_j\rangle.
\]
Its limit at $t_*$ is exactly \eqref{sbaudit:exit-charge}. The regular
training flow and all other probe contributions are smooth locally.
Continuity gives a positive interval before the exit; choose a compact
subinterval strictly before it and then a small open probe sector.
Intersect with the initial descent sector. The strict positive and
negative intervals cannot overlap, so they have the required ordering.
\end{proof}

In the planar notation $v_j=r_j(\cos\phi_j,\sin\phi_j)$,
$z_j=s_j(\cos\psi_j,\sin\psi_j)$, take
$\xi_*=(\sin\phi_j,-\cos\phi_j)$ and suppose $\dot\phi_j>0$.
Then $\dot q_j=-r_j\dot\phi_j$ and
$z_j^\top\xi_*=s_j\sin(\phi_j-\psi_j)$.
If the background were zero, the exit contribution would therefore be
\begin{equation}\label{sbaudit:angular-exit}
 r_js_j\dot\phi_j\sin(\phi_j-\psi_j).
\end{equation}
In the actual network the background in \eqref{sbaudit:exit-charge}
is retained exactly, including every unassigned neuron. Directional
estimates for that signed expression may be cheaper than bounding the
whole complementary mass. Specialization alone supplies neither an
outward exit with output-angle lag nor this signed background bound;
the proposition is a finite-window route to test, not a reached-exit
claim.

\subsection{What the re-audit does and does not decide}

The initialized decreasing side now has a direct high-probability noisy
many-neuron proof. It does not rely on the two-neuron example. The uniform
version covers a prescribed nonempty arc, with a quantitative derivative
margin, but its common initial time interval is only asserted to be positive.
Its width requirement
includes a $d$-dependent term, and is not a dimension-free width theorem.

The first remaining reversal-specific task is to reach, at a probe
covered by that decreasing-side result, either
\begin{equation}\label{sbaudit:later-witness}
\textbf{Continuation update.} Section~\ref{sbsep:scope} now proves a restricted noisy learned-strong-cohort endpoint theorem. The following Stage 0 status is historical; the original two-family goal remains open.

 \mathfrak J_\xi(t_+)<0
\end{equation}
at a regular later time, or the second strict secant inequality in
\eqref{sbaudit:two-gaps}. For example, once $t_1>0$ is in the initial
decreasing interval, $E(\xi,t_2)\ge E(\xi,0)$ is a stronger sufficient
endpoint target, but is not necessary and should not be imposed if the
weaker second secant can be proved. A desired post-specialization
decreasing interval would require its own witness; initial descent is
not automatically descent after specialization.

None of the supplied specialization conclusions explicitly bounds the
signed mixed pairing in \eqref{sbaudit:regulator-split}, supplies this
later witness, or supplies the second secant. Nor do the early feedback
and creation guards automatically extend to a later witness time.
Conversely, no argument here excludes deriving that one-sided quantity
from a short finite-window use of the signed equations.

To distinguish non-implication from impossibility, in the noiseless
planar limit put balanced neurons exactly on the two positive axes.
For arbitrary multiplicities the exact flow gives
$f(x)=m_1(x_1)_+e_1+m_2(x_2)_+e_2$, with
$m_j'=2\lambda_jm_j(1-m_j)$ and $0<m_j(0)<1$.
For every fixed probe its error is nonincreasing. This verifies that
perfect axis concentration, fixed gates, increasing progress and balanced
gradient flow do not, by themselves, force reversal. It is not an
isotropic-initialization counterexample, not a noisy-model counterexample,
and not a counterexample to the full basin-creation hypotheses.

Accordingly, neither O1 nor O2 is closed by this continuation, and O3
has not been justified. The Scope A pivot is not accepted on the basis
of the old grading obstruction. The reduced O2 route remains live:
one probe, one finite window, a later signed witness or second secant,
and the sector extension above. Proving that witness from the same
initialized event and closing specialization's common primitive region
are still required before asserting the original headline theorem.

\section{Quantitative descent and specialization-linked endpoint certificates}
\label{sbend:scope}

The original noisy many-neuron specialization-to-reversal goal remains open.
This section makes the initialized decreasing interval quantitative, isolates
the signed endpoint change needed for a post-specialization rebound, and
audits the angular lag without assuming that it decays monotonically.
All conclusions called endpoint certificates below are conditional on their
displayed reached-state hypotheses. No Scope A pivot is made.

\subsection{A deterministic decreasing interval without a minimum gate margin}

\begin{lemma}[Gate-robust control of the initial target pairing]
\label{sbend:pairing-motion}
Let $s_0=h\varepsilon^2\le1$, $C_x\le\overline C$, and
$\|\mathbb E_nxx^\top\|\le\overline\lambda$. On any interval where
$S(t)\le2s_0$, define
\[
 B(t)=\int_0^t\overline\lambda(1+S(v))\sqrt{S(v)}\,dv.
\]
For a unit planar probe and the exact kernel of
\eqref{sbaudit:kernel}, put
$G_\xi(t)=\mathbb E_n[\xi^\top\mathsf K_t(\xi,x)x]$.
Uniformly in the probe and compatible selectors,
\begin{equation}\label{sbend:pairing-motion-bound}
 |G_\xi(t)-G_\xi(0)|
 \le(36+6\sqrt2)\overline C\,\overline\lambda s_0t
 <45\overline C\,\overline\lambda s_0t.
\end{equation}
Also
\begin{equation}\label{sbend:residual-cost}
 |\dot E(\xi,t)+G_\xi(t)|
 \le2\overline C S(t)^2(2+S(t))
\end{equation}
almost everywhere. These estimates do not require fixed training or probe
gates or a lower bound on their initial margins.
\end{lemma}
\begin{proof}
Use stacked matrices $U,W$ whose columns are the full neuron vectors.
Balance gives $\|U\|_F=\|W\|_F=\sqrt S$. The selected full residual
matrix for each neuron has norm at most
$\overline\lambda(1+S)$: its target part is bounded by
$\overline\lambda$, and Cauchy--Schwarz with
$\|f\|_n\le\sqrt{\overline\lambda}S$ bounds the learned part by
$\overline\lambda S$. Thus both parameter displacements have Frobenius
norm at most $B(t)$, and $\|W-U\|_F\le2B(t)$.

Set
\[
 I_\xi(U)=\mathbb E_n\left[(\xi^\top x)
             \sum_i(u_i^\top\xi)_+(u_i^\top x)_+\right].
\]
ReLU Lipschitzness and Cauchy--Schwarz give
\[
 |I_\xi(U(t))-I_\xi(U(0))|
 \le\overline C(\sqrt{S(t)}+\sqrt{s_0})B(t).
\]
In the second kernel term replace
$(w_i^\top\xi)(w_i^\top x)$ by
$(u_i^\top\xi)(u_i^\top x)$, retaining the actual two selectors.
Their product with the latter expression is exactly the ReLU product,
also at zero preactivations. Therefore
\[
 |G_\xi(t)-2I_\xi(U(t))|
 \le\overline C(\|W\|_F+\|U\|_F)\|W-U\|_F
 \le4\overline C\sqrt{S(t)}B(t).
\]
At initialization the difference is zero. Since
$B(t)\le3\overline\lambda\sqrt{2s_0}t$, combining the two estimates
proves \eqref{sbend:pairing-motion-bound}. Finally
$\|f(x)\|\le S\|x\|$ and
$\|\mathsf K_t(\xi,x)\|\le2S\|x\|$ bound the cost of replacing both
residuals in \eqref{sbaudit:alignment} by their targets, proving
\eqref{sbend:residual-cost}.
\end{proof}

\begin{theorem}[Explicit uniform initial descent time and snapshot drop]
\label{sbend:descent-clock}
On the initialization event proved in
Corollary~\ref{sbaudit:uniform-descent}, with its primitive width and
initialization conditions, define
\begin{equation}\label{sbend:tau-dec}
 \tau_{\rm dec}:=\min\left\{
 \frac{\log(5/4)}{6\overline\lambda},
 \frac{J_*}{720d\overline C\,\overline\lambda}\right\},
 \qquad c_{\rm dec}:=\frac{s_0J_*}{4d}.
\end{equation}
Then every probe on the deterministic arc $\mathcal I$ satisfies
\[
 \dot E(\xi,t)\le-c_{\rm dec}
 \quad\text{for almost every }0\le t\le\tau_{\rm dec}.
\]
In particular, the deterministic time and uniform drop
\begin{equation}\label{sbend:early-drop}
 t_{\rm early}:=\tau_{\rm dec}/2,\qquad
 a_{\rm dec}:=\frac{s_0J_*\tau_{\rm dec}}{8d}>0
\end{equation}
obey $E(\xi,0)-E(\xi,t_{\rm early})\ge a_{\rm dec}$.
The arc's angular width is exactly
\[
 \Delta_{\rm init}=\arcsin r_1-\arcsin r_0.
\]
For $r_j=\kappa_je$, it lies between
$(\kappa_1-\kappa_0)e$ and
$(\kappa_1-\kappa_0)e/\sqrt{1-r_1^2}$.
\end{theorem}
\begin{proof}
Before the residual correction, the uniform-grid event gives
$G_\xi(0)\ge5s_0J_*/(8d)$ for every probe. The grid fluctuation costs
$J_*/4$ and the transport costs $J_*/8$ from a target lower bound $J_*$.
The installed exponential estimate in Lemma~\ref{sbwindow:mass} gives
\[
 S(t)\le\alpha s_0,\qquad \alpha=(5/4)^{1/3}<2
 \quad(t\le\tau_{\rm dec}).
\]
The initial residual cost obeys
$2\overline C s_0^2(2+s_0)\le s_0J_*/(4d)$ by the primitive
initialization cap. Its bound at time $t$ is at most $\alpha^3$ times
this value, hence at most $5s_0J_*/(16d)$.
Lemma~\ref{sbend:pairing-motion} costs at most $s_0J_*/(16d)$ on the
chosen interval. Subtracting these two costs from $5s_0J_*/(8d)$ leaves
$s_0J_*/(4d)$. Integrate the almost-everywhere derivative bound; the
absolute continuity statement includes gate changes. The two arc-width
bounds follow by integrating $(\arcsin r)'=(1-r^2)^{-1/2}$.
\end{proof}

\subsection{The correct time bookkeeping for an O2b endpoint result}

Initial descent supplies a drop at $t_{\rm early}$. It does not supply
$E(0)>E(t_{\rm spec})$. Consequently an initial drop and a distinct
post-specialization rebound naturally use four snapshots, not a tacit
identification of $t_{\rm early}$ with a specialization time.

\begin{proposition}[Four snapshots and an explicit final sector]
\label{sbend:four-snapshots}
Suppose $0<t_{\rm early}<t_{\rm spec}\le t_a<t_b\le H$,
and $\xi_0$ is a probe in the central half of the angular arc
$\mathcal I$. In addition to Theorem~\ref{sbend:descent-clock}, suppose
\[
 E(\xi_0,t_b)-E(\xi_0,t_a)\ge b_*>0.
\]
Put
\begin{align}
 S_H&:=\min\{s_0e^{2\overline\lambda H},
                         s_0+\overline C H/2\},\label{sbend:mass-cap}\\
 \rho_*&:=\min\left\{1,
       \frac{\min(a_{\rm dec},b_*)}{4(1+S_H)^2}\right\},\notag\\
 \Delta_{\rm swing}&:=\min\{\Delta_{\rm init}/2,
                                  4\arcsin(\rho_*/2)\}>0.
 \label{sbend:sector-width}
\end{align}
There is an open sector centered at $\xi_0$ of angular width
$\Delta_{\rm swing}$ on which the initial drop is at least
$a_{\rm dec}/2$ and the post-specialization rebound is at least $b_*/2$.
Each probe error has an interior minimum on $[0,t_b]$ and a rebound from
that minimum of at least $b_*/2$. This is an O2b-endpoint certificate when
the stipulated specialization state is actually certified. It does not
place that minimum after $t_{\rm spec}$.
\end{proposition}
\begin{proof}
The exponential and linear mass bounds, respectively
Lemmas~\ref{sbwindow:mass} and~\ref{sbnew:global}, give $S\le S_H$.
Neuronwise balance implies
$\sum_i\|w_i\|\|u_i\|=S$, so this also bounds the snapshot Lipschitz
constant. Apply the proof of Proposition~\ref{sbaudit:secants} to the
two separate pairs of times. A chord radius $\rho_*$ is an angular
half-width $2\arcsin(\rho_*/2)$. Intersect with the initial arc; the
central-half assumption leaves an angular half-width at least
$\Delta_{\rm init}/4$.
The global minimum is interior because the left endpoint exceeds the
value at $t_{\rm early}$ and the right endpoint exceeds the value at
$t_a$, even though those two intermediate values need not be ordered.
\end{proof}

\subsection{Which transverse output removal can actually cause a rebound?}

For a fixed $\xi=re_1-ke_2$, $k=\sqrt{1-r^2}$, put
\[
 X=f_1(\xi)-r,\qquad \Lambda=-f_2(\xi),\qquad P=f_\perp(\xi).
\]
The exact snapshot identity is
\begin{equation}\label{sbend:leakage-secants}
 E_b-E_a=(\Lambda_a-\Lambda_b)
          \left[k-\frac{\Lambda_a+\Lambda_b}{2}\right]
       +\frac{X_b^2-X_a^2}{2}
       +\frac{\|P_b\|^2-\|P_a\|^2}{2}.
\end{equation}
In active-matrix notation, this scalar is
\[
 \Lambda_j=k(M_j)_{22}-r(M_j)_{21}.
\]
Every neuron is included and the active sets may differ at the two times.
For the noise-scaled normal form it is
$\Lambda=\eta^2(kd-zc)$, not either cross entry separately.

\begin{proposition}[Minimal transverse snapshot certificate]
\label{sbend:leakage-certificate}
At two reached times $t_{\rm spec}\le t_a<t_b\le H$, suppose
$|\Lambda_a|,|\Lambda_b|\le\ell<k$, $|X_a|\le\alpha$, and
$\Lambda_a-\Lambda_b\ge d_\Lambda>0$. Then
\begin{equation}\label{sbend:leakage-budget}
 E_b-E_a\ge (k-\ell)d_\Lambda
                          -\frac{\alpha^2+s_0S_H}{2}.
\end{equation}
A strictly positive right side supplies $b_*$ in
Proposition~\ref{sbend:four-snapshots}.
\end{proposition}
\begin{proof}
Only the possibly negative longitudinal difference is bounded by
$-X_a^2/2$. The installed ambient contraction gives
$\|P_a\|\le\|W_\perp(t_a)\|_F\sqrt{S(t_a)}
\le\sqrt{s_0S_H}$. Discard the nonnegative final ambient square.
The remaining estimate is exactly \eqref{sbend:leakage-secants}.
\end{proof}

Thus removing a \emph{negative} coordinate-2 output that helped reconstruct
this probe can supply a rebound. In contrast, if $f_2\ge0$ and only this
coordinate decreases while the others remain fixed, the error decreases:
the derivative of its squared residual is $(k+f_2)\dot f_2\le0$.
Likewise ambient output has no target component here; reducing its norm
alone improves reconstruction. The contraction of the stacked ambient
weights does not by itself imply pointwise contraction of $P(\xi)$ because
input activations also move. These signs prevent using unspecified
``leakage decay'' as the missing later witness.

For example, a held positive strong family has $z_{i2}\ge0$. Its own
contribution to $\Lambda$ is nonpositive at every probe, irrespective of
the input gates. A nearly positive-$e_2$ weak family is inactive on this
negative-$e_2$ sector whenever its input angles exceed $\arcsin r$.
Beneficial negative transverse output, if used, must therefore be tracked
in the actual active complement or in a strong family not yet known to
have positive output coordinates. Positivity of the aggregate training
coefficient $c_{21}^O$ does not settle this probe-selected quantity.

\subsection{A zero-lag endpoint mechanism: removing a useful probe response}

Small lag is not an obstruction to the snapshot route. A perfectly
matched coherent planar block with unit direction $n(t)$ and mass $m(t)$
has $f_B(\xi)=m n(n^\top\xi)_+$ and, with no background,
\begin{equation}\label{sbend:matched-benefit}
 E(\xi,t)=\frac12-\frac12m(t)(2-m(t))(n(t)^\top\xi)_+^2.
\end{equation}
For $0<m\le1$, losing the positive probe activation can return the
error toward $1/2$ even with identically zero input--output angular lag.
The gate-exit derivative vanishes at the exact zero-lag exit, but the
finite endpoint gain in \eqref{sbend:matched-benefit} need not vanish.
This calculation is algebraic; a matched rotating block is not asserted
to follow the actual gradient flow.

\begin{proposition}[Post-specialization gate removal with signed background]
\label{sbend:sweep-certificate}
Fix a cohort $B$, two reached times $t_{\rm spec}\le t_a<t_b\le H$,
and a probe $\xi=(\sin\theta,-\cos\theta)$. Let $0<g\le1/4$.
At $t_a$ suppose every $i\in B$ has nonzero projected radii and
\[
 g\le\theta-\phi_i\le5g/4,\qquad
 g\le\theta-\psi_i\le5g/4,\qquad
 m_*\le M_B:=\sum_{i\in B}r_is_i\le1.
\]
At $t_b$ suppose all its probe gates are off. Let $G_j$ denote the full
output of $B^c$ at the probe at time $t_j$, let
$F_a=f_B(\xi,t_a)$, and define the signed background cost
\begin{equation}\label{sbend:background-cost}
 \mathfrak C_B:=\langle G_a,F_a\rangle
 +\frac12\|G_a-\xi\|^2-\frac12\|G_b-\xi\|^2
 +\frac12\|P_{\mathcal S^\perp}F_a\|^2.
\end{equation}
If $\mathfrak C_B\le7m_*g^2/96$, then
\begin{equation}\label{sbend:sweep-gap}
 E(\xi,t_b)-E(\xi,t_a)\ge m_*g^2/8>0.
\end{equation}
This conclusion requires no lower bound on angular lag. The background
contains every neuron outside $B$, including weak and unassigned neurons.
\end{proposition}
\begin{proof}
The two angle bounds give
$\xi^\top F_a\ge M_B\sin^2g$ and
$\|P_{\mathcal S}F_a\|\le M_B\sin(5g/4)\le5M_Bg/4$.
Since $M_B\le1$ and
$\sin^2g\ge g^2-g^4/3\ge47g^2/48$,
\[
 \xi^\top F_a-\frac12\|P_{\mathcal S}F_a\|^2
 \ge M_Bg^2\left(\frac{47}{48}-\frac{25}{32}\right)
 =\frac{19}{96}M_Bg^2.
\]
The exact endpoint expansion, using $f_B(\xi,t_b)=0$, is this
benefit minus $\mathfrak C_B$, leaving at least
$(19-7)m_*g^2/96=m_*g^2/8$.
\end{proof}

\subsection{Exact lag profile and its conditional viability}

Let $n_i=v_i/r_i$, $t_i=n_i^\angle$, $q_i=s_i/r_i$, and
$\zeta_i=\psi_i-\phi_i$. In this subsection $q_i$ is the radial ratio,
not the longitudinal ratio used in Section~\ref{sbpos:scope}.
For the actual compatible selected matrix write
\[
 a_i=n_i^\top A_i n_i,\quad b_i=n_i^\top A_i t_i,\quad
 c_i=t_i^\top A_i n_i,\quad d_i=t_i^\top A_i t_i,\quad
 p_i=t_i^\top P_i/r_i.
\]
The exact polar equations are
\[
 \dot\phi_i=q_i(b_i\cos\zeta_i+d_i\sin\zeta_i)-p_i,\qquad
 \dot\psi_i=q_i^{-1}(c_i\cos\zeta_i-a_i\sin\zeta_i).
\]
Thus, with
$\alpha_i=q_i^{-1}a_i+q_id_i$ and
$\beta_i=q_i^{-1}c_i-q_ib_i$,
\begin{equation}\label{sbend:lag-exact}
 \dot\zeta_i=-\alpha_i\sin\zeta_i+\beta_i\cos\zeta_i+p_i.
\end{equation}
There is no extra boundary approximation here: $A_i$ already uses the
actual compatible selectors. In particular
\begin{equation}\label{sbend:lag-skew}
 \beta_i=X_i+(q_i^{-1}-1)c_i-(q_i-1)b_i,\qquad
 X_i=\mathbb E_n[(f_1x_2-f_2x_1)\alpha_i(x)].
\end{equation}
For $q_i=1$ and zero ambient forcing, the forcing at zero lag is
exactly $X_i$, while the damping coefficient is $\operatorname{tr}A_i$.
Both depend on the learned residual; neither has an automatic late sign.

Put $y_i=\sin\zeta_i$, $k_i=\alpha_i\cos\zeta_i$, and
$F_i=\beta_i\cos^2\zeta_i+p_i\cos\zeta_i$. Then the exact, scalar
variation-of-constants identity is
\begin{equation}\label{sbend:lag-profile}
 y_i(t)=e^{-\int_a^t k_i}y_i(a)
             +\int_a^t e^{-\int_v^t k_i}F_i(v)\,dv.
\end{equation}
For example, on a stopped interval with
$0<\underline k\le k_i\le\overline k$ and $|F_i|\le F_*$,
\[
 |y_i(t)|\le e^{-\underline k(t-a)}|y_i(a)|
       +\frac{F_*}{\underline k}(1-e^{-\underline k(t-a)}).
\]
If instead $F_i\le-F_-<0$ there, the one-sided estimate is
\begin{equation}\label{sbend:lag-lower}
 -y_i(t)\ge
 \frac{F_-}{\overline k}(1-e^{-\overline k(t-a)})
                 -e^{-\underline k(t-a)}[y_i(a)]_+.
\end{equation}
Thus forcing of order $e$ can sustain lag of order $e$ over a fixed
window. Monotone lag decay is not a consequence of the mismatch damping
term. Conversely, a nonzero lower lag bound requires a signed forcing
estimate such as the one in \eqref{sbend:lag-lower}; an upper Cartesian
mismatch bound cannot supply it.

At the outward exit normal, the zero-background contribution is
$-r_is_i\dot\phi_i y_i$. The full contribution plus background is
\[
 -r_is_i\dot\phi_i y_i
 -r_i\dot\phi_i\langle G,w_i\rangle
 +\langle G-\xi,\dot G\rangle.
\]
The last two terms are signed. An $O(e)$ lag and $O(e)$ angular velocity
only give an $O(r_is_i e^2)$ individual exit contribution; they do not
bound the rest of the network. This is why the endpoint route remains
primary even after the lag audit.

\section{A learned-mass endpoint theorem in an extreme-separation region}
\label{sbsep:scope}

This section proves a restricted noisy many-neuron result from isotropic
initialization, with a learned strong cohort of mass at least $1/2$ before
the rebound. It does not prove learned weak-family specialization. The
region permits $\mu_2/\mu_1$ and the noise to be very small relative to the
initialization scale. Consequently this is a candidate one-concept
O2b-endpoint theorem, not a closure of the original two-family theorem or
of its fixed concept-ratio regime. Both restrictions are theorem-scope
choices; neither is evidence of an obstruction to the original goal.

\subsection{Directional concentration removes an artificial width conflict}

\begin{lemma}[Directional initial fluctuation bound]
\label{sbsep:directional-variance}
In Theorem~\ref{sbaudit:initial-descent}, set
$C_\xi=\mathbb E_n[|\xi^\top x|\|x\|]$. One may replace $C_x$ by
$C_\xi$ in the definition of $b_h$, while retaining $C_x$ in the residual
cost. If $C_\xi\le C_{\rm dir}$ on a deterministic arc, the uniform-grid
proof likewise permits replacing $\overline C$ by $C_{\rm dir}$ in its
width requirement, but not in its net radius or initialization cap.
\end{lemma}
\begin{proof}
Write $v_x=x/\|x\|$ and use weights
$|\xi^\top x|\|x\|/C_\xi$ in Jensen's inequality. Then
\[
 |Y_i|\le C_\xi,\qquad
 \mathbb EY_i^2\le C_\xi^2\mathbb E_n^{\rm weighted}
 \mathbb E[(\widehat u^\top\xi)^2(\widehat u^\top v_x)^2]
 \le\frac{3C_\xi^2}{d(d+2)}.
\]
The Bernstein and net arguments are unchanged. In particular a target
margin of order $r$ and $C_{\rm dir}=O(r)$ do not force $h\gtrsim r^{-2}$.
\end{proof}

\subsection{Exact all-neuron rank-one reference and uniform comparison}

Set $\mu_1=1$. The reference dataset has equal mass at $e_1$ and zero;
its time variable is $\tau=t/2$. Let $I_+=\{i:u_{i1}(0)>0\}$.
On these neurons write the input matrix as $(a^\top,B^\top)^\top$ and
the output matrix as $(c^\top,D^\top)^\top$, separating $e_1$ from its
orthogonal complement. Here $a,c\in\mathbb R^{|I_+|}$ and
$B,D\in\mathbb R^{(d-1)\times|I_+|}$. While the strong training gates
are retained, $B$ is constant and
\begin{equation}\label{sbsep:exact-system}
 a'=(1-a^\top c)c-D^\top Da,\qquad
 c'=(1-a^\top c)a,\qquad D'=-Daa^\top.
\end{equation}
All neurons in $I_+^c$ are stationary. Balanced isotropic initialization
gives $a(0)=c(0)=a_0$, $D(0)=B$. Therefore
\begin{equation}\label{sbsep:gram}
 cc^\top+D^\top D-aa^\top=B^\top B,\qquad
 a'=c-\|a\|^2a-B^\top Ba,\qquad \|D\|_F\le\|B\|_F.
\end{equation}
This reference includes every neuron, not a selected two-neuron subsystem.

Put $A_0=\|a_0\|$, $v=a_0/A_0$, and define
\begin{equation}\label{sbsep:scalar-reference}
 A'=A(1-A^2),\quad A(0)=A_0,\qquad
 d_0(\tau)=\sqrt{\frac{1-A(\tau)^2}{1-A_0^2}},\qquad
 D^\dagger=B[I-(1-d_0)vv^\top].
\end{equation}
The superscript $\dagger$ below means $a^\dagger=c^\dagger=Av$ and
$D^\dagger$ as above, with the actual constant $B$ and the actual inactive
neurons. It is an approximation to \eqref{sbsep:exact-system}, not an
assertion that all input or output directions coincide.

\begin{lemma}[Relative comparison through learned mass]
\label{sbsep:relative-comparison}
Suppose $\|B\|_F^2\le s_0$, $A_0^2\le s_0$, $v_i\ge v_*>0$,
and $0<\nu\le1/64$. On a scalar-reference interval $0\le\tau\le T\le H$
with $A^2\le1-\nu$, set
\[
 \rho=\frac{3s_0H}{\nu},\qquad
 \delta_c=4\rho+\frac{4s_0}{\nu},\qquad
 D_* =\frac{36s_0^{3/2}H^2}{\nu}.
\]
If
\begin{equation}\label{sbsep:comparison-caps}
 s_0\le\frac{\nu^2}{64(H+1)},\qquad \delta_c\le v_*/4,
\end{equation}
then all strong gates are retained throughout and
\begin{align}
 \|a-Av\|&\le4A\rho,&\quad\|c-Av\|&\le A\delta_c,
 &\quad\|D-D^\dagger\|_F&\le D_* .\label{sbsep:comparison-bounds}
\end{align}
The total reference mass is at most $1+2s_0$. If $A^2\ge3/4$, then
$a^\top c\ge2/3$ and $a_i,c_i\ge v_*/2$.
\end{lemma}
\begin{proof}
Let $u=\|a\|$, $p=c-a$, $l=\|p\|/u$ and
$\vartheta=a^\top p/u^2$. Stop before a strong gate changes or either
$u^2=1-\nu/2$ or $l=\nu/8$. Equations \eqref{sbsep:gram} give
\[
 \frac{u'}u=1-u^2+\vartheta-\widehat a^\top B^\top B\widehat a=:g,
 \qquad p'=-(1-a^\top c)p+D^\top Da.
\]
On the stopped interval $g\ge\nu/4$, and
$1-a^\top c+g\ge\nu/2$. In the upper-Dini sense,
$l'\le-(\nu/2)l+s_0$, hence $l\le2s_0/\nu$.
Writing $q=\log(u/A)$ yields
$q'=-A^2(e^{2q}-1)+\epsilon_q$, $|\epsilon_q|\le3s_0/\nu$.
The restoring sign gives $|q|\le3s_0\tau/\nu$ without an exponential
in $H$. Also $\|\widehat a'\|\le l+s_0\le3s_0/\nu$, so
$\|\widehat a-v\|\le3s_0\tau/\nu$.
These bounds give the first two inequalities in
\eqref{sbsep:comparison-bounds}. The caps make
$\rho<\nu/8$ and exclude both radial stopping faces:
$(1-\nu)e^{\nu/4}<1-\nu/2$ and $2s_0/\nu<\nu/8$.
They also give $a_i,c_i\ge3Av_i/4$, which excludes gate exit and
proves continuation from the initially positive gates.

The error $Z=D-D^\dagger$ satisfies
\[
 Z'=-Zaa^\top-D^\dagger(aa^\top-A^2vv^\top).
\]
The first term dissipates $\|Z\|_F^2$.
Since $\|aa^\top-A^2vv^\top\|\le12A^2(3s_0\tau/\nu)$,
integration gives the stated $D_*$ (with a factor-two reserve).
Finally $\|a\|^2<1$, $\|B\|_F^2\le s_0$, and the inactive input
mass is at most $s_0$. This gives the full mass bound. At $A^2\ge3/4$,
\[
 a^\top c\ge A^2e^{-2\rho}(1-2s_0/\nu)
 \ge\tfrac34(63/64)^2>2/3,
\]
and $a_i,c_i\ge3Av_i/4\ge v_*/2$.
\end{proof}

\subsection{The signed endpoint gain}

Let $B_2$ be the row of $B$ corresponding to $e_2$ and set
$b=Bv$, $b_2=B_2v$, $\eta=|b_2|>0$, and $s=\operatorname{sign}b_2$.
Choose $r=8\sqrt{s_0}/v_*\le1/8$, $k=\sqrt{1-r^2}$,
$\xi_0=re_1+sk e_2$, and $z=r/\eta$.
Assume $\eta\ge\eta_*>0$, let $Z_*=r/\eta_*$, and take
$\nu=1/(512Z_*^2)$. In particular $8\le z\le Z_*$.
Define scalar-reference times by
\begin{equation}\label{sbsep:snapshot-times}
 A(\tau_{\rm spec})^2=3/4,\qquad
 (A(\tau_a)^2-1)z+A(\tau_a)k=0,\qquad
 A(\tau_b)^2=1-\nu.
\end{equation}
Then $\tau_{\rm spec}<\tau_a<\tau_b$. These times may depend on
initialization; all will have a deterministic upper bound below.
For $\tau\ge\tau_{\rm spec}$ every $I_+$ probe gate is strictly on.
Indeed $ra_i\ge rv_*/2=4\sqrt{s_0}>|B_{2i}|$.

Write $G$ for the stationary output of $I_+^c$ at $\xi_0$, so
$\|G\|\le s_0$. Put $C=BB_2^\top-bb_2$, with $\|C\|\le2s_0$.
The approximate full output is exactly
\begin{equation}\label{sbsep:probe-reference}
 f^\dagger(\xi_0)=
 e_1(A^2r+A\eta k)+(0,skC)+\eta(0,b)T+G,
 \qquad T=d_0(Az+k).
\end{equation}
The transverse component along $b$ initially helps reconstruct the
probe and subsequently decreases; the inactive complement $G$ is
retained in the endpoint squares.

\begin{lemma}[Endpoint rebound at learned mass]
\label{sbsep:reference-rebound}
In addition to Lemma~\ref{sbsep:relative-comparison}, suppose
\begin{equation}\label{sbsep:rebound-caps}
 s_0\le\min\{1/256,1/(32Z_*^2)\},\quad
 s_0^{3/2}\le\eta_*/16,\quad s_0^2\le\eta_*^2/8,
\end{equation}
and
\begin{equation}\label{sbsep:output-cap}
 \epsilon_f:=4r\delta_c+3rD_*+4r\sqrt{s_0}\rho
 \le\eta_*^2/128.
\end{equation}
Then the exact rank-one reference obeys
$E(\xi_0,2\tau_b)-E(\xi_0,2\tau_a)\ge3\eta^2/8$.
Here $E(\xi_0,2\tau)$ denotes the reference error at physical time
$t=2\tau$.
\end{lemma}
\begin{proof}
The unique root defining $\tau_a$ has $A_a>9/10$: evaluate its quadratic
at $A=9/10$ and use $z\ge8$. Moreover
$1-A_a^2=A_ak/z>\nu$, proving the strict ordering.
At the two endpoints,
\[
 T_a\ge\tfrac9{16}\sqrt z,\qquad
 T_b\le\sqrt{2\nu}(Z_*+1)\le1/8,
 \qquad T_a-T_b\ge1,\quad 0\le T_a,T_b\le2Z_*.
\]
Set $Q=skC+G_\perp-sk e_2$, viewing $e_2$ in the $(d-1)$-dimensional
transverse space. Then
\[
 Q^\top b=-\eta k+(skC+G_\perp)^\top b,
 \qquad |(skC+G_\perp)^\top b|\le3s_0^{3/2}\le\eta k/4.
\]
Thus the transverse error secant in \eqref{sbsep:probe-reference} is
\[
 \eta(T_b-T_a)Q^\top b+
 \tfrac12\eta^2\|b\|^2(T_b^2-T_a^2)
 \ge\tfrac9{16}\eta^2-2\eta^2s_0Z_*^2
 \ge\tfrac8{16}\eta^2.
\]
At $\tau_a$ the longitudinal residual of the variable part is zero.
Its full initial longitudinal residual is $G_1$, so its secant is
bounded below by $-G_1^2/2\ge-\eta^2/16$.
The approximate full secant is therefore at least $7\eta^2/16$.

At either endpoint, \eqref{sbsep:comparison-bounds} and
$\|ra+skB_2^\top\|\le2r+\sqrt{s_0}$ give
$\|f-f^\dagger\|\le\epsilon_f$; the first-coordinate cost is at most
$\delta_c(3r+\sqrt{s_0})\le4r\delta_c$, and the transverse cost is
at most $3rD_*+4r\sqrt{s_0}\rho$.
The exact reference has mass at most $1+2s_0$, and $\epsilon_f<1/16$,
so an error perturbation of norm $\epsilon_f$ changes each error square
by at most $4\epsilon_f$. The two-endpoint cost is at most
$8\epsilon_f\le\eta^2/16$, as required.
\end{proof}

\subsection{Nonempty primitive noisy region and high-probability initialization}

The following explicit choices are deliberately conservative. They certify
nonemptiness, not a useful experimental parameter range. All constants
are primitive; none presupposes a trajectory event.
Fix $d\ge2$, $0<\delta<1/4$, and an even total sample count $N\ge2$.
Set
\[
 G_*=2+\lceil256\pi d\rceil,\qquad \ell=\log(8G_*/\delta),
\]
and choose the integer width so that
\begin{equation}\label{sbsep:width}
 h\ge\max\{1536(16\pi)^2\ell,
           (1024\pi/3)d\ell,32\log(16/\delta)\}.
\end{equation}
Define
\begin{align*}
 \beta&=\frac{\delta}{32h\sqrt d},&
 v_*&=\frac\beta{\sqrt h},&
 b_*&=\frac{\delta^3}{262144\,d\,h^{3/2}},\\
 s_0&=h\varepsilon^2,& \eta_*&=\varepsilon b_*,&
 r&=\frac{8\sqrt{s_0}}{v_*},\\
 Z_*&=\frac r{\eta_*}=\frac{8\sqrt h}{v_*b_*},&
 \nu&=\frac1{512Z_*^2},&
 H&=\frac12\log\frac4{\nu s_0\beta^2}.
\end{align*}
Use $\rho,\delta_c,D_*,\epsilon_f$ from the preceding lemmas with
this $H$. Choose $\varepsilon>0$ sufficiently small that all of
\eqref{sbsep:comparison-caps}, \eqref{sbsep:rebound-caps}, and
\eqref{sbsep:output-cap} hold, together with
\begin{equation}\label{sbsep:remaining-caps}
 r\le1/8,\qquad s_0\le\frac r{384\pi d},\qquad
 \frac{8\sqrt{s_0}\rho}{v_*^2}
       +\frac{3(\delta_c+D_*)}{v_*}\le\eta_*/40.
\end{equation}
Finally set
\begin{align}
 \chi_*&=\min\{1/32,\varepsilon\beta/8,
                   \eta_*^2/1024,\eta_*v_*/240\},\notag\\
 q_D&=\chi_*e^{-20H},\label{sbsep:data-caps}\\
 0<\mu_2^2&\le\min\{r^2/100,q_D/256\},\notag\\
 0<a_D&\le\min\{\mu_2/4,r^2/100,
                           \varepsilon\beta/16,q_D/160\},\notag\\
 0<\sigma_1,\sigma_2&\le
                   a_D/\sqrt{2\log(4N/\delta)}.\notag
\end{align}
Here $\sigma_k$ are standard deviations of planar isotropic Gaussian
noise, there are $N/2$ independent samples per cluster, and $\mu_1=1$.

\begin{theorem}[Noisy many-neuron strong-cohort specialization and endpoint rebound]
\label{sbsep:theorem}
For the primitive region above, balanced isotropic independent
initialization $u_i(0)=w_i(0)=\varepsilon\widehat u_i$ has the following
properties with probability at least $1-\delta$ jointly with the samples.
There is a fixed initialization cohort $I_+$ of at least $h/4$ neurons,
and times
\[
 0<t_{\rm early}<t_{\rm spec}<t_a<t_b\le2H,
\]
such that throughout $[t_{\rm spec},t_b]$ its input and output vectors
lie within angle $r/2$ of $e_1$ and its learned longitudinal mass satisfies
\begin{equation}\label{sbsep:learned-mass}
 \sum_{i\in I_+}u_{i1}(t)w_{i1}(t)\ge1/2.
\end{equation}
The entire network has $S(t)\le4$ on $[0,t_b]$.
On each of the two deterministic arcs
\[
 \mathcal I_\pm=\{\sin\theta\,e_1\pm\cos\theta\,e_2:
       |\theta-\arcsin r|\le r/4\},\qquad
 \Delta_{\rm init}=r/2,
\]
there is uniform initial decrease with the explicit
$\tau_{\rm dec},c_{\rm dec},t_{\rm early},a_{\rm dec}$ of
Theorem~\ref{sbend:descent-clock}, taking
$J_*=r/(8\pi)$ and $\overline C=\overline\lambda=2$.
At one of their centers $\xi_0$ the later secant is
\begin{equation}\label{sbsep:noisy-rebound}
 E(\xi_0,t_b)-E(\xi_0,t_a)\ge b_{\rm reb}:=\eta_*^2/8>0.
\end{equation}
On an open sector, the initial drop is at least $a_{\rm dec}/2$.
The rebound after specialization is at least $b_{\rm reb}/2$.
An explicit lower bound on its angular width is
\begin{equation}\label{sbsep:final-sector}
 \min\{r/4,4\arcsin(\rho_{\rm swing}/2)\},\qquad
 \rho_{\rm swing}=\min\{1,\min(a_{\rm dec},b_{\rm reb})/100\}.
\end{equation}
Every such error has an interior minimum on $[0,t_b]$.
The conclusion does not locate the minimum after $t_{\rm spec}$, fit the
weak concept, or create a second learned family.
\end{theorem}
\begin{proof}
\emph{The initialization geometry.}
Let $R_i^2=\widehat u_{i1}^2+\widehat u_{i2}^2$ and
$R_*^2=\delta/(32dh)$. With failure probability less than $\delta/4$,
\begin{equation}\label{sbsep:geometry-event}
 |I_+|\ge h/4,\quad
 \min_i|\widehat u_{i1}|\ge\beta,\quad
 \min_iR_i^2\ge R_*^2,\quad |B_2v|\ge\varepsilon b_*.
\end{equation}
Indeed the first failure is at most $e^{-h/8}\le\delta/32$.
The density of one spherical coordinate on $[-\beta,\beta]$ is at most
$\sqrt d$, giving a union bound $2h\sqrt d\beta=\delta/16$.
For $d>2$, $R_i^2$ has distribution function
$1-(1-q)^{(d-2)/2}\le dq$ for $0\le q\le1/2$; for $d=2$ it is one.
The radius failure is at most $\delta/32$.
Condition on all radii, signs of the first coordinates, and all but the
first positive neuron's planar angle. Its angle is uniform on a
half-circle. Its contribution to
$\sum_{I_+}\widehat u_{i1}\widehat u_{i2}$ is
$(R_i^2/2)\sin(2\theta_i)$. The elementary arcsine small-ball bound
\[
 \sup_z\Pr\{|\sin(2\theta_i)-z|\le u\}\le4\sqrt u
\]
follows by integrating its density on an interval, with the largest
endpoint singularity proportional to $1/\sqrt{1-|x|}$.
Taking $q=R_*^2\delta^2/8192$ bounds the probability that this sum has
absolute value less than $q$ by $\delta/16$.
Since $A_0\le\varepsilon\sqrt h$, this proves the last bound in
\eqref{sbsep:geometry-event}. This conditioning is on signs and radii,
not on the minimum-coordinate event; a union bound combines them.
Moreover
\[
 s_0\beta^2/4\le A_0^2\le s_0,\qquad v_i\ge\beta/\sqrt h=v_*.
\]
The explicit logistic solution gives
$\tau_b=\tfrac12\log[(1-\nu)(1-A_0^2)/(\nu A_0^2)]\le H$.
Apply the preceding reference lemmas up to this time.

\emph{Uniform noisy comparison, including arbitrary weak gates.}
On the sample event $\max\|x^{(k)}-\mu_ke_k\|\le a_D$, the failure
probability is at most $\delta/4$ by the planar Gaussian radial tail
and a union bound over $N$ samples. Couple the real and reference
networks with the same initialization. Stop while both full input and
output Frobenius norms are at most $2$, and all strong-sample gates
agree with $I_+$. The reference vector field, in physical time and
combined parameter Frobenius norm, is $10$-Lipschitz there.
For example its longitudinal-input and output difference bounds have
coefficients $2$ and $9/2$, whose block norm is at most $13/2<10$.
At a fixed parameter point the strong-sample covariance changes by at
most $a_D+a_D^2/2$. Since $\|I-WU_{I_+}^\top\|\le5$, the combined
strong-gradient cost is at most $30a_D$. The entire weak gradient,
with its actual arbitrary selectors, has combined norm at most
$10(\mu_2+a_D)^2$. Thus a sufficient forcing bound is
\[
 P_D=40[a_D+(\mu_2+a_D)^2].
\]
The parameter distance up to physical time $2H$ is at most
\begin{equation}\label{sbsep:noisy-distance}
 \chi:=P_De^{20H}\le\chi_*.
\end{equation}
This is a finite-window comparison with the vanishing data perturbation,
not a Gronwall estimate on the $B^\top B$ amplification error; that error
was controlled relatively in Lemma~\ref{sbsep:relative-comparison}.
The caps in \eqref{sbsep:data-caps} give $P_D<q_D$.

Reference strong margins are at least $\varepsilon\beta/2$ for $I_+$,
and at least $\varepsilon\beta$ in absolute value for $I_+^c$.
The perturbations cost at most $\chi+2a_D\le\varepsilon\beta/4$.
Reference parameter norms are at most $\sqrt{1+2s_0}$, so adding
$\chi\le1/32$ also excludes the norm stopping face. The comparison
therefore continues to $2\tau_b\le2H$. In particular $S\le4$
through the last snapshot. The deterministic $2H$ is used only as an
upper bound in the comparison integral; no radial estimate past
$2\tau_b$ is needed.

At the reached times $t_j=2\tau_j$, the learned strong mass changes
by at most $4\chi$ from its reference value $2/3$, leaving at least
$1/2$. Its longitudinal coordinates are at least $v_*/2-\chi$ and its
transverse input/output norms at most $\sqrt{s_0}+\chi$, giving angles
at most $4\sqrt{s_0}/v_*=r/2$. At either rebound snapshot the full
output changes by at most $4\chi$, hence its error changes by at most
$16\chi$. The secant loses at most $32\chi\le\eta_*^2/32$, which is
less than the reserve in Lemma~\ref{sbsep:reference-rebound}.
This proves \eqref{sbsep:noisy-rebound}.
The limitation concerning the weak concept is quantitative: throughout this
window, transverse input and output Frobenius norms are each at most
$\sqrt{s_0}+\chi$, so
$|f_2(e_2,t)|\le(\sqrt{s_0}+\chi)^2\le2s_0$.
By homogeneity this is also the coordinate-2 fitting fraction at
$\mu_2e_2$. Thus the theorem does not hide a second learned family.

\emph{The initialized decreasing event on the same primitive region.}
Uniformly on both arcs, the first coordinate of the probe is in
$[3r/4,5r/4]$. The bounded-noise event and the data caps give
$C_x\le2$, $C_\xi\le2r$ and $J_\xi\ge r/(8\pi)$.
For the last assertion use $cF(c)\ge c$ on cluster 1 and the crude
negative bound $cF(c)\ge-1$ on cluster 2. Its numerator is at least
\[
 [(3r/4)(1-a_D)-a_D](1-a_D)-(\mu_2+a_D)^2\ge r/4.
\]
This estimate holds for both signs of the second probe coordinate.
A net with chord mesh $r/(512\pi d)$ requires at most $G_*$ points
per arc. Lemma~\ref{sbsep:directional-variance}, the width condition,
and a union bound over both nets give failure probability at most
$\delta/4$. Keep $\overline C=2$ for grid transport and the initial
residual cost. The initialization cap in
\eqref{sbsep:remaining-caps} is precisely the required one.
Theorem~\ref{sbend:descent-clock} now gives the claimed deterministic
initial drop. Since $A_0^2\le1/256$, $2\tau_{\rm spec}>1$ whereas
$t_{\rm early}<1$, proving the strict ordering.
The geometry, data and descent failures sum to less than $\delta$;
no independence between their derived events is assumed.

Finally apply the four-snapshot argument with the proven mass cap $4$
at the four snapshots. It gives \eqref{sbsep:final-sector}.
All unit probes in these sectors are at distance at least $3/4$ from
every training sample under the displayed data caps, so the statement
concerns extrapolation probes even if the selected sign is positive.

\emph{Nonemptiness.}
Choose $d,\delta,N,h$ first. The constants $v_*,b_*,Z_*,\nu$ do not
depend on $\varepsilon$. As $\varepsilon\downarrow0$,
$H=O(\log(1/\varepsilon))$, $\rho,\delta_c=O(\varepsilon^2H)$,
$D_*=O(\varepsilon^3H^2)$ and
$\epsilon_f=O(\varepsilon^3H)$, whereas $\eta_*^2$ is a positive
constant times $\varepsilon^2$. Every initialization cap is therefore
satisfied for sufficiently small positive $\varepsilon$. Then choose
positive $\mu_2,a_D,\sigma_1,\sigma_2$ successively from
\eqref{sbsep:data-caps}. The region is nonempty and all time, probability
and endpoint budgets hold there simultaneously.
\end{proof}

\subsection{The lag remains of initialization scale after learned mass}

This proof also identifies the lag profile; it does not assume damping
has removed it. For $i\in I_+$ and $\tau_{\rm spec}\le\tau\le\tau_b$,
put $q_i=B_{2i}/(Av_i)$. In the approximation,
\begin{equation}\label{sbsep:lag-profile}
 \zeta_i^\dagger(\tau)=
 \arctan\left(q_i-\frac{(1-d_0)b_2}{A}\right)-\arctan q_i.
\end{equation}
Here $|q_i|\le r/4$, $d_0\le\sqrt{(1/4)/(1-s_0)}$, and
$A\ge\sqrt{3/4}$. The mean-value theorem gives
\[
 \eta/5\le-s\zeta_i^\dagger\le2\eta.
\]
The reference comparison and the noisy parameter perturbation give
\[
 |\zeta_i-\zeta_i^\dagger|
 \le\frac{8\sqrt{s_0}\rho}{v_*^2}
       +\frac{3(\delta_c+D_*)}{v_*}+\frac{6\chi}{v_*}
 \le\eta_*/20.
\]
Consequently the actual planar lag obeys the uniform learned-window bound
\begin{equation}\label{sbsep:lag-lower}
 \eta/8\le-s\zeta_i(t)\le3\eta
 \qquad(t_{\rm spec}\le t\le t_b).
\end{equation}
Thus this trajectory retains a nonzero lag of order $|B_2v|$, inherited
from initialization, after learning order-one strong mass. It is not a
noise-generated lag. No gate-exit witness is used: all strong probe gates
are on during the two endpoint snapshots. The decaying signed transverse
coefficient in \eqref{sbsep:probe-reference} supplies the rebound.

The scope question left by this result is exact: whether an initialized
noisy many-neuron theorem with one learned strong cohort and
initialization-dependent extreme concept separation is an acceptable
narrowing, or whether the paper requires two learned families and/or a
concept ratio bounded away from zero. The theorem above establishes the
former mathematical result; it does not choose it as the paper's headline.

\section{Quantitative audit: typical initialization, residual stability, and a signed endpoint limitation}
\label{sbquant:scope}

The preceding restricted theorem is retained unchanged as R0. This section
does not assert R1 or R2. It proves two initialization improvements, a
replacement for the exponential data-comparison bound, and an exact signed
limitation of the all-active endpoint calculation. The robust-core event
is not substituted for an all-neuron gate event without a proof.

\subsection{A core retaining almost all initial longitudinal weight}

Write $t_i=\widehat u_{i1}$, $y_i=\widehat u_{i2}$,
$Q=\sum_i(t_i)_+^2$, $v_i=(t_i)_+/\sqrt Q$ on $I_+$, and set
\[
 I_{\rm core}=\{i:t_i\ge(8\sqrt d)^{-1}\}.
\]
\begin{lemma}[Robust initialization geometry]\label{sbquant:core}
With failure probability at most $2e^{-h/64}+e^{-h/192}$,
\begin{equation}\label{sbquant:core-event}
 \frac{h}{4d}\le Q\le\frac hd,\qquad
 |I_{\rm core}|\ge\frac h{48},\qquad
 v_i\ge\frac1{8\sqrt h}\ (i\in I_{\rm core}),\qquad
 \sum_{I_{\rm core}}v_i^2\ge\frac{15}{16}.
\end{equation}
In particular $s_0/(4d)\le A_0^2\le s_0/d$.
For $0<\alpha<1$, with additional failure at most $\alpha$,
\begin{equation}\label{sbquant:individual}
 \max_i|B_{2i}|\le\varepsilon b_{\max},\qquad
 b_{\max}=\min\left\{1,
 2\sqrt{\frac{\log(\sqrt2h/\alpha)}d}\right\}.
\end{equation}
The deterministic bound $\|B_{\cdot i}\|\le\varepsilon$ also holds.
\end{lemma}
\begin{proof}
For $Y=d(t)_+^2$, $\mathbb EY=1/2$, $\mathbb EY^2\le3/2$.
Spherical moments give
$\mathbb E\exp(\lambda d t^2)\le(1-2\lambda)^{-1/2}$ for
$0\le\lambda<1/2$. At $\lambda=1/10$ this implies
$\Pr(Q\ge h/d)\le e^{-h/32}$.
The elementary bound $e^{-x}\le1-x+x^2/2$ gives
$\mathbb Ee^{-Y/10}\le0.9575$, whence
$\Pr(Q\le h/(4d))\le(e^{1/40}0.9575)^h\le e^{-h/64}$.
For $Z=dt^2$, $\mathbb EZ=1$, $\mathbb EZ^2\le3$.
Paley--Zygmund and symmetry imply
$\Pr(t\ge1/\sqrt{2d})\ge1/24$.
This event is contained in the core. The multiplicative Chernoff bound
therefore gives its count with failure at most $e^{-h/192}$.
The stated coordinate bound follows from the upper bound for $Q$.
The omitted longitudinal weight is at most
$h/(64dQ)\le1/16$. Finally
$\mathbb E\exp(dy_i^2/4)\le\sqrt2$ and a union bound prove
\eqref{sbquant:individual}.
\end{proof}

\subsection{Anti-concentration of the actual dependent coefficient}

\begin{lemma}[Finite-width typical transverse coefficient]
\label{sbquant:smallball}
For $u>0$, the actual initialization coefficient satisfies
\begin{equation}\label{sbquant:smallball-bound}
 \Pr\{|B_2v|<\varepsilon u\}
 \le 2e^{-h/64}
 +2u\sqrt{\frac{d+2}{\pi}}
 +\frac{8\sqrt2}{\pi\sqrt h}.
\end{equation}
Here $Q=0$ is counted as a failure. For example, with
$h\ge(64/\delta)^2$ and $0<\delta<1/4$, the choice
\[
 \eta_* =\frac{\delta\varepsilon}{32\sqrt{d+2}}
\]
has failure less than $\delta/8$. There is no negative power of $h$ in
this lower bound.
\end{lemma}
\begin{proof}
Use independent scalar summands, rather than incorrectly making $B_2$
independent of $v$:
\[
 X_i=t_i y_i\mathbf1_{t_i>0},\qquad
 B_2v=\varepsilon\frac{\sum_iX_i}{\sqrt Q}.
\]
Reflection in coordinate 2 and spherical moments give
\[
 \mathbb EX_i=0,\quad
 \omega^2:=\mathbb EX_i^2=\frac1{2d(d+2)},\quad
 \mathbb E|X_i|^3=\frac4{\pi d(d+2)(d+4)}.
\]
Consequently $\mathbb E|X_i|^3/\omega^3\le8\sqrt2/\pi$.
The iid Berry--Esseen inequality with the conservative constant $1/2$
gives Kolmogorov distance at most $4\sqrt2/(\pi\sqrt h)$.
On $Q\le h/d$, the small-ball event is contained in
\[
 \left\{\left|\frac{\sum_iX_i}{\omega\sqrt h}\right|
 <u\sqrt{2(d+2)}\right\}.
\]
The standard normal density is at most $(2\pi)^{-1/2}$.
Two distribution-function errors and Lemma~\ref{sbquant:core} prove
\eqref{sbquant:smallball-bound}. The numerical corollary follows by
substitution. The Berry--Esseen constant used here is weaker than the
proved bound $0.4756$ in I. Shevtsova,
\emph{On the absolute constants in the Berry--Esseen type inequalities
for identically distributed summands} (2011),
\href{https://arxiv.org/abs/1111.6554}{arXiv:1111.6554}.
\end{proof}

If a dynamical comparison supplies $a_i\ge Av_i/2$ on the core and
$A\ge\sqrt3/2$, its probe gates are positive as soon as
\begin{equation}\label{sbquant:core-radius}
 r\ge32\varepsilon\sqrt h\,b_{\max}.
\end{equation}
This is a core-only statement. It says neither that the other probe
gates are on nor that their training gates persist. Core output-angle
control additionally requires a columnwise bound on $D$; an input
coordinate estimate alone does not establish it.

\subsection{Training-residual stability replaces the generic exponential}

\begin{lemma}[One-sided comparison on a common strong training pattern]
\label{sbquant:stability}
Let $\widehat\theta=(\widehat W,\widehat a)$ solve the rank-one reference
flow in physical time, and let $\theta=(W,a)$ solve that same vector field
plus a forcing of combined Frobenius norm at most $P(t)$. Include passive
coordinates in the distance if they are forced. Assume the same fixed
strong training pattern, and put
\[
 e=\|\theta-\widehat\theta\|,
 \widehat R=\|\widehat W\widehat a-e_1\|.
\]
Then, in the upper-Dini sense,
\begin{equation}\label{sbquant:distance-ode}
 e'\le\tfrac12\widehat R e+\tfrac1{32}e^3+P(t).
\end{equation}
While $e\le\chi_*$, define
\[
 K(t)=\exp\left(\frac12\int_0^t\widehat R(v)\,dv
                     +\frac{\chi_*^2t}{32}\right).
\]
Then $e(t)\le K(t)[e(0)+\int_0^tP(v)/K(v)\,dv]$.
This conclusion makes no assertion about changing strong training gates.
\end{lemma}
\begin{proof}
Set $F=Wa$, $\widehat F=\widehat W\widehat a$,
$\Delta F=F-\widehat F$ and $Q_e=(W-\widehat W)(a-\widehat a)$.
For $L(W,a)=\frac12\|Wa-e_1\|^2$, exact expansion gives
\[
 \langle\theta-\widehat\theta,\nabla L(\theta)-\nabla L(\widehat\theta)\rangle
 =\|\Delta F\|^2+
   \langle2(\widehat F-e_1)+\Delta F,Q_e\rangle.
\]
Since $\|Q_e\|\le e^2/2$, completing the square in $\|\Delta F\|$
gives, in normalized reference time,
$\frac12(e^2)'\le\widehat R e^2+e^4/16$.
Physical time contributes the factor $1/2$. Add the forcing and divide
by $e$, using the usual regularization at $e=0$.
\end{proof}

Under Lemma~\ref{sbsep:relative-comparison}, let
\[
 \Gamma=3\delta_c+\sqrt{s_0}(1+4\rho)+D_*.
\]
The reference residual obeys
\[
 \widehat R(2\tau)\le1-A(\tau)^2+\Gamma,
 \qquad
 \exp\left(\frac12\int_0^{2\tau}\widehat R\right)
 \le\frac{A(\tau)}{A_0}e^{\Gamma\tau}.
\]
Indeed the longitudinal mass error is at most $3\delta_c$ and
$\|Da\|\le\sqrt{s_0}(1+4\rho)+D_*$.
Thus if $\Gamma H+\chi_*^2H/16\le\log2$, $e(0)=0$, and
$P(t)\le P$, a sufficient bootstrap bound through $2\tau_b\le2H$ is
\begin{equation}\label{sbquant:polynomial-comparison}
 e(t)\le\frac{4HP}{A_0}
 \le\frac{8H\sqrt d}{\sqrt{s_0}}P.
\end{equation}
In the R0 common-gate argument the separate forcing allowances are
\[
 P_s\le30a_D,\qquad P_w\le10(\mu_2+a_D)^2.
\]
The new estimate therefore improves the data comparison itself, not just
the constant in $e^{20H}$. Strong-noise and weak-signal allowances remain
different: the former enters linearly in the maximum sample displacement,
the latter quadratically in its amplitude. All gate and stopping guards
must still be closed.

\subsection{An explicitly instantiated improvement that is still not R1}

The following conservative variant of R0 is useful for locating the next
loss. Keep its width condition and every reference comparison/rebound cap,
but use
\[
 \beta=\frac{\delta}{32h\sqrt d},\quad
 v_* =\frac{\delta}{32h^{3/2}},\quad
 b_* =\frac{\delta}{32\sqrt{d+2}},\quad
 \eta_* =\varepsilon b_*,\quad r=\frac{32\varepsilon}{v_*},
\]
\[
 Z_* =\frac{32}{v_*b_*},\quad
 \nu=\frac1{512Z_*^2},\quad
 H=\frac12\log\frac{4d}{\nu s_0}.
\]
Here the minimum-coordinate event is deliberately retained only for this
fully proved variant. It has failure at most $\delta/16$; combine it with
Lemmas~\ref{sbquant:core}--\ref{sbquant:smallball} by a union bound.
Require also $D_*\le\varepsilon$,
$\Gamma H+\chi_*^2H/16\le\log2$, and set
\[
 \chi_* =\min\{1/32,\varepsilon\beta/8,\eta_*^2/1024,
                   \eta_*v_*/240,\varepsilon/16\},\qquad
 q_D=\frac{\chi_*\sqrt{s_0}}{16H\sqrt d}.
\]
Use the data caps of R0 with this $q_D$. The all-positive input probe
gates follow from $ra_i\ge16\varepsilon>|B_{2i}|$.
For output angles use
$\|D_{\cdot i}\|\le\varepsilon+\sqrt{s_0}v_i+D_*$,
not the global norm. Together with $a_i,c_i\ge3Av_i/4$ and the
noisy perturbation cap, these give the stated $r/2$ cones.
All output comparison estimates in the preceding proof remain valid
because $r\ge8\sqrt{s_0}$. The original lag statement is not needed for
this variant. Its four-snapshot, mass, and sector conclusions are the
same as R0; the original theorem and its lag corollary are unchanged.

For example the following values, rounded down for the data allowances,
pass this entire modified ledger:
\[
\begin{array}{c|c@{\qquad}c|c}
 d&2&\delta&0.1\\
 N&2000&h&10^8\\
 \varepsilon&10^{-77}&\mu_1&1\\
 \mu_2&6\cdot10^{-122}&\sigma_1=\sigma_2&10^{-243}\\
 a_D&10^{-242}&H&215.574142\text{ (upper bound)}\\
 \eta_*&1.5625\cdot10^{-80}&b_{\rm reb}&3.0517578125\cdot10^{-161}
\end{array}
\]
One may certify $\Delta_{\rm swing}\ge3.6\cdot10^{-278}$ radians.
The ratios are $\mu_2/\mu_1=6\cdot10^{-122}$,
$\sigma/\mu_1=10^{-243}$, and $\sigma/\varepsilon=10^{-166}$.
This is explicitly \emph{not} a useful noisy R1 regime. It shows that
removing the generic exponential and repairing anti-concentration alone
do not repair the endpoint proof. In this variant the binding initialization
loss is the scalar-reference output cap, not the data exponential.

\subsection{The inactive background has an adverse sign}

For the old all-active probe let $g=-G_1$. In dimension two,
$C=\|B_2\|^2-\eta^2\ge0$, and $sG_2\ge0$ because a stationary neuron
with negative first coordinate can activate only with transverse sign $s$.
The approximate endpoint secant therefore has the following \emph{upper}
bound, not merely a failed lower bound.

\begin{proposition}[Signed limitation of the all-active scalar endpoints]
\label{sbquant:signed-obstruction}
Suppose $d=2$, $s_0\le1/256$, $r\le1/8$, $z=r/\eta\ge8$, and take
the two scalar endpoints of \eqref{sbsep:snapshot-times}. For the complete
all-active approximate output \eqref{sbsep:probe-reference},
\begin{equation}\label{sbquant:signed-upper}
 E_b^\dagger-E_a^\dagger
 \le-\frac78g\eta+3\eta^{3/2}\sqrt r.
\end{equation}
In particular these endpoints have \emph{negative} secant if
$g>4\sqrt{\eta r}$. This proposition concerns the stated scalar
approximation and endpoints; it is not an impossibility result for the
true network or other probe sectors.
\end{proposition}
\begin{proof}
At the first endpoint the full longitudinal residual is $G_1=-g$.
Its increment at the second is
$d_L=\eta(A_bk-\nu z)$, with $7\eta/8\le d_L\le\eta$.
Its squared-error secant is exactly $-gd_L+d_L^2/2$.
The transverse secant is
\[
 \eta^2(T_a-T_b)(k-kC-sG_2)
 +\tfrac12\eta^4(T_b^2-T_a^2)
 \le\eta^2T_a.
\]
The first endpoint identity gives
$T_a^2=A_ak(A_az+k)^2/[z(1-A_0^2)]\le4z$.
Combining these estimates and $z\ge8$ proves
\eqref{sbquant:signed-upper}.
\end{proof}

The scale of $g$ is not hypothetical. For isotropic initialization in any
$d\ge2$ and either sign, rotational moments imply
\[
 \mathbb E[t\mathbf1_{t<0}(rt+sky)_+]
 \le-\frac{k}{2\pi d}+\frac r{2d}
 \le-\frac1{4\pi d}\quad(0<r\le1/8).
\]
The first inequality compares with $r=0$ and uses the ReLU Lipschitz
bound; $\mathbb E|ty|=2/(\pi d)$. Each summand lies in $[-1/2,0]$.
Hoeffding and a union bound over the two signs consequently give
\begin{equation}\label{sbquant:background-size}
 g\ge\frac{s_0}{8\pi d}
 \quad\text{with failure at most }2e^{-h/(8\pi^2d^2)}.
\end{equation}
Thus the all-active scalar endpoint requires a comparison of
$s_0/d$ against $\sqrt{\eta r}$, even after retaining the complement
\emph{with its sign}. With $\eta\asymp\varepsilon$ and
$r\asymp\sqrt{s_0}$ this is a genuine $s_0\lesssim h^{-1/2}$ scale
constraint (up to the displayed constants). Increasing width at fixed
$s_0$ does not remove it. This is different from proving the full noisy
theorem impossible.

\subsection{What changes when only the core's probe gates are fixed}

For the same approximate parameters, put
$q_i=rAv_i+skB_{2i}$ and $\widehat w_i=(Av_i,D^\dagger_{\cdot i})$.
Without assuming all positive neurons are active, the exact forward
correction to \eqref{sbsep:probe-reference} is
\begin{equation}\label{sbquant:clipping-correction}
 \Gamma_{\rm clip}(A)=\sum_{i\in I_+}\widehat w_i(-q_i)_+.
\end{equation}
The core contributes zero when \eqref{sbquant:core-radius} holds.
The correction to the endpoint secant is exactly
\begin{align}\label{sbquant:clipping-secant}
 &\langle f_b^\dagger-\xi,\Gamma_{{\rm clip},b}\rangle
 -\langle f_a^\dagger-\xi,\Gamma_{{\rm clip},a}\rangle\\
 &\hspace{15mm}
 +\tfrac12(\|\Gamma_{{\rm clip},b}\|^2-
                    \|\Gamma_{{\rm clip},a}\|^2).\nonumber
\end{align}
Neither the core count nor its $15/16$ longitudinal share determines this
signed quantity. Dropping it is not a legitimate robust-core repair.

There is, however, an exact alternative mechanism illustrating why this
term should not automatically be treated as error. Consider a planar
reflection-symmetric initialization for which $Bv=0$. On the rank-one
dataset the exact solution is $a=c=Av$, $D=B$. Define
\[
 P=\sum_{I_+}v_i(B_{2i})_+.
\]
Reflection in the first coordinate gives $G_1=-A_0P$ at $\xi=e_2$,
and the full output at this probe is exactly
\begin{equation}\label{sbquant:coherent-output}
 f(e_2,A)=P(A-A_0)e_1+K e_2,
 \qquad K=\sum_i(B_{2i})_+^2.
\end{equation}
For $A_0\le1/4$, $A_a=\sqrt{3/4}$ and $A_b=\sqrt{9/10}$,
\begin{equation}\label{sbquant:coherent-rebound}
 E_b-E_a
 =\frac{P^2}{2}(A_b-A_a)(A_b+A_a-2A_0)
 \ge\frac{P^2}{20}>0.
\end{equation}
All positive neurons already carry mass $A^2\ge3/4$ in this window.
For isotropic angular averaging $P=\sqrt{s_0}/(2\pi)$ in dimension two,
so the gain has scale $s_0=h\varepsilon^2$, rather than $\varepsilon^2$.
This exact symmetric-reference calculation is \emph{not} yet a
high-probability noisy initialized theorem. The pure axis probe also
does not have strict initial descent by this calculation alone. Extending
to a positive-first-coordinate arc and transferring the reference to iid
initialization/noisy data are separate obligations.

The user has authorized pursuing the partly active probe mechanism with
its $h\varepsilon^2$ scale, while retaining R0. The next section develops
this route. It is not Scope A, O3, or a requirement that the weak family
be learned before the rebound.

\subsection{A rank-one comparison that permits marginal gate changes}

\begin{lemma}[Transverse-residual comparison without a minimum coordinate]
\label{sbquant:residual-comparison}
For the exact rank-one dataset, retain all initially positive neurons,
including those that subsequently reach zero. Put
\[
 R=Da,\quad u=\|a\|,\quad p=c-a,\quad
 z=\|R\|/u,\quad l=\|p\|/u,\quad W=z+l,
 \quad b_0=\|Bv\|.
\]
Suppose $s_0\le1$, $A^2\le1-\nu$ on $[0,H]$, and define
\[
 \overline W=b_0e^{(\sqrt{s_0}+0.02)H},\qquad
 q_*=2H\overline W.
\]
If $\overline W<\min\{0.01,\nu/16\}$ and $q_*<\nu/8$, then
through this interval
\begin{align}
 W&\le\overline W,&
 |\log(u/A)|&\le q_*,&
 \|a/u-v\|&\le q_*,\label{sbquant:residual-controls}\\
 \|a-Av\|&\le3Aq_*,&
 \|c-Av\|&\le A(3q_*+2\overline W),&
 \|D-B\|_F&\le H\overline W.\nonumber
\end{align}
No positive minimum of the initial coordinates is assumed. The comparison
retains every neuron and permits zero strong training coordinates.
\end{lemma}
\begin{proof}
An initially negative first coordinate is stationary on the rank-one
dataset. An initially positive one cannot become negative: the derivative
of its negative part is zero almost everywhere on the negative half-line.
Thus $a\ge0$ on $I_+$, even if some coordinates reach zero. Let
$\mathcal A$ be the diagonal compatible selector. Almost everywhere
$\mathcal A a=a$, and the exact system is
\[
 a'=\mathcal A(kc-D^\top R),\quad c'=ka,\quad D'=-Ra^\top,
 \qquad k=1-a^\top c.
\]
The Frobenius norm of $D$ is nonincreasing. With
$\theta=a^\top p/u^2$, direct differentiation gives
\begin{align*}
 g:=u'/u&=k(1+\theta)-z^2,\\
 p'&=-k\mathcal A p+\mathcal A D^\top R,\\
 R'&=(k-u^2)R+kD\mathcal A p-D\mathcal A D^\top R.
\end{align*}
While $0\le k\le1$ and $g\ge0$, the dissipative terms give
\[
 l'\le-gl+\sqrt{s_0}z,\qquad
 z'\le(-u^2-k\theta+z^2)z+k\sqrt{s_0}l.
\]
Consequently $W'\le\sqrt{s_0}W+W^2+W^3$ and $W(0)=b_0$.
For $W\le0.01$, scalar comparison gives $W\le\overline W$.
Moreover
\[
 g=1-u^2+(1-2u^2)\theta-u^2\theta^2-z^2.
\]
On $u^2\le1$, the error after $1-u^2$ has magnitude at most
$W+2W^2\le2\overline W$.
The restoring term in $(\log(u/A))'$ therefore gives the stated bound.
Also
$\|(a/u)'\|\le l+\sqrt{s_0}z\le2\overline W$.
The estimates on $a,c$ follow, and integrating
$\|D'\|_F=u^2z$ proves the last bound.

For completeness stop also at $u^2=1-\nu/2$, $k=0$, or $g=0$.
The logarithmic bound excludes the first face because
$(1-\nu)e^{\nu/4}<1-\nu/2$.
Then $k=1-u^2(1+\theta)\ge\nu/2-\overline W>\nu/4$,
and $g\ge k(1-\overline W)-\overline W^2>0$.
The strict caps exclude every stopping face. At a zero coordinate there
is no stopping condition and no division by that coordinate.
\end{proof}

\begin{proposition}[Full-network coherent endpoint with explicit charges]
\label{sbquant:coherent-finite}
In dimension two, under Lemma~\ref{sbquant:residual-comparison}, set
\[
 P=\sum_{I_+}v_i(B_{2i})_+,\quad
 C_0=\varepsilon^2\sum_i t_i(y_i)_+,\quad
 K=\varepsilon^2\sum_i(y_i)_+^2,\quad
 B_+=\left(\sum_{I_+}(B_{2i})_+^2\right)^{1/2}.
\]
Let $A_a=\sqrt{3/4}$ and $A_b=\sqrt{9/10}$, with $A_0<A_a$.
At $\xi=e_2$, define
\[
 L_j=P(A_j-A_0)+C_0,\quad
 e_{1,j}=A_j(3q_*+2\overline W)B_+,\quad
 e_2=H\overline W B_+.
\]
The exact rank-one network satisfies
\begin{align}\label{sbquant:coherent-finite-gap}
 E_b-E_a\ \ge\ &\frac{P^2}{2}(A_b-A_a)(A_b+A_a-2A_0)
                  -|C_0|P(A_b-A_a)\\
 &-\sum_{j=a,b}\left(|L_j|e_{1,j}+\frac{e_{1,j}^2}{2}
                   +|K-1|e_2+\frac{e_2^2}{2}\right).\nonumber
\end{align}
This includes all inactive neurons and every changing strong gate.
\end{proposition}
\begin{proof}
The input transverse coordinates are constant, so probe activations at
$e_2$ are exactly $(B_{2i})_+$, even when strong training gates change.
The comparison output $(L_j,K)$ follows from the definitions; the inactive
first-coordinate contribution is $C_0-A_0P$.
The two coordinate errors are bounded separately by
Lemma~\ref{sbquant:residual-comparison}.
Expand the two squared residuals coordinate by coordinate. In particular
the first-coordinate comparison pays $|L_j|e_{1,j}$, rather than replacing
its small residual by the full transverse residual of order one.
\end{proof}

\section{A noisy coherent-endpoint theorem at fixed aggregate initialization}
\label{sbwide:scope}

This continuation uses probes inside the strong family's angular spread.
It preserves the preceding fluctuation-scale theorem. Its width is still
very large and its allowed data noise is small; those limitations are
reported numerically, rather than hidden in an asymptotic existence claim.
All statements in this section concern $d=2$ and equal cluster weights.
Time $\tau=t/2$ is used in the dynamical estimates.

\subsection{A gate-robust noisy comparison}

Write $a,c$ for the full first-coordinate input and output vectors and
$B,D$ for the second-coordinate rows. Set
\begin{gather*}
 a_+=(a_i)_+,\qquad u=\|a_+\|,\qquad p=c-a,\qquad
 m=c^\top a_+,\qquad k=1-m,\\
 R=Da_+,\qquad z=|R|/u,\qquad l=\|p\|/u,\qquad W=z+l.
\end{gather*}
Thus the actual output at $e_1$ is $(m,R)$, without any block reduction.

\begin{lemma}[Axis comparison with actual noisy training selectors]
\label{sbwide:axis-system}
Suppose every strong sample is within $a_D\le10^{-3}$ of $e_1$, every
weak sample within $a_D$ of $\mu_2e_2$, and $0<\mu_2\le10^{-2}$.
On a stopped interval with $S\le2$, $0\le k\le1$, and $|R|\le1$, put
\[
 P=30\{a_D+(\mu_2+a_D)^2\}.
\]
There is a diagonal matrix $\mathcal A\in[0,1]^{h\times h}$ such that
almost everywhere
\begin{align}\label{sbwide:perturbed-axis}
 a'&=ka_++k\mathcal A p-\mathcal A D^\top R+E_a,
 &c'&=ka_++E_c,\\
 D'&=-Ra_+^\top+E_D,
 &B'&=E_B,\nonumber
\end{align}
where every error has Frobenius norm at most $P$ and every neuronwise
error has norm at most $P\|u_i\|$. Furthermore
\[
 \|(\mathcal A-I)a_+\|\le a_D\sqrt S.
\]
No common strong-sample mask is assumed. The weak selectors are arbitrary.
\end{lemma}
\begin{proof}
Let $\mathcal A_{ii}$ be the empirical strong-cluster mean of its actual
selector. Write $f(e_1)=(m,R)$. ReLU Lipschitzness gives
\[
 \|f(x)-f(e_1)\|\le Sa_D,\qquad
 |(u_i^\top x)_+-(a_i)_+|\le a_D\|u_i\|.
\]
The target-minus-output residual changes by at most $(1+S)a_D$.
At $e_1$ its norm is at most $\sqrt2$ on the stopped interval.
The output-gradient error on the strong cluster is therefore at most
$[3(1+a_D)+\sqrt2]a_D\|u_i\|<5a_D\|u_i\|$.
The first input-gradient error before replacing $\mathcal A a$ by $a_+$
has the same bound. A changed gate implies $|a_i|\le a_D\|u_i\|$;
thus $|\mathcal A_{ii}a_i-(a_i)_+|\le a_D\|u_i\|$.
The transverse input-gradient norm is at most
$(\sqrt2+3a_D)a_D\|u_i\|$.
On a weak sample, $\|x-f(x)\|\le3\|x\|$, so each input or output
gradient norm is at most $3(\mu_2+a_D)^2\|u_i\|$.
Full balance permits the same bound for input and output weights.
Combining these estimates gives each error at most
$9\{a_D+(\mu_2+a_D)^2\}\|u_i\|$.
Summing squares and using $\sqrt S\le\sqrt2$ proves both asserted
bounds with the stated reserve. The selector estimate follows from the
same changed-gate implication.
\end{proof}

Define the scalar logistic solution $A'=A(1-A^2)$ with
$A_0=\|(a(0))_+\|$, and put $v=(a(0))_+/A_0$.
The coordinates of $v$ outside $I_+$ are zero. All initial neurons have
norm $\varepsilon$, and $s_0=h\varepsilon^2$.

\begin{lemma}[Relative noisy bounds without a minimum coordinate]
\label{sbwide:relative}
Let $A_{0,-}\le A_0$, and let $A^2\le1-\nu$ through a terminal time
$T\le H$. Define
\begin{align*}
 \overline B&=\sqrt{s_0}+PH,\qquad b_0=\|B(0)v\|,\\
 \overline W&=(b_0+10PH/A_{0,-})e^{(\overline B+0.02)H},\\
 q_*&=H(2\overline W+3P/A_{0,-}),\\
 C_c&=3q_*+(H+1)\overline W+PH.
\end{align*}
Suppose $\overline B\le0.3$, $\overline W<\min\{0.01,\nu/16\}$,
$q_*<\nu/8$, $3P/A_{0,-}<\nu/16$, and
\begin{equation}\label{sbwide:mass-stop}
 1-\nu/2+
 \{\sqrt{s_0}+H[(1+\overline B)\overline W+P]\}^2
 +\overline B^2<2.
\end{equation}
Then, up to $T$,
\begin{align}\label{sbwide:relative-bounds}
 W&\le\overline W,& |\log(u/A)|&\le q_*,&
 \|a_+/u-v\|&\le q_*,\\
 \|a_+-Av\|&\le3Aq_*,&
 \|c-(Av+a(0)_-)\|&\le C_c,&
 \|D-B(0)\|&\le H(\overline W+P),\nonumber\\
 \|B-B(0)\|&\le PH,&\|D\|&\le\overline B,& S&<2.\nonumber
\end{align}
Here $a(0)_-=(\min(a_i(0),0))_i$. The same bounds hold with $b_0$
replaced by any deterministic upper bound.
\end{lemma}
\begin{proof}
Let $\mathcal G$ denote the diagonal selector for $a_+$, so
$\mathcal G a_+=a_+$. The chain rule for absolutely continuous ReLU
coordinates is used almost everywhere. Equations
\eqref{sbwide:perturbed-axis} imply
\[
 p'=-k\mathcal A p+\mathcal A D^\top R+E_c-E_a,
\]
\[
 R'=(k-u^2)R+kD\mathcal G\mathcal A p
 -D\mathcal G\mathcal A D^\top R+E_Da_++D\mathcal G E_a.
\]
The last quadratic term is dissipative. With
$\theta=a_+^\top p/u^2$ and $g=u'/u$, the selector bound gives
\[
 g=k(1+\theta)-z^2+E_g,
 \qquad |E_g|\le P(1+l+z)/u.
\]
In particular
\[
 g=1-u^2+(1-2u^2)\theta-u^2\theta^2-z^2+E_g.
\]
While $g\ge0$, $u\ge A_{0,-}/2$, $u^2\le1-\nu/2$, and
$W\le0.01$, the same dissipative calculation as in
Lemma~\ref{sbquant:residual-comparison} gives
\[
 W'\le(\overline B+0.02)W+10P/A_{0,-}.
\]
The initial value is $b_0$. The forcing after $1-u^2$ in $g$ is at most
$2\overline W+3P/A_{0,-}$ in absolute value. The restoring logarithmic
comparison and the projected directional equation give the bounds for
$\log(u/A)$ and $a_+/u$.
The bounds exclude $u^2=1-\nu/2$, $u=A_{0,-}/2$, $k=0$, and $g=0$
exactly as in the preceding rank-one proof; at the last face use the
additional $3P/A_{0,-}$ reserve.

The contraction in $D'=-Ra_+^\top+E_D$ gives
$\|D\|\le\sqrt{s_0}+PH$, and integration gives the two displacement
bounds. For the full $c$ comparison, on $c_i<0$ one has
$(a_i)_+\le|p_i|$. Therefore
$\|c_- -a(0)_-\|\le H(\overline W+P)$, while
$\|c_+-a_+\|\le\|p\|\le\overline W$.
This proves the asserted $C_c$ bound without an extra factor of $H$.
On $a_i<0$, its derivative in \eqref{sbwide:perturbed-axis} has no
$ka_+$ term. Hence
\[
 \|a_-\|\le\sqrt{s_0}+H[(1+\overline B)\overline W+P].
\]
Together with $\|B\|\le\overline B$, condition
\eqref{sbwide:mass-stop} excludes the final norm stopping face.
All estimates include neurons whose coordinate signs change.
\end{proof}

\subsection{A fixed initialization cohort at learned mass}

In the planar initialization write $\widehat u_i=(t_i,y_i)$ and take
$I_C=\{t_i\ge|y_i|\}$. Its angular aperture is $\pi/2$.
Define the following deterministic charges from Lemma~\ref{sbwide:relative}:
\begin{align*}
 M_*&=4q_*+2\overline W,& V_*&=H(M_*+\overline W+P),\\
 J_D&=e^{V_*}H(\overline W+P)/A_{0,-},&
 J_B&=e^{V_*}HP/A_{0,-},\\
 J_p&=H(1+J_D)\overline W+2e^{V_*}HP/A_{0,-},\\
 \Gamma_C&=e^{M_*H}-1+
 \sqrt2e^{2M_*H}H\{J_p+(1+J_D)\overline W+Pe^{V_*}/A_{0,-}\}.
\end{align*}

\begin{lemma}[Core holding from per-neuron estimates]
\label{sbwide:core-holding}
Suppose the preceding estimates hold, $\Gamma_C<1/4$, and
\[
 \frac{1-\Gamma_C}{\sqrt2}>a_De^{V_*}.
\]
Then all strong noisy training gates of $I_C$ stay active and
\[
 a_i\ge(1-\Gamma_C)A\varepsilon t_i/A_0,\qquad
 |p_i|\le\varepsilon J_p,\qquad
 |B_i|\le\varepsilon(|y_i|+J_B),\qquad
 |D_i|\le\varepsilon(|y_i|+J_D).
\]
If $\sum_{I_C}v_i^2\ge\omega_C$, its longitudinal learned mass is at least
\begin{equation}\label{sbwide:core-mass}
 A^2\omega_C(1-\Gamma_C)^2-\sqrt{s_0}J_p.
\end{equation}
On $A\ge A_{\rm spec}$, its input and output tangent-angle bounds are
respectively
\[
 \frac{A_0}{A_{\rm spec}}
 \frac{1+\sqrt2J_B}{1-\Gamma_C},\qquad
 \frac{A_0}{A_{\rm spec}}
 \frac{1+\sqrt2J_D}
 {1-\Gamma_C-\sqrt2 A_0J_p/A_{\rm spec}},
\]
provided the last denominator is positive.
\end{lemma}
\begin{proof}
The neuronwise residual matrix has norm at most
$k+|R|+P$ on the stopped interval: the strong residual contributes
$k+|R|$, and the same data estimates as in
Lemma~\ref{sbwide:axis-system} absorb the remaining terms.
Full balance thus gives
\[
 \|u_i\|=\|w_i\|\le\varepsilon\frac A{A_0}e^{V_*}.
\]
Here $|m-A^2|\le M_*$ follows from \eqref{sbwide:relative-bounds},
and $\int(1-A^2)=\log(A/A_0)$.
Integrating the per-neuron $B,D,p$ equations gives $J_B,J_D,J_p$.
Stop the core only when a strong noisy gate first ceases to be active.
Before that time,
$a_i'=ka_i+kp_i-D_iR+E_{a,i}$.
The integrating factor for $k$ is between
$(A/A_0)e^{-M_*H}$ and $(A/A_0)e^{M_*H}$.
Variation of constants and $t_i\ge1/\sqrt2$ give the displayed
relative error $\Gamma_C$.
The gate reserve follows because its preactivation is at least
$a_i-a_D\|u_i\|$; the stated inequality makes it strictly positive.
For the mass bound use $c_i=a_i+p_i$ and
$\sum_{I_C}|a_i|\le\sqrt h\|a_+\|\le\sqrt h$.
Finally $|y_i|\le t_i$ on this cohort proves the two angle bounds.
\end{proof}

\subsection{One common, explicitly checked parameter regime}

\begin{theorem}[Noisy planar many-neuron coherent rebound]
\label{sbwide:theorem}
Let $h\ge10^{10}$, initialize independently and uniformly on the planar
unit circle with balanced weights of norm
$\varepsilon=\sqrt{0.04/h}$, and train on $N=2000$ independent samples,
half from each of the two isotropic planar Gaussian clusters with
\[
 \mu_1=1,\qquad \mu_2=10^{-5},\qquad
 \sigma_1=\sigma_2=10^{-11}.
\]
With probability at least $0.9$, there is a fixed initialization cohort
$I_C$ of at least $0.24h$ neurons and times
\[
 0<t_{\rm early}<t_{\rm spec}<t_a<t_b<7
\]
with the following properties.

Throughout $[t_{\rm spec},t_b]$, every cohort member's input and output
angle has absolute value at most $0.13$ radians, and
\[
 \sum_{i\in I_C}u_{i1}w_{i1}\ge0.55.
\]
The whole-network mass is at most $2$ through $t_b$, and
$|f_2(e_2,t)|<0.011$ there; in particular the weak concept is not learned
to order-one fitting fraction.

On the deterministic probe arc
\[
 \xi_\theta=(\sin\theta,\cos\theta),\qquad
 |\theta-\arcsin(2\cdot10^{-5})|\le5\cdot10^{-6},
\]
the initial error derivative is at most $-10^{-8}$ almost everywhere
for $0\le t\le10^{-8}$. Thus $t_{\rm early}=5\cdot10^{-9}$ supplies
a uniform drop $a_{\rm dec}=5\cdot10^{-17}$.
At the center of this arc,
\[
 E(\xi,t_b)-E(\xi,t_a)\ge3\cdot10^{-5}.
\]
On an open arc of angular width $2\cdot10^{-18}$ centered there,
the initial drop is at least $2.5\cdot10^{-17}$ and the
post-specialization rebound at least $1.5\cdot10^{-5}$.
Every such error has an interior minimum on $[0,t_b]$; this does not
assert that its minimum occurs after specialization.
\end{theorem}

\begin{proof}
\emph{Common initialization and sample event.}
Set $s_0=0.04$ and use the following planar empirical events:
\begin{align*}
 &Q/h\in[0.249,0.251],\qquad
 h^{-1}\sum(t_i)_+(y_i)_+\in[0.077,0.082],\\
 &\left|h^{-1}\sum t_i(y_i)_+\right|\le10^{-4},\qquad
 h^{-1}\sum(y_i)_+^2\in[0.249,0.251],\\
 &|I_C|\ge0.24h,\qquad h^{-1}\sum_{I_C}t_i^2\ge0.202,
 \qquad |B(0)v|\le2\varepsilon.
\end{align*}
The respective population means are $1/4$, $1/(4\pi)$, zero, $1/4$,
$1/4$, and $1/8+1/(4\pi)$ for the first six quantities.
For the first six events the scalar ranges are at most $2$ and their
mean-to-threshold distances are at least $10^{-4}$. Hoeffding bounds
their total failure by $12e^{-h\,10^{-8}/2}<0.001$ at the stated
widths. For the last event use the same iid scalar sum as in
Lemma~\ref{sbquant:smallball}, now for an upper tail. On $Q\ge0.249h$,
the threshold for its variance-normalized sum is
$8\sqrt{0.249}>3.99$. The normal two-sided tail plus two
Berry--Esseen errors is below $0.001$. No independence between these
events is assumed. The cohort's normalized longitudinal share exceeds
$0.202/0.251>0.80$.

The sample event $\max\|x^{(k)}-\mu_ke_k\|\le a_D=10^{-10}$ has failure
at most $2000e^{-50}<0.001$, by the planar Gaussian radial tail.
The additional initialization event for initial descent below costs
$0.005$. The union of all failures is less than $0.1$.

\emph{Dynamical ledger and learned mass.}
Take $H=3.5$, $\nu=0.1$,
$A_{0,-}=\sqrt{0.249s_0}$ and $A_{0,+}=\sqrt{0.251s_0}$.
Define the reached times by $A^2=0.7,0.75,0.9$, respectively;
the last normalized time is at most
$\frac12\log\{9(1-A_{0,-}^2)/A_{0,-}^2\}<3.5$.
All subsequent bounds improve as $h$ increases, so it suffices to use
$h=10^{10}$ and $b_0\le2\varepsilon=4\cdot10^{-6}$ in the ledger.
Direct substitution in the preceding lemmas gives the safe upper bounds
\[
\begin{array}{c|c@{\qquad}c|c}
 P&6.001\cdot10^{-9}&\overline W&1.319\cdot10^{-5}\\
 q_*&9.30\cdot10^{-5}&C_c&3.382\cdot10^{-4}\\
 M_*&3.985\cdot10^{-4}&V_*&1.442\cdot10^{-3}\\
 J_D&4.64\cdot10^{-4}&J_B&2.11\cdot10^{-7}\\
 J_p&4.67\cdot10^{-5}&\Gamma_C&1.695\cdot10^{-3}.
\end{array}
\]
Every stopping reserve is strict. The left side of
\eqref{sbwide:mass-stop} is less than $1.031$.
Equation~\eqref{sbwide:core-mass} is greater than $0.558$, and the
larger angle bound is less than $0.120$ radians. These imply the stated
mass and cone with reserve, throughout the whole learned window.

\emph{Endpoint, including every neuron.}
Let $P_\xi=\sum_{I_+}v_i(B_i(0))_+$,
$C_0=\varepsilon^2\sum_i t_i(y_i)_+$, and
$K=\varepsilon^2\sum_i(y_i)_+^2$.
The comparison output at $e_2$ is
\[
 (L_j,K),\qquad L_j=P_\xi(A_j-A_0)+C_0.
\]
The full actual coordinate errors at either snapshot are bounded by
\[
 e_1=C_c\sqrt{K_{\max}}+\sqrt2PH<3.392\cdot10^{-5},
 \quad
 e_2=H(\overline W+P)\sqrt{K_{\max}}+\overline BPH
       <4.64\cdot10^{-6},
\]
where $K_{\max}=0.251s_0$. These bounds use the actual input row
$B(\tau)$ and its Lipschitz ReLU change, not fixed probe gates.
The moment event gives
\[
 P_{\xi,-}=\sqrt{s_0}\frac{0.077}{\sqrt{0.251}},\qquad
 P_{\xi,+}=\sqrt{s_0}\frac{0.082}{\sqrt{0.249}},\qquad
 |C_0|\le4\cdot10^{-6}.
\]
Expanding coordinatewise as in
Proposition~\ref{sbquant:coherent-finite}, the uncharged endpoint gain
exceeds $6.302\cdot10^{-5}$, and its complete coordinate-error charge
is less than $1.109\cdot10^{-5}$. Hence the actual secant at $e_2$
exceeds $5.193\cdot10^{-5}$.
The same output estimate at every time, with $A\le\sqrt{0.9}$,
gives $|f_2(e_2,t)|\le K_{\max}+e_2<0.011$.

At either endpoint and $0\le\theta\le2.5\cdot10^{-5}$,
\[
 |f_1(\xi_\theta)-\sin\theta|<0.029,
 \qquad |f_2(\xi_\theta)-\cos\theta|<1.011.
\]
The spatial derivatives of these residual coordinates have absolute
bounds $3$ and $\overline B\sqrt2+2.5\cdot10^{-5}<0.283$, respectively.
Thus $|\partial_\theta E|<0.4$ almost everywhere. Spatial absolute
continuity transfers the two-snapshot gap to the prescribed center,
losing at most $0.8\arcsin(2\cdot10^{-5})<1.601\cdot10^{-5}$.
The remaining gap exceeds $3\cdot10^{-5}$.

\emph{Uniform initial decrease on the same arc.}
Put $m_0=s_0/4=0.01$ and $\gamma=0.99$.
For a deterministic unit direction $v$, the scalar
$Y_j=\widehat u_j(\widehat u^\top v)_+$ has mean $v_j/4$,
variance $(1+v_j^2)/16\le1/8$, and centered absolute bound $2$.
Apply two-sided Bernstein to both coordinates at 11 equally spaced
angles on the probe arc and at $e_1,e_2$. With
$\ell=\log(10400)$, failure is at most $0.005$ and the vector error
at these points is at most
\[
 \sqrt2s_0\left\{\sqrt{\frac\ell{4h}}+
                         \frac{4\ell}{3h}\right\}<8.61\cdot10^{-7}.
\]
The residual map $f_0(v)-m_0v$ is $s_0+m_0=0.05$ Lipschitz.
Transport from the grid therefore bounds its error by $9\cdot10^{-7}$
on the full probe arc and the normalized strong samples. For the
normalized weak samples the bound is $1.9\cdot10^{-6}$, since their
distance from $e_2$ is at most $2a_D/\mu_2$.

The mass bound $S(t)\le s_0e^{2\lambda t}$ with $\lambda\le1$
gives $S\le0.04000001$ on $0\le t\le10^{-8}$.
Each of the two parameter displacements has norm at most $0.21t$;
hence $\|W-U\|_F\le0.42t$, $\|f_t(v)-f_0(v)\|\le0.1t$ for unit $v$,
and each ReLU-product sum below changes by at most $0.1t$.
For unit $\xi,v$ write
\[
 c=\xi^\top v,\qquad
 L=\sum_i(u_i^\top\xi)_+(u_i^\top v)_+.
\]
Replacing $W$ by $U$ only in the second term of the exact instantaneous
kernel gives, if both residual errors relative to $\gamma I$ are at most
$\delta_f$,
\begin{align*}
 \langle\xi-f_t(\xi),K_t(\xi,v)(v-f_t(v))\rangle
 \ge\ &2\gamma^2cL-(L+cS)(2\gamma\delta_f+\delta_f^2)\\
 &-c(\|W\|_F+\|U\|_F)\|W-U\|_F(\gamma+\delta_f)^2.
\end{align*}
This identity-based bound uses the factor $c$ on the second kernel term;
it does not pay a global $S^2$ residual cost.
For strong samples, $c\ge1.49\cdot10^{-5}$, $L\ge0.003$, and
$\delta_f\le9.02\cdot10^{-7}$ throughout this time interval. Substitution
gives a strong-cluster contribution, including its weight $1/2$ and
sample norm squared, greater than $4\cdot10^{-8}$.
For weak samples $c>0.999$, $L>0.0095$, and
$\delta_f<2\cdot10^{-6}$, so their contribution is nonnegative.
The claimed bound $\dot E\le-10^{-8}$ follows with reserve, uniformly
over the arc. Gate changes are allowed in this argument.

\emph{Time ordering and final sector.}
Since $A_0^2\le0.01004$, the specialization time is greater than one
in physical time and precedes both endpoint times strictly.
Integrating the initial derivative gives the stated $a_{\rm dec}$.
Use the four-snapshot argument with the proven mass bound $2$.
Its chord radius is at least $a_{\rm dec}/36$ and its angular width
exceeds $2.7\cdot10^{-18}$. The smaller reported arc stays inside the
initial probe arc. All these probes have distance greater than $0.9$
from every training sample on the same sample event.
\end{proof}

At the smallest certified width, the concrete primitive tuple is
\[
 (d,\delta,N,h,\varepsilon,\mu_1,\mu_2,\sigma,H)
 =(2,0.1,2000,10^{10},2\cdot10^{-6},1,10^{-5},10^{-11},3.5).
\]
Here $H$ is normalized time; physical time is less than $7$.
The ratios are $\mu_2/\mu_1=10^{-5}$,
$\sigma/\mu_1=10^{-11}$, and $\sigma/\varepsilon=5\cdot10^{-6}$.
The rebound is certified at $3\cdot10^{-5}$ at the chosen center.
The theorem permits fixed means and fixed noise as $h$ increases with
$h\varepsilon^2=0.04$; it does not assert a fixed-width limit
$\varepsilon\downarrow0$ at those same data parameters.

\section{Sharper descent, a common sector, and a ten-million-neuron witness}
\label{sbsharp:scope}

This section preserves R0 and Theorem~\ref{sbwide:theorem}. It strengthens
that theorem's initial descent and common sector, and supplies a second
primitive regime with smaller width and larger noise. The learned object
is one dominant family. Neither result is a two-learned-family theorem.
A signed derivative estimate also locates every global minimum on the
certified interval \emph{before} the declared specialization time; no
claim that specialization causes the minimum is made.

\subsection{Optimizing the actual instantaneous-kernel certificate}

All times in this subsection are physical. Write $s=h\varepsilon^2$,
$\gamma=1-s/4$, and suppose the empirical covariance and initial loss give
$\lambda\le\lambda_0$. Then the installed norm envelope implies
\begin{gather*}
 S(t)\le s e^{2\lambda_0t},\quad
 \|U(t)-U(0)\|_F,\|W(t)-W(0)\|_F
       \le\sqrt{s}(e^{\lambda_0t}-1),\\
 F(t):=s(e^{2\lambda_0t}-1),\qquad
 D_*(t):=4s e^{\lambda_0t}(e^{\lambda_0t}-1).
\end{gather*}
Here $F$ bounds both the change of the output on any unit probe and the
change of any ReLU-product sum $L$. The second kernel term's
$W$-to-$U$ replacement costs at most $D_*$ before its factor
$\xi^\top v$. These estimates follow by integrating
$\|\dot U\|_F,\|\dot W\|_F\le\lambda_0\sqrt S$ and applying
ReLU Lipschitzness twice; they require no fixed gates.

Let the deterministic probe arc be
$\xi_\theta=(\sin\theta,\cos\theta)$, $\alpha\le\theta\le\beta$.
Suppose its initial radial error, and that on normalized strong samples,
is at most $d_0$. Suppose the initial strong ReLU-product sum is at least
$L_0$, and $\xi^\top v\ge c_0>0$ on that cluster. Define
\[
 d(t)=d_0+F(t),\quad L_-(t)=L_0-F(t),\quad S_+(t)=se^{2\lambda_0t}.
\]
\begin{align*}
 \mathcal J(t)&=2\gamma^2c_0 L_-(t)
 -(L_-(t)+c_0S_+(t))(2\gamma d(t)+d(t)^2)\\
 &\hspace{2em}-c_0D_*(t)(\gamma+d(t))^2.
\end{align*}

\begin{lemma}[Explicit maximal interval for a specified kernel budget]
\label{sbsharp:descent}
Suppose the corresponding weak-cluster budget is nonnegative.
While $L_->0$, $2\gamma^2c_0>2\gamma d+d^2$, and
$2\gamma^2L_--S_+(2\gamma d+d^2)-D_*(\gamma+d)^2>0$,
\[
 \dot E(\xi_\theta,t)\le-\frac{(1-a_D)^2}{2}\mathcal J(t).
\]
The function $\mathcal J$ is strictly decreasing until its first zero
$T_0$ in this regime. For every $0<T<T_0$ the uniform descent rate and
integrated drop are
\[
 c(T)=\frac{(1-a_D)^2}{2}\mathcal J(T)>0,\qquad
 a(T)=\frac{(1-a_D)^2}{2}\int_0^T\mathcal J(t)\,dt>0.
\]
Thus $T_0$ is the supremum of the strictly negative-derivative intervals
certified by \emph{this budget}, not a claimed optimal interval for the
actual flow. The integral is an explicit polynomial antiderivative in
$x=e^{\lambda_0t}$: integrate its constant term as $T$ and each $x^j$
term as $(e^{j\lambda_0T}-1)/(j\lambda_0)$.
\end{lemma}
\begin{proof}
Use the exact kernel bound in Theorem~\ref{sbwide:theorem} with the
new displacement estimates. The two displayed positivity conditions
make that bound increasing in both $c$ and $L$, permitting their lower
bounds to be substituted without replacing a negative term incorrectly.
Multiply by the strong-cluster weight and its squared sample norm.
For monotonicity, write the first two terms as
$L_-(2\gamma^2c_0-2\gamma d-d^2)-c_0S_+(2\gamma d+d^2)$.
The first product decreases while its factors are positive, and both
subtracted terms increase. Integration proves the claims.
\end{proof}

For the original parameters of Theorem~\ref{sbwide:theorem}, use
the radial-error refinement of $\lambda_0$ below, the slightly larger arc
$[1.5\cdot10^{-5},2.5\cdot10^{-5}]$, the same eleven-point net, and
\begin{gather*}
 d_0=\sqrt2s\left\{\sqrt{\frac{\log10400}{4h}}
             +\frac{4\log10400}{3h}\right\}
             +0.05(5\cdot10^{-7}),\\
 c_0=\sin(1.5\cdot10^{-5})-2a_D,\qquad
 L_0=0.077s\cos(2.5\cdot10^{-5})-2a_Ds.
\end{gather*}
Let $d_{\rm train}$ be the net bound plus $0.05(2a_D/\mu_2)$.
The radial initialization event and loss monotonicity sharpen the rate to
\[
 \lambda_0=(\gamma+d_{\rm train})
        \frac{(1+a_D)^2+(\mu_2+a_D)^2}{2},
\]
since $\lambda\le\sqrt{2\mathcal L(0)}
\sqrt{\mathbb E_n\|x\|^2}$.
At $h=10^{10}$, $\lambda_0<0.495001$ and $T_0$ lies between
$0.00035080$ and $0.00035081$.
At $T=0.00035$ the rate is greater than $9\cdot10^{-11}$ and
\begin{equation}\label{sbsharp:old-drop}
 E(\xi,0)-E(\xi,0.00035)>7.45\cdot10^{-12}.
\end{equation}
The weak budget is positive throughout. These bounds improve with $h$.
The stronger integrated estimate, rather than $Tc(T)$, is used here.
For reference the constant-rate-product optimum is approximately
$0.00017522$, determined by $\mathcal J+T\mathcal J'=0$.

There is also no need to transfer this small gap spatially. The initial
drop is already uniform on the whole arc. The old axis secant exceeds
$5.193\cdot10^{-5}$; its spatial loss on the whole arc is at most
$0.8(2.5\cdot10^{-5})$. Consequently the later secant exceeds
$3\cdot10^{-5}$ \emph{on that same entire arc}. Its open interior has
width $10^{-5}$ radians. This replaces the artificial $2\cdot10^{-18}$
sector restriction without spending any initial-drop margin.

\subsection{A weighted residual comparison}

Return to normalized time $\tau=t/2$ and the notation of
Lemma~\ref{sbwide:relative}. The following refinement is important at
smaller width: the initial residual grows with $\varepsilon$, so simply
relaxing concentration intervals is insufficient.
Set $q_c=10^{-3}$, $\beta_B=\sqrt{s}+PH$, and suppose $A^2\le1-\nu$.
Define
\begin{align*}
 I_*&=\operatorname{arctanh}\sqrt{1-\nu}
                 -\operatorname{arctanh}A_{0,-},\\
 J_*&=\tfrac12\log\frac{1-A_{0,-}^2}{\nu},\\
 Z_*&=\{b_0+PH(1+6\beta_B/A_{0,-})\}
                  e^{(\beta_B^2+0.01)H},\\
 l_*&=\beta_B Z_*e^{2q_c}\frac{I_*}{\sqrt{1-\nu}}
               +\frac{2PH e^{q_c}}{A_{0,-}},\qquad W_*=Z_*+l_*,\\
 q_*&=H\max\{l_*+l_*^2+Z_*^2+2P/A_{0,-},\,
                     l_*+\beta_BZ_*+2P/A_{0,-}\},\\
 C_c^*&=3q_*+(H+1)l_*+PH,\qquad
 D_*^{\rm mov}=Z_*e^{2q_*}J_*+PH.
\end{align*}

\begin{lemma}[Cancellation between normalized residual and mismatch]
\label{sbsharp:weighted}
Suppose $\beta_B\le0.3$, $W_*<0.005$, $q_*<q_c$,
$3P/A_{0,-}<0.001$, and the strict stopping conditions of
Lemma~\ref{sbwide:relative} hold with $W_*,q_*$.
Then its conclusions and the core-holding lemma hold with these new
constants, and with the improved bounds
\[
 z\le Z_*,\qquad l\le l_*,\qquad
 \|c-(Av+a(0)_-)\|\le C_c^*,\qquad
 \|D-B(0)\|\le D_*^{\rm mov}.
\]
The core charges $M_*,V_*,J_D,J_B,J_p,\Gamma_C$ retain their original
form, using $W_*,q_*$. No neuron is discarded.
\end{lemma}
\begin{proof}
Stop additionally at $z+l=0.005$ and $|\log(u/A)|=q_c$.
The exact residual and mismatch equations imply, almost everywhere,
\begin{align*}
 l'&\le-gl+\beta_Bz+2P/u,\\
 z'&\le(-u^2-k\theta+z^2+|E_g|)z
                    +k\beta_Bl+P(1+\beta_B/u),
 \qquad |\theta|\le l.
\end{align*}
For $Z=z+\beta_Bl$ the $k\beta_Bl$ term cancels the matching part of
$-\beta_Bgl$ because $k-g=-k\theta+z^2-E_g$. Thus
\[
 Z'\le(\beta_B^2+l+z^2+|E_g|)Z+P(1+3\beta_B/u)
 \le(\beta_B^2+0.01)Z+P(1+6\beta_B/A_{0,-}).
\]
The last inequality uses the stopping bounds and
$|E_g|\le P(1+l+z)/u$. This proves $Z\le Z_*$.
Independently, dissipativity in the Cartesian mismatch equation gives
$(\|p\|)'\le\beta_Buz+2P$. Since
$\int_0^\tau A\,ds=\operatorname{arctanh}A(\tau)
-\operatorname{arctanh}A_0$, division by $u\ge Ae^{-q_c}$ gives $l_*$.
The function $(\operatorname{arctanh}A-\operatorname{arctanh}A_0)/A$
is increasing in $A\ge A_0$, justifying the uniform terminal bound.

In the exact radial identity, the forcing after $1-u^2$ is at most
$l_*+l_*^2+Z_*^2+2P/A_{0,-}$. The directional speed is at most
$l_*+\beta_BZ_*+2P/A_{0,-}$. The restoring logarithmic comparison
therefore gives $q_*$. These strict bounds close both added faces;
the original norm, $k$, $g$, and radial faces close as before.
The negative-part estimate for $c$ now uses $\|p\|\le l_*$, proving
$C_c^*$. Finally $\|D'\|\le u^2z+P$ and
$\int_0^\tau A^2\,ds=\frac12\log[(1-A_0^2)/(1-A(\tau)^2)]$
give $D_*^{\rm mov}$. All per-neuron core proofs use only $W_*,q_*$
and consequently remain valid with these stronger estimates.
\end{proof}

\subsection{The revised primitive witness}

\begin{theorem}[A smaller-width noisy coherent rebound and its timing]
\label{sbsharp:theorem}
Let $d=2$, $h\ge10^7$, $\varepsilon=\sqrt{0.04/h}$, and let $N=2000$
consist of independent samples with exactly half in each isotropic
Gaussian cluster, with
\[
 \mu_1=1,\qquad\mu_2=10^{-4},\qquad\sigma_1=\sigma_2=10^{-9}.
\]
Use independent isotropic balanced initialization. With probability at
least $0.97$, one fixed cohort has at least $0.24h$ neurons, input and
output angles at most $0.13$ radians, and longitudinal learned mass at
least $0.55$ throughout $[t_{\rm spec},t_b]$. With reference levels
$A^2=0.69,0.71,0.715,0.95$ defining $t_-,t_{\rm spec},t_a,t_b$,
\[
 0<0.00017<0.0055<t_-<t_{\rm spec}<t_a<t_b<7.6.
\]
The full mass is below $2$ through $t_b$ and $|f_2(e_2,t)|<0.011$.
On the entire deterministic open arc
\[
 \xi_\theta=(\sin\theta,\cos\theta),\qquad
 3.5\cdot10^{-5}<\theta<4\cdot10^{-4},
\]
whose angular width is $3.65\cdot10^{-4}$ radians,
\begin{align*}
 \dot E(\xi_\theta,t)&\le-1.5\cdot10^{-9}
                         &&(0\le t\le0.00017),\\
 E(\xi_\theta,0)-E(\xi_\theta,0.00017)&>2.04\cdot10^{-12},\\
 E(\xi_\theta,t_b)-E(\xi_\theta,t_a)&>7.2\cdot10^{-6},\\
 \dot E(\xi_\theta,t)&>3.56\cdot10^{-6}
                         &&(t_-\le t\le t_b)
\end{align*}
almost everywhere where derivatives are stated. On the inner arc
$3.5\cdot10^{-5}<\theta<4.5\cdot10^{-5}$ the rebound exceeds
$2\cdot10^{-5}$. On the upper subarc
$3\cdot10^{-4}<\theta<4\cdot10^{-4}$, of width $10^{-4}$,
\[
 \dot E\le-2.67\cdot10^{-8}\quad(0\le t\le0.0055),\qquad
 E(\xi,0)-E(\xi,0.0055)>2.15\cdot10^{-9}.
\]
It retains the full-arc rebound $7.2\cdot10^{-6}$.
Every global minimum
on $[0,t_b]$ belongs to $(0,t_-]$, hence occurs strictly before the
certified specialization time. The statement certifies an initial
transient and subsequent degradation of a learned dominant-family state;
it does not locate a reversal caused by specialization.
\end{theorem}

\begin{proof}
\emph{One event and one ledger.}
Use empirical windows
\[
 Q/h,K/s\in[0.248,0.252],\qquad
 \left|h^{-1}\sum t_i^+y_i^+-\frac1{4\pi}\right|\le0.002,
\]
$|h^{-1}\sum t_i y_i^+|\le0.001$, core count at least $0.24h$,
and core squared-first-coordinate mean at least
$1/8+1/(4\pi)-0.002$. Their scalar ranges are respectively
$1,1,1/2,1,1,1$. Applying Hoeffding separately gives total failure at most
\[
 6e^{-2h(0.002)^2}+2e^{-8h(0.002)^2}
             +2e^{-2h(0.001)^2}+2e^{-2h(0.01)^2}<4.13\cdot10^{-9}.
\]
On $Q\ge0.248h$, the event $b_0\le1.25\varepsilon$ has
variance-normalized threshold $5\sqrt{0.248}$. The same
Berry--Esseen bound as before and the elementary Gaussian tail bound
give failure less than $0.015575$ for $h\ge10^7$.
The eleven-probe plus two-axis Bernstein event costs $0.005$.
With $a_D=10^{-8}$, the sample event costs $2000e^{-50}$.
The total is below $0.020575<0.03$; no independence of these events is
used.

Set $H=3.8$, $\nu=0.05$, $A_{0,\pm}=\sqrt{0.04(0.25\pm0.002)}$.
The normalized terminal time is at most $3.773836<H$. Direct
substitution into Lemma~\ref{sbsharp:weighted} gives the following
outward-rounded upper bounds (the minimum-width case is worst):
\[
\begin{array}{c|c@{\quad}c|c}
 b_0&7.906\cdot10^{-5}&P&6.001\cdot10^{-7}\\
 \beta_B&0.200003&Z_*&1.316\cdot10^{-4}\\
 l_*&1.021\cdot10^{-4}&W_*&2.337\cdot10^{-4}\\
 q_*&5.337\cdot10^{-4}&C_c^*&0.002094\\
 D_*^{\rm mov}&1.990\cdot10^{-4}&M_*&0.002602\\
 V_*&0.010777&J_D&0.009034\\
 J_B&2.315\cdot10^{-5}&J_p&9.422\cdot10^{-4}\\
 \Gamma_C&0.016426&&
\end{array}
\]
All stopping conditions are strict. The mass-stop expression is below
$1.056$; the learned core mass exceeds $0.55196$, its angle is below
$0.12210$, and its normalized share exceeds $0.80387$.
The calculation uses the exact definitions, not successive substitution
of the displayed rounded values.

\emph{Endpoint on the inner arc.}
Put $K_+=0.252s$, $C_+=0.001s$ and
\[
 P_-=\sqrt{s}\frac{1/(4\pi)-0.002}{\sqrt{0.252}},\qquad
 P_+=\sqrt{s}\frac{1/(4\pi)+0.002}{\sqrt{0.248}}.
\]
The two coordinate errors are
\[
 e_1=C_c^*\sqrt{K_+}+\sqrt2PH<0.000213365,\qquad
 e_2=D_*^{\rm mov}\sqrt{K_+}+\beta_BPH<0.000020428.
\]
The exact uncharged secant lower bound is
\[
 \frac{P_-^2}{2}(A_b-A_a)(A_b+A_a-2A_{0,+})
       -C_+P_+(A_b-A_a)>9.9694\cdot10^{-5}.
\]
With $L_j^+=P_+(A_j-A_{0,-})+C_+$, its charge is
$(L_a^++L_b^+)e_1+e_1^2+2e_2+e_2^2<5.2251\cdot10^{-5}$.
The exact difference exceeds $4.74439\cdot10^{-5}$.
Throughout $0\le\theta\le4.5\cdot10^{-5}$, the residual coordinate
bounds are $0.030$ and $1.011$, with spatial derivatives at most
$3$ and $\beta_B\sqrt2+\theta$. Thus $|\partial_\theta E|<0.4$.
The inner-arc secant consequently exceeds
$4.74439\cdot10^{-5}-0.8(4.5\cdot10^{-5})>1.1\cdot10^{-5}$.
Also $K_++e_2<0.011$ bounds weak fitting throughout.

\emph{Initial decrease.}
Use Lemma~\ref{sbsharp:descent} with the radial-error rate $\lambda_0<0.495019$ above,
$c_0=\sin(3.5\cdot10^{-5})-2a_D$,
$L_0=s(1/(4\pi)-0.002)\cos(4.5\cdot10^{-5})-2a_Ds$,
and
\[
 d_0=\sqrt2s\left\{\sqrt{\frac{\log10400}{4h}}
        +\frac{4\log10400}{3h}\right\}+0.05(5\cdot10^{-7}).
\]
Its worst value is less than $2.729707\cdot10^{-5}$.
For the weak budget use
$c_w=\cos\beta-2a_D/\mu_2$,
$L_{w,0}=s(0.248)-s(\beta+2a_D/\mu_2)$ and radial error equal to the
net bound plus $0.05(2a_D/\mu_2)$. This budget stays positive.
The strong budget's zero belongs to
$(0.00018270,0.00018271)$, and at $T=0.00017$ its weighted rate exceeds
$1.5604\cdot10^{-9}$. The exact polynomial antiderivative gives a
drop greater than $2.04297\cdot10^{-12}$, proving the reported bounds.
All substitutions respect the two monotonicity conditions in the lemma.
To extend to the larger arc, anchor the initial residual at
$\alpha=3.5\cdot10^{-5}$. Its Lipschitz constant $0.05$ gives the
pointwise radial budget $d_0+0.05(\theta-\alpha)$, while use
$c(\theta)=\sin\theta-2a_D$ and
$L_0(\theta)=s(1/(4\pi)-0.002)\cos\theta-2a_Ds$.
For $\alpha\le\theta\le4\cdot10^{-4}$ and $t\le0.00017$, direct
partial differentiation gives
$\mathcal J_c>0.0060$, $|\mathcal J_d|<0.0063$, and
$|\mathcal J_L L_0'(\theta)|<2\cdot10^{-9}$.
Thus $\partial_\theta\mathcal J>0.0060\cos(4\cdot10^{-4})
-0.05(0.0063)-2\cdot10^{-9}>0.005$.
The weakest strong budget remains the left endpoint's, so the same
drop and rate apply to the larger arc without any new probability cost.
For the weak budget use the larger endpoint radial error if necessary;
it remains strictly positive.

For the longer initial transient on the upper subarc, keep the same
probability event and use the \emph{uniform} radial budget
$d_0+0.05(4\cdot10^{-4}-3.5\cdot10^{-5})$, cosine at
$4\cdot10^{-4}$ in $L_0$, and $c_0=\sin(3\cdot10^{-4})-2a_D$.
This is merely Lipschitz transport from the existing anchor.
The corresponding budget zero lies in $(0.00571064,0.00571066)$.
At $T=0.0055$, direct polynomial integration gives a drop greater than
$2.15495\cdot10^{-9}$ and a uniform derivative magnitude greater than
$2.6743\cdot10^{-8}$. The weak budget is still positive. Thus the longer
transient and the later rebound coexist on this same $10^{-4}$-wide arc.

\emph{Signed increase and the minimum's location.}
The following direct calculation is valid at every time with
$A^2\in[0.69,0.95]$, including times before the declared specialization.
Put $\bar\theta=4\cdot10^{-4}$, $A_-^\dagger=\sqrt{0.69}$,
$k_-=0.05-M_*$, and
\begin{align*}
 T_-&=A_-^\dagger P_- -3q_*\sqrt{K_+}-PH,\\
 T_+&=\sqrt{0.95}P_++3q_*\sqrt{K_+}+PH,\\
 F_-&=P_-(A_-^\dagger-A_{0,+})-C_+-e_1-3\bar\theta,\\
 V_1&=k_-\cos\bar\theta\,T_-
   -\bar\theta\{l_*+\sqrt2(l_*+\beta_BZ_*+P)\}-2\sqrt2P,\\
 V_2&=Z_*(T_++\bar\theta)+\bar\theta(T_++D_*^{\rm mov})
   +\bar\theta\beta_B(l_*+\beta_BZ_*+P)+P(\sqrt2+\beta_B).
\end{align*}
For $q_i=a_i\sin\theta+B_i\cos\theta$, differentiate the actual
all-neuron sum $(c,D)(q_i)_+$. The leading first-coordinate term is
$k\sum a_i^+(q_i)_+\ge k\cos\theta\sum a_i^+(B_i)_+$.
The other potentially negative longitudinal leading term satisfies
$k\sin\theta\sum c_i a_i^+\mathbf1_{q_i>0}\ge-\sin\theta\,l_*$,
because $c_i=a_i+p_i$. The remaining input terms cost at most
$\sin\theta\sqrt2(l_*+\beta_BZ_*+P)$, and the two output/input error
terms cost $2\sqrt2P$. Thus $f_1'\ge V_1$.
For the second coordinate, retain the leading signed term
$-R\sum a_i^+(q_i)_+$ and then bound its magnitude by
$Z_*(T_++\bar\theta)$. For the input-leading term, preserve its sign:
\[
 \sum D_i a_i^+\mathbf1_{q_i>0}
 \le\sum (B_i(0))_+a_i^++\|D-B(0)\|u
 \le T_++D_*^{\rm mov}.
\]
Its upper bound is consequently $\bar\theta(T_++D_*^{\rm mov})$.
The remaining terms cost
$\bar\theta\beta_B(l_*+\beta_BZ_*+P)$, and the two error terms cost
$P(\sqrt2+\beta_B)$. Hence $f_2'\le V_2$; an absolute bound is neither
claimed nor needed. Since $f_2-\cos\theta<0$, this one-sided bound has
the correct sign in the error derivative.
These formulas permit both training and probe gate changes.
We also have $f_1-\sin\theta\ge F_->0$ and
$|f_2-\cos\theta|<1.011$. Therefore, in physical time,
\[
 \dot E\ge\tfrac12(F_-V_1-1.011V_2)>3.5618\cdot10^{-6}.
\]
The exact formulas give the strict inequality. The physical endpoint
separation is
$\log\frac{(0.95)(0.285)}{(0.05)(0.715)}>2.0246$; integration gives a
full-arc rebound exceeding $7.2\cdot10^{-6}$. Substitution of
$\bar\theta=4.5\cdot10^{-5}$ gives a rate above $10^{-5}$ and an
inner-arc rebound above $2\cdot10^{-5}$. Absolute continuity gives strict increase on $[t_-,t_b]$.
Initial descent excludes a global minimum at zero; strict increase
excludes every time after $t_-$. Compactness supplies a minimum, proving
its stated location. The probes are more than $0.9$ from every training
sample on the same event.
\end{proof}

\subsection{What these improvements do and do not resolve}

The original theorem's descent interval improves by a factor $35000$,
and its certified drop by more than $149000$, on the same parameters.
Its common sector improves from $2\cdot10^{-18}$ to $10^{-5}$ radians.
The new theorem reduces width by $1000$, increases $\mu_2/\mu_1$ by $10$
and noise by $100$, while keeping the learned-mass and angular guarantees.
At its smallest width, $\sigma/\varepsilon=1.58114\cdot10^{-5}$.
These remain a very small noise, strongly unequal concept regime.
Choosing the upper subarc further raises the rigorous early drop to
$2.15\cdot10^{-9}$ while retaining rebound $7.2\cdot10^{-6}$;
these phases still have substantially different certified amplitudes.

The signed-increase calculation also applies on the old theorem's event:
use its own moment bounds by replacing $P_-,P_+$ with
$\sqrt{s}(0.077)/\sqrt{0.252}$ and
$\sqrt{s}(0.082)/\sqrt{0.248}$, and $C_+$ with $4\cdot10^{-6}$.
Its $b_0\le4\cdot10^{-6}$ and forcing are smaller than the new
comparison bounds. On its old arc, contained in
$0<\theta<2.50001\cdot10^{-5}$, the same formula gives a derivative
above $9\cdot10^{-6}$. Consequently every old-theorem
global minimum on $[0,t_b]$ also precedes its specialization time.
This is a \emph{construction-level timing limitation}, not evidence
against other noisy many-neuron reversal mechanisms.

For comparison, in the exactly isotropic continuum rank-one reference,
let $P=\sqrt{s}/(2\pi)$, $A_0=\sqrt{s}/2$, $r=\sin\theta$,
$k_r=\sqrt{1-r^2}$, and $x=2Ar/(k_r\sqrt{s})$.
Its evolving positive-first-coordinate contribution is exactly
\[
 F_1^+(A,r)=\frac{A\sqrt{s}k_r}{2\pi}
       \{1+x(\pi/2+\arctan x)\},\qquad
 F_2^+(A,r)=\frac{s k_r}{4\pi}(\pi/2+\arctan x).
\]
The full output is $(s/4)\xi+F^+(A,r)-F^+(A_0,r)$, so no inactive
background is omitted. Direct differentiation gives
$E_A(A_0,r)=-2\gamma Pr+O(r^2)$ and
$E_{AA}(A_0,0)=P^2$. The implicit-function theorem therefore gives a
local minimum with
$A-A_0=2\gamma r/P+O(r^2)$ and local improvement
$2\gamma^2r^2+O(r^3)$, for fixed $s>0$ as $r\downarrow0$.
This explains an early, small transient in that reference construction.
It is \emph{not} an upper bound on the finite-width noisy trajectory's
maximum improvement: the current comparison errors do not justify that
transfer. The actual improvement/rebound amplitude asymmetry beyond the
improved certificates remains unresolved.

The weak-fitting upper bound rules out meaningful weak learning within
this same window. Reaching such learning, or a minimum after a learned
specialization state, needs an additional argument and is not a modest
restatement of the present theorem.

\subsection{Numerical validation at the certified primitive point}

After the preceding analytic ledger closed, the exact smallest-width
point was instantiated with $h=10^7$, all $2000$ noisy training samples,
and random seed $20260918$. Its initialization and sample events passed
without resampling. The empirical gradients were evaluated using the
actual masks: neurons with a mask uniform across a cluster were summed
as matrices, and all boundary neurons were evaluated sample by sample.
An all-pairs preflight containing both types of boundary neuron agreed
in outputs and gradients to less than $4\cdot10^{-16}$.

Classical RK4 runs with maximum physical steps $0.02$ and $0.01$,
smaller initial steps, and exact snapshot stops gave the following
\emph{numerical} values:
\[
\begin{array}{c|r}
 t_{\rm spec}&5.489569515\\
 t_a&5.513978830\\
 t_b&7.538624447\\
 \text{core mass at }t_{\rm spec}&0.581020472\\
 \text{largest core angle there}&0.118179854\\
 \max S&0.979990746\\
 \max |f_2(e_2)|&0.009992690
\end{array}
\]
For $\theta=4\cdot10^{-5}$, the observed secant was
$1.05870369\cdot10^{-4}$. A sampled derivative crossing occurred between
$0.04617$ and $0.04717$, and the sampled maximum initial improvement was
$2.76713\cdot10^{-9}$.
For $\theta=3.5\cdot10^{-4}$, the observed secant was
$1.04807195\cdot10^{-4}$, the sampled crossing bracket was
$[0.38017,0.40017]$, and the sampled improvement was
$2.36472\cdot10^{-7}$. Both are early transients well before the
specialization declaration.

At $407$ common snapshots, halving the step changed the inner-probe
error by at most $5.08\cdot10^{-14}$ and core mass by at most
$9.18\cdot10^{-11}$. All $22$ instrumented bounds on mass, angles,
relative residual, mismatch, comparison errors, core holding, and
initial/late derivatives passed at recorded times. These are numerical
consistency checks, not rigorous integration-error bounds, continuous-time
certificates, minimum-location proofs, or proof premises. The wider-angle
probe was illustrated at step $0.02$; the refinement check concerns the
same training trajectory observed at the inner probe.

\end{document}
